\documentclass[12pt,fontset=fandol]{ctexart}
\makeatletter
\def\input@path{{../}{../../}}
\makeatother
\usepackage{researchnotes}
\title{数学分析与高等数学本科习题及解答}
\author{开放资源中文选编与训练扩展}
\date{2026年10月9日}
\begin{document}
\maketitle
\section*{编写说明与使用方法}
\addcontentsline{toc}{section}{编写说明与使用方法}
本书面向已经学习或正在学习本科数学分析、高等数学的读者，采用中文题面与中文解答。选题参考加拿大不列颠哥伦比亚大学（UBC）的 CLP 微积分开放教材与习题册。原作者为 Joel Feldman、Andrew Rechnitzer 和 Elyse Yeager。资料的原文、完整题库与许可链接见书末来源说明。检索日期为2026年10月9日。

本书是一份\keyterm{选译、改编与补充训练集}，不是某一本教材的逐题译本或官方答案。标注“选译”的题目来自已核对的原题；标注“改编训练”的题目按所列开放教材章节的主题和方法重新组织，并包含参数变式；标注“补充证明”的题目是为数学分析课程补写的常见论证。不同参数变式分别编号，主要用于练习同一方法的代数细节，不应把这些变式理解为彼此独立的新题型。中文解答重新组织了推导与条件说明，涉及编写者的翻译、改编和扩展，原作者并未审定这些改动。

全书共有45章、1076道编号习题，含同类参数变式；每一道均有中文参考解答。每章先给出习题，再集中给出对应解答。题目后的页码链接可跳至解答，解答标题中的题号可跳回题目。建议首次学习时按章节顺序完成各类的第一题；掌握方法后，再用同类后续题检验熟练度。基础计算路线可重点阅读极限、求导、积分、级数与多元微积分各章；数学专业读者还应完成连续性、收敛性和补充证明章节。证明题要交代假设与量词，计算题要保留定义域、积分限、单位和方向等信息。阅读答案后，应合上答案重新写出主要步骤。

全书默认实变量与实值函数，未另行说明的有限区间 $[a,b]$ 满足 $a<b$，$\log$ 表示自然对数，三角函数使用弧度制；$C$ 为任意积分常数，不同连通区间上的常数可以独立。自然数下标从1开始，另有说明者除外；幂级数按题面注明的起点求和。$O(h^k)$ 表示绝对值被常数倍 $|h|^k$ 控制，$o(h^k)$ 表示除以 $h^k$ 后趋零。曲线积分的方向和曲面积分的法向以题面为准。

内容覆盖本科微积分的主要计算与分析方法；并不以完备性公理、一般度量空间或测度论为完整理论主线。每道题均附参考解答，但未经独立教师逐题审稿，宜作为学习与核对材料。复核重点是题目条件、关键推导和最终结果；不能将文件成功编译等同于全部数学命题已经由计算机证明。

本书中据 CLP 改编的内容与新增中文编排按\href{https://creativecommons.org/licenses/by-nc-sa/4.0/}{知识共享署名—非商业性使用—相同方式共享4.0国际许可（CC BY-NC-SA 4.0）}提供。再分发时应保留原作者署名、来源、许可链接及修改说明，不得暗示原作者认可本书。
\tableofcontents
\clearpage
\section{数列极限与收敛的基本训练}\label{sec:ch1}
数列的极限（limit of a sequence）描述尾部的统一行为。计算极限前先识别主导项；递推数列则必须先证明有界性与单调性，不能只解一个形式上的极限方程。以下题目中 $n$ 均取正整数。

\textbf{资源对应：}CLP-2 \S~3.1。本章共 24 道编号习题。
\subsection{习题}
\begin{example}[同阶多项式的比]\label{ex:c1q1}
求数列 $a_n=\dfrac{2n^2+1n+1}{3n^2+n+2}$ 的极限，并给出误差为 $O(1/n)$ 的说明。
\par\textit{改编训练。}\par\hyperref[ans:c1q1]{解答见第~\pageref{ans:c1q1} 页。}
\end{example}
\begin{example}[同阶多项式的比]\label{ex:c1q2}
求数列 $a_n=\dfrac{3n^2+(-2)n+1}{4n^2+n+2}$ 的极限，并给出误差为 $O(1/n)$ 的说明。
\par\textit{改编训练。}\par\hyperref[ans:c1q2]{解答见第~\pageref{ans:c1q2} 页。}
\end{example}
\begin{example}[同阶多项式的比]\label{ex:c1q3}
求数列 $a_n=\dfrac{1n^2+5n+1}{2n^2+n+2}$ 的极限，并给出误差为 $O(1/n)$ 的说明。
\par\textit{改编训练。}\par\hyperref[ans:c1q3]{解答见第~\pageref{ans:c1q3} 页。}
\end{example}
\begin{example}[同阶多项式的比]\label{ex:c1q4}
求数列 $a_n=\dfrac{4n^2+3n+1}{7n^2+n+2}$ 的极限，并给出误差为 $O(1/n)$ 的说明。
\par\textit{改编训练。}\par\hyperref[ans:c1q4]{解答见第~\pageref{ans:c1q4} 页。}
\end{example}
\begin{example}[同阶多项式的比]\label{ex:c1q5}
求数列 $a_n=\dfrac{5n^2+(-1)n+1}{6n^2+n+2}$ 的极限，并给出误差为 $O(1/n)$ 的说明。
\par\textit{改编训练。}\par\hyperref[ans:c1q5]{解答见第~\pageref{ans:c1q5} 页。}
\end{example}
\begin{example}[同阶多项式的比]\label{ex:c1q6}
求数列 $a_n=\dfrac{7n^2+2n+1}{3n^2+n+2}$ 的极限，并给出误差为 $O(1/n)$ 的说明。
\par\textit{改编训练。}\par\hyperref[ans:c1q6]{解答见第~\pageref{ans:c1q6} 页。}
\end{example}
\begin{example}[根式差的有理化]\label{ex:c1q7}
求 $a_n=\sqrt{n^2+1n+2}-n$ 的极限，并说明它与 $\sqrt{n^2+1n+2}/n$ 的极限有何不同。
\par\textit{改编训练。}\par\hyperref[ans:c1q7]{解答见第~\pageref{ans:c1q7} 页。}
\end{example}
\begin{example}[根式差的有理化]\label{ex:c1q8}
求 $a_n=\sqrt{n^2+2n+5}-n$ 的极限，并说明它与 $\sqrt{n^2+2n+5}/n$ 的极限有何不同。
\par\textit{改编训练。}\par\hyperref[ans:c1q8]{解答见第~\pageref{ans:c1q8} 页。}
\end{example}
\begin{example}[根式差的有理化]\label{ex:c1q9}
求 $a_n=\sqrt{n^2+3n+1}-n$ 的极限，并说明它与 $\sqrt{n^2+3n+1}/n$ 的极限有何不同。
\par\textit{改编训练。}\par\hyperref[ans:c1q9]{解答见第~\pageref{ans:c1q9} 页。}
\end{example}
\begin{example}[根式差的有理化]\label{ex:c1q10}
求 $a_n=\sqrt{n^2+4n+7}-n$ 的极限，并说明它与 $\sqrt{n^2+4n+7}/n$ 的极限有何不同。
\par\textit{改编训练。}\par\hyperref[ans:c1q10]{解答见第~\pageref{ans:c1q10} 页。}
\end{example}
\begin{example}[根式差的有理化]\label{ex:c1q11}
求 $a_n=\sqrt{n^2+5n+3}-n$ 的极限，并说明它与 $\sqrt{n^2+5n+3}/n$ 的极限有何不同。
\par\textit{改编训练。}\par\hyperref[ans:c1q11]{解答见第~\pageref{ans:c1q11} 页。}
\end{example}
\begin{example}[根式差的有理化]\label{ex:c1q12}
求 $a_n=\sqrt{n^2+6n+9}-n$ 的极限，并说明它与 $\sqrt{n^2+6n+9}/n$ 的极限有何不同。
\par\textit{改编训练。}\par\hyperref[ans:c1q12]{解答见第~\pageref{ans:c1q12} 页。}
\end{example}
\begin{example}[指数型数列]\label{ex:c1q13}
求 $a_n=(1+1/n)^{1n}$ 的极限，并判断 $\log a_n$ 是否收敛。
\par\textit{改编训练。}\par\hyperref[ans:c1q13]{解答见第~\pageref{ans:c1q13} 页。}
\end{example}
\begin{example}[指数型数列]\label{ex:c1q14}
求 $a_n=(1+2/n)^{1n}$ 的极限，并判断 $\log a_n$ 是否收敛。
\par\textit{改编训练。}\par\hyperref[ans:c1q14]{解答见第~\pageref{ans:c1q14} 页。}
\end{example}
\begin{example}[指数型数列]\label{ex:c1q15}
求 $a_n=(1+1/n)^{3n}$ 的极限，并判断 $\log a_n$ 是否收敛。
\par\textit{改编训练。}\par\hyperref[ans:c1q15]{解答见第~\pageref{ans:c1q15} 页。}
\end{example}
\begin{example}[指数型数列]\label{ex:c1q16}
求 $a_n=(1+3/n)^{2n}$ 的极限，并判断 $\log a_n$ 是否收敛。
\par\textit{改编训练。}\par\hyperref[ans:c1q16]{解答见第~\pageref{ans:c1q16} 页。}
\end{example}
\begin{example}[指数型数列]\label{ex:c1q17}
求 $a_n=(1+2/n)^{3n}$ 的极限，并判断 $\log a_n$ 是否收敛。
\par\textit{改编训练。}\par\hyperref[ans:c1q17]{解答见第~\pageref{ans:c1q17} 页。}
\end{example}
\begin{example}[指数型数列]\label{ex:c1q18}
求 $a_n=(1+4/n)^{1n}$ 的极限，并判断 $\log a_n$ 是否收敛。
\par\textit{改编训练。}\par\hyperref[ans:c1q18]{解答见第~\pageref{ans:c1q18} 页。}
\end{example}
\begin{example}[根式递推与单调有界]\label{ex:c1q19}
设 $a_1=1$，$a_{n+1}=\sqrt{2+a_n}$。证明数列收敛并求极限。
\par\textit{改编训练。}\par\hyperref[ans:c1q19]{解答见第~\pageref{ans:c1q19} 页。}
\end{example}
\begin{example}[根式递推与单调有界]\label{ex:c1q20}
设 $a_1=1$，$a_{n+1}=\sqrt{6+a_n}$。证明数列收敛并求极限。
\par\textit{改编训练。}\par\hyperref[ans:c1q20]{解答见第~\pageref{ans:c1q20} 页。}
\end{example}
\begin{example}[根式递推与单调有界]\label{ex:c1q21}
设 $a_1=1$，$a_{n+1}=\sqrt{12+a_n}$。证明数列收敛并求极限。
\par\textit{改编训练。}\par\hyperref[ans:c1q21]{解答见第~\pageref{ans:c1q21} 页。}
\end{example}
\begin{example}[根式递推与单调有界]\label{ex:c1q22}
设 $a_1=1$，$a_{n+1}=\sqrt{20+a_n}$。证明数列收敛并求极限。
\par\textit{改编训练。}\par\hyperref[ans:c1q22]{解答见第~\pageref{ans:c1q22} 页。}
\end{example}
\begin{example}[根式递推与单调有界]\label{ex:c1q23}
设 $a_1=1$，$a_{n+1}=\sqrt{30+a_n}$。证明数列收敛并求极限。
\par\textit{改编训练。}\par\hyperref[ans:c1q23]{解答见第~\pageref{ans:c1q23} 页。}
\end{example}
\begin{example}[根式递推与单调有界]\label{ex:c1q24}
设 $a_1=1$，$a_{n+1}=\sqrt{42+a_n}$。证明数列收敛并求极限。
\par\textit{改编训练。}\par\hyperref[ans:c1q24]{解答见第~\pageref{ans:c1q24} 页。}
\end{example}
\subsection{解答}
\begin{solve}[习题~\ref{ex:c1q1}：同阶多项式的比]\phantomsection\label{ans:c1q1}
分子、分母同除以 $n^2$，得到
\[
a_n=\frac{2+1/n+1/n^2}{3+1/n+2/n^2}\longrightarrow \frac{2}{3}.
\]
分母趋向正数 $3$，因而可以使用商的极限法则。若 $L=\frac{2}{3}$，则精确地有
\[
a_n-L=\frac{(1)n+(-1)}{3(3n^2+n+2)}.
\]
因此 $|a_n-L|\le\dfrac{2}{9n}$。这同时给出收敛速度；这里的 $O(1/n)$ 表示误差绝对值被某个固定常数乘 $1/n$ 控制。
\end{solve}

\begin{solve}[习题~\ref{ex:c1q2}：同阶多项式的比]\phantomsection\label{ans:c1q2}
分子、分母同除以 $n^2$，得到
\[
a_n=\frac{3+(-2)/n+1/n^2}{4+1/n+2/n^2}\longrightarrow \frac{3}{4}.
\]
分母趋向正数 $4$，因而可以使用商的极限法则。若 $L=\frac{3}{4}$，则精确地有
\[
a_n-L=\frac{(-11)n+(-2)}{4(4n^2+n+2)}.
\]
因此 $|a_n-L|\le\dfrac{13}{16n}$。这同时给出收敛速度；这里的 $O(1/n)$ 表示误差绝对值被某个固定常数乘 $1/n$ 控制。
\end{solve}

\begin{solve}[习题~\ref{ex:c1q3}：同阶多项式的比]\phantomsection\label{ans:c1q3}
分子、分母同除以 $n^2$，得到
\[
a_n=\frac{1+5/n+1/n^2}{2+1/n+2/n^2}\longrightarrow \frac{1}{2}.
\]
分母趋向正数 $2$，因而可以使用商的极限法则。若 $L=\frac{1}{2}$，则精确地有
\[
a_n-L=\frac{(9)n+(0)}{2(2n^2+n+2)}.
\]
因此 $|a_n-L|\le\dfrac{9}{4n}$。这同时给出收敛速度；这里的 $O(1/n)$ 表示误差绝对值被某个固定常数乘 $1/n$ 控制。
\end{solve}

\begin{solve}[习题~\ref{ex:c1q4}：同阶多项式的比]\phantomsection\label{ans:c1q4}
分子、分母同除以 $n^2$，得到
\[
a_n=\frac{4+3/n+1/n^2}{7+1/n+2/n^2}\longrightarrow \frac{4}{7}.
\]
分母趋向正数 $7$，因而可以使用商的极限法则。若 $L=\frac{4}{7}$，则精确地有
\[
a_n-L=\frac{(17)n+(-1)}{7(7n^2+n+2)}.
\]
因此 $|a_n-L|\le\dfrac{18}{49n}$。这同时给出收敛速度；这里的 $O(1/n)$ 表示误差绝对值被某个固定常数乘 $1/n$ 控制。
\end{solve}

\begin{solve}[习题~\ref{ex:c1q5}：同阶多项式的比]\phantomsection\label{ans:c1q5}
分子、分母同除以 $n^2$，得到
\[
a_n=\frac{5+(-1)/n+1/n^2}{6+1/n+2/n^2}\longrightarrow \frac{5}{6}.
\]
分母趋向正数 $6$，因而可以使用商的极限法则。若 $L=\frac{5}{6}$，则精确地有
\[
a_n-L=\frac{(-11)n+(-4)}{6(6n^2+n+2)}.
\]
因此 $|a_n-L|\le\dfrac{15}{36n}$。这同时给出收敛速度；这里的 $O(1/n)$ 表示误差绝对值被某个固定常数乘 $1/n$ 控制。
\end{solve}

\begin{solve}[习题~\ref{ex:c1q6}：同阶多项式的比]\phantomsection\label{ans:c1q6}
分子、分母同除以 $n^2$，得到
\[
a_n=\frac{7+2/n+1/n^2}{3+1/n+2/n^2}\longrightarrow \frac{7}{3}.
\]
分母趋向正数 $3$，因而可以使用商的极限法则。若 $L=\frac{7}{3}$，则精确地有
\[
a_n-L=\frac{(-1)n+(-11)}{3(3n^2+n+2)}.
\]
因此 $|a_n-L|\le\dfrac{12}{9n}$。这同时给出收敛速度；这里的 $O(1/n)$ 表示误差绝对值被某个固定常数乘 $1/n$ 控制。
\end{solve}

\begin{solve}[习题~\ref{ex:c1q7}：根式差的有理化]\phantomsection\label{ans:c1q7}
原式是两个无穷大量之差，不能分别取极限后相减。乘以共轭式可把差转化为商：
\[
a_n=\frac{1n+2}{\sqrt{n^2+1n+2}+n}=\frac{1+2/n}{\sqrt{1+1/n+2/n^2}+1}\longrightarrow \frac{1}{2}.
\]
另一方面，
\[
\frac{\sqrt{n^2+1n+2}}n=\sqrt{1+1/n+2/n^2}\longrightarrow1.
\]
后一结论只说明两个量的比趋于 $1$，不能单独确定它们的差的极限。本题中被根号项为正，有理化所用分母也始终为正。
\end{solve}

\begin{solve}[习题~\ref{ex:c1q8}：根式差的有理化]\phantomsection\label{ans:c1q8}
原式是两个无穷大量之差，不能分别取极限后相减。乘以共轭式可把差转化为商：
\[
a_n=\frac{2n+5}{\sqrt{n^2+2n+5}+n}=\frac{2+5/n}{\sqrt{1+2/n+5/n^2}+1}\longrightarrow 1.
\]
另一方面，
\[
\frac{\sqrt{n^2+2n+5}}n=\sqrt{1+2/n+5/n^2}\longrightarrow1.
\]
后一结论只说明两个量的比趋于 $1$，不能单独确定它们的差的极限。本题中被根号项为正，有理化所用分母也始终为正。
\end{solve}

\begin{solve}[习题~\ref{ex:c1q9}：根式差的有理化]\phantomsection\label{ans:c1q9}
原式是两个无穷大量之差，不能分别取极限后相减。乘以共轭式可把差转化为商：
\[
a_n=\frac{3n+1}{\sqrt{n^2+3n+1}+n}=\frac{3+1/n}{\sqrt{1+3/n+1/n^2}+1}\longrightarrow \frac{3}{2}.
\]
另一方面，
\[
\frac{\sqrt{n^2+3n+1}}n=\sqrt{1+3/n+1/n^2}\longrightarrow1.
\]
后一结论只说明两个量的比趋于 $1$，不能单独确定它们的差的极限。本题中被根号项为正，有理化所用分母也始终为正。
\end{solve}

\begin{solve}[习题~\ref{ex:c1q10}：根式差的有理化]\phantomsection\label{ans:c1q10}
原式是两个无穷大量之差，不能分别取极限后相减。乘以共轭式可把差转化为商：
\[
a_n=\frac{4n+7}{\sqrt{n^2+4n+7}+n}=\frac{4+7/n}{\sqrt{1+4/n+7/n^2}+1}\longrightarrow 2.
\]
另一方面，
\[
\frac{\sqrt{n^2+4n+7}}n=\sqrt{1+4/n+7/n^2}\longrightarrow1.
\]
后一结论只说明两个量的比趋于 $1$，不能单独确定它们的差的极限。本题中被根号项为正，有理化所用分母也始终为正。
\end{solve}

\begin{solve}[习题~\ref{ex:c1q11}：根式差的有理化]\phantomsection\label{ans:c1q11}
原式是两个无穷大量之差，不能分别取极限后相减。乘以共轭式可把差转化为商：
\[
a_n=\frac{5n+3}{\sqrt{n^2+5n+3}+n}=\frac{5+3/n}{\sqrt{1+5/n+3/n^2}+1}\longrightarrow \frac{5}{2}.
\]
另一方面，
\[
\frac{\sqrt{n^2+5n+3}}n=\sqrt{1+5/n+3/n^2}\longrightarrow1.
\]
后一结论只说明两个量的比趋于 $1$，不能单独确定它们的差的极限。本题中被根号项为正，有理化所用分母也始终为正。
\end{solve}

\begin{solve}[习题~\ref{ex:c1q12}：根式差的有理化]\phantomsection\label{ans:c1q12}
原式是两个无穷大量之差，不能分别取极限后相减。乘以共轭式可把差转化为商：
\[
a_n=\frac{6n+9}{\sqrt{n^2+6n+9}+n}=\frac{6+9/n}{\sqrt{1+6/n+9/n^2}+1}\longrightarrow 3.
\]
另一方面，
\[
\frac{\sqrt{n^2+6n+9}}n=\sqrt{1+6/n+9/n^2}\longrightarrow1.
\]
后一结论只说明两个量的比趋于 $1$，不能单独确定它们的差的极限。本题中被根号项为正，有理化所用分母也始终为正。
\end{solve}

\begin{solve}[习题~\ref{ex:c1q13}：指数型数列]\phantomsection\label{ans:c1q13}
各项均为正，先取自然对数。令 $t=1/n$，则 $t\to0^+$，由 $\lim_{t\to0}\log(1+t)/t=1$ 得
\[
\log a_n=1n\log(1+1/n)=1\frac{\log(1+1/n)}{1/n}\longrightarrow 1.
\]
指数函数连续，所以 $a_n\to e^{1}$。还可以用 $t-t^2/2\le\log(1+t)\le t$（$t\ge0$）直接夹逼：
\[
1-\frac{\frac{1}{2}}n\le\log a_n\le 1.
\]
该不等式来自对 $1/(1+t)$ 的积分估计，避免把未定式 $1^\infty$ 当成通常的乘方。
\end{solve}

\begin{solve}[习题~\ref{ex:c1q14}：指数型数列]\phantomsection\label{ans:c1q14}
各项均为正，先取自然对数。令 $t=2/n$，则 $t\to0^+$，由 $\lim_{t\to0}\log(1+t)/t=1$ 得
\[
\log a_n=1n\log(1+2/n)=2\frac{\log(1+2/n)}{2/n}\longrightarrow 2.
\]
指数函数连续，所以 $a_n\to e^{2}$。还可以用 $t-t^2/2\le\log(1+t)\le t$（$t\ge0$）直接夹逼：
\[
2-\frac{2}n\le\log a_n\le 2.
\]
该不等式来自对 $1/(1+t)$ 的积分估计，避免把未定式 $1^\infty$ 当成通常的乘方。
\end{solve}

\begin{solve}[习题~\ref{ex:c1q15}：指数型数列]\phantomsection\label{ans:c1q15}
各项均为正，先取自然对数。令 $t=1/n$，则 $t\to0^+$，由 $\lim_{t\to0}\log(1+t)/t=1$ 得
\[
\log a_n=3n\log(1+1/n)=3\frac{\log(1+1/n)}{1/n}\longrightarrow 3.
\]
指数函数连续，所以 $a_n\to e^{3}$。还可以用 $t-t^2/2\le\log(1+t)\le t$（$t\ge0$）直接夹逼：
\[
3-\frac{\frac{3}{2}}n\le\log a_n\le 3.
\]
该不等式来自对 $1/(1+t)$ 的积分估计，避免把未定式 $1^\infty$ 当成通常的乘方。
\end{solve}

\begin{solve}[习题~\ref{ex:c1q16}：指数型数列]\phantomsection\label{ans:c1q16}
各项均为正，先取自然对数。令 $t=3/n$，则 $t\to0^+$，由 $\lim_{t\to0}\log(1+t)/t=1$ 得
\[
\log a_n=2n\log(1+3/n)=6\frac{\log(1+3/n)}{3/n}\longrightarrow 6.
\]
指数函数连续，所以 $a_n\to e^{6}$。还可以用 $t-t^2/2\le\log(1+t)\le t$（$t\ge0$）直接夹逼：
\[
6-\frac{9}n\le\log a_n\le 6.
\]
该不等式来自对 $1/(1+t)$ 的积分估计，避免把未定式 $1^\infty$ 当成通常的乘方。
\end{solve}

\begin{solve}[习题~\ref{ex:c1q17}：指数型数列]\phantomsection\label{ans:c1q17}
各项均为正，先取自然对数。令 $t=2/n$，则 $t\to0^+$，由 $\lim_{t\to0}\log(1+t)/t=1$ 得
\[
\log a_n=3n\log(1+2/n)=6\frac{\log(1+2/n)}{2/n}\longrightarrow 6.
\]
指数函数连续，所以 $a_n\to e^{6}$。还可以用 $t-t^2/2\le\log(1+t)\le t$（$t\ge0$）直接夹逼：
\[
6-\frac{6}n\le\log a_n\le 6.
\]
该不等式来自对 $1/(1+t)$ 的积分估计，避免把未定式 $1^\infty$ 当成通常的乘方。
\end{solve}

\begin{solve}[习题~\ref{ex:c1q18}：指数型数列]\phantomsection\label{ans:c1q18}
各项均为正，先取自然对数。令 $t=4/n$，则 $t\to0^+$，由 $\lim_{t\to0}\log(1+t)/t=1$ 得
\[
\log a_n=1n\log(1+4/n)=4\frac{\log(1+4/n)}{4/n}\longrightarrow 4.
\]
指数函数连续，所以 $a_n\to e^{4}$。还可以用 $t-t^2/2\le\log(1+t)\le t$（$t\ge0$）直接夹逼：
\[
4-\frac{8}n\le\log a_n\le 4.
\]
该不等式来自对 $1/(1+t)$ 的积分估计，避免把未定式 $1^\infty$ 当成通常的乘方。
\end{solve}

\begin{solve}[习题~\ref{ex:c1q19}：根式递推与单调有界]\phantomsection\label{ans:c1q19}
先用归纳法证明 $1\le a_n<2$：初值满足要求；若 $a_n<2$，则 $a_{n+1}<\sqrt{2+2}=2$，下界也被保持。比较相邻项，有
\[
a_{n+1}-a_n=\frac{2+a_n-a_n^2}{\sqrt{2+a_n}+a_n}=\frac{(2-a_n)(a_n+1)}{\sqrt{2+a_n}+a_n}>0.
\]
故数列严格递增且有上界，单调有界定理保证它收敛。设极限为 $\ell$，对递推式取极限，
\[
\ell^2-\ell-2=0,\qquad(\ell-2)(\ell+1)=0.
\]
因 $\ell\ge1$，排除负根，得到 $\ell=2$。先证明收敛再解方程是本题不可省略的顺序。
\end{solve}

\begin{solve}[习题~\ref{ex:c1q20}：根式递推与单调有界]\phantomsection\label{ans:c1q20}
先用归纳法证明 $1\le a_n<3$：初值满足要求；若 $a_n<3$，则 $a_{n+1}<\sqrt{6+3}=3$，下界也被保持。比较相邻项，有
\[
a_{n+1}-a_n=\frac{6+a_n-a_n^2}{\sqrt{6+a_n}+a_n}=\frac{(3-a_n)(a_n+2)}{\sqrt{6+a_n}+a_n}>0.
\]
故数列严格递增且有上界，单调有界定理保证它收敛。设极限为 $\ell$，对递推式取极限，
\[
\ell^2-\ell-6=0,\qquad(\ell-3)(\ell+2)=0.
\]
因 $\ell\ge1$，排除负根，得到 $\ell=3$。先证明收敛再解方程是本题不可省略的顺序。
\end{solve}

\begin{solve}[习题~\ref{ex:c1q21}：根式递推与单调有界]\phantomsection\label{ans:c1q21}
先用归纳法证明 $1\le a_n<4$：初值满足要求；若 $a_n<4$，则 $a_{n+1}<\sqrt{12+4}=4$，下界也被保持。比较相邻项，有
\[
a_{n+1}-a_n=\frac{12+a_n-a_n^2}{\sqrt{12+a_n}+a_n}=\frac{(4-a_n)(a_n+3)}{\sqrt{12+a_n}+a_n}>0.
\]
故数列严格递增且有上界，单调有界定理保证它收敛。设极限为 $\ell$，对递推式取极限，
\[
\ell^2-\ell-12=0,\qquad(\ell-4)(\ell+3)=0.
\]
因 $\ell\ge1$，排除负根，得到 $\ell=4$。先证明收敛再解方程是本题不可省略的顺序。
\end{solve}

\begin{solve}[习题~\ref{ex:c1q22}：根式递推与单调有界]\phantomsection\label{ans:c1q22}
先用归纳法证明 $1\le a_n<5$：初值满足要求；若 $a_n<5$，则 $a_{n+1}<\sqrt{20+5}=5$，下界也被保持。比较相邻项，有
\[
a_{n+1}-a_n=\frac{20+a_n-a_n^2}{\sqrt{20+a_n}+a_n}=\frac{(5-a_n)(a_n+4)}{\sqrt{20+a_n}+a_n}>0.
\]
故数列严格递增且有上界，单调有界定理保证它收敛。设极限为 $\ell$，对递推式取极限，
\[
\ell^2-\ell-20=0,\qquad(\ell-5)(\ell+4)=0.
\]
因 $\ell\ge1$，排除负根，得到 $\ell=5$。先证明收敛再解方程是本题不可省略的顺序。
\end{solve}

\begin{solve}[习题~\ref{ex:c1q23}：根式递推与单调有界]\phantomsection\label{ans:c1q23}
先用归纳法证明 $1\le a_n<6$：初值满足要求；若 $a_n<6$，则 $a_{n+1}<\sqrt{30+6}=6$，下界也被保持。比较相邻项，有
\[
a_{n+1}-a_n=\frac{30+a_n-a_n^2}{\sqrt{30+a_n}+a_n}=\frac{(6-a_n)(a_n+5)}{\sqrt{30+a_n}+a_n}>0.
\]
故数列严格递增且有上界，单调有界定理保证它收敛。设极限为 $\ell$，对递推式取极限，
\[
\ell^2-\ell-30=0,\qquad(\ell-6)(\ell+5)=0.
\]
因 $\ell\ge1$，排除负根，得到 $\ell=6$。先证明收敛再解方程是本题不可省略的顺序。
\end{solve}

\begin{solve}[习题~\ref{ex:c1q24}：根式递推与单调有界]\phantomsection\label{ans:c1q24}
先用归纳法证明 $1\le a_n<7$：初值满足要求；若 $a_n<7$，则 $a_{n+1}<\sqrt{42+7}=7$，下界也被保持。比较相邻项，有
\[
a_{n+1}-a_n=\frac{42+a_n-a_n^2}{\sqrt{42+a_n}+a_n}=\frac{(7-a_n)(a_n+6)}{\sqrt{42+a_n}+a_n}>0.
\]
故数列严格递增且有上界，单调有界定理保证它收敛。设极限为 $\ell$，对递推式取极限，
\[
\ell^2-\ell-42=0,\qquad(\ell-7)(\ell+6)=0.
\]
因 $\ell\ge1$，排除负根，得到 $\ell=7$。先证明收敛再解方程是本题不可省略的顺序。
\end{solve}

\clearpage
\section{实数完备性与数列证明专题}\label{sec:ch2}
这一章从定义与实数完备性出发，补上计算练习背后的证明能力。证明中允许使用实数的上确界公理、自然数的 Archimedes 性质及此前已建立的基本极限定理；若换用其他完备性等价形式，应避免循环论证。

\textbf{资源对应：}编者补充证明；与 CLP-2 \S~3.1 数列主题衔接。本章共 20 道编号习题。
\subsection{习题}
\begin{example}[上确界的逼近性质]\label{ex:c2q1}
设非空集合 $A\subset\mathbb R$ 有上界，$s=\sup A$。证明：对每个 $\varepsilon>0$，存在 $a\in A$ 使 $s-\varepsilon<a\le s$。并说明逆命题所需的条件。
\par\textit{补充证明。来源：编者补充：本科数学分析标准命题与反例。}\par\hyperref[ans:c2q1]{解答见第~\pageref{ans:c2q1} 页。}
\end{example}
\begin{example}[从上确界构造逼近数列]\label{ex:c2q2}
设 $s=\sup A$ 为有限实数。构造 $A$ 中的数列收敛到 $s$，并判断能否要求该数列严格递增。
\par\textit{补充证明。来源：编者补充：本科数学分析标准命题与反例。}\par\hyperref[ans:c2q2]{解答见第~\pageref{ans:c2q2} 页。}
\end{example}
\begin{example}[单调有界定理的证明]\label{ex:c2q3}
从上确界公理证明：单调不减且有上界的实数列收敛。
\par\textit{补充证明。来源：编者补充：本科数学分析标准命题与反例。}\par\hyperref[ans:c2q3]{解答见第~\pageref{ans:c2q3} 页。}
\end{example}
\begin{example}[收敛数列有界]\label{ex:c2q4}
证明每个收敛的实数列有界，并给出有界但不收敛的例子。
\par\textit{补充证明。来源：编者补充：本科数学分析标准命题与反例。}\par\hyperref[ans:c2q4]{解答见第~\pageref{ans:c2q4} 页。}
\end{example}
\begin{example}[极限的唯一性]\label{ex:c2q5}
直接由定义证明实数列的极限唯一。
\par\textit{补充证明。来源：编者补充：本科数学分析标准命题与反例。}\par\hyperref[ans:c2q5]{解答见第~\pageref{ans:c2q5} 页。}
\end{example}
\begin{example}[保号性与严格不等式]\label{ex:c2q6}
若 $a_n\to a>0$，证明从某项起 $a_n>0$。若每项 $a_n>0$，能否推出 $a>0$？
\par\textit{补充证明。来源：编者补充：本科数学分析标准命题与反例。}\par\hyperref[ans:c2q6]{解答见第~\pageref{ans:c2q6} 页。}
\end{example}
\begin{example}[收敛子列与原列]\label{ex:c2q7}
证明收敛数列的每个子列收敛到同一个极限。一个子列收敛是否足以推出原列收敛？
\par\textit{补充证明。来源：编者补充：本科数学分析标准命题与反例。}\par\hyperref[ans:c2q7]{解答见第~\pageref{ans:c2q7} 页。}
\end{example}
\begin{example}[奇偶子列判别]\label{ex:c2q8}
若 $a_{2n}\to L$ 且 $a_{2n-1}\to L$，证明 $a_n\to L$。若两个极限不同呢？
\par\textit{补充证明。来源：编者补充：本科数学分析标准命题与反例。}\par\hyperref[ans:c2q8]{解答见第~\pageref{ans:c2q8} 页。}
\end{example}
\begin{example}[Cauchy 数列必有界]\label{ex:c2q9}
设实数列满足 Cauchy 条件：对每个 $\varepsilon>0$，存在 $N$ 使 $m,n\ge N$ 时 $|a_m-a_n|<\varepsilon$。证明它有界。
\par\textit{补充证明。来源：编者补充：本科数学分析标准命题与反例。}\par\hyperref[ans:c2q9]{解答见第~\pageref{ans:c2q9} 页。}
\end{example}
\begin{example}[Cauchy 判别的充分性]\label{ex:c2q10}
已知 Bolzano--Weierstrass 定理，证明实 Cauchy 数列收敛。
\par\textit{补充证明。来源：编者补充：本科数学分析标准命题与反例。}\par\hyperref[ans:c2q10]{解答见第~\pageref{ans:c2q10} 页。}
\end{example}
\begin{example}[嵌套区间定理]\label{ex:c2q11}
设闭区间 $[a_n,b_n]$ 嵌套递减且 $b_n-a_n\to0$。证明交集恰含一个实数。
\par\textit{补充证明。来源：编者补充：本科数学分析标准命题与反例。}\par\hyperref[ans:c2q11]{解答见第~\pageref{ans:c2q11} 页。}
\end{example}
\begin{example}[有界数列必有收敛子列]\label{ex:c2q12}
用二分嵌套区间法证明 Bolzano--Weierstrass 定理。
\par\textit{补充证明。来源：编者补充：本科数学分析标准命题与反例。}\par\hyperref[ans:c2q12]{解答见第~\pageref{ans:c2q12} 页。}
\end{example}
\begin{example}[Cesàro 平均的收敛]\label{ex:c2q13}
若 $a_n\to L$，证明 $(a_1+\cdots+a_n)/n\to L$。
\par\textit{补充证明。来源：编者补充：本科数学分析标准命题与反例。}\par\hyperref[ans:c2q13]{解答见第~\pageref{ans:c2q13} 页。}
\end{example}
\begin{example}[相邻差趋零并不保证收敛]\label{ex:c2q14}
举例说明 $a_{n+1}-a_n\to0$ 不足以保证数列收敛，也不足以保证有界。
\par\textit{补充证明。来源：编者补充：本科数学分析标准命题与反例。}\par\hyperref[ans:c2q14]{解答见第~\pageref{ans:c2q14} 页。}
\end{example}
\begin{example}[乘积数列的极限]\label{ex:c2q15}
直接给出 $a_n\to a,b_n\to b$ 时 $a_nb_n\to ab$ 的误差估计。
\par\textit{补充证明。来源：编者补充：本科数学分析标准命题与反例。}\par\hyperref[ans:c2q15]{解答见第~\pageref{ans:c2q15} 页。}
\end{example}
\begin{example}[倒数数列的极限]\label{ex:c2q16}
若 $a_n\to a\ne0$，证明从某项起 $a_n\ne0$ 且 $1/a_n\to1/a$。
\par\textit{补充证明。来源：编者补充：本科数学分析标准命题与反例。}\par\hyperref[ans:c2q16]{解答见第~\pageref{ans:c2q16} 页。}
\end{example}
\begin{example}[收敛数列与极值取得]\label{ex:c2q17}
若 $a_n\to L$ 且存在某项 $a_k>L$，证明数列的上确界由某个项取得。
\par\textit{补充证明。来源：编者补充：本科数学分析标准命题与反例。}\par\hyperref[ans:c2q17]{解答见第~\pageref{ans:c2q17} 页。}
\end{example}
\begin{example}[递推压缩的几何误差]\label{ex:c2q18}
设 $0<q<1$ 且数列满足 $|a_{n+1}-L|\le q|a_n-L|$。证明收敛并给出误差界。
\par\textit{补充证明。来源：编者补充：本科数学分析标准命题与反例。}\par\hyperref[ans:c2q18]{解答见第~\pageref{ans:c2q18} 页。}
\end{example}
\begin{example}[递推收敛与不动点存在]\label{ex:c2q19}
设 $T:[a,b]\to[a,b]$ 满足 $|T(x)-T(y)|\le q|x-y|$，其中 $0<q<1$。从任意 $x_0\in[a,b]$ 出发，令 $x_{n+1}=T(x_n)$，证明收敛到唯一不动点。
\par\textit{补充证明。来源：编者补充：本科数学分析标准命题与反例。}\par\hyperref[ans:c2q19]{解答见第~\pageref{ans:c2q19} 页。}
\end{example}
\begin{example}[上、下极限的简单计算]\label{ex:c2q20}
求数列 $a_n=(-1)^n+1/n$ 的上极限与下极限，并说明为何二者不等就不收敛。
\par\textit{补充证明。来源：编者补充：本科数学分析标准命题与反例。}\par\hyperref[ans:c2q20]{解答见第~\pageref{ans:c2q20} 页。}
\end{example}
\subsection{解答}
\begin{solve}[习题~\ref{ex:c2q1}：上确界的逼近性质]\phantomsection\label{ans:c2q1}
因为 $s$ 为上界，所有 $a\in A$ 都满足 $a\le s$。如果某个 $\varepsilon>0$ 对应的区间内没有集合元素，那么所有 $a\in A$ 均满足 $a\le s-\varepsilon$，即 $s-\varepsilon$ 也是上界，与 $s$ 是最小上界矛盾。反过来，若 $s$ 是上界且对每个正 $\varepsilon$ 都存在这样的 $a$，则任何 $u<s$ 都不是上界：取 $\varepsilon=s-u$，便有 $a>u$。所以 $s$ 确为上确界。逆命题不能删去“$s$ 为上界”，因为仅有集合在 $s$ 左侧任意逼近，并不排除集合中还有比 $s$ 大的元素。
\end{solve}

\begin{solve}[习题~\ref{ex:c2q2}：从上确界构造逼近数列]\phantomsection\label{ans:c2q2}
对每个正整数 $n$，由上确界逼近性质选 $a_n\in A$，使 $s-1/n<a_n\le s$。于是 $0\le s-a_n<1/n$，夹逼定理给出 $a_n\to s$。选出的数列未必单调。定义 $b_n=\max(a_1,\ldots,a_n)$，因有限最大值等于其中某个元素，仍有 $b_n\in A$；且 $b_n$ 单调不减、$a_n\le b_n\le s$，故也收敛到 $s$。一般不能要求严格递增，例如 $A=\{s\}$ 只有一个元素。若 $s\notin A$，则可以递归选取比前项更接近 $s$ 的元素，得到严格递增的逼近列。
\end{solve}

\begin{solve}[习题~\ref{ex:c2q3}：单调有界定理的证明]\phantomsection\label{ans:c2q3}
设 $s=\sup\{a_n:n\ge1\}$。给定 $\varepsilon>0$，由上确界性质存在 $N$ 使 $a_N>s-\varepsilon$。对所有 $n\ge N$，单调性给出 $a_n\ge a_N$，上界性质给出 $a_n\le s$，所以 $s-\varepsilon<a_n\le s$，即 $|a_n-s|<\varepsilon$。这正是数列收敛定义。上界只保证上确界有限；单调性把一个接近上确界的项推广为整个尾部都接近它。两者缺一不可：$n$ 单调但无界，$(-1)^n$ 有界但不单调，均不收敛。
\end{solve}

\begin{solve}[习题~\ref{ex:c2q4}：收敛数列有界]\phantomsection\label{ans:c2q4}
设 $a_n\to a$。在定义中取 $\varepsilon=1$，存在 $N$ 使 $n\ge N$ 时 $|a_n-a|<1$，于是 $|a_n|<|a|+1$。前面只有有限项，令 $M=\max(|a_1|,\ldots,|a_{N-1}|,|a|+1)$，便对全部 $n$ 有 $|a_n|\le M$。证明必须处理前面有限项，不能只给出尾部有界就结束。逆命题不成立：$a_n=(-1)^n$ 有界，但偶数子列恒为1、奇数子列恒为负1，因此不可能有同一个极限。
\end{solve}

\begin{solve}[习题~\ref{ex:c2q5}：极限的唯一性]\phantomsection\label{ans:c2q5}
假设 $a_n\to a$ 且 $a_n\to b$，而 $a\ne b$。取 $\varepsilon=|a-b|/3$。存在 $N_1,N_2$，使相应尾部满足 $|a_n-a|<\varepsilon$ 与 $|a_n-b|<\varepsilon$。取 $n\ge\max(N_1,N_2)$，由三角不等式得 $|a-b|\le|a-a_n|+|a_n-b|<2|a-b|/3$，矛盾。这里两个收敛条件可能各有不同阈值，因此必须取最大值，使同一个下标同时满足二者。
\end{solve}

\begin{solve}[习题~\ref{ex:c2q6}：保号性与严格不等式]\phantomsection\label{ans:c2q6}
取 $\varepsilon=a/2$。从某项起 $|a_n-a|<a/2$，因此 $a_n>a/2>0$。反方向只能推出 $a\ge0$，不能推出严格为正；例如 $a_n=1/n>0$，极限却为0。若想保证极限严格为正，一个足够的附加条件是所有项都满足 $a_n\ge c>0$，此时极限满足 $a\ge c$。严格不等号在取极限时可能变成等号，这是极限计算中需要经常检查的边界现象。
\end{solve}

\begin{solve}[习题~\ref{ex:c2q7}：收敛子列与原列]\phantomsection\label{ans:c2q7}
设 $a_n\to a$，子列下标 $n_k$ 严格递增，因此 $n_k\ge k$。给定 $\varepsilon>0$，选 $N$ 使 $n\ge N$ 时 $|a_n-a|<\varepsilon$。当 $k\ge N$ 时，$n_k\ge N$，所以 $|a_{n_k}-a|<\varepsilon$，子列收敛到 $a$。逆命题不成立：$(-1)^n$ 的偶数子列收敛到1，但原列不收敛。只有一个“好”的子列说明原列在部分位置接近某个数，并不能控制所有充分大的下标。
\end{solve}

\begin{solve}[习题~\ref{ex:c2q8}：奇偶子列判别]\phantomsection\label{ans:c2q8}
给定 $\varepsilon>0$，分别选阈值 $N_1,N_2$ 控制偶、奇子列，令 $N=2\max(N_1,N_2)+1$。任何 $n\ge N$ 必为足够大的偶数或奇数，所以都有 $|a_n-L|<\varepsilon$，从而原列收敛到 $L$。若两子列极限不同，原列若收敛，其全部子列必须同极限，产生矛盾，因此原列不收敛。本题成立的关键是两条子列覆盖原列的所有下标。
\end{solve}

\begin{solve}[习题~\ref{ex:c2q9}：Cauchy 数列必有界]\phantomsection\label{ans:c2q9}
取 $\varepsilon=1$，并固定一个满足条件的下标 $N$。对 $n\ge N$ 有 $|a_n|\le|a_N|+|a_n-a_N|<|a_N|+1$。再把前面有限项的绝对值与 $|a_N|+1$ 取最大值，即得全列的上界。Cauchy 条件虽然没有预先给出极限候选值，仍能通过固定一个尾部项把整个尾部限制在一个有界区间。
\end{solve}

\begin{solve}[习题~\ref{ex:c2q10}：Cauchy 判别的充分性]\phantomsection\label{ans:c2q10}
Cauchy 数列有界，所以存在收敛子列 $a_{n_k}\to L$。给定 $\varepsilon>0$，Cauchy 条件给出 $N$，使 $m,n\ge N$ 时距离小于 $\varepsilon/2$。选足够大的 $k$，使 $n_k\ge N$ 且 $|a_{n_k}-L|<\varepsilon/2$。对所有 $n\ge N$，三角不等式给出 $|a_n-L|\le|a_n-a_{n_k}|+|a_{n_k}-L|<\varepsilon$，所以原列收敛。子列提供极限候选值，Cauchy 条件则把其余尾部项一并拉近该候选值。
\end{solve}

\begin{solve}[习题~\ref{ex:c2q11}：嵌套区间定理]\phantomsection\label{ans:c2q11}
由嵌套性，$a_n$ 单调不减、$b_n$ 单调不增，且 $a_1\le a_n\le b_n\le b_1$。两列由单调有界定理收敛，设极限为 $a,b$。长度趋零说明 $a=b=:c$。对每个固定的 $n$，嵌套性和取极限给出 $a_n\le c\le b_n$，所以 $c$ 属于每个区间。若另有 $d$ 也在全部区间内，则 $|c-d|\le b_n-a_n\to0$，故 $d=c$。闭区间条件保证极限点仍在区间内；对开区间如 $(0,1/n)$，交集可能为空。
\end{solve}

\begin{solve}[习题~\ref{ex:c2q12}：有界数列必有收敛子列]\phantomsection\label{ans:c2q12}
把有界数列全部放入某个闭区间。将区间二等分，至少有一个半区间含有无穷多个数列项；保留它并重复。得到闭区间嵌套列，长度每次减半，且每个区间含无穷多个项。递归选择严格递增下标 $n_k$，使 $a_{n_k}$ 位于第 $k$ 个区间中。由嵌套区间定理交集为唯一点 $c$，且 $|a_{n_k}-c|$ 不超过第 $k$ 个区间长度，故子列收敛到 $c$。每个区间含无穷多个项保证总能选到比前次更大的下标。
\end{solve}

\begin{solve}[习题~\ref{ex:c2q13}：Cesàro 平均的收敛]\phantomsection\label{ans:c2q13}
令 $b_n=a_n-L$，只须证明其平均趋零。给定 $\varepsilon>0$，选 $N$ 使 $k\ge N$ 时 $|b_k|<\varepsilon/2$。令 $C=\sum_{k=1}^{N-1}|b_k|$，则 $n\ge N$ 时，$|n^{-1}\sum_{k=1}^nb_k|\le C/n+\varepsilon/2$。再取 $n>2C/\varepsilon$，右侧小于 $\varepsilon$。有限个早期大误差被分母 $n$ 稀释，尾部每项则已经统一小。逆命题不成立，例如交错的 $\pm1$ 数列不收敛，其算术平均却趋零。
\end{solve}

\begin{solve}[习题~\ref{ex:c2q14}：相邻差趋零并不保证收敛]\phantomsection\label{ans:c2q14}
取 $a_n=\log n$。则 $a_{n+1}-a_n=\log(1+1/n)\to0$，但 $a_n\to\infty$，所以无界且不收敛。另取 $a_n=\sin\sqrt n$，由正弦的 Lipschitz 估计有 $|a_{n+1}-a_n|\le\sqrt{n+1}-\sqrt n\to0$，它虽有界却不收敛；可选整数 $n_k$ 最接近 $(\pi/2+2k\pi)^2$，使一条子列趋于1，再选最接近 $(3\pi/2+2k\pi)^2$ 的整数得到趋于负1的子列。相邻差只控制一步变化，不能控制任意远的尾部项之间的距离。
\end{solve}

\begin{solve}[习题~\ref{ex:c2q15}：乘积数列的极限]\phantomsection\label{ans:c2q15}
由于收敛数列有界，可取 $M\ge1$ 使 $|a_n|\le M$。分拆误差为 $a_nb_n-ab=a_n(b_n-b)+b(a_n-a)$，故其绝对值不超过 $M|b_n-b|+|b||a_n-a|$。分别把两个误差控制在 $\varepsilon/(2M)$ 与 $\varepsilon/(2(|b|+1))$ 内，再取共同阈值，就使总误差小于 $\varepsilon$。这里先用有界性控制一个因子，避免两个随下标变化的误差相互耦合。
\end{solve}

\begin{solve}[习题~\ref{ex:c2q16}：倒数数列的极限]\phantomsection\label{ans:c2q16}
由 $|a_n-a|<|a|/2$ 可知从某项起 $|a_n|\ge|a|-|a_n-a|>|a|/2$，特别地分母非零。在该尾部上，$|1/a_n-1/a|=|a_n-a|/(|a_n||a|)\le2|a_n-a|/|a|^2\to0$。条件 $a\ne0$ 至关重要：若 $a_n=1/n\to0$，倒数趋于无穷；若 $a_n=(-1)^n/n$，倒数甚至在正负方向之间振荡。
\end{solve}

\begin{solve}[习题~\ref{ex:c2q17}：收敛数列与极值取得]\phantomsection\label{ans:c2q17}
令 $d=a_k-L>0$。从某项起 $|a_n-L|<d/2$，所以尾部所有项都小于 $L+d/2<a_k$。因此上确界只可能来自前面有限多个项，而有限集合必有最大值。这个最大值就是全列上确界。若没有任何项大于极限，结论可能失败：$a_n=1-1/n$ 的上确界为1，但没有一项等于1。
\end{solve}

\begin{solve}[习题~\ref{ex:c2q18}：递推压缩的几何误差]\phantomsection\label{ans:c2q18}
逐次迭代该不等式得到 $|a_n-L|\le q^{n-1}|a_1-L|$。由于 $q^{n-1}\to0$，夹逼定理给出 $a_n\to L$。若初始误差非零，为使误差小于 $\varepsilon$，可选 $n$ 满足 $(n-1)\log q<\log(\varepsilon/|a_1-L|)$；因 $\log q<0$，除以它时不等号方向必须反转。这类几何误差估计同时给出收敛证明与达到给定精度的迭代次数上界。
\end{solve}

\begin{solve}[习题~\ref{ex:c2q19}：递推收敛与不动点存在]\phantomsection\label{ans:c2q19}
首先所有迭代都留在闭区间内。相邻差满足 $|x_{n+1}-x_n|\le q^n|x_1-x_0|$，对 $m>n$ 求和得 $|x_m-x_n|\le q^n|x_1-x_0|/(1-q)$，因此为 Cauchy 数列。实数完备性保证它收敛到 $L$，闭区间保证 $L\in[a,b]$。Lipschitz 条件蕴含连续性，所以 $L=T(L)$。若 $K$ 为另一个不动点，则 $|L-K|=|T(L)-T(K)|\le q|L-K|$，迫使 $L=K$。取 $m\to\infty$ 还得到误差界 $|L-x_n|\le q^n|x_1-x_0|/(1-q)$。
\end{solve}

\begin{solve}[习题~\ref{ex:c2q20}：上、下极限的简单计算]\phantomsection\label{ans:c2q20}
偶数子列为 $1+1/(2n)$，趋于1；奇数子列为 $-1+1/(2n-1)$，趋于负1。任意尾部都同时含有这两类项，尾部上确界趋于1、尾部下确界趋于负1，因此 $\limsup a_n=1$、$\liminf a_n=-1$。若一个有界数列收敛到 $L$，其尾部最终全部落入 $(L-\varepsilon,L+\varepsilon)$，从而尾部上下确界也趋于 $L$。本题两者不同，所以原列不收敛。
\end{solve}

\clearpage
\section{函数极限与无穷小比较}\label{sec:ch3}
函数极限只涉及去心邻域内的取值。遇到 $0/0$ 型表达式，应寻找公因子、有理化或局部展开；两个趋零函数之商可以趋于任意有限数、无穷或根本没有极限。

\textbf{资源对应：}CLP-1 \S~1.4、\S~1.5、\S~3.7。本章共 25 道编号习题。
\subsection{习题}
\begin{example}[零比零不能直接运算]\label{ex:c3q1}
构造在 $x\to3$ 时均趋于 $0$ 的函数 $f,g$，使 $f/g$ 分别趋于 $10$、$0$、$+\infty$，再构造一个商没有极限的例子。
\par\textit{选译并补充。来源：CLP-1 \S~1.4，标签 s1.4zeroes1--3 及其后题。}\par\hyperref[ans:c3q1]{解答见第~\pageref{ans:c3q1} 页。}
\end{example}
\begin{example}[差幂因式分解]\label{ex:c3q2}
求 $\displaystyle\lim_{x\to 1}\dfrac{x^{3}-1}{x-1}$，要求不用洛必达法则。
\par\textit{改编训练。}\par\hyperref[ans:c3q2]{解答见第~\pageref{ans:c3q2} 页。}
\end{example}
\begin{example}[差幂因式分解]\label{ex:c3q3}
求 $\displaystyle\lim_{x\to 2}\dfrac{x^{3}-8}{x-2}$，要求不用洛必达法则。
\par\textit{改编训练。}\par\hyperref[ans:c3q3]{解答见第~\pageref{ans:c3q3} 页。}
\end{example}
\begin{example}[差幂因式分解]\label{ex:c3q4}
求 $\displaystyle\lim_{x\to 1}\dfrac{x^{4}-1}{x-1}$，要求不用洛必达法则。
\par\textit{改编训练。}\par\hyperref[ans:c3q4]{解答见第~\pageref{ans:c3q4} 页。}
\end{example}
\begin{example}[差幂因式分解]\label{ex:c3q5}
求 $\displaystyle\lim_{x\to 2}\dfrac{x^{4}-16}{x-2}$，要求不用洛必达法则。
\par\textit{改编训练。}\par\hyperref[ans:c3q5]{解答见第~\pageref{ans:c3q5} 页。}
\end{example}
\begin{example}[差幂因式分解]\label{ex:c3q6}
求 $\displaystyle\lim_{x\to 3}\dfrac{x^{3}-27}{x-3}$，要求不用洛必达法则。
\par\textit{改编训练。}\par\hyperref[ans:c3q6]{解答见第~\pageref{ans:c3q6} 页。}
\end{example}
\begin{example}[差幂因式分解]\label{ex:c3q7}
求 $\displaystyle\lim_{x\to 2}\dfrac{x^{5}-32}{x-2}$，要求不用洛必达法则。
\par\textit{改编训练。}\par\hyperref[ans:c3q7]{解答见第~\pageref{ans:c3q7} 页。}
\end{example}
\begin{example}[三角函数的局部比较]\label{ex:c3q8}
求 $\displaystyle\lim_{x\to0}\frac{\sin(1x)}{\tan(2x)}$ 以及 $\displaystyle\lim_{x\to0}\frac{1-\cos(1x)}{x^2}$。
\par\textit{改编训练。}\par\hyperref[ans:c3q8]{解答见第~\pageref{ans:c3q8} 页。}
\end{example}
\begin{example}[三角函数的局部比较]\label{ex:c3q9}
求 $\displaystyle\lim_{x\to0}\frac{\sin(2x)}{\tan(3x)}$ 以及 $\displaystyle\lim_{x\to0}\frac{1-\cos(2x)}{x^2}$。
\par\textit{改编训练。}\par\hyperref[ans:c3q9]{解答见第~\pageref{ans:c3q9} 页。}
\end{example}
\begin{example}[三角函数的局部比较]\label{ex:c3q10}
求 $\displaystyle\lim_{x\to0}\frac{\sin(3x)}{\tan(5x)}$ 以及 $\displaystyle\lim_{x\to0}\frac{1-\cos(3x)}{x^2}$。
\par\textit{改编训练。}\par\hyperref[ans:c3q10]{解答见第~\pageref{ans:c3q10} 页。}
\end{example}
\begin{example}[三角函数的局部比较]\label{ex:c3q11}
求 $\displaystyle\lim_{x\to0}\frac{\sin(4x)}{\tan(1x)}$ 以及 $\displaystyle\lim_{x\to0}\frac{1-\cos(4x)}{x^2}$。
\par\textit{改编训练。}\par\hyperref[ans:c3q11]{解答见第~\pageref{ans:c3q11} 页。}
\end{example}
\begin{example}[三角函数的局部比较]\label{ex:c3q12}
求 $\displaystyle\lim_{x\to0}\frac{\sin(5x)}{\tan(2x)}$ 以及 $\displaystyle\lim_{x\to0}\frac{1-\cos(5x)}{x^2}$。
\par\textit{改编训练。}\par\hyperref[ans:c3q12]{解答见第~\pageref{ans:c3q12} 页。}
\end{example}
\begin{example}[三角函数的局部比较]\label{ex:c3q13}
求 $\displaystyle\lim_{x\to0}\frac{\sin(6x)}{\tan(7x)}$ 以及 $\displaystyle\lim_{x\to0}\frac{1-\cos(6x)}{x^2}$。
\par\textit{改编训练。}\par\hyperref[ans:c3q13]{解答见第~\pageref{ans:c3q13} 页。}
\end{example}
\begin{example}[对数与指数复合的极限]\label{ex:c3q14}
求 $\displaystyle\lim_{x\to0}\frac{e^{1x}-1}{\log(1+2x)}$，并写出使对数有定义的局部条件。
\par\textit{改编训练。}\par\hyperref[ans:c3q14]{解答见第~\pageref{ans:c3q14} 页。}
\end{example}
\begin{example}[对数与指数复合的极限]\label{ex:c3q15}
求 $\displaystyle\lim_{x\to0}\frac{e^{2x}-1}{\log(1+1x)}$，并写出使对数有定义的局部条件。
\par\textit{改编训练。}\par\hyperref[ans:c3q15]{解答见第~\pageref{ans:c3q15} 页。}
\end{example}
\begin{example}[对数与指数复合的极限]\label{ex:c3q16}
求 $\displaystyle\lim_{x\to0}\frac{e^{3x}-1}{\log(1+2x)}$，并写出使对数有定义的局部条件。
\par\textit{改编训练。}\par\hyperref[ans:c3q16]{解答见第~\pageref{ans:c3q16} 页。}
\end{example}
\begin{example}[对数与指数复合的极限]\label{ex:c3q17}
求 $\displaystyle\lim_{x\to0}\frac{e^{2x}-1}{\log(1+5x)}$，并写出使对数有定义的局部条件。
\par\textit{改编训练。}\par\hyperref[ans:c3q17]{解答见第~\pageref{ans:c3q17} 页。}
\end{example}
\begin{example}[对数与指数复合的极限]\label{ex:c3q18}
求 $\displaystyle\lim_{x\to0}\frac{e^{4x}-1}{\log(1+3x)}$，并写出使对数有定义的局部条件。
\par\textit{改编训练。}\par\hyperref[ans:c3q18]{解答见第~\pageref{ans:c3q18} 页。}
\end{example}
\begin{example}[对数与指数复合的极限]\label{ex:c3q19}
求 $\displaystyle\lim_{x\to0}\frac{e^{5x}-1}{\log(1+4x)}$，并写出使对数有定义的局部条件。
\par\textit{改编训练。}\par\hyperref[ans:c3q19]{解答见第~\pageref{ans:c3q19} 页。}
\end{example}
\begin{example}[消去低阶项后的极限]\label{ex:c3q20}
求 $\displaystyle\lim_{x\to0}\frac{e^{1x}-1-1x}{x^2}$ 和 $\displaystyle\lim_{x\to0}\frac{\sin(1x)-1x}{x^3}$。
\par\textit{改编训练。}\par\hyperref[ans:c3q20]{解答见第~\pageref{ans:c3q20} 页。}
\end{example}
\begin{example}[消去低阶项后的极限]\label{ex:c3q21}
求 $\displaystyle\lim_{x\to0}\frac{e^{2x}-1-2x}{x^2}$ 和 $\displaystyle\lim_{x\to0}\frac{\sin(2x)-2x}{x^3}$。
\par\textit{改编训练。}\par\hyperref[ans:c3q21]{解答见第~\pageref{ans:c3q21} 页。}
\end{example}
\begin{example}[消去低阶项后的极限]\label{ex:c3q22}
求 $\displaystyle\lim_{x\to0}\frac{e^{3x}-1-3x}{x^2}$ 和 $\displaystyle\lim_{x\to0}\frac{\sin(3x)-3x}{x^3}$。
\par\textit{改编训练。}\par\hyperref[ans:c3q22]{解答见第~\pageref{ans:c3q22} 页。}
\end{example}
\begin{example}[消去低阶项后的极限]\label{ex:c3q23}
求 $\displaystyle\lim_{x\to0}\frac{e^{4x}-1-4x}{x^2}$ 和 $\displaystyle\lim_{x\to0}\frac{\sin(4x)-4x}{x^3}$。
\par\textit{改编训练。}\par\hyperref[ans:c3q23]{解答见第~\pageref{ans:c3q23} 页。}
\end{example}
\begin{example}[消去低阶项后的极限]\label{ex:c3q24}
求 $\displaystyle\lim_{x\to0}\frac{e^{5x}-1-5x}{x^2}$ 和 $\displaystyle\lim_{x\to0}\frac{\sin(5x)-5x}{x^3}$。
\par\textit{改编训练。}\par\hyperref[ans:c3q24]{解答见第~\pageref{ans:c3q24} 页。}
\end{example}
\begin{example}[消去低阶项后的极限]\label{ex:c3q25}
求 $\displaystyle\lim_{x\to0}\frac{e^{6x}-1-6x}{x^2}$ 和 $\displaystyle\lim_{x\to0}\frac{\sin(6x)-6x}{x^3}$。
\par\textit{改编训练。}\par\hyperref[ans:c3q25]{解答见第~\pageref{ans:c3q25} 页。}
\end{example}
\subsection{解答}
\begin{solve}[习题~\ref{ex:c3q1}：零比零不能直接运算]\phantomsection\label{ans:c3q1}
前三种情形分别取 $(f,g)=(10(x-3),x-3)$、$((x-3)^2,x-3)$、$(x-3,(x-3)^3)$，其商在 $x\ne3$ 时分别为 $10$、$x-3$、$1/(x-3)^2$。第四种情形取 $f(x)=(x-3)\sin(1/(x-3))$、$g(x)=x-3$；由夹逼定理 $f\to0$，但其商沿两列趋向 $3$ 的点分别取 $1$ 与 $-1$。这些例子说明 $0/0$ 不是一个可直接运算的数。
\end{solve}

\begin{solve}[习题~\ref{ex:c3q2}：差幂因式分解]\phantomsection\label{ans:c3q2}
利用恒等式 $x^m-a^m=(x-a)\sum_{j=0}^{m-1}x^{m-1-j}a^j$。只在去心邻域内约去 $x-a$，得到
\[
\frac{x^{3}-1}{x-1}=\sum_{j=0}^{2}x^{2-j}1^j\quad(x\ne 1).
\]
右侧是多项式，代入 $x=1$ 后每一项均为 $1$，共有 $3$ 项，所以
\[
\lim_{x\to 1}\frac{x^{3}-1}{x-1}=3.
\]
原函数在目标点不一定有定义，但极限不受该点本身取值的影响。这一计算也是幂函数导数公式的一个直接验证。
\end{solve}

\begin{solve}[习题~\ref{ex:c3q3}：差幂因式分解]\phantomsection\label{ans:c3q3}
利用恒等式 $x^m-a^m=(x-a)\sum_{j=0}^{m-1}x^{m-1-j}a^j$。只在去心邻域内约去 $x-a$，得到
\[
\frac{x^{3}-8}{x-2}=\sum_{j=0}^{2}x^{2-j}2^j\quad(x\ne 2).
\]
右侧是多项式，代入 $x=2$ 后每一项均为 $4$，共有 $3$ 项，所以
\[
\lim_{x\to 2}\frac{x^{3}-8}{x-2}=12.
\]
原函数在目标点不一定有定义，但极限不受该点本身取值的影响。这一计算也是幂函数导数公式的一个直接验证。
\end{solve}

\begin{solve}[习题~\ref{ex:c3q4}：差幂因式分解]\phantomsection\label{ans:c3q4}
利用恒等式 $x^m-a^m=(x-a)\sum_{j=0}^{m-1}x^{m-1-j}a^j$。只在去心邻域内约去 $x-a$，得到
\[
\frac{x^{4}-1}{x-1}=\sum_{j=0}^{3}x^{3-j}1^j\quad(x\ne 1).
\]
右侧是多项式，代入 $x=1$ 后每一项均为 $1$，共有 $4$ 项，所以
\[
\lim_{x\to 1}\frac{x^{4}-1}{x-1}=4.
\]
原函数在目标点不一定有定义，但极限不受该点本身取值的影响。这一计算也是幂函数导数公式的一个直接验证。
\end{solve}

\begin{solve}[习题~\ref{ex:c3q5}：差幂因式分解]\phantomsection\label{ans:c3q5}
利用恒等式 $x^m-a^m=(x-a)\sum_{j=0}^{m-1}x^{m-1-j}a^j$。只在去心邻域内约去 $x-a$，得到
\[
\frac{x^{4}-16}{x-2}=\sum_{j=0}^{3}x^{3-j}2^j\quad(x\ne 2).
\]
右侧是多项式，代入 $x=2$ 后每一项均为 $8$，共有 $4$ 项，所以
\[
\lim_{x\to 2}\frac{x^{4}-16}{x-2}=32.
\]
原函数在目标点不一定有定义，但极限不受该点本身取值的影响。这一计算也是幂函数导数公式的一个直接验证。
\end{solve}

\begin{solve}[习题~\ref{ex:c3q6}：差幂因式分解]\phantomsection\label{ans:c3q6}
利用恒等式 $x^m-a^m=(x-a)\sum_{j=0}^{m-1}x^{m-1-j}a^j$。只在去心邻域内约去 $x-a$，得到
\[
\frac{x^{3}-27}{x-3}=\sum_{j=0}^{2}x^{2-j}3^j\quad(x\ne 3).
\]
右侧是多项式，代入 $x=3$ 后每一项均为 $9$，共有 $3$ 项，所以
\[
\lim_{x\to 3}\frac{x^{3}-27}{x-3}=27.
\]
原函数在目标点不一定有定义，但极限不受该点本身取值的影响。这一计算也是幂函数导数公式的一个直接验证。
\end{solve}

\begin{solve}[习题~\ref{ex:c3q7}：差幂因式分解]\phantomsection\label{ans:c3q7}
利用恒等式 $x^m-a^m=(x-a)\sum_{j=0}^{m-1}x^{m-1-j}a^j$。只在去心邻域内约去 $x-a$，得到
\[
\frac{x^{5}-32}{x-2}=\sum_{j=0}^{4}x^{4-j}2^j\quad(x\ne 2).
\]
右侧是多项式，代入 $x=2$ 后每一项均为 $16$，共有 $5$ 项，所以
\[
\lim_{x\to 2}\frac{x^{5}-32}{x-2}=80.
\]
原函数在目标点不一定有定义，但极限不受该点本身取值的影响。这一计算也是幂函数导数公式的一个直接验证。
\end{solve}

\begin{solve}[习题~\ref{ex:c3q8}：三角函数的局部比较]\phantomsection\label{ans:c3q8}
把各表达式拆成标准极限的乘积。首先
\[
\frac{\sin(1x)}{\tan(2x)}=\frac{1}{2}\frac{\sin(1x)}{1x}\frac{2x}{\sin(2x)}\cos(2x)\longrightarrow \frac{1}{2}.
\]
其次使用半角公式 $1-\cos t=2\sin^2(t/2)$，
\[
\frac{1-\cos(1x)}{x^2}=\frac{1}2\left(\frac{\sin(1x/2)}{1x/2}\right)^2\longrightarrow \frac{1}{2}.
\]
计算采用弧度制。第二式中的分子是二阶无穷小，若错误地把它按一阶无穷小处理，就会丢失所求常数。
\end{solve}

\begin{solve}[习题~\ref{ex:c3q9}：三角函数的局部比较]\phantomsection\label{ans:c3q9}
把各表达式拆成标准极限的乘积。首先
\[
\frac{\sin(2x)}{\tan(3x)}=\frac{2}{3}\frac{\sin(2x)}{2x}\frac{3x}{\sin(3x)}\cos(3x)\longrightarrow \frac{2}{3}.
\]
其次使用半角公式 $1-\cos t=2\sin^2(t/2)$，
\[
\frac{1-\cos(2x)}{x^2}=\frac{4}2\left(\frac{\sin(2x/2)}{2x/2}\right)^2\longrightarrow 2.
\]
计算采用弧度制。第二式中的分子是二阶无穷小，若错误地把它按一阶无穷小处理，就会丢失所求常数。
\end{solve}

\begin{solve}[习题~\ref{ex:c3q10}：三角函数的局部比较]\phantomsection\label{ans:c3q10}
把各表达式拆成标准极限的乘积。首先
\[
\frac{\sin(3x)}{\tan(5x)}=\frac{3}{5}\frac{\sin(3x)}{3x}\frac{5x}{\sin(5x)}\cos(5x)\longrightarrow \frac{3}{5}.
\]
其次使用半角公式 $1-\cos t=2\sin^2(t/2)$，
\[
\frac{1-\cos(3x)}{x^2}=\frac{9}2\left(\frac{\sin(3x/2)}{3x/2}\right)^2\longrightarrow \frac{9}{2}.
\]
计算采用弧度制。第二式中的分子是二阶无穷小，若错误地把它按一阶无穷小处理，就会丢失所求常数。
\end{solve}

\begin{solve}[习题~\ref{ex:c3q11}：三角函数的局部比较]\phantomsection\label{ans:c3q11}
把各表达式拆成标准极限的乘积。首先
\[
\frac{\sin(4x)}{\tan(1x)}=\frac{4}{1}\frac{\sin(4x)}{4x}\frac{1x}{\sin(1x)}\cos(1x)\longrightarrow 4.
\]
其次使用半角公式 $1-\cos t=2\sin^2(t/2)$，
\[
\frac{1-\cos(4x)}{x^2}=\frac{16}2\left(\frac{\sin(4x/2)}{4x/2}\right)^2\longrightarrow 8.
\]
计算采用弧度制。第二式中的分子是二阶无穷小，若错误地把它按一阶无穷小处理，就会丢失所求常数。
\end{solve}

\begin{solve}[习题~\ref{ex:c3q12}：三角函数的局部比较]\phantomsection\label{ans:c3q12}
把各表达式拆成标准极限的乘积。首先
\[
\frac{\sin(5x)}{\tan(2x)}=\frac{5}{2}\frac{\sin(5x)}{5x}\frac{2x}{\sin(2x)}\cos(2x)\longrightarrow \frac{5}{2}.
\]
其次使用半角公式 $1-\cos t=2\sin^2(t/2)$，
\[
\frac{1-\cos(5x)}{x^2}=\frac{25}2\left(\frac{\sin(5x/2)}{5x/2}\right)^2\longrightarrow \frac{25}{2}.
\]
计算采用弧度制。第二式中的分子是二阶无穷小，若错误地把它按一阶无穷小处理，就会丢失所求常数。
\end{solve}

\begin{solve}[习题~\ref{ex:c3q13}：三角函数的局部比较]\phantomsection\label{ans:c3q13}
把各表达式拆成标准极限的乘积。首先
\[
\frac{\sin(6x)}{\tan(7x)}=\frac{6}{7}\frac{\sin(6x)}{6x}\frac{7x}{\sin(7x)}\cos(7x)\longrightarrow \frac{6}{7}.
\]
其次使用半角公式 $1-\cos t=2\sin^2(t/2)$，
\[
\frac{1-\cos(6x)}{x^2}=\frac{36}2\left(\frac{\sin(6x/2)}{6x/2}\right)^2\longrightarrow 18.
\]
计算采用弧度制。第二式中的分子是二阶无穷小，若错误地把它按一阶无穷小处理，就会丢失所求常数。
\end{solve}

\begin{solve}[习题~\ref{ex:c3q14}：对数与指数复合的极限]\phantomsection\label{ans:c3q14}
在 $|x|<1/4$ 时有 $1+2x>0$，所以两侧足够靠近 $0$ 的点均在定义域内。把分式写成
\[
\frac{e^{1x}-1}{\log(1+2x)}=\frac{1}{2}\frac{e^{1x}-1}{1x}\frac{2x}{\log(1+2x)}\longrightarrow \frac{1}{2}.
\]
所用两个标准极限为 $(e^t-1)/t\to1$ 和 $\log(1+t)/t\to1$。这也可由二者在 $0$ 处的导数定义得到。等价无穷小的替换在这里用于乘除结构，不能不加检查地移入相减表达式中。
\end{solve}

\begin{solve}[习题~\ref{ex:c3q15}：对数与指数复合的极限]\phantomsection\label{ans:c3q15}
在 $|x|<1/2$ 时有 $1+1x>0$，所以两侧足够靠近 $0$ 的点均在定义域内。把分式写成
\[
\frac{e^{2x}-1}{\log(1+1x)}=\frac{2}{1}\frac{e^{2x}-1}{2x}\frac{1x}{\log(1+1x)}\longrightarrow 2.
\]
所用两个标准极限为 $(e^t-1)/t\to1$ 和 $\log(1+t)/t\to1$。这也可由二者在 $0$ 处的导数定义得到。等价无穷小的替换在这里用于乘除结构，不能不加检查地移入相减表达式中。
\end{solve}

\begin{solve}[习题~\ref{ex:c3q16}：对数与指数复合的极限]\phantomsection\label{ans:c3q16}
在 $|x|<1/4$ 时有 $1+2x>0$，所以两侧足够靠近 $0$ 的点均在定义域内。把分式写成
\[
\frac{e^{3x}-1}{\log(1+2x)}=\frac{3}{2}\frac{e^{3x}-1}{3x}\frac{2x}{\log(1+2x)}\longrightarrow \frac{3}{2}.
\]
所用两个标准极限为 $(e^t-1)/t\to1$ 和 $\log(1+t)/t\to1$。这也可由二者在 $0$ 处的导数定义得到。等价无穷小的替换在这里用于乘除结构，不能不加检查地移入相减表达式中。
\end{solve}

\begin{solve}[习题~\ref{ex:c3q17}：对数与指数复合的极限]\phantomsection\label{ans:c3q17}
在 $|x|<1/10$ 时有 $1+5x>0$，所以两侧足够靠近 $0$ 的点均在定义域内。把分式写成
\[
\frac{e^{2x}-1}{\log(1+5x)}=\frac{2}{5}\frac{e^{2x}-1}{2x}\frac{5x}{\log(1+5x)}\longrightarrow \frac{2}{5}.
\]
所用两个标准极限为 $(e^t-1)/t\to1$ 和 $\log(1+t)/t\to1$。这也可由二者在 $0$ 处的导数定义得到。等价无穷小的替换在这里用于乘除结构，不能不加检查地移入相减表达式中。
\end{solve}

\begin{solve}[习题~\ref{ex:c3q18}：对数与指数复合的极限]\phantomsection\label{ans:c3q18}
在 $|x|<1/6$ 时有 $1+3x>0$，所以两侧足够靠近 $0$ 的点均在定义域内。把分式写成
\[
\frac{e^{4x}-1}{\log(1+3x)}=\frac{4}{3}\frac{e^{4x}-1}{4x}\frac{3x}{\log(1+3x)}\longrightarrow \frac{4}{3}.
\]
所用两个标准极限为 $(e^t-1)/t\to1$ 和 $\log(1+t)/t\to1$。这也可由二者在 $0$ 处的导数定义得到。等价无穷小的替换在这里用于乘除结构，不能不加检查地移入相减表达式中。
\end{solve}

\begin{solve}[习题~\ref{ex:c3q19}：对数与指数复合的极限]\phantomsection\label{ans:c3q19}
在 $|x|<1/8$ 时有 $1+4x>0$，所以两侧足够靠近 $0$ 的点均在定义域内。把分式写成
\[
\frac{e^{5x}-1}{\log(1+4x)}=\frac{5}{4}\frac{e^{5x}-1}{5x}\frac{4x}{\log(1+4x)}\longrightarrow \frac{5}{4}.
\]
所用两个标准极限为 $(e^t-1)/t\to1$ 和 $\log(1+t)/t\to1$。这也可由二者在 $0$ 处的导数定义得到。等价无穷小的替换在这里用于乘除结构，不能不加检查地移入相减表达式中。
\end{solve}

\begin{solve}[习题~\ref{ex:c3q20}：消去低阶项后的极限]\phantomsection\label{ans:c3q20}
减法会消去展开式的首项，必须展开到第一个没有消失的阶数。由带 Peano 余项的泰勒公式，
\[
e^{1x}=1+1x+\frac{1}{2}x^2+o(x^2),\qquad\sin(1x)=1x-\frac{1}{6}x^3+o(x^3).
\]
代入并分别约去 $x^2$ 与 $x^3$，得到
\[
\frac{e^{1x}-1-1x}{x^2}\to \frac{1}{2},\qquad\frac{\sin(1x)-1x}{x^3}\to-\frac{1}{6}.
\]
这里 $o(x^m)/x^m\to0$。也可用相应次数的洛必达法则验证，但每次使用前都应确认分子、分母仍属于适用的未定型。
\end{solve}

\begin{solve}[习题~\ref{ex:c3q21}：消去低阶项后的极限]\phantomsection\label{ans:c3q21}
减法会消去展开式的首项，必须展开到第一个没有消失的阶数。由带 Peano 余项的泰勒公式，
\[
e^{2x}=1+2x+2x^2+o(x^2),\qquad\sin(2x)=2x-\frac{4}{3}x^3+o(x^3).
\]
代入并分别约去 $x^2$ 与 $x^3$，得到
\[
\frac{e^{2x}-1-2x}{x^2}\to 2,\qquad\frac{\sin(2x)-2x}{x^3}\to-\frac{4}{3}.
\]
这里 $o(x^m)/x^m\to0$。也可用相应次数的洛必达法则验证，但每次使用前都应确认分子、分母仍属于适用的未定型。
\end{solve}

\begin{solve}[习题~\ref{ex:c3q22}：消去低阶项后的极限]\phantomsection\label{ans:c3q22}
减法会消去展开式的首项，必须展开到第一个没有消失的阶数。由带 Peano 余项的泰勒公式，
\[
e^{3x}=1+3x+\frac{9}{2}x^2+o(x^2),\qquad\sin(3x)=3x-\frac{9}{2}x^3+o(x^3).
\]
代入并分别约去 $x^2$ 与 $x^3$，得到
\[
\frac{e^{3x}-1-3x}{x^2}\to \frac{9}{2},\qquad\frac{\sin(3x)-3x}{x^3}\to-\frac{9}{2}.
\]
这里 $o(x^m)/x^m\to0$。也可用相应次数的洛必达法则验证，但每次使用前都应确认分子、分母仍属于适用的未定型。
\end{solve}

\begin{solve}[习题~\ref{ex:c3q23}：消去低阶项后的极限]\phantomsection\label{ans:c3q23}
减法会消去展开式的首项，必须展开到第一个没有消失的阶数。由带 Peano 余项的泰勒公式，
\[
e^{4x}=1+4x+8x^2+o(x^2),\qquad\sin(4x)=4x-\frac{32}{3}x^3+o(x^3).
\]
代入并分别约去 $x^2$ 与 $x^3$，得到
\[
\frac{e^{4x}-1-4x}{x^2}\to 8,\qquad\frac{\sin(4x)-4x}{x^3}\to-\frac{32}{3}.
\]
这里 $o(x^m)/x^m\to0$。也可用相应次数的洛必达法则验证，但每次使用前都应确认分子、分母仍属于适用的未定型。
\end{solve}

\begin{solve}[习题~\ref{ex:c3q24}：消去低阶项后的极限]\phantomsection\label{ans:c3q24}
减法会消去展开式的首项，必须展开到第一个没有消失的阶数。由带 Peano 余项的泰勒公式，
\[
e^{5x}=1+5x+\frac{25}{2}x^2+o(x^2),\qquad\sin(5x)=5x-\frac{125}{6}x^3+o(x^3).
\]
代入并分别约去 $x^2$ 与 $x^3$，得到
\[
\frac{e^{5x}-1-5x}{x^2}\to \frac{25}{2},\qquad\frac{\sin(5x)-5x}{x^3}\to-\frac{125}{6}.
\]
这里 $o(x^m)/x^m\to0$。也可用相应次数的洛必达法则验证，但每次使用前都应确认分子、分母仍属于适用的未定型。
\end{solve}

\begin{solve}[习题~\ref{ex:c3q25}：消去低阶项后的极限]\phantomsection\label{ans:c3q25}
减法会消去展开式的首项，必须展开到第一个没有消失的阶数。由带 Peano 余项的泰勒公式，
\[
e^{6x}=1+6x+18x^2+o(x^2),\qquad\sin(6x)=6x-36x^3+o(x^3).
\]
代入并分别约去 $x^2$ 与 $x^3$，得到
\[
\frac{e^{6x}-1-6x}{x^2}\to 18,\qquad\frac{\sin(6x)-6x}{x^3}\to-36.
\]
这里 $o(x^m)/x^m\to0$。也可用相应次数的洛必达法则验证，但每次使用前都应确认分子、分母仍属于适用的未定型。
\end{solve}

\clearpage
\section{无穷远极限与渐近线}\label{sec:ch4}
无穷远极限常由增长速度和最高次项决定。水平、斜渐近线描述函数与一条直线之差趋零，而不仅是二者的比趋于1；两种描述不能混用。

\textbf{资源对应：}CLP-1 \S~1.5、\S~3.6、\S~3.7。本章共 24 道编号习题。
\subsection{习题}
\begin{example}[二次根式的两端渐近线]\label{ex:c4q1}
求 $f(x)=\sqrt{x^2+1x+5}$ 在 $x\to+\infty$ 与 $x\to-\infty$ 时的斜渐近线。
\par\textit{改编训练。}\par\hyperref[ans:c4q1]{解答见第~\pageref{ans:c4q1} 页。}
\end{example}
\begin{example}[二次根式的两端渐近线]\label{ex:c4q2}
求 $f(x)=\sqrt{x^2+2x+5}$ 在 $x\to+\infty$ 与 $x\to-\infty$ 时的斜渐近线。
\par\textit{改编训练。}\par\hyperref[ans:c4q2]{解答见第~\pageref{ans:c4q2} 页。}
\end{example}
\begin{example}[二次根式的两端渐近线]\label{ex:c4q3}
求 $f(x)=\sqrt{x^2+3x+13}$ 在 $x\to+\infty$ 与 $x\to-\infty$ 时的斜渐近线。
\par\textit{改编训练。}\par\hyperref[ans:c4q3]{解答见第~\pageref{ans:c4q3} 页。}
\end{example}
\begin{example}[二次根式的两端渐近线]\label{ex:c4q4}
求 $f(x)=\sqrt{x^2+2x+29}$ 在 $x\to+\infty$ 与 $x\to-\infty$ 时的斜渐近线。
\par\textit{改编训练。}\par\hyperref[ans:c4q4]{解答见第~\pageref{ans:c4q4} 页。}
\end{example}
\begin{example}[二次根式的两端渐近线]\label{ex:c4q5}
求 $f(x)=\sqrt{x^2+4x+25}$ 在 $x\to+\infty$ 与 $x\to-\infty$ 时的斜渐近线。
\par\textit{改编训练。}\par\hyperref[ans:c4q5]{解答见第~\pageref{ans:c4q5} 页。}
\end{example}
\begin{example}[二次根式的两端渐近线]\label{ex:c4q6}
求 $f(x)=\sqrt{x^2+5x+41}$ 在 $x\to+\infty$ 与 $x\to-\infty$ 时的斜渐近线。
\par\textit{改编训练。}\par\hyperref[ans:c4q6]{解答见第~\pageref{ans:c4q6} 页。}
\end{example}
\begin{example}[指数增长压过多项式]\label{ex:c4q7}
求 $\displaystyle\lim_{x\to\infty}x^{1}e^{-1x}$，要求给出一个不依赖反复求导的估计。
\par\textit{改编训练。}\par\hyperref[ans:c4q7]{解答见第~\pageref{ans:c4q7} 页。}
\end{example}
\begin{example}[指数增长压过多项式]\label{ex:c4q8}
求 $\displaystyle\lim_{x\to\infty}x^{2}e^{-2x}$，要求给出一个不依赖反复求导的估计。
\par\textit{改编训练。}\par\hyperref[ans:c4q8]{解答见第~\pageref{ans:c4q8} 页。}
\end{example}
\begin{example}[指数增长压过多项式]\label{ex:c4q9}
求 $\displaystyle\lim_{x\to\infty}x^{3}e^{-1x}$，要求给出一个不依赖反复求导的估计。
\par\textit{改编训练。}\par\hyperref[ans:c4q9]{解答见第~\pageref{ans:c4q9} 页。}
\end{example}
\begin{example}[指数增长压过多项式]\label{ex:c4q10}
求 $\displaystyle\lim_{x\to\infty}x^{2}e^{-3x}$，要求给出一个不依赖反复求导的估计。
\par\textit{改编训练。}\par\hyperref[ans:c4q10]{解答见第~\pageref{ans:c4q10} 页。}
\end{example}
\begin{example}[指数增长压过多项式]\label{ex:c4q11}
求 $\displaystyle\lim_{x\to\infty}x^{4}e^{-2x}$，要求给出一个不依赖反复求导的估计。
\par\textit{改编训练。}\par\hyperref[ans:c4q11]{解答见第~\pageref{ans:c4q11} 页。}
\end{example}
\begin{example}[指数增长压过多项式]\label{ex:c4q12}
求 $\displaystyle\lim_{x\to\infty}x^{5}e^{-1x}$，要求给出一个不依赖反复求导的估计。
\par\textit{改编训练。}\par\hyperref[ans:c4q12]{解答见第~\pageref{ans:c4q12} 页。}
\end{example}
\begin{example}[对数增长慢于正幂]\label{ex:c4q13}
求 $\displaystyle\lim_{x\to\infty}\frac{\log x}{x^{1/1}}$，并说明在 $x\to0^+$ 时 $x^{1/1}\log x$ 的极限。
\par\textit{改编训练。}\par\hyperref[ans:c4q13]{解答见第~\pageref{ans:c4q13} 页。}
\end{example}
\begin{example}[对数增长慢于正幂]\label{ex:c4q14}
求 $\displaystyle\lim_{x\to\infty}\frac{\log x}{x^{1/2}}$，并说明在 $x\to0^+$ 时 $x^{1/2}\log x$ 的极限。
\par\textit{改编训练。}\par\hyperref[ans:c4q14]{解答见第~\pageref{ans:c4q14} 页。}
\end{example}
\begin{example}[对数增长慢于正幂]\label{ex:c4q15}
求 $\displaystyle\lim_{x\to\infty}\frac{\log x}{x^{1/3}}$，并说明在 $x\to0^+$ 时 $x^{1/3}\log x$ 的极限。
\par\textit{改编训练。}\par\hyperref[ans:c4q15]{解答见第~\pageref{ans:c4q15} 页。}
\end{example}
\begin{example}[对数增长慢于正幂]\label{ex:c4q16}
求 $\displaystyle\lim_{x\to\infty}\frac{\log x}{x^{1/4}}$，并说明在 $x\to0^+$ 时 $x^{1/4}\log x$ 的极限。
\par\textit{改编训练。}\par\hyperref[ans:c4q16]{解答见第~\pageref{ans:c4q16} 页。}
\end{example}
\begin{example}[对数增长慢于正幂]\label{ex:c4q17}
求 $\displaystyle\lim_{x\to\infty}\frac{\log x}{x^{1/5}}$，并说明在 $x\to0^+$ 时 $x^{1/5}\log x$ 的极限。
\par\textit{改编训练。}\par\hyperref[ans:c4q17]{解答见第~\pageref{ans:c4q17} 页。}
\end{example}
\begin{example}[对数增长慢于正幂]\label{ex:c4q18}
求 $\displaystyle\lim_{x\to\infty}\frac{\log x}{x^{1/6}}$，并说明在 $x\to0^+$ 时 $x^{1/6}\log x$ 的极限。
\par\textit{改编训练。}\par\hyperref[ans:c4q18]{解答见第~\pageref{ans:c4q18} 页。}
\end{example}
\begin{example}[多项式除法求渐近线]\label{ex:c4q19}
求 $f(x)=\dfrac{x^2+1x+2}{x+1}$ 的斜渐近线，并讨论 $x=-1$ 是否为竖直渐近线。
\par\textit{改编训练。}\par\hyperref[ans:c4q19]{解答见第~\pageref{ans:c4q19} 页。}
\end{example}
\begin{example}[多项式除法求渐近线]\label{ex:c4q20}
求 $f(x)=\dfrac{x^2+2x+1}{x+1}$ 的斜渐近线，并讨论 $x=-1$ 是否为竖直渐近线。
\par\textit{改编训练。}\par\hyperref[ans:c4q20]{解答见第~\pageref{ans:c4q20} 页。}
\end{example}
\begin{example}[多项式除法求渐近线]\label{ex:c4q21}
求 $f(x)=\dfrac{x^2+3x+2}{x+1}$ 的斜渐近线，并讨论 $x=-1$ 是否为竖直渐近线。
\par\textit{改编训练。}\par\hyperref[ans:c4q21]{解答见第~\pageref{ans:c4q21} 页。}
\end{example}
\begin{example}[多项式除法求渐近线]\label{ex:c4q22}
求 $f(x)=\dfrac{x^2+2x+5}{x+1}$ 的斜渐近线，并讨论 $x=-1$ 是否为竖直渐近线。
\par\textit{改编训练。}\par\hyperref[ans:c4q22]{解答见第~\pageref{ans:c4q22} 页。}
\end{example}
\begin{example}[多项式除法求渐近线]\label{ex:c4q23}
求 $f(x)=\dfrac{x^2+4x+3}{x+1}$ 的斜渐近线，并讨论 $x=-1$ 是否为竖直渐近线。
\par\textit{改编训练。}\par\hyperref[ans:c4q23]{解答见第~\pageref{ans:c4q23} 页。}
\end{example}
\begin{example}[多项式除法求渐近线]\label{ex:c4q24}
求 $f(x)=\dfrac{x^2+5x+4}{x+1}$ 的斜渐近线，并讨论 $x=-1$ 是否为竖直渐近线。
\par\textit{改编训练。}\par\hyperref[ans:c4q24]{解答见第~\pageref{ans:c4q24} 页。}
\end{example}
\subsection{解答}
\begin{solve}[习题~\ref{ex:c4q1}：二次根式的两端渐近线]\phantomsection\label{ans:c4q1}
被根式的判别式为负，故函数处处有定义。正无穷处作有理化，
\[
f(x)-x=\frac{1x+5}{f(x)+x}\longrightarrow \frac{1}{2}.
\]
故正向渐近线为 $y=x+\frac{1}{2}$。负无穷处需使用 $|x|=-x$；同理
\[
f(x)+x=\frac{1x+5}{f(x)-x}\longrightarrow-\frac{1}{2},
\]
故负向渐近线为 $y=-x-\frac{1}{2}$。把 $\sqrt{x^2}$ 在整个实轴都写成 $x$，会错误地把两条渐近线合并。
\end{solve}

\begin{solve}[习题~\ref{ex:c4q2}：二次根式的两端渐近线]\phantomsection\label{ans:c4q2}
被根式的判别式为负，故函数处处有定义。正无穷处作有理化，
\[
f(x)-x=\frac{2x+5}{f(x)+x}\longrightarrow 1.
\]
故正向渐近线为 $y=x+1$。负无穷处需使用 $|x|=-x$；同理
\[
f(x)+x=\frac{2x+5}{f(x)-x}\longrightarrow-1,
\]
故负向渐近线为 $y=-x-1$。把 $\sqrt{x^2}$ 在整个实轴都写成 $x$，会错误地把两条渐近线合并。
\end{solve}

\begin{solve}[习题~\ref{ex:c4q3}：二次根式的两端渐近线]\phantomsection\label{ans:c4q3}
被根式的判别式为负，故函数处处有定义。正无穷处作有理化，
\[
f(x)-x=\frac{3x+13}{f(x)+x}\longrightarrow \frac{3}{2}.
\]
故正向渐近线为 $y=x+\frac{3}{2}$。负无穷处需使用 $|x|=-x$；同理
\[
f(x)+x=\frac{3x+13}{f(x)-x}\longrightarrow-\frac{3}{2},
\]
故负向渐近线为 $y=-x-\frac{3}{2}$。把 $\sqrt{x^2}$ 在整个实轴都写成 $x$，会错误地把两条渐近线合并。
\end{solve}

\begin{solve}[习题~\ref{ex:c4q4}：二次根式的两端渐近线]\phantomsection\label{ans:c4q4}
被根式的判别式为负，故函数处处有定义。正无穷处作有理化，
\[
f(x)-x=\frac{2x+29}{f(x)+x}\longrightarrow 1.
\]
故正向渐近线为 $y=x+1$。负无穷处需使用 $|x|=-x$；同理
\[
f(x)+x=\frac{2x+29}{f(x)-x}\longrightarrow-1,
\]
故负向渐近线为 $y=-x-1$。把 $\sqrt{x^2}$ 在整个实轴都写成 $x$，会错误地把两条渐近线合并。
\end{solve}

\begin{solve}[习题~\ref{ex:c4q5}：二次根式的两端渐近线]\phantomsection\label{ans:c4q5}
被根式的判别式为负，故函数处处有定义。正无穷处作有理化，
\[
f(x)-x=\frac{4x+25}{f(x)+x}\longrightarrow 2.
\]
故正向渐近线为 $y=x+2$。负无穷处需使用 $|x|=-x$；同理
\[
f(x)+x=\frac{4x+25}{f(x)-x}\longrightarrow-2,
\]
故负向渐近线为 $y=-x-2$。把 $\sqrt{x^2}$ 在整个实轴都写成 $x$，会错误地把两条渐近线合并。
\end{solve}

\begin{solve}[习题~\ref{ex:c4q6}：二次根式的两端渐近线]\phantomsection\label{ans:c4q6}
被根式的判别式为负，故函数处处有定义。正无穷处作有理化，
\[
f(x)-x=\frac{5x+41}{f(x)+x}\longrightarrow \frac{5}{2}.
\]
故正向渐近线为 $y=x+\frac{5}{2}$。负无穷处需使用 $|x|=-x$；同理
\[
f(x)+x=\frac{5x+41}{f(x)-x}\longrightarrow-\frac{5}{2},
\]
故负向渐近线为 $y=-x-\frac{5}{2}$。把 $\sqrt{x^2}$ 在整个实轴都写成 $x$，会错误地把两条渐近线合并。
\end{solve}

\begin{solve}[习题~\ref{ex:c4q7}：指数增长压过多项式]\phantomsection\label{ans:c4q7}
对 $x>0$，指数级数各项非负，故
\[
e^{1x}\ge\frac{(1x)^{2}}{2}.
\]
从而
\[
0\le x^{1}e^{-1x}\le 2\frac1x\longrightarrow0.
\]
夹逼定理给出所求极限为0。关键是选择指数级数中次数比原多项式高的一项；如果只选择同次数项，只能得到有界性，不能证明趋零。
\end{solve}

\begin{solve}[习题~\ref{ex:c4q8}：指数增长压过多项式]\phantomsection\label{ans:c4q8}
对 $x>0$，指数级数各项非负，故
\[
e^{2x}\ge\frac{(2x)^{3}}{6}.
\]
从而
\[
0\le x^{2}e^{-2x}\le \frac{3}{4}\frac1x\longrightarrow0.
\]
夹逼定理给出所求极限为0。关键是选择指数级数中次数比原多项式高的一项；如果只选择同次数项，只能得到有界性，不能证明趋零。
\end{solve}

\begin{solve}[习题~\ref{ex:c4q9}：指数增长压过多项式]\phantomsection\label{ans:c4q9}
对 $x>0$，指数级数各项非负，故
\[
e^{1x}\ge\frac{(1x)^{4}}{24}.
\]
从而
\[
0\le x^{3}e^{-1x}\le 24\frac1x\longrightarrow0.
\]
夹逼定理给出所求极限为0。关键是选择指数级数中次数比原多项式高的一项；如果只选择同次数项，只能得到有界性，不能证明趋零。
\end{solve}

\begin{solve}[习题~\ref{ex:c4q10}：指数增长压过多项式]\phantomsection\label{ans:c4q10}
对 $x>0$，指数级数各项非负，故
\[
e^{3x}\ge\frac{(3x)^{3}}{6}.
\]
从而
\[
0\le x^{2}e^{-3x}\le \frac{2}{9}\frac1x\longrightarrow0.
\]
夹逼定理给出所求极限为0。关键是选择指数级数中次数比原多项式高的一项；如果只选择同次数项，只能得到有界性，不能证明趋零。
\end{solve}

\begin{solve}[习题~\ref{ex:c4q11}：指数增长压过多项式]\phantomsection\label{ans:c4q11}
对 $x>0$，指数级数各项非负，故
\[
e^{2x}\ge\frac{(2x)^{5}}{120}.
\]
从而
\[
0\le x^{4}e^{-2x}\le \frac{15}{4}\frac1x\longrightarrow0.
\]
夹逼定理给出所求极限为0。关键是选择指数级数中次数比原多项式高的一项；如果只选择同次数项，只能得到有界性，不能证明趋零。
\end{solve}

\begin{solve}[习题~\ref{ex:c4q12}：指数增长压过多项式]\phantomsection\label{ans:c4q12}
对 $x>0$，指数级数各项非负，故
\[
e^{1x}\ge\frac{(1x)^{6}}{720}.
\]
从而
\[
0\le x^{5}e^{-1x}\le 720\frac1x\longrightarrow0.
\]
夹逼定理给出所求极限为0。关键是选择指数级数中次数比原多项式高的一项；如果只选择同次数项，只能得到有界性，不能证明趋零。
\end{solve}

\begin{solve}[习题~\ref{ex:c4q13}：对数增长慢于正幂]\phantomsection\label{ans:c4q13}
第一式令 $t=\log x$，则 $t\to\infty$，原式为 $te^{-t/1}$，由指数压过多项式得极限0。第二式令 $t=-\log x$，则
\[
x^{1/1}\log x=-te^{-t/1}\longrightarrow0.
\]
两个结果都等于0，但趋近方向和符号不同：第一式最终为正，第二式在 $0<x<1$ 上为负。用变量替换可以把有限端点与无穷远的增长比较统一起来，也说明不应把 $0\cdot(-\infty)$ 当作普通乘积。
\end{solve}

\begin{solve}[习题~\ref{ex:c4q14}：对数增长慢于正幂]\phantomsection\label{ans:c4q14}
第一式令 $t=\log x$，则 $t\to\infty$，原式为 $te^{-t/2}$，由指数压过多项式得极限0。第二式令 $t=-\log x$，则
\[
x^{1/2}\log x=-te^{-t/2}\longrightarrow0.
\]
两个结果都等于0，但趋近方向和符号不同：第一式最终为正，第二式在 $0<x<1$ 上为负。用变量替换可以把有限端点与无穷远的增长比较统一起来，也说明不应把 $0\cdot(-\infty)$ 当作普通乘积。
\end{solve}

\begin{solve}[习题~\ref{ex:c4q15}：对数增长慢于正幂]\phantomsection\label{ans:c4q15}
第一式令 $t=\log x$，则 $t\to\infty$，原式为 $te^{-t/3}$，由指数压过多项式得极限0。第二式令 $t=-\log x$，则
\[
x^{1/3}\log x=-te^{-t/3}\longrightarrow0.
\]
两个结果都等于0，但趋近方向和符号不同：第一式最终为正，第二式在 $0<x<1$ 上为负。用变量替换可以把有限端点与无穷远的增长比较统一起来，也说明不应把 $0\cdot(-\infty)$ 当作普通乘积。
\end{solve}

\begin{solve}[习题~\ref{ex:c4q16}：对数增长慢于正幂]\phantomsection\label{ans:c4q16}
第一式令 $t=\log x$，则 $t\to\infty$，原式为 $te^{-t/4}$，由指数压过多项式得极限0。第二式令 $t=-\log x$，则
\[
x^{1/4}\log x=-te^{-t/4}\longrightarrow0.
\]
两个结果都等于0，但趋近方向和符号不同：第一式最终为正，第二式在 $0<x<1$ 上为负。用变量替换可以把有限端点与无穷远的增长比较统一起来，也说明不应把 $0\cdot(-\infty)$ 当作普通乘积。
\end{solve}

\begin{solve}[习题~\ref{ex:c4q17}：对数增长慢于正幂]\phantomsection\label{ans:c4q17}
第一式令 $t=\log x$，则 $t\to\infty$，原式为 $te^{-t/5}$，由指数压过多项式得极限0。第二式令 $t=-\log x$，则
\[
x^{1/5}\log x=-te^{-t/5}\longrightarrow0.
\]
两个结果都等于0，但趋近方向和符号不同：第一式最终为正，第二式在 $0<x<1$ 上为负。用变量替换可以把有限端点与无穷远的增长比较统一起来，也说明不应把 $0\cdot(-\infty)$ 当作普通乘积。
\end{solve}

\begin{solve}[习题~\ref{ex:c4q18}：对数增长慢于正幂]\phantomsection\label{ans:c4q18}
第一式令 $t=\log x$，则 $t\to\infty$，原式为 $te^{-t/6}$，由指数压过多项式得极限0。第二式令 $t=-\log x$，则
\[
x^{1/6}\log x=-te^{-t/6}\longrightarrow0.
\]
两个结果都等于0，但趋近方向和符号不同：第一式最终为正，第二式在 $0<x<1$ 上为负。用变量替换可以把有限端点与无穷远的增长比较统一起来，也说明不应把 $0\cdot(-\infty)$ 当作普通乘积。
\end{solve}

\begin{solve}[习题~\ref{ex:c4q19}：多项式除法求渐近线]\phantomsection\label{ans:c4q19}
长除法给出
\[
f(x)=x+(0)+\frac{2}{x+1}.
\]
余项在 $x\to\pm\infty$ 时趋零，所以斜渐近线为 $y=x+(0)$。在 $x=-1$ 处分子为 $2\ne0$，故商的绝对值趋于无穷，$x=-1$ 是竖直渐近线。分母为零只是候选位置；必须检查是否与分子零点抵消。
\end{solve}

\begin{solve}[习题~\ref{ex:c4q20}：多项式除法求渐近线]\phantomsection\label{ans:c4q20}
长除法给出
\[
f(x)=x+(1)+\frac{0}{x+1}.
\]
余项在 $x\to\pm\infty$ 时趋零，所以斜渐近线为 $y=x+(1)$。在 $x=-1$ 处分子也为零，约去公因子后是一个一次函数，因此该点为可去间断点，不是竖直渐近线。分母为零只是候选位置；必须检查是否与分子零点抵消。
\end{solve}

\begin{solve}[习题~\ref{ex:c4q21}：多项式除法求渐近线]\phantomsection\label{ans:c4q21}
长除法给出
\[
f(x)=x+(2)+\frac{0}{x+1}.
\]
余项在 $x\to\pm\infty$ 时趋零，所以斜渐近线为 $y=x+(2)$。在 $x=-1$ 处分子也为零，约去公因子后是一个一次函数，因此该点为可去间断点，不是竖直渐近线。分母为零只是候选位置；必须检查是否与分子零点抵消。
\end{solve}

\begin{solve}[习题~\ref{ex:c4q22}：多项式除法求渐近线]\phantomsection\label{ans:c4q22}
长除法给出
\[
f(x)=x+(1)+\frac{4}{x+1}.
\]
余项在 $x\to\pm\infty$ 时趋零，所以斜渐近线为 $y=x+(1)$。在 $x=-1$ 处分子为 $4\ne0$，故商的绝对值趋于无穷，$x=-1$ 是竖直渐近线。分母为零只是候选位置；必须检查是否与分子零点抵消。
\end{solve}

\begin{solve}[习题~\ref{ex:c4q23}：多项式除法求渐近线]\phantomsection\label{ans:c4q23}
长除法给出
\[
f(x)=x+(3)+\frac{0}{x+1}.
\]
余项在 $x\to\pm\infty$ 时趋零，所以斜渐近线为 $y=x+(3)$。在 $x=-1$ 处分子也为零，约去公因子后是一个一次函数，因此该点为可去间断点，不是竖直渐近线。分母为零只是候选位置；必须检查是否与分子零点抵消。
\end{solve}

\begin{solve}[习题~\ref{ex:c4q24}：多项式除法求渐近线]\phantomsection\label{ans:c4q24}
长除法给出
\[
f(x)=x+(4)+\frac{0}{x+1}.
\]
余项在 $x\to\pm\infty$ 时趋零，所以斜渐近线为 $y=x+(4)$。在 $x=-1$ 处分子也为零，约去公因子后是一个一次函数，因此该点为可去间断点，不是竖直渐近线。分母为零只是候选位置；必须检查是否与分子零点抵消。
\end{solve}

\clearpage
\section{复合极限与渐近展开综合}\label{sec:ch5}
幂指型极限先取对数，差消型极限先确定需要保留的阶数。写出余项符号可以说明舍弃的部分在除以分母后确实消失；只写“约等于”不能替代极限论证。

\textbf{资源对应：}CLP-1 \S~1.4、\S~3.4、\S~3.7；CLP-2 \S~3.1。本章共 24 道编号习题。
\subsection{习题}
\begin{example}[幂指型函数极限]\label{ex:c5q1}
求 $\displaystyle\lim_{x\to0}(1+1x)^{2/x}$，解释为什么两侧足够近时都可取实对数。
\par\textit{改编训练。}\par\hyperref[ans:c5q1]{解答见第~\pageref{ans:c5q1} 页。}
\end{example}
\begin{example}[幂指型函数极限]\label{ex:c5q2}
求 $\displaystyle\lim_{x\to0}(1+2x)^{1/x}$，解释为什么两侧足够近时都可取实对数。
\par\textit{改编训练。}\par\hyperref[ans:c5q2]{解答见第~\pageref{ans:c5q2} 页。}
\end{example}
\begin{example}[幂指型函数极限]\label{ex:c5q3}
求 $\displaystyle\lim_{x\to0}(1+2x)^{3/x}$，解释为什么两侧足够近时都可取实对数。
\par\textit{改编训练。}\par\hyperref[ans:c5q3]{解答见第~\pageref{ans:c5q3} 页。}
\end{example}
\begin{example}[幂指型函数极限]\label{ex:c5q4}
求 $\displaystyle\lim_{x\to0}(1+3x)^{2/x}$，解释为什么两侧足够近时都可取实对数。
\par\textit{改编训练。}\par\hyperref[ans:c5q4]{解答见第~\pageref{ans:c5q4} 页。}
\end{example}
\begin{example}[幂指型函数极限]\label{ex:c5q5}
求 $\displaystyle\lim_{x\to0}(1+4x)^{1/x}$，解释为什么两侧足够近时都可取实对数。
\par\textit{改编训练。}\par\hyperref[ans:c5q5]{解答见第~\pageref{ans:c5q5} 页。}
\end{example}
\begin{example}[幂指型函数极限]\label{ex:c5q6}
求 $\displaystyle\lim_{x\to0}(1+3x)^{4/x}$，解释为什么两侧足够近时都可取实对数。
\par\textit{改编训练。}\par\hyperref[ans:c5q6]{解答见第~\pageref{ans:c5q6} 页。}
\end{example}
\begin{example}[缩放后的根式增量]\label{ex:c5q7}
求 $\displaystyle\lim_{n\to\infty}n\left(\sqrt{1+1/n}-1\right)$，再求其减去极限后乘以 $n$ 的极限。
\par\textit{改编训练。}\par\hyperref[ans:c5q7]{解答见第~\pageref{ans:c5q7} 页。}
\end{example}
\begin{example}[缩放后的根式增量]\label{ex:c5q8}
求 $\displaystyle\lim_{n\to\infty}n\left(\sqrt{1+2/n}-1\right)$，再求其减去极限后乘以 $n$ 的极限。
\par\textit{改编训练。}\par\hyperref[ans:c5q8]{解答见第~\pageref{ans:c5q8} 页。}
\end{example}
\begin{example}[缩放后的根式增量]\label{ex:c5q9}
求 $\displaystyle\lim_{n\to\infty}n\left(\sqrt{1+3/n}-1\right)$，再求其减去极限后乘以 $n$ 的极限。
\par\textit{改编训练。}\par\hyperref[ans:c5q9]{解答见第~\pageref{ans:c5q9} 页。}
\end{example}
\begin{example}[缩放后的根式增量]\label{ex:c5q10}
求 $\displaystyle\lim_{n\to\infty}n\left(\sqrt{1+4/n}-1\right)$，再求其减去极限后乘以 $n$ 的极限。
\par\textit{改编训练。}\par\hyperref[ans:c5q10]{解答见第~\pageref{ans:c5q10} 页。}
\end{example}
\begin{example}[缩放后的根式增量]\label{ex:c5q11}
求 $\displaystyle\lim_{n\to\infty}n\left(\sqrt{1+5/n}-1\right)$，再求其减去极限后乘以 $n$ 的极限。
\par\textit{改编训练。}\par\hyperref[ans:c5q11]{解答见第~\pageref{ans:c5q11} 页。}
\end{example}
\begin{example}[缩放后的根式增量]\label{ex:c5q12}
求 $\displaystyle\lim_{n\to\infty}n\left(\sqrt{1+6/n}-1\right)$，再求其减去极限后乘以 $n$ 的极限。
\par\textit{改编训练。}\par\hyperref[ans:c5q12]{解答见第~\pageref{ans:c5q12} 页。}
\end{example}
\begin{example}[三角函数差的三阶项]\label{ex:c5q13}
求 $\displaystyle\lim_{x\to0}\frac{\tan(1x)-\sin(1x)}{x^3}$。
\par\textit{改编训练。}\par\hyperref[ans:c5q13]{解答见第~\pageref{ans:c5q13} 页。}
\end{example}
\begin{example}[三角函数差的三阶项]\label{ex:c5q14}
求 $\displaystyle\lim_{x\to0}\frac{\tan(2x)-\sin(2x)}{x^3}$。
\par\textit{改编训练。}\par\hyperref[ans:c5q14]{解答见第~\pageref{ans:c5q14} 页。}
\end{example}
\begin{example}[三角函数差的三阶项]\label{ex:c5q15}
求 $\displaystyle\lim_{x\to0}\frac{\tan(3x)-\sin(3x)}{x^3}$。
\par\textit{改编训练。}\par\hyperref[ans:c5q15]{解答见第~\pageref{ans:c5q15} 页。}
\end{example}
\begin{example}[三角函数差的三阶项]\label{ex:c5q16}
求 $\displaystyle\lim_{x\to0}\frac{\tan(4x)-\sin(4x)}{x^3}$。
\par\textit{改编训练。}\par\hyperref[ans:c5q16]{解答见第~\pageref{ans:c5q16} 页。}
\end{example}
\begin{example}[三角函数差的三阶项]\label{ex:c5q17}
求 $\displaystyle\lim_{x\to0}\frac{\tan(5x)-\sin(5x)}{x^3}$。
\par\textit{改编训练。}\par\hyperref[ans:c5q17]{解答见第~\pageref{ans:c5q17} 页。}
\end{example}
\begin{example}[三角函数差的三阶项]\label{ex:c5q18}
求 $\displaystyle\lim_{x\to0}\frac{\tan(6x)-\sin(6x)}{x^3}$。
\par\textit{改编训练。}\par\hyperref[ans:c5q18]{解答见第~\pageref{ans:c5q18} 页。}
\end{example}
\begin{example}[对数减去线性项]\label{ex:c5q19}
求 $\displaystyle\lim_{x\to0}\frac{\log(1+1x)-1x}{x^2}$，并判断分子在0附近的符号。
\par\textit{改编训练。}\par\hyperref[ans:c5q19]{解答见第~\pageref{ans:c5q19} 页。}
\end{example}
\begin{example}[对数减去线性项]\label{ex:c5q20}
求 $\displaystyle\lim_{x\to0}\frac{\log(1+2x)-2x}{x^2}$，并判断分子在0附近的符号。
\par\textit{改编训练。}\par\hyperref[ans:c5q20]{解答见第~\pageref{ans:c5q20} 页。}
\end{example}
\begin{example}[对数减去线性项]\label{ex:c5q21}
求 $\displaystyle\lim_{x\to0}\frac{\log(1+3x)-3x}{x^2}$，并判断分子在0附近的符号。
\par\textit{改编训练。}\par\hyperref[ans:c5q21]{解答见第~\pageref{ans:c5q21} 页。}
\end{example}
\begin{example}[对数减去线性项]\label{ex:c5q22}
求 $\displaystyle\lim_{x\to0}\frac{\log(1+4x)-4x}{x^2}$，并判断分子在0附近的符号。
\par\textit{改编训练。}\par\hyperref[ans:c5q22]{解答见第~\pageref{ans:c5q22} 页。}
\end{example}
\begin{example}[对数减去线性项]\label{ex:c5q23}
求 $\displaystyle\lim_{x\to0}\frac{\log(1+5x)-5x}{x^2}$，并判断分子在0附近的符号。
\par\textit{改编训练。}\par\hyperref[ans:c5q23]{解答见第~\pageref{ans:c5q23} 页。}
\end{example}
\begin{example}[对数减去线性项]\label{ex:c5q24}
求 $\displaystyle\lim_{x\to0}\frac{\log(1+6x)-6x}{x^2}$，并判断分子在0附近的符号。
\par\textit{改编训练。}\par\hyperref[ans:c5q24]{解答见第~\pageref{ans:c5q24} 页。}
\end{example}
\subsection{解答}
\begin{solve}[习题~\ref{ex:c5q1}：幂指型函数极限]\phantomsection\label{ans:c5q1}
取 $|x|<1/2$，则 $1+1x>0$。设原式为 $Y(x)$，有
\[
\log Y(x)=\frac{2}x\log(1+1x)=2\frac{\log(1+1x)}{1x}\longrightarrow 2.
\]
由指数函数连续，$Y(x)\to e^{2}$。底数趋于1并不能直接说明结果为1，因为指数无界；取对数把这种相互竞争转化为一阶无穷小的比较。
\end{solve}

\begin{solve}[习题~\ref{ex:c5q2}：幂指型函数极限]\phantomsection\label{ans:c5q2}
取 $|x|<1/4$，则 $1+2x>0$。设原式为 $Y(x)$，有
\[
\log Y(x)=\frac{1}x\log(1+2x)=2\frac{\log(1+2x)}{2x}\longrightarrow 2.
\]
由指数函数连续，$Y(x)\to e^{2}$。底数趋于1并不能直接说明结果为1，因为指数无界；取对数把这种相互竞争转化为一阶无穷小的比较。
\end{solve}

\begin{solve}[习题~\ref{ex:c5q3}：幂指型函数极限]\phantomsection\label{ans:c5q3}
取 $|x|<1/4$，则 $1+2x>0$。设原式为 $Y(x)$，有
\[
\log Y(x)=\frac{3}x\log(1+2x)=6\frac{\log(1+2x)}{2x}\longrightarrow 6.
\]
由指数函数连续，$Y(x)\to e^{6}$。底数趋于1并不能直接说明结果为1，因为指数无界；取对数把这种相互竞争转化为一阶无穷小的比较。
\end{solve}

\begin{solve}[习题~\ref{ex:c5q4}：幂指型函数极限]\phantomsection\label{ans:c5q4}
取 $|x|<1/6$，则 $1+3x>0$。设原式为 $Y(x)$，有
\[
\log Y(x)=\frac{2}x\log(1+3x)=6\frac{\log(1+3x)}{3x}\longrightarrow 6.
\]
由指数函数连续，$Y(x)\to e^{6}$。底数趋于1并不能直接说明结果为1，因为指数无界；取对数把这种相互竞争转化为一阶无穷小的比较。
\end{solve}

\begin{solve}[习题~\ref{ex:c5q5}：幂指型函数极限]\phantomsection\label{ans:c5q5}
取 $|x|<1/8$，则 $1+4x>0$。设原式为 $Y(x)$，有
\[
\log Y(x)=\frac{1}x\log(1+4x)=4\frac{\log(1+4x)}{4x}\longrightarrow 4.
\]
由指数函数连续，$Y(x)\to e^{4}$。底数趋于1并不能直接说明结果为1，因为指数无界；取对数把这种相互竞争转化为一阶无穷小的比较。
\end{solve}

\begin{solve}[习题~\ref{ex:c5q6}：幂指型函数极限]\phantomsection\label{ans:c5q6}
取 $|x|<1/6$，则 $1+3x>0$。设原式为 $Y(x)$，有
\[
\log Y(x)=\frac{4}x\log(1+3x)=12\frac{\log(1+3x)}{3x}\longrightarrow 12.
\]
由指数函数连续，$Y(x)\to e^{12}$。底数趋于1并不能直接说明结果为1，因为指数无界；取对数把这种相互竞争转化为一阶无穷小的比较。
\end{solve}

\begin{solve}[习题~\ref{ex:c5q7}：缩放后的根式增量]\phantomsection\label{ans:c5q7}
第一步有理化得到 $n(\sqrt{1+1/n}-1)=1/(\sqrt{1+1/n}+1)\to \frac{1}{2}$。为求下一阶项，使用
\[
\sqrt{1+1/n}=1+\frac{1}{2n}-\frac{1}{8n^2}+O(n^{-3}).
\]
因此
\[
n\left[n\left(\sqrt{1+1/n}-1\right)-\frac{1}{2}\right]\longrightarrow-\frac{1}{8}.
\]
第一个极限只需要一阶项，第二个极限消去该项后必须保留二阶项；精度要求由分母或缩放因子的阶数决定。
\end{solve}

\begin{solve}[习题~\ref{ex:c5q8}：缩放后的根式增量]\phantomsection\label{ans:c5q8}
第一步有理化得到 $n(\sqrt{1+2/n}-1)=2/(\sqrt{1+2/n}+1)\to 1$。为求下一阶项，使用
\[
\sqrt{1+2/n}=1+\frac{2}{2n}-\frac{4}{8n^2}+O(n^{-3}).
\]
因此
\[
n\left[n\left(\sqrt{1+2/n}-1\right)-1\right]\longrightarrow-\frac{1}{2}.
\]
第一个极限只需要一阶项，第二个极限消去该项后必须保留二阶项；精度要求由分母或缩放因子的阶数决定。
\end{solve}

\begin{solve}[习题~\ref{ex:c5q9}：缩放后的根式增量]\phantomsection\label{ans:c5q9}
第一步有理化得到 $n(\sqrt{1+3/n}-1)=3/(\sqrt{1+3/n}+1)\to \frac{3}{2}$。为求下一阶项，使用
\[
\sqrt{1+3/n}=1+\frac{3}{2n}-\frac{9}{8n^2}+O(n^{-3}).
\]
因此
\[
n\left[n\left(\sqrt{1+3/n}-1\right)-\frac{3}{2}\right]\longrightarrow-\frac{9}{8}.
\]
第一个极限只需要一阶项，第二个极限消去该项后必须保留二阶项；精度要求由分母或缩放因子的阶数决定。
\end{solve}

\begin{solve}[习题~\ref{ex:c5q10}：缩放后的根式增量]\phantomsection\label{ans:c5q10}
第一步有理化得到 $n(\sqrt{1+4/n}-1)=4/(\sqrt{1+4/n}+1)\to 2$。为求下一阶项，使用
\[
\sqrt{1+4/n}=1+\frac{4}{2n}-\frac{16}{8n^2}+O(n^{-3}).
\]
因此
\[
n\left[n\left(\sqrt{1+4/n}-1\right)-2\right]\longrightarrow-2.
\]
第一个极限只需要一阶项，第二个极限消去该项后必须保留二阶项；精度要求由分母或缩放因子的阶数决定。
\end{solve}

\begin{solve}[习题~\ref{ex:c5q11}：缩放后的根式增量]\phantomsection\label{ans:c5q11}
第一步有理化得到 $n(\sqrt{1+5/n}-1)=5/(\sqrt{1+5/n}+1)\to \frac{5}{2}$。为求下一阶项，使用
\[
\sqrt{1+5/n}=1+\frac{5}{2n}-\frac{25}{8n^2}+O(n^{-3}).
\]
因此
\[
n\left[n\left(\sqrt{1+5/n}-1\right)-\frac{5}{2}\right]\longrightarrow-\frac{25}{8}.
\]
第一个极限只需要一阶项，第二个极限消去该项后必须保留二阶项；精度要求由分母或缩放因子的阶数决定。
\end{solve}

\begin{solve}[习题~\ref{ex:c5q12}：缩放后的根式增量]\phantomsection\label{ans:c5q12}
第一步有理化得到 $n(\sqrt{1+6/n}-1)=6/(\sqrt{1+6/n}+1)\to 3$。为求下一阶项，使用
\[
\sqrt{1+6/n}=1+\frac{6}{2n}-\frac{36}{8n^2}+O(n^{-3}).
\]
因此
\[
n\left[n\left(\sqrt{1+6/n}-1\right)-3\right]\longrightarrow-\frac{9}{2}.
\]
第一个极限只需要一阶项，第二个极限消去该项后必须保留二阶项；精度要求由分母或缩放因子的阶数决定。
\end{solve}

\begin{solve}[习题~\ref{ex:c5q13}：三角函数差的三阶项]\phantomsection\label{ans:c5q13}
两个函数的一阶项相同，会相互消去。标准展开为 $\tan t=t+t^3/3+o(t^3)$ 与 $\sin t=t-t^3/6+o(t^3)$，所以
\[
\tan(1x)-\sin(1x)=\frac{1}{2}x^3+o(x^3).
\]
所求极限为 $\frac{1}{2}$。也可先写成 $\sin(1x)(1-\cos(1x))/\cos(1x)$，分别使用一阶与二阶标准极限，得到相同结果。这一代数变形避免了三次反复使用洛必达法则。
\end{solve}

\begin{solve}[习题~\ref{ex:c5q14}：三角函数差的三阶项]\phantomsection\label{ans:c5q14}
两个函数的一阶项相同，会相互消去。标准展开为 $\tan t=t+t^3/3+o(t^3)$ 与 $\sin t=t-t^3/6+o(t^3)$，所以
\[
\tan(2x)-\sin(2x)=4x^3+o(x^3).
\]
所求极限为 $4$。也可先写成 $\sin(2x)(1-\cos(2x))/\cos(2x)$，分别使用一阶与二阶标准极限，得到相同结果。这一代数变形避免了三次反复使用洛必达法则。
\end{solve}

\begin{solve}[习题~\ref{ex:c5q15}：三角函数差的三阶项]\phantomsection\label{ans:c5q15}
两个函数的一阶项相同，会相互消去。标准展开为 $\tan t=t+t^3/3+o(t^3)$ 与 $\sin t=t-t^3/6+o(t^3)$，所以
\[
\tan(3x)-\sin(3x)=\frac{27}{2}x^3+o(x^3).
\]
所求极限为 $\frac{27}{2}$。也可先写成 $\sin(3x)(1-\cos(3x))/\cos(3x)$，分别使用一阶与二阶标准极限，得到相同结果。这一代数变形避免了三次反复使用洛必达法则。
\end{solve}

\begin{solve}[习题~\ref{ex:c5q16}：三角函数差的三阶项]\phantomsection\label{ans:c5q16}
两个函数的一阶项相同，会相互消去。标准展开为 $\tan t=t+t^3/3+o(t^3)$ 与 $\sin t=t-t^3/6+o(t^3)$，所以
\[
\tan(4x)-\sin(4x)=32x^3+o(x^3).
\]
所求极限为 $32$。也可先写成 $\sin(4x)(1-\cos(4x))/\cos(4x)$，分别使用一阶与二阶标准极限，得到相同结果。这一代数变形避免了三次反复使用洛必达法则。
\end{solve}

\begin{solve}[习题~\ref{ex:c5q17}：三角函数差的三阶项]\phantomsection\label{ans:c5q17}
两个函数的一阶项相同，会相互消去。标准展开为 $\tan t=t+t^3/3+o(t^3)$ 与 $\sin t=t-t^3/6+o(t^3)$，所以
\[
\tan(5x)-\sin(5x)=\frac{125}{2}x^3+o(x^3).
\]
所求极限为 $\frac{125}{2}$。也可先写成 $\sin(5x)(1-\cos(5x))/\cos(5x)$，分别使用一阶与二阶标准极限，得到相同结果。这一代数变形避免了三次反复使用洛必达法则。
\end{solve}

\begin{solve}[习题~\ref{ex:c5q18}：三角函数差的三阶项]\phantomsection\label{ans:c5q18}
两个函数的一阶项相同，会相互消去。标准展开为 $\tan t=t+t^3/3+o(t^3)$ 与 $\sin t=t-t^3/6+o(t^3)$，所以
\[
\tan(6x)-\sin(6x)=108x^3+o(x^3).
\]
所求极限为 $108$。也可先写成 $\sin(6x)(1-\cos(6x))/\cos(6x)$，分别使用一阶与二阶标准极限，得到相同结果。这一代数变形避免了三次反复使用洛必达法则。
\end{solve}

\begin{solve}[习题~\ref{ex:c5q19}：对数减去线性项]\phantomsection\label{ans:c5q19}
由 $\log(1+t)=t-t^2/2+o(t^2)$，有
\[
\log(1+1x)-1x=-\frac{1}{2}x^2+o(x^2).
\]
所以极限为 $-\frac{1}{2}$。实际上，当 $1+1x>0$ 时，由 $\log u\le u-1$ 取 $u=1+1x$，可知分子始终不大于0，且仅在 $x=0$ 取等。局部展开给出精确阶数，整体不等式则提供符号检查。
\end{solve}

\begin{solve}[习题~\ref{ex:c5q20}：对数减去线性项]\phantomsection\label{ans:c5q20}
由 $\log(1+t)=t-t^2/2+o(t^2)$，有
\[
\log(1+2x)-2x=-2x^2+o(x^2).
\]
所以极限为 $-2$。实际上，当 $1+2x>0$ 时，由 $\log u\le u-1$ 取 $u=1+2x$，可知分子始终不大于0，且仅在 $x=0$ 取等。局部展开给出精确阶数，整体不等式则提供符号检查。
\end{solve}

\begin{solve}[习题~\ref{ex:c5q21}：对数减去线性项]\phantomsection\label{ans:c5q21}
由 $\log(1+t)=t-t^2/2+o(t^2)$，有
\[
\log(1+3x)-3x=-\frac{9}{2}x^2+o(x^2).
\]
所以极限为 $-\frac{9}{2}$。实际上，当 $1+3x>0$ 时，由 $\log u\le u-1$ 取 $u=1+3x$，可知分子始终不大于0，且仅在 $x=0$ 取等。局部展开给出精确阶数，整体不等式则提供符号检查。
\end{solve}

\begin{solve}[习题~\ref{ex:c5q22}：对数减去线性项]\phantomsection\label{ans:c5q22}
由 $\log(1+t)=t-t^2/2+o(t^2)$，有
\[
\log(1+4x)-4x=-8x^2+o(x^2).
\]
所以极限为 $-8$。实际上，当 $1+4x>0$ 时，由 $\log u\le u-1$ 取 $u=1+4x$，可知分子始终不大于0，且仅在 $x=0$ 取等。局部展开给出精确阶数，整体不等式则提供符号检查。
\end{solve}

\begin{solve}[习题~\ref{ex:c5q23}：对数减去线性项]\phantomsection\label{ans:c5q23}
由 $\log(1+t)=t-t^2/2+o(t^2)$，有
\[
\log(1+5x)-5x=-\frac{25}{2}x^2+o(x^2).
\]
所以极限为 $-\frac{25}{2}$。实际上，当 $1+5x>0$ 时，由 $\log u\le u-1$ 取 $u=1+5x$，可知分子始终不大于0，且仅在 $x=0$ 取等。局部展开给出精确阶数，整体不等式则提供符号检查。
\end{solve}

\begin{solve}[习题~\ref{ex:c5q24}：对数减去线性项]\phantomsection\label{ans:c5q24}
由 $\log(1+t)=t-t^2/2+o(t^2)$，有
\[
\log(1+6x)-6x=-18x^2+o(x^2).
\]
所以极限为 $-18$。实际上，当 $1+6x>0$ 时，由 $\log u\le u-1$ 取 $u=1+6x$，可知分子始终不大于0，且仅在 $x=0$ 取等。局部展开给出精确阶数，整体不等式则提供符号检查。
\end{solve}

\clearpage
\section{连续性、零点与定义检验}\label{sec:ch6}
连续（continuity）要求极限等于函数在该点的值；一致连续（uniform continuity）则要求同一个距离尺度能对定义域内所有点生效。零点存在性和唯一性通常由不同的理由保证。

\textbf{资源对应：}CLP-1 \S~1.6、\S~2.13；补充证明训练。本章共 24 道编号习题。
\subsection{习题}
\begin{example}[可去间断点的补定义]\label{ex:c6q1}
在 $x\ne0$ 时设 $f(x)=\sin(1x)/x$，在 $x=0$ 时设 $f(0)=c$。求使 $f$ 连续的 $c$，并求连续补定义后 $f^{\prime}(0)$。
\par\textit{改编训练。}\par\hyperref[ans:c6q1]{解答见第~\pageref{ans:c6q1} 页。}
\end{example}
\begin{example}[可去间断点的补定义]\label{ex:c6q2}
在 $x\ne0$ 时设 $f(x)=\sin(2x)/x$，在 $x=0$ 时设 $f(0)=c$。求使 $f$ 连续的 $c$，并求连续补定义后 $f^{\prime}(0)$。
\par\textit{改编训练。}\par\hyperref[ans:c6q2]{解答见第~\pageref{ans:c6q2} 页。}
\end{example}
\begin{example}[可去间断点的补定义]\label{ex:c6q3}
在 $x\ne0$ 时设 $f(x)=\sin(3x)/x$，在 $x=0$ 时设 $f(0)=c$。求使 $f$ 连续的 $c$，并求连续补定义后 $f^{\prime}(0)$。
\par\textit{改编训练。}\par\hyperref[ans:c6q3]{解答见第~\pageref{ans:c6q3} 页。}
\end{example}
\begin{example}[可去间断点的补定义]\label{ex:c6q4}
在 $x\ne0$ 时设 $f(x)=\sin(4x)/x$，在 $x=0$ 时设 $f(0)=c$。求使 $f$ 连续的 $c$，并求连续补定义后 $f^{\prime}(0)$。
\par\textit{改编训练。}\par\hyperref[ans:c6q4]{解答见第~\pageref{ans:c6q4} 页。}
\end{example}
\begin{example}[可去间断点的补定义]\label{ex:c6q5}
在 $x\ne0$ 时设 $f(x)=\sin(5x)/x$，在 $x=0$ 时设 $f(0)=c$。求使 $f$ 连续的 $c$，并求连续补定义后 $f^{\prime}(0)$。
\par\textit{改编训练。}\par\hyperref[ans:c6q5]{解答见第~\pageref{ans:c6q5} 页。}
\end{example}
\begin{example}[可去间断点的补定义]\label{ex:c6q6}
在 $x\ne0$ 时设 $f(x)=\sin(6x)/x$，在 $x=0$ 时设 $f(0)=c$。求使 $f$ 连续的 $c$，并求连续补定义后 $f^{\prime}(0)$。
\par\textit{改编训练。}\par\hyperref[ans:c6q6]{解答见第~\pageref{ans:c6q6} 页。}
\end{example}
\begin{example}[介值定理与单调性]\label{ex:c6q7}
证明方程 $x^3+1x-1=0$ 在 $(0,1)$ 中恰有一个实根，并说明如何用二分法缩小根的区间。
\par\textit{改编训练。}\par\hyperref[ans:c6q7]{解答见第~\pageref{ans:c6q7} 页。}
\end{example}
\begin{example}[介值定理与单调性]\label{ex:c6q8}
证明方程 $x^3+2x-1=0$ 在 $(0,1)$ 中恰有一个实根，并说明如何用二分法缩小根的区间。
\par\textit{改编训练。}\par\hyperref[ans:c6q8]{解答见第~\pageref{ans:c6q8} 页。}
\end{example}
\begin{example}[介值定理与单调性]\label{ex:c6q9}
证明方程 $x^3+3x-1=0$ 在 $(0,1)$ 中恰有一个实根，并说明如何用二分法缩小根的区间。
\par\textit{改编训练。}\par\hyperref[ans:c6q9]{解答见第~\pageref{ans:c6q9} 页。}
\end{example}
\begin{example}[介值定理与单调性]\label{ex:c6q10}
证明方程 $x^3+4x-1=0$ 在 $(0,1)$ 中恰有一个实根，并说明如何用二分法缩小根的区间。
\par\textit{改编训练。}\par\hyperref[ans:c6q10]{解答见第~\pageref{ans:c6q10} 页。}
\end{example}
\begin{example}[介值定理与单调性]\label{ex:c6q11}
证明方程 $x^3+5x-1=0$ 在 $(0,1)$ 中恰有一个实根，并说明如何用二分法缩小根的区间。
\par\textit{改编训练。}\par\hyperref[ans:c6q11]{解答见第~\pageref{ans:c6q11} 页。}
\end{example}
\begin{example}[介值定理与单调性]\label{ex:c6q12}
证明方程 $x^3+6x-1=0$ 在 $(0,1)$ 中恰有一个实根，并说明如何用二分法缩小根的区间。
\par\textit{改编训练。}\par\hyperref[ans:c6q12]{解答见第~\pageref{ans:c6q12} 页。}
\end{example}
\begin{example}[一致连续与无界导数]\label{ex:c6q13}
证明 $f(x)=1\sqrt x$ 在 $[0,\infty)$ 上一致连续。说明导数在 $0$ 附近无界为什么不构成反例。
\par\textit{改编训练。}\par\hyperref[ans:c6q13]{解答见第~\pageref{ans:c6q13} 页。}
\end{example}
\begin{example}[一致连续与无界导数]\label{ex:c6q14}
证明 $f(x)=2\sqrt x$ 在 $[0,\infty)$ 上一致连续。说明导数在 $0$ 附近无界为什么不构成反例。
\par\textit{改编训练。}\par\hyperref[ans:c6q14]{解答见第~\pageref{ans:c6q14} 页。}
\end{example}
\begin{example}[一致连续与无界导数]\label{ex:c6q15}
证明 $f(x)=3\sqrt x$ 在 $[0,\infty)$ 上一致连续。说明导数在 $0$ 附近无界为什么不构成反例。
\par\textit{改编训练。}\par\hyperref[ans:c6q15]{解答见第~\pageref{ans:c6q15} 页。}
\end{example}
\begin{example}[一致连续与无界导数]\label{ex:c6q16}
证明 $f(x)=4\sqrt x$ 在 $[0,\infty)$ 上一致连续。说明导数在 $0$ 附近无界为什么不构成反例。
\par\textit{改编训练。}\par\hyperref[ans:c6q16]{解答见第~\pageref{ans:c6q16} 页。}
\end{example}
\begin{example}[一致连续与无界导数]\label{ex:c6q17}
证明 $f(x)=5\sqrt x$ 在 $[0,\infty)$ 上一致连续。说明导数在 $0$ 附近无界为什么不构成反例。
\par\textit{改编训练。}\par\hyperref[ans:c6q17]{解答见第~\pageref{ans:c6q17} 页。}
\end{example}
\begin{example}[一致连续与无界导数]\label{ex:c6q18}
证明 $f(x)=6\sqrt x$ 在 $[0,\infty)$ 上一致连续。说明导数在 $0$ 附近无界为什么不构成反例。
\par\textit{改编训练。}\par\hyperref[ans:c6q18]{解答见第~\pageref{ans:c6q18} 页。}
\end{example}
\begin{example}[用点列否定一致连续]\label{ex:c6q19}
证明 $f(x)=1x^2$ 在 $[0,\infty)$ 上不一致连续，但在每个有界闭区间 $[0,M]$ 上一致连续。
\par\textit{改编训练。}\par\hyperref[ans:c6q19]{解答见第~\pageref{ans:c6q19} 页。}
\end{example}
\begin{example}[用点列否定一致连续]\label{ex:c6q20}
证明 $f(x)=2x^2$ 在 $[0,\infty)$ 上不一致连续，但在每个有界闭区间 $[0,M]$ 上一致连续。
\par\textit{改编训练。}\par\hyperref[ans:c6q20]{解答见第~\pageref{ans:c6q20} 页。}
\end{example}
\begin{example}[用点列否定一致连续]\label{ex:c6q21}
证明 $f(x)=3x^2$ 在 $[0,\infty)$ 上不一致连续，但在每个有界闭区间 $[0,M]$ 上一致连续。
\par\textit{改编训练。}\par\hyperref[ans:c6q21]{解答见第~\pageref{ans:c6q21} 页。}
\end{example}
\begin{example}[用点列否定一致连续]\label{ex:c6q22}
证明 $f(x)=4x^2$ 在 $[0,\infty)$ 上不一致连续，但在每个有界闭区间 $[0,M]$ 上一致连续。
\par\textit{改编训练。}\par\hyperref[ans:c6q22]{解答见第~\pageref{ans:c6q22} 页。}
\end{example}
\begin{example}[用点列否定一致连续]\label{ex:c6q23}
证明 $f(x)=5x^2$ 在 $[0,\infty)$ 上不一致连续，但在每个有界闭区间 $[0,M]$ 上一致连续。
\par\textit{改编训练。}\par\hyperref[ans:c6q23]{解答见第~\pageref{ans:c6q23} 页。}
\end{example}
\begin{example}[用点列否定一致连续]\label{ex:c6q24}
证明 $f(x)=6x^2$ 在 $[0,\infty)$ 上不一致连续，但在每个有界闭区间 $[0,M]$ 上一致连续。
\par\textit{改编训练。}\par\hyperref[ans:c6q24]{解答见第~\pageref{ans:c6q24} 页。}
\end{example}
\subsection{解答}
\begin{solve}[习题~\ref{ex:c6q1}：可去间断点的补定义]\phantomsection\label{ans:c6q1}
先求函数极限：
\[
\lim_{x\to0}\frac{\sin(1x)}x=1\lim_{x\to0}\frac{\sin(1x)}{1x}=1.
\]
连续的充要条件为 $c=1$。此时用导数定义而不是在分式上直接代入 $0$：
\[
f^{\prime}(0)=\lim_{h\to0}\frac{\sin(1h)/h-1}h=\lim_{h\to0}\frac{\sin(1h)-1h}{h^2}=0.
\]
最后一步可以由三阶展开得到。连续补定义不只是指定任意一个数；只有与两侧共同极限一致的值才有效。
\end{solve}

\begin{solve}[习题~\ref{ex:c6q2}：可去间断点的补定义]\phantomsection\label{ans:c6q2}
先求函数极限：
\[
\lim_{x\to0}\frac{\sin(2x)}x=2\lim_{x\to0}\frac{\sin(2x)}{2x}=2.
\]
连续的充要条件为 $c=2$。此时用导数定义而不是在分式上直接代入 $0$：
\[
f^{\prime}(0)=\lim_{h\to0}\frac{\sin(2h)/h-2}h=\lim_{h\to0}\frac{\sin(2h)-2h}{h^2}=0.
\]
最后一步可以由三阶展开得到。连续补定义不只是指定任意一个数；只有与两侧共同极限一致的值才有效。
\end{solve}

\begin{solve}[习题~\ref{ex:c6q3}：可去间断点的补定义]\phantomsection\label{ans:c6q3}
先求函数极限：
\[
\lim_{x\to0}\frac{\sin(3x)}x=3\lim_{x\to0}\frac{\sin(3x)}{3x}=3.
\]
连续的充要条件为 $c=3$。此时用导数定义而不是在分式上直接代入 $0$：
\[
f^{\prime}(0)=\lim_{h\to0}\frac{\sin(3h)/h-3}h=\lim_{h\to0}\frac{\sin(3h)-3h}{h^2}=0.
\]
最后一步可以由三阶展开得到。连续补定义不只是指定任意一个数；只有与两侧共同极限一致的值才有效。
\end{solve}

\begin{solve}[习题~\ref{ex:c6q4}：可去间断点的补定义]\phantomsection\label{ans:c6q4}
先求函数极限：
\[
\lim_{x\to0}\frac{\sin(4x)}x=4\lim_{x\to0}\frac{\sin(4x)}{4x}=4.
\]
连续的充要条件为 $c=4$。此时用导数定义而不是在分式上直接代入 $0$：
\[
f^{\prime}(0)=\lim_{h\to0}\frac{\sin(4h)/h-4}h=\lim_{h\to0}\frac{\sin(4h)-4h}{h^2}=0.
\]
最后一步可以由三阶展开得到。连续补定义不只是指定任意一个数；只有与两侧共同极限一致的值才有效。
\end{solve}

\begin{solve}[习题~\ref{ex:c6q5}：可去间断点的补定义]\phantomsection\label{ans:c6q5}
先求函数极限：
\[
\lim_{x\to0}\frac{\sin(5x)}x=5\lim_{x\to0}\frac{\sin(5x)}{5x}=5.
\]
连续的充要条件为 $c=5$。此时用导数定义而不是在分式上直接代入 $0$：
\[
f^{\prime}(0)=\lim_{h\to0}\frac{\sin(5h)/h-5}h=\lim_{h\to0}\frac{\sin(5h)-5h}{h^2}=0.
\]
最后一步可以由三阶展开得到。连续补定义不只是指定任意一个数；只有与两侧共同极限一致的值才有效。
\end{solve}

\begin{solve}[习题~\ref{ex:c6q6}：可去间断点的补定义]\phantomsection\label{ans:c6q6}
先求函数极限：
\[
\lim_{x\to0}\frac{\sin(6x)}x=6\lim_{x\to0}\frac{\sin(6x)}{6x}=6.
\]
连续的充要条件为 $c=6$。此时用导数定义而不是在分式上直接代入 $0$：
\[
f^{\prime}(0)=\lim_{h\to0}\frac{\sin(6h)/h-6}h=\lim_{h\to0}\frac{\sin(6h)-6h}{h^2}=0.
\]
最后一步可以由三阶展开得到。连续补定义不只是指定任意一个数；只有与两侧共同极限一致的值才有效。
\end{solve}

\begin{solve}[习题~\ref{ex:c6q7}：介值定理与单调性]\phantomsection\label{ans:c6q7}
设 $F(x)=x^3+1x-1$。它在 $[0,1]$ 上连续，且 $F(0)=-1<0$、$F(1)=1>0$，故介值定理给出一个零点。另一方面，
\[
F^{\prime}(x)=3x^2+1>0\qquad(x\in\mathbb R),
\]
所以函数严格递增，零点至多一个。二分法从 $[a_0,b_0]=[0,1]$ 开始，每次计算中点 $c_k$；若函数值非零，就保留两端函数值异号的半区间。第 $k$ 次后区间长为 $2^{-k}$，其中点与根的距离不超过 $2^{-k-1}$。存在性来自连续性，唯一性来自严格单调性，误差估计来自区间长度。
\end{solve}

\begin{solve}[习题~\ref{ex:c6q8}：介值定理与单调性]\phantomsection\label{ans:c6q8}
设 $F(x)=x^3+2x-1$。它在 $[0,1]$ 上连续，且 $F(0)=-1<0$、$F(1)=2>0$，故介值定理给出一个零点。另一方面，
\[
F^{\prime}(x)=3x^2+2>0\qquad(x\in\mathbb R),
\]
所以函数严格递增，零点至多一个。二分法从 $[a_0,b_0]=[0,1]$ 开始，每次计算中点 $c_k$；若函数值非零，就保留两端函数值异号的半区间。第 $k$ 次后区间长为 $2^{-k}$，其中点与根的距离不超过 $2^{-k-1}$。存在性来自连续性，唯一性来自严格单调性，误差估计来自区间长度。
\end{solve}

\begin{solve}[习题~\ref{ex:c6q9}：介值定理与单调性]\phantomsection\label{ans:c6q9}
设 $F(x)=x^3+3x-1$。它在 $[0,1]$ 上连续，且 $F(0)=-1<0$、$F(1)=3>0$，故介值定理给出一个零点。另一方面，
\[
F^{\prime}(x)=3x^2+3>0\qquad(x\in\mathbb R),
\]
所以函数严格递增，零点至多一个。二分法从 $[a_0,b_0]=[0,1]$ 开始，每次计算中点 $c_k$；若函数值非零，就保留两端函数值异号的半区间。第 $k$ 次后区间长为 $2^{-k}$，其中点与根的距离不超过 $2^{-k-1}$。存在性来自连续性，唯一性来自严格单调性，误差估计来自区间长度。
\end{solve}

\begin{solve}[习题~\ref{ex:c6q10}：介值定理与单调性]\phantomsection\label{ans:c6q10}
设 $F(x)=x^3+4x-1$。它在 $[0,1]$ 上连续，且 $F(0)=-1<0$、$F(1)=4>0$，故介值定理给出一个零点。另一方面，
\[
F^{\prime}(x)=3x^2+4>0\qquad(x\in\mathbb R),
\]
所以函数严格递增，零点至多一个。二分法从 $[a_0,b_0]=[0,1]$ 开始，每次计算中点 $c_k$；若函数值非零，就保留两端函数值异号的半区间。第 $k$ 次后区间长为 $2^{-k}$，其中点与根的距离不超过 $2^{-k-1}$。存在性来自连续性，唯一性来自严格单调性，误差估计来自区间长度。
\end{solve}

\begin{solve}[习题~\ref{ex:c6q11}：介值定理与单调性]\phantomsection\label{ans:c6q11}
设 $F(x)=x^3+5x-1$。它在 $[0,1]$ 上连续，且 $F(0)=-1<0$、$F(1)=5>0$，故介值定理给出一个零点。另一方面，
\[
F^{\prime}(x)=3x^2+5>0\qquad(x\in\mathbb R),
\]
所以函数严格递增，零点至多一个。二分法从 $[a_0,b_0]=[0,1]$ 开始，每次计算中点 $c_k$；若函数值非零，就保留两端函数值异号的半区间。第 $k$ 次后区间长为 $2^{-k}$，其中点与根的距离不超过 $2^{-k-1}$。存在性来自连续性，唯一性来自严格单调性，误差估计来自区间长度。
\end{solve}

\begin{solve}[习题~\ref{ex:c6q12}：介值定理与单调性]\phantomsection\label{ans:c6q12}
设 $F(x)=x^3+6x-1$。它在 $[0,1]$ 上连续，且 $F(0)=-1<0$、$F(1)=6>0$，故介值定理给出一个零点。另一方面，
\[
F^{\prime}(x)=3x^2+6>0\qquad(x\in\mathbb R),
\]
所以函数严格递增，零点至多一个。二分法从 $[a_0,b_0]=[0,1]$ 开始，每次计算中点 $c_k$；若函数值非零，就保留两端函数值异号的半区间。第 $k$ 次后区间长为 $2^{-k}$，其中点与根的距离不超过 $2^{-k-1}$。存在性来自连续性，唯一性来自严格单调性，误差估计来自区间长度。
\end{solve}

\begin{solve}[习题~\ref{ex:c6q13}：一致连续与无界导数]\phantomsection\label{ans:c6q13}
对于任意 $x,y\ge0$，不妨设 $x\ge y$。因为 $\sqrt x+\sqrt y\ge\sqrt{x-y}$，有
\[
|1\sqrt x-1\sqrt y|=1\frac{|x-y|}{\sqrt x+\sqrt y}\le 1\sqrt{|x-y|}.
\]
给定 $\varepsilon>0$，取 $\delta=(\varepsilon/1)^2$。只要 $|x-y|<\delta$，就有 $|f(x)-f(y)|<\varepsilon$；$\delta$ 与点的位置无关，因此满足一致连续定义。对 $x>0$ 有 $f^{\prime}(x)=1/(2\sqrt x)$，导数确实无界。导数有界是保证一致连续的一个充分条件，并不是必要条件。
\end{solve}

\begin{solve}[习题~\ref{ex:c6q14}：一致连续与无界导数]\phantomsection\label{ans:c6q14}
对于任意 $x,y\ge0$，不妨设 $x\ge y$。因为 $\sqrt x+\sqrt y\ge\sqrt{x-y}$，有
\[
|2\sqrt x-2\sqrt y|=2\frac{|x-y|}{\sqrt x+\sqrt y}\le 2\sqrt{|x-y|}.
\]
给定 $\varepsilon>0$，取 $\delta=(\varepsilon/2)^2$。只要 $|x-y|<\delta$，就有 $|f(x)-f(y)|<\varepsilon$；$\delta$ 与点的位置无关，因此满足一致连续定义。对 $x>0$ 有 $f^{\prime}(x)=2/(2\sqrt x)$，导数确实无界。导数有界是保证一致连续的一个充分条件，并不是必要条件。
\end{solve}

\begin{solve}[习题~\ref{ex:c6q15}：一致连续与无界导数]\phantomsection\label{ans:c6q15}
对于任意 $x,y\ge0$，不妨设 $x\ge y$。因为 $\sqrt x+\sqrt y\ge\sqrt{x-y}$，有
\[
|3\sqrt x-3\sqrt y|=3\frac{|x-y|}{\sqrt x+\sqrt y}\le 3\sqrt{|x-y|}.
\]
给定 $\varepsilon>0$，取 $\delta=(\varepsilon/3)^2$。只要 $|x-y|<\delta$，就有 $|f(x)-f(y)|<\varepsilon$；$\delta$ 与点的位置无关，因此满足一致连续定义。对 $x>0$ 有 $f^{\prime}(x)=3/(2\sqrt x)$，导数确实无界。导数有界是保证一致连续的一个充分条件，并不是必要条件。
\end{solve}

\begin{solve}[习题~\ref{ex:c6q16}：一致连续与无界导数]\phantomsection\label{ans:c6q16}
对于任意 $x,y\ge0$，不妨设 $x\ge y$。因为 $\sqrt x+\sqrt y\ge\sqrt{x-y}$，有
\[
|4\sqrt x-4\sqrt y|=4\frac{|x-y|}{\sqrt x+\sqrt y}\le 4\sqrt{|x-y|}.
\]
给定 $\varepsilon>0$，取 $\delta=(\varepsilon/4)^2$。只要 $|x-y|<\delta$，就有 $|f(x)-f(y)|<\varepsilon$；$\delta$ 与点的位置无关，因此满足一致连续定义。对 $x>0$ 有 $f^{\prime}(x)=4/(2\sqrt x)$，导数确实无界。导数有界是保证一致连续的一个充分条件，并不是必要条件。
\end{solve}

\begin{solve}[习题~\ref{ex:c6q17}：一致连续与无界导数]\phantomsection\label{ans:c6q17}
对于任意 $x,y\ge0$，不妨设 $x\ge y$。因为 $\sqrt x+\sqrt y\ge\sqrt{x-y}$，有
\[
|5\sqrt x-5\sqrt y|=5\frac{|x-y|}{\sqrt x+\sqrt y}\le 5\sqrt{|x-y|}.
\]
给定 $\varepsilon>0$，取 $\delta=(\varepsilon/5)^2$。只要 $|x-y|<\delta$，就有 $|f(x)-f(y)|<\varepsilon$；$\delta$ 与点的位置无关，因此满足一致连续定义。对 $x>0$ 有 $f^{\prime}(x)=5/(2\sqrt x)$，导数确实无界。导数有界是保证一致连续的一个充分条件，并不是必要条件。
\end{solve}

\begin{solve}[习题~\ref{ex:c6q18}：一致连续与无界导数]\phantomsection\label{ans:c6q18}
对于任意 $x,y\ge0$，不妨设 $x\ge y$。因为 $\sqrt x+\sqrt y\ge\sqrt{x-y}$，有
\[
|6\sqrt x-6\sqrt y|=6\frac{|x-y|}{\sqrt x+\sqrt y}\le 6\sqrt{|x-y|}.
\]
给定 $\varepsilon>0$，取 $\delta=(\varepsilon/6)^2$。只要 $|x-y|<\delta$，就有 $|f(x)-f(y)|<\varepsilon$；$\delta$ 与点的位置无关，因此满足一致连续定义。对 $x>0$ 有 $f^{\prime}(x)=6/(2\sqrt x)$，导数确实无界。导数有界是保证一致连续的一个充分条件，并不是必要条件。
\end{solve}

\begin{solve}[习题~\ref{ex:c6q19}：用点列否定一致连续]\phantomsection\label{ans:c6q19}
取 $x_n=n$、$y_n=n+1/n$，则 $|x_n-y_n|=1/n\to0$，而
\[
|f(y_n)-f(x_n)|=1\left(2+\frac1{n^2}\right)\longrightarrow 2>0.
\]
若函数一致连续，输入距离趋零应迫使输出距离趋零，出现矛盾。具体可取 $\varepsilon=1$ 来否定定义。在 $[0,M]$ 上，当 $M>0$ 时，
\[
|f(x)-f(y)|=1|x+y||x-y|\le 2M|x-y|.
\]
取 $\delta=\varepsilon/(2M)$ 即可；$M=0$ 时结论平凡成立。定义域是否有界会实质性改变结论。
\end{solve}

\begin{solve}[习题~\ref{ex:c6q20}：用点列否定一致连续]\phantomsection\label{ans:c6q20}
取 $x_n=n$、$y_n=n+1/n$，则 $|x_n-y_n|=1/n\to0$，而
\[
|f(y_n)-f(x_n)|=2\left(2+\frac1{n^2}\right)\longrightarrow 4>0.
\]
若函数一致连续，输入距离趋零应迫使输出距离趋零，出现矛盾。具体可取 $\varepsilon=2$ 来否定定义。在 $[0,M]$ 上，当 $M>0$ 时，
\[
|f(x)-f(y)|=2|x+y||x-y|\le 4M|x-y|.
\]
取 $\delta=\varepsilon/(4M)$ 即可；$M=0$ 时结论平凡成立。定义域是否有界会实质性改变结论。
\end{solve}

\begin{solve}[习题~\ref{ex:c6q21}：用点列否定一致连续]\phantomsection\label{ans:c6q21}
取 $x_n=n$、$y_n=n+1/n$，则 $|x_n-y_n|=1/n\to0$，而
\[
|f(y_n)-f(x_n)|=3\left(2+\frac1{n^2}\right)\longrightarrow 6>0.
\]
若函数一致连续，输入距离趋零应迫使输出距离趋零，出现矛盾。具体可取 $\varepsilon=3$ 来否定定义。在 $[0,M]$ 上，当 $M>0$ 时，
\[
|f(x)-f(y)|=3|x+y||x-y|\le 6M|x-y|.
\]
取 $\delta=\varepsilon/(6M)$ 即可；$M=0$ 时结论平凡成立。定义域是否有界会实质性改变结论。
\end{solve}

\begin{solve}[习题~\ref{ex:c6q22}：用点列否定一致连续]\phantomsection\label{ans:c6q22}
取 $x_n=n$、$y_n=n+1/n$，则 $|x_n-y_n|=1/n\to0$，而
\[
|f(y_n)-f(x_n)|=4\left(2+\frac1{n^2}\right)\longrightarrow 8>0.
\]
若函数一致连续，输入距离趋零应迫使输出距离趋零，出现矛盾。具体可取 $\varepsilon=4$ 来否定定义。在 $[0,M]$ 上，当 $M>0$ 时，
\[
|f(x)-f(y)|=4|x+y||x-y|\le 8M|x-y|.
\]
取 $\delta=\varepsilon/(8M)$ 即可；$M=0$ 时结论平凡成立。定义域是否有界会实质性改变结论。
\end{solve}

\begin{solve}[习题~\ref{ex:c6q23}：用点列否定一致连续]\phantomsection\label{ans:c6q23}
取 $x_n=n$、$y_n=n+1/n$，则 $|x_n-y_n|=1/n\to0$，而
\[
|f(y_n)-f(x_n)|=5\left(2+\frac1{n^2}\right)\longrightarrow 10>0.
\]
若函数一致连续，输入距离趋零应迫使输出距离趋零，出现矛盾。具体可取 $\varepsilon=5$ 来否定定义。在 $[0,M]$ 上，当 $M>0$ 时，
\[
|f(x)-f(y)|=5|x+y||x-y|\le 10M|x-y|.
\]
取 $\delta=\varepsilon/(10M)$ 即可；$M=0$ 时结论平凡成立。定义域是否有界会实质性改变结论。
\end{solve}

\begin{solve}[习题~\ref{ex:c6q24}：用点列否定一致连续]\phantomsection\label{ans:c6q24}
取 $x_n=n$、$y_n=n+1/n$，则 $|x_n-y_n|=1/n\to0$，而
\[
|f(y_n)-f(x_n)|=6\left(2+\frac1{n^2}\right)\longrightarrow 12>0.
\]
若函数一致连续，输入距离趋零应迫使输出距离趋零，出现矛盾。具体可取 $\varepsilon=6$ 来否定定义。在 $[0,M]$ 上，当 $M>0$ 时，
\[
|f(x)-f(y)|=6|x+y||x-y|\le 12M|x-y|.
\]
取 $\delta=\varepsilon/(12M)$ 即可；$M=0$ 时结论平凡成立。定义域是否有界会实质性改变结论。
\end{solve}

\clearpage
\section{连续函数与微分定理的证明专题}\label{sec:ch7}
定理的结论依赖全部假设。以下证明重点说明紧性、连续性、可导性分别在哪一步起作用，并通过反例检验删去条件后是否仍然成立。

\textbf{资源对应：}编者补充证明；与 CLP-1 \S~1.6、\S~2.13 衔接。本章共 20 道编号习题。
\subsection{习题}
\begin{example}[连续函数的序列判别]\label{ex:c7q1}
证明函数 $f$ 在 $a$ 连续，当且仅当对定义域中每个趋于 $a$ 的数列 $x_n$，都有 $f(x_n)\to f(a)$。
\par\textit{补充证明。来源：编者补充：本科数学分析标准命题与反例。}\par\hyperref[ans:c7q1]{解答见第~\pageref{ans:c7q1} 页。}
\end{example}
\begin{example}[闭区间连续函数有界]\label{ex:c7q2}
用数列法证明：连续函数在有界闭区间上有界。
\par\textit{补充证明。来源：编者补充：本科数学分析标准命题与反例。}\par\hyperref[ans:c7q2]{解答见第~\pageref{ans:c7q2} 页。}
\end{example}
\begin{example}[最大值定理的证明]\label{ex:c7q3}
设 $f$ 在 $[a,b]$ 上连续，证明它取得最大值与最小值。
\par\textit{补充证明。来源：编者补充：本科数学分析标准命题与反例。}\par\hyperref[ans:c7q3]{解答见第~\pageref{ans:c7q3} 页。}
\end{example}
\begin{example}[介值定理的上确界证明]\label{ex:c7q4}
设 $f$ 在 $[a,b]$ 连续，$f(a)<0<f(b)$。用上确界法证明存在零点。
\par\textit{补充证明。来源：编者补充：本科数学分析标准命题与反例。}\par\hyperref[ans:c7q4]{解答见第~\pageref{ans:c7q4} 页。}
\end{example}
\begin{example}[闭区间上一致连续]\label{ex:c7q5}
证明闭区间上的连续函数一致连续。
\par\textit{补充证明。来源：编者补充：本科数学分析标准命题与反例。}\par\hyperref[ans:c7q5]{解答见第~\pageref{ans:c7q5} 页。}
\end{example}
\begin{example}[一致连续延拓到端点]\label{ex:c7q6}
若 $f$ 在 $(a,b)$ 上一致连续，证明它在两个有限端点都有有限单侧极限，从而可连续延拓到 $[a,b]$。
\par\textit{补充证明。来源：编者补充：本科数学分析标准命题与反例。}\par\hyperref[ans:c7q6]{解答见第~\pageref{ans:c7q6} 页。}
\end{example}
\begin{example}[可导推出连续]\label{ex:c7q7}
从导数定义证明：在某点可导必在该点连续。
\par\textit{补充证明。来源：编者补充：本科数学分析标准命题与反例。}\par\hyperref[ans:c7q7]{解答见第~\pageref{ans:c7q7} 页。}
\end{example}
\begin{example}[Fermat 必要条件]\label{ex:c7q8}
设 $f$ 在内点 $a$ 可导，并在该点取得局部极大值。证明 $f^{\prime}(a)=0$。
\par\textit{补充证明。来源：编者补充：本科数学分析标准命题与反例。}\par\hyperref[ans:c7q8]{解答见第~\pageref{ans:c7q8} 页。}
\end{example}
\begin{example}[Rolle 定理的证明]\label{ex:c7q9}
证明 Rolle 定理，并说明常函数情形如何处理。
\par\textit{补充证明。来源：编者补充：本科数学分析标准命题与反例。}\par\hyperref[ans:c7q9]{解答见第~\pageref{ans:c7q9} 页。}
\end{example}
\begin{example}[从 Rolle 推出 Lagrange]\label{ex:c7q10}
由 Rolle 定理证明 Lagrange 中值定理。
\par\textit{补充证明。来源：编者补充：本科数学分析标准命题与反例。}\par\hyperref[ans:c7q10]{解答见第~\pageref{ans:c7q10} 页。}
\end{example}
\begin{example}[导数恒零推出常函数]\label{ex:c7q11}
证明若 $f$ 在一个区间内处处可导且导数恒为零，则 $f$ 为常函数。说明定义域为不连通集合时的变化。
\par\textit{补充证明。来源：编者补充：本科数学分析标准命题与反例。}\par\hyperref[ans:c7q11]{解答见第~\pageref{ans:c7q11} 页。}
\end{example}
\begin{example}[严格递增的充分条件并非必要]\label{ex:c7q12}
证明 $f^{\prime}>0$ 蕴含函数严格递增，并举例说明严格递增不一定处处有正导数。
\par\textit{补充证明。来源：编者补充：本科数学分析标准命题与反例。}\par\hyperref[ans:c7q12]{解答见第~\pageref{ans:c7q12} 页。}
\end{example}
\begin{example}[凸函数的切线下界]\label{ex:c7q13}
设 $f$ 二阶可导且 $f^{\prime\prime}\ge0$。证明对任意 $x,a$，$f(x)\ge f(a)+f^{\prime}(a)(x-a)$。
\par\textit{补充证明。来源：编者补充：本科数学分析标准命题与反例。}\par\hyperref[ans:c7q13]{解答见第~\pageref{ans:c7q13} 页。}
\end{example}
\begin{example}[指数与对数的基本不等式]\label{ex:c7q14}
证明 $e^x\ge1+x$ 对所有实数成立，且 $\log x\le x-1$ 对所有 $x>0$ 成立，并说明等号条件。
\par\textit{补充证明。来源：编者补充：本科数学分析标准命题与反例。}\par\hyperref[ans:c7q14]{解答见第~\pageref{ans:c7q14} 页。}
\end{example}
\begin{example}[导函数的介值性]\label{ex:c7q15}
说明如何证明导函数具有 Darboux 介值性质：若 $f$ 在 $[a,b]$ 的邻域可导，且 $f^{\prime}(a)<\lambda<f^{\prime}(b)$，则内部存在 $c$ 使 $f^{\prime}(c)=\lambda$。
\par\textit{补充证明。来源：编者补充：本科数学分析标准命题与反例。}\par\hyperref[ans:c7q15]{解答见第~\pageref{ans:c7q15} 页。}
\end{example}
\begin{example}[导数不连续的可导函数]\label{ex:c7q16}
证明 $f(x)=x^2\sin(1/x)$（$x\ne0$）、$f(0)=0$ 处处可导，但导数在0不连续。
\par\textit{补充证明。来源：编者补充：本科数学分析标准命题与反例。}\par\hyperref[ans:c7q16]{解答见第~\pageref{ans:c7q16} 页。}
\end{example}
\begin{example}[两个原函数之差]\label{ex:c7q17}
若 $F^{\prime}=G^{\prime}=f$ 在一个区间上成立，证明 $F-G$ 是常数。
\par\textit{补充证明。来源：编者补充：本科数学分析标准命题与反例。}\par\hyperref[ans:c7q17]{解答见第~\pageref{ans:c7q17} 页。}
\end{example}
\begin{example}[利用 Rolle 控制零点个数]\label{ex:c7q18}
设 $f$ 在区间上二阶可导且 $f^{\prime\prime}>0$。证明 $f$ 最多有两个不同零点。
\par\textit{补充证明。来源：编者补充：本科数学分析标准命题与反例。}\par\hyperref[ans:c7q18]{解答见第~\pageref{ans:c7q18} 页。}
\end{example}
\begin{example}[Taylor 余项的局部估计]\label{ex:c7q19}
设 $f$ 在 $a$ 附近有连续的三阶导数，且 $|f^{(3)}|\le M$。证明二阶 Taylor 近似的误差不超过 $M|x-a|^3/6$。
\par\textit{补充证明。来源：编者补充：本科数学分析标准命题与反例。}\par\hyperref[ans:c7q19]{解答见第~\pageref{ans:c7q19} 页。}
\end{example}
\begin{example}[逆函数的导数条件]\label{ex:c7q20}
设 $f$ 在 $a$ 附近连续严格单调，在 $a$ 可导且 $f^{\prime}(a)\ne0$。证明逆函数在 $b=f(a)$ 的导数为 $1/f^{\prime}(a)$。
\par\textit{补充证明。来源：编者补充：本科数学分析标准命题与反例。}\par\hyperref[ans:c7q20]{解答见第~\pageref{ans:c7q20} 页。}
\end{example}
\subsection{解答}
\begin{solve}[习题~\ref{ex:c7q1}：连续函数的序列判别]\phantomsection\label{ans:c7q1}
若连续，给定输出误差 $\varepsilon$，连续定义提供输入距离 $\delta$；数列收敛保证尾部输入距离小于 $\delta$，故输出误差小于 $\varepsilon$。反过来，若不连续，则存在 $\varepsilon_0>0$，对每个 $\delta>0$ 都能找到定义域中的 $x$，使 $|x-a|<\delta$ 而 $|f(x)-f(a)|\ge\varepsilon_0$。令 $\delta=1/n$ 逐次选择 $x_n$，则 $x_n\to a$，但输出不趋于 $f(a)$，与序列条件矛盾。
\end{solve}

\begin{solve}[习题~\ref{ex:c7q2}：闭区间连续函数有界]\phantomsection\label{ans:c7q2}
反设 $f$ 在 $[a,b]$ 上无界，则对每个 $n$ 可取 $x_n\in[a,b]$ 使 $|f(x_n)|>n$。由有界数列子列定理，可选 $x_{n_k}\to x_*$。闭区间保证 $x_*\in[a,b]$，连续性保证 $f(x_{n_k})\to f(x_*)$，所以该输出子列有界；但其绝对值大于 $n_k\to\infty$，矛盾。闭性保证极限点留在定义域，有界性用于产生收敛子列，连续性则把输入收敛转为输出收敛。
\end{solve}

\begin{solve}[习题~\ref{ex:c7q3}：最大值定理的证明]\phantomsection\label{ans:c7q3}
先由有界性定理设 $M=\sup f([a,b])$。对每个 $n$ 选 $x_n$ 使 $M-1/n<f(x_n)\le M$。从 $x_n$ 中取收敛子列，极限为区间内某点 $x_*$。由连续性，$f(x_{n_k})\to f(x_*)$；另一方面由夹逼，该子列函数值趋于 $M$。极限唯一性给出 $f(x_*)=M$，所以最大值被取得。对 $-f$ 应用同一论证即得最小值。仅有有界性不足以取得上确界，例如 $(0,1)$ 上的 $f(x)=x$。
\end{solve}

\begin{solve}[习题~\ref{ex:c7q4}：介值定理的上确界证明]\phantomsection\label{ans:c7q4}
令 $A=\{x\in[a,b]:f(x)<0\}$，它非空且有上界。设 $c=\sup A$。由 $f(b)>0$ 及连续性，$b$ 左侧某邻域上函数仍为正，所以 $c<b$；同理由 $f(a)<0$ 可知 $c>a$。若 $f(c)<0$，连续性保证 $c$ 右侧仍有负值，违反上界性；若 $f(c)>0$，连续性保证 $c$ 左侧某邻域无负值，违反 $A$ 从左侧任意逼近上确界的性质。因此只能有 $f(c)=0$。
\end{solve}

\begin{solve}[习题~\ref{ex:c7q5}：闭区间上一致连续]\phantomsection\label{ans:c7q5}
若不一致连续，存在 $\varepsilon_0>0$，对每个 $n$ 可选 $x_n,y_n\in[a,b]$，使 $|x_n-y_n|<1/n$，但 $|f(x_n)-f(y_n)|\ge\varepsilon_0$。从 $x_n$ 取收敛子列 $x_{n_k}\to c\in[a,b]$，由两点距离趋零也有 $y_{n_k}\to c$。连续性给出两列函数值都趋于 $f(c)$，因此其差趋零，矛盾。紧性让可能在不同位置失控的点列集中到同一个极限点，而连续性排除该处的失控。
\end{solve}

\begin{solve}[习题~\ref{ex:c7q6}：一致连续延拓到端点]\phantomsection\label{ans:c7q6}
先取任意 $x_n\in(a,b)$ 且 $x_n\to a^+$。它是 Cauchy 数列，一致连续使 $f(x_n)$ 也为 Cauchy 数列，故有有限极限。若另取 $y_n\to a^+$，则 $|x_n-y_n|\to0$，一致连续推出输出差趋零，所以不同趋近数列的输出极限相同。由序列判别，右极限存在且有限。定义 $f(a)$ 为这个极限；右端点同理。只要求通常连续则不够，例如 $1/(x-a)$ 在开区间连续却无法有限延拓到左端点。
\end{solve}

\begin{solve}[习题~\ref{ex:c7q7}：可导推出连续]\phantomsection\label{ans:c7q7}
设 $f$ 在 $a$ 可导。对 $h\ne0$，写成 $f(a+h)-f(a)=h\,[f(a+h)-f(a)]/h$。第二因子趋于有限数 $f^{\prime}(a)$，因此在0附近有界；第一因子趋零，于是乘积趋零。这说明 $f(a+h)\to f(a)$。逆命题不成立，$|x|$ 在0连续，但左右差商分别为负1和1，导数不存在。可导性包含对一阶变化率的要求，比仅要求函数值接近更强。
\end{solve}

\begin{solve}[习题~\ref{ex:c7q8}：Fermat 必要条件]\phantomsection\label{ans:c7q8}
局部极大意味着充分小的 $h$ 满足 $f(a+h)-f(a)\le0$。当 $h>0$ 时差商不大于0，因此右导数不大于0；当 $h<0$ 时除以负数，差商不小于0，因此左导数不小于0。可导性使两侧极限相同，只能等于0。内点条件保证可从两侧趋近，可导条件保证左右导数相等。端点最值如 $[0,1]$ 上的 $f(x)=x$ 在1处最大，不要求导数为零。
\end{solve}

\begin{solve}[习题~\ref{ex:c7q9}：Rolle 定理的证明]\phantomsection\label{ans:c7q9}
设 $f$ 在 $[a,b]$ 连续、在 $(a,b)$ 可导，且 $f(a)=f(b)$。由最值定理，最大值与最小值均取得。若二者相等，函数处处为常数，任意内点的导数都为零。若二者不同，则至少一个最值不同于端点公共值，因此在内点取得；由 Fermat 必要条件，该内点导数为零。连续性用于取得最值，可导性用于内点极值的导数结论，端点相等用于迫使非平凡极值进入内部。
\end{solve}

\begin{solve}[习题~\ref{ex:c7q10}：从 Rolle 推出 Lagrange]\phantomsection\label{ans:c7q10}
设割线斜率 $k=[f(b)-f(a)]/(b-a)$，构造 $g(x)=f(x)-f(a)-k(x-a)$。它继承原函数在闭区间的连续性和开区间的可导性，且 $g(a)=g(b)=0$。Rolle 定理给出 $\xi\in(a,b)$，满足 $g^{\prime}(\xi)=0$，即 $f^{\prime}(\xi)=k$。辅助函数的意义是从原图像中减去端点割线，使两个端点被拉到同一高度。
\end{solve}

\begin{solve}[习题~\ref{ex:c7q11}：导数恒零推出常函数]\phantomsection\label{ans:c7q11}
任取区间中的 $x<y$，闭线段 $[x,y]$ 仍在定义域内。可导性蕴含连续性，可对该闭线段使用中值定理，得到 $f(y)-f(x)=f^{\prime}(\xi)(y-x)=0$。因为两点任意，函数为常数。如果定义域不连通，结论只在每个区间分支上成立；例如在 $(-\infty,0)\cup(0,\infty)$ 上，左支恒为0、右支恒为1，导数处处为零却不是全域同一常数。
\end{solve}

\begin{solve}[习题~\ref{ex:c7q12}：严格递增的充分条件并非必要]\phantomsection\label{ans:c7q12}
任意 $x<y$，中值定理给出 $f(y)-f(x)=f^{\prime}(\xi)(y-x)>0$，所以严格递增。逆命题不成立，$f(x)=x^3$ 在实轴严格递增，但 $f^{\prime}(0)=0$。对可导且递增函数，只能由正向差商取极限得到 $f^{\prime}\ge0$，严格性在取极限时可能消失。
\end{solve}

\begin{solve}[习题~\ref{ex:c7q13}：凸函数的切线下界]\phantomsection\label{ans:c7q13}
二阶导数非负说明一阶导数单调不减。若 $x>a$，中值定理给出差 $f(x)-f(a)=f^{\prime}(\xi)(x-a)$，其中 $\xi>a$，故不小于 $f^{\prime}(a)(x-a)$。若 $x<a$，中间点 $\xi<a$，有 $f^{\prime}(\xi)\le f^{\prime}(a)$，再乘负数 $x-a$ 后不等号反向，仍得到同一结论。$x=a$ 时等号成立。分别处理左右两侧可以避免符号反转错误。
\end{solve}

\begin{solve}[习题~\ref{ex:c7q14}：指数与对数的基本不等式]\phantomsection\label{ans:c7q14}
设 $g(x)=e^x-1-x$，则 $g^{\prime}=e^x-1$ 在负半轴为负、正半轴为正，所以 $g$ 在0取得唯一最小值0，得到第一式且仅在0取等。设 $h(x)=x-1-\log x$，导数为 $(x-1)/x$，在 $(0,1)$ 为负、$(1,\infty)$ 为正，所以在1处取得唯一最小值0，得到第二式。两者也可看作凸函数位于切线上方的实例。
\end{solve}

\begin{solve}[习题~\ref{ex:c7q15}：导函数的介值性]\phantomsection\label{ans:c7q15}
设 $g(x)=f(x)-\lambda x$，则 $g^{\prime}(a)<0$、$g^{\prime}(b)>0$。由端点导数的符号，$a$ 右侧有点使 $g$ 小于 $g(a)$，$b$ 左侧也有点使 $g$ 小于 $g(b)$，所以两个端点都不是闭区间上的最小点。连续函数在闭区间取得最小值，最小点必位于内部；Fermat 条件给出 $g^{\prime}(c)=0$，即所求。证明没有假设导函数连续，说明导函数即使不连续也不能发生跳跃间断。
\end{solve}

\begin{solve}[习题~\ref{ex:c7q16}：导数不连续的可导函数]\phantomsection\label{ans:c7q16}
在0处由定义，$[f(h)-f(0)]/h=h\sin(1/h)\to0$，故 $f^{\prime}(0)=0$。非零点用乘积和链式法则，$f^{\prime}(x)=2x\sin(1/x)-\cos(1/x)$。取 $x_n=1/(2n\pi)$，导数为负1；取 $y_n=1/((2n+1)\pi)$，导数为1。两列都趋于0，所以导函数在0没有极限，更不连续。存在导数与导数连续是两个不同层次的正则性条件。
\end{solve}

\begin{solve}[习题~\ref{ex:c7q17}：两个原函数之差]\phantomsection\label{ans:c7q17}
差函数 $H=F-G$ 在区间上可导且导数恒为零，由中值定理推出它为常函数，所以 $F=G+C$。这解释了不定积分中的任意常数。若定义域分成多个互不连通区间，则各区间可以取不同常数，例如 $1/x$ 的原函数在负半轴与正半轴上的常数不必相同。
\end{solve}

\begin{solve}[习题~\ref{ex:c7q18}：利用 Rolle 控制零点个数]\phantomsection\label{ans:c7q18}
反设有三个不同零点 $x_1<x_2<x_3$。分别在两段相邻零点之间使用 Rolle 定理，得到 $u\in(x_1,x_2)$、$v\in(x_2,x_3)$，使 $f^{\prime}(u)=f^{\prime}(v)=0$。再对导函数在 $[u,v]$ 上用 Rolle 定理，存在 $w$ 使 $f^{\prime\prime}(w)=0$，与严格为正矛盾。二阶可导保证 $f^{\prime}$ 连续，满足第二次应用的闭区间连续条件。
\end{solve}

\begin{solve}[习题~\ref{ex:c7q19}：Taylor 余项的局部估计]\phantomsection\label{ans:c7q19}
在连接 $a$ 与 $x$ 的闭线段上使用带 Lagrange 余项的 Taylor 公式，存在中间点 $\xi$ 使 $f(x)=f(a)+f^{\prime}(a)(x-a)+f^{\prime\prime}(a)(x-a)^2/2+f^{(3)}(\xi)(x-a)^3/6$。取绝对值并用导数上界即得结论。上界必须覆盖整条线段，不能只知道 $|f^{(3)}(a)|$ 就替代 $M$。当 $x<a$ 时三次幂为负，但绝对误差仍由 $|x-a|^3$ 控制。
\end{solve}

\begin{solve}[习题~\ref{ex:c7q20}：逆函数的导数条件]\phantomsection\label{ans:c7q20}
设逆函数为 $g$，由连续严格单调性，$g$ 在 $b$ 连续。令 $y\to b$、$x=g(y)\to a$，则 $y=f(x)$，且去心时 $x\ne a$。差商为 $[g(y)-g(b)]/(y-b)=(x-a)/(f(x)-f(a))$，极限是非零导数的倒数。条件 $f^{\prime}(a)\ne0$ 不能省略；例如 $f(x)=x^3$ 在0导数为0，其逆函数 $x^{1/3}$ 在0没有有限导数。
\end{solve}

\clearpage
\section{导数与复合函数求导}\label{sec:ch8}
导数（derivative）是局部线性变化率。乘法、商法与链式法则可以组合使用；在核对公式时，应把“外层函数的导数”和“内层函数的导数”分别写清楚。

\textbf{资源对应：}CLP-1 \S~2.4、\S~2.6--\S~2.10。本章共 24 道编号习题。
\subsection{习题}
\begin{example}[多项式与指数的乘积]\label{ex:c8q1}
求 $f(x)=x^{2}e^{1x}$ 的一阶、二阶导数，并求 $f^{\prime}(1)$。
\par\textit{改编训练。}\par\hyperref[ans:c8q1]{解答见第~\pageref{ans:c8q1} 页。}
\end{example}
\begin{example}[多项式与指数的乘积]\label{ex:c8q2}
求 $f(x)=x^{3}e^{2x}$ 的一阶、二阶导数，并求 $f^{\prime}(1)$。
\par\textit{改编训练。}\par\hyperref[ans:c8q2]{解答见第~\pageref{ans:c8q2} 页。}
\end{example}
\begin{example}[多项式与指数的乘积]\label{ex:c8q3}
求 $f(x)=x^{2}e^{3x}$ 的一阶、二阶导数，并求 $f^{\prime}(1)$。
\par\textit{改编训练。}\par\hyperref[ans:c8q3]{解答见第~\pageref{ans:c8q3} 页。}
\end{example}
\begin{example}[多项式与指数的乘积]\label{ex:c8q4}
求 $f(x)=x^{4}e^{2x}$ 的一阶、二阶导数，并求 $f^{\prime}(1)$。
\par\textit{改编训练。}\par\hyperref[ans:c8q4]{解答见第~\pageref{ans:c8q4} 页。}
\end{example}
\begin{example}[多项式与指数的乘积]\label{ex:c8q5}
求 $f(x)=x^{3}e^{4x}$ 的一阶、二阶导数，并求 $f^{\prime}(1)$。
\par\textit{改编训练。}\par\hyperref[ans:c8q5]{解答见第~\pageref{ans:c8q5} 页。}
\end{example}
\begin{example}[多项式与指数的乘积]\label{ex:c8q6}
求 $f(x)=x^{5}e^{3x}$ 的一阶、二阶导数，并求 $f^{\prime}(1)$。
\par\textit{改编训练。}\par\hyperref[ans:c8q6]{解答见第~\pageref{ans:c8q6} 页。}
\end{example}
\begin{example}[幂函数的链式法则]\label{ex:c8q7}
求 $f(x)=(1x^2+2)^{2}$ 的导数，写出 $x=1$ 处的切线，并判断此处函数增减方向。
\par\textit{改编训练。}\par\hyperref[ans:c8q7]{解答见第~\pageref{ans:c8q7} 页。}
\end{example}
\begin{example}[幂函数的链式法则]\label{ex:c8q8}
求 $f(x)=(2x^2+1)^{3}$ 的导数，写出 $x=1$ 处的切线，并判断此处函数增减方向。
\par\textit{改编训练。}\par\hyperref[ans:c8q8]{解答见第~\pageref{ans:c8q8} 页。}
\end{example}
\begin{example}[幂函数的链式法则]\label{ex:c8q9}
求 $f(x)=(3x^2+2)^{2}$ 的导数，写出 $x=1$ 处的切线，并判断此处函数增减方向。
\par\textit{改编训练。}\par\hyperref[ans:c8q9]{解答见第~\pageref{ans:c8q9} 页。}
\end{example}
\begin{example}[幂函数的链式法则]\label{ex:c8q10}
求 $f(x)=(2x^2+3)^{4}$ 的导数，写出 $x=1$ 处的切线，并判断此处函数增减方向。
\par\textit{改编训练。}\par\hyperref[ans:c8q10]{解答见第~\pageref{ans:c8q10} 页。}
\end{example}
\begin{example}[幂函数的链式法则]\label{ex:c8q11}
求 $f(x)=(4x^2+1)^{3}$ 的导数，写出 $x=1$ 处的切线，并判断此处函数增减方向。
\par\textit{改编训练。}\par\hyperref[ans:c8q11]{解答见第~\pageref{ans:c8q11} 页。}
\end{example}
\begin{example}[幂函数的链式法则]\label{ex:c8q12}
求 $f(x)=(3x^2+1)^{5}$ 的导数，写出 $x=1$ 处的切线，并判断此处函数增减方向。
\par\textit{改编训练。}\par\hyperref[ans:c8q12]{解答见第~\pageref{ans:c8q12} 页。}
\end{example}
\begin{example}[对数复合与二阶导数]\label{ex:c8q13}
设 $f(x)=\log(1x^2+1)$。求 $f^{\prime}$、$f^{\prime\prime}$，并写出导数为零的点。
\par\textit{改编训练。}\par\hyperref[ans:c8q13]{解答见第~\pageref{ans:c8q13} 页。}
\end{example}
\begin{example}[对数复合与二阶导数]\label{ex:c8q14}
设 $f(x)=\log(2x^2+1)$。求 $f^{\prime}$、$f^{\prime\prime}$，并写出导数为零的点。
\par\textit{改编训练。}\par\hyperref[ans:c8q14]{解答见第~\pageref{ans:c8q14} 页。}
\end{example}
\begin{example}[对数复合与二阶导数]\label{ex:c8q15}
设 $f(x)=\log(1x^2+3)$。求 $f^{\prime}$、$f^{\prime\prime}$，并写出导数为零的点。
\par\textit{改编训练。}\par\hyperref[ans:c8q15]{解答见第~\pageref{ans:c8q15} 页。}
\end{example}
\begin{example}[对数复合与二阶导数]\label{ex:c8q16}
设 $f(x)=\log(3x^2+2)$。求 $f^{\prime}$、$f^{\prime\prime}$，并写出导数为零的点。
\par\textit{改编训练。}\par\hyperref[ans:c8q16]{解答见第~\pageref{ans:c8q16} 页。}
\end{example}
\begin{example}[对数复合与二阶导数]\label{ex:c8q17}
设 $f(x)=\log(2x^2+5)$。求 $f^{\prime}$、$f^{\prime\prime}$，并写出导数为零的点。
\par\textit{改编训练。}\par\hyperref[ans:c8q17]{解答见第~\pageref{ans:c8q17} 页。}
\end{example}
\begin{example}[对数复合与二阶导数]\label{ex:c8q18}
设 $f(x)=\log(4x^2+3)$。求 $f^{\prime}$、$f^{\prime\prime}$，并写出导数为零的点。
\par\textit{改编训练。}\par\hyperref[ans:c8q18]{解答见第~\pageref{ans:c8q18} 页。}
\end{example}
\begin{example}[指数与三角函数的复合求导]\label{ex:c8q19}
求 $f(x)=e^{1x}\sin(2x)$ 的一阶、二阶导数，并给出 $f^{\prime}(0)$ 与 $f^{\prime\prime}(0)$。
\par\textit{改编训练。}\par\hyperref[ans:c8q19]{解答见第~\pageref{ans:c8q19} 页。}
\end{example}
\begin{example}[指数与三角函数的复合求导]\label{ex:c8q20}
求 $f(x)=e^{2x}\sin(1x)$ 的一阶、二阶导数，并给出 $f^{\prime}(0)$ 与 $f^{\prime\prime}(0)$。
\par\textit{改编训练。}\par\hyperref[ans:c8q20]{解答见第~\pageref{ans:c8q20} 页。}
\end{example}
\begin{example}[指数与三角函数的复合求导]\label{ex:c8q21}
求 $f(x)=e^{2x}\sin(3x)$ 的一阶、二阶导数，并给出 $f^{\prime}(0)$ 与 $f^{\prime\prime}(0)$。
\par\textit{改编训练。}\par\hyperref[ans:c8q21]{解答见第~\pageref{ans:c8q21} 页。}
\end{example}
\begin{example}[指数与三角函数的复合求导]\label{ex:c8q22}
求 $f(x)=e^{3x}\sin(2x)$ 的一阶、二阶导数，并给出 $f^{\prime}(0)$ 与 $f^{\prime\prime}(0)$。
\par\textit{改编训练。}\par\hyperref[ans:c8q22]{解答见第~\pageref{ans:c8q22} 页。}
\end{example}
\begin{example}[指数与三角函数的复合求导]\label{ex:c8q23}
求 $f(x)=e^{4x}\sin(1x)$ 的一阶、二阶导数，并给出 $f^{\prime}(0)$ 与 $f^{\prime\prime}(0)$。
\par\textit{改编训练。}\par\hyperref[ans:c8q23]{解答见第~\pageref{ans:c8q23} 页。}
\end{example}
\begin{example}[指数与三角函数的复合求导]\label{ex:c8q24}
求 $f(x)=e^{1x}\sin(5x)$ 的一阶、二阶导数，并给出 $f^{\prime}(0)$ 与 $f^{\prime\prime}(0)$。
\par\textit{改编训练。}\par\hyperref[ans:c8q24]{解答见第~\pageref{ans:c8q24} 页。}
\end{example}
\subsection{解答}
\begin{solve}[习题~\ref{ex:c8q1}：多项式与指数的乘积]\phantomsection\label{ans:c8q1}
先用乘积法则，再用链式法则对指数函数求导：
\[
f^{\prime}(x)=e^{1x}(2x^{1}+1x^{2}).
\]
继续对两个因子求导并合并同类项，
\[
f^{\prime\prime}(x)=e^{1x}\left(2x^{0}+4x^{1}+1x^{2}\right).
\]
在 $x=1$ 处得到 $f^{\prime}(1)=3e^{1}$。二阶导数中间项的系数来自两次相同的交叉贡献；只对原乘积的两个因子分别求二阶导数会漏掉这一项。
\end{solve}

\begin{solve}[习题~\ref{ex:c8q2}：多项式与指数的乘积]\phantomsection\label{ans:c8q2}
先用乘积法则，再用链式法则对指数函数求导：
\[
f^{\prime}(x)=e^{2x}(3x^{2}+2x^{3}).
\]
继续对两个因子求导并合并同类项，
\[
f^{\prime\prime}(x)=e^{2x}\left(6x^{1}+12x^{2}+4x^{3}\right).
\]
在 $x=1$ 处得到 $f^{\prime}(1)=5e^{2}$。二阶导数中间项的系数来自两次相同的交叉贡献；只对原乘积的两个因子分别求二阶导数会漏掉这一项。
\end{solve}

\begin{solve}[习题~\ref{ex:c8q3}：多项式与指数的乘积]\phantomsection\label{ans:c8q3}
先用乘积法则，再用链式法则对指数函数求导：
\[
f^{\prime}(x)=e^{3x}(2x^{1}+3x^{2}).
\]
继续对两个因子求导并合并同类项，
\[
f^{\prime\prime}(x)=e^{3x}\left(2x^{0}+12x^{1}+9x^{2}\right).
\]
在 $x=1$ 处得到 $f^{\prime}(1)=5e^{3}$。二阶导数中间项的系数来自两次相同的交叉贡献；只对原乘积的两个因子分别求二阶导数会漏掉这一项。
\end{solve}

\begin{solve}[习题~\ref{ex:c8q4}：多项式与指数的乘积]\phantomsection\label{ans:c8q4}
先用乘积法则，再用链式法则对指数函数求导：
\[
f^{\prime}(x)=e^{2x}(4x^{3}+2x^{4}).
\]
继续对两个因子求导并合并同类项，
\[
f^{\prime\prime}(x)=e^{2x}\left(12x^{2}+16x^{3}+4x^{4}\right).
\]
在 $x=1$ 处得到 $f^{\prime}(1)=6e^{2}$。二阶导数中间项的系数来自两次相同的交叉贡献；只对原乘积的两个因子分别求二阶导数会漏掉这一项。
\end{solve}

\begin{solve}[习题~\ref{ex:c8q5}：多项式与指数的乘积]\phantomsection\label{ans:c8q5}
先用乘积法则，再用链式法则对指数函数求导：
\[
f^{\prime}(x)=e^{4x}(3x^{2}+4x^{3}).
\]
继续对两个因子求导并合并同类项，
\[
f^{\prime\prime}(x)=e^{4x}\left(6x^{1}+24x^{2}+16x^{3}\right).
\]
在 $x=1$ 处得到 $f^{\prime}(1)=7e^{4}$。二阶导数中间项的系数来自两次相同的交叉贡献；只对原乘积的两个因子分别求二阶导数会漏掉这一项。
\end{solve}

\begin{solve}[习题~\ref{ex:c8q6}：多项式与指数的乘积]\phantomsection\label{ans:c8q6}
先用乘积法则，再用链式法则对指数函数求导：
\[
f^{\prime}(x)=e^{3x}(5x^{4}+3x^{5}).
\]
继续对两个因子求导并合并同类项，
\[
f^{\prime\prime}(x)=e^{3x}\left(20x^{3}+30x^{4}+9x^{5}\right).
\]
在 $x=1$ 处得到 $f^{\prime}(1)=8e^{3}$。二阶导数中间项的系数来自两次相同的交叉贡献；只对原乘积的两个因子分别求二阶导数会漏掉这一项。
\end{solve}

\begin{solve}[习题~\ref{ex:c8q7}：幂函数的链式法则]\phantomsection\label{ans:c8q7}
外层为 $u^{2}$，内层为 $u=1x^2+2$。因此
\[
f^{\prime}(x)=4x(1x^2+2)^{1}.
\]
函数值为 $f(1)=9$，斜率为 $f^{\prime}(1)=12$，切线方程为
\[
y=9+12(x-1).
\]
该斜率为正，且导数在此处连续，故函数在 $1$ 的一个邻域内严格递增。切线给的是局部线性近似，不能据此断言函数在整个实轴上递增。
\end{solve}

\begin{solve}[习题~\ref{ex:c8q8}：幂函数的链式法则]\phantomsection\label{ans:c8q8}
外层为 $u^{3}$，内层为 $u=2x^2+1$。因此
\[
f^{\prime}(x)=12x(2x^2+1)^{2}.
\]
函数值为 $f(1)=27$，斜率为 $f^{\prime}(1)=108$，切线方程为
\[
y=27+108(x-1).
\]
该斜率为正，且导数在此处连续，故函数在 $1$ 的一个邻域内严格递增。切线给的是局部线性近似，不能据此断言函数在整个实轴上递增。
\end{solve}

\begin{solve}[习题~\ref{ex:c8q9}：幂函数的链式法则]\phantomsection\label{ans:c8q9}
外层为 $u^{2}$，内层为 $u=3x^2+2$。因此
\[
f^{\prime}(x)=12x(3x^2+2)^{1}.
\]
函数值为 $f(1)=25$，斜率为 $f^{\prime}(1)=60$，切线方程为
\[
y=25+60(x-1).
\]
该斜率为正，且导数在此处连续，故函数在 $1$ 的一个邻域内严格递增。切线给的是局部线性近似，不能据此断言函数在整个实轴上递增。
\end{solve}

\begin{solve}[习题~\ref{ex:c8q10}：幂函数的链式法则]\phantomsection\label{ans:c8q10}
外层为 $u^{4}$，内层为 $u=2x^2+3$。因此
\[
f^{\prime}(x)=16x(2x^2+3)^{3}.
\]
函数值为 $f(1)=625$，斜率为 $f^{\prime}(1)=2000$，切线方程为
\[
y=625+2000(x-1).
\]
该斜率为正，且导数在此处连续，故函数在 $1$ 的一个邻域内严格递增。切线给的是局部线性近似，不能据此断言函数在整个实轴上递增。
\end{solve}

\begin{solve}[习题~\ref{ex:c8q11}：幂函数的链式法则]\phantomsection\label{ans:c8q11}
外层为 $u^{3}$，内层为 $u=4x^2+1$。因此
\[
f^{\prime}(x)=24x(4x^2+1)^{2}.
\]
函数值为 $f(1)=125$，斜率为 $f^{\prime}(1)=600$，切线方程为
\[
y=125+600(x-1).
\]
该斜率为正，且导数在此处连续，故函数在 $1$ 的一个邻域内严格递增。切线给的是局部线性近似，不能据此断言函数在整个实轴上递增。
\end{solve}

\begin{solve}[习题~\ref{ex:c8q12}：幂函数的链式法则]\phantomsection\label{ans:c8q12}
外层为 $u^{5}$，内层为 $u=3x^2+1$。因此
\[
f^{\prime}(x)=30x(3x^2+1)^{4}.
\]
函数值为 $f(1)=1024$，斜率为 $f^{\prime}(1)=7680$，切线方程为
\[
y=1024+7680(x-1).
\]
该斜率为正，且导数在此处连续，故函数在 $1$ 的一个邻域内严格递增。切线给的是局部线性近似，不能据此断言函数在整个实轴上递增。
\end{solve}

\begin{solve}[习题~\ref{ex:c8q13}：对数复合与二阶导数]\phantomsection\label{ans:c8q13}
由于 $1x^2+1>0$，定义域为全体实数。链式法则给出
\[
f^{\prime}(x)=\frac{2x}{1x^2+1}.
\]
使用商法则，
\[
f^{\prime\prime}(x)=\frac{2(1x^2+1)-4x^2}{(1x^2+1)^2}=\frac{2-2x^2}{(1x^2+1)^2}.
\]
一阶导数的分母处处为正，所以唯一驻点为 $x=0$。导数在负半轴为负、正半轴为正，因此该点还是全局最小点，最小值为 $\log 1$。求导前先确认对数真数为正，可以避免错误扩张定义域。
\end{solve}

\begin{solve}[习题~\ref{ex:c8q14}：对数复合与二阶导数]\phantomsection\label{ans:c8q14}
由于 $2x^2+1>0$，定义域为全体实数。链式法则给出
\[
f^{\prime}(x)=\frac{4x}{2x^2+1}.
\]
使用商法则，
\[
f^{\prime\prime}(x)=\frac{4(2x^2+1)-16x^2}{(2x^2+1)^2}=\frac{4-8x^2}{(2x^2+1)^2}.
\]
一阶导数的分母处处为正，所以唯一驻点为 $x=0$。导数在负半轴为负、正半轴为正，因此该点还是全局最小点，最小值为 $\log 1$。求导前先确认对数真数为正，可以避免错误扩张定义域。
\end{solve}

\begin{solve}[习题~\ref{ex:c8q15}：对数复合与二阶导数]\phantomsection\label{ans:c8q15}
由于 $1x^2+3>0$，定义域为全体实数。链式法则给出
\[
f^{\prime}(x)=\frac{2x}{1x^2+3}.
\]
使用商法则，
\[
f^{\prime\prime}(x)=\frac{2(1x^2+3)-4x^2}{(1x^2+3)^2}=\frac{6-2x^2}{(1x^2+3)^2}.
\]
一阶导数的分母处处为正，所以唯一驻点为 $x=0$。导数在负半轴为负、正半轴为正，因此该点还是全局最小点，最小值为 $\log 3$。求导前先确认对数真数为正，可以避免错误扩张定义域。
\end{solve}

\begin{solve}[习题~\ref{ex:c8q16}：对数复合与二阶导数]\phantomsection\label{ans:c8q16}
由于 $3x^2+2>0$，定义域为全体实数。链式法则给出
\[
f^{\prime}(x)=\frac{6x}{3x^2+2}.
\]
使用商法则，
\[
f^{\prime\prime}(x)=\frac{6(3x^2+2)-36x^2}{(3x^2+2)^2}=\frac{12-18x^2}{(3x^2+2)^2}.
\]
一阶导数的分母处处为正，所以唯一驻点为 $x=0$。导数在负半轴为负、正半轴为正，因此该点还是全局最小点，最小值为 $\log 2$。求导前先确认对数真数为正，可以避免错误扩张定义域。
\end{solve}

\begin{solve}[习题~\ref{ex:c8q17}：对数复合与二阶导数]\phantomsection\label{ans:c8q17}
由于 $2x^2+5>0$，定义域为全体实数。链式法则给出
\[
f^{\prime}(x)=\frac{4x}{2x^2+5}.
\]
使用商法则，
\[
f^{\prime\prime}(x)=\frac{4(2x^2+5)-16x^2}{(2x^2+5)^2}=\frac{20-8x^2}{(2x^2+5)^2}.
\]
一阶导数的分母处处为正，所以唯一驻点为 $x=0$。导数在负半轴为负、正半轴为正，因此该点还是全局最小点，最小值为 $\log 5$。求导前先确认对数真数为正，可以避免错误扩张定义域。
\end{solve}

\begin{solve}[习题~\ref{ex:c8q18}：对数复合与二阶导数]\phantomsection\label{ans:c8q18}
由于 $4x^2+3>0$，定义域为全体实数。链式法则给出
\[
f^{\prime}(x)=\frac{8x}{4x^2+3}.
\]
使用商法则，
\[
f^{\prime\prime}(x)=\frac{8(4x^2+3)-64x^2}{(4x^2+3)^2}=\frac{24-32x^2}{(4x^2+3)^2}.
\]
一阶导数的分母处处为正，所以唯一驻点为 $x=0$。导数在负半轴为负、正半轴为正，因此该点还是全局最小点，最小值为 $\log 3$。求导前先确认对数真数为正，可以避免错误扩张定义域。
\end{solve}

\begin{solve}[习题~\ref{ex:c8q19}：指数与三角函数的复合求导]\phantomsection\label{ans:c8q19}
乘积法则与链式法则给出
\[
f^{\prime}(x)=e^{1x}\bigl(1\sin(2x)+2\cos(2x)\bigr).
\]
再次求导并合并正弦、余弦项，
\[
f^{\prime\prime}(x)=e^{1x}\bigl(-3\sin(2x)+4\cos(2x)\bigr).
\]
故 $f^{\prime}(0)=2$、$f^{\prime\prime}(0)=4$。检验二阶导数时应特别注意余弦求导带来的负号，这正是正弦项系数中出现 $a^2-b^2$ 的原因。
\end{solve}

\begin{solve}[习题~\ref{ex:c8q20}：指数与三角函数的复合求导]\phantomsection\label{ans:c8q20}
乘积法则与链式法则给出
\[
f^{\prime}(x)=e^{2x}\bigl(2\sin(1x)+1\cos(1x)\bigr).
\]
再次求导并合并正弦、余弦项，
\[
f^{\prime\prime}(x)=e^{2x}\bigl(3\sin(1x)+4\cos(1x)\bigr).
\]
故 $f^{\prime}(0)=1$、$f^{\prime\prime}(0)=4$。检验二阶导数时应特别注意余弦求导带来的负号，这正是正弦项系数中出现 $a^2-b^2$ 的原因。
\end{solve}

\begin{solve}[习题~\ref{ex:c8q21}：指数与三角函数的复合求导]\phantomsection\label{ans:c8q21}
乘积法则与链式法则给出
\[
f^{\prime}(x)=e^{2x}\bigl(2\sin(3x)+3\cos(3x)\bigr).
\]
再次求导并合并正弦、余弦项，
\[
f^{\prime\prime}(x)=e^{2x}\bigl(-5\sin(3x)+12\cos(3x)\bigr).
\]
故 $f^{\prime}(0)=3$、$f^{\prime\prime}(0)=12$。检验二阶导数时应特别注意余弦求导带来的负号，这正是正弦项系数中出现 $a^2-b^2$ 的原因。
\end{solve}

\begin{solve}[习题~\ref{ex:c8q22}：指数与三角函数的复合求导]\phantomsection\label{ans:c8q22}
乘积法则与链式法则给出
\[
f^{\prime}(x)=e^{3x}\bigl(3\sin(2x)+2\cos(2x)\bigr).
\]
再次求导并合并正弦、余弦项，
\[
f^{\prime\prime}(x)=e^{3x}\bigl(5\sin(2x)+12\cos(2x)\bigr).
\]
故 $f^{\prime}(0)=2$、$f^{\prime\prime}(0)=12$。检验二阶导数时应特别注意余弦求导带来的负号，这正是正弦项系数中出现 $a^2-b^2$ 的原因。
\end{solve}

\begin{solve}[习题~\ref{ex:c8q23}：指数与三角函数的复合求导]\phantomsection\label{ans:c8q23}
乘积法则与链式法则给出
\[
f^{\prime}(x)=e^{4x}\bigl(4\sin(1x)+1\cos(1x)\bigr).
\]
再次求导并合并正弦、余弦项，
\[
f^{\prime\prime}(x)=e^{4x}\bigl(15\sin(1x)+8\cos(1x)\bigr).
\]
故 $f^{\prime}(0)=1$、$f^{\prime\prime}(0)=8$。检验二阶导数时应特别注意余弦求导带来的负号，这正是正弦项系数中出现 $a^2-b^2$ 的原因。
\end{solve}

\begin{solve}[习题~\ref{ex:c8q24}：指数与三角函数的复合求导]\phantomsection\label{ans:c8q24}
乘积法则与链式法则给出
\[
f^{\prime}(x)=e^{1x}\bigl(1\sin(5x)+5\cos(5x)\bigr).
\]
再次求导并合并正弦、余弦项，
\[
f^{\prime\prime}(x)=e^{1x}\bigl(-24\sin(5x)+10\cos(5x)\bigr).
\]
故 $f^{\prime}(0)=5$、$f^{\prime\prime}(0)=10$。检验二阶导数时应特别注意余弦求导带来的负号，这正是正弦项系数中出现 $a^2-b^2$ 的原因。
\end{solve}

\clearpage
\section{隐函数、参数方程与高阶导数}\label{sec:ch9}
隐函数求导必须把 $y$ 看成 $x$ 的函数；参数曲线求导则用 $dy/dx=(dy/dt)/(dx/dt)$，适用点需要 $dx/dt\ne0$。高阶导数可以利用归纳法或 Leibniz 公式。

\textbf{资源对应：}CLP-1 \S~2.11、\S~2.14；CLP-3 \S~1.6。本章共 24 道编号习题。
\subsection{习题}
\begin{example}[圆的隐函数导数]\label{ex:c9q1}
曲线 $x^2+y^2=5$ 经过 $P=(1,2)$。求 $P$ 处的切线，以及局部函数 $y(x)$ 在该点的一阶、二阶导数。
\par\textit{改编训练。}\par\hyperref[ans:c9q1]{解答见第~\pageref{ans:c9q1} 页。}
\end{example}
\begin{example}[圆的隐函数导数]\label{ex:c9q2}
曲线 $x^2+y^2=13$ 经过 $P=(2,3)$。求 $P$ 处的切线，以及局部函数 $y(x)$ 在该点的一阶、二阶导数。
\par\textit{改编训练。}\par\hyperref[ans:c9q2]{解答见第~\pageref{ans:c9q2} 页。}
\end{example}
\begin{example}[圆的隐函数导数]\label{ex:c9q3}
曲线 $x^2+y^2=10$ 经过 $P=(3,1)$。求 $P$ 处的切线，以及局部函数 $y(x)$ 在该点的一阶、二阶导数。
\par\textit{改编训练。}\par\hyperref[ans:c9q3]{解答见第~\pageref{ans:c9q3} 页。}
\end{example}
\begin{example}[圆的隐函数导数]\label{ex:c9q4}
曲线 $x^2+y^2=29$ 经过 $P=(2,5)$。求 $P$ 处的切线，以及局部函数 $y(x)$ 在该点的一阶、二阶导数。
\par\textit{改编训练。}\par\hyperref[ans:c9q4]{解答见第~\pageref{ans:c9q4} 页。}
\end{example}
\begin{example}[圆的隐函数导数]\label{ex:c9q5}
曲线 $x^2+y^2=25$ 经过 $P=(4,3)$。求 $P$ 处的切线，以及局部函数 $y(x)$ 在该点的一阶、二阶导数。
\par\textit{改编训练。}\par\hyperref[ans:c9q5]{解答见第~\pageref{ans:c9q5} 页。}
\end{example}
\begin{example}[圆的隐函数导数]\label{ex:c9q6}
曲线 $x^2+y^2=34$ 经过 $P=(3,5)$。求 $P$ 处的切线，以及局部函数 $y(x)$ 在该点的一阶、二阶导数。
\par\textit{改编训练。}\par\hyperref[ans:c9q6]{解答见第~\pageref{ans:c9q6} 页。}
\end{example}
\begin{example}[幂参数曲线]\label{ex:c9q7}
设 $x=1t^2$、$y=2t^3$。求 $t=1$ 处的 $dy/dx$、$d^2y/dx^2$ 和切线方程。
\par\textit{改编训练。}\par\hyperref[ans:c9q7]{解答见第~\pageref{ans:c9q7} 页。}
\end{example}
\begin{example}[幂参数曲线]\label{ex:c9q8}
设 $x=2t^2$、$y=1t^3$。求 $t=1$ 处的 $dy/dx$、$d^2y/dx^2$ 和切线方程。
\par\textit{改编训练。}\par\hyperref[ans:c9q8]{解答见第~\pageref{ans:c9q8} 页。}
\end{example}
\begin{example}[幂参数曲线]\label{ex:c9q9}
设 $x=2t^2$、$y=3t^3$。求 $t=1$ 处的 $dy/dx$、$d^2y/dx^2$ 和切线方程。
\par\textit{改编训练。}\par\hyperref[ans:c9q9]{解答见第~\pageref{ans:c9q9} 页。}
\end{example}
\begin{example}[幂参数曲线]\label{ex:c9q10}
设 $x=3t^2$、$y=2t^3$。求 $t=1$ 处的 $dy/dx$、$d^2y/dx^2$ 和切线方程。
\par\textit{改编训练。}\par\hyperref[ans:c9q10]{解答见第~\pageref{ans:c9q10} 页。}
\end{example}
\begin{example}[幂参数曲线]\label{ex:c9q11}
设 $x=4t^2$、$y=1t^3$。求 $t=1$ 处的 $dy/dx$、$d^2y/dx^2$ 和切线方程。
\par\textit{改编训练。}\par\hyperref[ans:c9q11]{解答见第~\pageref{ans:c9q11} 页。}
\end{example}
\begin{example}[幂参数曲线]\label{ex:c9q12}
设 $x=1t^2$、$y=4t^3$。求 $t=1$ 处的 $dy/dx$、$d^2y/dx^2$ 和切线方程。
\par\textit{改编训练。}\par\hyperref[ans:c9q12]{解答见第~\pageref{ans:c9q12} 页。}
\end{example}
\begin{example}[对数的高阶导数]\label{ex:c9q13}
设 $f(x)=\log(x+1)$，$x>-1$。证明其 $m$ 阶导数公式，并求 $f^{(3)}(0)$。
\par\textit{改编训练。}\par\hyperref[ans:c9q13]{解答见第~\pageref{ans:c9q13} 页。}
\end{example}
\begin{example}[对数的高阶导数]\label{ex:c9q14}
设 $f(x)=\log(x+2)$，$x>-2$。证明其 $m$ 阶导数公式，并求 $f^{(4)}(0)$。
\par\textit{改编训练。}\par\hyperref[ans:c9q14]{解答见第~\pageref{ans:c9q14} 页。}
\end{example}
\begin{example}[对数的高阶导数]\label{ex:c9q15}
设 $f(x)=\log(x+3)$，$x>-3$。证明其 $m$ 阶导数公式，并求 $f^{(5)}(0)$。
\par\textit{改编训练。}\par\hyperref[ans:c9q15]{解答见第~\pageref{ans:c9q15} 页。}
\end{example}
\begin{example}[对数的高阶导数]\label{ex:c9q16}
设 $f(x)=\log(x+4)$，$x>-4$。证明其 $m$ 阶导数公式，并求 $f^{(3)}(0)$。
\par\textit{改编训练。}\par\hyperref[ans:c9q16]{解答见第~\pageref{ans:c9q16} 页。}
\end{example}
\begin{example}[对数的高阶导数]\label{ex:c9q17}
设 $f(x)=\log(x+2)$，$x>-2$。证明其 $m$ 阶导数公式，并求 $f^{(6)}(0)$。
\par\textit{改编训练。}\par\hyperref[ans:c9q17]{解答见第~\pageref{ans:c9q17} 页。}
\end{example}
\begin{example}[对数的高阶导数]\label{ex:c9q18}
设 $f(x)=\log(x+5)$，$x>-5$。证明其 $m$ 阶导数公式，并求 $f^{(4)}(0)$。
\par\textit{改编训练。}\par\hyperref[ans:c9q18]{解答见第~\pageref{ans:c9q18} 页。}
\end{example}
\begin{example}[Leibniz 公式的短和]\label{ex:c9q19}
求 $f(x)=x^2e^{1x}$ 的 $3$ 阶导数及其在 $0$ 处的值。
\par\textit{改编训练。}\par\hyperref[ans:c9q19]{解答见第~\pageref{ans:c9q19} 页。}
\end{example}
\begin{example}[Leibniz 公式的短和]\label{ex:c9q20}
求 $f(x)=x^2e^{2x}$ 的 $4$ 阶导数及其在 $0$ 处的值。
\par\textit{改编训练。}\par\hyperref[ans:c9q20]{解答见第~\pageref{ans:c9q20} 页。}
\end{example}
\begin{example}[Leibniz 公式的短和]\label{ex:c9q21}
求 $f(x)=x^2e^{3x}$ 的 $2$ 阶导数及其在 $0$ 处的值。
\par\textit{改编训练。}\par\hyperref[ans:c9q21]{解答见第~\pageref{ans:c9q21} 页。}
\end{example}
\begin{example}[Leibniz 公式的短和]\label{ex:c9q22}
求 $f(x)=x^2e^{2x}$ 的 $5$ 阶导数及其在 $0$ 处的值。
\par\textit{改编训练。}\par\hyperref[ans:c9q22]{解答见第~\pageref{ans:c9q22} 页。}
\end{example}
\begin{example}[Leibniz 公式的短和]\label{ex:c9q23}
求 $f(x)=x^2e^{4x}$ 的 $3$ 阶导数及其在 $0$ 处的值。
\par\textit{改编训练。}\par\hyperref[ans:c9q23]{解答见第~\pageref{ans:c9q23} 页。}
\end{example}
\begin{example}[Leibniz 公式的短和]\label{ex:c9q24}
求 $f(x)=x^2e^{3x}$ 的 $6$ 阶导数及其在 $0$ 处的值。
\par\textit{改编训练。}\par\hyperref[ans:c9q24]{解答见第~\pageref{ans:c9q24} 页。}
\end{example}
\subsection{解答}
\begin{solve}[习题~\ref{ex:c9q1}：圆的隐函数导数]\phantomsection\label{ans:c9q1}
因 $P$ 处 $y=2\ne0$，可以在附近把曲线写成可微函数。对原式求导，
\[
2x+2yy^{{\prime}}=0,\qquad y^{{\prime}}=-x/y.
\]
代入得 $y^{\prime}(1)=-\frac{1}{2}$，切线为
\[
1(x-1)+2(y-2)=0.
\]
再求导得到 $1+(y^{\prime})^2+yy^{\prime\prime}=0$，所以
\[
y^{\prime\prime}=-\frac{x^2+y^2}{y^3}=-\frac{5}{y^3},\qquad y^{\prime\prime}(1)=-\frac{5}{8}.
\]
切线公式也可由半径与切线垂直得到，这提供一个独立的几何检验。
\end{solve}

\begin{solve}[习题~\ref{ex:c9q2}：圆的隐函数导数]\phantomsection\label{ans:c9q2}
因 $P$ 处 $y=3\ne0$，可以在附近把曲线写成可微函数。对原式求导，
\[
2x+2yy^{{\prime}}=0,\qquad y^{{\prime}}=-x/y.
\]
代入得 $y^{\prime}(2)=-\frac{2}{3}$，切线为
\[
2(x-2)+3(y-3)=0.
\]
再求导得到 $1+(y^{\prime})^2+yy^{\prime\prime}=0$，所以
\[
y^{\prime\prime}=-\frac{x^2+y^2}{y^3}=-\frac{13}{y^3},\qquad y^{\prime\prime}(2)=-\frac{13}{27}.
\]
切线公式也可由半径与切线垂直得到，这提供一个独立的几何检验。
\end{solve}

\begin{solve}[习题~\ref{ex:c9q3}：圆的隐函数导数]\phantomsection\label{ans:c9q3}
因 $P$ 处 $y=1\ne0$，可以在附近把曲线写成可微函数。对原式求导，
\[
2x+2yy^{{\prime}}=0,\qquad y^{{\prime}}=-x/y.
\]
代入得 $y^{\prime}(3)=-3$，切线为
\[
3(x-3)+1(y-1)=0.
\]
再求导得到 $1+(y^{\prime})^2+yy^{\prime\prime}=0$，所以
\[
y^{\prime\prime}=-\frac{x^2+y^2}{y^3}=-\frac{10}{y^3},\qquad y^{\prime\prime}(3)=-10.
\]
切线公式也可由半径与切线垂直得到，这提供一个独立的几何检验。
\end{solve}

\begin{solve}[习题~\ref{ex:c9q4}：圆的隐函数导数]\phantomsection\label{ans:c9q4}
因 $P$ 处 $y=5\ne0$，可以在附近把曲线写成可微函数。对原式求导，
\[
2x+2yy^{{\prime}}=0,\qquad y^{{\prime}}=-x/y.
\]
代入得 $y^{\prime}(2)=-\frac{2}{5}$，切线为
\[
2(x-2)+5(y-5)=0.
\]
再求导得到 $1+(y^{\prime})^2+yy^{\prime\prime}=0$，所以
\[
y^{\prime\prime}=-\frac{x^2+y^2}{y^3}=-\frac{29}{y^3},\qquad y^{\prime\prime}(2)=-\frac{29}{125}.
\]
切线公式也可由半径与切线垂直得到，这提供一个独立的几何检验。
\end{solve}

\begin{solve}[习题~\ref{ex:c9q5}：圆的隐函数导数]\phantomsection\label{ans:c9q5}
因 $P$ 处 $y=3\ne0$，可以在附近把曲线写成可微函数。对原式求导，
\[
2x+2yy^{{\prime}}=0,\qquad y^{{\prime}}=-x/y.
\]
代入得 $y^{\prime}(4)=-\frac{4}{3}$，切线为
\[
4(x-4)+3(y-3)=0.
\]
再求导得到 $1+(y^{\prime})^2+yy^{\prime\prime}=0$，所以
\[
y^{\prime\prime}=-\frac{x^2+y^2}{y^3}=-\frac{25}{y^3},\qquad y^{\prime\prime}(4)=-\frac{25}{27}.
\]
切线公式也可由半径与切线垂直得到，这提供一个独立的几何检验。
\end{solve}

\begin{solve}[习题~\ref{ex:c9q6}：圆的隐函数导数]\phantomsection\label{ans:c9q6}
因 $P$ 处 $y=5\ne0$，可以在附近把曲线写成可微函数。对原式求导，
\[
2x+2yy^{{\prime}}=0,\qquad y^{{\prime}}=-x/y.
\]
代入得 $y^{\prime}(3)=-\frac{3}{5}$，切线为
\[
3(x-3)+5(y-5)=0.
\]
再求导得到 $1+(y^{\prime})^2+yy^{\prime\prime}=0$，所以
\[
y^{\prime\prime}=-\frac{x^2+y^2}{y^3}=-\frac{34}{y^3},\qquad y^{\prime\prime}(3)=-\frac{34}{125}.
\]
切线公式也可由半径与切线垂直得到，这提供一个独立的几何检验。
\end{solve}

\begin{solve}[习题~\ref{ex:c9q7}：幂参数曲线]\phantomsection\label{ans:c9q7}
先求参数导数 $x^{\prime}(t)=2t$、$y^{\prime}(t)=6t^2$。当 $t\ne0$ 时，
\[
\frac{dy}{dx}=\frac{6t^2}{2t}=3t.
\]
二阶导数必须先对 $t$ 求导，再除以 $dx/dt$，不能直接把 $y^{\prime\prime}(t)$ 与 $x^{\prime\prime}(t)$ 相除：
\[
\frac{d^2y}{dx^2}=\frac{d(dy/dx)/dt}{dx/dt}=\frac{\frac{3}{2}}t.
\]
在 $t=1$ 处，点为 $(1,2)$，切线为 $y-2=3(x-1)$，二阶导数为 $\frac{3}{2}$。
\end{solve}

\begin{solve}[习题~\ref{ex:c9q8}：幂参数曲线]\phantomsection\label{ans:c9q8}
先求参数导数 $x^{\prime}(t)=4t$、$y^{\prime}(t)=3t^2$。当 $t\ne0$ 时，
\[
\frac{dy}{dx}=\frac{3t^2}{4t}=\frac{3}{4}t.
\]
二阶导数必须先对 $t$ 求导，再除以 $dx/dt$，不能直接把 $y^{\prime\prime}(t)$ 与 $x^{\prime\prime}(t)$ 相除：
\[
\frac{d^2y}{dx^2}=\frac{d(dy/dx)/dt}{dx/dt}=\frac{\frac{3}{16}}t.
\]
在 $t=1$ 处，点为 $(2,1)$，切线为 $y-1=\frac{3}{4}(x-2)$，二阶导数为 $\frac{3}{16}$。
\end{solve}

\begin{solve}[习题~\ref{ex:c9q9}：幂参数曲线]\phantomsection\label{ans:c9q9}
先求参数导数 $x^{\prime}(t)=4t$、$y^{\prime}(t)=9t^2$。当 $t\ne0$ 时，
\[
\frac{dy}{dx}=\frac{9t^2}{4t}=\frac{9}{4}t.
\]
二阶导数必须先对 $t$ 求导，再除以 $dx/dt$，不能直接把 $y^{\prime\prime}(t)$ 与 $x^{\prime\prime}(t)$ 相除：
\[
\frac{d^2y}{dx^2}=\frac{d(dy/dx)/dt}{dx/dt}=\frac{\frac{9}{16}}t.
\]
在 $t=1$ 处，点为 $(2,3)$，切线为 $y-3=\frac{9}{4}(x-2)$，二阶导数为 $\frac{9}{16}$。
\end{solve}

\begin{solve}[习题~\ref{ex:c9q10}：幂参数曲线]\phantomsection\label{ans:c9q10}
先求参数导数 $x^{\prime}(t)=6t$、$y^{\prime}(t)=6t^2$。当 $t\ne0$ 时，
\[
\frac{dy}{dx}=\frac{6t^2}{6t}=1t.
\]
二阶导数必须先对 $t$ 求导，再除以 $dx/dt$，不能直接把 $y^{\prime\prime}(t)$ 与 $x^{\prime\prime}(t)$ 相除：
\[
\frac{d^2y}{dx^2}=\frac{d(dy/dx)/dt}{dx/dt}=\frac{\frac{1}{6}}t.
\]
在 $t=1$ 处，点为 $(3,2)$，切线为 $y-2=1(x-3)$，二阶导数为 $\frac{1}{6}$。
\end{solve}

\begin{solve}[习题~\ref{ex:c9q11}：幂参数曲线]\phantomsection\label{ans:c9q11}
先求参数导数 $x^{\prime}(t)=8t$、$y^{\prime}(t)=3t^2$。当 $t\ne0$ 时，
\[
\frac{dy}{dx}=\frac{3t^2}{8t}=\frac{3}{8}t.
\]
二阶导数必须先对 $t$ 求导，再除以 $dx/dt$，不能直接把 $y^{\prime\prime}(t)$ 与 $x^{\prime\prime}(t)$ 相除：
\[
\frac{d^2y}{dx^2}=\frac{d(dy/dx)/dt}{dx/dt}=\frac{\frac{3}{64}}t.
\]
在 $t=1$ 处，点为 $(4,1)$，切线为 $y-1=\frac{3}{8}(x-4)$，二阶导数为 $\frac{3}{64}$。
\end{solve}

\begin{solve}[习题~\ref{ex:c9q12}：幂参数曲线]\phantomsection\label{ans:c9q12}
先求参数导数 $x^{\prime}(t)=2t$、$y^{\prime}(t)=12t^2$。当 $t\ne0$ 时，
\[
\frac{dy}{dx}=\frac{12t^2}{2t}=6t.
\]
二阶导数必须先对 $t$ 求导，再除以 $dx/dt$，不能直接把 $y^{\prime\prime}(t)$ 与 $x^{\prime\prime}(t)$ 相除：
\[
\frac{d^2y}{dx^2}=\frac{d(dy/dx)/dt}{dx/dt}=\frac{3}t.
\]
在 $t=1$ 处，点为 $(1,4)$，切线为 $y-4=6(x-1)$，二阶导数为 $3$。
\end{solve}

\begin{solve}[习题~\ref{ex:c9q13}：对数的高阶导数]\phantomsection\label{ans:c9q13}
一阶导数为 $f^{\prime}(x)=1/(x+1)$。归纳猜想为
\[
f^{(m)}(x)=(-1)^{m-1}\frac{(m-1)!}{(x+1)^m},\quad m\ge1.
\]
若公式对 $m$ 成立，则再求一次导数，系数乘以 $-m$，分母指数增为 $m+1$，恰好得到下一阶公式。故归纳完成。代入给定阶数与 $x=0$，
\[
f^{(3)}(0)=2.
\]
导数公式的区间始终为对数原函数的定义域；分母形式本身并不能授权把原函数扩展到真数为负的区间。
\end{solve}

\begin{solve}[习题~\ref{ex:c9q14}：对数的高阶导数]\phantomsection\label{ans:c9q14}
一阶导数为 $f^{\prime}(x)=1/(x+2)$。归纳猜想为
\[
f^{(m)}(x)=(-1)^{m-1}\frac{(m-1)!}{(x+2)^m},\quad m\ge1.
\]
若公式对 $m$ 成立，则再求一次导数，系数乘以 $-m$，分母指数增为 $m+1$，恰好得到下一阶公式。故归纳完成。代入给定阶数与 $x=0$，
\[
f^{(4)}(0)=-\frac{3}{8}.
\]
导数公式的区间始终为对数原函数的定义域；分母形式本身并不能授权把原函数扩展到真数为负的区间。
\end{solve}

\begin{solve}[习题~\ref{ex:c9q15}：对数的高阶导数]\phantomsection\label{ans:c9q15}
一阶导数为 $f^{\prime}(x)=1/(x+3)$。归纳猜想为
\[
f^{(m)}(x)=(-1)^{m-1}\frac{(m-1)!}{(x+3)^m},\quad m\ge1.
\]
若公式对 $m$ 成立，则再求一次导数，系数乘以 $-m$，分母指数增为 $m+1$，恰好得到下一阶公式。故归纳完成。代入给定阶数与 $x=0$，
\[
f^{(5)}(0)=\frac{8}{81}.
\]
导数公式的区间始终为对数原函数的定义域；分母形式本身并不能授权把原函数扩展到真数为负的区间。
\end{solve}

\begin{solve}[习题~\ref{ex:c9q16}：对数的高阶导数]\phantomsection\label{ans:c9q16}
一阶导数为 $f^{\prime}(x)=1/(x+4)$。归纳猜想为
\[
f^{(m)}(x)=(-1)^{m-1}\frac{(m-1)!}{(x+4)^m},\quad m\ge1.
\]
若公式对 $m$ 成立，则再求一次导数，系数乘以 $-m$，分母指数增为 $m+1$，恰好得到下一阶公式。故归纳完成。代入给定阶数与 $x=0$，
\[
f^{(3)}(0)=\frac{1}{32}.
\]
导数公式的区间始终为对数原函数的定义域；分母形式本身并不能授权把原函数扩展到真数为负的区间。
\end{solve}

\begin{solve}[习题~\ref{ex:c9q17}：对数的高阶导数]\phantomsection\label{ans:c9q17}
一阶导数为 $f^{\prime}(x)=1/(x+2)$。归纳猜想为
\[
f^{(m)}(x)=(-1)^{m-1}\frac{(m-1)!}{(x+2)^m},\quad m\ge1.
\]
若公式对 $m$ 成立，则再求一次导数，系数乘以 $-m$，分母指数增为 $m+1$，恰好得到下一阶公式。故归纳完成。代入给定阶数与 $x=0$，
\[
f^{(6)}(0)=-\frac{15}{8}.
\]
导数公式的区间始终为对数原函数的定义域；分母形式本身并不能授权把原函数扩展到真数为负的区间。
\end{solve}

\begin{solve}[习题~\ref{ex:c9q18}：对数的高阶导数]\phantomsection\label{ans:c9q18}
一阶导数为 $f^{\prime}(x)=1/(x+5)$。归纳猜想为
\[
f^{(m)}(x)=(-1)^{m-1}\frac{(m-1)!}{(x+5)^m},\quad m\ge1.
\]
若公式对 $m$ 成立，则再求一次导数，系数乘以 $-m$，分母指数增为 $m+1$，恰好得到下一阶公式。故归纳完成。代入给定阶数与 $x=0$，
\[
f^{(4)}(0)=-\frac{6}{625}.
\]
导数公式的区间始终为对数原函数的定义域；分母形式本身并不能授权把原函数扩展到真数为负的区间。
\end{solve}

\begin{solve}[习题~\ref{ex:c9q19}：Leibniz 公式的短和]\phantomsection\label{ans:c9q19}
多项式 $x^2$ 的三阶及更高阶导数为零，所以 Leibniz 公式只留下三项：
\[
f^{(3)}(x)=e^{1x}\left(1x^2+6x+6\right).
\]
这里三项分别对应对多项式求 $0$ 次、$1$ 次和 $2$ 次导数，组合系数为 $1,3,\binom{3}2$。因此
\[
f^{(3)}(0)=6.
\]
可以从 $x^2e^{ax}=\sum_{k=0}^\infty a^kx^{k+2}/k!$ 的系数再作验证：第 $m$ 阶导数在 $0$ 处等于 $x^m$ 项系数乘以 $m!$。
\end{solve}

\begin{solve}[习题~\ref{ex:c9q20}：Leibniz 公式的短和]\phantomsection\label{ans:c9q20}
多项式 $x^2$ 的三阶及更高阶导数为零，所以 Leibniz 公式只留下三项：
\[
f^{(4)}(x)=e^{2x}\left(16x^2+64x+48\right).
\]
这里三项分别对应对多项式求 $0$ 次、$1$ 次和 $2$ 次导数，组合系数为 $1,4,\binom{4}2$。因此
\[
f^{(4)}(0)=48.
\]
可以从 $x^2e^{ax}=\sum_{k=0}^\infty a^kx^{k+2}/k!$ 的系数再作验证：第 $m$ 阶导数在 $0$ 处等于 $x^m$ 项系数乘以 $m!$。
\end{solve}

\begin{solve}[习题~\ref{ex:c9q21}：Leibniz 公式的短和]\phantomsection\label{ans:c9q21}
多项式 $x^2$ 的三阶及更高阶导数为零，所以 Leibniz 公式只留下三项：
\[
f^{(2)}(x)=e^{3x}\left(9x^2+12x+2\right).
\]
这里三项分别对应对多项式求 $0$ 次、$1$ 次和 $2$ 次导数，组合系数为 $1,2,\binom{2}2$。因此
\[
f^{(2)}(0)=2.
\]
可以从 $x^2e^{ax}=\sum_{k=0}^\infty a^kx^{k+2}/k!$ 的系数再作验证：第 $m$ 阶导数在 $0$ 处等于 $x^m$ 项系数乘以 $m!$。
\end{solve}

\begin{solve}[习题~\ref{ex:c9q22}：Leibniz 公式的短和]\phantomsection\label{ans:c9q22}
多项式 $x^2$ 的三阶及更高阶导数为零，所以 Leibniz 公式只留下三项：
\[
f^{(5)}(x)=e^{2x}\left(32x^2+160x+160\right).
\]
这里三项分别对应对多项式求 $0$ 次、$1$ 次和 $2$ 次导数，组合系数为 $1,5,\binom{5}2$。因此
\[
f^{(5)}(0)=160.
\]
可以从 $x^2e^{ax}=\sum_{k=0}^\infty a^kx^{k+2}/k!$ 的系数再作验证：第 $m$ 阶导数在 $0$ 处等于 $x^m$ 项系数乘以 $m!$。
\end{solve}

\begin{solve}[习题~\ref{ex:c9q23}：Leibniz 公式的短和]\phantomsection\label{ans:c9q23}
多项式 $x^2$ 的三阶及更高阶导数为零，所以 Leibniz 公式只留下三项：
\[
f^{(3)}(x)=e^{4x}\left(64x^2+96x+24\right).
\]
这里三项分别对应对多项式求 $0$ 次、$1$ 次和 $2$ 次导数，组合系数为 $1,3,\binom{3}2$。因此
\[
f^{(3)}(0)=24.
\]
可以从 $x^2e^{ax}=\sum_{k=0}^\infty a^kx^{k+2}/k!$ 的系数再作验证：第 $m$ 阶导数在 $0$ 处等于 $x^m$ 项系数乘以 $m!$。
\end{solve}

\begin{solve}[习题~\ref{ex:c9q24}：Leibniz 公式的短和]\phantomsection\label{ans:c9q24}
多项式 $x^2$ 的三阶及更高阶导数为零，所以 Leibniz 公式只留下三项：
\[
f^{(6)}(x)=e^{3x}\left(729x^2+2916x+2430\right).
\]
这里三项分别对应对多项式求 $0$ 次、$1$ 次和 $2$ 次导数，组合系数为 $1,6,\binom{6}2$。因此
\[
f^{(6)}(0)=2430.
\]
可以从 $x^2e^{ax}=\sum_{k=0}^\infty a^kx^{k+2}/k!$ 的系数再作验证：第 $m$ 阶导数在 $0$ 处等于 $x^m$ 项系数乘以 $m!$。
\end{solve}

\clearpage
\section{中值定理与泰勒近似}\label{sec:ch10}
Rolle 定理和 Lagrange 中值定理连接区间上的整体变化与某一点的导数。泰勒公式把函数拆成多项式和余项；计算近似值时，误差界与多项式本身同样重要。

\textbf{资源对应：}CLP-1 \S~2.13、\S~3.4。本章共 25 道编号习题。
\subsection{习题}
\begin{example}[端点连续性不能省略]\label{ex:c10q1}
构造在 $(0,10)$ 上可导、满足 $f(0)=0$、$f(10)=10$，但在 $(0,10)$ 内处处有 $f^{\prime}=0$ 的函数。它为什么不违反中值定理？
\par\textit{选译并补充。来源：CLP-1 \S~2.13，对应端点条件反例题。}\par\hyperref[ans:c10q1]{解答见第~\pageref{ans:c10q1} 页。}
\end{example}
\begin{example}[中值点的精确计算]\label{ex:c10q2}
对 $f(x)=x^{2}$ 在 $[0,1]$ 上应用 Lagrange 中值定理，求所有满足结论的 $\xi$。
\par\textit{改编训练。}\par\hyperref[ans:c10q2]{解答见第~\pageref{ans:c10q2} 页。}
\end{example}
\begin{example}[中值点的精确计算]\label{ex:c10q3}
对 $f(x)=x^{3}$ 在 $[0,1]$ 上应用 Lagrange 中值定理，求所有满足结论的 $\xi$。
\par\textit{改编训练。}\par\hyperref[ans:c10q3]{解答见第~\pageref{ans:c10q3} 页。}
\end{example}
\begin{example}[中值点的精确计算]\label{ex:c10q4}
对 $f(x)=x^{4}$ 在 $[0,1]$ 上应用 Lagrange 中值定理，求所有满足结论的 $\xi$。
\par\textit{改编训练。}\par\hyperref[ans:c10q4]{解答见第~\pageref{ans:c10q4} 页。}
\end{example}
\begin{example}[中值点的精确计算]\label{ex:c10q5}
对 $f(x)=x^{5}$ 在 $[0,1]$ 上应用 Lagrange 中值定理，求所有满足结论的 $\xi$。
\par\textit{改编训练。}\par\hyperref[ans:c10q5]{解答见第~\pageref{ans:c10q5} 页。}
\end{example}
\begin{example}[中值点的精确计算]\label{ex:c10q6}
对 $f(x)=x^{6}$ 在 $[0,1]$ 上应用 Lagrange 中值定理，求所有满足结论的 $\xi$。
\par\textit{改编训练。}\par\hyperref[ans:c10q6]{解答见第~\pageref{ans:c10q6} 页。}
\end{example}
\begin{example}[中值点的精确计算]\label{ex:c10q7}
对 $f(x)=x^{7}$ 在 $[0,1]$ 上应用 Lagrange 中值定理，求所有满足结论的 $\xi$。
\par\textit{改编训练。}\par\hyperref[ans:c10q7]{解答见第~\pageref{ans:c10q7} 页。}
\end{example}
\begin{example}[导数界转为函数差的界]\label{ex:c10q8}
证明任意 $x,y\in\mathbb R$ 满足 $|\sin(1x)-\sin(1y)|\le 1|x-y|$，并据此证明该函数一致连续。
\par\textit{改编训练。}\par\hyperref[ans:c10q8]{解答见第~\pageref{ans:c10q8} 页。}
\end{example}
\begin{example}[导数界转为函数差的界]\label{ex:c10q9}
证明任意 $x,y\in\mathbb R$ 满足 $|\sin(2x)-\sin(2y)|\le 2|x-y|$，并据此证明该函数一致连续。
\par\textit{改编训练。}\par\hyperref[ans:c10q9]{解答见第~\pageref{ans:c10q9} 页。}
\end{example}
\begin{example}[导数界转为函数差的界]\label{ex:c10q10}
证明任意 $x,y\in\mathbb R$ 满足 $|\sin(3x)-\sin(3y)|\le 3|x-y|$，并据此证明该函数一致连续。
\par\textit{改编训练。}\par\hyperref[ans:c10q10]{解答见第~\pageref{ans:c10q10} 页。}
\end{example}
\begin{example}[导数界转为函数差的界]\label{ex:c10q11}
证明任意 $x,y\in\mathbb R$ 满足 $|\sin(4x)-\sin(4y)|\le 4|x-y|$，并据此证明该函数一致连续。
\par\textit{改编训练。}\par\hyperref[ans:c10q11]{解答见第~\pageref{ans:c10q11} 页。}
\end{example}
\begin{example}[导数界转为函数差的界]\label{ex:c10q12}
证明任意 $x,y\in\mathbb R$ 满足 $|\sin(5x)-\sin(5y)|\le 5|x-y|$，并据此证明该函数一致连续。
\par\textit{改编训练。}\par\hyperref[ans:c10q12]{解答见第~\pageref{ans:c10q12} 页。}
\end{example}
\begin{example}[导数界转为函数差的界]\label{ex:c10q13}
证明任意 $x,y\in\mathbb R$ 满足 $|\sin(6x)-\sin(6y)|\le 6|x-y|$，并据此证明该函数一致连续。
\par\textit{改编训练。}\par\hyperref[ans:c10q13]{解答见第~\pageref{ans:c10q13} 页。}
\end{example}
\begin{example}[指数函数的泰勒多项式与余项]\label{ex:c10q14}
求 $e^{1x}$ 在 $0$ 处的 $3$ 次泰勒多项式，并给出 $|x|\le1/10$ 时的余项上界。
\par\textit{改编训练。}\par\hyperref[ans:c10q14]{解答见第~\pageref{ans:c10q14} 页。}
\end{example}
\begin{example}[指数函数的泰勒多项式与余项]\label{ex:c10q15}
求 $e^{2x}$ 在 $0$ 处的 $3$ 次泰勒多项式，并给出 $|x|\le1/10$ 时的余项上界。
\par\textit{改编训练。}\par\hyperref[ans:c10q15]{解答见第~\pageref{ans:c10q15} 页。}
\end{example}
\begin{example}[指数函数的泰勒多项式与余项]\label{ex:c10q16}
求 $e^{1x}$ 在 $0$ 处的 $4$ 次泰勒多项式，并给出 $|x|\le1/10$ 时的余项上界。
\par\textit{改编训练。}\par\hyperref[ans:c10q16]{解答见第~\pageref{ans:c10q16} 页。}
\end{example}
\begin{example}[指数函数的泰勒多项式与余项]\label{ex:c10q17}
求 $e^{2x}$ 在 $0$ 处的 $4$ 次泰勒多项式，并给出 $|x|\le1/10$ 时的余项上界。
\par\textit{改编训练。}\par\hyperref[ans:c10q17]{解答见第~\pageref{ans:c10q17} 页。}
\end{example}
\begin{example}[指数函数的泰勒多项式与余项]\label{ex:c10q18}
求 $e^{3x}$ 在 $0$ 处的 $3$ 次泰勒多项式，并给出 $|x|\le1/10$ 时的余项上界。
\par\textit{改编训练。}\par\hyperref[ans:c10q18]{解答见第~\pageref{ans:c10q18} 页。}
\end{example}
\begin{example}[指数函数的泰勒多项式与余项]\label{ex:c10q19}
求 $e^{1x}$ 在 $0$ 处的 $5$ 次泰勒多项式，并给出 $|x|\le1/10$ 时的余项上界。
\par\textit{改编训练。}\par\hyperref[ans:c10q19]{解答见第~\pageref{ans:c10q19} 页。}
\end{example}
\begin{example}[由泰勒系数恢复导数]\label{ex:c10q20}
某函数在 $0$ 处的三次泰勒多项式为 $1+2x^2+\frac{1}{2}x^3$。求其前四个 Taylor 数据 $f(0),f^{\prime}(0),f^{\prime\prime}(0),f^{(3)}(0)$。
\par\textit{改编训练。来源：CLP-1 \S~3.4，2015Q 泰勒系数题的参数改编。}\par\hyperref[ans:c10q20]{解答见第~\pageref{ans:c10q20} 页。}
\end{example}
\begin{example}[由泰勒系数恢复导数]\label{ex:c10q21}
某函数在 $0$ 处的三次泰勒多项式为 $2+3x^2+1x^3$。求其前四个 Taylor 数据 $f(0),f^{\prime}(0),f^{\prime\prime}(0),f^{(3)}(0)$。
\par\textit{改编训练。来源：CLP-1 \S~3.4，2015Q 泰勒系数题的参数改编。}\par\hyperref[ans:c10q21]{解答见第~\pageref{ans:c10q21} 页。}
\end{example}
\begin{example}[由泰勒系数恢复导数]\label{ex:c10q22}
某函数在 $0$ 处的三次泰勒多项式为 $3+4x^2+\frac{3}{2}x^3$。求其前四个 Taylor 数据 $f(0),f^{\prime}(0),f^{\prime\prime}(0),f^{(3)}(0)$。
\par\textit{改编训练。来源：CLP-1 \S~3.4，2015Q 泰勒系数题的参数改编。}\par\hyperref[ans:c10q22]{解答见第~\pageref{ans:c10q22} 页。}
\end{example}
\begin{example}[由泰勒系数恢复导数]\label{ex:c10q23}
某函数在 $0$ 处的三次泰勒多项式为 $4+5x^2+2x^3$。求其前四个 Taylor 数据 $f(0),f^{\prime}(0),f^{\prime\prime}(0),f^{(3)}(0)$。
\par\textit{改编训练。来源：CLP-1 \S~3.4，2015Q 泰勒系数题的参数改编。}\par\hyperref[ans:c10q23]{解答见第~\pageref{ans:c10q23} 页。}
\end{example}
\begin{example}[由泰勒系数恢复导数]\label{ex:c10q24}
某函数在 $0$ 处的三次泰勒多项式为 $5+6x^2+\frac{5}{2}x^3$。求其前四个 Taylor 数据 $f(0),f^{\prime}(0),f^{\prime\prime}(0),f^{(3)}(0)$。
\par\textit{改编训练。来源：CLP-1 \S~3.4，2015Q 泰勒系数题的参数改编。}\par\hyperref[ans:c10q24]{解答见第~\pageref{ans:c10q24} 页。}
\end{example}
\begin{example}[由泰勒系数恢复导数]\label{ex:c10q25}
某函数在 $0$ 处的三次泰勒多项式为 $6+7x^2+3x^3$。求其前四个 Taylor 数据 $f(0),f^{\prime}(0),f^{\prime\prime}(0),f^{(3)}(0)$。
\par\textit{改编训练。来源：CLP-1 \S~3.4，2015Q 泰勒系数题的参数改编。}\par\hyperref[ans:c10q25]{解答见第~\pageref{ans:c10q25} 页。}
\end{example}
\subsection{解答}
\begin{solve}[习题~\ref{ex:c10q1}：端点连续性不能省略]\phantomsection\label{ans:c10q1}
取 $f(x)=0$（$0\le x<10$），并定义 $f(10)=10$。它在开区间内是常函数，所以导数恒为零，但两端差商为 $1$。在 $x=10$ 处，左极限为 $0$ 而函数值为 $10$，故不满足闭区间上的连续性假设。中值定理要求闭区间连续、开区间可导，少一个条件就可能失效。
\end{solve}

\begin{solve}[习题~\ref{ex:c10q2}：中值点的精确计算]\phantomsection\label{ans:c10q2}
幂函数在闭区间连续、开区间可导，所需假设均满足。区间两端差商为 $1$，故中值点应满足
\[
2\xi^{1}=\frac{f(1)-f(0)}{1-0}=1.
\]
因此 $\xi=(\frac{1}{2})^{1/1}$。此数严格介于 $0$ 与 $1$ 之间；而函数 $x\mapsto 2x^{1}$ 在 $(0,1)$ 上严格递增，所以这是唯一中值点。中值定理本身只保证至少存在一个点，唯一性需要这里额外使用的单调性。
\end{solve}

\begin{solve}[习题~\ref{ex:c10q3}：中值点的精确计算]\phantomsection\label{ans:c10q3}
幂函数在闭区间连续、开区间可导，所需假设均满足。区间两端差商为 $1$，故中值点应满足
\[
3\xi^{2}=\frac{f(1)-f(0)}{1-0}=1.
\]
因此 $\xi=(\frac{1}{3})^{1/2}$。此数严格介于 $0$ 与 $1$ 之间；而函数 $x\mapsto 3x^{2}$ 在 $(0,1)$ 上严格递增，所以这是唯一中值点。中值定理本身只保证至少存在一个点，唯一性需要这里额外使用的单调性。
\end{solve}

\begin{solve}[习题~\ref{ex:c10q4}：中值点的精确计算]\phantomsection\label{ans:c10q4}
幂函数在闭区间连续、开区间可导，所需假设均满足。区间两端差商为 $1$，故中值点应满足
\[
4\xi^{3}=\frac{f(1)-f(0)}{1-0}=1.
\]
因此 $\xi=(\frac{1}{4})^{1/3}$。此数严格介于 $0$ 与 $1$ 之间；而函数 $x\mapsto 4x^{3}$ 在 $(0,1)$ 上严格递增，所以这是唯一中值点。中值定理本身只保证至少存在一个点，唯一性需要这里额外使用的单调性。
\end{solve}

\begin{solve}[习题~\ref{ex:c10q5}：中值点的精确计算]\phantomsection\label{ans:c10q5}
幂函数在闭区间连续、开区间可导，所需假设均满足。区间两端差商为 $1$，故中值点应满足
\[
5\xi^{4}=\frac{f(1)-f(0)}{1-0}=1.
\]
因此 $\xi=(\frac{1}{5})^{1/4}$。此数严格介于 $0$ 与 $1$ 之间；而函数 $x\mapsto 5x^{4}$ 在 $(0,1)$ 上严格递增，所以这是唯一中值点。中值定理本身只保证至少存在一个点，唯一性需要这里额外使用的单调性。
\end{solve}

\begin{solve}[习题~\ref{ex:c10q6}：中值点的精确计算]\phantomsection\label{ans:c10q6}
幂函数在闭区间连续、开区间可导，所需假设均满足。区间两端差商为 $1$，故中值点应满足
\[
6\xi^{5}=\frac{f(1)-f(0)}{1-0}=1.
\]
因此 $\xi=(\frac{1}{6})^{1/5}$。此数严格介于 $0$ 与 $1$ 之间；而函数 $x\mapsto 6x^{5}$ 在 $(0,1)$ 上严格递增，所以这是唯一中值点。中值定理本身只保证至少存在一个点，唯一性需要这里额外使用的单调性。
\end{solve}

\begin{solve}[习题~\ref{ex:c10q7}：中值点的精确计算]\phantomsection\label{ans:c10q7}
幂函数在闭区间连续、开区间可导，所需假设均满足。区间两端差商为 $1$，故中值点应满足
\[
7\xi^{6}=\frac{f(1)-f(0)}{1-0}=1.
\]
因此 $\xi=(\frac{1}{7})^{1/6}$。此数严格介于 $0$ 与 $1$ 之间；而函数 $x\mapsto 7x^{6}$ 在 $(0,1)$ 上严格递增，所以这是唯一中值点。中值定理本身只保证至少存在一个点，唯一性需要这里额外使用的单调性。
\end{solve}

\begin{solve}[习题~\ref{ex:c10q8}：导数界转为函数差的界]\phantomsection\label{ans:c10q8}
若 $x=y$ 结论成立。否则对两点之间的闭区间使用中值定理，存在其中的 $\xi$，使
\[
\sin(1x)-\sin(1y)=1\cos(1\xi)(x-y).
\]
由 $|\cos(1\xi)|\le1$ 得所需估计。给定 $\varepsilon>0$，取 $\delta=\varepsilon/1$，则 $|x-y|<\delta$ 蕴含函数差小于 $\varepsilon$。该估计是 Lipschitz 条件；它比一致连续更强，而且常数与 $x,y$ 无关。证明的关键是导数在整个定义域上的统一上界。
\end{solve}

\begin{solve}[习题~\ref{ex:c10q9}：导数界转为函数差的界]\phantomsection\label{ans:c10q9}
若 $x=y$ 结论成立。否则对两点之间的闭区间使用中值定理，存在其中的 $\xi$，使
\[
\sin(2x)-\sin(2y)=2\cos(2\xi)(x-y).
\]
由 $|\cos(2\xi)|\le1$ 得所需估计。给定 $\varepsilon>0$，取 $\delta=\varepsilon/2$，则 $|x-y|<\delta$ 蕴含函数差小于 $\varepsilon$。该估计是 Lipschitz 条件；它比一致连续更强，而且常数与 $x,y$ 无关。证明的关键是导数在整个定义域上的统一上界。
\end{solve}

\begin{solve}[习题~\ref{ex:c10q10}：导数界转为函数差的界]\phantomsection\label{ans:c10q10}
若 $x=y$ 结论成立。否则对两点之间的闭区间使用中值定理，存在其中的 $\xi$，使
\[
\sin(3x)-\sin(3y)=3\cos(3\xi)(x-y).
\]
由 $|\cos(3\xi)|\le1$ 得所需估计。给定 $\varepsilon>0$，取 $\delta=\varepsilon/3$，则 $|x-y|<\delta$ 蕴含函数差小于 $\varepsilon$。该估计是 Lipschitz 条件；它比一致连续更强，而且常数与 $x,y$ 无关。证明的关键是导数在整个定义域上的统一上界。
\end{solve}

\begin{solve}[习题~\ref{ex:c10q11}：导数界转为函数差的界]\phantomsection\label{ans:c10q11}
若 $x=y$ 结论成立。否则对两点之间的闭区间使用中值定理，存在其中的 $\xi$，使
\[
\sin(4x)-\sin(4y)=4\cos(4\xi)(x-y).
\]
由 $|\cos(4\xi)|\le1$ 得所需估计。给定 $\varepsilon>0$，取 $\delta=\varepsilon/4$，则 $|x-y|<\delta$ 蕴含函数差小于 $\varepsilon$。该估计是 Lipschitz 条件；它比一致连续更强，而且常数与 $x,y$ 无关。证明的关键是导数在整个定义域上的统一上界。
\end{solve}

\begin{solve}[习题~\ref{ex:c10q12}：导数界转为函数差的界]\phantomsection\label{ans:c10q12}
若 $x=y$ 结论成立。否则对两点之间的闭区间使用中值定理，存在其中的 $\xi$，使
\[
\sin(5x)-\sin(5y)=5\cos(5\xi)(x-y).
\]
由 $|\cos(5\xi)|\le1$ 得所需估计。给定 $\varepsilon>0$，取 $\delta=\varepsilon/5$，则 $|x-y|<\delta$ 蕴含函数差小于 $\varepsilon$。该估计是 Lipschitz 条件；它比一致连续更强，而且常数与 $x,y$ 无关。证明的关键是导数在整个定义域上的统一上界。
\end{solve}

\begin{solve}[习题~\ref{ex:c10q13}：导数界转为函数差的界]\phantomsection\label{ans:c10q13}
若 $x=y$ 结论成立。否则对两点之间的闭区间使用中值定理，存在其中的 $\xi$，使
\[
\sin(6x)-\sin(6y)=6\cos(6\xi)(x-y).
\]
由 $|\cos(6\xi)|\le1$ 得所需估计。给定 $\varepsilon>0$，取 $\delta=\varepsilon/6$，则 $|x-y|<\delta$ 蕴含函数差小于 $\varepsilon$。该估计是 Lipschitz 条件；它比一致连续更强，而且常数与 $x,y$ 无关。证明的关键是导数在整个定义域上的统一上界。
\end{solve}

\begin{solve}[习题~\ref{ex:c10q14}：指数函数的泰勒多项式与余项]\phantomsection\label{ans:c10q14}
函数的第 $k$ 阶导数在 $0$ 处为 $1^k$，因此
\[
T_{3}(x)=\frac{1}{6}x^{3}+\frac{1}{2}x^{2}+x+1.
\]
由 Lagrange 余项，某个介于 $0$ 与 $x$ 之间的 $\xi$ 满足
\[
R_{3}(x)=\frac{1e^{1\xi}}{24}x^{4}.
\]
当 $|x|\le1/10$ 时，$e^{1\xi}\le e^{1/10}$，故
\[
|R_{3}(x)|\le \frac{1}{240000}e^{1/10}.
\]
这里的上界同时控制整个给定区间；不能把余项中的未知中间点随意替换成 $0$。
\end{solve}

\begin{solve}[习题~\ref{ex:c10q15}：指数函数的泰勒多项式与余项]\phantomsection\label{ans:c10q15}
函数的第 $k$ 阶导数在 $0$ 处为 $2^k$，因此
\[
T_{3}(x)=\frac{4}{3}x^{3}+2x^{2}+2x+1.
\]
由 Lagrange 余项，某个介于 $0$ 与 $x$ 之间的 $\xi$ 满足
\[
R_{3}(x)=\frac{16e^{2\xi}}{24}x^{4}.
\]
当 $|x|\le1/10$ 时，$e^{2\xi}\le e^{2/10}$，故
\[
|R_{3}(x)|\le \frac{1}{15000}e^{2/10}.
\]
这里的上界同时控制整个给定区间；不能把余项中的未知中间点随意替换成 $0$。
\end{solve}

\begin{solve}[习题~\ref{ex:c10q16}：指数函数的泰勒多项式与余项]\phantomsection\label{ans:c10q16}
函数的第 $k$ 阶导数在 $0$ 处为 $1^k$，因此
\[
T_{4}(x)=\frac{1}{24}x^{4}+\frac{1}{6}x^{3}+\frac{1}{2}x^{2}+x+1.
\]
由 Lagrange 余项，某个介于 $0$ 与 $x$ 之间的 $\xi$ 满足
\[
R_{4}(x)=\frac{1e^{1\xi}}{120}x^{5}.
\]
当 $|x|\le1/10$ 时，$e^{1\xi}\le e^{1/10}$，故
\[
|R_{4}(x)|\le \frac{1}{12000000}e^{1/10}.
\]
这里的上界同时控制整个给定区间；不能把余项中的未知中间点随意替换成 $0$。
\end{solve}

\begin{solve}[习题~\ref{ex:c10q17}：指数函数的泰勒多项式与余项]\phantomsection\label{ans:c10q17}
函数的第 $k$ 阶导数在 $0$ 处为 $2^k$，因此
\[
T_{4}(x)=\frac{2}{3}x^{4}+\frac{4}{3}x^{3}+2x^{2}+2x+1.
\]
由 Lagrange 余项，某个介于 $0$ 与 $x$ 之间的 $\xi$ 满足
\[
R_{4}(x)=\frac{32e^{2\xi}}{120}x^{5}.
\]
当 $|x|\le1/10$ 时，$e^{2\xi}\le e^{2/10}$，故
\[
|R_{4}(x)|\le \frac{1}{375000}e^{2/10}.
\]
这里的上界同时控制整个给定区间；不能把余项中的未知中间点随意替换成 $0$。
\end{solve}

\begin{solve}[习题~\ref{ex:c10q18}：指数函数的泰勒多项式与余项]\phantomsection\label{ans:c10q18}
函数的第 $k$ 阶导数在 $0$ 处为 $3^k$，因此
\[
T_{3}(x)=\frac{9}{2}x^{3}+\frac{9}{2}x^{2}+3x+1.
\]
由 Lagrange 余项，某个介于 $0$ 与 $x$ 之间的 $\xi$ 满足
\[
R_{3}(x)=\frac{81e^{3\xi}}{24}x^{4}.
\]
当 $|x|\le1/10$ 时，$e^{3\xi}\le e^{3/10}$，故
\[
|R_{3}(x)|\le \frac{27}{80000}e^{3/10}.
\]
这里的上界同时控制整个给定区间；不能把余项中的未知中间点随意替换成 $0$。
\end{solve}

\begin{solve}[习题~\ref{ex:c10q19}：指数函数的泰勒多项式与余项]\phantomsection\label{ans:c10q19}
函数的第 $k$ 阶导数在 $0$ 处为 $1^k$，因此
\[
T_{5}(x)=\frac{1}{120}x^{5}+\frac{1}{24}x^{4}+\frac{1}{6}x^{3}+\frac{1}{2}x^{2}+x+1.
\]
由 Lagrange 余项，某个介于 $0$ 与 $x$ 之间的 $\xi$ 满足
\[
R_{5}(x)=\frac{1e^{1\xi}}{720}x^{6}.
\]
当 $|x|\le1/10$ 时，$e^{1\xi}\le e^{1/10}$，故
\[
|R_{5}(x)|\le \frac{1}{720000000}e^{1/10}.
\]
这里的上界同时控制整个给定区间；不能把余项中的未知中间点随意替换成 $0$。
\end{solve}

\begin{solve}[习题~\ref{ex:c10q20}：由泰勒系数恢复导数]\phantomsection\label{ans:c10q20}
泰勒多项式的定义为
\[
T_3(x)=f(0)+f^{\prime}(0)x+\frac{f^{\prime\prime}(0)}2x^2+\frac{f^{(3)}(0)}6x^3.
\]
逐项比较系数得
\[
f(0)=1,\quad f^{\prime}(0)=0,\quad f^{\prime\prime}(0)=4,\quad f^{(3)}(0)=3.
\]
多项式中没有一次项，意味着该阶系数为零；二阶、三阶导数则要分别乘以 $2!$、$3!$。仅知道三次泰勒多项式，并不能确定函数在其他点的精确取值，也不能给出没有附加条件的余项数值界。
\end{solve}

\begin{solve}[习题~\ref{ex:c10q21}：由泰勒系数恢复导数]\phantomsection\label{ans:c10q21}
泰勒多项式的定义为
\[
T_3(x)=f(0)+f^{\prime}(0)x+\frac{f^{\prime\prime}(0)}2x^2+\frac{f^{(3)}(0)}6x^3.
\]
逐项比较系数得
\[
f(0)=2,\quad f^{\prime}(0)=0,\quad f^{\prime\prime}(0)=6,\quad f^{(3)}(0)=6.
\]
多项式中没有一次项，意味着该阶系数为零；二阶、三阶导数则要分别乘以 $2!$、$3!$。仅知道三次泰勒多项式，并不能确定函数在其他点的精确取值，也不能给出没有附加条件的余项数值界。
\end{solve}

\begin{solve}[习题~\ref{ex:c10q22}：由泰勒系数恢复导数]\phantomsection\label{ans:c10q22}
泰勒多项式的定义为
\[
T_3(x)=f(0)+f^{\prime}(0)x+\frac{f^{\prime\prime}(0)}2x^2+\frac{f^{(3)}(0)}6x^3.
\]
逐项比较系数得
\[
f(0)=3,\quad f^{\prime}(0)=0,\quad f^{\prime\prime}(0)=8,\quad f^{(3)}(0)=9.
\]
多项式中没有一次项，意味着该阶系数为零；二阶、三阶导数则要分别乘以 $2!$、$3!$。仅知道三次泰勒多项式，并不能确定函数在其他点的精确取值，也不能给出没有附加条件的余项数值界。
\end{solve}

\begin{solve}[习题~\ref{ex:c10q23}：由泰勒系数恢复导数]\phantomsection\label{ans:c10q23}
泰勒多项式的定义为
\[
T_3(x)=f(0)+f^{\prime}(0)x+\frac{f^{\prime\prime}(0)}2x^2+\frac{f^{(3)}(0)}6x^3.
\]
逐项比较系数得
\[
f(0)=4,\quad f^{\prime}(0)=0,\quad f^{\prime\prime}(0)=10,\quad f^{(3)}(0)=12.
\]
多项式中没有一次项，意味着该阶系数为零；二阶、三阶导数则要分别乘以 $2!$、$3!$。仅知道三次泰勒多项式，并不能确定函数在其他点的精确取值，也不能给出没有附加条件的余项数值界。
\end{solve}

\begin{solve}[习题~\ref{ex:c10q24}：由泰勒系数恢复导数]\phantomsection\label{ans:c10q24}
泰勒多项式的定义为
\[
T_3(x)=f(0)+f^{\prime}(0)x+\frac{f^{\prime\prime}(0)}2x^2+\frac{f^{(3)}(0)}6x^3.
\]
逐项比较系数得
\[
f(0)=5,\quad f^{\prime}(0)=0,\quad f^{\prime\prime}(0)=12,\quad f^{(3)}(0)=15.
\]
多项式中没有一次项，意味着该阶系数为零；二阶、三阶导数则要分别乘以 $2!$、$3!$。仅知道三次泰勒多项式，并不能确定函数在其他点的精确取值，也不能给出没有附加条件的余项数值界。
\end{solve}

\begin{solve}[习题~\ref{ex:c10q25}：由泰勒系数恢复导数]\phantomsection\label{ans:c10q25}
泰勒多项式的定义为
\[
T_3(x)=f(0)+f^{\prime}(0)x+\frac{f^{\prime\prime}(0)}2x^2+\frac{f^{(3)}(0)}6x^3.
\]
逐项比较系数得
\[
f(0)=6,\quad f^{\prime}(0)=0,\quad f^{\prime\prime}(0)=14,\quad f^{(3)}(0)=18.
\]
多项式中没有一次项，意味着该阶系数为零；二阶、三阶导数则要分别乘以 $2!$、$3!$。仅知道三次泰勒多项式，并不能确定函数在其他点的精确取值，也不能给出没有附加条件的余项数值界。
\end{solve}

\clearpage
\section{导数的应用与最优化}\label{sec:ch11}
求最值时，应先写出定义域与约束，再寻找驻点、不可导点和边界候选点。局部极值与全局最值不同；某点的导数为零也不自动保证该点是极值点。

\textbf{资源对应：}CLP-1 \S~3.2、\S~3.5、\S~3.6。本章共 25 道编号习题。
\subsection{习题}
\begin{example}[闭区间最值的候选点]\label{ex:c11q1}
求 $f(x)=x^5-5x+2$ 在 $[-2,0]$ 上的最大值和最小值。
\par\textit{选译。来源：CLP-1 \S~3.5.3，2015Q，标签 s3.5.3given1。}\par\hyperref[ans:c11q1]{解答见第~\pageref{ans:c11q1} 页。}
\end{example}
\begin{example}[三次函数的极值与凹凸]\label{ex:c11q2}
讨论 $f(x)=x^3-3x$ 的单调区间、局部极值和拐点。
\par\textit{改编训练。}\par\hyperref[ans:c11q2]{解答见第~\pageref{ans:c11q2} 页。}
\end{example}
\begin{example}[三次函数的极值与凹凸]\label{ex:c11q3}
讨论 $f(x)=x^3-12x$ 的单调区间、局部极值和拐点。
\par\textit{改编训练。}\par\hyperref[ans:c11q3]{解答见第~\pageref{ans:c11q3} 页。}
\end{example}
\begin{example}[三次函数的极值与凹凸]\label{ex:c11q4}
讨论 $f(x)=x^3-27x$ 的单调区间、局部极值和拐点。
\par\textit{改编训练。}\par\hyperref[ans:c11q4]{解答见第~\pageref{ans:c11q4} 页。}
\end{example}
\begin{example}[三次函数的极值与凹凸]\label{ex:c11q5}
讨论 $f(x)=x^3-48x$ 的单调区间、局部极值和拐点。
\par\textit{改编训练。}\par\hyperref[ans:c11q5]{解答见第~\pageref{ans:c11q5} 页。}
\end{example}
\begin{example}[三次函数的极值与凹凸]\label{ex:c11q6}
讨论 $f(x)=x^3-75x$ 的单调区间、局部极值和拐点。
\par\textit{改编训练。}\par\hyperref[ans:c11q6]{解答见第~\pageref{ans:c11q6} 页。}
\end{example}
\begin{example}[三次函数的极值与凹凸]\label{ex:c11q7}
讨论 $f(x)=x^3-108x$ 的单调区间、局部极值和拐点。
\par\textit{改编训练。}\par\hyperref[ans:c11q7]{解答见第~\pageref{ans:c11q7} 页。}
\end{example}
\begin{example}[固定周长矩形的最大面积]\label{ex:c11q8}
矩形周长为 $4$。设一边为 $x$，求面积的最大值，写出可行区间并证明最大值能够取得。
\par\textit{改编训练。}\par\hyperref[ans:c11q8]{解答见第~\pageref{ans:c11q8} 页。}
\end{example}
\begin{example}[固定周长矩形的最大面积]\label{ex:c11q9}
矩形周长为 $8$。设一边为 $x$，求面积的最大值，写出可行区间并证明最大值能够取得。
\par\textit{改编训练。}\par\hyperref[ans:c11q9]{解答见第~\pageref{ans:c11q9} 页。}
\end{example}
\begin{example}[固定周长矩形的最大面积]\label{ex:c11q10}
矩形周长为 $12$。设一边为 $x$，求面积的最大值，写出可行区间并证明最大值能够取得。
\par\textit{改编训练。}\par\hyperref[ans:c11q10]{解答见第~\pageref{ans:c11q10} 页。}
\end{example}
\begin{example}[固定周长矩形的最大面积]\label{ex:c11q11}
矩形周长为 $16$。设一边为 $x$，求面积的最大值，写出可行区间并证明最大值能够取得。
\par\textit{改编训练。}\par\hyperref[ans:c11q11]{解答见第~\pageref{ans:c11q11} 页。}
\end{example}
\begin{example}[固定周长矩形的最大面积]\label{ex:c11q12}
矩形周长为 $20$。设一边为 $x$，求面积的最大值，写出可行区间并证明最大值能够取得。
\par\textit{改编训练。}\par\hyperref[ans:c11q12]{解答见第~\pageref{ans:c11q12} 页。}
\end{example}
\begin{example}[固定周长矩形的最大面积]\label{ex:c11q13}
矩形周长为 $24$。设一边为 $x$，求面积的最大值，写出可行区间并证明最大值能够取得。
\par\textit{改编训练。}\par\hyperref[ans:c11q13]{解答见第~\pageref{ans:c11q13} 页。}
\end{example}
\begin{example}[指数衰减与内点最大值]\label{ex:c11q14}
求 $f(x)=xe^{-1x}$ 在 $[0,\infty)$ 上的最大值，并判断是否有最小值。
\par\textit{改编训练。}\par\hyperref[ans:c11q14]{解答见第~\pageref{ans:c11q14} 页。}
\end{example}
\begin{example}[指数衰减与内点最大值]\label{ex:c11q15}
求 $f(x)=xe^{-2x}$ 在 $[0,\infty)$ 上的最大值，并判断是否有最小值。
\par\textit{改编训练。}\par\hyperref[ans:c11q15]{解答见第~\pageref{ans:c11q15} 页。}
\end{example}
\begin{example}[指数衰减与内点最大值]\label{ex:c11q16}
求 $f(x)=xe^{-3x}$ 在 $[0,\infty)$ 上的最大值，并判断是否有最小值。
\par\textit{改编训练。}\par\hyperref[ans:c11q16]{解答见第~\pageref{ans:c11q16} 页。}
\end{example}
\begin{example}[指数衰减与内点最大值]\label{ex:c11q17}
求 $f(x)=xe^{-4x}$ 在 $[0,\infty)$ 上的最大值，并判断是否有最小值。
\par\textit{改编训练。}\par\hyperref[ans:c11q17]{解答见第~\pageref{ans:c11q17} 页。}
\end{example}
\begin{example}[指数衰减与内点最大值]\label{ex:c11q18}
求 $f(x)=xe^{-5x}$ 在 $[0,\infty)$ 上的最大值，并判断是否有最小值。
\par\textit{改编训练。}\par\hyperref[ans:c11q18]{解答见第~\pageref{ans:c11q18} 页。}
\end{example}
\begin{example}[指数衰减与内点最大值]\label{ex:c11q19}
求 $f(x)=xe^{-6x}$ 在 $[0,\infty)$ 上的最大值，并判断是否有最小值。
\par\textit{改编训练。}\par\hyperref[ans:c11q19]{解答见第~\pageref{ans:c11q19} 页。}
\end{example}
\begin{example}[相关变化率中的链式法则]\label{ex:c11q20}
球的半径以每秒 $1$ 厘米的速度增加。求半径为 $2$ 厘米时体积和表面积的变化率。
\par\textit{改编训练。}\par\hyperref[ans:c11q20]{解答见第~\pageref{ans:c11q20} 页。}
\end{example}
\begin{example}[相关变化率中的链式法则]\label{ex:c11q21}
球的半径以每秒 $2$ 厘米的速度增加。求半径为 $3$ 厘米时体积和表面积的变化率。
\par\textit{改编训练。}\par\hyperref[ans:c11q21]{解答见第~\pageref{ans:c11q21} 页。}
\end{example}
\begin{example}[相关变化率中的链式法则]\label{ex:c11q22}
球的半径以每秒 $3$ 厘米的速度增加。求半径为 $4$ 厘米时体积和表面积的变化率。
\par\textit{改编训练。}\par\hyperref[ans:c11q22]{解答见第~\pageref{ans:c11q22} 页。}
\end{example}
\begin{example}[相关变化率中的链式法则]\label{ex:c11q23}
球的半径以每秒 $4$ 厘米的速度增加。求半径为 $5$ 厘米时体积和表面积的变化率。
\par\textit{改编训练。}\par\hyperref[ans:c11q23]{解答见第~\pageref{ans:c11q23} 页。}
\end{example}
\begin{example}[相关变化率中的链式法则]\label{ex:c11q24}
球的半径以每秒 $5$ 厘米的速度增加。求半径为 $6$ 厘米时体积和表面积的变化率。
\par\textit{改编训练。}\par\hyperref[ans:c11q24]{解答见第~\pageref{ans:c11q24} 页。}
\end{example}
\begin{example}[相关变化率中的链式法则]\label{ex:c11q25}
球的半径以每秒 $6$ 厘米的速度增加。求半径为 $7$ 厘米时体积和表面积的变化率。
\par\textit{改编训练。}\par\hyperref[ans:c11q25]{解答见第~\pageref{ans:c11q25} 页。}
\end{example}
\subsection{解答}
\begin{solve}[习题~\ref{ex:c11q1}：闭区间最值的候选点]\phantomsection\label{ans:c11q1}
函数在闭区间上连续，最值一定存在。导数 $f^{\prime}(x)=5(x^4-1)$，实驻点为 $\pm1$，但只有 $-1$ 在区间内。因此只须比较 $f(-2)=-20$、$f(-1)=6$、$f(0)=2$。最大值为 $6$，在 $x=-1$ 处取得；最小值为 $-20$，在左端点取得。把区间外的驻点 $1$ 纳入比较，或者只比较驻点而遗漏端点，都会破坏解题逻辑。
\end{solve}

\begin{solve}[习题~\ref{ex:c11q2}：三次函数的极值与凹凸]\phantomsection\label{ans:c11q2}
求导并因式分解：
\[
f^{\prime}(x)=3(x-1)(x+1),\qquad f^{\prime\prime}(x)=6x.
\]
一阶导数在 $(-\infty,-1)$、$(1,\infty)$ 为正，在 $(-1,1)$ 为负，因此函数依次递增、递减、递增。在 $x=-1$ 取得局部最大值 $2$，在 $x=1$ 取得局部最小值 $-2$。二阶导数在 $0$ 两侧变号，所以拐点为 $(0,0)$。函数向两端分别趋于负无穷和正无穷，故这两个局部极值都不是全实轴上的全局最值。
\end{solve}

\begin{solve}[习题~\ref{ex:c11q3}：三次函数的极值与凹凸]\phantomsection\label{ans:c11q3}
求导并因式分解：
\[
f^{\prime}(x)=3(x-2)(x+2),\qquad f^{\prime\prime}(x)=6x.
\]
一阶导数在 $(-\infty,-2)$、$(2,\infty)$ 为正，在 $(-2,2)$ 为负，因此函数依次递增、递减、递增。在 $x=-2$ 取得局部最大值 $16$，在 $x=2$ 取得局部最小值 $-16$。二阶导数在 $0$ 两侧变号，所以拐点为 $(0,0)$。函数向两端分别趋于负无穷和正无穷，故这两个局部极值都不是全实轴上的全局最值。
\end{solve}

\begin{solve}[习题~\ref{ex:c11q4}：三次函数的极值与凹凸]\phantomsection\label{ans:c11q4}
求导并因式分解：
\[
f^{\prime}(x)=3(x-3)(x+3),\qquad f^{\prime\prime}(x)=6x.
\]
一阶导数在 $(-\infty,-3)$、$(3,\infty)$ 为正，在 $(-3,3)$ 为负，因此函数依次递增、递减、递增。在 $x=-3$ 取得局部最大值 $54$，在 $x=3$ 取得局部最小值 $-54$。二阶导数在 $0$ 两侧变号，所以拐点为 $(0,0)$。函数向两端分别趋于负无穷和正无穷，故这两个局部极值都不是全实轴上的全局最值。
\end{solve}

\begin{solve}[习题~\ref{ex:c11q5}：三次函数的极值与凹凸]\phantomsection\label{ans:c11q5}
求导并因式分解：
\[
f^{\prime}(x)=3(x-4)(x+4),\qquad f^{\prime\prime}(x)=6x.
\]
一阶导数在 $(-\infty,-4)$、$(4,\infty)$ 为正，在 $(-4,4)$ 为负，因此函数依次递增、递减、递增。在 $x=-4$ 取得局部最大值 $128$，在 $x=4$ 取得局部最小值 $-128$。二阶导数在 $0$ 两侧变号，所以拐点为 $(0,0)$。函数向两端分别趋于负无穷和正无穷，故这两个局部极值都不是全实轴上的全局最值。
\end{solve}

\begin{solve}[习题~\ref{ex:c11q6}：三次函数的极值与凹凸]\phantomsection\label{ans:c11q6}
求导并因式分解：
\[
f^{\prime}(x)=3(x-5)(x+5),\qquad f^{\prime\prime}(x)=6x.
\]
一阶导数在 $(-\infty,-5)$、$(5,\infty)$ 为正，在 $(-5,5)$ 为负，因此函数依次递增、递减、递增。在 $x=-5$ 取得局部最大值 $250$，在 $x=5$ 取得局部最小值 $-250$。二阶导数在 $0$ 两侧变号，所以拐点为 $(0,0)$。函数向两端分别趋于负无穷和正无穷，故这两个局部极值都不是全实轴上的全局最值。
\end{solve}

\begin{solve}[习题~\ref{ex:c11q7}：三次函数的极值与凹凸]\phantomsection\label{ans:c11q7}
求导并因式分解：
\[
f^{\prime}(x)=3(x-6)(x+6),\qquad f^{\prime\prime}(x)=6x.
\]
一阶导数在 $(-\infty,-6)$、$(6,\infty)$ 为正，在 $(-6,6)$ 为负，因此函数依次递增、递减、递增。在 $x=-6$ 取得局部最大值 $432$，在 $x=6$ 取得局部最小值 $-432$。二阶导数在 $0$ 两侧变号，所以拐点为 $(0,0)$。函数向两端分别趋于负无穷和正无穷，故这两个局部极值都不是全实轴上的全局最值。
\end{solve}

\begin{solve}[习题~\ref{ex:c11q8}：固定周长矩形的最大面积]\phantomsection\label{ans:c11q8}
另一边为 $2-x$，非退化矩形要求 $0<x<2$。面积函数为
\[
A(x)=x(2-x),\qquad A^{\prime}(x)=2-2x.
\]
唯一驻点是 $x=1$，在其左侧导数为正、右侧为负，故面积在这里取得全局最大值 $1$。也可配方得到
\[
A(x)=1-(x-1)^2\le 1,
\]
等号恰在正方形边长为 $1$ 时成立。可行域虽为开区间，但最大值在内部达到；本题不能仅凭闭区间最值定理来代替实际论证。
\end{solve}

\begin{solve}[习题~\ref{ex:c11q9}：固定周长矩形的最大面积]\phantomsection\label{ans:c11q9}
另一边为 $4-x$，非退化矩形要求 $0<x<4$。面积函数为
\[
A(x)=x(4-x),\qquad A^{\prime}(x)=4-2x.
\]
唯一驻点是 $x=2$，在其左侧导数为正、右侧为负，故面积在这里取得全局最大值 $4$。也可配方得到
\[
A(x)=4-(x-2)^2\le 4,
\]
等号恰在正方形边长为 $2$ 时成立。可行域虽为开区间，但最大值在内部达到；本题不能仅凭闭区间最值定理来代替实际论证。
\end{solve}

\begin{solve}[习题~\ref{ex:c11q10}：固定周长矩形的最大面积]\phantomsection\label{ans:c11q10}
另一边为 $6-x$，非退化矩形要求 $0<x<6$。面积函数为
\[
A(x)=x(6-x),\qquad A^{\prime}(x)=6-2x.
\]
唯一驻点是 $x=3$，在其左侧导数为正、右侧为负，故面积在这里取得全局最大值 $9$。也可配方得到
\[
A(x)=9-(x-3)^2\le 9,
\]
等号恰在正方形边长为 $3$ 时成立。可行域虽为开区间，但最大值在内部达到；本题不能仅凭闭区间最值定理来代替实际论证。
\end{solve}

\begin{solve}[习题~\ref{ex:c11q11}：固定周长矩形的最大面积]\phantomsection\label{ans:c11q11}
另一边为 $8-x$，非退化矩形要求 $0<x<8$。面积函数为
\[
A(x)=x(8-x),\qquad A^{\prime}(x)=8-2x.
\]
唯一驻点是 $x=4$，在其左侧导数为正、右侧为负，故面积在这里取得全局最大值 $16$。也可配方得到
\[
A(x)=16-(x-4)^2\le 16,
\]
等号恰在正方形边长为 $4$ 时成立。可行域虽为开区间，但最大值在内部达到；本题不能仅凭闭区间最值定理来代替实际论证。
\end{solve}

\begin{solve}[习题~\ref{ex:c11q12}：固定周长矩形的最大面积]\phantomsection\label{ans:c11q12}
另一边为 $10-x$，非退化矩形要求 $0<x<10$。面积函数为
\[
A(x)=x(10-x),\qquad A^{\prime}(x)=10-2x.
\]
唯一驻点是 $x=5$，在其左侧导数为正、右侧为负，故面积在这里取得全局最大值 $25$。也可配方得到
\[
A(x)=25-(x-5)^2\le 25,
\]
等号恰在正方形边长为 $5$ 时成立。可行域虽为开区间，但最大值在内部达到；本题不能仅凭闭区间最值定理来代替实际论证。
\end{solve}

\begin{solve}[习题~\ref{ex:c11q13}：固定周长矩形的最大面积]\phantomsection\label{ans:c11q13}
另一边为 $12-x$，非退化矩形要求 $0<x<12$。面积函数为
\[
A(x)=x(12-x),\qquad A^{\prime}(x)=12-2x.
\]
唯一驻点是 $x=6$，在其左侧导数为正、右侧为负，故面积在这里取得全局最大值 $36$。也可配方得到
\[
A(x)=36-(x-6)^2\le 36,
\]
等号恰在正方形边长为 $6$ 时成立。可行域虽为开区间，但最大值在内部达到；本题不能仅凭闭区间最值定理来代替实际论证。
\end{solve}

\begin{solve}[习题~\ref{ex:c11q14}：指数衰减与内点最大值]\phantomsection\label{ans:c11q14}
求导得
\[
f^{\prime}(x)=e^{-1x}(1-1x).
\]
指数因子为正，故导数在 $x<1/1$ 时为正，在 $x>1/1$ 时为负。函数先增后减，最大值为
\[
f(1/1)=\frac1{1e}.
\]
由于函数在该半轴上非负且 $f(0)=0$，最小值为 $0$，在左端点取得。虽然 $f(x)\to0$ 当 $x\to\infty$，无穷远不是定义域中的一个点；它只提供渐近信息，不能作为一个取得最值的位置。
\end{solve}

\begin{solve}[习题~\ref{ex:c11q15}：指数衰减与内点最大值]\phantomsection\label{ans:c11q15}
求导得
\[
f^{\prime}(x)=e^{-2x}(1-2x).
\]
指数因子为正，故导数在 $x<1/2$ 时为正，在 $x>1/2$ 时为负。函数先增后减，最大值为
\[
f(1/2)=\frac1{2e}.
\]
由于函数在该半轴上非负且 $f(0)=0$，最小值为 $0$，在左端点取得。虽然 $f(x)\to0$ 当 $x\to\infty$，无穷远不是定义域中的一个点；它只提供渐近信息，不能作为一个取得最值的位置。
\end{solve}

\begin{solve}[习题~\ref{ex:c11q16}：指数衰减与内点最大值]\phantomsection\label{ans:c11q16}
求导得
\[
f^{\prime}(x)=e^{-3x}(1-3x).
\]
指数因子为正，故导数在 $x<1/3$ 时为正，在 $x>1/3$ 时为负。函数先增后减，最大值为
\[
f(1/3)=\frac1{3e}.
\]
由于函数在该半轴上非负且 $f(0)=0$，最小值为 $0$，在左端点取得。虽然 $f(x)\to0$ 当 $x\to\infty$，无穷远不是定义域中的一个点；它只提供渐近信息，不能作为一个取得最值的位置。
\end{solve}

\begin{solve}[习题~\ref{ex:c11q17}：指数衰减与内点最大值]\phantomsection\label{ans:c11q17}
求导得
\[
f^{\prime}(x)=e^{-4x}(1-4x).
\]
指数因子为正，故导数在 $x<1/4$ 时为正，在 $x>1/4$ 时为负。函数先增后减，最大值为
\[
f(1/4)=\frac1{4e}.
\]
由于函数在该半轴上非负且 $f(0)=0$，最小值为 $0$，在左端点取得。虽然 $f(x)\to0$ 当 $x\to\infty$，无穷远不是定义域中的一个点；它只提供渐近信息，不能作为一个取得最值的位置。
\end{solve}

\begin{solve}[习题~\ref{ex:c11q18}：指数衰减与内点最大值]\phantomsection\label{ans:c11q18}
求导得
\[
f^{\prime}(x)=e^{-5x}(1-5x).
\]
指数因子为正，故导数在 $x<1/5$ 时为正，在 $x>1/5$ 时为负。函数先增后减，最大值为
\[
f(1/5)=\frac1{5e}.
\]
由于函数在该半轴上非负且 $f(0)=0$，最小值为 $0$，在左端点取得。虽然 $f(x)\to0$ 当 $x\to\infty$，无穷远不是定义域中的一个点；它只提供渐近信息，不能作为一个取得最值的位置。
\end{solve}

\begin{solve}[习题~\ref{ex:c11q19}：指数衰减与内点最大值]\phantomsection\label{ans:c11q19}
求导得
\[
f^{\prime}(x)=e^{-6x}(1-6x).
\]
指数因子为正，故导数在 $x<1/6$ 时为正，在 $x>1/6$ 时为负。函数先增后减，最大值为
\[
f(1/6)=\frac1{6e}.
\]
由于函数在该半轴上非负且 $f(0)=0$，最小值为 $0$，在左端点取得。虽然 $f(x)\to0$ 当 $x\to\infty$，无穷远不是定义域中的一个点；它只提供渐近信息，不能作为一个取得最值的位置。
\end{solve}

\begin{solve}[习题~\ref{ex:c11q20}：相关变化率中的链式法则]\phantomsection\label{ans:c11q20}
体积 $V=4\pi r^3/3$、表面积 $S=4\pi r^2$ 都通过 $r=r(t)$ 依赖时间。因此
\[
\frac{dV}{dt}=4\pi r^2\frac{dr}{dt},\qquad\frac{dS}{dt}=8\pi r\frac{dr}{dt}.
\]
将 $r=2$、$dr/dt=1$ 代入，得到
\[
\frac{dV}{dt}=16\pi\ \mathrm{cm}^3/\mathrm{s},\qquad\frac{dS}{dt}=16\pi\ \mathrm{cm}^2/\mathrm{s}.
\]
体积变化率与面积变化率的量纲不同；不能把对半径的导数直接当成对时间的导数。先建立变量之间的关系，再代入某一时刻的数据，可以避免把半径误当成固定常数。
\end{solve}

\begin{solve}[习题~\ref{ex:c11q21}：相关变化率中的链式法则]\phantomsection\label{ans:c11q21}
体积 $V=4\pi r^3/3$、表面积 $S=4\pi r^2$ 都通过 $r=r(t)$ 依赖时间。因此
\[
\frac{dV}{dt}=4\pi r^2\frac{dr}{dt},\qquad\frac{dS}{dt}=8\pi r\frac{dr}{dt}.
\]
将 $r=3$、$dr/dt=2$ 代入，得到
\[
\frac{dV}{dt}=72\pi\ \mathrm{cm}^3/\mathrm{s},\qquad\frac{dS}{dt}=48\pi\ \mathrm{cm}^2/\mathrm{s}.
\]
体积变化率与面积变化率的量纲不同；不能把对半径的导数直接当成对时间的导数。先建立变量之间的关系，再代入某一时刻的数据，可以避免把半径误当成固定常数。
\end{solve}

\begin{solve}[习题~\ref{ex:c11q22}：相关变化率中的链式法则]\phantomsection\label{ans:c11q22}
体积 $V=4\pi r^3/3$、表面积 $S=4\pi r^2$ 都通过 $r=r(t)$ 依赖时间。因此
\[
\frac{dV}{dt}=4\pi r^2\frac{dr}{dt},\qquad\frac{dS}{dt}=8\pi r\frac{dr}{dt}.
\]
将 $r=4$、$dr/dt=3$ 代入，得到
\[
\frac{dV}{dt}=192\pi\ \mathrm{cm}^3/\mathrm{s},\qquad\frac{dS}{dt}=96\pi\ \mathrm{cm}^2/\mathrm{s}.
\]
体积变化率与面积变化率的量纲不同；不能把对半径的导数直接当成对时间的导数。先建立变量之间的关系，再代入某一时刻的数据，可以避免把半径误当成固定常数。
\end{solve}

\begin{solve}[习题~\ref{ex:c11q23}：相关变化率中的链式法则]\phantomsection\label{ans:c11q23}
体积 $V=4\pi r^3/3$、表面积 $S=4\pi r^2$ 都通过 $r=r(t)$ 依赖时间。因此
\[
\frac{dV}{dt}=4\pi r^2\frac{dr}{dt},\qquad\frac{dS}{dt}=8\pi r\frac{dr}{dt}.
\]
将 $r=5$、$dr/dt=4$ 代入，得到
\[
\frac{dV}{dt}=400\pi\ \mathrm{cm}^3/\mathrm{s},\qquad\frac{dS}{dt}=160\pi\ \mathrm{cm}^2/\mathrm{s}.
\]
体积变化率与面积变化率的量纲不同；不能把对半径的导数直接当成对时间的导数。先建立变量之间的关系，再代入某一时刻的数据，可以避免把半径误当成固定常数。
\end{solve}

\begin{solve}[习题~\ref{ex:c11q24}：相关变化率中的链式法则]\phantomsection\label{ans:c11q24}
体积 $V=4\pi r^3/3$、表面积 $S=4\pi r^2$ 都通过 $r=r(t)$ 依赖时间。因此
\[
\frac{dV}{dt}=4\pi r^2\frac{dr}{dt},\qquad\frac{dS}{dt}=8\pi r\frac{dr}{dt}.
\]
将 $r=6$、$dr/dt=5$ 代入，得到
\[
\frac{dV}{dt}=720\pi\ \mathrm{cm}^3/\mathrm{s},\qquad\frac{dS}{dt}=240\pi\ \mathrm{cm}^2/\mathrm{s}.
\]
体积变化率与面积变化率的量纲不同；不能把对半径的导数直接当成对时间的导数。先建立变量之间的关系，再代入某一时刻的数据，可以避免把半径误当成固定常数。
\end{solve}

\begin{solve}[习题~\ref{ex:c11q25}：相关变化率中的链式法则]\phantomsection\label{ans:c11q25}
体积 $V=4\pi r^3/3$、表面积 $S=4\pi r^2$ 都通过 $r=r(t)$ 依赖时间。因此
\[
\frac{dV}{dt}=4\pi r^2\frac{dr}{dt},\qquad\frac{dS}{dt}=8\pi r\frac{dr}{dt}.
\]
将 $r=7$、$dr/dt=6$ 代入，得到
\[
\frac{dV}{dt}=1176\pi\ \mathrm{cm}^3/\mathrm{s},\qquad\frac{dS}{dt}=336\pi\ \mathrm{cm}^2/\mathrm{s}.
\]
体积变化率与面积变化率的量纲不同；不能把对半径的导数直接当成对时间的导数。先建立变量之间的关系，再代入某一时刻的数据，可以避免把半径误当成固定常数。
\end{solve}

\clearpage
\section{不定积分与第一换元法}\label{sec:ch12}
原函数（antiderivative）在每个连通区间上只相差一个常数。换元法是链式法则的逆过程：被积函数中应当同时出现一个复合结构及其内层导数。所有不定积分答案均理解为在被积函数有定义的区间内成立。

\textbf{资源对应：}CLP-2 \S~1.4、\S~1.13。本章共 24 道编号习题。
\subsection{习题}
\begin{example}[识别幂函数的内层导数]\label{ex:c12q1}
计算 $\displaystyle\int x(1x^2+1)^{2}\,dx$，并通过求导验证结果。
\par\textit{改编训练。}\par\hyperref[ans:c12q1]{解答见第~\pageref{ans:c12q1} 页。}
\end{example}
\begin{example}[识别幂函数的内层导数]\label{ex:c12q2}
计算 $\displaystyle\int x(2x^2+1)^{3}\,dx$，并通过求导验证结果。
\par\textit{改编训练。}\par\hyperref[ans:c12q2]{解答见第~\pageref{ans:c12q2} 页。}
\end{example}
\begin{example}[识别幂函数的内层导数]\label{ex:c12q3}
计算 $\displaystyle\int x(3x^2+1)^{2}\,dx$，并通过求导验证结果。
\par\textit{改编训练。}\par\hyperref[ans:c12q3]{解答见第~\pageref{ans:c12q3} 页。}
\end{example}
\begin{example}[识别幂函数的内层导数]\label{ex:c12q4}
计算 $\displaystyle\int x(2x^2+1)^{4}\,dx$，并通过求导验证结果。
\par\textit{改编训练。}\par\hyperref[ans:c12q4]{解答见第~\pageref{ans:c12q4} 页。}
\end{example}
\begin{example}[识别幂函数的内层导数]\label{ex:c12q5}
计算 $\displaystyle\int x(4x^2+1)^{3}\,dx$，并通过求导验证结果。
\par\textit{改编训练。}\par\hyperref[ans:c12q5]{解答见第~\pageref{ans:c12q5} 页。}
\end{example}
\begin{example}[识别幂函数的内层导数]\label{ex:c12q6}
计算 $\displaystyle\int x(3x^2+1)^{5}\,dx$，并通过求导验证结果。
\par\textit{改编训练。}\par\hyperref[ans:c12q6]{解答见第~\pageref{ans:c12q6} 页。}
\end{example}
\begin{example}[分母导数型积分]\label{ex:c12q7}
计算 $\displaystyle\int\frac{x}{1x^2+1}\,dx$，并说明是否需要绝对值符号。
\par\textit{改编训练。}\par\hyperref[ans:c12q7]{解答见第~\pageref{ans:c12q7} 页。}
\end{example}
\begin{example}[分母导数型积分]\label{ex:c12q8}
计算 $\displaystyle\int\frac{x}{2x^2+1}\,dx$，并说明是否需要绝对值符号。
\par\textit{改编训练。}\par\hyperref[ans:c12q8]{解答见第~\pageref{ans:c12q8} 页。}
\end{example}
\begin{example}[分母导数型积分]\label{ex:c12q9}
计算 $\displaystyle\int\frac{x}{3x^2+2}\,dx$，并说明是否需要绝对值符号。
\par\textit{改编训练。}\par\hyperref[ans:c12q9]{解答见第~\pageref{ans:c12q9} 页。}
\end{example}
\begin{example}[分母导数型积分]\label{ex:c12q10}
计算 $\displaystyle\int\frac{x}{4x^2+3}\,dx$，并说明是否需要绝对值符号。
\par\textit{改编训练。}\par\hyperref[ans:c12q10]{解答见第~\pageref{ans:c12q10} 页。}
\end{example}
\begin{example}[分母导数型积分]\label{ex:c12q11}
计算 $\displaystyle\int\frac{x}{2x^2+5}\,dx$，并说明是否需要绝对值符号。
\par\textit{改编训练。}\par\hyperref[ans:c12q11]{解答见第~\pageref{ans:c12q11} 页。}
\end{example}
\begin{example}[分母导数型积分]\label{ex:c12q12}
计算 $\displaystyle\int\frac{x}{5x^2+2}\,dx$，并说明是否需要绝对值符号。
\par\textit{改编训练。}\par\hyperref[ans:c12q12]{解答见第~\pageref{ans:c12q12} 页。}
\end{example}
\begin{example}[指数复合的反向链式法则]\label{ex:c12q13}
计算 $\displaystyle\int x e^{1x^2}\,dx$ 和 $\displaystyle\int e^{1x}\cos(e^{1x})\,dx$。
\par\textit{改编训练。}\par\hyperref[ans:c12q13]{解答见第~\pageref{ans:c12q13} 页。}
\end{example}
\begin{example}[指数复合的反向链式法则]\label{ex:c12q14}
计算 $\displaystyle\int x e^{2x^2}\,dx$ 和 $\displaystyle\int e^{2x}\cos(e^{2x})\,dx$。
\par\textit{改编训练。}\par\hyperref[ans:c12q14]{解答见第~\pageref{ans:c12q14} 页。}
\end{example}
\begin{example}[指数复合的反向链式法则]\label{ex:c12q15}
计算 $\displaystyle\int x e^{3x^2}\,dx$ 和 $\displaystyle\int e^{3x}\cos(e^{3x})\,dx$。
\par\textit{改编训练。}\par\hyperref[ans:c12q15]{解答见第~\pageref{ans:c12q15} 页。}
\end{example}
\begin{example}[指数复合的反向链式法则]\label{ex:c12q16}
计算 $\displaystyle\int x e^{4x^2}\,dx$ 和 $\displaystyle\int e^{4x}\cos(e^{4x})\,dx$。
\par\textit{改编训练。}\par\hyperref[ans:c12q16]{解答见第~\pageref{ans:c12q16} 页。}
\end{example}
\begin{example}[指数复合的反向链式法则]\label{ex:c12q17}
计算 $\displaystyle\int x e^{5x^2}\,dx$ 和 $\displaystyle\int e^{5x}\cos(e^{5x})\,dx$。
\par\textit{改编训练。}\par\hyperref[ans:c12q17]{解答见第~\pageref{ans:c12q17} 页。}
\end{example}
\begin{example}[指数复合的反向链式法则]\label{ex:c12q18}
计算 $\displaystyle\int x e^{6x^2}\,dx$ 和 $\displaystyle\int e^{6x}\cos(e^{6x})\,dx$。
\par\textit{改编训练。}\par\hyperref[ans:c12q18]{解答见第~\pageref{ans:c12q18} 页。}
\end{example}
\begin{example}[对数幂与变量范围]\label{ex:c12q19}
在 $x>0$ 上计算 $\displaystyle\int\frac{(\log x)^{2}}{1x}\,dx$，并求从 $1$ 到 $e$ 的定积分。
\par\textit{改编训练。}\par\hyperref[ans:c12q19]{解答见第~\pageref{ans:c12q19} 页。}
\end{example}
\begin{example}[对数幂与变量范围]\label{ex:c12q20}
在 $x>0$ 上计算 $\displaystyle\int\frac{(\log x)^{3}}{2x}\,dx$，并求从 $1$ 到 $e$ 的定积分。
\par\textit{改编训练。}\par\hyperref[ans:c12q20]{解答见第~\pageref{ans:c12q20} 页。}
\end{example}
\begin{example}[对数幂与变量范围]\label{ex:c12q21}
在 $x>0$ 上计算 $\displaystyle\int\frac{(\log x)^{2}}{3x}\,dx$，并求从 $1$ 到 $e$ 的定积分。
\par\textit{改编训练。}\par\hyperref[ans:c12q21]{解答见第~\pageref{ans:c12q21} 页。}
\end{example}
\begin{example}[对数幂与变量范围]\label{ex:c12q22}
在 $x>0$ 上计算 $\displaystyle\int\frac{(\log x)^{4}}{2x}\,dx$，并求从 $1$ 到 $e$ 的定积分。
\par\textit{改编训练。}\par\hyperref[ans:c12q22]{解答见第~\pageref{ans:c12q22} 页。}
\end{example}
\begin{example}[对数幂与变量范围]\label{ex:c12q23}
在 $x>0$ 上计算 $\displaystyle\int\frac{(\log x)^{3}}{4x}\,dx$，并求从 $1$ 到 $e$ 的定积分。
\par\textit{改编训练。}\par\hyperref[ans:c12q23]{解答见第~\pageref{ans:c12q23} 页。}
\end{example}
\begin{example}[对数幂与变量范围]\label{ex:c12q24}
在 $x>0$ 上计算 $\displaystyle\int\frac{(\log x)^{5}}{3x}\,dx$，并求从 $1$ 到 $e$ 的定积分。
\par\textit{改编训练。}\par\hyperref[ans:c12q24]{解答见第~\pageref{ans:c12q24} 页。}
\end{example}
\subsection{解答}
\begin{solve}[习题~\ref{ex:c12q1}：识别幂函数的内层导数]\phantomsection\label{ans:c12q1}
令 $u=1x^2+1$，则 $du=2x\,dx$。原积分变为
\[
\frac1{2}\int u^{2}\,du=\frac{1}{6}u^{3}+C.
\]
代回得 $F(x)=\frac{1}{6}(1x^2+1)^{3}+C$。求导时，外层带来 $3$，内层带来 $2x$，恰好抵消前面的系数，恢复原被积函数。换元过程中少除一个 $2$ 是常见的系数错误；不定积分中的 $C$ 可以吸收常数，但不能吸收乘在非恒定函数前的错误系数。
\end{solve}

\begin{solve}[习题~\ref{ex:c12q2}：识别幂函数的内层导数]\phantomsection\label{ans:c12q2}
令 $u=2x^2+1$，则 $du=4x\,dx$。原积分变为
\[
\frac1{4}\int u^{3}\,du=\frac{1}{16}u^{4}+C.
\]
代回得 $F(x)=\frac{1}{16}(2x^2+1)^{4}+C$。求导时，外层带来 $4$，内层带来 $4x$，恰好抵消前面的系数，恢复原被积函数。换元过程中少除一个 $4$ 是常见的系数错误；不定积分中的 $C$ 可以吸收常数，但不能吸收乘在非恒定函数前的错误系数。
\end{solve}

\begin{solve}[习题~\ref{ex:c12q3}：识别幂函数的内层导数]\phantomsection\label{ans:c12q3}
令 $u=3x^2+1$，则 $du=6x\,dx$。原积分变为
\[
\frac1{6}\int u^{2}\,du=\frac{1}{18}u^{3}+C.
\]
代回得 $F(x)=\frac{1}{18}(3x^2+1)^{3}+C$。求导时，外层带来 $3$，内层带来 $6x$，恰好抵消前面的系数，恢复原被积函数。换元过程中少除一个 $6$ 是常见的系数错误；不定积分中的 $C$ 可以吸收常数，但不能吸收乘在非恒定函数前的错误系数。
\end{solve}

\begin{solve}[习题~\ref{ex:c12q4}：识别幂函数的内层导数]\phantomsection\label{ans:c12q4}
令 $u=2x^2+1$，则 $du=4x\,dx$。原积分变为
\[
\frac1{4}\int u^{4}\,du=\frac{1}{20}u^{5}+C.
\]
代回得 $F(x)=\frac{1}{20}(2x^2+1)^{5}+C$。求导时，外层带来 $5$，内层带来 $4x$，恰好抵消前面的系数，恢复原被积函数。换元过程中少除一个 $4$ 是常见的系数错误；不定积分中的 $C$ 可以吸收常数，但不能吸收乘在非恒定函数前的错误系数。
\end{solve}

\begin{solve}[习题~\ref{ex:c12q5}：识别幂函数的内层导数]\phantomsection\label{ans:c12q5}
令 $u=4x^2+1$，则 $du=8x\,dx$。原积分变为
\[
\frac1{8}\int u^{3}\,du=\frac{1}{32}u^{4}+C.
\]
代回得 $F(x)=\frac{1}{32}(4x^2+1)^{4}+C$。求导时，外层带来 $4$，内层带来 $8x$，恰好抵消前面的系数，恢复原被积函数。换元过程中少除一个 $8$ 是常见的系数错误；不定积分中的 $C$ 可以吸收常数，但不能吸收乘在非恒定函数前的错误系数。
\end{solve}

\begin{solve}[习题~\ref{ex:c12q6}：识别幂函数的内层导数]\phantomsection\label{ans:c12q6}
令 $u=3x^2+1$，则 $du=6x\,dx$。原积分变为
\[
\frac1{6}\int u^{5}\,du=\frac{1}{36}u^{6}+C.
\]
代回得 $F(x)=\frac{1}{36}(3x^2+1)^{6}+C$。求导时，外层带来 $6$，内层带来 $6x$，恰好抵消前面的系数，恢复原被积函数。换元过程中少除一个 $6$ 是常见的系数错误；不定积分中的 $C$ 可以吸收常数，但不能吸收乘在非恒定函数前的错误系数。
\end{solve}

\begin{solve}[习题~\ref{ex:c12q7}：分母导数型积分]\phantomsection\label{ans:c12q7}
取 $u=1x^2+1$，则 $du=2x\,dx$。因此
\[
\int\frac{x}{1x^2+1}\,dx=\frac{1}{2}\int\frac{du}u=\frac{1}{2}\log(1x^2+1)+C.
\]
因为 $1,1>0$，分母在整个实轴上为正，故对数内可以省略绝对值。验证：
\[
\frac d{dx}\left[\frac{1}{2}\log(1x^2+1)\right]=\frac{x}{1x^2+1}.
\]
一般的 $\int u^{-1}du$ 应写成 $\log|u|+C$；本题省略绝对值是由已核对的正性保证的。
\end{solve}

\begin{solve}[习题~\ref{ex:c12q8}：分母导数型积分]\phantomsection\label{ans:c12q8}
取 $u=2x^2+1$，则 $du=4x\,dx$。因此
\[
\int\frac{x}{2x^2+1}\,dx=\frac{1}{4}\int\frac{du}u=\frac{1}{4}\log(2x^2+1)+C.
\]
因为 $2,1>0$，分母在整个实轴上为正，故对数内可以省略绝对值。验证：
\[
\frac d{dx}\left[\frac{1}{4}\log(2x^2+1)\right]=\frac{x}{2x^2+1}.
\]
一般的 $\int u^{-1}du$ 应写成 $\log|u|+C$；本题省略绝对值是由已核对的正性保证的。
\end{solve}

\begin{solve}[习题~\ref{ex:c12q9}：分母导数型积分]\phantomsection\label{ans:c12q9}
取 $u=3x^2+2$，则 $du=6x\,dx$。因此
\[
\int\frac{x}{3x^2+2}\,dx=\frac{1}{6}\int\frac{du}u=\frac{1}{6}\log(3x^2+2)+C.
\]
因为 $3,2>0$，分母在整个实轴上为正，故对数内可以省略绝对值。验证：
\[
\frac d{dx}\left[\frac{1}{6}\log(3x^2+2)\right]=\frac{x}{3x^2+2}.
\]
一般的 $\int u^{-1}du$ 应写成 $\log|u|+C$；本题省略绝对值是由已核对的正性保证的。
\end{solve}

\begin{solve}[习题~\ref{ex:c12q10}：分母导数型积分]\phantomsection\label{ans:c12q10}
取 $u=4x^2+3$，则 $du=8x\,dx$。因此
\[
\int\frac{x}{4x^2+3}\,dx=\frac{1}{8}\int\frac{du}u=\frac{1}{8}\log(4x^2+3)+C.
\]
因为 $4,3>0$，分母在整个实轴上为正，故对数内可以省略绝对值。验证：
\[
\frac d{dx}\left[\frac{1}{8}\log(4x^2+3)\right]=\frac{x}{4x^2+3}.
\]
一般的 $\int u^{-1}du$ 应写成 $\log|u|+C$；本题省略绝对值是由已核对的正性保证的。
\end{solve}

\begin{solve}[习题~\ref{ex:c12q11}：分母导数型积分]\phantomsection\label{ans:c12q11}
取 $u=2x^2+5$，则 $du=4x\,dx$。因此
\[
\int\frac{x}{2x^2+5}\,dx=\frac{1}{4}\int\frac{du}u=\frac{1}{4}\log(2x^2+5)+C.
\]
因为 $2,5>0$，分母在整个实轴上为正，故对数内可以省略绝对值。验证：
\[
\frac d{dx}\left[\frac{1}{4}\log(2x^2+5)\right]=\frac{x}{2x^2+5}.
\]
一般的 $\int u^{-1}du$ 应写成 $\log|u|+C$；本题省略绝对值是由已核对的正性保证的。
\end{solve}

\begin{solve}[习题~\ref{ex:c12q12}：分母导数型积分]\phantomsection\label{ans:c12q12}
取 $u=5x^2+2$，则 $du=10x\,dx$。因此
\[
\int\frac{x}{5x^2+2}\,dx=\frac{1}{10}\int\frac{du}u=\frac{1}{10}\log(5x^2+2)+C.
\]
因为 $5,2>0$，分母在整个实轴上为正，故对数内可以省略绝对值。验证：
\[
\frac d{dx}\left[\frac{1}{10}\log(5x^2+2)\right]=\frac{x}{5x^2+2}.
\]
一般的 $\int u^{-1}du$ 应写成 $\log|u|+C$；本题省略绝对值是由已核对的正性保证的。
\end{solve}

\begin{solve}[习题~\ref{ex:c12q13}：指数复合的反向链式法则]\phantomsection\label{ans:c12q13}
第一式令 $u=1x^2$，第二式令 $v=e^{1x}$。两个替换的微分分别为 $du=2x\,dx$ 与 $dv=1e^{1x}\,dx$，从而
\[
\int xe^{1x^2}\,dx=\frac{1}{2}e^{1x^2}+C_1,
\]

\[
\int e^{1x}\cos(e^{1x})\,dx=1\sin(e^{1x})+C_2.
\]
对右侧分别求导即可恢复左侧的被积函数。两个不定积分彼此独立，所以使用两个任意常数；它们没有理由必须相等。
\end{solve}

\begin{solve}[习题~\ref{ex:c12q14}：指数复合的反向链式法则]\phantomsection\label{ans:c12q14}
第一式令 $u=2x^2$，第二式令 $v=e^{2x}$。两个替换的微分分别为 $du=4x\,dx$ 与 $dv=2e^{2x}\,dx$，从而
\[
\int xe^{2x^2}\,dx=\frac{1}{4}e^{2x^2}+C_1,
\]

\[
\int e^{2x}\cos(e^{2x})\,dx=\frac{1}{2}\sin(e^{2x})+C_2.
\]
对右侧分别求导即可恢复左侧的被积函数。两个不定积分彼此独立，所以使用两个任意常数；它们没有理由必须相等。
\end{solve}

\begin{solve}[习题~\ref{ex:c12q15}：指数复合的反向链式法则]\phantomsection\label{ans:c12q15}
第一式令 $u=3x^2$，第二式令 $v=e^{3x}$。两个替换的微分分别为 $du=6x\,dx$ 与 $dv=3e^{3x}\,dx$，从而
\[
\int xe^{3x^2}\,dx=\frac{1}{6}e^{3x^2}+C_1,
\]

\[
\int e^{3x}\cos(e^{3x})\,dx=\frac{1}{3}\sin(e^{3x})+C_2.
\]
对右侧分别求导即可恢复左侧的被积函数。两个不定积分彼此独立，所以使用两个任意常数；它们没有理由必须相等。
\end{solve}

\begin{solve}[习题~\ref{ex:c12q16}：指数复合的反向链式法则]\phantomsection\label{ans:c12q16}
第一式令 $u=4x^2$，第二式令 $v=e^{4x}$。两个替换的微分分别为 $du=8x\,dx$ 与 $dv=4e^{4x}\,dx$，从而
\[
\int xe^{4x^2}\,dx=\frac{1}{8}e^{4x^2}+C_1,
\]

\[
\int e^{4x}\cos(e^{4x})\,dx=\frac{1}{4}\sin(e^{4x})+C_2.
\]
对右侧分别求导即可恢复左侧的被积函数。两个不定积分彼此独立，所以使用两个任意常数；它们没有理由必须相等。
\end{solve}

\begin{solve}[习题~\ref{ex:c12q17}：指数复合的反向链式法则]\phantomsection\label{ans:c12q17}
第一式令 $u=5x^2$，第二式令 $v=e^{5x}$。两个替换的微分分别为 $du=10x\,dx$ 与 $dv=5e^{5x}\,dx$，从而
\[
\int xe^{5x^2}\,dx=\frac{1}{10}e^{5x^2}+C_1,
\]

\[
\int e^{5x}\cos(e^{5x})\,dx=\frac{1}{5}\sin(e^{5x})+C_2.
\]
对右侧分别求导即可恢复左侧的被积函数。两个不定积分彼此独立，所以使用两个任意常数；它们没有理由必须相等。
\end{solve}

\begin{solve}[习题~\ref{ex:c12q18}：指数复合的反向链式法则]\phantomsection\label{ans:c12q18}
第一式令 $u=6x^2$，第二式令 $v=e^{6x}$。两个替换的微分分别为 $du=12x\,dx$ 与 $dv=6e^{6x}\,dx$，从而
\[
\int xe^{6x^2}\,dx=\frac{1}{12}e^{6x^2}+C_1,
\]

\[
\int e^{6x}\cos(e^{6x})\,dx=\frac{1}{6}\sin(e^{6x})+C_2.
\]
对右侧分别求导即可恢复左侧的被积函数。两个不定积分彼此独立，所以使用两个任意常数；它们没有理由必须相等。
\end{solve}

\begin{solve}[习题~\ref{ex:c12q19}：对数幂与变量范围]\phantomsection\label{ans:c12q19}
令 $u=\log x$，则 $du=dx/x$，于是
\[
\int\frac{(\log x)^{2}}{1x}\,dx=\frac{1}{3}(\log x)^{3}+C.
\]
计算定积分时，$x=1$ 对应 $u=0$，$x=e$ 对应 $u=1$，所以
\[
\int_1^e\frac{(\log x)^{2}}{1x}\,dx=\frac1{1}\int_0^1u^{2}\,du=\frac{1}{3}.
\]
定积分换元后必须同时改变上下限。若选择保留原上下限，则必须先把原函数代回 $x$，不能把新变量与旧上下限混用。
\end{solve}

\begin{solve}[习题~\ref{ex:c12q20}：对数幂与变量范围]\phantomsection\label{ans:c12q20}
令 $u=\log x$，则 $du=dx/x$，于是
\[
\int\frac{(\log x)^{3}}{2x}\,dx=\frac{1}{8}(\log x)^{4}+C.
\]
计算定积分时，$x=1$ 对应 $u=0$，$x=e$ 对应 $u=1$，所以
\[
\int_1^e\frac{(\log x)^{3}}{2x}\,dx=\frac1{2}\int_0^1u^{3}\,du=\frac{1}{8}.
\]
定积分换元后必须同时改变上下限。若选择保留原上下限，则必须先把原函数代回 $x$，不能把新变量与旧上下限混用。
\end{solve}

\begin{solve}[习题~\ref{ex:c12q21}：对数幂与变量范围]\phantomsection\label{ans:c12q21}
令 $u=\log x$，则 $du=dx/x$，于是
\[
\int\frac{(\log x)^{2}}{3x}\,dx=\frac{1}{9}(\log x)^{3}+C.
\]
计算定积分时，$x=1$ 对应 $u=0$，$x=e$ 对应 $u=1$，所以
\[
\int_1^e\frac{(\log x)^{2}}{3x}\,dx=\frac1{3}\int_0^1u^{2}\,du=\frac{1}{9}.
\]
定积分换元后必须同时改变上下限。若选择保留原上下限，则必须先把原函数代回 $x$，不能把新变量与旧上下限混用。
\end{solve}

\begin{solve}[习题~\ref{ex:c12q22}：对数幂与变量范围]\phantomsection\label{ans:c12q22}
令 $u=\log x$，则 $du=dx/x$，于是
\[
\int\frac{(\log x)^{4}}{2x}\,dx=\frac{1}{10}(\log x)^{5}+C.
\]
计算定积分时，$x=1$ 对应 $u=0$，$x=e$ 对应 $u=1$，所以
\[
\int_1^e\frac{(\log x)^{4}}{2x}\,dx=\frac1{2}\int_0^1u^{4}\,du=\frac{1}{10}.
\]
定积分换元后必须同时改变上下限。若选择保留原上下限，则必须先把原函数代回 $x$，不能把新变量与旧上下限混用。
\end{solve}

\begin{solve}[习题~\ref{ex:c12q23}：对数幂与变量范围]\phantomsection\label{ans:c12q23}
令 $u=\log x$，则 $du=dx/x$，于是
\[
\int\frac{(\log x)^{3}}{4x}\,dx=\frac{1}{16}(\log x)^{4}+C.
\]
计算定积分时，$x=1$ 对应 $u=0$，$x=e$ 对应 $u=1$，所以
\[
\int_1^e\frac{(\log x)^{3}}{4x}\,dx=\frac1{4}\int_0^1u^{3}\,du=\frac{1}{16}.
\]
定积分换元后必须同时改变上下限。若选择保留原上下限，则必须先把原函数代回 $x$，不能把新变量与旧上下限混用。
\end{solve}

\begin{solve}[习题~\ref{ex:c12q24}：对数幂与变量范围]\phantomsection\label{ans:c12q24}
令 $u=\log x$，则 $du=dx/x$，于是
\[
\int\frac{(\log x)^{5}}{3x}\,dx=\frac{1}{18}(\log x)^{6}+C.
\]
计算定积分时，$x=1$ 对应 $u=0$，$x=e$ 对应 $u=1$，所以
\[
\int_1^e\frac{(\log x)^{5}}{3x}\,dx=\frac1{3}\int_0^1u^{5}\,du=\frac{1}{18}.
\]
定积分换元后必须同时改变上下限。若选择保留原上下限，则必须先把原函数代回 $x$，不能把新变量与旧上下限混用。
\end{solve}

\clearpage
\section{分部积分与有理函数积分}\label{sec:ch13}
分部积分（integration by parts）来自乘积求导公式；合适的选择应当让剩余积分更简单。有理函数先比较分子、分母次数，再因式分解和分解部分分式。

\textbf{资源对应：}CLP-2 \S~1.7、\S~1.10。本章共 24 道编号习题。
\subsection{习题}
\begin{example}[一次多项式乘指数]\label{ex:c13q1}
计算 $\displaystyle\int xe^{1x}\,dx$，以及 $\displaystyle\int_0^1xe^{1x}\,dx$。
\par\textit{改编训练。}\par\hyperref[ans:c13q1]{解答见第~\pageref{ans:c13q1} 页。}
\end{example}
\begin{example}[一次多项式乘指数]\label{ex:c13q2}
计算 $\displaystyle\int xe^{2x}\,dx$，以及 $\displaystyle\int_0^1xe^{2x}\,dx$。
\par\textit{改编训练。}\par\hyperref[ans:c13q2]{解答见第~\pageref{ans:c13q2} 页。}
\end{example}
\begin{example}[一次多项式乘指数]\label{ex:c13q3}
计算 $\displaystyle\int xe^{3x}\,dx$，以及 $\displaystyle\int_0^1xe^{3x}\,dx$。
\par\textit{改编训练。}\par\hyperref[ans:c13q3]{解答见第~\pageref{ans:c13q3} 页。}
\end{example}
\begin{example}[一次多项式乘指数]\label{ex:c13q4}
计算 $\displaystyle\int xe^{4x}\,dx$，以及 $\displaystyle\int_0^1xe^{4x}\,dx$。
\par\textit{改编训练。}\par\hyperref[ans:c13q4]{解答见第~\pageref{ans:c13q4} 页。}
\end{example}
\begin{example}[一次多项式乘指数]\label{ex:c13q5}
计算 $\displaystyle\int xe^{5x}\,dx$，以及 $\displaystyle\int_0^1xe^{5x}\,dx$。
\par\textit{改编训练。}\par\hyperref[ans:c13q5]{解答见第~\pageref{ans:c13q5} 页。}
\end{example}
\begin{example}[一次多项式乘指数]\label{ex:c13q6}
计算 $\displaystyle\int xe^{6x}\,dx$，以及 $\displaystyle\int_0^1xe^{6x}\,dx$。
\par\textit{改编训练。}\par\hyperref[ans:c13q6]{解答见第~\pageref{ans:c13q6} 页。}
\end{example}
\begin{example}[幂函数乘对数]\label{ex:c13q7}
在 $x>0$ 上计算 $\displaystyle\int x^{1}\log x\,dx$，并计算 $\displaystyle\int_0^1x^{1}\log x\,dx$。
\par\textit{改编训练。}\par\hyperref[ans:c13q7]{解答见第~\pageref{ans:c13q7} 页。}
\end{example}
\begin{example}[幂函数乘对数]\label{ex:c13q8}
在 $x>0$ 上计算 $\displaystyle\int x^{2}\log x\,dx$，并计算 $\displaystyle\int_0^1x^{2}\log x\,dx$。
\par\textit{改编训练。}\par\hyperref[ans:c13q8]{解答见第~\pageref{ans:c13q8} 页。}
\end{example}
\begin{example}[幂函数乘对数]\label{ex:c13q9}
在 $x>0$ 上计算 $\displaystyle\int x^{3}\log x\,dx$，并计算 $\displaystyle\int_0^1x^{3}\log x\,dx$。
\par\textit{改编训练。}\par\hyperref[ans:c13q9]{解答见第~\pageref{ans:c13q9} 页。}
\end{example}
\begin{example}[幂函数乘对数]\label{ex:c13q10}
在 $x>0$ 上计算 $\displaystyle\int x^{4}\log x\,dx$，并计算 $\displaystyle\int_0^1x^{4}\log x\,dx$。
\par\textit{改编训练。}\par\hyperref[ans:c13q10]{解答见第~\pageref{ans:c13q10} 页。}
\end{example}
\begin{example}[幂函数乘对数]\label{ex:c13q11}
在 $x>0$ 上计算 $\displaystyle\int x^{5}\log x\,dx$，并计算 $\displaystyle\int_0^1x^{5}\log x\,dx$。
\par\textit{改编训练。}\par\hyperref[ans:c13q11]{解答见第~\pageref{ans:c13q11} 页。}
\end{example}
\begin{example}[幂函数乘对数]\label{ex:c13q12}
在 $x>0$ 上计算 $\displaystyle\int x^{6}\log x\,dx$，并计算 $\displaystyle\int_0^1x^{6}\log x\,dx$。
\par\textit{改编训练。}\par\hyperref[ans:c13q12]{解答见第~\pageref{ans:c13q12} 页。}
\end{example}
\begin{example}[互异一次因子的部分分式]\label{ex:c13q13}
计算 $\displaystyle\int\frac{dx}{(x+1)(x+2)}$，并写出 $[0,1]$ 上的定积分。
\par\textit{改编训练。}\par\hyperref[ans:c13q13]{解答见第~\pageref{ans:c13q13} 页。}
\end{example}
\begin{example}[互异一次因子的部分分式]\label{ex:c13q14}
计算 $\displaystyle\int\frac{dx}{(x+1)(x+3)}$，并写出 $[0,1]$ 上的定积分。
\par\textit{改编训练。}\par\hyperref[ans:c13q14]{解答见第~\pageref{ans:c13q14} 页。}
\end{example}
\begin{example}[互异一次因子的部分分式]\label{ex:c13q15}
计算 $\displaystyle\int\frac{dx}{(x+2)(x+3)}$，并写出 $[0,1]$ 上的定积分。
\par\textit{改编训练。}\par\hyperref[ans:c13q15]{解答见第~\pageref{ans:c13q15} 页。}
\end{example}
\begin{example}[互异一次因子的部分分式]\label{ex:c13q16}
计算 $\displaystyle\int\frac{dx}{(x+2)(x+5)}$，并写出 $[0,1]$ 上的定积分。
\par\textit{改编训练。}\par\hyperref[ans:c13q16]{解答见第~\pageref{ans:c13q16} 页。}
\end{example}
\begin{example}[互异一次因子的部分分式]\label{ex:c13q17}
计算 $\displaystyle\int\frac{dx}{(x+3)(x+4)}$，并写出 $[0,1]$ 上的定积分。
\par\textit{改编训练。}\par\hyperref[ans:c13q17]{解答见第~\pageref{ans:c13q17} 页。}
\end{example}
\begin{example}[互异一次因子的部分分式]\label{ex:c13q18}
计算 $\displaystyle\int\frac{dx}{(x+3)(x+7)}$，并写出 $[0,1]$ 上的定积分。
\par\textit{改编训练。}\par\hyperref[ans:c13q18]{解答见第~\pageref{ans:c13q18} 页。}
\end{example}
\begin{example}[不可约二次因子的积分]\label{ex:c13q19}
计算 $\displaystyle\int\frac{x+1}{x^2+1}\,dx$，并给出验证。
\par\textit{改编训练。}\par\hyperref[ans:c13q19]{解答见第~\pageref{ans:c13q19} 页。}
\end{example}
\begin{example}[不可约二次因子的积分]\label{ex:c13q20}
计算 $\displaystyle\int\frac{x+1}{x^2+4}\,dx$，并给出验证。
\par\textit{改编训练。}\par\hyperref[ans:c13q20]{解答见第~\pageref{ans:c13q20} 页。}
\end{example}
\begin{example}[不可约二次因子的积分]\label{ex:c13q21}
计算 $\displaystyle\int\frac{x+1}{x^2+9}\,dx$，并给出验证。
\par\textit{改编训练。}\par\hyperref[ans:c13q21]{解答见第~\pageref{ans:c13q21} 页。}
\end{example}
\begin{example}[不可约二次因子的积分]\label{ex:c13q22}
计算 $\displaystyle\int\frac{x+1}{x^2+16}\,dx$，并给出验证。
\par\textit{改编训练。}\par\hyperref[ans:c13q22]{解答见第~\pageref{ans:c13q22} 页。}
\end{example}
\begin{example}[不可约二次因子的积分]\label{ex:c13q23}
计算 $\displaystyle\int\frac{x+1}{x^2+25}\,dx$，并给出验证。
\par\textit{改编训练。}\par\hyperref[ans:c13q23]{解答见第~\pageref{ans:c13q23} 页。}
\end{example}
\begin{example}[不可约二次因子的积分]\label{ex:c13q24}
计算 $\displaystyle\int\frac{x+1}{x^2+36}\,dx$，并给出验证。
\par\textit{改编训练。}\par\hyperref[ans:c13q24]{解答见第~\pageref{ans:c13q24} 页。}
\end{example}
\subsection{解答}
\begin{solve}[习题~\ref{ex:c13q1}：一次多项式乘指数]\phantomsection\label{ans:c13q1}
选 $u=x$、$dv=e^{1x}dx$，则 $du=dx$、$v=e^{1x}/1$。分部积分给出
\[
\int xe^{1x}\,dx=\frac{xe^{1x}}{1}-\frac1{1}\int e^{1x}\,dx=e^{1x}\left(\frac x{1}-\frac1{1}\right)+C.
\]
定积分使用该原函数在两端的差：
\[
\int_0^1xe^{1x}\,dx=0e^{1}+1.
\]
若反过来把指数函数作为需要求导的部分，多项式次数会在积分时上升，通常不会简化问题。最后对原函数求导，两个不含 $x$ 的指数项相互抵消。
\end{solve}

\begin{solve}[习题~\ref{ex:c13q2}：一次多项式乘指数]\phantomsection\label{ans:c13q2}
选 $u=x$、$dv=e^{2x}dx$，则 $du=dx$、$v=e^{2x}/2$。分部积分给出
\[
\int xe^{2x}\,dx=\frac{xe^{2x}}{2}-\frac1{2}\int e^{2x}\,dx=e^{2x}\left(\frac x{2}-\frac1{4}\right)+C.
\]
定积分使用该原函数在两端的差：
\[
\int_0^1xe^{2x}\,dx=\frac{1}{4}e^{2}+\frac{1}{4}.
\]
若反过来把指数函数作为需要求导的部分，多项式次数会在积分时上升，通常不会简化问题。最后对原函数求导，两个不含 $x$ 的指数项相互抵消。
\end{solve}

\begin{solve}[习题~\ref{ex:c13q3}：一次多项式乘指数]\phantomsection\label{ans:c13q3}
选 $u=x$、$dv=e^{3x}dx$，则 $du=dx$、$v=e^{3x}/3$。分部积分给出
\[
\int xe^{3x}\,dx=\frac{xe^{3x}}{3}-\frac1{3}\int e^{3x}\,dx=e^{3x}\left(\frac x{3}-\frac1{9}\right)+C.
\]
定积分使用该原函数在两端的差：
\[
\int_0^1xe^{3x}\,dx=\frac{2}{9}e^{3}+\frac{1}{9}.
\]
若反过来把指数函数作为需要求导的部分，多项式次数会在积分时上升，通常不会简化问题。最后对原函数求导，两个不含 $x$ 的指数项相互抵消。
\end{solve}

\begin{solve}[习题~\ref{ex:c13q4}：一次多项式乘指数]\phantomsection\label{ans:c13q4}
选 $u=x$、$dv=e^{4x}dx$，则 $du=dx$、$v=e^{4x}/4$。分部积分给出
\[
\int xe^{4x}\,dx=\frac{xe^{4x}}{4}-\frac1{4}\int e^{4x}\,dx=e^{4x}\left(\frac x{4}-\frac1{16}\right)+C.
\]
定积分使用该原函数在两端的差：
\[
\int_0^1xe^{4x}\,dx=\frac{3}{16}e^{4}+\frac{1}{16}.
\]
若反过来把指数函数作为需要求导的部分，多项式次数会在积分时上升，通常不会简化问题。最后对原函数求导，两个不含 $x$ 的指数项相互抵消。
\end{solve}

\begin{solve}[习题~\ref{ex:c13q5}：一次多项式乘指数]\phantomsection\label{ans:c13q5}
选 $u=x$、$dv=e^{5x}dx$，则 $du=dx$、$v=e^{5x}/5$。分部积分给出
\[
\int xe^{5x}\,dx=\frac{xe^{5x}}{5}-\frac1{5}\int e^{5x}\,dx=e^{5x}\left(\frac x{5}-\frac1{25}\right)+C.
\]
定积分使用该原函数在两端的差：
\[
\int_0^1xe^{5x}\,dx=\frac{4}{25}e^{5}+\frac{1}{25}.
\]
若反过来把指数函数作为需要求导的部分，多项式次数会在积分时上升，通常不会简化问题。最后对原函数求导，两个不含 $x$ 的指数项相互抵消。
\end{solve}

\begin{solve}[习题~\ref{ex:c13q6}：一次多项式乘指数]\phantomsection\label{ans:c13q6}
选 $u=x$、$dv=e^{6x}dx$，则 $du=dx$、$v=e^{6x}/6$。分部积分给出
\[
\int xe^{6x}\,dx=\frac{xe^{6x}}{6}-\frac1{6}\int e^{6x}\,dx=e^{6x}\left(\frac x{6}-\frac1{36}\right)+C.
\]
定积分使用该原函数在两端的差：
\[
\int_0^1xe^{6x}\,dx=\frac{5}{36}e^{6}+\frac{1}{36}.
\]
若反过来把指数函数作为需要求导的部分，多项式次数会在积分时上升，通常不会简化问题。最后对原函数求导，两个不含 $x$ 的指数项相互抵消。
\end{solve}

\begin{solve}[习题~\ref{ex:c13q7}：幂函数乘对数]\phantomsection\label{ans:c13q7}
取 $u=\log x$、$dv=x^{1}dx$，可得
\[
\int x^{1}\log x\,dx=\frac{x^{2}\log x}{2}-\frac{x^{2}}{4}+C.
\]
第二个积分在 $0$ 处是反常积分，应先以 $\varepsilon>0$ 为下限。由于 $x^{2}\log x\to0$（令 $x=e^{-t}$ 可见），故
\[
\int_0^1x^{1}\log x\,dx=-\frac{1}{4}.
\]
被积函数在 $(0,1)$ 上非正，得到的负值与符号相符。端点处原函数存在有限极限，保证这一反常积分收敛。
\end{solve}

\begin{solve}[习题~\ref{ex:c13q8}：幂函数乘对数]\phantomsection\label{ans:c13q8}
取 $u=\log x$、$dv=x^{2}dx$，可得
\[
\int x^{2}\log x\,dx=\frac{x^{3}\log x}{3}-\frac{x^{3}}{9}+C.
\]
第二个积分在 $0$ 处是反常积分，应先以 $\varepsilon>0$ 为下限。由于 $x^{3}\log x\to0$（令 $x=e^{-t}$ 可见），故
\[
\int_0^1x^{2}\log x\,dx=-\frac{1}{9}.
\]
被积函数在 $(0,1)$ 上非正，得到的负值与符号相符。端点处原函数存在有限极限，保证这一反常积分收敛。
\end{solve}

\begin{solve}[习题~\ref{ex:c13q9}：幂函数乘对数]\phantomsection\label{ans:c13q9}
取 $u=\log x$、$dv=x^{3}dx$，可得
\[
\int x^{3}\log x\,dx=\frac{x^{4}\log x}{4}-\frac{x^{4}}{16}+C.
\]
第二个积分在 $0$ 处是反常积分，应先以 $\varepsilon>0$ 为下限。由于 $x^{4}\log x\to0$（令 $x=e^{-t}$ 可见），故
\[
\int_0^1x^{3}\log x\,dx=-\frac{1}{16}.
\]
被积函数在 $(0,1)$ 上非正，得到的负值与符号相符。端点处原函数存在有限极限，保证这一反常积分收敛。
\end{solve}

\begin{solve}[习题~\ref{ex:c13q10}：幂函数乘对数]\phantomsection\label{ans:c13q10}
取 $u=\log x$、$dv=x^{4}dx$，可得
\[
\int x^{4}\log x\,dx=\frac{x^{5}\log x}{5}-\frac{x^{5}}{25}+C.
\]
第二个积分在 $0$ 处是反常积分，应先以 $\varepsilon>0$ 为下限。由于 $x^{5}\log x\to0$（令 $x=e^{-t}$ 可见），故
\[
\int_0^1x^{4}\log x\,dx=-\frac{1}{25}.
\]
被积函数在 $(0,1)$ 上非正，得到的负值与符号相符。端点处原函数存在有限极限，保证这一反常积分收敛。
\end{solve}

\begin{solve}[习题~\ref{ex:c13q11}：幂函数乘对数]\phantomsection\label{ans:c13q11}
取 $u=\log x$、$dv=x^{5}dx$，可得
\[
\int x^{5}\log x\,dx=\frac{x^{6}\log x}{6}-\frac{x^{6}}{36}+C.
\]
第二个积分在 $0$ 处是反常积分，应先以 $\varepsilon>0$ 为下限。由于 $x^{6}\log x\to0$（令 $x=e^{-t}$ 可见），故
\[
\int_0^1x^{5}\log x\,dx=-\frac{1}{36}.
\]
被积函数在 $(0,1)$ 上非正，得到的负值与符号相符。端点处原函数存在有限极限，保证这一反常积分收敛。
\end{solve}

\begin{solve}[习题~\ref{ex:c13q12}：幂函数乘对数]\phantomsection\label{ans:c13q12}
取 $u=\log x$、$dv=x^{6}dx$，可得
\[
\int x^{6}\log x\,dx=\frac{x^{7}\log x}{7}-\frac{x^{7}}{49}+C.
\]
第二个积分在 $0$ 处是反常积分，应先以 $\varepsilon>0$ 为下限。由于 $x^{7}\log x\to0$（令 $x=e^{-t}$ 可见），故
\[
\int_0^1x^{6}\log x\,dx=-\frac{1}{49}.
\]
被积函数在 $(0,1)$ 上非正，得到的负值与符号相符。端点处原函数存在有限极限，保证这一反常积分收敛。
\end{solve}

\begin{solve}[习题~\ref{ex:c13q13}：互异一次因子的部分分式]\phantomsection\label{ans:c13q13}
设分解为 $A/(x+1)+B/(x+2)$，比较分子得 $A+B=0$、$2A+1B=1$，所以 $A=1$、$B=-1$。因此
\[
\int\frac{dx}{(x+1)(x+2)}=1\log\left|\frac{x+1}{x+2}\right|+C.
\]
在 $[0,1]$ 上两个因子都为正，可以直接代入端点，
\[
\int_0^1\frac{dx}{(x+1)(x+2)}=1\log\left(\frac{4}{3}\right).
\]
原函数分别定义在被两个极点分开的区间上，各区间的积分常数可以独立选择。
\end{solve}

\begin{solve}[习题~\ref{ex:c13q14}：互异一次因子的部分分式]\phantomsection\label{ans:c13q14}
设分解为 $A/(x+1)+B/(x+3)$，比较分子得 $A+B=0$、$3A+1B=1$，所以 $A=\frac{1}{2}$、$B=-\frac{1}{2}$。因此
\[
\int\frac{dx}{(x+1)(x+3)}=\frac{1}{2}\log\left|\frac{x+1}{x+3}\right|+C.
\]
在 $[0,1]$ 上两个因子都为正，可以直接代入端点，
\[
\int_0^1\frac{dx}{(x+1)(x+3)}=\frac{1}{2}\log\left(\frac{3}{2}\right).
\]
原函数分别定义在被两个极点分开的区间上，各区间的积分常数可以独立选择。
\end{solve}

\begin{solve}[习题~\ref{ex:c13q15}：互异一次因子的部分分式]\phantomsection\label{ans:c13q15}
设分解为 $A/(x+2)+B/(x+3)$，比较分子得 $A+B=0$、$3A+2B=1$，所以 $A=1$、$B=-1$。因此
\[
\int\frac{dx}{(x+2)(x+3)}=1\log\left|\frac{x+2}{x+3}\right|+C.
\]
在 $[0,1]$ 上两个因子都为正，可以直接代入端点，
\[
\int_0^1\frac{dx}{(x+2)(x+3)}=1\log\left(\frac{9}{8}\right).
\]
原函数分别定义在被两个极点分开的区间上，各区间的积分常数可以独立选择。
\end{solve}

\begin{solve}[习题~\ref{ex:c13q16}：互异一次因子的部分分式]\phantomsection\label{ans:c13q16}
设分解为 $A/(x+2)+B/(x+5)$，比较分子得 $A+B=0$、$5A+2B=1$，所以 $A=\frac{1}{3}$、$B=-\frac{1}{3}$。因此
\[
\int\frac{dx}{(x+2)(x+5)}=\frac{1}{3}\log\left|\frac{x+2}{x+5}\right|+C.
\]
在 $[0,1]$ 上两个因子都为正，可以直接代入端点，
\[
\int_0^1\frac{dx}{(x+2)(x+5)}=\frac{1}{3}\log\left(\frac{5}{4}\right).
\]
原函数分别定义在被两个极点分开的区间上，各区间的积分常数可以独立选择。
\end{solve}

\begin{solve}[习题~\ref{ex:c13q17}：互异一次因子的部分分式]\phantomsection\label{ans:c13q17}
设分解为 $A/(x+3)+B/(x+4)$，比较分子得 $A+B=0$、$4A+3B=1$，所以 $A=1$、$B=-1$。因此
\[
\int\frac{dx}{(x+3)(x+4)}=1\log\left|\frac{x+3}{x+4}\right|+C.
\]
在 $[0,1]$ 上两个因子都为正，可以直接代入端点，
\[
\int_0^1\frac{dx}{(x+3)(x+4)}=1\log\left(\frac{16}{15}\right).
\]
原函数分别定义在被两个极点分开的区间上，各区间的积分常数可以独立选择。
\end{solve}

\begin{solve}[习题~\ref{ex:c13q18}：互异一次因子的部分分式]\phantomsection\label{ans:c13q18}
设分解为 $A/(x+3)+B/(x+7)$，比较分子得 $A+B=0$、$7A+3B=1$，所以 $A=\frac{1}{4}$、$B=-\frac{1}{4}$。因此
\[
\int\frac{dx}{(x+3)(x+7)}=\frac{1}{4}\log\left|\frac{x+3}{x+7}\right|+C.
\]
在 $[0,1]$ 上两个因子都为正，可以直接代入端点，
\[
\int_0^1\frac{dx}{(x+3)(x+7)}=\frac{1}{4}\log\left(\frac{7}{6}\right).
\]
原函数分别定义在被两个极点分开的区间上，各区间的积分常数可以独立选择。
\end{solve}

\begin{solve}[习题~\ref{ex:c13q19}：不可约二次因子的积分]\phantomsection\label{ans:c13q19}
把分子拆成两部分，一部分匹配分母的导数，另一部分匹配反正切公式：
\[
\int\frac{x+1}{x^2+1}\,dx=\int\frac{x}{x^2+1}\,dx+\int\frac{dx}{x^2+1}.
\]
第一项令 $u=x^2+1$，第二项令 $v=x/1$，得到
\[
\frac12\log(x^2+1)+1\arctan(x/1)+C.
\]
求导后两项分别恢复 $x/(x^2+1)$ 与 $1/(x^2+1)$。二次多项式始终为正，因而对数不需要绝对值，且结果在整个实轴上成立。
\end{solve}

\begin{solve}[习题~\ref{ex:c13q20}：不可约二次因子的积分]\phantomsection\label{ans:c13q20}
把分子拆成两部分，一部分匹配分母的导数，另一部分匹配反正切公式：
\[
\int\frac{x+1}{x^2+4}\,dx=\int\frac{x}{x^2+4}\,dx+\int\frac{dx}{x^2+4}.
\]
第一项令 $u=x^2+4$，第二项令 $v=x/2$，得到
\[
\frac12\log(x^2+4)+\frac{1}{2}\arctan(x/2)+C.
\]
求导后两项分别恢复 $x/(x^2+4)$ 与 $1/(x^2+4)$。二次多项式始终为正，因而对数不需要绝对值，且结果在整个实轴上成立。
\end{solve}

\begin{solve}[习题~\ref{ex:c13q21}：不可约二次因子的积分]\phantomsection\label{ans:c13q21}
把分子拆成两部分，一部分匹配分母的导数，另一部分匹配反正切公式：
\[
\int\frac{x+1}{x^2+9}\,dx=\int\frac{x}{x^2+9}\,dx+\int\frac{dx}{x^2+9}.
\]
第一项令 $u=x^2+9$，第二项令 $v=x/3$，得到
\[
\frac12\log(x^2+9)+\frac{1}{3}\arctan(x/3)+C.
\]
求导后两项分别恢复 $x/(x^2+9)$ 与 $1/(x^2+9)$。二次多项式始终为正，因而对数不需要绝对值，且结果在整个实轴上成立。
\end{solve}

\begin{solve}[习题~\ref{ex:c13q22}：不可约二次因子的积分]\phantomsection\label{ans:c13q22}
把分子拆成两部分，一部分匹配分母的导数，另一部分匹配反正切公式：
\[
\int\frac{x+1}{x^2+16}\,dx=\int\frac{x}{x^2+16}\,dx+\int\frac{dx}{x^2+16}.
\]
第一项令 $u=x^2+16$，第二项令 $v=x/4$，得到
\[
\frac12\log(x^2+16)+\frac{1}{4}\arctan(x/4)+C.
\]
求导后两项分别恢复 $x/(x^2+16)$ 与 $1/(x^2+16)$。二次多项式始终为正，因而对数不需要绝对值，且结果在整个实轴上成立。
\end{solve}

\begin{solve}[习题~\ref{ex:c13q23}：不可约二次因子的积分]\phantomsection\label{ans:c13q23}
把分子拆成两部分，一部分匹配分母的导数，另一部分匹配反正切公式：
\[
\int\frac{x+1}{x^2+25}\,dx=\int\frac{x}{x^2+25}\,dx+\int\frac{dx}{x^2+25}.
\]
第一项令 $u=x^2+25$，第二项令 $v=x/5$，得到
\[
\frac12\log(x^2+25)+\frac{1}{5}\arctan(x/5)+C.
\]
求导后两项分别恢复 $x/(x^2+25)$ 与 $1/(x^2+25)$。二次多项式始终为正，因而对数不需要绝对值，且结果在整个实轴上成立。
\end{solve}

\begin{solve}[习题~\ref{ex:c13q24}：不可约二次因子的积分]\phantomsection\label{ans:c13q24}
把分子拆成两部分，一部分匹配分母的导数，另一部分匹配反正切公式：
\[
\int\frac{x+1}{x^2+36}\,dx=\int\frac{x}{x^2+36}\,dx+\int\frac{dx}{x^2+36}.
\]
第一项令 $u=x^2+36$，第二项令 $v=x/6$，得到
\[
\frac12\log(x^2+36)+\frac{1}{6}\arctan(x/6)+C.
\]
求导后两项分别恢复 $x/(x^2+36)$ 与 $1/(x^2+36)$。二次多项式始终为正，因而对数不需要绝对值，且结果在整个实轴上成立。
\end{solve}

\clearpage
\section{三角积分与根式换元}\label{sec:ch14}
三角积分需要在恒等变形和换元之间选择。奇次幂通常可分出一个微分因子；偶次幂可用降幂公式。根式换元还要说明新变量的区间，以确定平方根的符号。

\textbf{资源对应：}CLP-2 \S~1.8、\S~1.9。本章共 24 道编号习题。
\subsection{习题}
\begin{example}[正弦的一次因子]\label{ex:c14q1}
计算 $\displaystyle\int\sin x\cos^{1}x\,dx$ 和 $\displaystyle\int_0^{\pi/2}\sin x\cos^{1}x\,dx$。
\par\textit{改编训练。}\par\hyperref[ans:c14q1]{解答见第~\pageref{ans:c14q1} 页。}
\end{example}
\begin{example}[正弦的一次因子]\label{ex:c14q2}
计算 $\displaystyle\int\sin x\cos^{2}x\,dx$ 和 $\displaystyle\int_0^{\pi/2}\sin x\cos^{2}x\,dx$。
\par\textit{改编训练。}\par\hyperref[ans:c14q2]{解答见第~\pageref{ans:c14q2} 页。}
\end{example}
\begin{example}[正弦的一次因子]\label{ex:c14q3}
计算 $\displaystyle\int\sin x\cos^{3}x\,dx$ 和 $\displaystyle\int_0^{\pi/2}\sin x\cos^{3}x\,dx$。
\par\textit{改编训练。}\par\hyperref[ans:c14q3]{解答见第~\pageref{ans:c14q3} 页。}
\end{example}
\begin{example}[正弦的一次因子]\label{ex:c14q4}
计算 $\displaystyle\int\sin x\cos^{4}x\,dx$ 和 $\displaystyle\int_0^{\pi/2}\sin x\cos^{4}x\,dx$。
\par\textit{改编训练。}\par\hyperref[ans:c14q4]{解答见第~\pageref{ans:c14q4} 页。}
\end{example}
\begin{example}[正弦的一次因子]\label{ex:c14q5}
计算 $\displaystyle\int\sin x\cos^{5}x\,dx$ 和 $\displaystyle\int_0^{\pi/2}\sin x\cos^{5}x\,dx$。
\par\textit{改编训练。}\par\hyperref[ans:c14q5]{解答见第~\pageref{ans:c14q5} 页。}
\end{example}
\begin{example}[正弦的一次因子]\label{ex:c14q6}
计算 $\displaystyle\int\sin x\cos^{6}x\,dx$ 和 $\displaystyle\int_0^{\pi/2}\sin x\cos^{6}x\,dx$。
\par\textit{改编训练。}\par\hyperref[ans:c14q6]{解答见第~\pageref{ans:c14q6} 页。}
\end{example}
\begin{example}[降幂公式与平均值]\label{ex:c14q7}
计算 $\displaystyle\int\sin^2(1x)\,dx$，并求 $\sin^2(1x)$ 在 $[0,\pi]$ 上的平均值。
\par\textit{改编训练。}\par\hyperref[ans:c14q7]{解答见第~\pageref{ans:c14q7} 页。}
\end{example}
\begin{example}[降幂公式与平均值]\label{ex:c14q8}
计算 $\displaystyle\int\sin^2(2x)\,dx$，并求 $\sin^2(2x)$ 在 $[0,\pi]$ 上的平均值。
\par\textit{改编训练。}\par\hyperref[ans:c14q8]{解答见第~\pageref{ans:c14q8} 页。}
\end{example}
\begin{example}[降幂公式与平均值]\label{ex:c14q9}
计算 $\displaystyle\int\sin^2(3x)\,dx$，并求 $\sin^2(3x)$ 在 $[0,\pi]$ 上的平均值。
\par\textit{改编训练。}\par\hyperref[ans:c14q9]{解答见第~\pageref{ans:c14q9} 页。}
\end{example}
\begin{example}[降幂公式与平均值]\label{ex:c14q10}
计算 $\displaystyle\int\sin^2(4x)\,dx$，并求 $\sin^2(4x)$ 在 $[0,\pi]$ 上的平均值。
\par\textit{改编训练。}\par\hyperref[ans:c14q10]{解答见第~\pageref{ans:c14q10} 页。}
\end{example}
\begin{example}[降幂公式与平均值]\label{ex:c14q11}
计算 $\displaystyle\int\sin^2(5x)\,dx$，并求 $\sin^2(5x)$ 在 $[0,\pi]$ 上的平均值。
\par\textit{改编训练。}\par\hyperref[ans:c14q11]{解答见第~\pageref{ans:c14q11} 页。}
\end{example}
\begin{example}[降幂公式与平均值]\label{ex:c14q12}
计算 $\displaystyle\int\sin^2(6x)\,dx$，并求 $\sin^2(6x)$ 在 $[0,\pi]$ 上的平均值。
\par\textit{改编训练。}\par\hyperref[ans:c14q12]{解答见第~\pageref{ans:c14q12} 页。}
\end{example}
\begin{example}[圆根式与三角换元]\label{ex:c14q13}
计算 $\displaystyle\int\frac{dx}{\sqrt{1-x^2}}$，并求 $\displaystyle\int_0^{\frac{1}{2}}\frac{dx}{\sqrt{1-x^2}}$。
\par\textit{改编训练。}\par\hyperref[ans:c14q13]{解答见第~\pageref{ans:c14q13} 页。}
\end{example}
\begin{example}[圆根式与三角换元]\label{ex:c14q14}
计算 $\displaystyle\int\frac{dx}{\sqrt{4-x^2}}$，并求 $\displaystyle\int_0^{1}\frac{dx}{\sqrt{4-x^2}}$。
\par\textit{改编训练。}\par\hyperref[ans:c14q14]{解答见第~\pageref{ans:c14q14} 页。}
\end{example}
\begin{example}[圆根式与三角换元]\label{ex:c14q15}
计算 $\displaystyle\int\frac{dx}{\sqrt{9-x^2}}$，并求 $\displaystyle\int_0^{\frac{3}{2}}\frac{dx}{\sqrt{9-x^2}}$。
\par\textit{改编训练。}\par\hyperref[ans:c14q15]{解答见第~\pageref{ans:c14q15} 页。}
\end{example}
\begin{example}[圆根式与三角换元]\label{ex:c14q16}
计算 $\displaystyle\int\frac{dx}{\sqrt{16-x^2}}$，并求 $\displaystyle\int_0^{2}\frac{dx}{\sqrt{16-x^2}}$。
\par\textit{改编训练。}\par\hyperref[ans:c14q16]{解答见第~\pageref{ans:c14q16} 页。}
\end{example}
\begin{example}[圆根式与三角换元]\label{ex:c14q17}
计算 $\displaystyle\int\frac{dx}{\sqrt{25-x^2}}$，并求 $\displaystyle\int_0^{\frac{5}{2}}\frac{dx}{\sqrt{25-x^2}}$。
\par\textit{改编训练。}\par\hyperref[ans:c14q17]{解答见第~\pageref{ans:c14q17} 页。}
\end{example}
\begin{example}[圆根式与三角换元]\label{ex:c14q18}
计算 $\displaystyle\int\frac{dx}{\sqrt{36-x^2}}$，并求 $\displaystyle\int_0^{3}\frac{dx}{\sqrt{36-x^2}}$。
\par\textit{改编训练。}\par\hyperref[ans:c14q18]{解答见第~\pageref{ans:c14q18} 页。}
\end{example}
\begin{example}[双曲型根式的直接代换]\label{ex:c14q19}
计算 $\displaystyle\int\frac{x}{\sqrt{x^2+1}}\,dx$ 和 $\displaystyle\int_0^{1}\frac{x}{\sqrt{x^2+1}}\,dx$。
\par\textit{改编训练。}\par\hyperref[ans:c14q19]{解答见第~\pageref{ans:c14q19} 页。}
\end{example}
\begin{example}[双曲型根式的直接代换]\label{ex:c14q20}
计算 $\displaystyle\int\frac{x}{\sqrt{x^2+4}}\,dx$ 和 $\displaystyle\int_0^{2}\frac{x}{\sqrt{x^2+4}}\,dx$。
\par\textit{改编训练。}\par\hyperref[ans:c14q20]{解答见第~\pageref{ans:c14q20} 页。}
\end{example}
\begin{example}[双曲型根式的直接代换]\label{ex:c14q21}
计算 $\displaystyle\int\frac{x}{\sqrt{x^2+9}}\,dx$ 和 $\displaystyle\int_0^{3}\frac{x}{\sqrt{x^2+9}}\,dx$。
\par\textit{改编训练。}\par\hyperref[ans:c14q21]{解答见第~\pageref{ans:c14q21} 页。}
\end{example}
\begin{example}[双曲型根式的直接代换]\label{ex:c14q22}
计算 $\displaystyle\int\frac{x}{\sqrt{x^2+16}}\,dx$ 和 $\displaystyle\int_0^{4}\frac{x}{\sqrt{x^2+16}}\,dx$。
\par\textit{改编训练。}\par\hyperref[ans:c14q22]{解答见第~\pageref{ans:c14q22} 页。}
\end{example}
\begin{example}[双曲型根式的直接代换]\label{ex:c14q23}
计算 $\displaystyle\int\frac{x}{\sqrt{x^2+25}}\,dx$ 和 $\displaystyle\int_0^{5}\frac{x}{\sqrt{x^2+25}}\,dx$。
\par\textit{改编训练。}\par\hyperref[ans:c14q23]{解答见第~\pageref{ans:c14q23} 页。}
\end{example}
\begin{example}[双曲型根式的直接代换]\label{ex:c14q24}
计算 $\displaystyle\int\frac{x}{\sqrt{x^2+36}}\,dx$ 和 $\displaystyle\int_0^{6}\frac{x}{\sqrt{x^2+36}}\,dx$。
\par\textit{改编训练。}\par\hyperref[ans:c14q24]{解答见第~\pageref{ans:c14q24} 页。}
\end{example}
\subsection{解答}
\begin{solve}[习题~\ref{ex:c14q1}：正弦的一次因子]\phantomsection\label{ans:c14q1}
令 $u=\cos x$，则 $du=-\sin x\,dx$。因此
\[
\int\sin x\cos^{1}x\,dx=-\frac{\cos^{2}x}{2}+C.
\]
对于定积分，$x$ 从 $0$ 到 $\pi/2$ 时，$u$ 从 $1$ 降到 $0$；微分中的负号与上下限反向必须同时处理：
\[
\int_0^{\pi/2}\sin x\cos^{1}x\,dx=-\int_1^0u^{1}\,du=\frac{1}{2}.
\]
被积函数在该区间非负，所以定积分也应非负，这可以发现换元时漏掉负号造成的错误。
\end{solve}

\begin{solve}[习题~\ref{ex:c14q2}：正弦的一次因子]\phantomsection\label{ans:c14q2}
令 $u=\cos x$，则 $du=-\sin x\,dx$。因此
\[
\int\sin x\cos^{2}x\,dx=-\frac{\cos^{3}x}{3}+C.
\]
对于定积分，$x$ 从 $0$ 到 $\pi/2$ 时，$u$ 从 $1$ 降到 $0$；微分中的负号与上下限反向必须同时处理：
\[
\int_0^{\pi/2}\sin x\cos^{2}x\,dx=-\int_1^0u^{2}\,du=\frac{1}{3}.
\]
被积函数在该区间非负，所以定积分也应非负，这可以发现换元时漏掉负号造成的错误。
\end{solve}

\begin{solve}[习题~\ref{ex:c14q3}：正弦的一次因子]\phantomsection\label{ans:c14q3}
令 $u=\cos x$，则 $du=-\sin x\,dx$。因此
\[
\int\sin x\cos^{3}x\,dx=-\frac{\cos^{4}x}{4}+C.
\]
对于定积分，$x$ 从 $0$ 到 $\pi/2$ 时，$u$ 从 $1$ 降到 $0$；微分中的负号与上下限反向必须同时处理：
\[
\int_0^{\pi/2}\sin x\cos^{3}x\,dx=-\int_1^0u^{3}\,du=\frac{1}{4}.
\]
被积函数在该区间非负，所以定积分也应非负，这可以发现换元时漏掉负号造成的错误。
\end{solve}

\begin{solve}[习题~\ref{ex:c14q4}：正弦的一次因子]\phantomsection\label{ans:c14q4}
令 $u=\cos x$，则 $du=-\sin x\,dx$。因此
\[
\int\sin x\cos^{4}x\,dx=-\frac{\cos^{5}x}{5}+C.
\]
对于定积分，$x$ 从 $0$ 到 $\pi/2$ 时，$u$ 从 $1$ 降到 $0$；微分中的负号与上下限反向必须同时处理：
\[
\int_0^{\pi/2}\sin x\cos^{4}x\,dx=-\int_1^0u^{4}\,du=\frac{1}{5}.
\]
被积函数在该区间非负，所以定积分也应非负，这可以发现换元时漏掉负号造成的错误。
\end{solve}

\begin{solve}[习题~\ref{ex:c14q5}：正弦的一次因子]\phantomsection\label{ans:c14q5}
令 $u=\cos x$，则 $du=-\sin x\,dx$。因此
\[
\int\sin x\cos^{5}x\,dx=-\frac{\cos^{6}x}{6}+C.
\]
对于定积分，$x$ 从 $0$ 到 $\pi/2$ 时，$u$ 从 $1$ 降到 $0$；微分中的负号与上下限反向必须同时处理：
\[
\int_0^{\pi/2}\sin x\cos^{5}x\,dx=-\int_1^0u^{5}\,du=\frac{1}{6}.
\]
被积函数在该区间非负，所以定积分也应非负，这可以发现换元时漏掉负号造成的错误。
\end{solve}

\begin{solve}[习题~\ref{ex:c14q6}：正弦的一次因子]\phantomsection\label{ans:c14q6}
令 $u=\cos x$，则 $du=-\sin x\,dx$。因此
\[
\int\sin x\cos^{6}x\,dx=-\frac{\cos^{7}x}{7}+C.
\]
对于定积分，$x$ 从 $0$ 到 $\pi/2$ 时，$u$ 从 $1$ 降到 $0$；微分中的负号与上下限反向必须同时处理：
\[
\int_0^{\pi/2}\sin x\cos^{6}x\,dx=-\int_1^0u^{6}\,du=\frac{1}{7}.
\]
被积函数在该区间非负，所以定积分也应非负，这可以发现换元时漏掉负号造成的错误。
\end{solve}

\begin{solve}[习题~\ref{ex:c14q7}：降幂公式与平均值]\phantomsection\label{ans:c14q7}
使用 $\sin^2t=(1-\cos2t)/2$，
\[
\int\sin^2(1x)\,dx=\frac x2-\frac{\sin(2x)}{4}+C.
\]
因为 $1$ 为整数，$\sin(2\pi)=0$，所以定积分等于 $\pi/2$。平均值为
\[
\frac1\pi\int_0^\pi\sin^2(1x)\,dx=\frac12.
\]
平均值是函数积分除以区间长度，而不是把函数在两端的值取平均。这里两个端点函数值均为零，但内部多数点为正，平均值自然不是零。
\end{solve}

\begin{solve}[习题~\ref{ex:c14q8}：降幂公式与平均值]\phantomsection\label{ans:c14q8}
使用 $\sin^2t=(1-\cos2t)/2$，
\[
\int\sin^2(2x)\,dx=\frac x2-\frac{\sin(4x)}{8}+C.
\]
因为 $2$ 为整数，$\sin(4\pi)=0$，所以定积分等于 $\pi/2$。平均值为
\[
\frac1\pi\int_0^\pi\sin^2(2x)\,dx=\frac12.
\]
平均值是函数积分除以区间长度，而不是把函数在两端的值取平均。这里两个端点函数值均为零，但内部多数点为正，平均值自然不是零。
\end{solve}

\begin{solve}[习题~\ref{ex:c14q9}：降幂公式与平均值]\phantomsection\label{ans:c14q9}
使用 $\sin^2t=(1-\cos2t)/2$，
\[
\int\sin^2(3x)\,dx=\frac x2-\frac{\sin(6x)}{12}+C.
\]
因为 $3$ 为整数，$\sin(6\pi)=0$，所以定积分等于 $\pi/2$。平均值为
\[
\frac1\pi\int_0^\pi\sin^2(3x)\,dx=\frac12.
\]
平均值是函数积分除以区间长度，而不是把函数在两端的值取平均。这里两个端点函数值均为零，但内部多数点为正，平均值自然不是零。
\end{solve}

\begin{solve}[习题~\ref{ex:c14q10}：降幂公式与平均值]\phantomsection\label{ans:c14q10}
使用 $\sin^2t=(1-\cos2t)/2$，
\[
\int\sin^2(4x)\,dx=\frac x2-\frac{\sin(8x)}{16}+C.
\]
因为 $4$ 为整数，$\sin(8\pi)=0$，所以定积分等于 $\pi/2$。平均值为
\[
\frac1\pi\int_0^\pi\sin^2(4x)\,dx=\frac12.
\]
平均值是函数积分除以区间长度，而不是把函数在两端的值取平均。这里两个端点函数值均为零，但内部多数点为正，平均值自然不是零。
\end{solve}

\begin{solve}[习题~\ref{ex:c14q11}：降幂公式与平均值]\phantomsection\label{ans:c14q11}
使用 $\sin^2t=(1-\cos2t)/2$，
\[
\int\sin^2(5x)\,dx=\frac x2-\frac{\sin(10x)}{20}+C.
\]
因为 $5$ 为整数，$\sin(10\pi)=0$，所以定积分等于 $\pi/2$。平均值为
\[
\frac1\pi\int_0^\pi\sin^2(5x)\,dx=\frac12.
\]
平均值是函数积分除以区间长度，而不是把函数在两端的值取平均。这里两个端点函数值均为零，但内部多数点为正，平均值自然不是零。
\end{solve}

\begin{solve}[习题~\ref{ex:c14q12}：降幂公式与平均值]\phantomsection\label{ans:c14q12}
使用 $\sin^2t=(1-\cos2t)/2$，
\[
\int\sin^2(6x)\,dx=\frac x2-\frac{\sin(12x)}{24}+C.
\]
因为 $6$ 为整数，$\sin(12\pi)=0$，所以定积分等于 $\pi/2$。平均值为
\[
\frac1\pi\int_0^\pi\sin^2(6x)\,dx=\frac12.
\]
平均值是函数积分除以区间长度，而不是把函数在两端的值取平均。这里两个端点函数值均为零，但内部多数点为正，平均值自然不是零。
\end{solve}

\begin{solve}[习题~\ref{ex:c14q13}：圆根式与三角换元]\phantomsection\label{ans:c14q13}
在 $|x|<1$ 上令 $x=1\sin t$，选 $-\pi/2<t<\pi/2$，此时 $\cos t>0$，故平方根为 $1\cos t$。于是
\[
\int\frac{dx}{\sqrt{1-x^2}}=\int dt=\arcsin(x/1)+C.
\]
积分端点对应 $t=0$ 与 $t=\pi/6$，因此定积分为 $\pi/6$。如果不指定 $t$ 的范围，直接将 $\sqrt{\cos^2t}$ 写成 $\cos t$ 就可能丢掉绝对值。
\end{solve}

\begin{solve}[习题~\ref{ex:c14q14}：圆根式与三角换元]\phantomsection\label{ans:c14q14}
在 $|x|<2$ 上令 $x=2\sin t$，选 $-\pi/2<t<\pi/2$，此时 $\cos t>0$，故平方根为 $2\cos t$。于是
\[
\int\frac{dx}{\sqrt{4-x^2}}=\int dt=\arcsin(x/2)+C.
\]
积分端点对应 $t=0$ 与 $t=\pi/6$，因此定积分为 $\pi/6$。如果不指定 $t$ 的范围，直接将 $\sqrt{\cos^2t}$ 写成 $\cos t$ 就可能丢掉绝对值。
\end{solve}

\begin{solve}[习题~\ref{ex:c14q15}：圆根式与三角换元]\phantomsection\label{ans:c14q15}
在 $|x|<3$ 上令 $x=3\sin t$，选 $-\pi/2<t<\pi/2$，此时 $\cos t>0$，故平方根为 $3\cos t$。于是
\[
\int\frac{dx}{\sqrt{9-x^2}}=\int dt=\arcsin(x/3)+C.
\]
积分端点对应 $t=0$ 与 $t=\pi/6$，因此定积分为 $\pi/6$。如果不指定 $t$ 的范围，直接将 $\sqrt{\cos^2t}$ 写成 $\cos t$ 就可能丢掉绝对值。
\end{solve}

\begin{solve}[习题~\ref{ex:c14q16}：圆根式与三角换元]\phantomsection\label{ans:c14q16}
在 $|x|<4$ 上令 $x=4\sin t$，选 $-\pi/2<t<\pi/2$，此时 $\cos t>0$，故平方根为 $4\cos t$。于是
\[
\int\frac{dx}{\sqrt{16-x^2}}=\int dt=\arcsin(x/4)+C.
\]
积分端点对应 $t=0$ 与 $t=\pi/6$，因此定积分为 $\pi/6$。如果不指定 $t$ 的范围，直接将 $\sqrt{\cos^2t}$ 写成 $\cos t$ 就可能丢掉绝对值。
\end{solve}

\begin{solve}[习题~\ref{ex:c14q17}：圆根式与三角换元]\phantomsection\label{ans:c14q17}
在 $|x|<5$ 上令 $x=5\sin t$，选 $-\pi/2<t<\pi/2$，此时 $\cos t>0$，故平方根为 $5\cos t$。于是
\[
\int\frac{dx}{\sqrt{25-x^2}}=\int dt=\arcsin(x/5)+C.
\]
积分端点对应 $t=0$ 与 $t=\pi/6$，因此定积分为 $\pi/6$。如果不指定 $t$ 的范围，直接将 $\sqrt{\cos^2t}$ 写成 $\cos t$ 就可能丢掉绝对值。
\end{solve}

\begin{solve}[习题~\ref{ex:c14q18}：圆根式与三角换元]\phantomsection\label{ans:c14q18}
在 $|x|<6$ 上令 $x=6\sin t$，选 $-\pi/2<t<\pi/2$，此时 $\cos t>0$，故平方根为 $6\cos t$。于是
\[
\int\frac{dx}{\sqrt{36-x^2}}=\int dt=\arcsin(x/6)+C.
\]
积分端点对应 $t=0$ 与 $t=\pi/6$，因此定积分为 $\pi/6$。如果不指定 $t$ 的范围，直接将 $\sqrt{\cos^2t}$ 写成 $\cos t$ 就可能丢掉绝对值。
\end{solve}

\begin{solve}[习题~\ref{ex:c14q19}：双曲型根式的直接代换]\phantomsection\label{ans:c14q19}
不需要为每个根式都使用三角换元。这里分子与根号内表达式的导数匹配，令 $u=x^2+1$，则 $du=2x\,dx$，从而
\[
\int\frac{x}{\sqrt{x^2+1}}\,dx=\frac12\int u^{-1/2}du=\sqrt{x^2+1}+C.
\]
代入两端，
\[
\int_0^{1}\frac{x}{\sqrt{x^2+1}}\,dx=1(\sqrt2-1).
\]
所得原函数的导数等于原被积函数，且定积分为正，与区间上被积函数的符号一致。识别最短的代换比机械套用某一类公式更可靠。
\end{solve}

\begin{solve}[习题~\ref{ex:c14q20}：双曲型根式的直接代换]\phantomsection\label{ans:c14q20}
不需要为每个根式都使用三角换元。这里分子与根号内表达式的导数匹配，令 $u=x^2+4$，则 $du=2x\,dx$，从而
\[
\int\frac{x}{\sqrt{x^2+4}}\,dx=\frac12\int u^{-1/2}du=\sqrt{x^2+4}+C.
\]
代入两端，
\[
\int_0^{2}\frac{x}{\sqrt{x^2+4}}\,dx=2(\sqrt2-1).
\]
所得原函数的导数等于原被积函数，且定积分为正，与区间上被积函数的符号一致。识别最短的代换比机械套用某一类公式更可靠。
\end{solve}

\begin{solve}[习题~\ref{ex:c14q21}：双曲型根式的直接代换]\phantomsection\label{ans:c14q21}
不需要为每个根式都使用三角换元。这里分子与根号内表达式的导数匹配，令 $u=x^2+9$，则 $du=2x\,dx$，从而
\[
\int\frac{x}{\sqrt{x^2+9}}\,dx=\frac12\int u^{-1/2}du=\sqrt{x^2+9}+C.
\]
代入两端，
\[
\int_0^{3}\frac{x}{\sqrt{x^2+9}}\,dx=3(\sqrt2-1).
\]
所得原函数的导数等于原被积函数，且定积分为正，与区间上被积函数的符号一致。识别最短的代换比机械套用某一类公式更可靠。
\end{solve}

\begin{solve}[习题~\ref{ex:c14q22}：双曲型根式的直接代换]\phantomsection\label{ans:c14q22}
不需要为每个根式都使用三角换元。这里分子与根号内表达式的导数匹配，令 $u=x^2+16$，则 $du=2x\,dx$，从而
\[
\int\frac{x}{\sqrt{x^2+16}}\,dx=\frac12\int u^{-1/2}du=\sqrt{x^2+16}+C.
\]
代入两端，
\[
\int_0^{4}\frac{x}{\sqrt{x^2+16}}\,dx=4(\sqrt2-1).
\]
所得原函数的导数等于原被积函数，且定积分为正，与区间上被积函数的符号一致。识别最短的代换比机械套用某一类公式更可靠。
\end{solve}

\begin{solve}[习题~\ref{ex:c14q23}：双曲型根式的直接代换]\phantomsection\label{ans:c14q23}
不需要为每个根式都使用三角换元。这里分子与根号内表达式的导数匹配，令 $u=x^2+25$，则 $du=2x\,dx$，从而
\[
\int\frac{x}{\sqrt{x^2+25}}\,dx=\frac12\int u^{-1/2}du=\sqrt{x^2+25}+C.
\]
代入两端，
\[
\int_0^{5}\frac{x}{\sqrt{x^2+25}}\,dx=5(\sqrt2-1).
\]
所得原函数的导数等于原被积函数，且定积分为正，与区间上被积函数的符号一致。识别最短的代换比机械套用某一类公式更可靠。
\end{solve}

\begin{solve}[习题~\ref{ex:c14q24}：双曲型根式的直接代换]\phantomsection\label{ans:c14q24}
不需要为每个根式都使用三角换元。这里分子与根号内表达式的导数匹配，令 $u=x^2+36$，则 $du=2x\,dx$，从而
\[
\int\frac{x}{\sqrt{x^2+36}}\,dx=\frac12\int u^{-1/2}du=\sqrt{x^2+36}+C.
\]
代入两端，
\[
\int_0^{6}\frac{x}{\sqrt{x^2+36}}\,dx=6(\sqrt2-1).
\]
所得原函数的导数等于原被积函数，且定积分为正，与区间上被积函数的符号一致。识别最短的代换比机械套用某一类公式更可靠。
\end{solve}

\clearpage
\section{积分技巧的综合训练}\label{sec:ch15}
积分技巧常要组合使用：先作结构识别，再换元或分部积分，最后核对导数与定义域。本章同时要求写出关键中间式，避免只抄最终原函数。

\textbf{资源对应：}CLP-2 \S~1.7--\S~1.10、\S~1.13。本章共 24 道编号习题。
\subsection{习题}
\begin{example}[指数乘正弦的循环分部积分]\label{ex:c15q1}
计算 $I=\int e^{1x}\sin(1x)dx$。
\par\textit{改编训练。}\par\hyperref[ans:c15q1]{解答见第~\pageref{ans:c15q1} 页。}
\end{example}
\begin{example}[指数乘正弦的循环分部积分]\label{ex:c15q2}
计算 $I=\int e^{1x}\sin(2x)dx$。
\par\textit{改编训练。}\par\hyperref[ans:c15q2]{解答见第~\pageref{ans:c15q2} 页。}
\end{example}
\begin{example}[指数乘正弦的循环分部积分]\label{ex:c15q3}
计算 $I=\int e^{2x}\sin(1x)dx$。
\par\textit{改编训练。}\par\hyperref[ans:c15q3]{解答见第~\pageref{ans:c15q3} 页。}
\end{example}
\begin{example}[指数乘正弦的循环分部积分]\label{ex:c15q4}
计算 $I=\int e^{2x}\sin(3x)dx$。
\par\textit{改编训练。}\par\hyperref[ans:c15q4]{解答见第~\pageref{ans:c15q4} 页。}
\end{example}
\begin{example}[指数乘正弦的循环分部积分]\label{ex:c15q5}
计算 $I=\int e^{3x}\sin(2x)dx$。
\par\textit{改编训练。}\par\hyperref[ans:c15q5]{解答见第~\pageref{ans:c15q5} 页。}
\end{example}
\begin{example}[指数乘正弦的循环分部积分]\label{ex:c15q6}
计算 $I=\int e^{3x}\sin(4x)dx$。
\par\textit{改编训练。}\par\hyperref[ans:c15q6]{解答见第~\pageref{ans:c15q6} 页。}
\end{example}
\begin{example}[反三角函数的分部积分]\label{ex:c15q7}
计算 $\int\arctan(1x)dx$。
\par\textit{改编训练。}\par\hyperref[ans:c15q7]{解答见第~\pageref{ans:c15q7} 页。}
\end{example}
\begin{example}[反三角函数的分部积分]\label{ex:c15q8}
计算 $\int\arctan(2x)dx$。
\par\textit{改编训练。}\par\hyperref[ans:c15q8]{解答见第~\pageref{ans:c15q8} 页。}
\end{example}
\begin{example}[反三角函数的分部积分]\label{ex:c15q9}
计算 $\int\arctan(3x)dx$。
\par\textit{改编训练。}\par\hyperref[ans:c15q9]{解答见第~\pageref{ans:c15q9} 页。}
\end{example}
\begin{example}[反三角函数的分部积分]\label{ex:c15q10}
计算 $\int\arctan(4x)dx$。
\par\textit{改编训练。}\par\hyperref[ans:c15q10]{解答见第~\pageref{ans:c15q10} 页。}
\end{example}
\begin{example}[反三角函数的分部积分]\label{ex:c15q11}
计算 $\int\arctan(5x)dx$。
\par\textit{改编训练。}\par\hyperref[ans:c15q11]{解答见第~\pageref{ans:c15q11} 页。}
\end{example}
\begin{example}[反三角函数的分部积分]\label{ex:c15q12}
计算 $\int\arctan(6x)dx$。
\par\textit{改编训练。}\par\hyperref[ans:c15q12]{解答见第~\pageref{ans:c15q12} 页。}
\end{example}
\begin{example}[对数平方的分部积分]\label{ex:c15q13}
在 $x>0$ 上计算 $\displaystyle\int(\log x+1)^2dx$。
\par\textit{改编训练。}\par\hyperref[ans:c15q13]{解答见第~\pageref{ans:c15q13} 页。}
\end{example}
\begin{example}[对数平方的分部积分]\label{ex:c15q14}
在 $x>0$ 上计算 $\displaystyle\int(\log x+2)^2dx$。
\par\textit{改编训练。}\par\hyperref[ans:c15q14]{解答见第~\pageref{ans:c15q14} 页。}
\end{example}
\begin{example}[对数平方的分部积分]\label{ex:c15q15}
在 $x>0$ 上计算 $\displaystyle\int(\log x+3)^2dx$。
\par\textit{改编训练。}\par\hyperref[ans:c15q15]{解答见第~\pageref{ans:c15q15} 页。}
\end{example}
\begin{example}[对数平方的分部积分]\label{ex:c15q16}
在 $x>0$ 上计算 $\displaystyle\int(\log x+4)^2dx$。
\par\textit{改编训练。}\par\hyperref[ans:c15q16]{解答见第~\pageref{ans:c15q16} 页。}
\end{example}
\begin{example}[对数平方的分部积分]\label{ex:c15q17}
在 $x>0$ 上计算 $\displaystyle\int(\log x+5)^2dx$。
\par\textit{改编训练。}\par\hyperref[ans:c15q17]{解答见第~\pageref{ans:c15q17} 页。}
\end{example}
\begin{example}[对数平方的分部积分]\label{ex:c15q18}
在 $x>0$ 上计算 $\displaystyle\int(\log x+6)^2dx$。
\par\textit{改编训练。}\par\hyperref[ans:c15q18]{解答见第~\pageref{ans:c15q18} 页。}
\end{example}
\begin{example}[先拆分再积分的有理式]\label{ex:c15q19}
计算 $\displaystyle\int\frac{x^2}{x^2+1}dx$ 及其在 $[0,1]$ 上的定积分。
\par\textit{改编训练。}\par\hyperref[ans:c15q19]{解答见第~\pageref{ans:c15q19} 页。}
\end{example}
\begin{example}[先拆分再积分的有理式]\label{ex:c15q20}
计算 $\displaystyle\int\frac{x^2}{x^2+4}dx$ 及其在 $[0,2]$ 上的定积分。
\par\textit{改编训练。}\par\hyperref[ans:c15q20]{解答见第~\pageref{ans:c15q20} 页。}
\end{example}
\begin{example}[先拆分再积分的有理式]\label{ex:c15q21}
计算 $\displaystyle\int\frac{x^2}{x^2+9}dx$ 及其在 $[0,3]$ 上的定积分。
\par\textit{改编训练。}\par\hyperref[ans:c15q21]{解答见第~\pageref{ans:c15q21} 页。}
\end{example}
\begin{example}[先拆分再积分的有理式]\label{ex:c15q22}
计算 $\displaystyle\int\frac{x^2}{x^2+16}dx$ 及其在 $[0,4]$ 上的定积分。
\par\textit{改编训练。}\par\hyperref[ans:c15q22]{解答见第~\pageref{ans:c15q22} 页。}
\end{example}
\begin{example}[先拆分再积分的有理式]\label{ex:c15q23}
计算 $\displaystyle\int\frac{x^2}{x^2+25}dx$ 及其在 $[0,5]$ 上的定积分。
\par\textit{改编训练。}\par\hyperref[ans:c15q23]{解答见第~\pageref{ans:c15q23} 页。}
\end{example}
\begin{example}[先拆分再积分的有理式]\label{ex:c15q24}
计算 $\displaystyle\int\frac{x^2}{x^2+36}dx$ 及其在 $[0,6]$ 上的定积分。
\par\textit{改编训练。}\par\hyperref[ans:c15q24]{解答见第~\pageref{ans:c15q24} 页。}
\end{example}
\subsection{解答}
\begin{solve}[习题~\ref{ex:c15q1}：指数乘正弦的循环分部积分]\phantomsection\label{ans:c15q1}
同时记 $J=\int e^{1x}\cos(1x)dx$。分部积分给出 $I=e^{1x}\sin(1x)/1-(1)J$，以及 $J=e^{1x}\cos(1x)/1+(1)I$。代回并移项，得到
\[
I=\frac{e^{1x}}{2}\bigl(1\sin(1x)-1\cos(1x)\bigr)+C.
\]
两次分部积分回到原积分，这并不是计算失败；把它看作含 $I$ 的线性方程即可解出。最后求导可见余弦交叉项抵消，剩下所需被积函数。
\end{solve}

\begin{solve}[习题~\ref{ex:c15q2}：指数乘正弦的循环分部积分]\phantomsection\label{ans:c15q2}
同时记 $J=\int e^{1x}\cos(2x)dx$。分部积分给出 $I=e^{1x}\sin(2x)/1-(2)J$，以及 $J=e^{1x}\cos(2x)/1+(2)I$。代回并移项，得到
\[
I=\frac{e^{1x}}{5}\bigl(1\sin(2x)-2\cos(2x)\bigr)+C.
\]
两次分部积分回到原积分，这并不是计算失败；把它看作含 $I$ 的线性方程即可解出。最后求导可见余弦交叉项抵消，剩下所需被积函数。
\end{solve}

\begin{solve}[习题~\ref{ex:c15q3}：指数乘正弦的循环分部积分]\phantomsection\label{ans:c15q3}
同时记 $J=\int e^{2x}\cos(1x)dx$。分部积分给出 $I=e^{2x}\sin(1x)/2-(\frac{1}{2})J$，以及 $J=e^{2x}\cos(1x)/2+(\frac{1}{2})I$。代回并移项，得到
\[
I=\frac{e^{2x}}{5}\bigl(2\sin(1x)-1\cos(1x)\bigr)+C.
\]
两次分部积分回到原积分，这并不是计算失败；把它看作含 $I$ 的线性方程即可解出。最后求导可见余弦交叉项抵消，剩下所需被积函数。
\end{solve}

\begin{solve}[习题~\ref{ex:c15q4}：指数乘正弦的循环分部积分]\phantomsection\label{ans:c15q4}
同时记 $J=\int e^{2x}\cos(3x)dx$。分部积分给出 $I=e^{2x}\sin(3x)/2-(\frac{3}{2})J$，以及 $J=e^{2x}\cos(3x)/2+(\frac{3}{2})I$。代回并移项，得到
\[
I=\frac{e^{2x}}{13}\bigl(2\sin(3x)-3\cos(3x)\bigr)+C.
\]
两次分部积分回到原积分，这并不是计算失败；把它看作含 $I$ 的线性方程即可解出。最后求导可见余弦交叉项抵消，剩下所需被积函数。
\end{solve}

\begin{solve}[习题~\ref{ex:c15q5}：指数乘正弦的循环分部积分]\phantomsection\label{ans:c15q5}
同时记 $J=\int e^{3x}\cos(2x)dx$。分部积分给出 $I=e^{3x}\sin(2x)/3-(\frac{2}{3})J$，以及 $J=e^{3x}\cos(2x)/3+(\frac{2}{3})I$。代回并移项，得到
\[
I=\frac{e^{3x}}{13}\bigl(3\sin(2x)-2\cos(2x)\bigr)+C.
\]
两次分部积分回到原积分，这并不是计算失败；把它看作含 $I$ 的线性方程即可解出。最后求导可见余弦交叉项抵消，剩下所需被积函数。
\end{solve}

\begin{solve}[习题~\ref{ex:c15q6}：指数乘正弦的循环分部积分]\phantomsection\label{ans:c15q6}
同时记 $J=\int e^{3x}\cos(4x)dx$。分部积分给出 $I=e^{3x}\sin(4x)/3-(\frac{4}{3})J$，以及 $J=e^{3x}\cos(4x)/3+(\frac{4}{3})I$。代回并移项，得到
\[
I=\frac{e^{3x}}{25}\bigl(3\sin(4x)-4\cos(4x)\bigr)+C.
\]
两次分部积分回到原积分，这并不是计算失败；把它看作含 $I$ 的线性方程即可解出。最后求导可见余弦交叉项抵消，剩下所需被积函数。
\end{solve}

\begin{solve}[习题~\ref{ex:c15q7}：反三角函数的分部积分]\phantomsection\label{ans:c15q7}
取 $u=\arctan(1x)$、$dv=dx$，则 $du=1dx/(1+1x^2)$、$v=x$。因此
\[
\int\arctan(1x)dx=x\arctan(1x)-\int\frac{1x}{1+1x^2}dx.
\]
剩下的积分属于分母导数型，得到
\[
x\arctan(1x)-\frac{1}{2}\log(1+1x^2)+C.
\]
反三角函数自身不容易直接积分，但它的导数是有理函数，所以把它选作分部积分中被求导的部分会简化结构。结果在整个实轴上成立。
\end{solve}

\begin{solve}[习题~\ref{ex:c15q8}：反三角函数的分部积分]\phantomsection\label{ans:c15q8}
取 $u=\arctan(2x)$、$dv=dx$，则 $du=2dx/(1+4x^2)$、$v=x$。因此
\[
\int\arctan(2x)dx=x\arctan(2x)-\int\frac{2x}{1+4x^2}dx.
\]
剩下的积分属于分母导数型，得到
\[
x\arctan(2x)-\frac{1}{4}\log(1+4x^2)+C.
\]
反三角函数自身不容易直接积分，但它的导数是有理函数，所以把它选作分部积分中被求导的部分会简化结构。结果在整个实轴上成立。
\end{solve}

\begin{solve}[习题~\ref{ex:c15q9}：反三角函数的分部积分]\phantomsection\label{ans:c15q9}
取 $u=\arctan(3x)$、$dv=dx$，则 $du=3dx/(1+9x^2)$、$v=x$。因此
\[
\int\arctan(3x)dx=x\arctan(3x)-\int\frac{3x}{1+9x^2}dx.
\]
剩下的积分属于分母导数型，得到
\[
x\arctan(3x)-\frac{1}{6}\log(1+9x^2)+C.
\]
反三角函数自身不容易直接积分，但它的导数是有理函数，所以把它选作分部积分中被求导的部分会简化结构。结果在整个实轴上成立。
\end{solve}

\begin{solve}[习题~\ref{ex:c15q10}：反三角函数的分部积分]\phantomsection\label{ans:c15q10}
取 $u=\arctan(4x)$、$dv=dx$，则 $du=4dx/(1+16x^2)$、$v=x$。因此
\[
\int\arctan(4x)dx=x\arctan(4x)-\int\frac{4x}{1+16x^2}dx.
\]
剩下的积分属于分母导数型，得到
\[
x\arctan(4x)-\frac{1}{8}\log(1+16x^2)+C.
\]
反三角函数自身不容易直接积分，但它的导数是有理函数，所以把它选作分部积分中被求导的部分会简化结构。结果在整个实轴上成立。
\end{solve}

\begin{solve}[习题~\ref{ex:c15q11}：反三角函数的分部积分]\phantomsection\label{ans:c15q11}
取 $u=\arctan(5x)$、$dv=dx$，则 $du=5dx/(1+25x^2)$、$v=x$。因此
\[
\int\arctan(5x)dx=x\arctan(5x)-\int\frac{5x}{1+25x^2}dx.
\]
剩下的积分属于分母导数型，得到
\[
x\arctan(5x)-\frac{1}{10}\log(1+25x^2)+C.
\]
反三角函数自身不容易直接积分，但它的导数是有理函数，所以把它选作分部积分中被求导的部分会简化结构。结果在整个实轴上成立。
\end{solve}

\begin{solve}[习题~\ref{ex:c15q12}：反三角函数的分部积分]\phantomsection\label{ans:c15q12}
取 $u=\arctan(6x)$、$dv=dx$，则 $du=6dx/(1+36x^2)$、$v=x$。因此
\[
\int\arctan(6x)dx=x\arctan(6x)-\int\frac{6x}{1+36x^2}dx.
\]
剩下的积分属于分母导数型，得到
\[
x\arctan(6x)-\frac{1}{12}\log(1+36x^2)+C.
\]
反三角函数自身不容易直接积分，但它的导数是有理函数，所以把它选作分部积分中被求导的部分会简化结构。结果在整个实轴上成立。
\end{solve}

\begin{solve}[习题~\ref{ex:c15q13}：对数平方的分部积分]\phantomsection\label{ans:c15q13}
令 $u=(\log x+1)^2$、$dv=dx$，则 $du=2(\log x+1)dx/x$，从而
\[
I=x(\log x+1)^2-2\int(\log x+1)dx.
\]
利用 $\int\log x\,dx=x\log x-x$，化简为
\[
I=x\bigl[(\log x+1)^2-2(\log x+1)+2\bigr]+C.
\]
求导时，外侧 $x$ 的导数和括号内的导数同时出现，线性对数项与常数项恰好抵消。对数幂的分部积分会逐次降低幂次，可据此建立一般递推。
\end{solve}

\begin{solve}[习题~\ref{ex:c15q14}：对数平方的分部积分]\phantomsection\label{ans:c15q14}
令 $u=(\log x+2)^2$、$dv=dx$，则 $du=2(\log x+2)dx/x$，从而
\[
I=x(\log x+2)^2-2\int(\log x+2)dx.
\]
利用 $\int\log x\,dx=x\log x-x$，化简为
\[
I=x\bigl[(\log x+2)^2-2(\log x+2)+2\bigr]+C.
\]
求导时，外侧 $x$ 的导数和括号内的导数同时出现，线性对数项与常数项恰好抵消。对数幂的分部积分会逐次降低幂次，可据此建立一般递推。
\end{solve}

\begin{solve}[习题~\ref{ex:c15q15}：对数平方的分部积分]\phantomsection\label{ans:c15q15}
令 $u=(\log x+3)^2$、$dv=dx$，则 $du=2(\log x+3)dx/x$，从而
\[
I=x(\log x+3)^2-2\int(\log x+3)dx.
\]
利用 $\int\log x\,dx=x\log x-x$，化简为
\[
I=x\bigl[(\log x+3)^2-2(\log x+3)+2\bigr]+C.
\]
求导时，外侧 $x$ 的导数和括号内的导数同时出现，线性对数项与常数项恰好抵消。对数幂的分部积分会逐次降低幂次，可据此建立一般递推。
\end{solve}

\begin{solve}[习题~\ref{ex:c15q16}：对数平方的分部积分]\phantomsection\label{ans:c15q16}
令 $u=(\log x+4)^2$、$dv=dx$，则 $du=2(\log x+4)dx/x$，从而
\[
I=x(\log x+4)^2-2\int(\log x+4)dx.
\]
利用 $\int\log x\,dx=x\log x-x$，化简为
\[
I=x\bigl[(\log x+4)^2-2(\log x+4)+2\bigr]+C.
\]
求导时，外侧 $x$ 的导数和括号内的导数同时出现，线性对数项与常数项恰好抵消。对数幂的分部积分会逐次降低幂次，可据此建立一般递推。
\end{solve}

\begin{solve}[习题~\ref{ex:c15q17}：对数平方的分部积分]\phantomsection\label{ans:c15q17}
令 $u=(\log x+5)^2$、$dv=dx$，则 $du=2(\log x+5)dx/x$，从而
\[
I=x(\log x+5)^2-2\int(\log x+5)dx.
\]
利用 $\int\log x\,dx=x\log x-x$，化简为
\[
I=x\bigl[(\log x+5)^2-2(\log x+5)+2\bigr]+C.
\]
求导时，外侧 $x$ 的导数和括号内的导数同时出现，线性对数项与常数项恰好抵消。对数幂的分部积分会逐次降低幂次，可据此建立一般递推。
\end{solve}

\begin{solve}[习题~\ref{ex:c15q18}：对数平方的分部积分]\phantomsection\label{ans:c15q18}
令 $u=(\log x+6)^2$、$dv=dx$，则 $du=2(\log x+6)dx/x$，从而
\[
I=x(\log x+6)^2-2\int(\log x+6)dx.
\]
利用 $\int\log x\,dx=x\log x-x$，化简为
\[
I=x\bigl[(\log x+6)^2-2(\log x+6)+2\bigr]+C.
\]
求导时，外侧 $x$ 的导数和括号内的导数同时出现，线性对数项与常数项恰好抵消。对数幂的分部积分会逐次降低幂次，可据此建立一般递推。
\end{solve}

\begin{solve}[习题~\ref{ex:c15q19}：先拆分再积分的有理式]\phantomsection\label{ans:c15q19}
分子分母次数相同，先做代数拆分：
\[
\frac{x^2}{x^2+1}=1-\frac{1}{x^2+1}.
\]
所以原函数为
\[
F(x)=x-1\arctan(x/1)+C.
\]
定积分等于 $1(1-\pi/4)$。它为正且小于区间长度 $1$，与被积函数介于0和1之间一致。若不先拆分，直接尝试把分子当成分母导数，就会遇到次数不匹配的问题。
\end{solve}

\begin{solve}[习题~\ref{ex:c15q20}：先拆分再积分的有理式]\phantomsection\label{ans:c15q20}
分子分母次数相同，先做代数拆分：
\[
\frac{x^2}{x^2+4}=1-\frac{4}{x^2+4}.
\]
所以原函数为
\[
F(x)=x-2\arctan(x/2)+C.
\]
定积分等于 $2(1-\pi/4)$。它为正且小于区间长度 $2$，与被积函数介于0和1之间一致。若不先拆分，直接尝试把分子当成分母导数，就会遇到次数不匹配的问题。
\end{solve}

\begin{solve}[习题~\ref{ex:c15q21}：先拆分再积分的有理式]\phantomsection\label{ans:c15q21}
分子分母次数相同，先做代数拆分：
\[
\frac{x^2}{x^2+9}=1-\frac{9}{x^2+9}.
\]
所以原函数为
\[
F(x)=x-3\arctan(x/3)+C.
\]
定积分等于 $3(1-\pi/4)$。它为正且小于区间长度 $3$，与被积函数介于0和1之间一致。若不先拆分，直接尝试把分子当成分母导数，就会遇到次数不匹配的问题。
\end{solve}

\begin{solve}[习题~\ref{ex:c15q22}：先拆分再积分的有理式]\phantomsection\label{ans:c15q22}
分子分母次数相同，先做代数拆分：
\[
\frac{x^2}{x^2+16}=1-\frac{16}{x^2+16}.
\]
所以原函数为
\[
F(x)=x-4\arctan(x/4)+C.
\]
定积分等于 $4(1-\pi/4)$。它为正且小于区间长度 $4$，与被积函数介于0和1之间一致。若不先拆分，直接尝试把分子当成分母导数，就会遇到次数不匹配的问题。
\end{solve}

\begin{solve}[习题~\ref{ex:c15q23}：先拆分再积分的有理式]\phantomsection\label{ans:c15q23}
分子分母次数相同，先做代数拆分：
\[
\frac{x^2}{x^2+25}=1-\frac{25}{x^2+25}.
\]
所以原函数为
\[
F(x)=x-5\arctan(x/5)+C.
\]
定积分等于 $5(1-\pi/4)$。它为正且小于区间长度 $5$，与被积函数介于0和1之间一致。若不先拆分，直接尝试把分子当成分母导数，就会遇到次数不匹配的问题。
\end{solve}

\begin{solve}[习题~\ref{ex:c15q24}：先拆分再积分的有理式]\phantomsection\label{ans:c15q24}
分子分母次数相同，先做代数拆分：
\[
\frac{x^2}{x^2+36}=1-\frac{36}{x^2+36}.
\]
所以原函数为
\[
F(x)=x-6\arctan(x/6)+C.
\]
定积分等于 $6(1-\pi/4)$。它为正且小于区间长度 $6$，与被积函数介于0和1之间一致。若不先拆分，直接尝试把分子当成分母导数，就会遇到次数不匹配的问题。
\end{solve}

\clearpage
\section{三角积分递推与正交性}\label{sec:ch16}
递推关系可以把高次幂积分降到少数初值，正交性则使不同频率的乘积积分消失。递推前要核对边界项，正交公式中要区分零频率与非零频率。

\textbf{资源对应：}CLP-2 \S~1.7、\S~1.8；傅里叶背景为补充。本章共 24 道编号习题。
\subsection{习题}
\begin{example}[正弦偶次幂的递推]\label{ex:c16q1}
计算 $I_{2}=\int_0^{\pi/2}\sin^{2}x\,dx$，先推导一般递推式。
\par\textit{改编训练。}\par\hyperref[ans:c16q1]{解答见第~\pageref{ans:c16q1} 页。}
\end{example}
\begin{example}[正弦偶次幂的递推]\label{ex:c16q2}
计算 $I_{4}=\int_0^{\pi/2}\sin^{4}x\,dx$，先推导一般递推式。
\par\textit{改编训练。}\par\hyperref[ans:c16q2]{解答见第~\pageref{ans:c16q2} 页。}
\end{example}
\begin{example}[正弦偶次幂的递推]\label{ex:c16q3}
计算 $I_{6}=\int_0^{\pi/2}\sin^{6}x\,dx$，先推导一般递推式。
\par\textit{改编训练。}\par\hyperref[ans:c16q3]{解答见第~\pageref{ans:c16q3} 页。}
\end{example}
\begin{example}[正弦偶次幂的递推]\label{ex:c16q4}
计算 $I_{8}=\int_0^{\pi/2}\sin^{8}x\,dx$，先推导一般递推式。
\par\textit{改编训练。}\par\hyperref[ans:c16q4]{解答见第~\pageref{ans:c16q4} 页。}
\end{example}
\begin{example}[正弦偶次幂的递推]\label{ex:c16q5}
计算 $I_{10}=\int_0^{\pi/2}\sin^{10}x\,dx$，先推导一般递推式。
\par\textit{改编训练。}\par\hyperref[ans:c16q5]{解答见第~\pageref{ans:c16q5} 页。}
\end{example}
\begin{example}[正弦偶次幂的递推]\label{ex:c16q6}
计算 $I_{12}=\int_0^{\pi/2}\sin^{12}x\,dx$，先推导一般递推式。
\par\textit{改编训练。}\par\hyperref[ans:c16q6]{解答见第~\pageref{ans:c16q6} 页。}
\end{example}
\begin{example}[正弦奇次幂的递推]\label{ex:c16q7}
计算 $\int_0^{\pi/2}\sin^{3}x\,dx$，并与偶次幂积分的形式比较。
\par\textit{改编训练。}\par\hyperref[ans:c16q7]{解答见第~\pageref{ans:c16q7} 页。}
\end{example}
\begin{example}[正弦奇次幂的递推]\label{ex:c16q8}
计算 $\int_0^{\pi/2}\sin^{5}x\,dx$，并与偶次幂积分的形式比较。
\par\textit{改编训练。}\par\hyperref[ans:c16q8]{解答见第~\pageref{ans:c16q8} 页。}
\end{example}
\begin{example}[正弦奇次幂的递推]\label{ex:c16q9}
计算 $\int_0^{\pi/2}\sin^{7}x\,dx$，并与偶次幂积分的形式比较。
\par\textit{改编训练。}\par\hyperref[ans:c16q9]{解答见第~\pageref{ans:c16q9} 页。}
\end{example}
\begin{example}[正弦奇次幂的递推]\label{ex:c16q10}
计算 $\int_0^{\pi/2}\sin^{9}x\,dx$，并与偶次幂积分的形式比较。
\par\textit{改编训练。}\par\hyperref[ans:c16q10]{解答见第~\pageref{ans:c16q10} 页。}
\end{example}
\begin{example}[正弦奇次幂的递推]\label{ex:c16q11}
计算 $\int_0^{\pi/2}\sin^{11}x\,dx$，并与偶次幂积分的形式比较。
\par\textit{改编训练。}\par\hyperref[ans:c16q11]{解答见第~\pageref{ans:c16q11} 页。}
\end{example}
\begin{example}[正弦奇次幂的递推]\label{ex:c16q12}
计算 $\int_0^{\pi/2}\sin^{13}x\,dx$，并与偶次幂积分的形式比较。
\par\textit{改编训练。}\par\hyperref[ans:c16q12]{解答见第~\pageref{ans:c16q12} 页。}
\end{example}
\begin{example}[余弦函数的正交积分]\label{ex:c16q13}
计算 $\int_{-\pi}^\pi\cos(1x)\cos(2x)dx$，说明频率相同时为什么不同。
\par\textit{改编训练。}\par\hyperref[ans:c16q13]{解答见第~\pageref{ans:c16q13} 页。}
\end{example}
\begin{example}[余弦函数的正交积分]\label{ex:c16q14}
计算 $\int_{-\pi}^\pi\cos(2x)\cos(3x)dx$，说明频率相同时为什么不同。
\par\textit{改编训练。}\par\hyperref[ans:c16q14]{解答见第~\pageref{ans:c16q14} 页。}
\end{example}
\begin{example}[余弦函数的正交积分]\label{ex:c16q15}
计算 $\int_{-\pi}^\pi\cos(3x)\cos(3x)dx$，说明频率相同时为什么不同。
\par\textit{改编训练。}\par\hyperref[ans:c16q15]{解答见第~\pageref{ans:c16q15} 页。}
\end{example}
\begin{example}[余弦函数的正交积分]\label{ex:c16q16}
计算 $\int_{-\pi}^\pi\cos(4x)\cos(2x)dx$，说明频率相同时为什么不同。
\par\textit{改编训练。}\par\hyperref[ans:c16q16]{解答见第~\pageref{ans:c16q16} 页。}
\end{example}
\begin{example}[余弦函数的正交积分]\label{ex:c16q17}
计算 $\int_{-\pi}^\pi\cos(5x)\cos(5x)dx$，说明频率相同时为什么不同。
\par\textit{改编训练。}\par\hyperref[ans:c16q17]{解答见第~\pageref{ans:c16q17} 页。}
\end{example}
\begin{example}[余弦函数的正交积分]\label{ex:c16q18}
计算 $\int_{-\pi}^\pi\cos(6x)\cos(3x)dx$，说明频率相同时为什么不同。
\par\textit{改编训练。}\par\hyperref[ans:c16q18]{解答见第~\pageref{ans:c16q18} 页。}
\end{example}
\begin{example}[两种奇次幂的混合积分]\label{ex:c16q19}
计算 $\int_0^{\pi/2}\sin^{3}x\cos^{3}x\,dx$。
\par\textit{改编训练。}\par\hyperref[ans:c16q19]{解答见第~\pageref{ans:c16q19} 页。}
\end{example}
\begin{example}[两种奇次幂的混合积分]\label{ex:c16q20}
计算 $\int_0^{\pi/2}\sin^{3}x\cos^{5}x\,dx$。
\par\textit{改编训练。}\par\hyperref[ans:c16q20]{解答见第~\pageref{ans:c16q20} 页。}
\end{example}
\begin{example}[两种奇次幂的混合积分]\label{ex:c16q21}
计算 $\int_0^{\pi/2}\sin^{5}x\cos^{3}x\,dx$。
\par\textit{改编训练。}\par\hyperref[ans:c16q21]{解答见第~\pageref{ans:c16q21} 页。}
\end{example}
\begin{example}[两种奇次幂的混合积分]\label{ex:c16q22}
计算 $\int_0^{\pi/2}\sin^{5}x\cos^{5}x\,dx$。
\par\textit{改编训练。}\par\hyperref[ans:c16q22]{解答见第~\pageref{ans:c16q22} 页。}
\end{example}
\begin{example}[两种奇次幂的混合积分]\label{ex:c16q23}
计算 $\int_0^{\pi/2}\sin^{7}x\cos^{3}x\,dx$。
\par\textit{改编训练。}\par\hyperref[ans:c16q23]{解答见第~\pageref{ans:c16q23} 页。}
\end{example}
\begin{example}[两种奇次幂的混合积分]\label{ex:c16q24}
计算 $\int_0^{\pi/2}\sin^{3}x\cos^{7}x\,dx$。
\par\textit{改编训练。}\par\hyperref[ans:c16q24]{解答见第~\pageref{ans:c16q24} 页。}
\end{example}
\subsection{解答}
\begin{solve}[习题~\ref{ex:c16q1}：正弦偶次幂的递推]\phantomsection\label{ans:c16q1}
令 $I_k=\int_0^{\pi/2}\sin^kx\,dx$。当 $k\ge2$ 时，写成 $\sin^{k-1}x\sin x$ 并分部积分，边界项为零，得到 $I_k=(k-1)\int\sin^{k-2}x\cos^2x$，所以
\[
I_k=(k-1)(I_{k-2}-I_k),\qquad I_k=\frac{k-1}{k}I_{k-2}.
\]
以 $I_0=\pi/2$ 为起点递推，得到
\[
I_{2}=\frac{1}{4}\pi.
\]
偶次链最终落到 $I_0$，奇次链最终落到 $I_1=1$；混淆两个初值会引入多余的 $\pi$ 因子。
\end{solve}

\begin{solve}[习题~\ref{ex:c16q2}：正弦偶次幂的递推]\phantomsection\label{ans:c16q2}
令 $I_k=\int_0^{\pi/2}\sin^kx\,dx$。当 $k\ge2$ 时，写成 $\sin^{k-1}x\sin x$ 并分部积分，边界项为零，得到 $I_k=(k-1)\int\sin^{k-2}x\cos^2x$，所以
\[
I_k=(k-1)(I_{k-2}-I_k),\qquad I_k=\frac{k-1}{k}I_{k-2}.
\]
以 $I_0=\pi/2$ 为起点递推，得到
\[
I_{4}=\frac{3}{16}\pi.
\]
偶次链最终落到 $I_0$，奇次链最终落到 $I_1=1$；混淆两个初值会引入多余的 $\pi$ 因子。
\end{solve}

\begin{solve}[习题~\ref{ex:c16q3}：正弦偶次幂的递推]\phantomsection\label{ans:c16q3}
令 $I_k=\int_0^{\pi/2}\sin^kx\,dx$。当 $k\ge2$ 时，写成 $\sin^{k-1}x\sin x$ 并分部积分，边界项为零，得到 $I_k=(k-1)\int\sin^{k-2}x\cos^2x$，所以
\[
I_k=(k-1)(I_{k-2}-I_k),\qquad I_k=\frac{k-1}{k}I_{k-2}.
\]
以 $I_0=\pi/2$ 为起点递推，得到
\[
I_{6}=\frac{5}{32}\pi.
\]
偶次链最终落到 $I_0$，奇次链最终落到 $I_1=1$；混淆两个初值会引入多余的 $\pi$ 因子。
\end{solve}

\begin{solve}[习题~\ref{ex:c16q4}：正弦偶次幂的递推]\phantomsection\label{ans:c16q4}
令 $I_k=\int_0^{\pi/2}\sin^kx\,dx$。当 $k\ge2$ 时，写成 $\sin^{k-1}x\sin x$ 并分部积分，边界项为零，得到 $I_k=(k-1)\int\sin^{k-2}x\cos^2x$，所以
\[
I_k=(k-1)(I_{k-2}-I_k),\qquad I_k=\frac{k-1}{k}I_{k-2}.
\]
以 $I_0=\pi/2$ 为起点递推，得到
\[
I_{8}=\frac{35}{256}\pi.
\]
偶次链最终落到 $I_0$，奇次链最终落到 $I_1=1$；混淆两个初值会引入多余的 $\pi$ 因子。
\end{solve}

\begin{solve}[习题~\ref{ex:c16q5}：正弦偶次幂的递推]\phantomsection\label{ans:c16q5}
令 $I_k=\int_0^{\pi/2}\sin^kx\,dx$。当 $k\ge2$ 时，写成 $\sin^{k-1}x\sin x$ 并分部积分，边界项为零，得到 $I_k=(k-1)\int\sin^{k-2}x\cos^2x$，所以
\[
I_k=(k-1)(I_{k-2}-I_k),\qquad I_k=\frac{k-1}{k}I_{k-2}.
\]
以 $I_0=\pi/2$ 为起点递推，得到
\[
I_{10}=\frac{63}{512}\pi.
\]
偶次链最终落到 $I_0$，奇次链最终落到 $I_1=1$；混淆两个初值会引入多余的 $\pi$ 因子。
\end{solve}

\begin{solve}[习题~\ref{ex:c16q6}：正弦偶次幂的递推]\phantomsection\label{ans:c16q6}
令 $I_k=\int_0^{\pi/2}\sin^kx\,dx$。当 $k\ge2$ 时，写成 $\sin^{k-1}x\sin x$ 并分部积分，边界项为零，得到 $I_k=(k-1)\int\sin^{k-2}x\cos^2x$，所以
\[
I_k=(k-1)(I_{k-2}-I_k),\qquad I_k=\frac{k-1}{k}I_{k-2}.
\]
以 $I_0=\pi/2$ 为起点递推，得到
\[
I_{12}=\frac{231}{2048}\pi.
\]
偶次链最终落到 $I_0$，奇次链最终落到 $I_1=1$；混淆两个初值会引入多余的 $\pi$ 因子。
\end{solve}

\begin{solve}[习题~\ref{ex:c16q7}：正弦奇次幂的递推]\phantomsection\label{ans:c16q7}
分部积分递推 $I_k=(k-1)I_{k-2}/k$ 对 $k\ge2$ 成立。奇数次每次降低2，最终落到 $I_1=\int_0^{\pi/2}\sin x\,dx=1$，故
\[
I_{3}=\prod_{j=1}^{1}\frac{2j}{2j+1}=\frac{2}{3}.
\]
所以此处得到一个有理数。另一种方法是分出 $\sin x\,dx$，用 $u=\cos x$ 把剩下的偶次正弦幂改写成 $(1-u^2)^m$，展开积分也会给出同样答案。
\end{solve}

\begin{solve}[习题~\ref{ex:c16q8}：正弦奇次幂的递推]\phantomsection\label{ans:c16q8}
分部积分递推 $I_k=(k-1)I_{k-2}/k$ 对 $k\ge2$ 成立。奇数次每次降低2，最终落到 $I_1=\int_0^{\pi/2}\sin x\,dx=1$，故
\[
I_{5}=\prod_{j=1}^{2}\frac{2j}{2j+1}=\frac{8}{15}.
\]
所以此处得到一个有理数。另一种方法是分出 $\sin x\,dx$，用 $u=\cos x$ 把剩下的偶次正弦幂改写成 $(1-u^2)^m$，展开积分也会给出同样答案。
\end{solve}

\begin{solve}[习题~\ref{ex:c16q9}：正弦奇次幂的递推]\phantomsection\label{ans:c16q9}
分部积分递推 $I_k=(k-1)I_{k-2}/k$ 对 $k\ge2$ 成立。奇数次每次降低2，最终落到 $I_1=\int_0^{\pi/2}\sin x\,dx=1$，故
\[
I_{7}=\prod_{j=1}^{3}\frac{2j}{2j+1}=\frac{16}{35}.
\]
所以此处得到一个有理数。另一种方法是分出 $\sin x\,dx$，用 $u=\cos x$ 把剩下的偶次正弦幂改写成 $(1-u^2)^m$，展开积分也会给出同样答案。
\end{solve}

\begin{solve}[习题~\ref{ex:c16q10}：正弦奇次幂的递推]\phantomsection\label{ans:c16q10}
分部积分递推 $I_k=(k-1)I_{k-2}/k$ 对 $k\ge2$ 成立。奇数次每次降低2，最终落到 $I_1=\int_0^{\pi/2}\sin x\,dx=1$，故
\[
I_{9}=\prod_{j=1}^{4}\frac{2j}{2j+1}=\frac{128}{315}.
\]
所以此处得到一个有理数。另一种方法是分出 $\sin x\,dx$，用 $u=\cos x$ 把剩下的偶次正弦幂改写成 $(1-u^2)^m$，展开积分也会给出同样答案。
\end{solve}

\begin{solve}[习题~\ref{ex:c16q11}：正弦奇次幂的递推]\phantomsection\label{ans:c16q11}
分部积分递推 $I_k=(k-1)I_{k-2}/k$ 对 $k\ge2$ 成立。奇数次每次降低2，最终落到 $I_1=\int_0^{\pi/2}\sin x\,dx=1$，故
\[
I_{11}=\prod_{j=1}^{5}\frac{2j}{2j+1}=\frac{256}{693}.
\]
所以此处得到一个有理数。另一种方法是分出 $\sin x\,dx$，用 $u=\cos x$ 把剩下的偶次正弦幂改写成 $(1-u^2)^m$，展开积分也会给出同样答案。
\end{solve}

\begin{solve}[习题~\ref{ex:c16q12}：正弦奇次幂的递推]\phantomsection\label{ans:c16q12}
分部积分递推 $I_k=(k-1)I_{k-2}/k$ 对 $k\ge2$ 成立。奇数次每次降低2，最终落到 $I_1=\int_0^{\pi/2}\sin x\,dx=1$，故
\[
I_{13}=\prod_{j=1}^{6}\frac{2j}{2j+1}=\frac{1024}{3003}.
\]
所以此处得到一个有理数。另一种方法是分出 $\sin x\,dx$，用 $u=\cos x$ 把剩下的偶次正弦幂改写成 $(1-u^2)^m$，展开积分也会给出同样答案。
\end{solve}

\begin{solve}[习题~\ref{ex:c16q13}：余弦函数的正交积分]\phantomsection\label{ans:c16q13}
用积化和差公式 $2\cos(1x)\cos(2x)=\cos((-1)x)+\cos((3)x)$。非零整数频率 $k$ 满足 $\int_{-\pi}^\pi\cos(kx)dx=0$。本题两个频率不同，差频也非零，所以积分等于0。正交性是傅里叶系数公式的基础；不能在频率相等时继续使用只适用于非零频率的“积分为零”结论。
\end{solve}

\begin{solve}[习题~\ref{ex:c16q14}：余弦函数的正交积分]\phantomsection\label{ans:c16q14}
用积化和差公式 $2\cos(2x)\cos(3x)=\cos((-1)x)+\cos((5)x)$。非零整数频率 $k$ 满足 $\int_{-\pi}^\pi\cos(kx)dx=0$。本题两个频率不同，差频也非零，所以积分等于0。正交性是傅里叶系数公式的基础；不能在频率相等时继续使用只适用于非零频率的“积分为零”结论。
\end{solve}

\begin{solve}[习题~\ref{ex:c16q15}：余弦函数的正交积分]\phantomsection\label{ans:c16q15}
用积化和差公式 $2\cos(3x)\cos(3x)=\cos((0)x)+\cos((6)x)$。非零整数频率 $k$ 满足 $\int_{-\pi}^\pi\cos(kx)dx=0$。本题两个频率相同，差频为0，对应常函数1的积分为 $2\pi$；再乘外侧的 $1/2$，得到积分为 $\pi$。正交性是傅里叶系数公式的基础；不能在频率相等时继续使用只适用于非零频率的“积分为零”结论。
\end{solve}

\begin{solve}[习题~\ref{ex:c16q16}：余弦函数的正交积分]\phantomsection\label{ans:c16q16}
用积化和差公式 $2\cos(4x)\cos(2x)=\cos((2)x)+\cos((6)x)$。非零整数频率 $k$ 满足 $\int_{-\pi}^\pi\cos(kx)dx=0$。本题两个频率不同，差频也非零，所以积分等于0。正交性是傅里叶系数公式的基础；不能在频率相等时继续使用只适用于非零频率的“积分为零”结论。
\end{solve}

\begin{solve}[习题~\ref{ex:c16q17}：余弦函数的正交积分]\phantomsection\label{ans:c16q17}
用积化和差公式 $2\cos(5x)\cos(5x)=\cos((0)x)+\cos((10)x)$。非零整数频率 $k$ 满足 $\int_{-\pi}^\pi\cos(kx)dx=0$。本题两个频率相同，差频为0，对应常函数1的积分为 $2\pi$；再乘外侧的 $1/2$，得到积分为 $\pi$。正交性是傅里叶系数公式的基础；不能在频率相等时继续使用只适用于非零频率的“积分为零”结论。
\end{solve}

\begin{solve}[习题~\ref{ex:c16q18}：余弦函数的正交积分]\phantomsection\label{ans:c16q18}
用积化和差公式 $2\cos(6x)\cos(3x)=\cos((3)x)+\cos((9)x)$。非零整数频率 $k$ 满足 $\int_{-\pi}^\pi\cos(kx)dx=0$。本题两个频率不同，差频也非零，所以积分等于0。正交性是傅里叶系数公式的基础；不能在频率相等时继续使用只适用于非零频率的“积分为零”结论。
\end{solve}

\begin{solve}[习题~\ref{ex:c16q19}：两种奇次幂的混合积分]\phantomsection\label{ans:c16q19}
令 $u=\sin^2x$，则 $du=2\sin x\cos x\,dx$，剩下的幂分别化为 $u^{1}$ 和 $(1-u)^{1}$，所以
\[
I=\frac12\int_0^1u^{1}(1-u)^{1}du.
\]
有限次分部积分或多项式展开给出 $\int_0^1u^a(1-u)^bdu=a!b!/(a+b+1)!$（这里 $a,b$ 为非负整数），因此
\[
I=\frac{1}{12}.
\]
换元后上下限是0和1，且变量单调增加，不产生方向翻转。该整数情形完全可以通过初等积分验证，不需要预先定义特殊函数。
\end{solve}

\begin{solve}[习题~\ref{ex:c16q20}：两种奇次幂的混合积分]\phantomsection\label{ans:c16q20}
令 $u=\sin^2x$，则 $du=2\sin x\cos x\,dx$，剩下的幂分别化为 $u^{1}$ 和 $(1-u)^{2}$，所以
\[
I=\frac12\int_0^1u^{1}(1-u)^{2}du.
\]
有限次分部积分或多项式展开给出 $\int_0^1u^a(1-u)^bdu=a!b!/(a+b+1)!$（这里 $a,b$ 为非负整数），因此
\[
I=\frac{1}{24}.
\]
换元后上下限是0和1，且变量单调增加，不产生方向翻转。该整数情形完全可以通过初等积分验证，不需要预先定义特殊函数。
\end{solve}

\begin{solve}[习题~\ref{ex:c16q21}：两种奇次幂的混合积分]\phantomsection\label{ans:c16q21}
令 $u=\sin^2x$，则 $du=2\sin x\cos x\,dx$，剩下的幂分别化为 $u^{2}$ 和 $(1-u)^{1}$，所以
\[
I=\frac12\int_0^1u^{2}(1-u)^{1}du.
\]
有限次分部积分或多项式展开给出 $\int_0^1u^a(1-u)^bdu=a!b!/(a+b+1)!$（这里 $a,b$ 为非负整数），因此
\[
I=\frac{1}{24}.
\]
换元后上下限是0和1，且变量单调增加，不产生方向翻转。该整数情形完全可以通过初等积分验证，不需要预先定义特殊函数。
\end{solve}

\begin{solve}[习题~\ref{ex:c16q22}：两种奇次幂的混合积分]\phantomsection\label{ans:c16q22}
令 $u=\sin^2x$，则 $du=2\sin x\cos x\,dx$，剩下的幂分别化为 $u^{2}$ 和 $(1-u)^{2}$，所以
\[
I=\frac12\int_0^1u^{2}(1-u)^{2}du.
\]
有限次分部积分或多项式展开给出 $\int_0^1u^a(1-u)^bdu=a!b!/(a+b+1)!$（这里 $a,b$ 为非负整数），因此
\[
I=\frac{1}{60}.
\]
换元后上下限是0和1，且变量单调增加，不产生方向翻转。该整数情形完全可以通过初等积分验证，不需要预先定义特殊函数。
\end{solve}

\begin{solve}[习题~\ref{ex:c16q23}：两种奇次幂的混合积分]\phantomsection\label{ans:c16q23}
令 $u=\sin^2x$，则 $du=2\sin x\cos x\,dx$，剩下的幂分别化为 $u^{3}$ 和 $(1-u)^{1}$，所以
\[
I=\frac12\int_0^1u^{3}(1-u)^{1}du.
\]
有限次分部积分或多项式展开给出 $\int_0^1u^a(1-u)^bdu=a!b!/(a+b+1)!$（这里 $a,b$ 为非负整数），因此
\[
I=\frac{1}{40}.
\]
换元后上下限是0和1，且变量单调增加，不产生方向翻转。该整数情形完全可以通过初等积分验证，不需要预先定义特殊函数。
\end{solve}

\begin{solve}[习题~\ref{ex:c16q24}：两种奇次幂的混合积分]\phantomsection\label{ans:c16q24}
令 $u=\sin^2x$，则 $du=2\sin x\cos x\,dx$，剩下的幂分别化为 $u^{1}$ 和 $(1-u)^{3}$，所以
\[
I=\frac12\int_0^1u^{1}(1-u)^{3}du.
\]
有限次分部积分或多项式展开给出 $\int_0^1u^a(1-u)^bdu=a!b!/(a+b+1)!$（这里 $a,b$ 为非负整数），因此
\[
I=\frac{1}{40}.
\]
换元后上下限是0和1，且变量单调增加，不产生方向翻转。该整数情形完全可以通过初等积分验证，不需要预先定义特殊函数。
\end{solve}

\clearpage
\section{定积分的性质与微积分基本定理}\label{sec:ch17}
定积分（definite integral）既是累积量，也是黎曼和的极限。微积分基本定理把积分与原函数连接起来；变上限积分求导时要同时处理上下限的变化。

\textbf{资源对应：}CLP-2 \S~1.1--\S~1.4。本章共 26 道编号习题。
\subsection{习题}
\begin{example}[由原函数信息计算定积分]\label{ex:c17q1}
已知 $F(x)=e^{x^2-3}+1$ 是连续函数 $f$ 的原函数，求 $\int_1^{\sqrt5}f(x)\,dx$。
\par\textit{选译。来源：CLP-2 \S~1.3，2016Q2。}\par\hyperref[ans:c17q1]{解答见第~\pageref{ans:c17q1} 页。}
\end{example}
\begin{example}[纠正换元后的上下限]\label{ex:c17q2}
有人用 $u=\log t$ 把 $\int_1^\pi \cos(\log t)\,dt/t$ 写成 $\int_1^\pi\cos u\,du$。指出错误并求正确值。
\par\textit{选译。来源：CLP-2 \S~1.4，换元推理纠错题。}\par\hyperref[ans:c17q2]{解答见第~\pageref{ans:c17q2} 页。}
\end{example}
\begin{example}[黎曼和的识别]\label{ex:c17q3}
求 $\displaystyle\lim_{n\to\infty}\frac1n\sum_{k=1}^n\left(\frac kn\right)^{1}$，并给出积分解释。
\par\textit{改编训练。}\par\hyperref[ans:c17q3]{解答见第~\pageref{ans:c17q3} 页。}
\end{example}
\begin{example}[黎曼和的识别]\label{ex:c17q4}
求 $\displaystyle\lim_{n\to\infty}\frac1n\sum_{k=1}^n\left(\frac kn\right)^{2}$，并给出积分解释。
\par\textit{改编训练。}\par\hyperref[ans:c17q4]{解答见第~\pageref{ans:c17q4} 页。}
\end{example}
\begin{example}[黎曼和的识别]\label{ex:c17q5}
求 $\displaystyle\lim_{n\to\infty}\frac1n\sum_{k=1}^n\left(\frac kn\right)^{3}$，并给出积分解释。
\par\textit{改编训练。}\par\hyperref[ans:c17q5]{解答见第~\pageref{ans:c17q5} 页。}
\end{example}
\begin{example}[黎曼和的识别]\label{ex:c17q6}
求 $\displaystyle\lim_{n\to\infty}\frac1n\sum_{k=1}^n\left(\frac kn\right)^{4}$，并给出积分解释。
\par\textit{改编训练。}\par\hyperref[ans:c17q6]{解答见第~\pageref{ans:c17q6} 页。}
\end{example}
\begin{example}[黎曼和的识别]\label{ex:c17q7}
求 $\displaystyle\lim_{n\to\infty}\frac1n\sum_{k=1}^n\left(\frac kn\right)^{5}$，并给出积分解释。
\par\textit{改编训练。}\par\hyperref[ans:c17q7]{解答见第~\pageref{ans:c17q7} 页。}
\end{example}
\begin{example}[黎曼和的识别]\label{ex:c17q8}
求 $\displaystyle\lim_{n\to\infty}\frac1n\sum_{k=1}^n\left(\frac kn\right)^{6}$，并给出积分解释。
\par\textit{改编训练。}\par\hyperref[ans:c17q8]{解答见第~\pageref{ans:c17q8} 页。}
\end{example}
\begin{example}[两个积分端点都在变化]\label{ex:c17q9}
设 $F(x)=\displaystyle\int_{1x}^{2x}e^{t^2}dt$。求 $F^{\prime}(x)$、$F^{\prime}(0)$，并说明为什么无须求被积函数的初等原函数。
\par\textit{改编训练。}\par\hyperref[ans:c17q9]{解答见第~\pageref{ans:c17q9} 页。}
\end{example}
\begin{example}[两个积分端点都在变化]\label{ex:c17q10}
设 $F(x)=\displaystyle\int_{2x}^{3x}e^{t^2}dt$。求 $F^{\prime}(x)$、$F^{\prime}(0)$，并说明为什么无须求被积函数的初等原函数。
\par\textit{改编训练。}\par\hyperref[ans:c17q10]{解答见第~\pageref{ans:c17q10} 页。}
\end{example}
\begin{example}[两个积分端点都在变化]\label{ex:c17q11}
设 $F(x)=\displaystyle\int_{3x}^{4x}e^{t^2}dt$。求 $F^{\prime}(x)$、$F^{\prime}(0)$，并说明为什么无须求被积函数的初等原函数。
\par\textit{改编训练。}\par\hyperref[ans:c17q11]{解答见第~\pageref{ans:c17q11} 页。}
\end{example}
\begin{example}[两个积分端点都在变化]\label{ex:c17q12}
设 $F(x)=\displaystyle\int_{1x}^{4x}e^{t^2}dt$。求 $F^{\prime}(x)$、$F^{\prime}(0)$，并说明为什么无须求被积函数的初等原函数。
\par\textit{改编训练。}\par\hyperref[ans:c17q12]{解答见第~\pageref{ans:c17q12} 页。}
\end{example}
\begin{example}[两个积分端点都在变化]\label{ex:c17q13}
设 $F(x)=\displaystyle\int_{2x}^{5x}e^{t^2}dt$。求 $F^{\prime}(x)$、$F^{\prime}(0)$，并说明为什么无须求被积函数的初等原函数。
\par\textit{改编训练。}\par\hyperref[ans:c17q13]{解答见第~\pageref{ans:c17q13} 页。}
\end{example}
\begin{example}[两个积分端点都在变化]\label{ex:c17q14}
设 $F(x)=\displaystyle\int_{3x}^{5x}e^{t^2}dt$。求 $F^{\prime}(x)$、$F^{\prime}(0)$，并说明为什么无须求被积函数的初等原函数。
\par\textit{改编训练。}\par\hyperref[ans:c17q14]{解答见第~\pageref{ans:c17q14} 页。}
\end{example}
\begin{example}[对称区间上的奇偶分解]\label{ex:c17q15}
计算 $\displaystyle\int_{-1}^1\bigl(x^{1}+x^{2}\bigr)dx$，并比较这两个幂函数各自的贡献。
\par\textit{改编训练。}\par\hyperref[ans:c17q15]{解答见第~\pageref{ans:c17q15} 页。}
\end{example}
\begin{example}[对称区间上的奇偶分解]\label{ex:c17q16}
计算 $\displaystyle\int_{-1}^1\bigl(x^{3}+x^{4}\bigr)dx$，并比较这两个幂函数各自的贡献。
\par\textit{改编训练。}\par\hyperref[ans:c17q16]{解答见第~\pageref{ans:c17q16} 页。}
\end{example}
\begin{example}[对称区间上的奇偶分解]\label{ex:c17q17}
计算 $\displaystyle\int_{-1}^1\bigl(x^{5}+x^{6}\bigr)dx$，并比较这两个幂函数各自的贡献。
\par\textit{改编训练。}\par\hyperref[ans:c17q17]{解答见第~\pageref{ans:c17q17} 页。}
\end{example}
\begin{example}[对称区间上的奇偶分解]\label{ex:c17q18}
计算 $\displaystyle\int_{-1}^1\bigl(x^{7}+x^{8}\bigr)dx$，并比较这两个幂函数各自的贡献。
\par\textit{改编训练。}\par\hyperref[ans:c17q18]{解答见第~\pageref{ans:c17q18} 页。}
\end{example}
\begin{example}[对称区间上的奇偶分解]\label{ex:c17q19}
计算 $\displaystyle\int_{-1}^1\bigl(x^{9}+x^{10}\bigr)dx$，并比较这两个幂函数各自的贡献。
\par\textit{改编训练。}\par\hyperref[ans:c17q19]{解答见第~\pageref{ans:c17q19} 页。}
\end{example}
\begin{example}[对称区间上的奇偶分解]\label{ex:c17q20}
计算 $\displaystyle\int_{-1}^1\bigl(x^{11}+x^{12}\bigr)dx$，并比较这两个幂函数各自的贡献。
\par\textit{改编训练。}\par\hyperref[ans:c17q20]{解答见第~\pageref{ans:c17q20} 页。}
\end{example}
\begin{example}[对称代换消去复杂分母]\label{ex:c17q21}
求 $I=\displaystyle\int_0^{\pi/2}\frac{\sin^{1}x}{\sin^{1}x+\cos^{1}x}dx$。
\par\textit{改编训练。}\par\hyperref[ans:c17q21]{解答见第~\pageref{ans:c17q21} 页。}
\end{example}
\begin{example}[对称代换消去复杂分母]\label{ex:c17q22}
求 $I=\displaystyle\int_0^{\pi/2}\frac{\sin^{2}x}{\sin^{2}x+\cos^{2}x}dx$。
\par\textit{改编训练。}\par\hyperref[ans:c17q22]{解答见第~\pageref{ans:c17q22} 页。}
\end{example}
\begin{example}[对称代换消去复杂分母]\label{ex:c17q23}
求 $I=\displaystyle\int_0^{\pi/2}\frac{\sin^{3}x}{\sin^{3}x+\cos^{3}x}dx$。
\par\textit{改编训练。}\par\hyperref[ans:c17q23]{解答见第~\pageref{ans:c17q23} 页。}
\end{example}
\begin{example}[对称代换消去复杂分母]\label{ex:c17q24}
求 $I=\displaystyle\int_0^{\pi/2}\frac{\sin^{4}x}{\sin^{4}x+\cos^{4}x}dx$。
\par\textit{改编训练。}\par\hyperref[ans:c17q24]{解答见第~\pageref{ans:c17q24} 页。}
\end{example}
\begin{example}[对称代换消去复杂分母]\label{ex:c17q25}
求 $I=\displaystyle\int_0^{\pi/2}\frac{\sin^{5}x}{\sin^{5}x+\cos^{5}x}dx$。
\par\textit{改编训练。}\par\hyperref[ans:c17q25]{解答见第~\pageref{ans:c17q25} 页。}
\end{example}
\begin{example}[对称代换消去复杂分母]\label{ex:c17q26}
求 $I=\displaystyle\int_0^{\pi/2}\frac{\sin^{6}x}{\sin^{6}x+\cos^{6}x}dx$。
\par\textit{改编训练。}\par\hyperref[ans:c17q26]{解答见第~\pageref{ans:c17q26} 页。}
\end{example}
\subsection{解答}
\begin{solve}[习题~\ref{ex:c17q1}：由原函数信息计算定积分]\phantomsection\label{ans:c17q1}
微积分基本定理给出 $\int_1^{\sqrt5}f=F(\sqrt5)-F(1)$，所以结果为 $(e^2+1)-(e^{-2}+1)=e^2-e^{-2}$。这里不必先把 $f$ 显式求出来。原函数中增加的常数在作端点差时抵消。
\end{solve}

\begin{solve}[习题~\ref{ex:c17q2}：纠正换元后的上下限]\phantomsection\label{ans:c17q2}
被积函数与微分替换正确，但新变量端点没有替换。$t=1$ 对应 $u=0$，$t=\pi$ 对应 $u=\log\pi$，所以积分为 $\int_0^{\log\pi}\cos u\,du=\sin(\log\pi)$。也可以先求原函数 $\sin(\log t)$，再代入原变量的两个端点，得到相同结果。
\end{solve}

\begin{solve}[习题~\ref{ex:c17q3}：黎曼和的识别]\phantomsection\label{ans:c17q3}
函数 $f(x)=x^{1}$ 在 $[0,1]$ 上连续，因而黎曼可积。将区间等分为 $n$ 份，每份宽 $1/n$，右端点为 $k/n$。所给表达式正是右端点黎曼和，因此
\[
\lim_{n\to\infty}\frac1n\sum_{k=1}^n\left(\frac kn\right)^{1}=\int_0^1x^{1}dx=\frac{1}{2}.
\]
也可利用单调性把该和夹在左端点和与右端点和之间。两者之差为 $1/n$，进一步说明采用左右端点并不改变极限。外面的 $1/n$ 表示区间宽度，不能遗漏。
\end{solve}

\begin{solve}[习题~\ref{ex:c17q4}：黎曼和的识别]\phantomsection\label{ans:c17q4}
函数 $f(x)=x^{2}$ 在 $[0,1]$ 上连续，因而黎曼可积。将区间等分为 $n$ 份，每份宽 $1/n$，右端点为 $k/n$。所给表达式正是右端点黎曼和，因此
\[
\lim_{n\to\infty}\frac1n\sum_{k=1}^n\left(\frac kn\right)^{2}=\int_0^1x^{2}dx=\frac{1}{3}.
\]
也可利用单调性把该和夹在左端点和与右端点和之间。两者之差为 $1/n$，进一步说明采用左右端点并不改变极限。外面的 $1/n$ 表示区间宽度，不能遗漏。
\end{solve}

\begin{solve}[习题~\ref{ex:c17q5}：黎曼和的识别]\phantomsection\label{ans:c17q5}
函数 $f(x)=x^{3}$ 在 $[0,1]$ 上连续，因而黎曼可积。将区间等分为 $n$ 份，每份宽 $1/n$，右端点为 $k/n$。所给表达式正是右端点黎曼和，因此
\[
\lim_{n\to\infty}\frac1n\sum_{k=1}^n\left(\frac kn\right)^{3}=\int_0^1x^{3}dx=\frac{1}{4}.
\]
也可利用单调性把该和夹在左端点和与右端点和之间。两者之差为 $1/n$，进一步说明采用左右端点并不改变极限。外面的 $1/n$ 表示区间宽度，不能遗漏。
\end{solve}

\begin{solve}[习题~\ref{ex:c17q6}：黎曼和的识别]\phantomsection\label{ans:c17q6}
函数 $f(x)=x^{4}$ 在 $[0,1]$ 上连续，因而黎曼可积。将区间等分为 $n$ 份，每份宽 $1/n$，右端点为 $k/n$。所给表达式正是右端点黎曼和，因此
\[
\lim_{n\to\infty}\frac1n\sum_{k=1}^n\left(\frac kn\right)^{4}=\int_0^1x^{4}dx=\frac{1}{5}.
\]
也可利用单调性把该和夹在左端点和与右端点和之间。两者之差为 $1/n$，进一步说明采用左右端点并不改变极限。外面的 $1/n$ 表示区间宽度，不能遗漏。
\end{solve}

\begin{solve}[习题~\ref{ex:c17q7}：黎曼和的识别]\phantomsection\label{ans:c17q7}
函数 $f(x)=x^{5}$ 在 $[0,1]$ 上连续，因而黎曼可积。将区间等分为 $n$ 份，每份宽 $1/n$，右端点为 $k/n$。所给表达式正是右端点黎曼和，因此
\[
\lim_{n\to\infty}\frac1n\sum_{k=1}^n\left(\frac kn\right)^{5}=\int_0^1x^{5}dx=\frac{1}{6}.
\]
也可利用单调性把该和夹在左端点和与右端点和之间。两者之差为 $1/n$，进一步说明采用左右端点并不改变极限。外面的 $1/n$ 表示区间宽度，不能遗漏。
\end{solve}

\begin{solve}[习题~\ref{ex:c17q8}：黎曼和的识别]\phantomsection\label{ans:c17q8}
函数 $f(x)=x^{6}$ 在 $[0,1]$ 上连续，因而黎曼可积。将区间等分为 $n$ 份，每份宽 $1/n$，右端点为 $k/n$。所给表达式正是右端点黎曼和，因此
\[
\lim_{n\to\infty}\frac1n\sum_{k=1}^n\left(\frac kn\right)^{6}=\int_0^1x^{6}dx=\frac{1}{7}.
\]
也可利用单调性把该和夹在左端点和与右端点和之间。两者之差为 $1/n$，进一步说明采用左右端点并不改变极限。外面的 $1/n$ 表示区间宽度，不能遗漏。
\end{solve}

\begin{solve}[习题~\ref{ex:c17q9}：两个积分端点都在变化]\phantomsection\label{ans:c17q9}
取 $G(u)=\int_0^u e^{t^2}dt$，则由连续性可知 $G^{\prime}(u)=e^{u^2}$，且 $F(x)=G(bx)-G(ax)$。链式法则给出
\[
F^{\prime}(x)=2e^{4x^2}-1e^{1x^2},\qquad F^{\prime}(0)=1.
\]
积分函数本身就提供了一个可微原函数，不要求它能写成初等函数的有限组合。下限项带负号，上、下限各自的导数都必须乘上。
\end{solve}

\begin{solve}[习题~\ref{ex:c17q10}：两个积分端点都在变化]\phantomsection\label{ans:c17q10}
取 $G(u)=\int_0^u e^{t^2}dt$，则由连续性可知 $G^{\prime}(u)=e^{u^2}$，且 $F(x)=G(bx)-G(ax)$。链式法则给出
\[
F^{\prime}(x)=3e^{9x^2}-2e^{4x^2},\qquad F^{\prime}(0)=1.
\]
积分函数本身就提供了一个可微原函数，不要求它能写成初等函数的有限组合。下限项带负号，上、下限各自的导数都必须乘上。
\end{solve}

\begin{solve}[习题~\ref{ex:c17q11}：两个积分端点都在变化]\phantomsection\label{ans:c17q11}
取 $G(u)=\int_0^u e^{t^2}dt$，则由连续性可知 $G^{\prime}(u)=e^{u^2}$，且 $F(x)=G(bx)-G(ax)$。链式法则给出
\[
F^{\prime}(x)=4e^{16x^2}-3e^{9x^2},\qquad F^{\prime}(0)=1.
\]
积分函数本身就提供了一个可微原函数，不要求它能写成初等函数的有限组合。下限项带负号，上、下限各自的导数都必须乘上。
\end{solve}

\begin{solve}[习题~\ref{ex:c17q12}：两个积分端点都在变化]\phantomsection\label{ans:c17q12}
取 $G(u)=\int_0^u e^{t^2}dt$，则由连续性可知 $G^{\prime}(u)=e^{u^2}$，且 $F(x)=G(bx)-G(ax)$。链式法则给出
\[
F^{\prime}(x)=4e^{16x^2}-1e^{1x^2},\qquad F^{\prime}(0)=3.
\]
积分函数本身就提供了一个可微原函数，不要求它能写成初等函数的有限组合。下限项带负号，上、下限各自的导数都必须乘上。
\end{solve}

\begin{solve}[习题~\ref{ex:c17q13}：两个积分端点都在变化]\phantomsection\label{ans:c17q13}
取 $G(u)=\int_0^u e^{t^2}dt$，则由连续性可知 $G^{\prime}(u)=e^{u^2}$，且 $F(x)=G(bx)-G(ax)$。链式法则给出
\[
F^{\prime}(x)=5e^{25x^2}-2e^{4x^2},\qquad F^{\prime}(0)=3.
\]
积分函数本身就提供了一个可微原函数，不要求它能写成初等函数的有限组合。下限项带负号，上、下限各自的导数都必须乘上。
\end{solve}

\begin{solve}[习题~\ref{ex:c17q14}：两个积分端点都在变化]\phantomsection\label{ans:c17q14}
取 $G(u)=\int_0^u e^{t^2}dt$，则由连续性可知 $G^{\prime}(u)=e^{u^2}$，且 $F(x)=G(bx)-G(ax)$。链式法则给出
\[
F^{\prime}(x)=5e^{25x^2}-3e^{9x^2},\qquad F^{\prime}(0)=2.
\]
积分函数本身就提供了一个可微原函数，不要求它能写成初等函数的有限组合。下限项带负号，上、下限各自的导数都必须乘上。
\end{solve}

\begin{solve}[习题~\ref{ex:c17q15}：对称区间上的奇偶分解]\phantomsection\label{ans:c17q15}
函数 $x^{1}$ 为奇函数，其在对称区间上的积分为零；$x^{2}$ 为偶函数，因此
\[
\int_{-1}^1\bigl(x^{1}+x^{2}\bigr)dx=2\int_0^1x^{2}dx=\frac{2}{3}.
\]
奇函数积分为零是正负贡献的抵消，不意味着曲线与横轴围成的几何面积为零。若题目要求面积，必须先取绝对值，或者按正负区间分段积分。此处两个幂函数连续，所有定积分都正常存在。
\end{solve}

\begin{solve}[习题~\ref{ex:c17q16}：对称区间上的奇偶分解]\phantomsection\label{ans:c17q16}
函数 $x^{3}$ 为奇函数，其在对称区间上的积分为零；$x^{4}$ 为偶函数，因此
\[
\int_{-1}^1\bigl(x^{3}+x^{4}\bigr)dx=2\int_0^1x^{4}dx=\frac{2}{5}.
\]
奇函数积分为零是正负贡献的抵消，不意味着曲线与横轴围成的几何面积为零。若题目要求面积，必须先取绝对值，或者按正负区间分段积分。此处两个幂函数连续，所有定积分都正常存在。
\end{solve}

\begin{solve}[习题~\ref{ex:c17q17}：对称区间上的奇偶分解]\phantomsection\label{ans:c17q17}
函数 $x^{5}$ 为奇函数，其在对称区间上的积分为零；$x^{6}$ 为偶函数，因此
\[
\int_{-1}^1\bigl(x^{5}+x^{6}\bigr)dx=2\int_0^1x^{6}dx=\frac{2}{7}.
\]
奇函数积分为零是正负贡献的抵消，不意味着曲线与横轴围成的几何面积为零。若题目要求面积，必须先取绝对值，或者按正负区间分段积分。此处两个幂函数连续，所有定积分都正常存在。
\end{solve}

\begin{solve}[习题~\ref{ex:c17q18}：对称区间上的奇偶分解]\phantomsection\label{ans:c17q18}
函数 $x^{7}$ 为奇函数，其在对称区间上的积分为零；$x^{8}$ 为偶函数，因此
\[
\int_{-1}^1\bigl(x^{7}+x^{8}\bigr)dx=2\int_0^1x^{8}dx=\frac{2}{9}.
\]
奇函数积分为零是正负贡献的抵消，不意味着曲线与横轴围成的几何面积为零。若题目要求面积，必须先取绝对值，或者按正负区间分段积分。此处两个幂函数连续，所有定积分都正常存在。
\end{solve}

\begin{solve}[习题~\ref{ex:c17q19}：对称区间上的奇偶分解]\phantomsection\label{ans:c17q19}
函数 $x^{9}$ 为奇函数，其在对称区间上的积分为零；$x^{10}$ 为偶函数，因此
\[
\int_{-1}^1\bigl(x^{9}+x^{10}\bigr)dx=2\int_0^1x^{10}dx=\frac{2}{11}.
\]
奇函数积分为零是正负贡献的抵消，不意味着曲线与横轴围成的几何面积为零。若题目要求面积，必须先取绝对值，或者按正负区间分段积分。此处两个幂函数连续，所有定积分都正常存在。
\end{solve}

\begin{solve}[习题~\ref{ex:c17q20}：对称区间上的奇偶分解]\phantomsection\label{ans:c17q20}
函数 $x^{11}$ 为奇函数，其在对称区间上的积分为零；$x^{12}$ 为偶函数，因此
\[
\int_{-1}^1\bigl(x^{11}+x^{12}\bigr)dx=2\int_0^1x^{12}dx=\frac{2}{13}.
\]
奇函数积分为零是正负贡献的抵消，不意味着曲线与横轴围成的几何面积为零。若题目要求面积，必须先取绝对值，或者按正负区间分段积分。此处两个幂函数连续，所有定积分都正常存在。
\end{solve}

\begin{solve}[习题~\ref{ex:c17q21}：对称代换消去复杂分母]\phantomsection\label{ans:c17q21}
在给定区间内分母为正，被积函数连续。作代换 $x=\pi/2-t$，正弦与余弦交换，得
\[
I=\int_0^{\pi/2}\frac{\cos^{1}t}{\sin^{1}t+\cos^{1}t}dt.
\]
两式相加后分子恰等于分母，
\[
2I=\int_0^{\pi/2}1\,dx=\frac\pi2,\qquad I=\frac\pi4.
\]
积分变量是哑变量，所以可将第二式的 $t$ 重新写成 $x$ 后相加。这种对称法计算的是定积分，通常无需也不容易先求一个初等原函数。
\end{solve}

\begin{solve}[习题~\ref{ex:c17q22}：对称代换消去复杂分母]\phantomsection\label{ans:c17q22}
在给定区间内分母为正，被积函数连续。作代换 $x=\pi/2-t$，正弦与余弦交换，得
\[
I=\int_0^{\pi/2}\frac{\cos^{2}t}{\sin^{2}t+\cos^{2}t}dt.
\]
两式相加后分子恰等于分母，
\[
2I=\int_0^{\pi/2}1\,dx=\frac\pi2,\qquad I=\frac\pi4.
\]
积分变量是哑变量，所以可将第二式的 $t$ 重新写成 $x$ 后相加。这种对称法计算的是定积分，通常无需也不容易先求一个初等原函数。
\end{solve}

\begin{solve}[习题~\ref{ex:c17q23}：对称代换消去复杂分母]\phantomsection\label{ans:c17q23}
在给定区间内分母为正，被积函数连续。作代换 $x=\pi/2-t$，正弦与余弦交换，得
\[
I=\int_0^{\pi/2}\frac{\cos^{3}t}{\sin^{3}t+\cos^{3}t}dt.
\]
两式相加后分子恰等于分母，
\[
2I=\int_0^{\pi/2}1\,dx=\frac\pi2,\qquad I=\frac\pi4.
\]
积分变量是哑变量，所以可将第二式的 $t$ 重新写成 $x$ 后相加。这种对称法计算的是定积分，通常无需也不容易先求一个初等原函数。
\end{solve}

\begin{solve}[习题~\ref{ex:c17q24}：对称代换消去复杂分母]\phantomsection\label{ans:c17q24}
在给定区间内分母为正，被积函数连续。作代换 $x=\pi/2-t$，正弦与余弦交换，得
\[
I=\int_0^{\pi/2}\frac{\cos^{4}t}{\sin^{4}t+\cos^{4}t}dt.
\]
两式相加后分子恰等于分母，
\[
2I=\int_0^{\pi/2}1\,dx=\frac\pi2,\qquad I=\frac\pi4.
\]
积分变量是哑变量，所以可将第二式的 $t$ 重新写成 $x$ 后相加。这种对称法计算的是定积分，通常无需也不容易先求一个初等原函数。
\end{solve}

\begin{solve}[习题~\ref{ex:c17q25}：对称代换消去复杂分母]\phantomsection\label{ans:c17q25}
在给定区间内分母为正，被积函数连续。作代换 $x=\pi/2-t$，正弦与余弦交换，得
\[
I=\int_0^{\pi/2}\frac{\cos^{5}t}{\sin^{5}t+\cos^{5}t}dt.
\]
两式相加后分子恰等于分母，
\[
2I=\int_0^{\pi/2}1\,dx=\frac\pi2,\qquad I=\frac\pi4.
\]
积分变量是哑变量，所以可将第二式的 $t$ 重新写成 $x$ 后相加。这种对称法计算的是定积分，通常无需也不容易先求一个初等原函数。
\end{solve}

\begin{solve}[习题~\ref{ex:c17q26}：对称代换消去复杂分母]\phantomsection\label{ans:c17q26}
在给定区间内分母为正，被积函数连续。作代换 $x=\pi/2-t$，正弦与余弦交换，得
\[
I=\int_0^{\pi/2}\frac{\cos^{6}t}{\sin^{6}t+\cos^{6}t}dt.
\]
两式相加后分子恰等于分母，
\[
2I=\int_0^{\pi/2}1\,dx=\frac\pi2,\qquad I=\frac\pi4.
\]
积分变量是哑变量，所以可将第二式的 $t$ 重新写成 $x$ 后相加。这种对称法计算的是定积分，通常无需也不容易先求一个初等原函数。
\end{solve}

\clearpage
\section{数值积分与误差估计}\label{sec:ch18}
数值近似应当同时报告方法和误差。复合梯形公式、复合中点公式与 Simpson 公式的误差阶不同，所需导数条件也不同；误差估计给出保证，而非实际误差的必然大小。

\textbf{资源对应：}CLP-2 \S~1.11；误差推导为补充训练。本章共 24 道编号习题。
\subsection{习题}
\begin{example}[二次函数的梯形与中点公式]\label{ex:c18q1}
用 $2$ 等分的复合梯形公式与复合中点公式近似 $\int_0^1x^2dx$，求精确误差。
\par\textit{改编训练。}\par\hyperref[ans:c18q1]{解答见第~\pageref{ans:c18q1} 页。}
\end{example}
\begin{example}[二次函数的梯形与中点公式]\label{ex:c18q2}
用 $4$ 等分的复合梯形公式与复合中点公式近似 $\int_0^1x^2dx$，求精确误差。
\par\textit{改编训练。}\par\hyperref[ans:c18q2]{解答见第~\pageref{ans:c18q2} 页。}
\end{example}
\begin{example}[二次函数的梯形与中点公式]\label{ex:c18q3}
用 $6$ 等分的复合梯形公式与复合中点公式近似 $\int_0^1x^2dx$，求精确误差。
\par\textit{改编训练。}\par\hyperref[ans:c18q3]{解答见第~\pageref{ans:c18q3} 页。}
\end{example}
\begin{example}[二次函数的梯形与中点公式]\label{ex:c18q4}
用 $8$ 等分的复合梯形公式与复合中点公式近似 $\int_0^1x^2dx$，求精确误差。
\par\textit{改编训练。}\par\hyperref[ans:c18q4]{解答见第~\pageref{ans:c18q4} 页。}
\end{example}
\begin{example}[二次函数的梯形与中点公式]\label{ex:c18q5}
用 $10$ 等分的复合梯形公式与复合中点公式近似 $\int_0^1x^2dx$，求精确误差。
\par\textit{改编训练。}\par\hyperref[ans:c18q5]{解答见第~\pageref{ans:c18q5} 页。}
\end{example}
\begin{example}[二次函数的梯形与中点公式]\label{ex:c18q6}
用 $12$ 等分的复合梯形公式与复合中点公式近似 $\int_0^1x^2dx$，求精确误差。
\par\textit{改编训练。}\par\hyperref[ans:c18q6]{解答见第~\pageref{ans:c18q6} 页。}
\end{example}
\begin{example}[Simpson 公式的实际误差]\label{ex:c18q7}
用 $2$ 等分的复合 Simpson 公式近似 $\int_0^1x^4dx$，给出近似值与误差。
\par\textit{改编训练。}\par\hyperref[ans:c18q7]{解答见第~\pageref{ans:c18q7} 页。}
\end{example}
\begin{example}[Simpson 公式的实际误差]\label{ex:c18q8}
用 $4$ 等分的复合 Simpson 公式近似 $\int_0^1x^4dx$，给出近似值与误差。
\par\textit{改编训练。}\par\hyperref[ans:c18q8]{解答见第~\pageref{ans:c18q8} 页。}
\end{example}
\begin{example}[Simpson 公式的实际误差]\label{ex:c18q9}
用 $6$ 等分的复合 Simpson 公式近似 $\int_0^1x^4dx$，给出近似值与误差。
\par\textit{改编训练。}\par\hyperref[ans:c18q9]{解答见第~\pageref{ans:c18q9} 页。}
\end{example}
\begin{example}[Simpson 公式的实际误差]\label{ex:c18q10}
用 $8$ 等分的复合 Simpson 公式近似 $\int_0^1x^4dx$，给出近似值与误差。
\par\textit{改编训练。}\par\hyperref[ans:c18q10]{解答见第~\pageref{ans:c18q10} 页。}
\end{example}
\begin{example}[Simpson 公式的实际误差]\label{ex:c18q11}
用 $10$ 等分的复合 Simpson 公式近似 $\int_0^1x^4dx$，给出近似值与误差。
\par\textit{改编训练。}\par\hyperref[ans:c18q11]{解答见第~\pageref{ans:c18q11} 页。}
\end{example}
\begin{example}[Simpson 公式的实际误差]\label{ex:c18q12}
用 $12$ 等分的复合 Simpson 公式近似 $\int_0^1x^4dx$，给出近似值与误差。
\par\textit{改编训练。}\par\hyperref[ans:c18q12]{解答见第~\pageref{ans:c18q12} 页。}
\end{example}
\begin{example}[用导数界选择网格]\label{ex:c18q13}
对 $I=\int_0^1e^x dx$，求 $5$ 等分复合梯形法和复合中点法的误差上界，并比较所需精度。
\par\textit{改编训练。}\par\hyperref[ans:c18q13]{解答见第~\pageref{ans:c18q13} 页。}
\end{example}
\begin{example}[用导数界选择网格]\label{ex:c18q14}
对 $I=\int_0^1e^x dx$，求 $10$ 等分复合梯形法和复合中点法的误差上界，并比较所需精度。
\par\textit{改编训练。}\par\hyperref[ans:c18q14]{解答见第~\pageref{ans:c18q14} 页。}
\end{example}
\begin{example}[用导数界选择网格]\label{ex:c18q15}
对 $I=\int_0^1e^x dx$，求 $15$ 等分复合梯形法和复合中点法的误差上界，并比较所需精度。
\par\textit{改编训练。}\par\hyperref[ans:c18q15]{解答见第~\pageref{ans:c18q15} 页。}
\end{example}
\begin{example}[用导数界选择网格]\label{ex:c18q16}
对 $I=\int_0^1e^x dx$，求 $20$ 等分复合梯形法和复合中点法的误差上界，并比较所需精度。
\par\textit{改编训练。}\par\hyperref[ans:c18q16]{解答见第~\pageref{ans:c18q16} 页。}
\end{example}
\begin{example}[用导数界选择网格]\label{ex:c18q17}
对 $I=\int_0^1e^x dx$，求 $25$ 等分复合梯形法和复合中点法的误差上界，并比较所需精度。
\par\textit{改编训练。}\par\hyperref[ans:c18q17]{解答见第~\pageref{ans:c18q17} 页。}
\end{example}
\begin{example}[用导数界选择网格]\label{ex:c18q18}
对 $I=\int_0^1e^x dx$，求 $30$ 等分复合梯形法和复合中点法的误差上界，并比较所需精度。
\par\textit{改编训练。}\par\hyperref[ans:c18q18]{解答见第~\pageref{ans:c18q18} 页。}
\end{example}
\begin{example}[中心差分的截断误差]\label{ex:c18q19}
用 $D_h=\dfrac{\sin(1h)-\sin(-1h)}{2h}$ 近似 $f(x)=\sin(1x)$ 在 $0$ 处的导数。求误差的主项。
\par\textit{改编训练。}\par\hyperref[ans:c18q19]{解答见第~\pageref{ans:c18q19} 页。}
\end{example}
\begin{example}[中心差分的截断误差]\label{ex:c18q20}
用 $D_h=\dfrac{\sin(2h)-\sin(-2h)}{2h}$ 近似 $f(x)=\sin(2x)$ 在 $0$ 处的导数。求误差的主项。
\par\textit{改编训练。}\par\hyperref[ans:c18q20]{解答见第~\pageref{ans:c18q20} 页。}
\end{example}
\begin{example}[中心差分的截断误差]\label{ex:c18q21}
用 $D_h=\dfrac{\sin(3h)-\sin(-3h)}{2h}$ 近似 $f(x)=\sin(3x)$ 在 $0$ 处的导数。求误差的主项。
\par\textit{改编训练。}\par\hyperref[ans:c18q21]{解答见第~\pageref{ans:c18q21} 页。}
\end{example}
\begin{example}[中心差分的截断误差]\label{ex:c18q22}
用 $D_h=\dfrac{\sin(4h)-\sin(-4h)}{2h}$ 近似 $f(x)=\sin(4x)$ 在 $0$ 处的导数。求误差的主项。
\par\textit{改编训练。}\par\hyperref[ans:c18q22]{解答见第~\pageref{ans:c18q22} 页。}
\end{example}
\begin{example}[中心差分的截断误差]\label{ex:c18q23}
用 $D_h=\dfrac{\sin(5h)-\sin(-5h)}{2h}$ 近似 $f(x)=\sin(5x)$ 在 $0$ 处的导数。求误差的主项。
\par\textit{改编训练。}\par\hyperref[ans:c18q23]{解答见第~\pageref{ans:c18q23} 页。}
\end{example}
\begin{example}[中心差分的截断误差]\label{ex:c18q24}
用 $D_h=\dfrac{\sin(6h)-\sin(-6h)}{2h}$ 近似 $f(x)=\sin(6x)$ 在 $0$ 处的导数。求误差的主项。
\par\textit{改编训练。}\par\hyperref[ans:c18q24]{解答见第~\pageref{ans:c18q24} 页。}
\end{example}
\subsection{解答}
\begin{solve}[习题~\ref{ex:c18q1}：二次函数的梯形与中点公式]\phantomsection\label{ans:c18q1}
令 $h=1/2$。梯形公式为 $T=h(f(0)/2+\sum_{k=1}^{1}f(kh)+f(1)/2)$，中点公式为 $M=h\sum_{k=1}^{2}f((k-1/2)h)$。利用平方和公式整理可得
\[
T=\frac13+\frac1{24}=\frac{3}{8},\qquad M=\frac13-\frac1{48}=\frac{5}{16}.
\]
精确积分为 $1/3$，因此梯形高估 $\frac{1}{24}$，中点低估 $\frac{1}{48}$。这与 $x^2$ 的凸性一致：弦在图像上方，而对称小区间的中点高度偏低。
\end{solve}

\begin{solve}[习题~\ref{ex:c18q2}：二次函数的梯形与中点公式]\phantomsection\label{ans:c18q2}
令 $h=1/4$。梯形公式为 $T=h(f(0)/2+\sum_{k=1}^{3}f(kh)+f(1)/2)$，中点公式为 $M=h\sum_{k=1}^{4}f((k-1/2)h)$。利用平方和公式整理可得
\[
T=\frac13+\frac1{96}=\frac{11}{32},\qquad M=\frac13-\frac1{192}=\frac{21}{64}.
\]
精确积分为 $1/3$，因此梯形高估 $\frac{1}{96}$，中点低估 $\frac{1}{192}$。这与 $x^2$ 的凸性一致：弦在图像上方，而对称小区间的中点高度偏低。
\end{solve}

\begin{solve}[习题~\ref{ex:c18q3}：二次函数的梯形与中点公式]\phantomsection\label{ans:c18q3}
令 $h=1/6$。梯形公式为 $T=h(f(0)/2+\sum_{k=1}^{5}f(kh)+f(1)/2)$，中点公式为 $M=h\sum_{k=1}^{6}f((k-1/2)h)$。利用平方和公式整理可得
\[
T=\frac13+\frac1{216}=\frac{73}{216},\qquad M=\frac13-\frac1{432}=\frac{143}{432}.
\]
精确积分为 $1/3$，因此梯形高估 $\frac{1}{216}$，中点低估 $\frac{1}{432}$。这与 $x^2$ 的凸性一致：弦在图像上方，而对称小区间的中点高度偏低。
\end{solve}

\begin{solve}[习题~\ref{ex:c18q4}：二次函数的梯形与中点公式]\phantomsection\label{ans:c18q4}
令 $h=1/8$。梯形公式为 $T=h(f(0)/2+\sum_{k=1}^{7}f(kh)+f(1)/2)$，中点公式为 $M=h\sum_{k=1}^{8}f((k-1/2)h)$。利用平方和公式整理可得
\[
T=\frac13+\frac1{384}=\frac{43}{128},\qquad M=\frac13-\frac1{768}=\frac{85}{256}.
\]
精确积分为 $1/3$，因此梯形高估 $\frac{1}{384}$，中点低估 $\frac{1}{768}$。这与 $x^2$ 的凸性一致：弦在图像上方，而对称小区间的中点高度偏低。
\end{solve}

\begin{solve}[习题~\ref{ex:c18q5}：二次函数的梯形与中点公式]\phantomsection\label{ans:c18q5}
令 $h=1/10$。梯形公式为 $T=h(f(0)/2+\sum_{k=1}^{9}f(kh)+f(1)/2)$，中点公式为 $M=h\sum_{k=1}^{10}f((k-1/2)h)$。利用平方和公式整理可得
\[
T=\frac13+\frac1{600}=\frac{67}{200},\qquad M=\frac13-\frac1{1200}=\frac{133}{400}.
\]
精确积分为 $1/3$，因此梯形高估 $\frac{1}{600}$，中点低估 $\frac{1}{1200}$。这与 $x^2$ 的凸性一致：弦在图像上方，而对称小区间的中点高度偏低。
\end{solve}

\begin{solve}[习题~\ref{ex:c18q6}：二次函数的梯形与中点公式]\phantomsection\label{ans:c18q6}
令 $h=1/12$。梯形公式为 $T=h(f(0)/2+\sum_{k=1}^{11}f(kh)+f(1)/2)$，中点公式为 $M=h\sum_{k=1}^{12}f((k-1/2)h)$。利用平方和公式整理可得
\[
T=\frac13+\frac1{864}=\frac{289}{864},\qquad M=\frac13-\frac1{1728}=\frac{575}{1728}.
\]
精确积分为 $1/3$，因此梯形高估 $\frac{1}{864}$，中点低估 $\frac{1}{1728}$。这与 $x^2$ 的凸性一致：弦在图像上方，而对称小区间的中点高度偏低。
\end{solve}

\begin{solve}[习题~\ref{ex:c18q7}：Simpson 公式的实际误差]\phantomsection\label{ans:c18q7}
这里记 $n=2$，它为偶数，满足 Simpson 公式的分组要求。用 $h=1/n$；下式两组求和都只取 $1\le k\le n-1$，代入
\[
S=\frac h3\left[f(0)+f(1)+4\sum_{k\text{为奇数}}f(kh)+2\sum_{\substack{k\text{为偶数}\\0<k<n}}f(kh)\right].
\]
利用四次幂求和公式，或者逐点精确求和，得到
\[
S=\frac15+\frac{1}{120}=\frac{5}{24}.
\]
所以高估量为 $\frac{1}{120}$。一般误差界为 $\max|f^{(4)}|/(180n^4)$，本题四阶导数为24，恰与这一实际误差一致。
\end{solve}

\begin{solve}[习题~\ref{ex:c18q8}：Simpson 公式的实际误差]\phantomsection\label{ans:c18q8}
这里记 $n=4$，它为偶数，满足 Simpson 公式的分组要求。用 $h=1/n$；下式两组求和都只取 $1\le k\le n-1$，代入
\[
S=\frac h3\left[f(0)+f(1)+4\sum_{k\text{为奇数}}f(kh)+2\sum_{\substack{k\text{为偶数}\\0<k<n}}f(kh)\right].
\]
利用四次幂求和公式，或者逐点精确求和，得到
\[
S=\frac15+\frac{1}{1920}=\frac{77}{384}.
\]
所以高估量为 $\frac{1}{1920}$。一般误差界为 $\max|f^{(4)}|/(180n^4)$，本题四阶导数为24，恰与这一实际误差一致。
\end{solve}

\begin{solve}[习题~\ref{ex:c18q9}：Simpson 公式的实际误差]\phantomsection\label{ans:c18q9}
这里记 $n=6$，它为偶数，满足 Simpson 公式的分组要求。用 $h=1/n$；下式两组求和都只取 $1\le k\le n-1$，代入
\[
S=\frac h3\left[f(0)+f(1)+4\sum_{k\text{为奇数}}f(kh)+2\sum_{\substack{k\text{为偶数}\\0<k<n}}f(kh)\right].
\]
利用四次幂求和公式，或者逐点精确求和，得到
\[
S=\frac15+\frac{1}{9720}=\frac{389}{1944}.
\]
所以高估量为 $\frac{1}{9720}$。一般误差界为 $\max|f^{(4)}|/(180n^4)$，本题四阶导数为24，恰与这一实际误差一致。
\end{solve}

\begin{solve}[习题~\ref{ex:c18q10}：Simpson 公式的实际误差]\phantomsection\label{ans:c18q10}
这里记 $n=8$，它为偶数，满足 Simpson 公式的分组要求。用 $h=1/n$；下式两组求和都只取 $1\le k\le n-1$，代入
\[
S=\frac h3\left[f(0)+f(1)+4\sum_{k\text{为奇数}}f(kh)+2\sum_{\substack{k\text{为偶数}\\0<k<n}}f(kh)\right].
\]
利用四次幂求和公式，或者逐点精确求和，得到
\[
S=\frac15+\frac{1}{30720}=\frac{1229}{6144}.
\]
所以高估量为 $\frac{1}{30720}$。一般误差界为 $\max|f^{(4)}|/(180n^4)$，本题四阶导数为24，恰与这一实际误差一致。
\end{solve}

\begin{solve}[习题~\ref{ex:c18q11}：Simpson 公式的实际误差]\phantomsection\label{ans:c18q11}
这里记 $n=10$，它为偶数，满足 Simpson 公式的分组要求。用 $h=1/n$；下式两组求和都只取 $1\le k\le n-1$，代入
\[
S=\frac h3\left[f(0)+f(1)+4\sum_{k\text{为奇数}}f(kh)+2\sum_{\substack{k\text{为偶数}\\0<k<n}}f(kh)\right].
\]
利用四次幂求和公式，或者逐点精确求和，得到
\[
S=\frac15+\frac{1}{75000}=\frac{15001}{75000}.
\]
所以高估量为 $\frac{1}{75000}$。一般误差界为 $\max|f^{(4)}|/(180n^4)$，本题四阶导数为24，恰与这一实际误差一致。
\end{solve}

\begin{solve}[习题~\ref{ex:c18q12}：Simpson 公式的实际误差]\phantomsection\label{ans:c18q12}
这里记 $n=12$，它为偶数，满足 Simpson 公式的分组要求。用 $h=1/n$；下式两组求和都只取 $1\le k\le n-1$，代入
\[
S=\frac h3\left[f(0)+f(1)+4\sum_{k\text{为奇数}}f(kh)+2\sum_{\substack{k\text{为偶数}\\0<k<n}}f(kh)\right].
\]
利用四次幂求和公式，或者逐点精确求和，得到
\[
S=\frac15+\frac{1}{155520}=\frac{6221}{31104}.
\]
所以高估量为 $\frac{1}{155520}$。一般误差界为 $\max|f^{(4)}|/(180n^4)$，本题四阶导数为24，恰与这一实际误差一致。
\end{solve}

\begin{solve}[习题~\ref{ex:c18q13}：用导数界选择网格]\phantomsection\label{ans:c18q13}
函数二阶导数为 $e^x$，在 $[0,1]$ 上的最大值为 $e$。复合公式的标准误差估计分别为
\[
|I-T_{5}|\le\frac e{300},\qquad |I-M_{5}|\le\frac e{600}.
\]
这些是绝对误差的保证，不是近似结果本身。若要求梯形法误差小于 $10^{-4}$，一个充分条件是 $n>\sqrt{e/(12\cdot10^{-4})}$；中点法把分母12改为24即可。使用上界会偏保守，但无需事先知道精确积分就能选择网格。
\end{solve}

\begin{solve}[习题~\ref{ex:c18q14}：用导数界选择网格]\phantomsection\label{ans:c18q14}
函数二阶导数为 $e^x$，在 $[0,1]$ 上的最大值为 $e$。复合公式的标准误差估计分别为
\[
|I-T_{10}|\le\frac e{1200},\qquad |I-M_{10}|\le\frac e{2400}.
\]
这些是绝对误差的保证，不是近似结果本身。若要求梯形法误差小于 $10^{-4}$，一个充分条件是 $n>\sqrt{e/(12\cdot10^{-4})}$；中点法把分母12改为24即可。使用上界会偏保守，但无需事先知道精确积分就能选择网格。
\end{solve}

\begin{solve}[习题~\ref{ex:c18q15}：用导数界选择网格]\phantomsection\label{ans:c18q15}
函数二阶导数为 $e^x$，在 $[0,1]$ 上的最大值为 $e$。复合公式的标准误差估计分别为
\[
|I-T_{15}|\le\frac e{2700},\qquad |I-M_{15}|\le\frac e{5400}.
\]
这些是绝对误差的保证，不是近似结果本身。若要求梯形法误差小于 $10^{-4}$，一个充分条件是 $n>\sqrt{e/(12\cdot10^{-4})}$；中点法把分母12改为24即可。使用上界会偏保守，但无需事先知道精确积分就能选择网格。
\end{solve}

\begin{solve}[习题~\ref{ex:c18q16}：用导数界选择网格]\phantomsection\label{ans:c18q16}
函数二阶导数为 $e^x$，在 $[0,1]$ 上的最大值为 $e$。复合公式的标准误差估计分别为
\[
|I-T_{20}|\le\frac e{4800},\qquad |I-M_{20}|\le\frac e{9600}.
\]
这些是绝对误差的保证，不是近似结果本身。若要求梯形法误差小于 $10^{-4}$，一个充分条件是 $n>\sqrt{e/(12\cdot10^{-4})}$；中点法把分母12改为24即可。使用上界会偏保守，但无需事先知道精确积分就能选择网格。
\end{solve}

\begin{solve}[习题~\ref{ex:c18q17}：用导数界选择网格]\phantomsection\label{ans:c18q17}
函数二阶导数为 $e^x$，在 $[0,1]$ 上的最大值为 $e$。复合公式的标准误差估计分别为
\[
|I-T_{25}|\le\frac e{7500},\qquad |I-M_{25}|\le\frac e{15000}.
\]
这些是绝对误差的保证，不是近似结果本身。若要求梯形法误差小于 $10^{-4}$，一个充分条件是 $n>\sqrt{e/(12\cdot10^{-4})}$；中点法把分母12改为24即可。使用上界会偏保守，但无需事先知道精确积分就能选择网格。
\end{solve}

\begin{solve}[习题~\ref{ex:c18q18}：用导数界选择网格]\phantomsection\label{ans:c18q18}
函数二阶导数为 $e^x$，在 $[0,1]$ 上的最大值为 $e$。复合公式的标准误差估计分别为
\[
|I-T_{30}|\le\frac e{10800},\qquad |I-M_{30}|\le\frac e{21600}.
\]
这些是绝对误差的保证，不是近似结果本身。若要求梯形法误差小于 $10^{-4}$，一个充分条件是 $n>\sqrt{e/(12\cdot10^{-4})}$；中点法把分母12改为24即可。使用上界会偏保守，但无需事先知道精确积分就能选择网格。
\end{solve}

\begin{solve}[习题~\ref{ex:c18q19}：中心差分的截断误差]\phantomsection\label{ans:c18q19}
由奇性 $D_h=\sin(1h)/h$。对正弦展开至三阶，得到
\[
D_h=1-\frac{1}{6}h^2+O(h^4).
\]
真导数为 $1$，所以 $D_h-f^{\prime}(0)=-\frac{1}{6}h^2+O(h^4)$。中心差分误差为二阶，这源于对称相减消去了泰勒展开的偶次项。该式仅描述精确实数运算中的截断误差；计算机浮点运算还包含舍入误差，因而并非 $h$ 越小总是越好。
\end{solve}

\begin{solve}[习题~\ref{ex:c18q20}：中心差分的截断误差]\phantomsection\label{ans:c18q20}
由奇性 $D_h=\sin(2h)/h$。对正弦展开至三阶，得到
\[
D_h=2-\frac{4}{3}h^2+O(h^4).
\]
真导数为 $2$，所以 $D_h-f^{\prime}(0)=-\frac{4}{3}h^2+O(h^4)$。中心差分误差为二阶，这源于对称相减消去了泰勒展开的偶次项。该式仅描述精确实数运算中的截断误差；计算机浮点运算还包含舍入误差，因而并非 $h$ 越小总是越好。
\end{solve}

\begin{solve}[习题~\ref{ex:c18q21}：中心差分的截断误差]\phantomsection\label{ans:c18q21}
由奇性 $D_h=\sin(3h)/h$。对正弦展开至三阶，得到
\[
D_h=3-\frac{9}{2}h^2+O(h^4).
\]
真导数为 $3$，所以 $D_h-f^{\prime}(0)=-\frac{9}{2}h^2+O(h^4)$。中心差分误差为二阶，这源于对称相减消去了泰勒展开的偶次项。该式仅描述精确实数运算中的截断误差；计算机浮点运算还包含舍入误差，因而并非 $h$ 越小总是越好。
\end{solve}

\begin{solve}[习题~\ref{ex:c18q22}：中心差分的截断误差]\phantomsection\label{ans:c18q22}
由奇性 $D_h=\sin(4h)/h$。对正弦展开至三阶，得到
\[
D_h=4-\frac{32}{3}h^2+O(h^4).
\]
真导数为 $4$，所以 $D_h-f^{\prime}(0)=-\frac{32}{3}h^2+O(h^4)$。中心差分误差为二阶，这源于对称相减消去了泰勒展开的偶次项。该式仅描述精确实数运算中的截断误差；计算机浮点运算还包含舍入误差，因而并非 $h$ 越小总是越好。
\end{solve}

\begin{solve}[习题~\ref{ex:c18q23}：中心差分的截断误差]\phantomsection\label{ans:c18q23}
由奇性 $D_h=\sin(5h)/h$。对正弦展开至三阶，得到
\[
D_h=5-\frac{125}{6}h^2+O(h^4).
\]
真导数为 $5$，所以 $D_h-f^{\prime}(0)=-\frac{125}{6}h^2+O(h^4)$。中心差分误差为二阶，这源于对称相减消去了泰勒展开的偶次项。该式仅描述精确实数运算中的截断误差；计算机浮点运算还包含舍入误差，因而并非 $h$ 越小总是越好。
\end{solve}

\begin{solve}[习题~\ref{ex:c18q24}：中心差分的截断误差]\phantomsection\label{ans:c18q24}
由奇性 $D_h=\sin(6h)/h$。对正弦展开至三阶，得到
\[
D_h=6-36h^2+O(h^4).
\]
真导数为 $6$，所以 $D_h-f^{\prime}(0)=-36h^2+O(h^4)$。中心差分误差为二阶，这源于对称相减消去了泰勒展开的偶次项。该式仅描述精确实数运算中的截断误差；计算机浮点运算还包含舍入误差，因而并非 $h$ 越小总是越好。
\end{solve}

\clearpage
\section{反常积分与收敛性}\label{sec:ch19}
反常积分（improper integral）通过极限定义。若区间有多个奇异端点，必须分别检验；把两个发散项相互抵消所得的主值，不能直接当成通常的反常积分。

\textbf{资源对应：}CLP-2 \S~1.12。本章共 25 道编号习题。
\subsection{习题}
\begin{example}[比较判别的方向]\label{ex:c19q1}
若 $0\le f(x)\le g(x)$ 且 $\int_1^\infty f(x)dx$ 收敛，能否断定 $\int_1^\infty g(x)dx$ 收敛？给出证明或反例。
\par\textit{选译。来源：CLP-2 \S~1.12，2015A 判断题。}\par\hyperref[ans:c19q1]{解答见第~\pageref{ans:c19q1} 页。}
\end{example}
\begin{example}[无穷区间上的幂函数]\label{ex:c19q2}
判断 $\displaystyle\int_1^\infty x^{-\frac{1}{2}}dx$ 的收敛性；收敛时求其值。
\par\textit{改编训练。}\par\hyperref[ans:c19q2]{解答见第~\pageref{ans:c19q2} 页。}
\end{example}
\begin{example}[无穷区间上的幂函数]\label{ex:c19q3}
判断 $\displaystyle\int_1^\infty x^{-\frac{2}{3}}dx$ 的收敛性；收敛时求其值。
\par\textit{改编训练。}\par\hyperref[ans:c19q3]{解答见第~\pageref{ans:c19q3} 页。}
\end{example}
\begin{example}[无穷区间上的幂函数]\label{ex:c19q4}
判断 $\displaystyle\int_1^\infty x^{-1}dx$ 的收敛性；收敛时求其值。
\par\textit{改编训练。}\par\hyperref[ans:c19q4]{解答见第~\pageref{ans:c19q4} 页。}
\end{example}
\begin{example}[无穷区间上的幂函数]\label{ex:c19q5}
判断 $\displaystyle\int_1^\infty x^{-\frac{3}{2}}dx$ 的收敛性；收敛时求其值。
\par\textit{改编训练。}\par\hyperref[ans:c19q5]{解答见第~\pageref{ans:c19q5} 页。}
\end{example}
\begin{example}[无穷区间上的幂函数]\label{ex:c19q6}
判断 $\displaystyle\int_1^\infty x^{-2}dx$ 的收敛性；收敛时求其值。
\par\textit{改编训练。}\par\hyperref[ans:c19q6]{解答见第~\pageref{ans:c19q6} 页。}
\end{example}
\begin{example}[无穷区间上的幂函数]\label{ex:c19q7}
判断 $\displaystyle\int_1^\infty x^{-3}dx$ 的收敛性；收敛时求其值。
\par\textit{改编训练。}\par\hyperref[ans:c19q7]{解答见第~\pageref{ans:c19q7} 页。}
\end{example}
\begin{example}[有限端点的幂奇异性]\label{ex:c19q8}
判断 $\displaystyle\int_0^1x^{-\frac{1}{3}}dx$ 的收敛性；收敛时求值。
\par\textit{改编训练。}\par\hyperref[ans:c19q8]{解答见第~\pageref{ans:c19q8} 页。}
\end{example}
\begin{example}[有限端点的幂奇异性]\label{ex:c19q9}
判断 $\displaystyle\int_0^1x^{-\frac{1}{2}}dx$ 的收敛性；收敛时求值。
\par\textit{改编训练。}\par\hyperref[ans:c19q9]{解答见第~\pageref{ans:c19q9} 页。}
\end{example}
\begin{example}[有限端点的幂奇异性]\label{ex:c19q10}
判断 $\displaystyle\int_0^1x^{-\frac{2}{3}}dx$ 的收敛性；收敛时求值。
\par\textit{改编训练。}\par\hyperref[ans:c19q10]{解答见第~\pageref{ans:c19q10} 页。}
\end{example}
\begin{example}[有限端点的幂奇异性]\label{ex:c19q11}
判断 $\displaystyle\int_0^1x^{-1}dx$ 的收敛性；收敛时求值。
\par\textit{改编训练。}\par\hyperref[ans:c19q11]{解答见第~\pageref{ans:c19q11} 页。}
\end{example}
\begin{example}[有限端点的幂奇异性]\label{ex:c19q12}
判断 $\displaystyle\int_0^1x^{-\frac{3}{2}}dx$ 的收敛性；收敛时求值。
\par\textit{改编训练。}\par\hyperref[ans:c19q12]{解答见第~\pageref{ans:c19q12} 页。}
\end{example}
\begin{example}[有限端点的幂奇异性]\label{ex:c19q13}
判断 $\displaystyle\int_0^1x^{-2}dx$ 的收敛性；收敛时求值。
\par\textit{改编训练。}\par\hyperref[ans:c19q13]{解答见第~\pageref{ans:c19q13} 页。}
\end{example}
\begin{example}[对数修正的临界情形]\label{ex:c19q14}
判断 $\displaystyle\int_e^\infty\frac{dx}{x(\log x)^{\frac{1}{2}}}$ 的敛散性，收敛时求值。
\par\textit{改编训练。}\par\hyperref[ans:c19q14]{解答见第~\pageref{ans:c19q14} 页。}
\end{example}
\begin{example}[对数修正的临界情形]\label{ex:c19q15}
判断 $\displaystyle\int_e^\infty\frac{dx}{x(\log x)^{1}}$ 的敛散性，收敛时求值。
\par\textit{改编训练。}\par\hyperref[ans:c19q15]{解答见第~\pageref{ans:c19q15} 页。}
\end{example}
\begin{example}[对数修正的临界情形]\label{ex:c19q16}
判断 $\displaystyle\int_e^\infty\frac{dx}{x(\log x)^{\frac{3}{2}}}$ 的敛散性，收敛时求值。
\par\textit{改编训练。}\par\hyperref[ans:c19q16]{解答见第~\pageref{ans:c19q16} 页。}
\end{example}
\begin{example}[对数修正的临界情形]\label{ex:c19q17}
判断 $\displaystyle\int_e^\infty\frac{dx}{x(\log x)^{2}}$ 的敛散性，收敛时求值。
\par\textit{改编训练。}\par\hyperref[ans:c19q17]{解答见第~\pageref{ans:c19q17} 页。}
\end{example}
\begin{example}[对数修正的临界情形]\label{ex:c19q18}
判断 $\displaystyle\int_e^\infty\frac{dx}{x(\log x)^{3}}$ 的敛散性，收敛时求值。
\par\textit{改编训练。}\par\hyperref[ans:c19q18]{解答见第~\pageref{ans:c19q18} 页。}
\end{example}
\begin{example}[对数修正的临界情形]\label{ex:c19q19}
判断 $\displaystyle\int_e^\infty\frac{dx}{x(\log x)^{4}}$ 的敛散性，收敛时求值。
\par\textit{改编训练。}\par\hyperref[ans:c19q19]{解答见第~\pageref{ans:c19q19} 页。}
\end{example}
\begin{example}[指数衰减积分的递推]\label{ex:c19q20}
计算 $I=\displaystyle\int_0^\infty x^{0}e^{-1x}dx$，并证明所用边界项趋于零。
\par\textit{改编训练。}\par\hyperref[ans:c19q20]{解答见第~\pageref{ans:c19q20} 页。}
\end{example}
\begin{example}[指数衰减积分的递推]\label{ex:c19q21}
计算 $I=\displaystyle\int_0^\infty x^{1}e^{-2x}dx$，并证明所用边界项趋于零。
\par\textit{改编训练。}\par\hyperref[ans:c19q21]{解答见第~\pageref{ans:c19q21} 页。}
\end{example}
\begin{example}[指数衰减积分的递推]\label{ex:c19q22}
计算 $I=\displaystyle\int_0^\infty x^{2}e^{-1x}dx$，并证明所用边界项趋于零。
\par\textit{改编训练。}\par\hyperref[ans:c19q22]{解答见第~\pageref{ans:c19q22} 页。}
\end{example}
\begin{example}[指数衰减积分的递推]\label{ex:c19q23}
计算 $I=\displaystyle\int_0^\infty x^{2}e^{-3x}dx$，并证明所用边界项趋于零。
\par\textit{改编训练。}\par\hyperref[ans:c19q23]{解答见第~\pageref{ans:c19q23} 页。}
\end{example}
\begin{example}[指数衰减积分的递推]\label{ex:c19q24}
计算 $I=\displaystyle\int_0^\infty x^{3}e^{-2x}dx$，并证明所用边界项趋于零。
\par\textit{改编训练。}\par\hyperref[ans:c19q24]{解答见第~\pageref{ans:c19q24} 页。}
\end{example}
\begin{example}[指数衰减积分的递推]\label{ex:c19q25}
计算 $I=\displaystyle\int_0^\infty x^{4}e^{-1x}dx$，并证明所用边界项趋于零。
\par\textit{改编训练。}\par\hyperref[ans:c19q25]{解答见第~\pageref{ans:c19q25} 页。}
\end{example}
\subsection{解答}
\begin{solve}[习题~\ref{ex:c19q1}：比较判别的方向]\phantomsection\label{ans:c19q1}
不能。取 $f(x)=1/x^2$、$g(x)=1/x$。在 $[1,\infty)$ 上满足大小关系，但前者积分为 $1$，后者发散。比较判别正确的方向是：由较大非负函数的积分收敛推出较小者收敛；由较小非负函数的积分发散推出较大者发散。
\end{solve}

\begin{solve}[习题~\ref{ex:c19q2}：无穷区间上的幂函数]\phantomsection\label{ans:c19q2}
先截断积分上限，在有限区间 $[1,R]$ 上积分。
\[
\int_1^Rx^{-\frac{1}{2}}dx=\frac{R^{\frac{1}{2}}-1}{\frac{1}{2}}.
\]
当 $R\to\infty$ 时，右端趋于正无穷，故积分发散。无穷区间上 $p$ 型积分的阈值是 $p>1$，临界点 $p=1$ 也发散。判断首先依赖极限是否为有限实数，而不是是否能够写出原函数。
\end{solve}

\begin{solve}[习题~\ref{ex:c19q3}：无穷区间上的幂函数]\phantomsection\label{ans:c19q3}
先截断积分上限，在有限区间 $[1,R]$ 上积分。
\[
\int_1^Rx^{-\frac{2}{3}}dx=\frac{R^{\frac{1}{3}}-1}{\frac{1}{3}}.
\]
当 $R\to\infty$ 时，右端趋于正无穷，故积分发散。无穷区间上 $p$ 型积分的阈值是 $p>1$，临界点 $p=1$ 也发散。判断首先依赖极限是否为有限实数，而不是是否能够写出原函数。
\end{solve}

\begin{solve}[习题~\ref{ex:c19q4}：无穷区间上的幂函数]\phantomsection\label{ans:c19q4}
先截断积分上限，在有限区间 $[1,R]$ 上积分。
\[
\int_1^R\frac{dx}{x}=\log R.
\]
当 $R\to\infty$ 时，右端趋于正无穷，故积分发散。无穷区间上 $p$ 型积分的阈值是 $p>1$，临界点 $p=1$ 也发散。判断首先依赖极限是否为有限实数，而不是是否能够写出原函数。
\end{solve}

\begin{solve}[习题~\ref{ex:c19q5}：无穷区间上的幂函数]\phantomsection\label{ans:c19q5}
先截断积分上限，在有限区间 $[1,R]$ 上积分。
\[
\int_1^Rx^{-\frac{3}{2}}dx=\frac{R^{-\frac{1}{2}}-1}{-\frac{1}{2}}.
\]
当 $R\to\infty$ 时，幂项趋于零，故积分收敛，值为 $2$。无穷区间上 $p$ 型积分的阈值是 $p>1$，临界点 $p=1$ 也发散。判断首先依赖极限是否为有限实数，而不是是否能够写出原函数。
\end{solve}

\begin{solve}[习题~\ref{ex:c19q6}：无穷区间上的幂函数]\phantomsection\label{ans:c19q6}
先截断积分上限，在有限区间 $[1,R]$ 上积分。
\[
\int_1^Rx^{-2}dx=\frac{R^{-1}-1}{-1}.
\]
当 $R\to\infty$ 时，幂项趋于零，故积分收敛，值为 $1$。无穷区间上 $p$ 型积分的阈值是 $p>1$，临界点 $p=1$ 也发散。判断首先依赖极限是否为有限实数，而不是是否能够写出原函数。
\end{solve}

\begin{solve}[习题~\ref{ex:c19q7}：无穷区间上的幂函数]\phantomsection\label{ans:c19q7}
先截断积分上限，在有限区间 $[1,R]$ 上积分。
\[
\int_1^Rx^{-3}dx=\frac{R^{-2}-1}{-2}.
\]
当 $R\to\infty$ 时，幂项趋于零，故积分收敛，值为 $\frac{1}{2}$。无穷区间上 $p$ 型积分的阈值是 $p>1$，临界点 $p=1$ 也发散。判断首先依赖极限是否为有限实数，而不是是否能够写出原函数。
\end{solve}

\begin{solve}[习题~\ref{ex:c19q8}：有限端点的幂奇异性]\phantomsection\label{ans:c19q8}
把下限替换为 $\varepsilon>0$，最后令 $\varepsilon\to0^+$。
\[
\int_\varepsilon^1x^{-\frac{1}{3}}dx=\frac{1-\varepsilon^{\frac{2}{3}}}{\frac{2}{3}}.
\]
指数 $1-p>0$，所以极限有限，积分值为 $\frac{3}{2}$。与无穷远端的情形相反，有限端点 $0$ 附近的 $p$ 型奇异性要求 $p<1$ 才可积。不能把两个判别条件记反。
\end{solve}

\begin{solve}[习题~\ref{ex:c19q9}：有限端点的幂奇异性]\phantomsection\label{ans:c19q9}
把下限替换为 $\varepsilon>0$，最后令 $\varepsilon\to0^+$。
\[
\int_\varepsilon^1x^{-\frac{1}{2}}dx=\frac{1-\varepsilon^{\frac{1}{2}}}{\frac{1}{2}}.
\]
指数 $1-p>0$，所以极限有限，积分值为 $2$。与无穷远端的情形相反，有限端点 $0$ 附近的 $p$ 型奇异性要求 $p<1$ 才可积。不能把两个判别条件记反。
\end{solve}

\begin{solve}[习题~\ref{ex:c19q10}：有限端点的幂奇异性]\phantomsection\label{ans:c19q10}
把下限替换为 $\varepsilon>0$，最后令 $\varepsilon\to0^+$。
\[
\int_\varepsilon^1x^{-\frac{2}{3}}dx=\frac{1-\varepsilon^{\frac{1}{3}}}{\frac{1}{3}}.
\]
指数 $1-p>0$，所以极限有限，积分值为 $3$。与无穷远端的情形相反，有限端点 $0$ 附近的 $p$ 型奇异性要求 $p<1$ 才可积。不能把两个判别条件记反。
\end{solve}

\begin{solve}[习题~\ref{ex:c19q11}：有限端点的幂奇异性]\phantomsection\label{ans:c19q11}
把下限替换为 $\varepsilon>0$，最后令 $\varepsilon\to0^+$。
\[
\int_\varepsilon^1\frac{dx}{x}=-\log\varepsilon.
\]
右端在极限过程中趋于正无穷，积分发散。与无穷远端的情形相反，有限端点 $0$ 附近的 $p$ 型奇异性要求 $p<1$ 才可积。不能把两个判别条件记反。
\end{solve}

\begin{solve}[习题~\ref{ex:c19q12}：有限端点的幂奇异性]\phantomsection\label{ans:c19q12}
把下限替换为 $\varepsilon>0$，最后令 $\varepsilon\to0^+$。
\[
\int_\varepsilon^1x^{-\frac{3}{2}}dx=\frac{1-\varepsilon^{-\frac{1}{2}}}{-\frac{1}{2}}.
\]
右端在极限过程中趋于正无穷，积分发散。与无穷远端的情形相反，有限端点 $0$ 附近的 $p$ 型奇异性要求 $p<1$ 才可积。不能把两个判别条件记反。
\end{solve}

\begin{solve}[习题~\ref{ex:c19q13}：有限端点的幂奇异性]\phantomsection\label{ans:c19q13}
把下限替换为 $\varepsilon>0$，最后令 $\varepsilon\to0^+$。
\[
\int_\varepsilon^1x^{-2}dx=\frac{1-\varepsilon^{-1}}{-1}.
\]
右端在极限过程中趋于正无穷，积分发散。与无穷远端的情形相反，有限端点 $0$ 附近的 $p$ 型奇异性要求 $p<1$ 才可积。不能把两个判别条件记反。
\end{solve}

\begin{solve}[习题~\ref{ex:c19q14}：对数修正的临界情形]\phantomsection\label{ans:c19q14}
在截断区间上令 $u=\log x$，则 $du=dx/x$，$x=e$ 对应 $u=1$。因此原问题化为
\[
\int_e^\infty\frac{dx}{x(\log x)^{\frac{1}{2}}}=\int_1^\infty u^{-\frac{1}{2}}\,du.
\]
由 $p$ 型判别，积分发散。这一等式应理解为先在有限区间完成换元，再对上限取极限。额外的对数因子虽然增长缓慢，却足以在幂指数超过 $1$ 时改变调和型积分的收敛性。
\end{solve}

\begin{solve}[习题~\ref{ex:c19q15}：对数修正的临界情形]\phantomsection\label{ans:c19q15}
在截断区间上令 $u=\log x$，则 $du=dx/x$，$x=e$ 对应 $u=1$。因此原问题化为
\[
\int_e^\infty\frac{dx}{x(\log x)^{1}}=\int_1^\infty u^{-1}\,du.
\]
由 $p$ 型判别，积分发散。这一等式应理解为先在有限区间完成换元，再对上限取极限。额外的对数因子虽然增长缓慢，却足以在幂指数超过 $1$ 时改变调和型积分的收敛性。
\end{solve}

\begin{solve}[习题~\ref{ex:c19q16}：对数修正的临界情形]\phantomsection\label{ans:c19q16}
在截断区间上令 $u=\log x$，则 $du=dx/x$，$x=e$ 对应 $u=1$。因此原问题化为
\[
\int_e^\infty\frac{dx}{x(\log x)^{\frac{3}{2}}}=\int_1^\infty u^{-\frac{3}{2}}\,du.
\]
由 $p$ 型判别，积分收敛，值为 $2$。这一等式应理解为先在有限区间完成换元，再对上限取极限。额外的对数因子虽然增长缓慢，却足以在幂指数超过 $1$ 时改变调和型积分的收敛性。
\end{solve}

\begin{solve}[习题~\ref{ex:c19q17}：对数修正的临界情形]\phantomsection\label{ans:c19q17}
在截断区间上令 $u=\log x$，则 $du=dx/x$，$x=e$ 对应 $u=1$。因此原问题化为
\[
\int_e^\infty\frac{dx}{x(\log x)^{2}}=\int_1^\infty u^{-2}\,du.
\]
由 $p$ 型判别，积分收敛，值为 $1$。这一等式应理解为先在有限区间完成换元，再对上限取极限。额外的对数因子虽然增长缓慢，却足以在幂指数超过 $1$ 时改变调和型积分的收敛性。
\end{solve}

\begin{solve}[习题~\ref{ex:c19q18}：对数修正的临界情形]\phantomsection\label{ans:c19q18}
在截断区间上令 $u=\log x$，则 $du=dx/x$，$x=e$ 对应 $u=1$。因此原问题化为
\[
\int_e^\infty\frac{dx}{x(\log x)^{3}}=\int_1^\infty u^{-3}\,du.
\]
由 $p$ 型判别，积分收敛，值为 $\frac{1}{2}$。这一等式应理解为先在有限区间完成换元，再对上限取极限。额外的对数因子虽然增长缓慢，却足以在幂指数超过 $1$ 时改变调和型积分的收敛性。
\end{solve}

\begin{solve}[习题~\ref{ex:c19q19}：对数修正的临界情形]\phantomsection\label{ans:c19q19}
在截断区间上令 $u=\log x$，则 $du=dx/x$，$x=e$ 对应 $u=1$。因此原问题化为
\[
\int_e^\infty\frac{dx}{x(\log x)^{4}}=\int_1^\infty u^{-4}\,du.
\]
由 $p$ 型判别，积分收敛，值为 $\frac{1}{3}$。这一等式应理解为先在有限区间完成换元，再对上限取极限。额外的对数因子虽然增长缓慢，却足以在幂指数超过 $1$ 时改变调和型积分的收敛性。
\end{solve}

\begin{solve}[习题~\ref{ex:c19q20}：指数衰减积分的递推]\phantomsection\label{ans:c19q20}
令 $I_k=\int_0^\infty x^ke^{-1x}dx$。在有限区间上分部积分，再令上限趋于无穷。指数比任意固定次数的多项式增长更快，故 $x^ke^{-1x}\to0$，例如由指数级数中的 $k+1$ 阶项即可给出上界。于是
\[
I_0=1,\qquad I_k=\frac k{1}I_{k-1}\quad(k\ge1).
\]
递推得到 $I_k=k!/1^{k+1}$，所以本题
\[
I=1.
\]
每一步都先在有限区间使用分部积分，避免在尚未确认收敛时直接对无穷端点代值。
\end{solve}

\begin{solve}[习题~\ref{ex:c19q21}：指数衰减积分的递推]\phantomsection\label{ans:c19q21}
令 $I_k=\int_0^\infty x^ke^{-2x}dx$。在有限区间上分部积分，再令上限趋于无穷。指数比任意固定次数的多项式增长更快，故 $x^ke^{-2x}\to0$，例如由指数级数中的 $k+1$ 阶项即可给出上界。于是
\[
I_0=\frac{1}{2},\qquad I_k=\frac k{2}I_{k-1}\quad(k\ge1).
\]
递推得到 $I_k=k!/2^{k+1}$，所以本题
\[
I=\frac{1}{4}.
\]
每一步都先在有限区间使用分部积分，避免在尚未确认收敛时直接对无穷端点代值。
\end{solve}

\begin{solve}[习题~\ref{ex:c19q22}：指数衰减积分的递推]\phantomsection\label{ans:c19q22}
令 $I_k=\int_0^\infty x^ke^{-1x}dx$。在有限区间上分部积分，再令上限趋于无穷。指数比任意固定次数的多项式增长更快，故 $x^ke^{-1x}\to0$，例如由指数级数中的 $k+1$ 阶项即可给出上界。于是
\[
I_0=1,\qquad I_k=\frac k{1}I_{k-1}\quad(k\ge1).
\]
递推得到 $I_k=k!/1^{k+1}$，所以本题
\[
I=2.
\]
每一步都先在有限区间使用分部积分，避免在尚未确认收敛时直接对无穷端点代值。
\end{solve}

\begin{solve}[习题~\ref{ex:c19q23}：指数衰减积分的递推]\phantomsection\label{ans:c19q23}
令 $I_k=\int_0^\infty x^ke^{-3x}dx$。在有限区间上分部积分，再令上限趋于无穷。指数比任意固定次数的多项式增长更快，故 $x^ke^{-3x}\to0$，例如由指数级数中的 $k+1$ 阶项即可给出上界。于是
\[
I_0=\frac{1}{3},\qquad I_k=\frac k{3}I_{k-1}\quad(k\ge1).
\]
递推得到 $I_k=k!/3^{k+1}$，所以本题
\[
I=\frac{2}{27}.
\]
每一步都先在有限区间使用分部积分，避免在尚未确认收敛时直接对无穷端点代值。
\end{solve}

\begin{solve}[习题~\ref{ex:c19q24}：指数衰减积分的递推]\phantomsection\label{ans:c19q24}
令 $I_k=\int_0^\infty x^ke^{-2x}dx$。在有限区间上分部积分，再令上限趋于无穷。指数比任意固定次数的多项式增长更快，故 $x^ke^{-2x}\to0$，例如由指数级数中的 $k+1$ 阶项即可给出上界。于是
\[
I_0=\frac{1}{2},\qquad I_k=\frac k{2}I_{k-1}\quad(k\ge1).
\]
递推得到 $I_k=k!/2^{k+1}$，所以本题
\[
I=\frac{3}{8}.
\]
每一步都先在有限区间使用分部积分，避免在尚未确认收敛时直接对无穷端点代值。
\end{solve}

\begin{solve}[习题~\ref{ex:c19q25}：指数衰减积分的递推]\phantomsection\label{ans:c19q25}
令 $I_k=\int_0^\infty x^ke^{-1x}dx$。在有限区间上分部积分，再令上限趋于无穷。指数比任意固定次数的多项式增长更快，故 $x^ke^{-1x}\to0$，例如由指数级数中的 $k+1$ 阶项即可给出上界。于是
\[
I_0=1,\qquad I_k=\frac k{1}I_{k-1}\quad(k\ge1).
\]
递推得到 $I_k=k!/1^{k+1}$，所以本题
\[
I=24.
\]
每一步都先在有限区间使用分部积分，避免在尚未确认收敛时直接对无穷端点代值。
\end{solve}

\clearpage
\section{含参积分与 Beta 型训练}\label{sec:ch20}
含参积分中的参数是给定的常数，积分变量才是求和累积的对象。对参数求导、交换积分等操作需要控制条件；本章先用有限区间积分或明确的可积控制函数保证这些步骤。

\textbf{资源对应：}CLP-2 \S~1.7、\S~1.12、\S~1.13 的扩展训练。本章共 24 道编号习题。
\subsection{习题}
\begin{example}[有限区间的 Beta 型递推]\label{ex:c20q1}
求 $I=\displaystyle\int_0^1x^{1}(1-x)^2dx$，并由分部积分解释结果中的连续因子。
\par\textit{改编训练。}\par\hyperref[ans:c20q1]{解答见第~\pageref{ans:c20q1} 页。}
\end{example}
\begin{example}[有限区间的 Beta 型递推]\label{ex:c20q2}
求 $I=\displaystyle\int_0^1x^{2}(1-x)^2dx$，并由分部积分解释结果中的连续因子。
\par\textit{改编训练。}\par\hyperref[ans:c20q2]{解答见第~\pageref{ans:c20q2} 页。}
\end{example}
\begin{example}[有限区间的 Beta 型递推]\label{ex:c20q3}
求 $I=\displaystyle\int_0^1x^{3}(1-x)^2dx$，并由分部积分解释结果中的连续因子。
\par\textit{改编训练。}\par\hyperref[ans:c20q3]{解答见第~\pageref{ans:c20q3} 页。}
\end{example}
\begin{example}[有限区间的 Beta 型递推]\label{ex:c20q4}
求 $I=\displaystyle\int_0^1x^{4}(1-x)^2dx$，并由分部积分解释结果中的连续因子。
\par\textit{改编训练。}\par\hyperref[ans:c20q4]{解答见第~\pageref{ans:c20q4} 页。}
\end{example}
\begin{example}[有限区间的 Beta 型递推]\label{ex:c20q5}
求 $I=\displaystyle\int_0^1x^{5}(1-x)^2dx$，并由分部积分解释结果中的连续因子。
\par\textit{改编训练。}\par\hyperref[ans:c20q5]{解答见第~\pageref{ans:c20q5} 页。}
\end{example}
\begin{example}[有限区间的 Beta 型递推]\label{ex:c20q6}
求 $I=\displaystyle\int_0^1x^{6}(1-x)^2dx$，并由分部积分解释结果中的连续因子。
\par\textit{改编训练。}\par\hyperref[ans:c20q6]{解答见第~\pageref{ans:c20q6} 页。}
\end{example}
\begin{example}[对数权重的定积分]\label{ex:c20q7}
计算 $\displaystyle\int_0^1x^{1}(\log x)^2dx$，并说明端点收敛。
\par\textit{改编训练。}\par\hyperref[ans:c20q7]{解答见第~\pageref{ans:c20q7} 页。}
\end{example}
\begin{example}[对数权重的定积分]\label{ex:c20q8}
计算 $\displaystyle\int_0^1x^{2}(\log x)^2dx$，并说明端点收敛。
\par\textit{改编训练。}\par\hyperref[ans:c20q8]{解答见第~\pageref{ans:c20q8} 页。}
\end{example}
\begin{example}[对数权重的定积分]\label{ex:c20q9}
计算 $\displaystyle\int_0^1x^{3}(\log x)^2dx$，并说明端点收敛。
\par\textit{改编训练。}\par\hyperref[ans:c20q9]{解答见第~\pageref{ans:c20q9} 页。}
\end{example}
\begin{example}[对数权重的定积分]\label{ex:c20q10}
计算 $\displaystyle\int_0^1x^{4}(\log x)^2dx$，并说明端点收敛。
\par\textit{改编训练。}\par\hyperref[ans:c20q10]{解答见第~\pageref{ans:c20q10} 页。}
\end{example}
\begin{example}[对数权重的定积分]\label{ex:c20q11}
计算 $\displaystyle\int_0^1x^{5}(\log x)^2dx$，并说明端点收敛。
\par\textit{改编训练。}\par\hyperref[ans:c20q11]{解答见第~\pageref{ans:c20q11} 页。}
\end{example}
\begin{example}[对数权重的定积分]\label{ex:c20q12}
计算 $\displaystyle\int_0^1x^{6}(\log x)^2dx$，并说明端点收敛。
\par\textit{改编训练。}\par\hyperref[ans:c20q12]{解答见第~\pageref{ans:c20q12} 页。}
\end{example}
\begin{example}[两个衰减指数的消差]\label{ex:c20q13}
计算 $I=\displaystyle\int_0^\infty\frac{e^{-1x}-e^{-2x}}x dx$。先证明它收敛。
\par\textit{改编训练。}\par\hyperref[ans:c20q13]{解答见第~\pageref{ans:c20q13} 页。}
\end{example}
\begin{example}[两个衰减指数的消差]\label{ex:c20q14}
计算 $I=\displaystyle\int_0^\infty\frac{e^{-1x}-e^{-3x}}x dx$。先证明它收敛。
\par\textit{改编训练。}\par\hyperref[ans:c20q14]{解答见第~\pageref{ans:c20q14} 页。}
\end{example}
\begin{example}[两个衰减指数的消差]\label{ex:c20q15}
计算 $I=\displaystyle\int_0^\infty\frac{e^{-2x}-e^{-3x}}x dx$。先证明它收敛。
\par\textit{改编训练。}\par\hyperref[ans:c20q15]{解答见第~\pageref{ans:c20q15} 页。}
\end{example}
\begin{example}[两个衰减指数的消差]\label{ex:c20q16}
计算 $I=\displaystyle\int_0^\infty\frac{e^{-2x}-e^{-5x}}x dx$。先证明它收敛。
\par\textit{改编训练。}\par\hyperref[ans:c20q16]{解答见第~\pageref{ans:c20q16} 页。}
\end{example}
\begin{example}[两个衰减指数的消差]\label{ex:c20q17}
计算 $I=\displaystyle\int_0^\infty\frac{e^{-3x}-e^{-4x}}x dx$。先证明它收敛。
\par\textit{改编训练。}\par\hyperref[ans:c20q17]{解答见第~\pageref{ans:c20q17} 页。}
\end{example}
\begin{example}[两个衰减指数的消差]\label{ex:c20q18}
计算 $I=\displaystyle\int_0^\infty\frac{e^{-3x}-e^{-7x}}x dx$。先证明它收敛。
\par\textit{改编训练。}\par\hyperref[ans:c20q18]{解答见第~\pageref{ans:c20q18} 页。}
\end{example}
\begin{example}[简单含参积分的求导]\label{ex:c20q19}
定义 $F(t)=\int_0^1\log(1+tx)dx$（$t>0$）。求 $F(t)$、$F^{\prime}(t)$，并计算 $F(1)$。
\par\textit{改编训练。}\par\hyperref[ans:c20q19]{解答见第~\pageref{ans:c20q19} 页。}
\end{example}
\begin{example}[简单含参积分的求导]\label{ex:c20q20}
定义 $F(t)=\int_0^1\log(1+tx)dx$（$t>0$）。求 $F(t)$、$F^{\prime}(t)$，并计算 $F(2)$。
\par\textit{改编训练。}\par\hyperref[ans:c20q20]{解答见第~\pageref{ans:c20q20} 页。}
\end{example}
\begin{example}[简单含参积分的求导]\label{ex:c20q21}
定义 $F(t)=\int_0^1\log(1+tx)dx$（$t>0$）。求 $F(t)$、$F^{\prime}(t)$，并计算 $F(3)$。
\par\textit{改编训练。}\par\hyperref[ans:c20q21]{解答见第~\pageref{ans:c20q21} 页。}
\end{example}
\begin{example}[简单含参积分的求导]\label{ex:c20q22}
定义 $F(t)=\int_0^1\log(1+tx)dx$（$t>0$）。求 $F(t)$、$F^{\prime}(t)$，并计算 $F(4)$。
\par\textit{改编训练。}\par\hyperref[ans:c20q22]{解答见第~\pageref{ans:c20q22} 页。}
\end{example}
\begin{example}[简单含参积分的求导]\label{ex:c20q23}
定义 $F(t)=\int_0^1\log(1+tx)dx$（$t>0$）。求 $F(t)$、$F^{\prime}(t)$，并计算 $F(5)$。
\par\textit{改编训练。}\par\hyperref[ans:c20q23]{解答见第~\pageref{ans:c20q23} 页。}
\end{example}
\begin{example}[简单含参积分的求导]\label{ex:c20q24}
定义 $F(t)=\int_0^1\log(1+tx)dx$（$t>0$）。求 $F(t)$、$F^{\prime}(t)$，并计算 $F(6)$。
\par\textit{改编训练。}\par\hyperref[ans:c20q24]{解答见第~\pageref{ans:c20q24} 页。}
\end{example}
\subsection{解答}
\begin{solve}[习题~\ref{ex:c20q1}：有限区间的 Beta 型递推]\phantomsection\label{ans:c20q1}
直接展开得到 $I=1/2-2/3+1/4$。也可取 $u=(1-x)^2$、$dv=x^{1}dx$，两端边界项为零，故
\[
I=\frac2{2}\int_0^1x^{2}(1-x)dx=\frac2{2}\left(\frac1{3}-\frac1{4}\right).
\]
所以 $I=\frac{1}{12}$。分部积分逐次降低 $(1-x)$ 的次数，同时提高 $x$ 的次数，因而自然产生三个连续的分母因子。
\end{solve}

\begin{solve}[习题~\ref{ex:c20q2}：有限区间的 Beta 型递推]\phantomsection\label{ans:c20q2}
直接展开得到 $I=1/3-2/4+1/5$。也可取 $u=(1-x)^2$、$dv=x^{2}dx$，两端边界项为零，故
\[
I=\frac2{3}\int_0^1x^{3}(1-x)dx=\frac2{3}\left(\frac1{4}-\frac1{5}\right).
\]
所以 $I=\frac{1}{30}$。分部积分逐次降低 $(1-x)$ 的次数，同时提高 $x$ 的次数，因而自然产生三个连续的分母因子。
\end{solve}

\begin{solve}[习题~\ref{ex:c20q3}：有限区间的 Beta 型递推]\phantomsection\label{ans:c20q3}
直接展开得到 $I=1/4-2/5+1/6$。也可取 $u=(1-x)^2$、$dv=x^{3}dx$，两端边界项为零，故
\[
I=\frac2{4}\int_0^1x^{4}(1-x)dx=\frac2{4}\left(\frac1{5}-\frac1{6}\right).
\]
所以 $I=\frac{1}{60}$。分部积分逐次降低 $(1-x)$ 的次数，同时提高 $x$ 的次数，因而自然产生三个连续的分母因子。
\end{solve}

\begin{solve}[习题~\ref{ex:c20q4}：有限区间的 Beta 型递推]\phantomsection\label{ans:c20q4}
直接展开得到 $I=1/5-2/6+1/7$。也可取 $u=(1-x)^2$、$dv=x^{4}dx$，两端边界项为零，故
\[
I=\frac2{5}\int_0^1x^{5}(1-x)dx=\frac2{5}\left(\frac1{6}-\frac1{7}\right).
\]
所以 $I=\frac{1}{105}$。分部积分逐次降低 $(1-x)$ 的次数，同时提高 $x$ 的次数，因而自然产生三个连续的分母因子。
\end{solve}

\begin{solve}[习题~\ref{ex:c20q5}：有限区间的 Beta 型递推]\phantomsection\label{ans:c20q5}
直接展开得到 $I=1/6-2/7+1/8$。也可取 $u=(1-x)^2$、$dv=x^{5}dx$，两端边界项为零，故
\[
I=\frac2{6}\int_0^1x^{6}(1-x)dx=\frac2{6}\left(\frac1{7}-\frac1{8}\right).
\]
所以 $I=\frac{1}{168}$。分部积分逐次降低 $(1-x)$ 的次数，同时提高 $x$ 的次数，因而自然产生三个连续的分母因子。
\end{solve}

\begin{solve}[习题~\ref{ex:c20q6}：有限区间的 Beta 型递推]\phantomsection\label{ans:c20q6}
直接展开得到 $I=1/7-2/8+1/9$。也可取 $u=(1-x)^2$、$dv=x^{6}dx$，两端边界项为零，故
\[
I=\frac2{7}\int_0^1x^{7}(1-x)dx=\frac2{7}\left(\frac1{8}-\frac1{9}\right).
\]
所以 $I=\frac{1}{252}$。分部积分逐次降低 $(1-x)$ 的次数，同时提高 $x$ 的次数，因而自然产生三个连续的分母因子。
\end{solve}

\begin{solve}[习题~\ref{ex:c20q7}：对数权重的定积分]\phantomsection\label{ans:c20q7}
令 $x=e^{-t}$，$dx=-e^{-t}dt$，则积分转化为
\[
\int_0^\infty t^2e^{-2t}dt.
\]
这是指数衰减积分，对它分部积分两次得到
\[
\int_0^1x^{1}(\log x)^2dx=\frac{1}{4}.
\]
变换后的被积函数非负且尾部由指数控制，所以积分收敛。原变量中虽然对数在0处无界，其平方仍被 $x^a$ 的衰减压制到可积程度。也可对 $\int_0^1x^pdx=1/(p+1)$ 对参数求导两次，但那需要先证明交换求导与积分合法。
\end{solve}

\begin{solve}[习题~\ref{ex:c20q8}：对数权重的定积分]\phantomsection\label{ans:c20q8}
令 $x=e^{-t}$，$dx=-e^{-t}dt$，则积分转化为
\[
\int_0^\infty t^2e^{-3t}dt.
\]
这是指数衰减积分，对它分部积分两次得到
\[
\int_0^1x^{2}(\log x)^2dx=\frac{2}{27}.
\]
变换后的被积函数非负且尾部由指数控制，所以积分收敛。原变量中虽然对数在0处无界，其平方仍被 $x^a$ 的衰减压制到可积程度。也可对 $\int_0^1x^pdx=1/(p+1)$ 对参数求导两次，但那需要先证明交换求导与积分合法。
\end{solve}

\begin{solve}[习题~\ref{ex:c20q9}：对数权重的定积分]\phantomsection\label{ans:c20q9}
令 $x=e^{-t}$，$dx=-e^{-t}dt$，则积分转化为
\[
\int_0^\infty t^2e^{-4t}dt.
\]
这是指数衰减积分，对它分部积分两次得到
\[
\int_0^1x^{3}(\log x)^2dx=\frac{1}{32}.
\]
变换后的被积函数非负且尾部由指数控制，所以积分收敛。原变量中虽然对数在0处无界，其平方仍被 $x^a$ 的衰减压制到可积程度。也可对 $\int_0^1x^pdx=1/(p+1)$ 对参数求导两次，但那需要先证明交换求导与积分合法。
\end{solve}

\begin{solve}[习题~\ref{ex:c20q10}：对数权重的定积分]\phantomsection\label{ans:c20q10}
令 $x=e^{-t}$，$dx=-e^{-t}dt$，则积分转化为
\[
\int_0^\infty t^2e^{-5t}dt.
\]
这是指数衰减积分，对它分部积分两次得到
\[
\int_0^1x^{4}(\log x)^2dx=\frac{2}{125}.
\]
变换后的被积函数非负且尾部由指数控制，所以积分收敛。原变量中虽然对数在0处无界，其平方仍被 $x^a$ 的衰减压制到可积程度。也可对 $\int_0^1x^pdx=1/(p+1)$ 对参数求导两次，但那需要先证明交换求导与积分合法。
\end{solve}

\begin{solve}[习题~\ref{ex:c20q11}：对数权重的定积分]\phantomsection\label{ans:c20q11}
令 $x=e^{-t}$，$dx=-e^{-t}dt$，则积分转化为
\[
\int_0^\infty t^2e^{-6t}dt.
\]
这是指数衰减积分，对它分部积分两次得到
\[
\int_0^1x^{5}(\log x)^2dx=\frac{1}{108}.
\]
变换后的被积函数非负且尾部由指数控制，所以积分收敛。原变量中虽然对数在0处无界，其平方仍被 $x^a$ 的衰减压制到可积程度。也可对 $\int_0^1x^pdx=1/(p+1)$ 对参数求导两次，但那需要先证明交换求导与积分合法。
\end{solve}

\begin{solve}[习题~\ref{ex:c20q12}：对数权重的定积分]\phantomsection\label{ans:c20q12}
令 $x=e^{-t}$，$dx=-e^{-t}dt$，则积分转化为
\[
\int_0^\infty t^2e^{-7t}dt.
\]
这是指数衰减积分，对它分部积分两次得到
\[
\int_0^1x^{6}(\log x)^2dx=\frac{2}{343}.
\]
变换后的被积函数非负且尾部由指数控制，所以积分收敛。原变量中虽然对数在0处无界，其平方仍被 $x^a$ 的衰减压制到可积程度。也可对 $\int_0^1x^pdx=1/(p+1)$ 对参数求导两次，但那需要先证明交换求导与积分合法。
\end{solve}

\begin{solve}[习题~\ref{ex:c20q13}：两个衰减指数的消差]\phantomsection\label{ans:c20q13}
在0附近，分子为 $(1)x+O(x^2)$，商有有限极限；无穷远处由指数衰减控制，因此积分收敛。恒等式
\[
\frac{e^{-1x}-e^{-2x}}x=\int_{1}^{2}e^{-tx}dt
\]
把相减结构改为非负二重积分。由非负积分的换序定理，
\[
I=\int_{1}^{2}\left(\int_0^\infty e^{-tx}dx\right)dt=\int_{1}^{2}\frac{dt}t=\log(2).
\]
不能把原积分拆成两个单独积分相减，因为各自都在0处发散，只有消差后的整体可积。
\end{solve}

\begin{solve}[习题~\ref{ex:c20q14}：两个衰减指数的消差]\phantomsection\label{ans:c20q14}
在0附近，分子为 $(2)x+O(x^2)$，商有有限极限；无穷远处由指数衰减控制，因此积分收敛。恒等式
\[
\frac{e^{-1x}-e^{-3x}}x=\int_{1}^{3}e^{-tx}dt
\]
把相减结构改为非负二重积分。由非负积分的换序定理，
\[
I=\int_{1}^{3}\left(\int_0^\infty e^{-tx}dx\right)dt=\int_{1}^{3}\frac{dt}t=\log(3).
\]
不能把原积分拆成两个单独积分相减，因为各自都在0处发散，只有消差后的整体可积。
\end{solve}

\begin{solve}[习题~\ref{ex:c20q15}：两个衰减指数的消差]\phantomsection\label{ans:c20q15}
在0附近，分子为 $(1)x+O(x^2)$，商有有限极限；无穷远处由指数衰减控制，因此积分收敛。恒等式
\[
\frac{e^{-2x}-e^{-3x}}x=\int_{2}^{3}e^{-tx}dt
\]
把相减结构改为非负二重积分。由非负积分的换序定理，
\[
I=\int_{2}^{3}\left(\int_0^\infty e^{-tx}dx\right)dt=\int_{2}^{3}\frac{dt}t=\log(\frac{3}{2}).
\]
不能把原积分拆成两个单独积分相减，因为各自都在0处发散，只有消差后的整体可积。
\end{solve}

\begin{solve}[习题~\ref{ex:c20q16}：两个衰减指数的消差]\phantomsection\label{ans:c20q16}
在0附近，分子为 $(3)x+O(x^2)$，商有有限极限；无穷远处由指数衰减控制，因此积分收敛。恒等式
\[
\frac{e^{-2x}-e^{-5x}}x=\int_{2}^{5}e^{-tx}dt
\]
把相减结构改为非负二重积分。由非负积分的换序定理，
\[
I=\int_{2}^{5}\left(\int_0^\infty e^{-tx}dx\right)dt=\int_{2}^{5}\frac{dt}t=\log(\frac{5}{2}).
\]
不能把原积分拆成两个单独积分相减，因为各自都在0处发散，只有消差后的整体可积。
\end{solve}

\begin{solve}[习题~\ref{ex:c20q17}：两个衰减指数的消差]\phantomsection\label{ans:c20q17}
在0附近，分子为 $(1)x+O(x^2)$，商有有限极限；无穷远处由指数衰减控制，因此积分收敛。恒等式
\[
\frac{e^{-3x}-e^{-4x}}x=\int_{3}^{4}e^{-tx}dt
\]
把相减结构改为非负二重积分。由非负积分的换序定理，
\[
I=\int_{3}^{4}\left(\int_0^\infty e^{-tx}dx\right)dt=\int_{3}^{4}\frac{dt}t=\log(\frac{4}{3}).
\]
不能把原积分拆成两个单独积分相减，因为各自都在0处发散，只有消差后的整体可积。
\end{solve}

\begin{solve}[习题~\ref{ex:c20q18}：两个衰减指数的消差]\phantomsection\label{ans:c20q18}
在0附近，分子为 $(4)x+O(x^2)$，商有有限极限；无穷远处由指数衰减控制，因此积分收敛。恒等式
\[
\frac{e^{-3x}-e^{-7x}}x=\int_{3}^{7}e^{-tx}dt
\]
把相减结构改为非负二重积分。由非负积分的换序定理，
\[
I=\int_{3}^{7}\left(\int_0^\infty e^{-tx}dx\right)dt=\int_{3}^{7}\frac{dt}t=\log(\frac{7}{3}).
\]
不能把原积分拆成两个单独积分相减，因为各自都在0处发散，只有消差后的整体可积。
\end{solve}

\begin{solve}[习题~\ref{ex:c20q19}：简单含参积分的求导]\phantomsection\label{ans:c20q19}
对 $x$ 作代换 $u=1+tx$，得到
\[
F(t)=\frac{(1+t)\log(1+t)-t}{t}.
\]
在任意紧参数区间 $0<\alpha\le t\le\beta$ 上，被积函数对参数的导数 $x/(1+tx)$ 连续且被1控制，故可在积分号内求导：
\[
F^{\prime}(t)=\int_0^1\frac{x}{1+tx}dx=\frac{t-\log(1+t)}{t^2}.
\]
从而 $F(1)=2\log(2)-1$。从显式表达式直接求导也得到相同结果，提供独立核对。
\end{solve}

\begin{solve}[习题~\ref{ex:c20q20}：简单含参积分的求导]\phantomsection\label{ans:c20q20}
对 $x$ 作代换 $u=1+tx$，得到
\[
F(t)=\frac{(1+t)\log(1+t)-t}{t}.
\]
在任意紧参数区间 $0<\alpha\le t\le\beta$ 上，被积函数对参数的导数 $x/(1+tx)$ 连续且被1控制，故可在积分号内求导：
\[
F^{\prime}(t)=\int_0^1\frac{x}{1+tx}dx=\frac{t-\log(1+t)}{t^2}.
\]
从而 $F(2)=\frac{3}{2}\log(3)-1$。从显式表达式直接求导也得到相同结果，提供独立核对。
\end{solve}

\begin{solve}[习题~\ref{ex:c20q21}：简单含参积分的求导]\phantomsection\label{ans:c20q21}
对 $x$ 作代换 $u=1+tx$，得到
\[
F(t)=\frac{(1+t)\log(1+t)-t}{t}.
\]
在任意紧参数区间 $0<\alpha\le t\le\beta$ 上，被积函数对参数的导数 $x/(1+tx)$ 连续且被1控制，故可在积分号内求导：
\[
F^{\prime}(t)=\int_0^1\frac{x}{1+tx}dx=\frac{t-\log(1+t)}{t^2}.
\]
从而 $F(3)=\frac{4}{3}\log(4)-1$。从显式表达式直接求导也得到相同结果，提供独立核对。
\end{solve}

\begin{solve}[习题~\ref{ex:c20q22}：简单含参积分的求导]\phantomsection\label{ans:c20q22}
对 $x$ 作代换 $u=1+tx$，得到
\[
F(t)=\frac{(1+t)\log(1+t)-t}{t}.
\]
在任意紧参数区间 $0<\alpha\le t\le\beta$ 上，被积函数对参数的导数 $x/(1+tx)$ 连续且被1控制，故可在积分号内求导：
\[
F^{\prime}(t)=\int_0^1\frac{x}{1+tx}dx=\frac{t-\log(1+t)}{t^2}.
\]
从而 $F(4)=\frac{5}{4}\log(5)-1$。从显式表达式直接求导也得到相同结果，提供独立核对。
\end{solve}

\begin{solve}[习题~\ref{ex:c20q23}：简单含参积分的求导]\phantomsection\label{ans:c20q23}
对 $x$ 作代换 $u=1+tx$，得到
\[
F(t)=\frac{(1+t)\log(1+t)-t}{t}.
\]
在任意紧参数区间 $0<\alpha\le t\le\beta$ 上，被积函数对参数的导数 $x/(1+tx)$ 连续且被1控制，故可在积分号内求导：
\[
F^{\prime}(t)=\int_0^1\frac{x}{1+tx}dx=\frac{t-\log(1+t)}{t^2}.
\]
从而 $F(5)=\frac{6}{5}\log(6)-1$。从显式表达式直接求导也得到相同结果，提供独立核对。
\end{solve}

\begin{solve}[习题~\ref{ex:c20q24}：简单含参积分的求导]\phantomsection\label{ans:c20q24}
对 $x$ 作代换 $u=1+tx$，得到
\[
F(t)=\frac{(1+t)\log(1+t)-t}{t}.
\]
在任意紧参数区间 $0<\alpha\le t\le\beta$ 上，被积函数对参数的导数 $x/(1+tx)$ 连续且被1控制，故可在积分号内求导：
\[
F^{\prime}(t)=\int_0^1\frac{x}{1+tx}dx=\frac{t-\log(1+t)}{t^2}.
\]
从而 $F(6)=\frac{7}{6}\log(7)-1$。从显式表达式直接求导也得到相同结果，提供独立核对。
\end{solve}

\clearpage
\section{积分的几何与物理应用}\label{sec:ch21}
把一个总量写成积分，需要明确微元的含义。面积使用上下边界之差，旋转体积使用截面积，质心使用一阶矩除以质量；只有在密度为常数时才能把质量与几何面积直接对应。

\textbf{资源对应：}CLP-2 \S~1.5、\S~1.6、\S~2.1--\S~2.3。本章共 24 道编号习题。
\subsection{习题}
\begin{example}[两条幂曲线之间的面积]\label{ex:c21q1}
求曲线 $y=x$ 与 $y=x^{2}$ 在 $0\le x\le1$ 内围成区域的面积，说明上下边界如何判断。
\par\textit{改编训练。}\par\hyperref[ans:c21q1]{解答见第~\pageref{ans:c21q1} 页。}
\end{example}
\begin{example}[两条幂曲线之间的面积]\label{ex:c21q2}
求曲线 $y=x$ 与 $y=x^{3}$ 在 $0\le x\le1$ 内围成区域的面积，说明上下边界如何判断。
\par\textit{改编训练。}\par\hyperref[ans:c21q2]{解答见第~\pageref{ans:c21q2} 页。}
\end{example}
\begin{example}[两条幂曲线之间的面积]\label{ex:c21q3}
求曲线 $y=x$ 与 $y=x^{4}$ 在 $0\le x\le1$ 内围成区域的面积，说明上下边界如何判断。
\par\textit{改编训练。}\par\hyperref[ans:c21q3]{解答见第~\pageref{ans:c21q3} 页。}
\end{example}
\begin{example}[两条幂曲线之间的面积]\label{ex:c21q4}
求曲线 $y=x$ 与 $y=x^{5}$ 在 $0\le x\le1$ 内围成区域的面积，说明上下边界如何判断。
\par\textit{改编训练。}\par\hyperref[ans:c21q4]{解答见第~\pageref{ans:c21q4} 页。}
\end{example}
\begin{example}[两条幂曲线之间的面积]\label{ex:c21q5}
求曲线 $y=x$ 与 $y=x^{6}$ 在 $0\le x\le1$ 内围成区域的面积，说明上下边界如何判断。
\par\textit{改编训练。}\par\hyperref[ans:c21q5]{解答见第~\pageref{ans:c21q5} 页。}
\end{example}
\begin{example}[两条幂曲线之间的面积]\label{ex:c21q6}
求曲线 $y=x$ 与 $y=x^{7}$ 在 $0\le x\le1$ 内围成区域的面积，说明上下边界如何判断。
\par\textit{改编训练。}\par\hyperref[ans:c21q6]{解答见第~\pageref{ans:c21q6} 页。}
\end{example}
\begin{example}[同一区域绕不同轴旋转]\label{ex:c21q7}
将区域 $0\le x\le1$、$0\le y\le x^{1}$ 分别绕 $x$ 轴与 $y$ 轴旋转，求两个体积。
\par\textit{改编训练。}\par\hyperref[ans:c21q7]{解答见第~\pageref{ans:c21q7} 页。}
\end{example}
\begin{example}[同一区域绕不同轴旋转]\label{ex:c21q8}
将区域 $0\le x\le1$、$0\le y\le x^{2}$ 分别绕 $x$ 轴与 $y$ 轴旋转，求两个体积。
\par\textit{改编训练。}\par\hyperref[ans:c21q8]{解答见第~\pageref{ans:c21q8} 页。}
\end{example}
\begin{example}[同一区域绕不同轴旋转]\label{ex:c21q9}
将区域 $0\le x\le1$、$0\le y\le x^{3}$ 分别绕 $x$ 轴与 $y$ 轴旋转，求两个体积。
\par\textit{改编训练。}\par\hyperref[ans:c21q9]{解答见第~\pageref{ans:c21q9} 页。}
\end{example}
\begin{example}[同一区域绕不同轴旋转]\label{ex:c21q10}
将区域 $0\le x\le1$、$0\le y\le x^{4}$ 分别绕 $x$ 轴与 $y$ 轴旋转，求两个体积。
\par\textit{改编训练。}\par\hyperref[ans:c21q10]{解答见第~\pageref{ans:c21q10} 页。}
\end{example}
\begin{example}[同一区域绕不同轴旋转]\label{ex:c21q11}
将区域 $0\le x\le1$、$0\le y\le x^{5}$ 分别绕 $x$ 轴与 $y$ 轴旋转，求两个体积。
\par\textit{改编训练。}\par\hyperref[ans:c21q11]{解答见第~\pageref{ans:c21q11} 页。}
\end{example}
\begin{example}[同一区域绕不同轴旋转]\label{ex:c21q12}
将区域 $0\le x\le1$、$0\le y\le x^{6}$ 分别绕 $x$ 轴与 $y$ 轴旋转，求两个体积。
\par\textit{改编训练。}\par\hyperref[ans:c21q12]{解答见第~\pageref{ans:c21q12} 页。}
\end{example}
\begin{example}[曲边梯形的形心]\label{ex:c21q13}
密度恒为1的薄片占据 $0\le x\le1$、$0\le y\le x^{1}$。求面积及形心 $(\bar x,\bar y)$。
\par\textit{改编训练。}\par\hyperref[ans:c21q13]{解答见第~\pageref{ans:c21q13} 页。}
\end{example}
\begin{example}[曲边梯形的形心]\label{ex:c21q14}
密度恒为1的薄片占据 $0\le x\le1$、$0\le y\le x^{2}$。求面积及形心 $(\bar x,\bar y)$。
\par\textit{改编训练。}\par\hyperref[ans:c21q14]{解答见第~\pageref{ans:c21q14} 页。}
\end{example}
\begin{example}[曲边梯形的形心]\label{ex:c21q15}
密度恒为1的薄片占据 $0\le x\le1$、$0\le y\le x^{3}$。求面积及形心 $(\bar x,\bar y)$。
\par\textit{改编训练。}\par\hyperref[ans:c21q15]{解答见第~\pageref{ans:c21q15} 页。}
\end{example}
\begin{example}[曲边梯形的形心]\label{ex:c21q16}
密度恒为1的薄片占据 $0\le x\le1$、$0\le y\le x^{4}$。求面积及形心 $(\bar x,\bar y)$。
\par\textit{改编训练。}\par\hyperref[ans:c21q16]{解答见第~\pageref{ans:c21q16} 页。}
\end{example}
\begin{example}[曲边梯形的形心]\label{ex:c21q17}
密度恒为1的薄片占据 $0\le x\le1$、$0\le y\le x^{5}$。求面积及形心 $(\bar x,\bar y)$。
\par\textit{改编训练。}\par\hyperref[ans:c21q17]{解答见第~\pageref{ans:c21q17} 页。}
\end{example}
\begin{example}[曲边梯形的形心]\label{ex:c21q18}
密度恒为1的薄片占据 $0\le x\le1$、$0\le y\le x^{6}$。求面积及形心 $(\bar x,\bar y)$。
\par\textit{改编训练。}\par\hyperref[ans:c21q18]{解答见第~\pageref{ans:c21q18} 页。}
\end{example}
\begin{example}[变力做功]\label{ex:c21q19}
弹簧满足 Hooke 定律 $F(x)=2x$ 牛顿，其中 $x$ 以米计。把弹簧从自然长度拉长到 $1$ 米需要多少功？若从 $1$ 米再拉到 $2$ 米，需要多少额外功？
\par\textit{改编训练。}\par\hyperref[ans:c21q19]{解答见第~\pageref{ans:c21q19} 页。}
\end{example}
\begin{example}[变力做功]\label{ex:c21q20}
弹簧满足 Hooke 定律 $F(x)=3x$ 牛顿，其中 $x$ 以米计。把弹簧从自然长度拉长到 $2$ 米需要多少功？若从 $2$ 米再拉到 $3$ 米，需要多少额外功？
\par\textit{改编训练。}\par\hyperref[ans:c21q20]{解答见第~\pageref{ans:c21q20} 页。}
\end{example}
\begin{example}[变力做功]\label{ex:c21q21}
弹簧满足 Hooke 定律 $F(x)=4x$ 牛顿，其中 $x$ 以米计。把弹簧从自然长度拉长到 $3$ 米需要多少功？若从 $3$ 米再拉到 $4$ 米，需要多少额外功？
\par\textit{改编训练。}\par\hyperref[ans:c21q21]{解答见第~\pageref{ans:c21q21} 页。}
\end{example}
\begin{example}[变力做功]\label{ex:c21q22}
弹簧满足 Hooke 定律 $F(x)=5x$ 牛顿，其中 $x$ 以米计。把弹簧从自然长度拉长到 $2$ 米需要多少功？若从 $2$ 米再拉到 $3$ 米，需要多少额外功？
\par\textit{改编训练。}\par\hyperref[ans:c21q22]{解答见第~\pageref{ans:c21q22} 页。}
\end{example}
\begin{example}[变力做功]\label{ex:c21q23}
弹簧满足 Hooke 定律 $F(x)=6x$ 牛顿，其中 $x$ 以米计。把弹簧从自然长度拉长到 $4$ 米需要多少功？若从 $4$ 米再拉到 $5$ 米，需要多少额外功？
\par\textit{改编训练。}\par\hyperref[ans:c21q23]{解答见第~\pageref{ans:c21q23} 页。}
\end{example}
\begin{example}[变力做功]\label{ex:c21q24}
弹簧满足 Hooke 定律 $F(x)=8x$ 牛顿，其中 $x$ 以米计。把弹簧从自然长度拉长到 $3$ 米需要多少功？若从 $3$ 米再拉到 $4$ 米，需要多少额外功？
\par\textit{改编训练。}\par\hyperref[ans:c21q24]{解答见第~\pageref{ans:c21q24} 页。}
\end{example}
\subsection{解答}
\begin{solve}[习题~\ref{ex:c21q1}：两条幂曲线之间的面积]\phantomsection\label{ans:c21q1}
在 $0\le x\le1$ 上有 $x^{2}\le x$，两曲线交于 $x=0,1$。竖直薄片高度为 $x-x^{2}$，故面积为
\[
A=\int_0^1(x-x^{2})dx=\frac12-\frac1{3}=\frac{1}{6}.
\]
积分结果为正，符合面积非负的要求。若改变定义区间，上下边界可能交换，必须重新比较，不能仅按曲线写出的先后顺序相减。本题先确定交点和大小关系，再积分，不需要先画一个精确比例的图。
\end{solve}

\begin{solve}[习题~\ref{ex:c21q2}：两条幂曲线之间的面积]\phantomsection\label{ans:c21q2}
在 $0\le x\le1$ 上有 $x^{3}\le x$，两曲线交于 $x=0,1$。竖直薄片高度为 $x-x^{3}$，故面积为
\[
A=\int_0^1(x-x^{3})dx=\frac12-\frac1{4}=\frac{1}{4}.
\]
积分结果为正，符合面积非负的要求。若改变定义区间，上下边界可能交换，必须重新比较，不能仅按曲线写出的先后顺序相减。本题先确定交点和大小关系，再积分，不需要先画一个精确比例的图。
\end{solve}

\begin{solve}[习题~\ref{ex:c21q3}：两条幂曲线之间的面积]\phantomsection\label{ans:c21q3}
在 $0\le x\le1$ 上有 $x^{4}\le x$，两曲线交于 $x=0,1$。竖直薄片高度为 $x-x^{4}$，故面积为
\[
A=\int_0^1(x-x^{4})dx=\frac12-\frac1{5}=\frac{3}{10}.
\]
积分结果为正，符合面积非负的要求。若改变定义区间，上下边界可能交换，必须重新比较，不能仅按曲线写出的先后顺序相减。本题先确定交点和大小关系，再积分，不需要先画一个精确比例的图。
\end{solve}

\begin{solve}[习题~\ref{ex:c21q4}：两条幂曲线之间的面积]\phantomsection\label{ans:c21q4}
在 $0\le x\le1$ 上有 $x^{5}\le x$，两曲线交于 $x=0,1$。竖直薄片高度为 $x-x^{5}$，故面积为
\[
A=\int_0^1(x-x^{5})dx=\frac12-\frac1{6}=\frac{1}{3}.
\]
积分结果为正，符合面积非负的要求。若改变定义区间，上下边界可能交换，必须重新比较，不能仅按曲线写出的先后顺序相减。本题先确定交点和大小关系，再积分，不需要先画一个精确比例的图。
\end{solve}

\begin{solve}[习题~\ref{ex:c21q5}：两条幂曲线之间的面积]\phantomsection\label{ans:c21q5}
在 $0\le x\le1$ 上有 $x^{6}\le x$，两曲线交于 $x=0,1$。竖直薄片高度为 $x-x^{6}$，故面积为
\[
A=\int_0^1(x-x^{6})dx=\frac12-\frac1{7}=\frac{5}{14}.
\]
积分结果为正，符合面积非负的要求。若改变定义区间，上下边界可能交换，必须重新比较，不能仅按曲线写出的先后顺序相减。本题先确定交点和大小关系，再积分，不需要先画一个精确比例的图。
\end{solve}

\begin{solve}[习题~\ref{ex:c21q6}：两条幂曲线之间的面积]\phantomsection\label{ans:c21q6}
在 $0\le x\le1$ 上有 $x^{7}\le x$，两曲线交于 $x=0,1$。竖直薄片高度为 $x-x^{7}$，故面积为
\[
A=\int_0^1(x-x^{7})dx=\frac12-\frac1{8}=\frac{3}{8}.
\]
积分结果为正，符合面积非负的要求。若改变定义区间，上下边界可能交换，必须重新比较，不能仅按曲线写出的先后顺序相减。本题先确定交点和大小关系，再积分，不需要先画一个精确比例的图。
\end{solve}

\begin{solve}[习题~\ref{ex:c21q7}：同一区域绕不同轴旋转]\phantomsection\label{ans:c21q7}
绕 $x$ 轴时，截面是半径为 $x^m$ 的圆盘。绕 $y$ 轴时，使用半径为 $x$、高度为 $x^m$ 的圆柱壳最方便。因此
\[
V_x=\pi\int_0^1x^{2}dx=\frac{1}{3}\pi,
\]

\[
V_y=2\pi\int_0^1x\,x^{1}dx=\frac{2}{3}\pi.
\]
两种体积一般不同，原因是旋转轴不同导致微元的几何权重不同。圆盘法的半径需要平方，柱壳法的半径只作为周长中的一个因子。
\end{solve}

\begin{solve}[习题~\ref{ex:c21q8}：同一区域绕不同轴旋转]\phantomsection\label{ans:c21q8}
绕 $x$ 轴时，截面是半径为 $x^m$ 的圆盘。绕 $y$ 轴时，使用半径为 $x$、高度为 $x^m$ 的圆柱壳最方便。因此
\[
V_x=\pi\int_0^1x^{4}dx=\frac{1}{5}\pi,
\]

\[
V_y=2\pi\int_0^1x\,x^{2}dx=\frac{1}{2}\pi.
\]
两种体积一般不同，原因是旋转轴不同导致微元的几何权重不同。圆盘法的半径需要平方，柱壳法的半径只作为周长中的一个因子。
\end{solve}

\begin{solve}[习题~\ref{ex:c21q9}：同一区域绕不同轴旋转]\phantomsection\label{ans:c21q9}
绕 $x$ 轴时，截面是半径为 $x^m$ 的圆盘。绕 $y$ 轴时，使用半径为 $x$、高度为 $x^m$ 的圆柱壳最方便。因此
\[
V_x=\pi\int_0^1x^{6}dx=\frac{1}{7}\pi,
\]

\[
V_y=2\pi\int_0^1x\,x^{3}dx=\frac{2}{5}\pi.
\]
两种体积一般不同，原因是旋转轴不同导致微元的几何权重不同。圆盘法的半径需要平方，柱壳法的半径只作为周长中的一个因子。
\end{solve}

\begin{solve}[习题~\ref{ex:c21q10}：同一区域绕不同轴旋转]\phantomsection\label{ans:c21q10}
绕 $x$ 轴时，截面是半径为 $x^m$ 的圆盘。绕 $y$ 轴时，使用半径为 $x$、高度为 $x^m$ 的圆柱壳最方便。因此
\[
V_x=\pi\int_0^1x^{8}dx=\frac{1}{9}\pi,
\]

\[
V_y=2\pi\int_0^1x\,x^{4}dx=\frac{1}{3}\pi.
\]
两种体积一般不同，原因是旋转轴不同导致微元的几何权重不同。圆盘法的半径需要平方，柱壳法的半径只作为周长中的一个因子。
\end{solve}

\begin{solve}[习题~\ref{ex:c21q11}：同一区域绕不同轴旋转]\phantomsection\label{ans:c21q11}
绕 $x$ 轴时，截面是半径为 $x^m$ 的圆盘。绕 $y$ 轴时，使用半径为 $x$、高度为 $x^m$ 的圆柱壳最方便。因此
\[
V_x=\pi\int_0^1x^{10}dx=\frac{1}{11}\pi,
\]

\[
V_y=2\pi\int_0^1x\,x^{5}dx=\frac{2}{7}\pi.
\]
两种体积一般不同，原因是旋转轴不同导致微元的几何权重不同。圆盘法的半径需要平方，柱壳法的半径只作为周长中的一个因子。
\end{solve}

\begin{solve}[习题~\ref{ex:c21q12}：同一区域绕不同轴旋转]\phantomsection\label{ans:c21q12}
绕 $x$ 轴时，截面是半径为 $x^m$ 的圆盘。绕 $y$ 轴时，使用半径为 $x$、高度为 $x^m$ 的圆柱壳最方便。因此
\[
V_x=\pi\int_0^1x^{12}dx=\frac{1}{13}\pi,
\]

\[
V_y=2\pi\int_0^1x\,x^{6}dx=\frac{1}{4}\pi.
\]
两种体积一般不同，原因是旋转轴不同导致微元的几何权重不同。圆盘法的半径需要平方，柱壳法的半径只作为周长中的一个因子。
\end{solve}

\begin{solve}[习题~\ref{ex:c21q13}：曲边梯形的形心]\phantomsection\label{ans:c21q13}
总质量等于面积 $A=\int_0^1x^{1}dx=1/2$。对坐标轴的一阶矩分别为
\[
M_y=\int_0^1x\,x^{1}dx=\frac1{3},\qquad M_x=\int_0^1\frac{(x^{1})^2}2dx=\frac1{6}.
\]
于是
\[
\bar x=\frac{M_y}A=\frac{2}{3},\qquad\bar y=\frac{M_x}A=\frac{1}{3}.
\]
注意 $M_y$ 决定横坐标、$M_x$ 决定纵坐标；矩的下标表示取矩的轴，不是结果的坐标。形心位于区域的凸包内，与所得坐标的范围一致。
\end{solve}

\begin{solve}[习题~\ref{ex:c21q14}：曲边梯形的形心]\phantomsection\label{ans:c21q14}
总质量等于面积 $A=\int_0^1x^{2}dx=1/3$。对坐标轴的一阶矩分别为
\[
M_y=\int_0^1x\,x^{2}dx=\frac1{4},\qquad M_x=\int_0^1\frac{(x^{2})^2}2dx=\frac1{10}.
\]
于是
\[
\bar x=\frac{M_y}A=\frac{3}{4},\qquad\bar y=\frac{M_x}A=\frac{3}{10}.
\]
注意 $M_y$ 决定横坐标、$M_x$ 决定纵坐标；矩的下标表示取矩的轴，不是结果的坐标。形心位于区域的凸包内，与所得坐标的范围一致。
\end{solve}

\begin{solve}[习题~\ref{ex:c21q15}：曲边梯形的形心]\phantomsection\label{ans:c21q15}
总质量等于面积 $A=\int_0^1x^{3}dx=1/4$。对坐标轴的一阶矩分别为
\[
M_y=\int_0^1x\,x^{3}dx=\frac1{5},\qquad M_x=\int_0^1\frac{(x^{3})^2}2dx=\frac1{14}.
\]
于是
\[
\bar x=\frac{M_y}A=\frac{4}{5},\qquad\bar y=\frac{M_x}A=\frac{2}{7}.
\]
注意 $M_y$ 决定横坐标、$M_x$ 决定纵坐标；矩的下标表示取矩的轴，不是结果的坐标。形心位于区域的凸包内，与所得坐标的范围一致。
\end{solve}

\begin{solve}[习题~\ref{ex:c21q16}：曲边梯形的形心]\phantomsection\label{ans:c21q16}
总质量等于面积 $A=\int_0^1x^{4}dx=1/5$。对坐标轴的一阶矩分别为
\[
M_y=\int_0^1x\,x^{4}dx=\frac1{6},\qquad M_x=\int_0^1\frac{(x^{4})^2}2dx=\frac1{18}.
\]
于是
\[
\bar x=\frac{M_y}A=\frac{5}{6},\qquad\bar y=\frac{M_x}A=\frac{5}{18}.
\]
注意 $M_y$ 决定横坐标、$M_x$ 决定纵坐标；矩的下标表示取矩的轴，不是结果的坐标。形心位于区域的凸包内，与所得坐标的范围一致。
\end{solve}

\begin{solve}[习题~\ref{ex:c21q17}：曲边梯形的形心]\phantomsection\label{ans:c21q17}
总质量等于面积 $A=\int_0^1x^{5}dx=1/6$。对坐标轴的一阶矩分别为
\[
M_y=\int_0^1x\,x^{5}dx=\frac1{7},\qquad M_x=\int_0^1\frac{(x^{5})^2}2dx=\frac1{22}.
\]
于是
\[
\bar x=\frac{M_y}A=\frac{6}{7},\qquad\bar y=\frac{M_x}A=\frac{3}{11}.
\]
注意 $M_y$ 决定横坐标、$M_x$ 决定纵坐标；矩的下标表示取矩的轴，不是结果的坐标。形心位于区域的凸包内，与所得坐标的范围一致。
\end{solve}

\begin{solve}[习题~\ref{ex:c21q18}：曲边梯形的形心]\phantomsection\label{ans:c21q18}
总质量等于面积 $A=\int_0^1x^{6}dx=1/7$。对坐标轴的一阶矩分别为
\[
M_y=\int_0^1x\,x^{6}dx=\frac1{8},\qquad M_x=\int_0^1\frac{(x^{6})^2}2dx=\frac1{26}.
\]
于是
\[
\bar x=\frac{M_y}A=\frac{7}{8},\qquad\bar y=\frac{M_x}A=\frac{7}{26}.
\]
注意 $M_y$ 决定横坐标、$M_x$ 决定纵坐标；矩的下标表示取矩的轴，不是结果的坐标。形心位于区域的凸包内，与所得坐标的范围一致。
\end{solve}

\begin{solve}[习题~\ref{ex:c21q19}：变力做功]\phantomsection\label{ans:c21q19}
弹力随位移变化，不能使用末端力直接乘总位移。把过程划分成微小位移，功的微元为 $dW=F(x)dx$。因此
\[
W_1=\int_0^{1}2x\,dx=1\ \mathrm{J},
\]

\[
W_2=\int_{1}^{2}2x\,dx=3\ \mathrm{J}.
\]
后半段即使位移较短，也可能需要较大的功，因为弹力随伸长增大。结果的单位是牛顿乘米，即焦耳；模型适用范围以题设的线性弹性假设为准。
\end{solve}

\begin{solve}[习题~\ref{ex:c21q20}：变力做功]\phantomsection\label{ans:c21q20}
弹力随位移变化，不能使用末端力直接乘总位移。把过程划分成微小位移，功的微元为 $dW=F(x)dx$。因此
\[
W_1=\int_0^{2}3x\,dx=6\ \mathrm{J},
\]

\[
W_2=\int_{2}^{3}3x\,dx=\frac{15}{2}\ \mathrm{J}.
\]
后半段即使位移较短，也可能需要较大的功，因为弹力随伸长增大。结果的单位是牛顿乘米，即焦耳；模型适用范围以题设的线性弹性假设为准。
\end{solve}

\begin{solve}[习题~\ref{ex:c21q21}：变力做功]\phantomsection\label{ans:c21q21}
弹力随位移变化，不能使用末端力直接乘总位移。把过程划分成微小位移，功的微元为 $dW=F(x)dx$。因此
\[
W_1=\int_0^{3}4x\,dx=18\ \mathrm{J},
\]

\[
W_2=\int_{3}^{4}4x\,dx=14\ \mathrm{J}.
\]
后半段即使位移较短，也可能需要较大的功，因为弹力随伸长增大。结果的单位是牛顿乘米，即焦耳；模型适用范围以题设的线性弹性假设为准。
\end{solve}

\begin{solve}[习题~\ref{ex:c21q22}：变力做功]\phantomsection\label{ans:c21q22}
弹力随位移变化，不能使用末端力直接乘总位移。把过程划分成微小位移，功的微元为 $dW=F(x)dx$。因此
\[
W_1=\int_0^{2}5x\,dx=10\ \mathrm{J},
\]

\[
W_2=\int_{2}^{3}5x\,dx=\frac{25}{2}\ \mathrm{J}.
\]
后半段即使位移较短，也可能需要较大的功，因为弹力随伸长增大。结果的单位是牛顿乘米，即焦耳；模型适用范围以题设的线性弹性假设为准。
\end{solve}

\begin{solve}[习题~\ref{ex:c21q23}：变力做功]\phantomsection\label{ans:c21q23}
弹力随位移变化，不能使用末端力直接乘总位移。把过程划分成微小位移，功的微元为 $dW=F(x)dx$。因此
\[
W_1=\int_0^{4}6x\,dx=48\ \mathrm{J},
\]

\[
W_2=\int_{4}^{5}6x\,dx=27\ \mathrm{J}.
\]
后半段即使位移较短，也可能需要较大的功，因为弹力随伸长增大。结果的单位是牛顿乘米，即焦耳；模型适用范围以题设的线性弹性假设为准。
\end{solve}

\begin{solve}[习题~\ref{ex:c21q24}：变力做功]\phantomsection\label{ans:c21q24}
弹力随位移变化，不能使用末端力直接乘总位移。把过程划分成微小位移，功的微元为 $dW=F(x)dx$。因此
\[
W_1=\int_0^{3}8x\,dx=36\ \mathrm{J},
\]

\[
W_2=\int_{3}^{4}8x\,dx=28\ \mathrm{J}.
\]
后半段即使位移较短，也可能需要较大的功，因为弹力随伸长增大。结果的单位是牛顿乘米，即焦耳；模型适用范围以题设的线性弹性假设为准。
\end{solve}

\clearpage
\section{常微分方程的初步训练}\label{sec:ch22}
微分方程的解是函数，而不是一个常数。分离变量时若除以含未知函数的因子，必须补查被除掉的常值解。初值条件用于选定特解，并帮助确定实际使用的区间。

\textbf{资源对应：}CLP-2 \S~2.4；线性方程与二阶方程为补充训练。本章共 24 道编号习题。
\subsection{习题}
\begin{example}[指数增长初值问题]\label{ex:c22q1}
解 $y^{\prime}=1y$、$y(0)=2$，并求达到初值两倍所需的时间。
\par\textit{改编训练。}\par\hyperref[ans:c22q1]{解答见第~\pageref{ans:c22q1} 页。}
\end{example}
\begin{example}[指数增长初值问题]\label{ex:c22q2}
解 $y^{\prime}=2y$、$y(0)=1$，并求达到初值两倍所需的时间。
\par\textit{改编训练。}\par\hyperref[ans:c22q2]{解答见第~\pageref{ans:c22q2} 页。}
\end{example}
\begin{example}[指数增长初值问题]\label{ex:c22q3}
解 $y^{\prime}=3y$、$y(0)=2$，并求达到初值两倍所需的时间。
\par\textit{改编训练。}\par\hyperref[ans:c22q3]{解答见第~\pageref{ans:c22q3} 页。}
\end{example}
\begin{example}[指数增长初值问题]\label{ex:c22q4}
解 $y^{\prime}=2y$、$y(0)=3$，并求达到初值两倍所需的时间。
\par\textit{改编训练。}\par\hyperref[ans:c22q4]{解答见第~\pageref{ans:c22q4} 页。}
\end{example}
\begin{example}[指数增长初值问题]\label{ex:c22q5}
解 $y^{\prime}=4y$、$y(0)=1$，并求达到初值两倍所需的时间。
\par\textit{改编训练。}\par\hyperref[ans:c22q5]{解答见第~\pageref{ans:c22q5} 页。}
\end{example}
\begin{example}[指数增长初值问题]\label{ex:c22q6}
解 $y^{\prime}=3y$、$y(0)=4$，并求达到初值两倍所需的时间。
\par\textit{改编训练。}\par\hyperref[ans:c22q6]{解答见第~\pageref{ans:c22q6} 页。}
\end{example}
\begin{example}[分离变量与有限时刻发散]\label{ex:c22q7}
解 $y^{\prime}=1y^2$、$y(0)=1$，并给出包含 $0$ 的最大解区间。
\par\textit{改编训练。}\par\hyperref[ans:c22q7]{解答见第~\pageref{ans:c22q7} 页。}
\end{example}
\begin{example}[分离变量与有限时刻发散]\label{ex:c22q8}
解 $y^{\prime}=2y^2$、$y(0)=1$，并给出包含 $0$ 的最大解区间。
\par\textit{改编训练。}\par\hyperref[ans:c22q8]{解答见第~\pageref{ans:c22q8} 页。}
\end{example}
\begin{example}[分离变量与有限时刻发散]\label{ex:c22q9}
解 $y^{\prime}=1y^2$、$y(0)=2$，并给出包含 $0$ 的最大解区间。
\par\textit{改编训练。}\par\hyperref[ans:c22q9]{解答见第~\pageref{ans:c22q9} 页。}
\end{example}
\begin{example}[分离变量与有限时刻发散]\label{ex:c22q10}
解 $y^{\prime}=3y^2$、$y(0)=2$，并给出包含 $0$ 的最大解区间。
\par\textit{改编训练。}\par\hyperref[ans:c22q10]{解答见第~\pageref{ans:c22q10} 页。}
\end{example}
\begin{example}[分离变量与有限时刻发散]\label{ex:c22q11}
解 $y^{\prime}=2y^2$、$y(0)=3$，并给出包含 $0$ 的最大解区间。
\par\textit{改编训练。}\par\hyperref[ans:c22q11]{解答见第~\pageref{ans:c22q11} 页。}
\end{example}
\begin{example}[分离变量与有限时刻发散]\label{ex:c22q12}
解 $y^{\prime}=4y^2$、$y(0)=1$，并给出包含 $0$ 的最大解区间。
\par\textit{改编训练。}\par\hyperref[ans:c22q12]{解答见第~\pageref{ans:c22q12} 页。}
\end{example}
\begin{example}[积分因子解线性方程]\label{ex:c22q13}
解 $y^{\prime}+1y=2$、$y(0)=0$，并求 $x\to\infty$ 时的极限。
\par\textit{改编训练。}\par\hyperref[ans:c22q13]{解答见第~\pageref{ans:c22q13} 页。}
\end{example}
\begin{example}[积分因子解线性方程]\label{ex:c22q14}
解 $y^{\prime}+2y=1$、$y(0)=0$，并求 $x\to\infty$ 时的极限。
\par\textit{改编训练。}\par\hyperref[ans:c22q14]{解答见第~\pageref{ans:c22q14} 页。}
\end{example}
\begin{example}[积分因子解线性方程]\label{ex:c22q15}
解 $y^{\prime}+3y=2$、$y(0)=0$，并求 $x\to\infty$ 时的极限。
\par\textit{改编训练。}\par\hyperref[ans:c22q15]{解答见第~\pageref{ans:c22q15} 页。}
\end{example}
\begin{example}[积分因子解线性方程]\label{ex:c22q16}
解 $y^{\prime}+2y=3$、$y(0)=0$，并求 $x\to\infty$ 时的极限。
\par\textit{改编训练。}\par\hyperref[ans:c22q16]{解答见第~\pageref{ans:c22q16} 页。}
\end{example}
\begin{example}[积分因子解线性方程]\label{ex:c22q17}
解 $y^{\prime}+4y=1$、$y(0)=0$，并求 $x\to\infty$ 时的极限。
\par\textit{改编训练。}\par\hyperref[ans:c22q17]{解答见第~\pageref{ans:c22q17} 页。}
\end{example}
\begin{example}[积分因子解线性方程]\label{ex:c22q18}
解 $y^{\prime}+3y=5$、$y(0)=0$，并求 $x\to\infty$ 时的极限。
\par\textit{改编训练。}\par\hyperref[ans:c22q18]{解答见第~\pageref{ans:c22q18} 页。}
\end{example}
\begin{example}[简谐振动初值问题]\label{ex:c22q19}
解 $y^{\prime\prime}+1y=0$、$y(0)=2$、$y^{\prime}(0)=0$，并给出周期。
\par\textit{改编训练。}\par\hyperref[ans:c22q19]{解答见第~\pageref{ans:c22q19} 页。}
\end{example}
\begin{example}[简谐振动初值问题]\label{ex:c22q20}
解 $y^{\prime\prime}+4y=0$、$y(0)=1$、$y^{\prime}(0)=0$，并给出周期。
\par\textit{改编训练。}\par\hyperref[ans:c22q20]{解答见第~\pageref{ans:c22q20} 页。}
\end{example}
\begin{example}[简谐振动初值问题]\label{ex:c22q21}
解 $y^{\prime\prime}+9y=0$、$y(0)=2$、$y^{\prime}(0)=0$，并给出周期。
\par\textit{改编训练。}\par\hyperref[ans:c22q21]{解答见第~\pageref{ans:c22q21} 页。}
\end{example}
\begin{example}[简谐振动初值问题]\label{ex:c22q22}
解 $y^{\prime\prime}+4y=0$、$y(0)=3$、$y^{\prime}(0)=0$，并给出周期。
\par\textit{改编训练。}\par\hyperref[ans:c22q22]{解答见第~\pageref{ans:c22q22} 页。}
\end{example}
\begin{example}[简谐振动初值问题]\label{ex:c22q23}
解 $y^{\prime\prime}+16y=0$、$y(0)=1$、$y^{\prime}(0)=0$，并给出周期。
\par\textit{改编训练。}\par\hyperref[ans:c22q23]{解答见第~\pageref{ans:c22q23} 页。}
\end{example}
\begin{example}[简谐振动初值问题]\label{ex:c22q24}
解 $y^{\prime\prime}+9y=0$、$y(0)=4$、$y^{\prime}(0)=0$，并给出周期。
\par\textit{改编训练。}\par\hyperref[ans:c22q24]{解答见第~\pageref{ans:c22q24} 页。}
\end{example}
\subsection{解答}
\begin{solve}[习题~\ref{ex:c22q1}：指数增长初值问题]\phantomsection\label{ans:c22q1}
因为初值为正，局部可分离变量：$dy/y=1dx$。积分得 $\log|y|=1x+C$，结合初值得到
\[
y(x)=2e^{1x}.
\]
该解在全实轴上有定义，直接求导与代入初值都可验证。令 $2e^{1x}=4$，得到倍增时间
\[
x=\frac{\log2}{1}.
\]
倍增时间与初始规模无关，是常系数指数增长模型的特征；它依赖增长率而非当前数量。
\end{solve}

\begin{solve}[习题~\ref{ex:c22q2}：指数增长初值问题]\phantomsection\label{ans:c22q2}
因为初值为正，局部可分离变量：$dy/y=2dx$。积分得 $\log|y|=2x+C$，结合初值得到
\[
y(x)=1e^{2x}.
\]
该解在全实轴上有定义，直接求导与代入初值都可验证。令 $1e^{2x}=2$，得到倍增时间
\[
x=\frac{\log2}{2}.
\]
倍增时间与初始规模无关，是常系数指数增长模型的特征；它依赖增长率而非当前数量。
\end{solve}

\begin{solve}[习题~\ref{ex:c22q3}：指数增长初值问题]\phantomsection\label{ans:c22q3}
因为初值为正，局部可分离变量：$dy/y=3dx$。积分得 $\log|y|=3x+C$，结合初值得到
\[
y(x)=2e^{3x}.
\]
该解在全实轴上有定义，直接求导与代入初值都可验证。令 $2e^{3x}=4$，得到倍增时间
\[
x=\frac{\log2}{3}.
\]
倍增时间与初始规模无关，是常系数指数增长模型的特征；它依赖增长率而非当前数量。
\end{solve}

\begin{solve}[习题~\ref{ex:c22q4}：指数增长初值问题]\phantomsection\label{ans:c22q4}
因为初值为正，局部可分离变量：$dy/y=2dx$。积分得 $\log|y|=2x+C$，结合初值得到
\[
y(x)=3e^{2x}.
\]
该解在全实轴上有定义，直接求导与代入初值都可验证。令 $3e^{2x}=6$，得到倍增时间
\[
x=\frac{\log2}{2}.
\]
倍增时间与初始规模无关，是常系数指数增长模型的特征；它依赖增长率而非当前数量。
\end{solve}

\begin{solve}[习题~\ref{ex:c22q5}：指数增长初值问题]\phantomsection\label{ans:c22q5}
因为初值为正，局部可分离变量：$dy/y=4dx$。积分得 $\log|y|=4x+C$，结合初值得到
\[
y(x)=1e^{4x}.
\]
该解在全实轴上有定义，直接求导与代入初值都可验证。令 $1e^{4x}=2$，得到倍增时间
\[
x=\frac{\log2}{4}.
\]
倍增时间与初始规模无关，是常系数指数增长模型的特征；它依赖增长率而非当前数量。
\end{solve}

\begin{solve}[习题~\ref{ex:c22q6}：指数增长初值问题]\phantomsection\label{ans:c22q6}
因为初值为正，局部可分离变量：$dy/y=3dx$。积分得 $\log|y|=3x+C$，结合初值得到
\[
y(x)=4e^{3x}.
\]
该解在全实轴上有定义，直接求导与代入初值都可验证。令 $4e^{3x}=8$，得到倍增时间
\[
x=\frac{\log2}{3}.
\]
倍增时间与初始规模无关，是常系数指数增长模型的特征；它依赖增长率而非当前数量。
\end{solve}

\begin{solve}[习题~\ref{ex:c22q7}：分离变量与有限时刻发散]\phantomsection\label{ans:c22q7}
正初值对应的解不恒为零，可先分离变量，得到 $-1/y=1x+C$。由 $y(0)=1$ 可知 $C=-1/1$，因此
\[
y(x)=\frac{1}{1-1x}.
\]
分母在 $x=1/1$ 为零，所以包含初值点的最大区间为 $(-\infty,1/1)$。在右端点前，解趋于正无穷。零函数也是原方程的解，但不满足本题非零初值。求出一个代数表达式后仍须报告区间，不能把极点两侧直接当成同一个初值解。
\end{solve}

\begin{solve}[习题~\ref{ex:c22q8}：分离变量与有限时刻发散]\phantomsection\label{ans:c22q8}
正初值对应的解不恒为零，可先分离变量，得到 $-1/y=2x+C$。由 $y(0)=1$ 可知 $C=-1/1$，因此
\[
y(x)=\frac{1}{1-2x}.
\]
分母在 $x=1/2$ 为零，所以包含初值点的最大区间为 $(-\infty,1/2)$。在右端点前，解趋于正无穷。零函数也是原方程的解，但不满足本题非零初值。求出一个代数表达式后仍须报告区间，不能把极点两侧直接当成同一个初值解。
\end{solve}

\begin{solve}[习题~\ref{ex:c22q9}：分离变量与有限时刻发散]\phantomsection\label{ans:c22q9}
正初值对应的解不恒为零，可先分离变量，得到 $-1/y=1x+C$。由 $y(0)=2$ 可知 $C=-1/2$，因此
\[
y(x)=\frac{2}{1-2x}.
\]
分母在 $x=1/2$ 为零，所以包含初值点的最大区间为 $(-\infty,1/2)$。在右端点前，解趋于正无穷。零函数也是原方程的解，但不满足本题非零初值。求出一个代数表达式后仍须报告区间，不能把极点两侧直接当成同一个初值解。
\end{solve}

\begin{solve}[习题~\ref{ex:c22q10}：分离变量与有限时刻发散]\phantomsection\label{ans:c22q10}
正初值对应的解不恒为零，可先分离变量，得到 $-1/y=3x+C$。由 $y(0)=2$ 可知 $C=-1/2$，因此
\[
y(x)=\frac{2}{1-6x}.
\]
分母在 $x=1/6$ 为零，所以包含初值点的最大区间为 $(-\infty,1/6)$。在右端点前，解趋于正无穷。零函数也是原方程的解，但不满足本题非零初值。求出一个代数表达式后仍须报告区间，不能把极点两侧直接当成同一个初值解。
\end{solve}

\begin{solve}[习题~\ref{ex:c22q11}：分离变量与有限时刻发散]\phantomsection\label{ans:c22q11}
正初值对应的解不恒为零，可先分离变量，得到 $-1/y=2x+C$。由 $y(0)=3$ 可知 $C=-1/3$，因此
\[
y(x)=\frac{3}{1-6x}.
\]
分母在 $x=1/6$ 为零，所以包含初值点的最大区间为 $(-\infty,1/6)$。在右端点前，解趋于正无穷。零函数也是原方程的解，但不满足本题非零初值。求出一个代数表达式后仍须报告区间，不能把极点两侧直接当成同一个初值解。
\end{solve}

\begin{solve}[习题~\ref{ex:c22q12}：分离变量与有限时刻发散]\phantomsection\label{ans:c22q12}
正初值对应的解不恒为零，可先分离变量，得到 $-1/y=4x+C$。由 $y(0)=1$ 可知 $C=-1/1$，因此
\[
y(x)=\frac{1}{1-4x}.
\]
分母在 $x=1/4$ 为零，所以包含初值点的最大区间为 $(-\infty,1/4)$。在右端点前，解趋于正无穷。零函数也是原方程的解，但不满足本题非零初值。求出一个代数表达式后仍须报告区间，不能把极点两侧直接当成同一个初值解。
\end{solve}

\begin{solve}[习题~\ref{ex:c22q13}：积分因子解线性方程]\phantomsection\label{ans:c22q13}
积分因子为 $e^{1x}$，方程可写成
\[
\left(e^{1x}y\right)^{\prime}=2e^{1x}.
\]
积分得 $e^{1x}y=(2)e^{1x}+C$，初值给出 $C=-2$，于是
\[
y(x)=2(1-e^{-1x}),\qquad\lim_{x\to\infty}y(x)=2.
\]
极限等于令导数为零得到的平衡值；本题系数为正，指数扰动衰减，所以该平衡值具有吸引性。最后把解代回原方程可以同时检查符号和积分常数。
\end{solve}

\begin{solve}[习题~\ref{ex:c22q14}：积分因子解线性方程]\phantomsection\label{ans:c22q14}
积分因子为 $e^{2x}$，方程可写成
\[
\left(e^{2x}y\right)^{\prime}=1e^{2x}.
\]
积分得 $e^{2x}y=(\frac{1}{2})e^{2x}+C$，初值给出 $C=-\frac{1}{2}$，于是
\[
y(x)=\frac{1}{2}(1-e^{-2x}),\qquad\lim_{x\to\infty}y(x)=\frac{1}{2}.
\]
极限等于令导数为零得到的平衡值；本题系数为正，指数扰动衰减，所以该平衡值具有吸引性。最后把解代回原方程可以同时检查符号和积分常数。
\end{solve}

\begin{solve}[习题~\ref{ex:c22q15}：积分因子解线性方程]\phantomsection\label{ans:c22q15}
积分因子为 $e^{3x}$，方程可写成
\[
\left(e^{3x}y\right)^{\prime}=2e^{3x}.
\]
积分得 $e^{3x}y=(\frac{2}{3})e^{3x}+C$，初值给出 $C=-\frac{2}{3}$，于是
\[
y(x)=\frac{2}{3}(1-e^{-3x}),\qquad\lim_{x\to\infty}y(x)=\frac{2}{3}.
\]
极限等于令导数为零得到的平衡值；本题系数为正，指数扰动衰减，所以该平衡值具有吸引性。最后把解代回原方程可以同时检查符号和积分常数。
\end{solve}

\begin{solve}[习题~\ref{ex:c22q16}：积分因子解线性方程]\phantomsection\label{ans:c22q16}
积分因子为 $e^{2x}$，方程可写成
\[
\left(e^{2x}y\right)^{\prime}=3e^{2x}.
\]
积分得 $e^{2x}y=(\frac{3}{2})e^{2x}+C$，初值给出 $C=-\frac{3}{2}$，于是
\[
y(x)=\frac{3}{2}(1-e^{-2x}),\qquad\lim_{x\to\infty}y(x)=\frac{3}{2}.
\]
极限等于令导数为零得到的平衡值；本题系数为正，指数扰动衰减，所以该平衡值具有吸引性。最后把解代回原方程可以同时检查符号和积分常数。
\end{solve}

\begin{solve}[习题~\ref{ex:c22q17}：积分因子解线性方程]\phantomsection\label{ans:c22q17}
积分因子为 $e^{4x}$，方程可写成
\[
\left(e^{4x}y\right)^{\prime}=1e^{4x}.
\]
积分得 $e^{4x}y=(\frac{1}{4})e^{4x}+C$，初值给出 $C=-\frac{1}{4}$，于是
\[
y(x)=\frac{1}{4}(1-e^{-4x}),\qquad\lim_{x\to\infty}y(x)=\frac{1}{4}.
\]
极限等于令导数为零得到的平衡值；本题系数为正，指数扰动衰减，所以该平衡值具有吸引性。最后把解代回原方程可以同时检查符号和积分常数。
\end{solve}

\begin{solve}[习题~\ref{ex:c22q18}：积分因子解线性方程]\phantomsection\label{ans:c22q18}
积分因子为 $e^{3x}$，方程可写成
\[
\left(e^{3x}y\right)^{\prime}=5e^{3x}.
\]
积分得 $e^{3x}y=(\frac{5}{3})e^{3x}+C$，初值给出 $C=-\frac{5}{3}$，于是
\[
y(x)=\frac{5}{3}(1-e^{-3x}),\qquad\lim_{x\to\infty}y(x)=\frac{5}{3}.
\]
极限等于令导数为零得到的平衡值；本题系数为正，指数扰动衰减，所以该平衡值具有吸引性。最后把解代回原方程可以同时检查符号和积分常数。
\end{solve}

\begin{solve}[习题~\ref{ex:c22q19}：简谐振动初值问题]\phantomsection\label{ans:c22q19}
特征方程为 $r^2+1=0$，根为 $\pm 1i$。实通解为 $y=A\cos(1x)+B\sin(1x)$。第一初值给出 $A=2$，对通解求导并代入第二初值得到 $1B=0$，所以
\[
y(x)=2\cos(1x),\qquad T=\frac{2\pi}{1}.
\]
直接求二阶导数可以验证方程。振幅由初始位移确定，周期由方程系数确定；对于这里的线性模型，振幅并不影响周期。这一题作为高等数学中常系数二阶方程的补充。
\end{solve}

\begin{solve}[习题~\ref{ex:c22q20}：简谐振动初值问题]\phantomsection\label{ans:c22q20}
特征方程为 $r^2+4=0$，根为 $\pm 2i$。实通解为 $y=A\cos(2x)+B\sin(2x)$。第一初值给出 $A=1$，对通解求导并代入第二初值得到 $2B=0$，所以
\[
y(x)=1\cos(2x),\qquad T=\frac{2\pi}{2}.
\]
直接求二阶导数可以验证方程。振幅由初始位移确定，周期由方程系数确定；对于这里的线性模型，振幅并不影响周期。这一题作为高等数学中常系数二阶方程的补充。
\end{solve}

\begin{solve}[习题~\ref{ex:c22q21}：简谐振动初值问题]\phantomsection\label{ans:c22q21}
特征方程为 $r^2+9=0$，根为 $\pm 3i$。实通解为 $y=A\cos(3x)+B\sin(3x)$。第一初值给出 $A=2$，对通解求导并代入第二初值得到 $3B=0$，所以
\[
y(x)=2\cos(3x),\qquad T=\frac{2\pi}{3}.
\]
直接求二阶导数可以验证方程。振幅由初始位移确定，周期由方程系数确定；对于这里的线性模型，振幅并不影响周期。这一题作为高等数学中常系数二阶方程的补充。
\end{solve}

\begin{solve}[习题~\ref{ex:c22q22}：简谐振动初值问题]\phantomsection\label{ans:c22q22}
特征方程为 $r^2+4=0$，根为 $\pm 2i$。实通解为 $y=A\cos(2x)+B\sin(2x)$。第一初值给出 $A=3$，对通解求导并代入第二初值得到 $2B=0$，所以
\[
y(x)=3\cos(2x),\qquad T=\frac{2\pi}{2}.
\]
直接求二阶导数可以验证方程。振幅由初始位移确定，周期由方程系数确定；对于这里的线性模型，振幅并不影响周期。这一题作为高等数学中常系数二阶方程的补充。
\end{solve}

\begin{solve}[习题~\ref{ex:c22q23}：简谐振动初值问题]\phantomsection\label{ans:c22q23}
特征方程为 $r^2+16=0$，根为 $\pm 4i$。实通解为 $y=A\cos(4x)+B\sin(4x)$。第一初值给出 $A=1$，对通解求导并代入第二初值得到 $4B=0$，所以
\[
y(x)=1\cos(4x),\qquad T=\frac{2\pi}{4}.
\]
直接求二阶导数可以验证方程。振幅由初始位移确定，周期由方程系数确定；对于这里的线性模型，振幅并不影响周期。这一题作为高等数学中常系数二阶方程的补充。
\end{solve}

\begin{solve}[习题~\ref{ex:c22q24}：简谐振动初值问题]\phantomsection\label{ans:c22q24}
特征方程为 $r^2+9=0$，根为 $\pm 3i$。实通解为 $y=A\cos(3x)+B\sin(3x)$。第一初值给出 $A=4$，对通解求导并代入第二初值得到 $3B=0$，所以
\[
y(x)=4\cos(3x),\qquad T=\frac{2\pi}{3}.
\]
直接求二阶导数可以验证方程。振幅由初始位移确定，周期由方程系数确定；对于这里的线性模型，振幅并不影响周期。这一题作为高等数学中常系数二阶方程的补充。
\end{solve}

\clearpage
\section{初值问题、平衡解与共振}\label{sec:ch23}
常微分方程的求解必须同时满足方程、初值和定义区间。非线性方程可能有平衡解，线性强迫振动可能出现共振；这些结构会影响通解形式。

\textbf{资源对应：}CLP-2 \S~2.4；二阶方程为本科高等数学补充。本章共 24 道编号习题。
\subsection{习题}
\begin{example}[变量系数的可分离方程]\label{ex:c23q1}
解 $y^{\prime}=2xy$、$y(0)=2$，并求 $y(1)$。
\par\textit{改编训练。}\par\hyperref[ans:c23q1]{解答见第~\pageref{ans:c23q1} 页。}
\end{example}
\begin{example}[变量系数的可分离方程]\label{ex:c23q2}
解 $y^{\prime}=4xy$、$y(0)=1$，并求 $y(1)$。
\par\textit{改编训练。}\par\hyperref[ans:c23q2]{解答见第~\pageref{ans:c23q2} 页。}
\end{example}
\begin{example}[变量系数的可分离方程]\label{ex:c23q3}
解 $y^{\prime}=4xy$、$y(0)=3$，并求 $y(1)$。
\par\textit{改编训练。}\par\hyperref[ans:c23q3]{解答见第~\pageref{ans:c23q3} 页。}
\end{example}
\begin{example}[变量系数的可分离方程]\label{ex:c23q4}
解 $y^{\prime}=6xy$、$y(0)=2$，并求 $y(1)$。
\par\textit{改编训练。}\par\hyperref[ans:c23q4]{解答见第~\pageref{ans:c23q4} 页。}
\end{example}
\begin{example}[变量系数的可分离方程]\label{ex:c23q5}
解 $y^{\prime}=8xy$、$y(0)=1$，并求 $y(1)$。
\par\textit{改编训练。}\par\hyperref[ans:c23q5]{解答见第~\pageref{ans:c23q5} 页。}
\end{example}
\begin{example}[变量系数的可分离方程]\label{ex:c23q6}
解 $y^{\prime}=6xy$、$y(0)=4$，并求 $y(1)$。
\par\textit{改编训练。}\par\hyperref[ans:c23q6]{解答见第~\pageref{ans:c23q6} 页。}
\end{example}
\begin{example}[Logistic 初值问题]\label{ex:c23q7}
解 $y^{\prime}=1y(1-y/2)$、$y(0)=1$，并求长期极限。
\par\textit{改编训练。}\par\hyperref[ans:c23q7]{解答见第~\pageref{ans:c23q7} 页。}
\end{example}
\begin{example}[Logistic 初值问题]\label{ex:c23q8}
解 $y^{\prime}=2y(1-y/4)$、$y(0)=2$，并求长期极限。
\par\textit{改编训练。}\par\hyperref[ans:c23q8]{解答见第~\pageref{ans:c23q8} 页。}
\end{example}
\begin{example}[Logistic 初值问题]\label{ex:c23q9}
解 $y^{\prime}=3y(1-y/6)$、$y(0)=3$，并求长期极限。
\par\textit{改编训练。}\par\hyperref[ans:c23q9]{解答见第~\pageref{ans:c23q9} 页。}
\end{example}
\begin{example}[Logistic 初值问题]\label{ex:c23q10}
解 $y^{\prime}=2y(1-y/8)$、$y(0)=4$，并求长期极限。
\par\textit{改编训练。}\par\hyperref[ans:c23q10]{解答见第~\pageref{ans:c23q10} 页。}
\end{example}
\begin{example}[Logistic 初值问题]\label{ex:c23q11}
解 $y^{\prime}=4y(1-y/10)$、$y(0)=5$，并求长期极限。
\par\textit{改编训练。}\par\hyperref[ans:c23q11]{解答见第~\pageref{ans:c23q11} 页。}
\end{example}
\begin{example}[Logistic 初值问题]\label{ex:c23q12}
解 $y^{\prime}=3y(1-y/12)$、$y(0)=6$，并求长期极限。
\par\textit{改编训练。}\par\hyperref[ans:c23q12]{解答见第~\pageref{ans:c23q12} 页。}
\end{example}
\begin{example}[重复特征根]\label{ex:c23q13}
解 $y^{\prime\prime}-2y^{\prime}+1y=0$、$y(0)=1$、$y^{\prime}(0)=0$。
\par\textit{改编训练。}\par\hyperref[ans:c23q13]{解答见第~\pageref{ans:c23q13} 页。}
\end{example}
\begin{example}[重复特征根]\label{ex:c23q14}
解 $y^{\prime\prime}-4y^{\prime}+4y=0$、$y(0)=1$、$y^{\prime}(0)=0$。
\par\textit{改编训练。}\par\hyperref[ans:c23q14]{解答见第~\pageref{ans:c23q14} 页。}
\end{example}
\begin{example}[重复特征根]\label{ex:c23q15}
解 $y^{\prime\prime}-6y^{\prime}+9y=0$、$y(0)=1$、$y^{\prime}(0)=0$。
\par\textit{改编训练。}\par\hyperref[ans:c23q15]{解答见第~\pageref{ans:c23q15} 页。}
\end{example}
\begin{example}[重复特征根]\label{ex:c23q16}
解 $y^{\prime\prime}-8y^{\prime}+16y=0$、$y(0)=1$、$y^{\prime}(0)=0$。
\par\textit{改编训练。}\par\hyperref[ans:c23q16]{解答见第~\pageref{ans:c23q16} 页。}
\end{example}
\begin{example}[重复特征根]\label{ex:c23q17}
解 $y^{\prime\prime}-10y^{\prime}+25y=0$、$y(0)=1$、$y^{\prime}(0)=0$。
\par\textit{改编训练。}\par\hyperref[ans:c23q17]{解答见第~\pageref{ans:c23q17} 页。}
\end{example}
\begin{example}[重复特征根]\label{ex:c23q18}
解 $y^{\prime\prime}-12y^{\prime}+36y=0$、$y(0)=1$、$y^{\prime}(0)=0$。
\par\textit{改编训练。}\par\hyperref[ans:c23q18]{解答见第~\pageref{ans:c23q18} 页。}
\end{example}
\begin{example}[共振强迫振动]\label{ex:c23q19}
解 $y^{\prime\prime}+1y=2\cos(1x)$，$y(0)=y^{\prime}(0)=0$。
\par\textit{改编训练。}\par\hyperref[ans:c23q19]{解答见第~\pageref{ans:c23q19} 页。}
\end{example}
\begin{example}[共振强迫振动]\label{ex:c23q20}
解 $y^{\prime\prime}+4y=1\cos(2x)$，$y(0)=y^{\prime}(0)=0$。
\par\textit{改编训练。}\par\hyperref[ans:c23q20]{解答见第~\pageref{ans:c23q20} 页。}
\end{example}
\begin{example}[共振强迫振动]\label{ex:c23q21}
解 $y^{\prime\prime}+4y=3\cos(2x)$，$y(0)=y^{\prime}(0)=0$。
\par\textit{改编训练。}\par\hyperref[ans:c23q21]{解答见第~\pageref{ans:c23q21} 页。}
\end{example}
\begin{example}[共振强迫振动]\label{ex:c23q22}
解 $y^{\prime\prime}+9y=2\cos(3x)$，$y(0)=y^{\prime}(0)=0$。
\par\textit{改编训练。}\par\hyperref[ans:c23q22]{解答见第~\pageref{ans:c23q22} 页。}
\end{example}
\begin{example}[共振强迫振动]\label{ex:c23q23}
解 $y^{\prime\prime}+16y=1\cos(4x)$，$y(0)=y^{\prime}(0)=0$。
\par\textit{改编训练。}\par\hyperref[ans:c23q23]{解答见第~\pageref{ans:c23q23} 页。}
\end{example}
\begin{example}[共振强迫振动]\label{ex:c23q24}
解 $y^{\prime\prime}+9y=4\cos(3x)$，$y(0)=y^{\prime}(0)=0$。
\par\textit{改编训练。}\par\hyperref[ans:c23q24]{解答见第~\pageref{ans:c23q24} 页。}
\end{example}
\subsection{解答}
\begin{solve}[习题~\ref{ex:c23q1}：变量系数的可分离方程]\phantomsection\label{ans:c23q1}
非零初值对应的解在局部非零，分离变量后 $dy/y=2x\,dx$。积分得到 $\log|y|=1x^2+C$，于是
\[
y(x)=2e^{1x^2},\qquad y(1)=2e^{1}.
\]
这个表达式在全实轴有定义，直接求导验证即可将解扩展到全轴。原方程还有零解，但它不满足给定初值；分离变量除以 $y$ 的步骤必须明确这一点。
\end{solve}

\begin{solve}[习题~\ref{ex:c23q2}：变量系数的可分离方程]\phantomsection\label{ans:c23q2}
非零初值对应的解在局部非零，分离变量后 $dy/y=4x\,dx$。积分得到 $\log|y|=2x^2+C$，于是
\[
y(x)=1e^{2x^2},\qquad y(1)=1e^{2}.
\]
这个表达式在全实轴有定义，直接求导验证即可将解扩展到全轴。原方程还有零解，但它不满足给定初值；分离变量除以 $y$ 的步骤必须明确这一点。
\end{solve}

\begin{solve}[习题~\ref{ex:c23q3}：变量系数的可分离方程]\phantomsection\label{ans:c23q3}
非零初值对应的解在局部非零，分离变量后 $dy/y=4x\,dx$。积分得到 $\log|y|=2x^2+C$，于是
\[
y(x)=3e^{2x^2},\qquad y(1)=3e^{2}.
\]
这个表达式在全实轴有定义，直接求导验证即可将解扩展到全轴。原方程还有零解，但它不满足给定初值；分离变量除以 $y$ 的步骤必须明确这一点。
\end{solve}

\begin{solve}[习题~\ref{ex:c23q4}：变量系数的可分离方程]\phantomsection\label{ans:c23q4}
非零初值对应的解在局部非零，分离变量后 $dy/y=6x\,dx$。积分得到 $\log|y|=3x^2+C$，于是
\[
y(x)=2e^{3x^2},\qquad y(1)=2e^{3}.
\]
这个表达式在全实轴有定义，直接求导验证即可将解扩展到全轴。原方程还有零解，但它不满足给定初值；分离变量除以 $y$ 的步骤必须明确这一点。
\end{solve}

\begin{solve}[习题~\ref{ex:c23q5}：变量系数的可分离方程]\phantomsection\label{ans:c23q5}
非零初值对应的解在局部非零，分离变量后 $dy/y=8x\,dx$。积分得到 $\log|y|=4x^2+C$，于是
\[
y(x)=1e^{4x^2},\qquad y(1)=1e^{4}.
\]
这个表达式在全实轴有定义，直接求导验证即可将解扩展到全轴。原方程还有零解，但它不满足给定初值；分离变量除以 $y$ 的步骤必须明确这一点。
\end{solve}

\begin{solve}[习题~\ref{ex:c23q6}：变量系数的可分离方程]\phantomsection\label{ans:c23q6}
非零初值对应的解在局部非零，分离变量后 $dy/y=6x\,dx$。积分得到 $\log|y|=3x^2+C$，于是
\[
y(x)=4e^{3x^2},\qquad y(1)=4e^{3}.
\]
这个表达式在全实轴有定义，直接求导验证即可将解扩展到全轴。原方程还有零解，但它不满足给定初值；分离变量除以 $y$ 的步骤必须明确这一点。
\end{solve}

\begin{solve}[习题~\ref{ex:c23q7}：Logistic 初值问题]\phantomsection\label{ans:c23q7}
两个平衡解为 $y=0$ 和 $y=2$。对本题初值处于两者之间的解，分离变量并分解部分分式，得 $\log(y/(2-y))=1x+C$。初值使 $C=0$，所以
\[
y(x)=\frac{2}{1+e^{-1x}},\qquad\lim_{x\to\infty}y(x)=2.
\]
分母始终为正，解在全实轴上存在并保持 $0<y<2$。它先近似指数增长，后来逐渐接近平衡容量。验证时求导得到 $2e^{-1x}/(1+e^{-1x})^2$，与代入方程右边一致。
\end{solve}

\begin{solve}[习题~\ref{ex:c23q8}：Logistic 初值问题]\phantomsection\label{ans:c23q8}
两个平衡解为 $y=0$ 和 $y=4$。对本题初值处于两者之间的解，分离变量并分解部分分式，得 $\log(y/(4-y))=2x+C$。初值使 $C=0$，所以
\[
y(x)=\frac{4}{1+e^{-2x}},\qquad\lim_{x\to\infty}y(x)=4.
\]
分母始终为正，解在全实轴上存在并保持 $0<y<4$。它先近似指数增长，后来逐渐接近平衡容量。验证时求导得到 $8e^{-2x}/(1+e^{-2x})^2$，与代入方程右边一致。
\end{solve}

\begin{solve}[习题~\ref{ex:c23q9}：Logistic 初值问题]\phantomsection\label{ans:c23q9}
两个平衡解为 $y=0$ 和 $y=6$。对本题初值处于两者之间的解，分离变量并分解部分分式，得 $\log(y/(6-y))=3x+C$。初值使 $C=0$，所以
\[
y(x)=\frac{6}{1+e^{-3x}},\qquad\lim_{x\to\infty}y(x)=6.
\]
分母始终为正，解在全实轴上存在并保持 $0<y<6$。它先近似指数增长，后来逐渐接近平衡容量。验证时求导得到 $18e^{-3x}/(1+e^{-3x})^2$，与代入方程右边一致。
\end{solve}

\begin{solve}[习题~\ref{ex:c23q10}：Logistic 初值问题]\phantomsection\label{ans:c23q10}
两个平衡解为 $y=0$ 和 $y=8$。对本题初值处于两者之间的解，分离变量并分解部分分式，得 $\log(y/(8-y))=2x+C$。初值使 $C=0$，所以
\[
y(x)=\frac{8}{1+e^{-2x}},\qquad\lim_{x\to\infty}y(x)=8.
\]
分母始终为正，解在全实轴上存在并保持 $0<y<8$。它先近似指数增长，后来逐渐接近平衡容量。验证时求导得到 $16e^{-2x}/(1+e^{-2x})^2$，与代入方程右边一致。
\end{solve}

\begin{solve}[习题~\ref{ex:c23q11}：Logistic 初值问题]\phantomsection\label{ans:c23q11}
两个平衡解为 $y=0$ 和 $y=10$。对本题初值处于两者之间的解，分离变量并分解部分分式，得 $\log(y/(10-y))=4x+C$。初值使 $C=0$，所以
\[
y(x)=\frac{10}{1+e^{-4x}},\qquad\lim_{x\to\infty}y(x)=10.
\]
分母始终为正，解在全实轴上存在并保持 $0<y<10$。它先近似指数增长，后来逐渐接近平衡容量。验证时求导得到 $40e^{-4x}/(1+e^{-4x})^2$，与代入方程右边一致。
\end{solve}

\begin{solve}[习题~\ref{ex:c23q12}：Logistic 初值问题]\phantomsection\label{ans:c23q12}
两个平衡解为 $y=0$ 和 $y=12$。对本题初值处于两者之间的解，分离变量并分解部分分式，得 $\log(y/(12-y))=3x+C$。初值使 $C=0$，所以
\[
y(x)=\frac{12}{1+e^{-3x}},\qquad\lim_{x\to\infty}y(x)=12.
\]
分母始终为正，解在全实轴上存在并保持 $0<y<12$。它先近似指数增长，后来逐渐接近平衡容量。验证时求导得到 $36e^{-3x}/(1+e^{-3x})^2$，与代入方程右边一致。
\end{solve}

\begin{solve}[习题~\ref{ex:c23q13}：重复特征根]\phantomsection\label{ans:c23q13}
特征方程为 $(r-1)^2=0$，出现二重根，通解应写成 $(C_1+C_2x)e^{1x}$。初值 $y(0)=1$ 给出 $C_1=1$；求导后 $y^{\prime}(0)=C_2+1C_1=0$，因此
\[
y(x)=(1-1x)e^{1x}.
\]
二重根情形的第二个线性无关解是 $xe^{ax}$，不能重复写两个相同指数并给它们不同系数，否则实际上只得到一个自由常数，无法满足任意两条初值条件。
\end{solve}

\begin{solve}[习题~\ref{ex:c23q14}：重复特征根]\phantomsection\label{ans:c23q14}
特征方程为 $(r-2)^2=0$，出现二重根，通解应写成 $(C_1+C_2x)e^{2x}$。初值 $y(0)=1$ 给出 $C_1=1$；求导后 $y^{\prime}(0)=C_2+2C_1=0$，因此
\[
y(x)=(1-2x)e^{2x}.
\]
二重根情形的第二个线性无关解是 $xe^{ax}$，不能重复写两个相同指数并给它们不同系数，否则实际上只得到一个自由常数，无法满足任意两条初值条件。
\end{solve}

\begin{solve}[习题~\ref{ex:c23q15}：重复特征根]\phantomsection\label{ans:c23q15}
特征方程为 $(r-3)^2=0$，出现二重根，通解应写成 $(C_1+C_2x)e^{3x}$。初值 $y(0)=1$ 给出 $C_1=1$；求导后 $y^{\prime}(0)=C_2+3C_1=0$，因此
\[
y(x)=(1-3x)e^{3x}.
\]
二重根情形的第二个线性无关解是 $xe^{ax}$，不能重复写两个相同指数并给它们不同系数，否则实际上只得到一个自由常数，无法满足任意两条初值条件。
\end{solve}

\begin{solve}[习题~\ref{ex:c23q16}：重复特征根]\phantomsection\label{ans:c23q16}
特征方程为 $(r-4)^2=0$，出现二重根，通解应写成 $(C_1+C_2x)e^{4x}$。初值 $y(0)=1$ 给出 $C_1=1$；求导后 $y^{\prime}(0)=C_2+4C_1=0$，因此
\[
y(x)=(1-4x)e^{4x}.
\]
二重根情形的第二个线性无关解是 $xe^{ax}$，不能重复写两个相同指数并给它们不同系数，否则实际上只得到一个自由常数，无法满足任意两条初值条件。
\end{solve}

\begin{solve}[习题~\ref{ex:c23q17}：重复特征根]\phantomsection\label{ans:c23q17}
特征方程为 $(r-5)^2=0$，出现二重根，通解应写成 $(C_1+C_2x)e^{5x}$。初值 $y(0)=1$ 给出 $C_1=1$；求导后 $y^{\prime}(0)=C_2+5C_1=0$，因此
\[
y(x)=(1-5x)e^{5x}.
\]
二重根情形的第二个线性无关解是 $xe^{ax}$，不能重复写两个相同指数并给它们不同系数，否则实际上只得到一个自由常数，无法满足任意两条初值条件。
\end{solve}

\begin{solve}[习题~\ref{ex:c23q18}：重复特征根]\phantomsection\label{ans:c23q18}
特征方程为 $(r-6)^2=0$，出现二重根，通解应写成 $(C_1+C_2x)e^{6x}$。初值 $y(0)=1$ 给出 $C_1=1$；求导后 $y^{\prime}(0)=C_2+6C_1=0$，因此
\[
y(x)=(1-6x)e^{6x}.
\]
二重根情形的第二个线性无关解是 $xe^{ax}$，不能重复写两个相同指数并给它们不同系数，否则实际上只得到一个自由常数，无法满足任意两条初值条件。
\end{solve}

\begin{solve}[习题~\ref{ex:c23q19}：共振强迫振动]\phantomsection\label{ans:c23q19}
强迫频率等于齐次方程的自然频率，不能直接试取余弦倍数作特解。取 $y_p=A x\sin(1x)$，计算可得
\[
y_p^{\prime\prime}+1y_p=2A\cos(1x).
\]
令 $A=1$。该特解在0处本身和一阶导数都为0，故已经满足初值，答案为
\[
y(x)=1x\sin(1x).
\]
前面的线性因子使振幅随时间增长，这是无阻尼线性模型中的共振表现。它来自与齐次解重合的强迫频率，而非积分常数选择不当。
\end{solve}

\begin{solve}[习题~\ref{ex:c23q20}：共振强迫振动]\phantomsection\label{ans:c23q20}
强迫频率等于齐次方程的自然频率，不能直接试取余弦倍数作特解。取 $y_p=A x\sin(2x)$，计算可得
\[
y_p^{\prime\prime}+4y_p=4A\cos(2x).
\]
令 $A=\frac{1}{4}$。该特解在0处本身和一阶导数都为0，故已经满足初值，答案为
\[
y(x)=\frac{1}{4}x\sin(2x).
\]
前面的线性因子使振幅随时间增长，这是无阻尼线性模型中的共振表现。它来自与齐次解重合的强迫频率，而非积分常数选择不当。
\end{solve}

\begin{solve}[习题~\ref{ex:c23q21}：共振强迫振动]\phantomsection\label{ans:c23q21}
强迫频率等于齐次方程的自然频率，不能直接试取余弦倍数作特解。取 $y_p=A x\sin(2x)$，计算可得
\[
y_p^{\prime\prime}+4y_p=4A\cos(2x).
\]
令 $A=\frac{3}{4}$。该特解在0处本身和一阶导数都为0，故已经满足初值，答案为
\[
y(x)=\frac{3}{4}x\sin(2x).
\]
前面的线性因子使振幅随时间增长，这是无阻尼线性模型中的共振表现。它来自与齐次解重合的强迫频率，而非积分常数选择不当。
\end{solve}

\begin{solve}[习题~\ref{ex:c23q22}：共振强迫振动]\phantomsection\label{ans:c23q22}
强迫频率等于齐次方程的自然频率，不能直接试取余弦倍数作特解。取 $y_p=A x\sin(3x)$，计算可得
\[
y_p^{\prime\prime}+9y_p=6A\cos(3x).
\]
令 $A=\frac{1}{3}$。该特解在0处本身和一阶导数都为0，故已经满足初值，答案为
\[
y(x)=\frac{1}{3}x\sin(3x).
\]
前面的线性因子使振幅随时间增长，这是无阻尼线性模型中的共振表现。它来自与齐次解重合的强迫频率，而非积分常数选择不当。
\end{solve}

\begin{solve}[习题~\ref{ex:c23q23}：共振强迫振动]\phantomsection\label{ans:c23q23}
强迫频率等于齐次方程的自然频率，不能直接试取余弦倍数作特解。取 $y_p=A x\sin(4x)$，计算可得
\[
y_p^{\prime\prime}+16y_p=8A\cos(4x).
\]
令 $A=\frac{1}{8}$。该特解在0处本身和一阶导数都为0，故已经满足初值，答案为
\[
y(x)=\frac{1}{8}x\sin(4x).
\]
前面的线性因子使振幅随时间增长，这是无阻尼线性模型中的共振表现。它来自与齐次解重合的强迫频率，而非积分常数选择不当。
\end{solve}

\begin{solve}[习题~\ref{ex:c23q24}：共振强迫振动]\phantomsection\label{ans:c23q24}
强迫频率等于齐次方程的自然频率，不能直接试取余弦倍数作特解。取 $y_p=A x\sin(3x)$，计算可得
\[
y_p^{\prime\prime}+9y_p=6A\cos(3x).
\]
令 $A=\frac{2}{3}$。该特解在0处本身和一阶导数都为0，故已经满足初值，答案为
\[
y(x)=\frac{2}{3}x\sin(3x).
\]
前面的线性因子使振幅随时间增长，这是无阻尼线性模型中的共振表现。它来自与齐次解重合的强迫频率，而非积分常数选择不当。
\end{solve}

\clearpage
\section{正项级数与判别方法}\label{sec:ch24}
级数（series）的收敛是部分和数列的收敛。通项趋零是必要条件而非充分条件；正项级数可以借助比较、比值或根值判别，但临界值等于1时需要另作分析。

\textbf{资源对应：}CLP-2 \S~3.2、\S~3.3。本章共 24 道编号习题。
\subsection{习题}
\begin{example}[几何级数与尾项]\label{ex:c24q1}
求 $\displaystyle\sum_{n=0}^\infty(\frac{1}{2})^n$ 的和，写出保留到 $n=N$ 后的误差。
\par\textit{改编训练。}\par\hyperref[ans:c24q1]{解答见第~\pageref{ans:c24q1} 页。}
\end{example}
\begin{example}[几何级数与尾项]\label{ex:c24q2}
求 $\displaystyle\sum_{n=0}^\infty(\frac{2}{3})^n$ 的和，写出保留到 $n=N$ 后的误差。
\par\textit{改编训练。}\par\hyperref[ans:c24q2]{解答见第~\pageref{ans:c24q2} 页。}
\end{example}
\begin{example}[几何级数与尾项]\label{ex:c24q3}
求 $\displaystyle\sum_{n=0}^\infty(\frac{3}{4})^n$ 的和，写出保留到 $n=N$ 后的误差。
\par\textit{改编训练。}\par\hyperref[ans:c24q3]{解答见第~\pageref{ans:c24q3} 页。}
\end{example}
\begin{example}[几何级数与尾项]\label{ex:c24q4}
求 $\displaystyle\sum_{n=0}^\infty(\frac{1}{3})^n$ 的和，写出保留到 $n=N$ 后的误差。
\par\textit{改编训练。}\par\hyperref[ans:c24q4]{解答见第~\pageref{ans:c24q4} 页。}
\end{example}
\begin{example}[几何级数与尾项]\label{ex:c24q5}
求 $\displaystyle\sum_{n=0}^\infty(\frac{2}{5})^n$ 的和，写出保留到 $n=N$ 后的误差。
\par\textit{改编训练。}\par\hyperref[ans:c24q5]{解答见第~\pageref{ans:c24q5} 页。}
\end{example}
\begin{example}[几何级数与尾项]\label{ex:c24q6}
求 $\displaystyle\sum_{n=0}^\infty(\frac{3}{7})^n$ 的和，写出保留到 $n=N$ 后的误差。
\par\textit{改编训练。}\par\hyperref[ans:c24q6]{解答见第~\pageref{ans:c24q6} 页。}
\end{example}
\begin{example}[裂项求和]\label{ex:c24q7}
求 $\displaystyle\sum_{n=1}^\infty\frac{1}{(n+1)(n+2)}$，并给出部分和。
\par\textit{改编训练。}\par\hyperref[ans:c24q7]{解答见第~\pageref{ans:c24q7} 页。}
\end{example}
\begin{example}[裂项求和]\label{ex:c24q8}
求 $\displaystyle\sum_{n=1}^\infty\frac{1}{(n+2)(n+3)}$，并给出部分和。
\par\textit{改编训练。}\par\hyperref[ans:c24q8]{解答见第~\pageref{ans:c24q8} 页。}
\end{example}
\begin{example}[裂项求和]\label{ex:c24q9}
求 $\displaystyle\sum_{n=1}^\infty\frac{1}{(n+3)(n+4)}$，并给出部分和。
\par\textit{改编训练。}\par\hyperref[ans:c24q9]{解答见第~\pageref{ans:c24q9} 页。}
\end{example}
\begin{example}[裂项求和]\label{ex:c24q10}
求 $\displaystyle\sum_{n=1}^\infty\frac{1}{(n+4)(n+5)}$，并给出部分和。
\par\textit{改编训练。}\par\hyperref[ans:c24q10]{解答见第~\pageref{ans:c24q10} 页。}
\end{example}
\begin{example}[裂项求和]\label{ex:c24q11}
求 $\displaystyle\sum_{n=1}^\infty\frac{1}{(n+5)(n+6)}$，并给出部分和。
\par\textit{改编训练。}\par\hyperref[ans:c24q11]{解答见第~\pageref{ans:c24q11} 页。}
\end{example}
\begin{example}[裂项求和]\label{ex:c24q12}
求 $\displaystyle\sum_{n=1}^\infty\frac{1}{(n+6)(n+7)}$，并给出部分和。
\par\textit{改编训练。}\par\hyperref[ans:c24q12]{解答见第~\pageref{ans:c24q12} 页。}
\end{example}
\begin{example}[阶乘压过指数]\label{ex:c24q13}
证明 $\displaystyle\sum_{n=0}^\infty\frac{1^n}{n!}$ 收敛，并用比值说明它的尾部如何估计。
\par\textit{改编训练。}\par\hyperref[ans:c24q13]{解答见第~\pageref{ans:c24q13} 页。}
\end{example}
\begin{example}[阶乘压过指数]\label{ex:c24q14}
证明 $\displaystyle\sum_{n=0}^\infty\frac{2^n}{n!}$ 收敛，并用比值说明它的尾部如何估计。
\par\textit{改编训练。}\par\hyperref[ans:c24q14]{解答见第~\pageref{ans:c24q14} 页。}
\end{example}
\begin{example}[阶乘压过指数]\label{ex:c24q15}
证明 $\displaystyle\sum_{n=0}^\infty\frac{3^n}{n!}$ 收敛，并用比值说明它的尾部如何估计。
\par\textit{改编训练。}\par\hyperref[ans:c24q15]{解答见第~\pageref{ans:c24q15} 页。}
\end{example}
\begin{example}[阶乘压过指数]\label{ex:c24q16}
证明 $\displaystyle\sum_{n=0}^\infty\frac{4^n}{n!}$ 收敛，并用比值说明它的尾部如何估计。
\par\textit{改编训练。}\par\hyperref[ans:c24q16]{解答见第~\pageref{ans:c24q16} 页。}
\end{example}
\begin{example}[阶乘压过指数]\label{ex:c24q17}
证明 $\displaystyle\sum_{n=0}^\infty\frac{5^n}{n!}$ 收敛，并用比值说明它的尾部如何估计。
\par\textit{改编训练。}\par\hyperref[ans:c24q17]{解答见第~\pageref{ans:c24q17} 页。}
\end{example}
\begin{example}[阶乘压过指数]\label{ex:c24q18}
证明 $\displaystyle\sum_{n=0}^\infty\frac{6^n}{n!}$ 收敛，并用比值说明它的尾部如何估计。
\par\textit{改编训练。}\par\hyperref[ans:c24q18]{解答见第~\pageref{ans:c24q18} 页。}
\end{example}
\begin{example}[极限比较与主导次数]\label{ex:c24q19}
判断 $\displaystyle\sum_{n=1}^\infty\frac{1n+1}{n^3+1}$ 的敛散性，并指出不能直接使用的结论。
\par\textit{改编训练。}\par\hyperref[ans:c24q19]{解答见第~\pageref{ans:c24q19} 页。}
\end{example}
\begin{example}[极限比较与主导次数]\label{ex:c24q20}
判断 $\displaystyle\sum_{n=1}^\infty\frac{2n+1}{n^3+2}$ 的敛散性，并指出不能直接使用的结论。
\par\textit{改编训练。}\par\hyperref[ans:c24q20]{解答见第~\pageref{ans:c24q20} 页。}
\end{example}
\begin{example}[极限比较与主导次数]\label{ex:c24q21}
判断 $\displaystyle\sum_{n=1}^\infty\frac{3n+1}{n^3+3}$ 的敛散性，并指出不能直接使用的结论。
\par\textit{改编训练。}\par\hyperref[ans:c24q21]{解答见第~\pageref{ans:c24q21} 页。}
\end{example}
\begin{example}[极限比较与主导次数]\label{ex:c24q22}
判断 $\displaystyle\sum_{n=1}^\infty\frac{4n+1}{n^3+4}$ 的敛散性，并指出不能直接使用的结论。
\par\textit{改编训练。}\par\hyperref[ans:c24q22]{解答见第~\pageref{ans:c24q22} 页。}
\end{example}
\begin{example}[极限比较与主导次数]\label{ex:c24q23}
判断 $\displaystyle\sum_{n=1}^\infty\frac{5n+1}{n^3+5}$ 的敛散性，并指出不能直接使用的结论。
\par\textit{改编训练。}\par\hyperref[ans:c24q23]{解答见第~\pageref{ans:c24q23} 页。}
\end{example}
\begin{example}[极限比较与主导次数]\label{ex:c24q24}
判断 $\displaystyle\sum_{n=1}^\infty\frac{6n+1}{n^3+6}$ 的敛散性，并指出不能直接使用的结论。
\par\textit{改编训练。}\par\hyperref[ans:c24q24]{解答见第~\pageref{ans:c24q24} 页。}
\end{example}
\subsection{解答}
\begin{solve}[习题~\ref{ex:c24q1}：几何级数与尾项]\phantomsection\label{ans:c24q1}
公比 $q=\frac{1}{2}$ 满足 $0<q<1$。有限和恒等式给出
\[
S_N=\frac{1-(\frac{1}{2})^{N+1}}{1-\frac{1}{2}}.
\]
当 $N\to\infty$，幂项趋零，故总和为 $2$。余项为
\[
R_N=\frac{(\frac{1}{2})^{N+1}}{1-\frac{1}{2}}=2(\frac{1}{2})^{N+1}.
\]
正项级数的部分和从下方逼近总和。若要求误差小于某个容许值，可以直接解这个几何衰减不等式，而无需逐项试加。
\end{solve}

\begin{solve}[习题~\ref{ex:c24q2}：几何级数与尾项]\phantomsection\label{ans:c24q2}
公比 $q=\frac{2}{3}$ 满足 $0<q<1$。有限和恒等式给出
\[
S_N=\frac{1-(\frac{2}{3})^{N+1}}{1-\frac{2}{3}}.
\]
当 $N\to\infty$，幂项趋零，故总和为 $3$。余项为
\[
R_N=\frac{(\frac{2}{3})^{N+1}}{1-\frac{2}{3}}=3(\frac{2}{3})^{N+1}.
\]
正项级数的部分和从下方逼近总和。若要求误差小于某个容许值，可以直接解这个几何衰减不等式，而无需逐项试加。
\end{solve}

\begin{solve}[习题~\ref{ex:c24q3}：几何级数与尾项]\phantomsection\label{ans:c24q3}
公比 $q=\frac{3}{4}$ 满足 $0<q<1$。有限和恒等式给出
\[
S_N=\frac{1-(\frac{3}{4})^{N+1}}{1-\frac{3}{4}}.
\]
当 $N\to\infty$，幂项趋零，故总和为 $4$。余项为
\[
R_N=\frac{(\frac{3}{4})^{N+1}}{1-\frac{3}{4}}=4(\frac{3}{4})^{N+1}.
\]
正项级数的部分和从下方逼近总和。若要求误差小于某个容许值，可以直接解这个几何衰减不等式，而无需逐项试加。
\end{solve}

\begin{solve}[习题~\ref{ex:c24q4}：几何级数与尾项]\phantomsection\label{ans:c24q4}
公比 $q=\frac{1}{3}$ 满足 $0<q<1$。有限和恒等式给出
\[
S_N=\frac{1-(\frac{1}{3})^{N+1}}{1-\frac{1}{3}}.
\]
当 $N\to\infty$，幂项趋零，故总和为 $\frac{3}{2}$。余项为
\[
R_N=\frac{(\frac{1}{3})^{N+1}}{1-\frac{1}{3}}=\frac{3}{2}(\frac{1}{3})^{N+1}.
\]
正项级数的部分和从下方逼近总和。若要求误差小于某个容许值，可以直接解这个几何衰减不等式，而无需逐项试加。
\end{solve}

\begin{solve}[习题~\ref{ex:c24q5}：几何级数与尾项]\phantomsection\label{ans:c24q5}
公比 $q=\frac{2}{5}$ 满足 $0<q<1$。有限和恒等式给出
\[
S_N=\frac{1-(\frac{2}{5})^{N+1}}{1-\frac{2}{5}}.
\]
当 $N\to\infty$，幂项趋零，故总和为 $\frac{5}{3}$。余项为
\[
R_N=\frac{(\frac{2}{5})^{N+1}}{1-\frac{2}{5}}=\frac{5}{3}(\frac{2}{5})^{N+1}.
\]
正项级数的部分和从下方逼近总和。若要求误差小于某个容许值，可以直接解这个几何衰减不等式，而无需逐项试加。
\end{solve}

\begin{solve}[习题~\ref{ex:c24q6}：几何级数与尾项]\phantomsection\label{ans:c24q6}
公比 $q=\frac{3}{7}$ 满足 $0<q<1$。有限和恒等式给出
\[
S_N=\frac{1-(\frac{3}{7})^{N+1}}{1-\frac{3}{7}}.
\]
当 $N\to\infty$，幂项趋零，故总和为 $\frac{7}{4}$。余项为
\[
R_N=\frac{(\frac{3}{7})^{N+1}}{1-\frac{3}{7}}=\frac{7}{4}(\frac{3}{7})^{N+1}.
\]
正项级数的部分和从下方逼近总和。若要求误差小于某个容许值，可以直接解这个几何衰减不等式，而无需逐项试加。
\end{solve}

\begin{solve}[习题~\ref{ex:c24q7}：裂项求和]\phantomsection\label{ans:c24q7}
先作部分分式分解：
\[
\frac1{(n+1)(n+2)}=\frac1{n+1}-\frac1{n+2}.
\]
写出有限部分和后，中间项逐项抵消，得到
\[
S_N=\frac1{2}-\frac1{N+2},\qquad S=\frac{1}{2}.
\]
这种消去应当先对有限和进行；不能把两个分别发散的调和型无穷级数直接相减。余项精确等于 $1/(N+a+1)$，从而还可读出收敛速度为 $O(1/N)$。
\end{solve}

\begin{solve}[习题~\ref{ex:c24q8}：裂项求和]\phantomsection\label{ans:c24q8}
先作部分分式分解：
\[
\frac1{(n+2)(n+3)}=\frac1{n+2}-\frac1{n+3}.
\]
写出有限部分和后，中间项逐项抵消，得到
\[
S_N=\frac1{3}-\frac1{N+3},\qquad S=\frac{1}{3}.
\]
这种消去应当先对有限和进行；不能把两个分别发散的调和型无穷级数直接相减。余项精确等于 $1/(N+a+1)$，从而还可读出收敛速度为 $O(1/N)$。
\end{solve}

\begin{solve}[习题~\ref{ex:c24q9}：裂项求和]\phantomsection\label{ans:c24q9}
先作部分分式分解：
\[
\frac1{(n+3)(n+4)}=\frac1{n+3}-\frac1{n+4}.
\]
写出有限部分和后，中间项逐项抵消，得到
\[
S_N=\frac1{4}-\frac1{N+4},\qquad S=\frac{1}{4}.
\]
这种消去应当先对有限和进行；不能把两个分别发散的调和型无穷级数直接相减。余项精确等于 $1/(N+a+1)$，从而还可读出收敛速度为 $O(1/N)$。
\end{solve}

\begin{solve}[习题~\ref{ex:c24q10}：裂项求和]\phantomsection\label{ans:c24q10}
先作部分分式分解：
\[
\frac1{(n+4)(n+5)}=\frac1{n+4}-\frac1{n+5}.
\]
写出有限部分和后，中间项逐项抵消，得到
\[
S_N=\frac1{5}-\frac1{N+5},\qquad S=\frac{1}{5}.
\]
这种消去应当先对有限和进行；不能把两个分别发散的调和型无穷级数直接相减。余项精确等于 $1/(N+a+1)$，从而还可读出收敛速度为 $O(1/N)$。
\end{solve}

\begin{solve}[习题~\ref{ex:c24q11}：裂项求和]\phantomsection\label{ans:c24q11}
先作部分分式分解：
\[
\frac1{(n+5)(n+6)}=\frac1{n+5}-\frac1{n+6}.
\]
写出有限部分和后，中间项逐项抵消，得到
\[
S_N=\frac1{6}-\frac1{N+6},\qquad S=\frac{1}{6}.
\]
这种消去应当先对有限和进行；不能把两个分别发散的调和型无穷级数直接相减。余项精确等于 $1/(N+a+1)$，从而还可读出收敛速度为 $O(1/N)$。
\end{solve}

\begin{solve}[习题~\ref{ex:c24q12}：裂项求和]\phantomsection\label{ans:c24q12}
先作部分分式分解：
\[
\frac1{(n+6)(n+7)}=\frac1{n+6}-\frac1{n+7}.
\]
写出有限部分和后，中间项逐项抵消，得到
\[
S_N=\frac1{7}-\frac1{N+7},\qquad S=\frac{1}{7}.
\]
这种消去应当先对有限和进行；不能把两个分别发散的调和型无穷级数直接相减。余项精确等于 $1/(N+a+1)$，从而还可读出收敛速度为 $O(1/N)$。
\end{solve}

\begin{solve}[习题~\ref{ex:c24q13}：阶乘压过指数]\phantomsection\label{ans:c24q13}
记 $u_n=1^n/n!$，则
\[
\frac{u_{n+1}}{u_n}=\frac{1}{n+1}\longrightarrow0<1.
\]
由比值判别法级数收敛。当 $n+1\ge 2$ 时，比值不超过 $1/2$，所以在足够大的 $N$ 后有
\[
\sum_{n=N+1}^\infty u_n\le u_{N+1}\sum_{j=0}^\infty2^{-j}=2u_{N+1}.
\]
结合指数函数的幂级数可进一步识别总和为 $e^{1}$。这里收敛性由比值判别独立保证，不需要先把所求和认成一个已知函数。
\end{solve}

\begin{solve}[习题~\ref{ex:c24q14}：阶乘压过指数]\phantomsection\label{ans:c24q14}
记 $u_n=2^n/n!$，则
\[
\frac{u_{n+1}}{u_n}=\frac{2}{n+1}\longrightarrow0<1.
\]
由比值判别法级数收敛。当 $n+1\ge 4$ 时，比值不超过 $1/2$，所以在足够大的 $N$ 后有
\[
\sum_{n=N+1}^\infty u_n\le u_{N+1}\sum_{j=0}^\infty2^{-j}=2u_{N+1}.
\]
结合指数函数的幂级数可进一步识别总和为 $e^{2}$。这里收敛性由比值判别独立保证，不需要先把所求和认成一个已知函数。
\end{solve}

\begin{solve}[习题~\ref{ex:c24q15}：阶乘压过指数]\phantomsection\label{ans:c24q15}
记 $u_n=3^n/n!$，则
\[
\frac{u_{n+1}}{u_n}=\frac{3}{n+1}\longrightarrow0<1.
\]
由比值判别法级数收敛。当 $n+1\ge 6$ 时，比值不超过 $1/2$，所以在足够大的 $N$ 后有
\[
\sum_{n=N+1}^\infty u_n\le u_{N+1}\sum_{j=0}^\infty2^{-j}=2u_{N+1}.
\]
结合指数函数的幂级数可进一步识别总和为 $e^{3}$。这里收敛性由比值判别独立保证，不需要先把所求和认成一个已知函数。
\end{solve}

\begin{solve}[习题~\ref{ex:c24q16}：阶乘压过指数]\phantomsection\label{ans:c24q16}
记 $u_n=4^n/n!$，则
\[
\frac{u_{n+1}}{u_n}=\frac{4}{n+1}\longrightarrow0<1.
\]
由比值判别法级数收敛。当 $n+1\ge 8$ 时，比值不超过 $1/2$，所以在足够大的 $N$ 后有
\[
\sum_{n=N+1}^\infty u_n\le u_{N+1}\sum_{j=0}^\infty2^{-j}=2u_{N+1}.
\]
结合指数函数的幂级数可进一步识别总和为 $e^{4}$。这里收敛性由比值判别独立保证，不需要先把所求和认成一个已知函数。
\end{solve}

\begin{solve}[习题~\ref{ex:c24q17}：阶乘压过指数]\phantomsection\label{ans:c24q17}
记 $u_n=5^n/n!$，则
\[
\frac{u_{n+1}}{u_n}=\frac{5}{n+1}\longrightarrow0<1.
\]
由比值判别法级数收敛。当 $n+1\ge 10$ 时，比值不超过 $1/2$，所以在足够大的 $N$ 后有
\[
\sum_{n=N+1}^\infty u_n\le u_{N+1}\sum_{j=0}^\infty2^{-j}=2u_{N+1}.
\]
结合指数函数的幂级数可进一步识别总和为 $e^{5}$。这里收敛性由比值判别独立保证，不需要先把所求和认成一个已知函数。
\end{solve}

\begin{solve}[习题~\ref{ex:c24q18}：阶乘压过指数]\phantomsection\label{ans:c24q18}
记 $u_n=6^n/n!$，则
\[
\frac{u_{n+1}}{u_n}=\frac{6}{n+1}\longrightarrow0<1.
\]
由比值判别法级数收敛。当 $n+1\ge 12$ 时，比值不超过 $1/2$，所以在足够大的 $N$ 后有
\[
\sum_{n=N+1}^\infty u_n\le u_{N+1}\sum_{j=0}^\infty2^{-j}=2u_{N+1}.
\]
结合指数函数的幂级数可进一步识别总和为 $e^{6}$。这里收敛性由比值判别独立保证，不需要先把所求和认成一个已知函数。
\end{solve}

\begin{solve}[习题~\ref{ex:c24q19}：极限比较与主导次数]\phantomsection\label{ans:c24q19}
通项为正，与收敛的 $\sum1/n^2$ 比较。比值的极限为
\[
\lim_{n\to\infty}\frac{(1n+1)/(n^3+1)}{1/n^2}=\lim_{n\to\infty}\frac{1+1/n}{1+1/n^3}=1.
\]
极限为正的有限数，极限比较判别说明两个正项级数同敛散，因此原级数收敛。仅有通项趋零不足以得到这一结论；此外，比较判别只能给出敛散性，不能据此声称原级数的总和等于比较级数总和乘以该极限。
\end{solve}

\begin{solve}[习题~\ref{ex:c24q20}：极限比较与主导次数]\phantomsection\label{ans:c24q20}
通项为正，与收敛的 $\sum1/n^2$ 比较。比值的极限为
\[
\lim_{n\to\infty}\frac{(2n+1)/(n^3+2)}{1/n^2}=\lim_{n\to\infty}\frac{2+1/n}{1+2/n^3}=2.
\]
极限为正的有限数，极限比较判别说明两个正项级数同敛散，因此原级数收敛。仅有通项趋零不足以得到这一结论；此外，比较判别只能给出敛散性，不能据此声称原级数的总和等于比较级数总和乘以该极限。
\end{solve}

\begin{solve}[习题~\ref{ex:c24q21}：极限比较与主导次数]\phantomsection\label{ans:c24q21}
通项为正，与收敛的 $\sum1/n^2$ 比较。比值的极限为
\[
\lim_{n\to\infty}\frac{(3n+1)/(n^3+3)}{1/n^2}=\lim_{n\to\infty}\frac{3+1/n}{1+3/n^3}=3.
\]
极限为正的有限数，极限比较判别说明两个正项级数同敛散，因此原级数收敛。仅有通项趋零不足以得到这一结论；此外，比较判别只能给出敛散性，不能据此声称原级数的总和等于比较级数总和乘以该极限。
\end{solve}

\begin{solve}[习题~\ref{ex:c24q22}：极限比较与主导次数]\phantomsection\label{ans:c24q22}
通项为正，与收敛的 $\sum1/n^2$ 比较。比值的极限为
\[
\lim_{n\to\infty}\frac{(4n+1)/(n^3+4)}{1/n^2}=\lim_{n\to\infty}\frac{4+1/n}{1+4/n^3}=4.
\]
极限为正的有限数，极限比较判别说明两个正项级数同敛散，因此原级数收敛。仅有通项趋零不足以得到这一结论；此外，比较判别只能给出敛散性，不能据此声称原级数的总和等于比较级数总和乘以该极限。
\end{solve}

\begin{solve}[习题~\ref{ex:c24q23}：极限比较与主导次数]\phantomsection\label{ans:c24q23}
通项为正，与收敛的 $\sum1/n^2$ 比较。比值的极限为
\[
\lim_{n\to\infty}\frac{(5n+1)/(n^3+5)}{1/n^2}=\lim_{n\to\infty}\frac{5+1/n}{1+5/n^3}=5.
\]
极限为正的有限数，极限比较判别说明两个正项级数同敛散，因此原级数收敛。仅有通项趋零不足以得到这一结论；此外，比较判别只能给出敛散性，不能据此声称原级数的总和等于比较级数总和乘以该极限。
\end{solve}

\begin{solve}[习题~\ref{ex:c24q24}：极限比较与主导次数]\phantomsection\label{ans:c24q24}
通项为正，与收敛的 $\sum1/n^2$ 比较。比值的极限为
\[
\lim_{n\to\infty}\frac{(6n+1)/(n^3+6)}{1/n^2}=\lim_{n\to\infty}\frac{6+1/n}{1+6/n^3}=6.
\]
极限为正的有限数，极限比较判别说明两个正项级数同敛散，因此原级数收敛。仅有通项趋零不足以得到这一结论；此外，比较判别只能给出敛散性，不能据此声称原级数的总和等于比较级数总和乘以该极限。
\end{solve}

\clearpage
\section{级数敛散性的进一步辨析}\label{sec:ch25}
比值或根值判别得到1时，不是“收敛”，也不是“发散”，而是该判别没有给出结论。需要转用比较、积分判别或分组法。改变有限项不影响敛散性，但重排条件收敛级数可能改变其和。

\textbf{资源对应：}CLP-2 \S~3.3、\S~3.4 的扩展训练。本章共 24 道编号习题。
\subsection{习题}
\begin{example}[调和级数的对数修正]\label{ex:c25q1}
判断 $\displaystyle\sum_{n=2}^\infty\frac1{n(\log n)^{\frac{1}{2}}}$ 的敛散性。
\par\textit{改编训练。}\par\hyperref[ans:c25q1]{解答见第~\pageref{ans:c25q1} 页。}
\end{example}
\begin{example}[调和级数的对数修正]\label{ex:c25q2}
判断 $\displaystyle\sum_{n=2}^\infty\frac1{n(\log n)^{1}}$ 的敛散性。
\par\textit{改编训练。}\par\hyperref[ans:c25q2]{解答见第~\pageref{ans:c25q2} 页。}
\end{example}
\begin{example}[调和级数的对数修正]\label{ex:c25q3}
判断 $\displaystyle\sum_{n=2}^\infty\frac1{n(\log n)^{\frac{3}{2}}}$ 的敛散性。
\par\textit{改编训练。}\par\hyperref[ans:c25q3]{解答见第~\pageref{ans:c25q3} 页。}
\end{example}
\begin{example}[调和级数的对数修正]\label{ex:c25q4}
判断 $\displaystyle\sum_{n=2}^\infty\frac1{n(\log n)^{2}}$ 的敛散性。
\par\textit{改编训练。}\par\hyperref[ans:c25q4]{解答见第~\pageref{ans:c25q4} 页。}
\end{example}
\begin{example}[调和级数的对数修正]\label{ex:c25q5}
判断 $\displaystyle\sum_{n=2}^\infty\frac1{n(\log n)^{3}}$ 的敛散性。
\par\textit{改编训练。}\par\hyperref[ans:c25q5]{解答见第~\pageref{ans:c25q5} 页。}
\end{example}
\begin{example}[调和级数的对数修正]\label{ex:c25q6}
判断 $\displaystyle\sum_{n=2}^\infty\frac1{n(\log n)^{4}}$ 的敛散性。
\par\textit{改编训练。}\par\hyperref[ans:c25q6]{解答见第~\pageref{ans:c25q6} 页。}
\end{example}
\begin{example}[根值判别的直接使用]\label{ex:c25q7}
判断 $\displaystyle\sum_{n=1}^\infty\left(\frac{1n+1}{2n+1}\right)^n$ 的敛散性。
\par\textit{改编训练。}\par\hyperref[ans:c25q7]{解答见第~\pageref{ans:c25q7} 页。}
\end{example}
\begin{example}[根值判别的直接使用]\label{ex:c25q8}
判断 $\displaystyle\sum_{n=1}^\infty\left(\frac{2n+1}{3n+1}\right)^n$ 的敛散性。
\par\textit{改编训练。}\par\hyperref[ans:c25q8]{解答见第~\pageref{ans:c25q8} 页。}
\end{example}
\begin{example}[根值判别的直接使用]\label{ex:c25q9}
判断 $\displaystyle\sum_{n=1}^\infty\left(\frac{3n+1}{4n+1}\right)^n$ 的敛散性。
\par\textit{改编训练。}\par\hyperref[ans:c25q9]{解答见第~\pageref{ans:c25q9} 页。}
\end{example}
\begin{example}[根值判别的直接使用]\label{ex:c25q10}
判断 $\displaystyle\sum_{n=1}^\infty\left(\frac{2n+1}{1n+1}\right)^n$ 的敛散性。
\par\textit{改编训练。}\par\hyperref[ans:c25q10]{解答见第~\pageref{ans:c25q10} 页。}
\end{example}
\begin{example}[根值判别的直接使用]\label{ex:c25q11}
判断 $\displaystyle\sum_{n=1}^\infty\left(\frac{3n+1}{2n+1}\right)^n$ 的敛散性。
\par\textit{改编训练。}\par\hyperref[ans:c25q11]{解答见第~\pageref{ans:c25q11} 页。}
\end{example}
\begin{example}[根值判别的直接使用]\label{ex:c25q12}
判断 $\displaystyle\sum_{n=1}^\infty\left(\frac{4n+1}{3n+1}\right)^n$ 的敛散性。
\par\textit{改编训练。}\par\hyperref[ans:c25q12]{解答见第~\pageref{ans:c25q12} 页。}
\end{example}
\begin{example}[通项趋零与调和型发散]\label{ex:c25q13}
判断 $\displaystyle\sum_{n=1}^\infty\frac{1n+1}{n^2+1}$，并说明通项的极限。
\par\textit{改编训练。}\par\hyperref[ans:c25q13]{解答见第~\pageref{ans:c25q13} 页。}
\end{example}
\begin{example}[通项趋零与调和型发散]\label{ex:c25q14}
判断 $\displaystyle\sum_{n=1}^\infty\frac{2n+1}{n^2+1}$，并说明通项的极限。
\par\textit{改编训练。}\par\hyperref[ans:c25q14]{解答见第~\pageref{ans:c25q14} 页。}
\end{example}
\begin{example}[通项趋零与调和型发散]\label{ex:c25q15}
判断 $\displaystyle\sum_{n=1}^\infty\frac{3n+1}{n^2+1}$，并说明通项的极限。
\par\textit{改编训练。}\par\hyperref[ans:c25q15]{解答见第~\pageref{ans:c25q15} 页。}
\end{example}
\begin{example}[通项趋零与调和型发散]\label{ex:c25q16}
判断 $\displaystyle\sum_{n=1}^\infty\frac{4n+1}{n^2+1}$，并说明通项的极限。
\par\textit{改编训练。}\par\hyperref[ans:c25q16]{解答见第~\pageref{ans:c25q16} 页。}
\end{example}
\begin{example}[通项趋零与调和型发散]\label{ex:c25q17}
判断 $\displaystyle\sum_{n=1}^\infty\frac{5n+1}{n^2+1}$，并说明通项的极限。
\par\textit{改编训练。}\par\hyperref[ans:c25q17]{解答见第~\pageref{ans:c25q17} 页。}
\end{example}
\begin{example}[通项趋零与调和型发散]\label{ex:c25q18}
判断 $\displaystyle\sum_{n=1}^\infty\frac{6n+1}{n^2+1}$，并说明通项的极限。
\par\textit{改编训练。}\par\hyperref[ans:c25q18]{解答见第~\pageref{ans:c25q18} 页。}
\end{example}
\begin{example}[正项级数的尾和积分界]\label{ex:c25q19}
对 $S=\sum_{n=1}^\infty1/n^{2}$，给出保留到 $n=N$ 后的误差上下界。
\par\textit{改编训练。}\par\hyperref[ans:c25q19]{解答见第~\pageref{ans:c25q19} 页。}
\end{example}
\begin{example}[正项级数的尾和积分界]\label{ex:c25q20}
对 $S=\sum_{n=1}^\infty1/n^{3}$，给出保留到 $n=N$ 后的误差上下界。
\par\textit{改编训练。}\par\hyperref[ans:c25q20]{解答见第~\pageref{ans:c25q20} 页。}
\end{example}
\begin{example}[正项级数的尾和积分界]\label{ex:c25q21}
对 $S=\sum_{n=1}^\infty1/n^{4}$，给出保留到 $n=N$ 后的误差上下界。
\par\textit{改编训练。}\par\hyperref[ans:c25q21]{解答见第~\pageref{ans:c25q21} 页。}
\end{example}
\begin{example}[正项级数的尾和积分界]\label{ex:c25q22}
对 $S=\sum_{n=1}^\infty1/n^{5}$，给出保留到 $n=N$ 后的误差上下界。
\par\textit{改编训练。}\par\hyperref[ans:c25q22]{解答见第~\pageref{ans:c25q22} 页。}
\end{example}
\begin{example}[正项级数的尾和积分界]\label{ex:c25q23}
对 $S=\sum_{n=1}^\infty1/n^{6}$，给出保留到 $n=N$ 后的误差上下界。
\par\textit{改编训练。}\par\hyperref[ans:c25q23]{解答见第~\pageref{ans:c25q23} 页。}
\end{example}
\begin{example}[正项级数的尾和积分界]\label{ex:c25q24}
对 $S=\sum_{n=1}^\infty1/n^{7}$，给出保留到 $n=N$ 后的误差上下界。
\par\textit{改编训练。}\par\hyperref[ans:c25q24]{解答见第~\pageref{ans:c25q24} 页。}
\end{example}
\subsection{解答}
\begin{solve}[习题~\ref{ex:c25q1}：调和级数的对数修正]\phantomsection\label{ans:c25q1}
函数 $f(x)=1/(x(\log x)^{\frac{1}{2}})$ 在 $x\ge2$ 上为正、连续且递减，适用积分判别。作 $u=\log x$，
\[
\int_2^\infty\frac{dx}{x(\log x)^{\frac{1}{2}}}=\int_{\log2}^\infty u^{-\frac{1}{2}}du.
\]
幂指数不超过1，因此积分与级数都发散。积分判别给出的是同敛散性，不表示级数和等于积分值；二者通常存在一个非零差。
\end{solve}

\begin{solve}[习题~\ref{ex:c25q2}：调和级数的对数修正]\phantomsection\label{ans:c25q2}
函数 $f(x)=1/(x(\log x)^{1})$ 在 $x\ge2$ 上为正、连续且递减，适用积分判别。作 $u=\log x$，
\[
\int_2^\infty\frac{dx}{x(\log x)^{1}}=\int_{\log2}^\infty u^{-1}du.
\]
幂指数不超过1，因此积分与级数都发散。积分判别给出的是同敛散性，不表示级数和等于积分值；二者通常存在一个非零差。
\end{solve}

\begin{solve}[习题~\ref{ex:c25q3}：调和级数的对数修正]\phantomsection\label{ans:c25q3}
函数 $f(x)=1/(x(\log x)^{\frac{3}{2}})$ 在 $x\ge2$ 上为正、连续且递减，适用积分判别。作 $u=\log x$，
\[
\int_2^\infty\frac{dx}{x(\log x)^{\frac{3}{2}}}=\int_{\log2}^\infty u^{-\frac{3}{2}}du.
\]
幂指数大于1，因此积分与级数都收敛。积分判别给出的是同敛散性，不表示级数和等于积分值；二者通常存在一个非零差。
\end{solve}

\begin{solve}[习题~\ref{ex:c25q4}：调和级数的对数修正]\phantomsection\label{ans:c25q4}
函数 $f(x)=1/(x(\log x)^{2})$ 在 $x\ge2$ 上为正、连续且递减，适用积分判别。作 $u=\log x$，
\[
\int_2^\infty\frac{dx}{x(\log x)^{2}}=\int_{\log2}^\infty u^{-2}du.
\]
幂指数大于1，因此积分与级数都收敛。积分判别给出的是同敛散性，不表示级数和等于积分值；二者通常存在一个非零差。
\end{solve}

\begin{solve}[习题~\ref{ex:c25q5}：调和级数的对数修正]\phantomsection\label{ans:c25q5}
函数 $f(x)=1/(x(\log x)^{3})$ 在 $x\ge2$ 上为正、连续且递减，适用积分判别。作 $u=\log x$，
\[
\int_2^\infty\frac{dx}{x(\log x)^{3}}=\int_{\log2}^\infty u^{-3}du.
\]
幂指数大于1，因此积分与级数都收敛。积分判别给出的是同敛散性，不表示级数和等于积分值；二者通常存在一个非零差。
\end{solve}

\begin{solve}[习题~\ref{ex:c25q6}：调和级数的对数修正]\phantomsection\label{ans:c25q6}
函数 $f(x)=1/(x(\log x)^{4})$ 在 $x\ge2$ 上为正、连续且递减，适用积分判别。作 $u=\log x$，
\[
\int_2^\infty\frac{dx}{x(\log x)^{4}}=\int_{\log2}^\infty u^{-4}du.
\]
幂指数大于1，因此积分与级数都收敛。积分判别给出的是同敛散性，不表示级数和等于积分值；二者通常存在一个非零差。
\end{solve}

\begin{solve}[习题~\ref{ex:c25q7}：根值判别的直接使用]\phantomsection\label{ans:c25q7}
通项非负，取 $n$ 次方根后恰好消去外层幂：
\[
\sqrt[n]{u_n}=\frac{1n+1}{2n+1}\longrightarrow \frac{1}{2}.
\]
极限小于1，根值判别说明级数收敛。可在极限值与1之间选一个 $q$，从某项起用 $q^n$ 控制通项。有 $n$ 次幂的正项结构往往优先考虑根值判别，比计算相邻项的完整比值简洁。
\end{solve}

\begin{solve}[习题~\ref{ex:c25q8}：根值判别的直接使用]\phantomsection\label{ans:c25q8}
通项非负，取 $n$ 次方根后恰好消去外层幂：
\[
\sqrt[n]{u_n}=\frac{2n+1}{3n+1}\longrightarrow \frac{2}{3}.
\]
极限小于1，根值判别说明级数收敛。可在极限值与1之间选一个 $q$，从某项起用 $q^n$ 控制通项。有 $n$ 次幂的正项结构往往优先考虑根值判别，比计算相邻项的完整比值简洁。
\end{solve}

\begin{solve}[习题~\ref{ex:c25q9}：根值判别的直接使用]\phantomsection\label{ans:c25q9}
通项非负，取 $n$ 次方根后恰好消去外层幂：
\[
\sqrt[n]{u_n}=\frac{3n+1}{4n+1}\longrightarrow \frac{3}{4}.
\]
极限小于1，根值判别说明级数收敛。可在极限值与1之间选一个 $q$，从某项起用 $q^n$ 控制通项。有 $n$ 次幂的正项结构往往优先考虑根值判别，比计算相邻项的完整比值简洁。
\end{solve}

\begin{solve}[习题~\ref{ex:c25q10}：根值判别的直接使用]\phantomsection\label{ans:c25q10}
通项非负，取 $n$ 次方根后恰好消去外层幂：
\[
\sqrt[n]{u_n}=\frac{2n+1}{1n+1}\longrightarrow 2.
\]
极限大于1，根值判别说明级数发散；事实上通项从某项起至少按某个大于1的几何比增长，根本不趋零。有 $n$ 次幂的正项结构往往优先考虑根值判别，比计算相邻项的完整比值简洁。
\end{solve}

\begin{solve}[习题~\ref{ex:c25q11}：根值判别的直接使用]\phantomsection\label{ans:c25q11}
通项非负，取 $n$ 次方根后恰好消去外层幂：
\[
\sqrt[n]{u_n}=\frac{3n+1}{2n+1}\longrightarrow \frac{3}{2}.
\]
极限大于1，根值判别说明级数发散；事实上通项从某项起至少按某个大于1的几何比增长，根本不趋零。有 $n$ 次幂的正项结构往往优先考虑根值判别，比计算相邻项的完整比值简洁。
\end{solve}

\begin{solve}[习题~\ref{ex:c25q12}：根值判别的直接使用]\phantomsection\label{ans:c25q12}
通项非负，取 $n$ 次方根后恰好消去外层幂：
\[
\sqrt[n]{u_n}=\frac{4n+1}{3n+1}\longrightarrow \frac{4}{3}.
\]
极限大于1，根值判别说明级数发散；事实上通项从某项起至少按某个大于1的几何比增长，根本不趋零。有 $n$ 次幂的正项结构往往优先考虑根值判别，比计算相邻项的完整比值简洁。
\end{solve}

\begin{solve}[习题~\ref{ex:c25q13}：通项趋零与调和型发散]\phantomsection\label{ans:c25q13}
通项的分子次数比分母低1，所以通项趋零。但与调和级数比较，
\[
\lim_{n\to\infty}\frac{(1n+1)/(n^2+1)}{1/n}=\lim_{n\to\infty}\frac{1+1/n}{1+1/n^2}=1>0.
\]
由极限比较，原级数与调和级数同为发散。无限多个趋零小量的累积可以无界，正是“数列趋零”与“级数收敛”之间的根本区别。
\end{solve}

\begin{solve}[习题~\ref{ex:c25q14}：通项趋零与调和型发散]\phantomsection\label{ans:c25q14}
通项的分子次数比分母低1，所以通项趋零。但与调和级数比较，
\[
\lim_{n\to\infty}\frac{(2n+1)/(n^2+1)}{1/n}=\lim_{n\to\infty}\frac{2+1/n}{1+1/n^2}=2>0.
\]
由极限比较，原级数与调和级数同为发散。无限多个趋零小量的累积可以无界，正是“数列趋零”与“级数收敛”之间的根本区别。
\end{solve}

\begin{solve}[习题~\ref{ex:c25q15}：通项趋零与调和型发散]\phantomsection\label{ans:c25q15}
通项的分子次数比分母低1，所以通项趋零。但与调和级数比较，
\[
\lim_{n\to\infty}\frac{(3n+1)/(n^2+1)}{1/n}=\lim_{n\to\infty}\frac{3+1/n}{1+1/n^2}=3>0.
\]
由极限比较，原级数与调和级数同为发散。无限多个趋零小量的累积可以无界，正是“数列趋零”与“级数收敛”之间的根本区别。
\end{solve}

\begin{solve}[习题~\ref{ex:c25q16}：通项趋零与调和型发散]\phantomsection\label{ans:c25q16}
通项的分子次数比分母低1，所以通项趋零。但与调和级数比较，
\[
\lim_{n\to\infty}\frac{(4n+1)/(n^2+1)}{1/n}=\lim_{n\to\infty}\frac{4+1/n}{1+1/n^2}=4>0.
\]
由极限比较，原级数与调和级数同为发散。无限多个趋零小量的累积可以无界，正是“数列趋零”与“级数收敛”之间的根本区别。
\end{solve}

\begin{solve}[习题~\ref{ex:c25q17}：通项趋零与调和型发散]\phantomsection\label{ans:c25q17}
通项的分子次数比分母低1，所以通项趋零。但与调和级数比较，
\[
\lim_{n\to\infty}\frac{(5n+1)/(n^2+1)}{1/n}=\lim_{n\to\infty}\frac{5+1/n}{1+1/n^2}=5>0.
\]
由极限比较，原级数与调和级数同为发散。无限多个趋零小量的累积可以无界，正是“数列趋零”与“级数收敛”之间的根本区别。
\end{solve}

\begin{solve}[习题~\ref{ex:c25q18}：通项趋零与调和型发散]\phantomsection\label{ans:c25q18}
通项的分子次数比分母低1，所以通项趋零。但与调和级数比较，
\[
\lim_{n\to\infty}\frac{(6n+1)/(n^2+1)}{1/n}=\lim_{n\to\infty}\frac{6+1/n}{1+1/n^2}=6>0.
\]
由极限比较，原级数与调和级数同为发散。无限多个趋零小量的累积可以无界，正是“数列趋零”与“级数收敛”之间的根本区别。
\end{solve}

\begin{solve}[习题~\ref{ex:c25q19}：正项级数的尾和积分界]\phantomsection\label{ans:c25q19}
函数 $f(x)=x^{-2}$ 正且递减，因此每个矩形面积可与左右相邻区间的积分比较。累加后得到
\[
\int_{N+1}^\infty x^{-2}dx\le R_N=\sum_{n=N+1}^\infty\frac1{n^{2}}\le\int_N^\infty x^{-2}dx.
\]
计算积分即得
\[
\frac1{1(N+1)^{1}}\le R_N\le\frac1{1N^{1}}.
\]
上下界同阶，既给出误差保证，也说明尾和的实际衰减级别。把第一项 $1/(N+1)^{a+1}$ 当成整个正项尾和的上界是不正确的。
\end{solve}

\begin{solve}[习题~\ref{ex:c25q20}：正项级数的尾和积分界]\phantomsection\label{ans:c25q20}
函数 $f(x)=x^{-3}$ 正且递减，因此每个矩形面积可与左右相邻区间的积分比较。累加后得到
\[
\int_{N+1}^\infty x^{-3}dx\le R_N=\sum_{n=N+1}^\infty\frac1{n^{3}}\le\int_N^\infty x^{-3}dx.
\]
计算积分即得
\[
\frac1{2(N+1)^{2}}\le R_N\le\frac1{2N^{2}}.
\]
上下界同阶，既给出误差保证，也说明尾和的实际衰减级别。把第一项 $1/(N+1)^{a+1}$ 当成整个正项尾和的上界是不正确的。
\end{solve}

\begin{solve}[习题~\ref{ex:c25q21}：正项级数的尾和积分界]\phantomsection\label{ans:c25q21}
函数 $f(x)=x^{-4}$ 正且递减，因此每个矩形面积可与左右相邻区间的积分比较。累加后得到
\[
\int_{N+1}^\infty x^{-4}dx\le R_N=\sum_{n=N+1}^\infty\frac1{n^{4}}\le\int_N^\infty x^{-4}dx.
\]
计算积分即得
\[
\frac1{3(N+1)^{3}}\le R_N\le\frac1{3N^{3}}.
\]
上下界同阶，既给出误差保证，也说明尾和的实际衰减级别。把第一项 $1/(N+1)^{a+1}$ 当成整个正项尾和的上界是不正确的。
\end{solve}

\begin{solve}[习题~\ref{ex:c25q22}：正项级数的尾和积分界]\phantomsection\label{ans:c25q22}
函数 $f(x)=x^{-5}$ 正且递减，因此每个矩形面积可与左右相邻区间的积分比较。累加后得到
\[
\int_{N+1}^\infty x^{-5}dx\le R_N=\sum_{n=N+1}^\infty\frac1{n^{5}}\le\int_N^\infty x^{-5}dx.
\]
计算积分即得
\[
\frac1{4(N+1)^{4}}\le R_N\le\frac1{4N^{4}}.
\]
上下界同阶，既给出误差保证，也说明尾和的实际衰减级别。把第一项 $1/(N+1)^{a+1}$ 当成整个正项尾和的上界是不正确的。
\end{solve}

\begin{solve}[习题~\ref{ex:c25q23}：正项级数的尾和积分界]\phantomsection\label{ans:c25q23}
函数 $f(x)=x^{-6}$ 正且递减，因此每个矩形面积可与左右相邻区间的积分比较。累加后得到
\[
\int_{N+1}^\infty x^{-6}dx\le R_N=\sum_{n=N+1}^\infty\frac1{n^{6}}\le\int_N^\infty x^{-6}dx.
\]
计算积分即得
\[
\frac1{5(N+1)^{5}}\le R_N\le\frac1{5N^{5}}.
\]
上下界同阶，既给出误差保证，也说明尾和的实际衰减级别。把第一项 $1/(N+1)^{a+1}$ 当成整个正项尾和的上界是不正确的。
\end{solve}

\begin{solve}[习题~\ref{ex:c25q24}：正项级数的尾和积分界]\phantomsection\label{ans:c25q24}
函数 $f(x)=x^{-7}$ 正且递减，因此每个矩形面积可与左右相邻区间的积分比较。累加后得到
\[
\int_{N+1}^\infty x^{-7}dx\le R_N=\sum_{n=N+1}^\infty\frac1{n^{7}}\le\int_N^\infty x^{-7}dx.
\]
计算积分即得
\[
\frac1{6(N+1)^{6}}\le R_N\le\frac1{6N^{6}}.
\]
上下界同阶，既给出误差保证，也说明尾和的实际衰减级别。把第一项 $1/(N+1)^{a+1}$ 当成整个正项尾和的上界是不正确的。
\end{solve}

\clearpage
\section{交错级数与幂级数}\label{sec:ch26}
绝对收敛（absolute convergence）比通常的收敛更强。幂级数在收敛半径以内可以逐项求导和积分，但端点必须单独代入，不能仅凭半径来决定。

\textbf{资源对应：}CLP-2 \S~3.4--\S~3.6。本章共 25 道编号习题。
\subsection{习题}
\begin{example}[交错收敛不推出绝对收敛]\label{ex:c26q1}
判断：若 $\sum_{n=1}^\infty(-1)^{n+1}b_n$ 收敛，则 $\sum_{n=1}^\infty b_n$ 也收敛。
\par\textit{选译。来源：CLP-2 \S~3.4，2015A 判断题。}\par\hyperref[ans:c26q1]{解答见第~\pageref{ans:c26q1} 页。}
\end{example}
\begin{example}[交错幂级数的端点]\label{ex:c26q2}
求 $\displaystyle\sum_{n=1}^\infty\frac{x^n}{n^{\frac{1}{3}}}$ 的收敛半径与收敛区间，并说明绝对收敛范围。
\par\textit{改编训练。}\par\hyperref[ans:c26q2]{解答见第~\pageref{ans:c26q2} 页。}
\end{example}
\begin{example}[交错幂级数的端点]\label{ex:c26q3}
求 $\displaystyle\sum_{n=1}^\infty\frac{x^n}{n^{\frac{1}{2}}}$ 的收敛半径与收敛区间，并说明绝对收敛范围。
\par\textit{改编训练。}\par\hyperref[ans:c26q3]{解答见第~\pageref{ans:c26q3} 页。}
\end{example}
\begin{example}[交错幂级数的端点]\label{ex:c26q4}
求 $\displaystyle\sum_{n=1}^\infty\frac{x^n}{n^{\frac{2}{3}}}$ 的收敛半径与收敛区间，并说明绝对收敛范围。
\par\textit{改编训练。}\par\hyperref[ans:c26q4]{解答见第~\pageref{ans:c26q4} 页。}
\end{example}
\begin{example}[交错幂级数的端点]\label{ex:c26q5}
求 $\displaystyle\sum_{n=1}^\infty\frac{x^n}{n^{1}}$ 的收敛半径与收敛区间，并说明绝对收敛范围。
\par\textit{改编训练。}\par\hyperref[ans:c26q5]{解答见第~\pageref{ans:c26q5} 页。}
\end{example}
\begin{example}[交错幂级数的端点]\label{ex:c26q6}
求 $\displaystyle\sum_{n=1}^\infty\frac{x^n}{n^{\frac{3}{2}}}$ 的收敛半径与收敛区间，并说明绝对收敛范围。
\par\textit{改编训练。}\par\hyperref[ans:c26q6]{解答见第~\pageref{ans:c26q6} 页。}
\end{example}
\begin{example}[交错幂级数的端点]\label{ex:c26q7}
求 $\displaystyle\sum_{n=1}^\infty\frac{x^n}{n^{2}}$ 的收敛半径与收敛区间，并说明绝对收敛范围。
\par\textit{改编训练。}\par\hyperref[ans:c26q7]{解答见第~\pageref{ans:c26q7} 页。}
\end{example}
\begin{example}[几何级数导出的和函数]\label{ex:c26q8}
求 $\displaystyle\sum_{n=1}^\infty n\left(\frac x{1}\right)^n$ 的收敛域和和函数。
\par\textit{改编训练。}\par\hyperref[ans:c26q8]{解答见第~\pageref{ans:c26q8} 页。}
\end{example}
\begin{example}[几何级数导出的和函数]\label{ex:c26q9}
求 $\displaystyle\sum_{n=1}^\infty n\left(\frac x{2}\right)^n$ 的收敛域和和函数。
\par\textit{改编训练。}\par\hyperref[ans:c26q9]{解答见第~\pageref{ans:c26q9} 页。}
\end{example}
\begin{example}[几何级数导出的和函数]\label{ex:c26q10}
求 $\displaystyle\sum_{n=1}^\infty n\left(\frac x{3}\right)^n$ 的收敛域和和函数。
\par\textit{改编训练。}\par\hyperref[ans:c26q10]{解答见第~\pageref{ans:c26q10} 页。}
\end{example}
\begin{example}[几何级数导出的和函数]\label{ex:c26q11}
求 $\displaystyle\sum_{n=1}^\infty n\left(\frac x{4}\right)^n$ 的收敛域和和函数。
\par\textit{改编训练。}\par\hyperref[ans:c26q11]{解答见第~\pageref{ans:c26q11} 页。}
\end{example}
\begin{example}[几何级数导出的和函数]\label{ex:c26q12}
求 $\displaystyle\sum_{n=1}^\infty n\left(\frac x{5}\right)^n$ 的收敛域和和函数。
\par\textit{改编训练。}\par\hyperref[ans:c26q12]{解答见第~\pageref{ans:c26q12} 页。}
\end{example}
\begin{example}[几何级数导出的和函数]\label{ex:c26q13}
求 $\displaystyle\sum_{n=1}^\infty n\left(\frac x{6}\right)^n$ 的收敛域和和函数。
\par\textit{改编训练。}\par\hyperref[ans:c26q13]{解答见第~\pageref{ans:c26q13} 页。}
\end{example}
\begin{example}[对数和函数与端点]\label{ex:c26q14}
求 $\displaystyle\sum_{n=1}^\infty\frac{x^n}{n1^n}$ 的和函数与收敛区间。
\par\textit{改编训练。}\par\hyperref[ans:c26q14]{解答见第~\pageref{ans:c26q14} 页。}
\end{example}
\begin{example}[对数和函数与端点]\label{ex:c26q15}
求 $\displaystyle\sum_{n=1}^\infty\frac{x^n}{n2^n}$ 的和函数与收敛区间。
\par\textit{改编训练。}\par\hyperref[ans:c26q15]{解答见第~\pageref{ans:c26q15} 页。}
\end{example}
\begin{example}[对数和函数与端点]\label{ex:c26q16}
求 $\displaystyle\sum_{n=1}^\infty\frac{x^n}{n3^n}$ 的和函数与收敛区间。
\par\textit{改编训练。}\par\hyperref[ans:c26q16]{解答见第~\pageref{ans:c26q16} 页。}
\end{example}
\begin{example}[对数和函数与端点]\label{ex:c26q17}
求 $\displaystyle\sum_{n=1}^\infty\frac{x^n}{n4^n}$ 的和函数与收敛区间。
\par\textit{改编训练。}\par\hyperref[ans:c26q17]{解答见第~\pageref{ans:c26q17} 页。}
\end{example}
\begin{example}[对数和函数与端点]\label{ex:c26q18}
求 $\displaystyle\sum_{n=1}^\infty\frac{x^n}{n5^n}$ 的和函数与收敛区间。
\par\textit{改编训练。}\par\hyperref[ans:c26q18]{解答见第~\pageref{ans:c26q18} 页。}
\end{example}
\begin{example}[对数和函数与端点]\label{ex:c26q19}
求 $\displaystyle\sum_{n=1}^\infty\frac{x^n}{n6^n}$ 的和函数与收敛区间。
\par\textit{改编训练。}\par\hyperref[ans:c26q19]{解答见第~\pageref{ans:c26q19} 页。}
\end{example}
\begin{example}[交错级数的误差控制]\label{ex:c26q20}
用前 $N$ 项近似 $\displaystyle S=\sum_{n=1}^\infty\frac{(-1)^{n-1}}{n^{1}}$。给出误差上界，并给出误差小于 $10^{-3}$ 的一个充分条件。
\par\textit{改编训练。}\par\hyperref[ans:c26q20]{解答见第~\pageref{ans:c26q20} 页。}
\end{example}
\begin{example}[交错级数的误差控制]\label{ex:c26q21}
用前 $N$ 项近似 $\displaystyle S=\sum_{n=1}^\infty\frac{(-1)^{n-1}}{n^{2}}$。给出误差上界，并给出误差小于 $10^{-3}$ 的一个充分条件。
\par\textit{改编训练。}\par\hyperref[ans:c26q21]{解答见第~\pageref{ans:c26q21} 页。}
\end{example}
\begin{example}[交错级数的误差控制]\label{ex:c26q22}
用前 $N$ 项近似 $\displaystyle S=\sum_{n=1}^\infty\frac{(-1)^{n-1}}{n^{3}}$。给出误差上界，并给出误差小于 $10^{-3}$ 的一个充分条件。
\par\textit{改编训练。}\par\hyperref[ans:c26q22]{解答见第~\pageref{ans:c26q22} 页。}
\end{example}
\begin{example}[交错级数的误差控制]\label{ex:c26q23}
用前 $N$ 项近似 $\displaystyle S=\sum_{n=1}^\infty\frac{(-1)^{n-1}}{n^{4}}$。给出误差上界，并给出误差小于 $10^{-3}$ 的一个充分条件。
\par\textit{改编训练。}\par\hyperref[ans:c26q23]{解答见第~\pageref{ans:c26q23} 页。}
\end{example}
\begin{example}[交错级数的误差控制]\label{ex:c26q24}
用前 $N$ 项近似 $\displaystyle S=\sum_{n=1}^\infty\frac{(-1)^{n-1}}{n^{5}}$。给出误差上界，并给出误差小于 $10^{-3}$ 的一个充分条件。
\par\textit{改编训练。}\par\hyperref[ans:c26q24]{解答见第~\pageref{ans:c26q24} 页。}
\end{example}
\begin{example}[交错级数的误差控制]\label{ex:c26q25}
用前 $N$ 项近似 $\displaystyle S=\sum_{n=1}^\infty\frac{(-1)^{n-1}}{n^{6}}$。给出误差上界，并给出误差小于 $10^{-3}$ 的一个充分条件。
\par\textit{改编训练。}\par\hyperref[ans:c26q25]{解答见第~\pageref{ans:c26q25} 页。}
\end{example}
\subsection{解答}
\begin{solve}[习题~\ref{ex:c26q1}：交错收敛不推出绝对收敛]\phantomsection\label{ans:c26q1}
命题错误。取 $b_n=1/n$，它递减趋零，所以交错调和级数由 Leibniz 判别收敛；但调和级数发散。该例区分了条件收敛与绝对收敛。若已知的是 $\sum|b_n|$ 收敛，则可以推出加任意正负号后的级数都收敛。
\end{solve}

\begin{solve}[习题~\ref{ex:c26q2}：交错幂级数的端点]\phantomsection\label{ans:c26q2}
系数的 $n$ 次方根趋于 $1$，故收敛半径 $R=1$。当 $|x|<1$ 时，根值判别给出绝对收敛；当 $|x|>1$ 时通项不趋零。在 $x=1$ 处得到 $\sum1/n^{\frac{1}{3}}$，因幂指数不超过1而发散。在 $x=-1$ 处得到交错级数，$n^{-\frac{1}{3}}$ 单调趋零，故总是收敛。其绝对值级数发散，所以左端点条件收敛，收敛区间为 $[-1,1)$，绝对收敛区间为 $(-1,1)$。端点检查是半径计算之后的独立步骤。
\end{solve}

\begin{solve}[习题~\ref{ex:c26q3}：交错幂级数的端点]\phantomsection\label{ans:c26q3}
系数的 $n$ 次方根趋于 $1$，故收敛半径 $R=1$。当 $|x|<1$ 时，根值判别给出绝对收敛；当 $|x|>1$ 时通项不趋零。在 $x=1$ 处得到 $\sum1/n^{\frac{1}{2}}$，因幂指数不超过1而发散。在 $x=-1$ 处得到交错级数，$n^{-\frac{1}{2}}$ 单调趋零，故总是收敛。其绝对值级数发散，所以左端点条件收敛，收敛区间为 $[-1,1)$，绝对收敛区间为 $(-1,1)$。端点检查是半径计算之后的独立步骤。
\end{solve}

\begin{solve}[习题~\ref{ex:c26q4}：交错幂级数的端点]\phantomsection\label{ans:c26q4}
系数的 $n$ 次方根趋于 $1$，故收敛半径 $R=1$。当 $|x|<1$ 时，根值判别给出绝对收敛；当 $|x|>1$ 时通项不趋零。在 $x=1$ 处得到 $\sum1/n^{\frac{2}{3}}$，因幂指数不超过1而发散。在 $x=-1$ 处得到交错级数，$n^{-\frac{2}{3}}$ 单调趋零，故总是收敛。其绝对值级数发散，所以左端点条件收敛，收敛区间为 $[-1,1)$，绝对收敛区间为 $(-1,1)$。端点检查是半径计算之后的独立步骤。
\end{solve}

\begin{solve}[习题~\ref{ex:c26q5}：交错幂级数的端点]\phantomsection\label{ans:c26q5}
系数的 $n$ 次方根趋于 $1$，故收敛半径 $R=1$。当 $|x|<1$ 时，根值判别给出绝对收敛；当 $|x|>1$ 时通项不趋零。在 $x=1$ 处得到 $\sum1/n^{1}$，因幂指数不超过1而发散。在 $x=-1$ 处得到交错级数，$n^{-1}$ 单调趋零，故总是收敛。其绝对值级数发散，所以左端点条件收敛，收敛区间为 $[-1,1)$，绝对收敛区间为 $(-1,1)$。端点检查是半径计算之后的独立步骤。
\end{solve}

\begin{solve}[习题~\ref{ex:c26q6}：交错幂级数的端点]\phantomsection\label{ans:c26q6}
系数的 $n$ 次方根趋于 $1$，故收敛半径 $R=1$。当 $|x|<1$ 时，根值判别给出绝对收敛；当 $|x|>1$ 时通项不趋零。在 $x=1$ 处得到 $\sum1/n^{\frac{3}{2}}$，因幂指数大于1而收敛。在 $x=-1$ 处得到交错级数，$n^{-\frac{3}{2}}$ 单调趋零，故总是收敛。其绝对值级数也收敛，所以两个端点均绝对收敛，收敛区间为 $[-1,1]$。端点检查是半径计算之后的独立步骤。
\end{solve}

\begin{solve}[习题~\ref{ex:c26q7}：交错幂级数的端点]\phantomsection\label{ans:c26q7}
系数的 $n$ 次方根趋于 $1$，故收敛半径 $R=1$。当 $|x|<1$ 时，根值判别给出绝对收敛；当 $|x|>1$ 时通项不趋零。在 $x=1$ 处得到 $\sum1/n^{2}$，因幂指数大于1而收敛。在 $x=-1$ 处得到交错级数，$n^{-2}$ 单调趋零，故总是收敛。其绝对值级数也收敛，所以两个端点均绝对收敛，收敛区间为 $[-1,1]$。端点检查是半径计算之后的独立步骤。
\end{solve}

\begin{solve}[习题~\ref{ex:c26q8}：几何级数导出的和函数]\phantomsection\label{ans:c26q8}
令 $t=x/1$。在 $|t|<1$ 内，几何级数 $\sum_{n=0}^\infty t^n=1/(1-t)$ 可以逐项求导，再乘以 $t$，得到
\[
\sum_{n=1}^\infty nt^n=\frac{t}{(1-t)^2},\qquad S(x)=\frac{1x}{(1-x)^2}.
\]
所以 $|x|<1$ 时收敛。在 $x=\pm 1$ 时通项分别为 $n$ 与 $n(-1)^n$，均不趋零，故两端都发散。收敛域为 $(-1,1)$。逐项求导只在半径内部使用，不能把所得有理函数在其他点的有限值解释为原级数的和。
\end{solve}

\begin{solve}[习题~\ref{ex:c26q9}：几何级数导出的和函数]\phantomsection\label{ans:c26q9}
令 $t=x/2$。在 $|t|<1$ 内，几何级数 $\sum_{n=0}^\infty t^n=1/(1-t)$ 可以逐项求导，再乘以 $t$，得到
\[
\sum_{n=1}^\infty nt^n=\frac{t}{(1-t)^2},\qquad S(x)=\frac{2x}{(2-x)^2}.
\]
所以 $|x|<2$ 时收敛。在 $x=\pm 2$ 时通项分别为 $n$ 与 $n(-1)^n$，均不趋零，故两端都发散。收敛域为 $(-2,2)$。逐项求导只在半径内部使用，不能把所得有理函数在其他点的有限值解释为原级数的和。
\end{solve}

\begin{solve}[习题~\ref{ex:c26q10}：几何级数导出的和函数]\phantomsection\label{ans:c26q10}
令 $t=x/3$。在 $|t|<1$ 内，几何级数 $\sum_{n=0}^\infty t^n=1/(1-t)$ 可以逐项求导，再乘以 $t$，得到
\[
\sum_{n=1}^\infty nt^n=\frac{t}{(1-t)^2},\qquad S(x)=\frac{3x}{(3-x)^2}.
\]
所以 $|x|<3$ 时收敛。在 $x=\pm 3$ 时通项分别为 $n$ 与 $n(-1)^n$，均不趋零，故两端都发散。收敛域为 $(-3,3)$。逐项求导只在半径内部使用，不能把所得有理函数在其他点的有限值解释为原级数的和。
\end{solve}

\begin{solve}[习题~\ref{ex:c26q11}：几何级数导出的和函数]\phantomsection\label{ans:c26q11}
令 $t=x/4$。在 $|t|<1$ 内，几何级数 $\sum_{n=0}^\infty t^n=1/(1-t)$ 可以逐项求导，再乘以 $t$，得到
\[
\sum_{n=1}^\infty nt^n=\frac{t}{(1-t)^2},\qquad S(x)=\frac{4x}{(4-x)^2}.
\]
所以 $|x|<4$ 时收敛。在 $x=\pm 4$ 时通项分别为 $n$ 与 $n(-1)^n$，均不趋零，故两端都发散。收敛域为 $(-4,4)$。逐项求导只在半径内部使用，不能把所得有理函数在其他点的有限值解释为原级数的和。
\end{solve}

\begin{solve}[习题~\ref{ex:c26q12}：几何级数导出的和函数]\phantomsection\label{ans:c26q12}
令 $t=x/5$。在 $|t|<1$ 内，几何级数 $\sum_{n=0}^\infty t^n=1/(1-t)$ 可以逐项求导，再乘以 $t$，得到
\[
\sum_{n=1}^\infty nt^n=\frac{t}{(1-t)^2},\qquad S(x)=\frac{5x}{(5-x)^2}.
\]
所以 $|x|<5$ 时收敛。在 $x=\pm 5$ 时通项分别为 $n$ 与 $n(-1)^n$，均不趋零，故两端都发散。收敛域为 $(-5,5)$。逐项求导只在半径内部使用，不能把所得有理函数在其他点的有限值解释为原级数的和。
\end{solve}

\begin{solve}[习题~\ref{ex:c26q13}：几何级数导出的和函数]\phantomsection\label{ans:c26q13}
令 $t=x/6$。在 $|t|<1$ 内，几何级数 $\sum_{n=0}^\infty t^n=1/(1-t)$ 可以逐项求导，再乘以 $t$，得到
\[
\sum_{n=1}^\infty nt^n=\frac{t}{(1-t)^2},\qquad S(x)=\frac{6x}{(6-x)^2}.
\]
所以 $|x|<6$ 时收敛。在 $x=\pm 6$ 时通项分别为 $n$ 与 $n(-1)^n$，均不趋零，故两端都发散。收敛域为 $(-6,6)$。逐项求导只在半径内部使用，不能把所得有理函数在其他点的有限值解释为原级数的和。
\end{solve}

\begin{solve}[习题~\ref{ex:c26q14}：对数和函数与端点]\phantomsection\label{ans:c26q14}
令 $t=x/1$。在 $|t|<1$ 上对几何级数逐项积分，可得
\[
\sum_{n=1}^\infty\frac{t^n}n=-\log(1-t),\qquad S(x)=-\log(1-x/1).
\]
在右端点 $x=1$ 处得到调和级数，发散；左端点 $x=-1$ 处是交错调和级数，收敛。因此收敛区间为 $[-1,1)$。左端点的和为 $-\log2$，可通过 Abel 定理或对交错调和级数的独立求和得到。区间内部的逐项积分不足以单独证明端点等式，需要补上这一端点依据。
\end{solve}

\begin{solve}[习题~\ref{ex:c26q15}：对数和函数与端点]\phantomsection\label{ans:c26q15}
令 $t=x/2$。在 $|t|<1$ 上对几何级数逐项积分，可得
\[
\sum_{n=1}^\infty\frac{t^n}n=-\log(1-t),\qquad S(x)=-\log(1-x/2).
\]
在右端点 $x=2$ 处得到调和级数，发散；左端点 $x=-2$ 处是交错调和级数，收敛。因此收敛区间为 $[-2,2)$。左端点的和为 $-\log2$，可通过 Abel 定理或对交错调和级数的独立求和得到。区间内部的逐项积分不足以单独证明端点等式，需要补上这一端点依据。
\end{solve}

\begin{solve}[习题~\ref{ex:c26q16}：对数和函数与端点]\phantomsection\label{ans:c26q16}
令 $t=x/3$。在 $|t|<1$ 上对几何级数逐项积分，可得
\[
\sum_{n=1}^\infty\frac{t^n}n=-\log(1-t),\qquad S(x)=-\log(1-x/3).
\]
在右端点 $x=3$ 处得到调和级数，发散；左端点 $x=-3$ 处是交错调和级数，收敛。因此收敛区间为 $[-3,3)$。左端点的和为 $-\log2$，可通过 Abel 定理或对交错调和级数的独立求和得到。区间内部的逐项积分不足以单独证明端点等式，需要补上这一端点依据。
\end{solve}

\begin{solve}[习题~\ref{ex:c26q17}：对数和函数与端点]\phantomsection\label{ans:c26q17}
令 $t=x/4$。在 $|t|<1$ 上对几何级数逐项积分，可得
\[
\sum_{n=1}^\infty\frac{t^n}n=-\log(1-t),\qquad S(x)=-\log(1-x/4).
\]
在右端点 $x=4$ 处得到调和级数，发散；左端点 $x=-4$ 处是交错调和级数，收敛。因此收敛区间为 $[-4,4)$。左端点的和为 $-\log2$，可通过 Abel 定理或对交错调和级数的独立求和得到。区间内部的逐项积分不足以单独证明端点等式，需要补上这一端点依据。
\end{solve}

\begin{solve}[习题~\ref{ex:c26q18}：对数和函数与端点]\phantomsection\label{ans:c26q18}
令 $t=x/5$。在 $|t|<1$ 上对几何级数逐项积分，可得
\[
\sum_{n=1}^\infty\frac{t^n}n=-\log(1-t),\qquad S(x)=-\log(1-x/5).
\]
在右端点 $x=5$ 处得到调和级数，发散；左端点 $x=-5$ 处是交错调和级数，收敛。因此收敛区间为 $[-5,5)$。左端点的和为 $-\log2$，可通过 Abel 定理或对交错调和级数的独立求和得到。区间内部的逐项积分不足以单独证明端点等式，需要补上这一端点依据。
\end{solve}

\begin{solve}[习题~\ref{ex:c26q19}：对数和函数与端点]\phantomsection\label{ans:c26q19}
令 $t=x/6$。在 $|t|<1$ 上对几何级数逐项积分，可得
\[
\sum_{n=1}^\infty\frac{t^n}n=-\log(1-t),\qquad S(x)=-\log(1-x/6).
\]
在右端点 $x=6$ 处得到调和级数，发散；左端点 $x=-6$ 处是交错调和级数，收敛。因此收敛区间为 $[-6,6)$。左端点的和为 $-\log2$，可通过 Abel 定理或对交错调和级数的独立求和得到。区间内部的逐项积分不足以单独证明端点等式，需要补上这一端点依据。
\end{solve}

\begin{solve}[习题~\ref{ex:c26q20}：交错级数的误差控制]\phantomsection\label{ans:c26q20}
因 $b_n=n^{-1}$ 正、单调递减且趋零，Leibniz 判别和余项估计适用。
\[
|S-S_N|\le\frac1{(N+1)^{1}}.
\]
因此只需 $(N+1)^{1}>1000$，即 $N+1>1000^{1/1}$。余项与第一个被舍去项同号，所以奇数次部分和在总和上方、偶数次部分和在总和下方。本题条件收敛；尽管绝对值级数发散，交错结构仍然给出有效的误差控制。
\end{solve}

\begin{solve}[习题~\ref{ex:c26q21}：交错级数的误差控制]\phantomsection\label{ans:c26q21}
因 $b_n=n^{-2}$ 正、单调递减且趋零，Leibniz 判别和余项估计适用。
\[
|S-S_N|\le\frac1{(N+1)^{2}}.
\]
因此只需 $(N+1)^{2}>1000$，即 $N+1>1000^{1/2}$。余项与第一个被舍去项同号，所以奇数次部分和在总和上方、偶数次部分和在总和下方。本题还绝对收敛，但交错余项估计通常比直接估计绝对值尾和更方便。
\end{solve}

\begin{solve}[习题~\ref{ex:c26q22}：交错级数的误差控制]\phantomsection\label{ans:c26q22}
因 $b_n=n^{-3}$ 正、单调递减且趋零，Leibniz 判别和余项估计适用。
\[
|S-S_N|\le\frac1{(N+1)^{3}}.
\]
因此只需 $(N+1)^{3}>1000$，即 $N+1>1000^{1/3}$。余项与第一个被舍去项同号，所以奇数次部分和在总和上方、偶数次部分和在总和下方。本题还绝对收敛，但交错余项估计通常比直接估计绝对值尾和更方便。
\end{solve}

\begin{solve}[习题~\ref{ex:c26q23}：交错级数的误差控制]\phantomsection\label{ans:c26q23}
因 $b_n=n^{-4}$ 正、单调递减且趋零，Leibniz 判别和余项估计适用。
\[
|S-S_N|\le\frac1{(N+1)^{4}}.
\]
因此只需 $(N+1)^{4}>1000$，即 $N+1>1000^{1/4}$。余项与第一个被舍去项同号，所以奇数次部分和在总和上方、偶数次部分和在总和下方。本题还绝对收敛，但交错余项估计通常比直接估计绝对值尾和更方便。
\end{solve}

\begin{solve}[习题~\ref{ex:c26q24}：交错级数的误差控制]\phantomsection\label{ans:c26q24}
因 $b_n=n^{-5}$ 正、单调递减且趋零，Leibniz 判别和余项估计适用。
\[
|S-S_N|\le\frac1{(N+1)^{5}}.
\]
因此只需 $(N+1)^{5}>1000$，即 $N+1>1000^{1/5}$。余项与第一个被舍去项同号，所以奇数次部分和在总和上方、偶数次部分和在总和下方。本题还绝对收敛，但交错余项估计通常比直接估计绝对值尾和更方便。
\end{solve}

\begin{solve}[习题~\ref{ex:c26q25}：交错级数的误差控制]\phantomsection\label{ans:c26q25}
因 $b_n=n^{-6}$ 正、单调递减且趋零，Leibniz 判别和余项估计适用。
\[
|S-S_N|\le\frac1{(N+1)^{6}}.
\]
因此只需 $(N+1)^{6}>1000$，即 $N+1>1000^{1/6}$。余项与第一个被舍去项同号，所以奇数次部分和在总和上方、偶数次部分和在总和下方。本题还绝对收敛，但交错余项估计通常比直接估计绝对值尾和更方便。
\end{solve}

\clearpage
\section{幂级数求和与生成函数技巧}\label{sec:ch27}
本章在收敛圆内部使用逐项求导与积分。引入 $D=t\,d/dt$ 可以把几何级数的系数从1依次变为 $n,n^2,n^3$；这一操作对有限次求导在半径内部合法。

\textbf{资源对应：}CLP-2 \S~3.5、\S~3.6。本章共 24 道编号习题。
\subsection{习题}
\begin{example}[平方系数的几何和]\label{ex:c27q1}
求 $\displaystyle\sum_{n=1}^\infty n^2(\frac{1}{2})^n$。
\par\textit{改编训练。}\par\hyperref[ans:c27q1]{解答见第~\pageref{ans:c27q1} 页。}
\end{example}
\begin{example}[平方系数的几何和]\label{ex:c27q2}
求 $\displaystyle\sum_{n=1}^\infty n^2(\frac{1}{3})^n$。
\par\textit{改编训练。}\par\hyperref[ans:c27q2]{解答见第~\pageref{ans:c27q2} 页。}
\end{example}
\begin{example}[平方系数的几何和]\label{ex:c27q3}
求 $\displaystyle\sum_{n=1}^\infty n^2(\frac{2}{3})^n$。
\par\textit{改编训练。}\par\hyperref[ans:c27q3]{解答见第~\pageref{ans:c27q3} 页。}
\end{example}
\begin{example}[平方系数的几何和]\label{ex:c27q4}
求 $\displaystyle\sum_{n=1}^\infty n^2(\frac{1}{4})^n$。
\par\textit{改编训练。}\par\hyperref[ans:c27q4]{解答见第~\pageref{ans:c27q4} 页。}
\end{example}
\begin{example}[平方系数的几何和]\label{ex:c27q5}
求 $\displaystyle\sum_{n=1}^\infty n^2(\frac{2}{5})^n$。
\par\textit{改编训练。}\par\hyperref[ans:c27q5]{解答见第~\pageref{ans:c27q5} 页。}
\end{example}
\begin{example}[平方系数的几何和]\label{ex:c27q6}
求 $\displaystyle\sum_{n=1}^\infty n^2(\frac{3}{4})^n$。
\par\textit{改编训练。}\par\hyperref[ans:c27q6]{解答见第~\pageref{ans:c27q6} 页。}
\end{example}
\begin{example}[立方系数的几何和]\label{ex:c27q7}
求 $\displaystyle\sum_{n=1}^\infty n^3(\frac{1}{2})^n$。
\par\textit{改编训练。}\par\hyperref[ans:c27q7]{解答见第~\pageref{ans:c27q7} 页。}
\end{example}
\begin{example}[立方系数的几何和]\label{ex:c27q8}
求 $\displaystyle\sum_{n=1}^\infty n^3(\frac{1}{3})^n$。
\par\textit{改编训练。}\par\hyperref[ans:c27q8]{解答见第~\pageref{ans:c27q8} 页。}
\end{example}
\begin{example}[立方系数的几何和]\label{ex:c27q9}
求 $\displaystyle\sum_{n=1}^\infty n^3(\frac{2}{3})^n$。
\par\textit{改编训练。}\par\hyperref[ans:c27q9]{解答见第~\pageref{ans:c27q9} 页。}
\end{example}
\begin{example}[立方系数的几何和]\label{ex:c27q10}
求 $\displaystyle\sum_{n=1}^\infty n^3(\frac{1}{4})^n$。
\par\textit{改编训练。}\par\hyperref[ans:c27q10]{解答见第~\pageref{ans:c27q10} 页。}
\end{example}
\begin{example}[立方系数的几何和]\label{ex:c27q11}
求 $\displaystyle\sum_{n=1}^\infty n^3(\frac{2}{5})^n$。
\par\textit{改编训练。}\par\hyperref[ans:c27q11]{解答见第~\pageref{ans:c27q11} 页。}
\end{example}
\begin{example}[立方系数的几何和]\label{ex:c27q12}
求 $\displaystyle\sum_{n=1}^\infty n^3(\frac{3}{4})^n$。
\par\textit{改编训练。}\par\hyperref[ans:c27q12]{解答见第~\pageref{ans:c27q12} 页。}
\end{example}
\begin{example}[两次积分产生的系数]\label{ex:c27q13}
求 $S(x)=\displaystyle\sum_{n=0}^\infty\frac{x^n}{(n+1)(n+2)1^n}$ 在 $|x|<1$ 时的和函数，并说明 $x=0$ 处如何理解。
\par\textit{改编训练。}\par\hyperref[ans:c27q13]{解答见第~\pageref{ans:c27q13} 页。}
\end{example}
\begin{example}[两次积分产生的系数]\label{ex:c27q14}
求 $S(x)=\displaystyle\sum_{n=0}^\infty\frac{x^n}{(n+1)(n+2)2^n}$ 在 $|x|<2$ 时的和函数，并说明 $x=0$ 处如何理解。
\par\textit{改编训练。}\par\hyperref[ans:c27q14]{解答见第~\pageref{ans:c27q14} 页。}
\end{example}
\begin{example}[两次积分产生的系数]\label{ex:c27q15}
求 $S(x)=\displaystyle\sum_{n=0}^\infty\frac{x^n}{(n+1)(n+2)3^n}$ 在 $|x|<3$ 时的和函数，并说明 $x=0$ 处如何理解。
\par\textit{改编训练。}\par\hyperref[ans:c27q15]{解答见第~\pageref{ans:c27q15} 页。}
\end{example}
\begin{example}[两次积分产生的系数]\label{ex:c27q16}
求 $S(x)=\displaystyle\sum_{n=0}^\infty\frac{x^n}{(n+1)(n+2)4^n}$ 在 $|x|<4$ 时的和函数，并说明 $x=0$ 处如何理解。
\par\textit{改编训练。}\par\hyperref[ans:c27q16]{解答见第~\pageref{ans:c27q16} 页。}
\end{example}
\begin{example}[两次积分产生的系数]\label{ex:c27q17}
求 $S(x)=\displaystyle\sum_{n=0}^\infty\frac{x^n}{(n+1)(n+2)5^n}$ 在 $|x|<5$ 时的和函数，并说明 $x=0$ 处如何理解。
\par\textit{改编训练。}\par\hyperref[ans:c27q17]{解答见第~\pageref{ans:c27q17} 页。}
\end{example}
\begin{example}[两次积分产生的系数]\label{ex:c27q18}
求 $S(x)=\displaystyle\sum_{n=0}^\infty\frac{x^n}{(n+1)(n+2)6^n}$ 在 $|x|<6$ 时的和函数，并说明 $x=0$ 处如何理解。
\par\textit{改编训练。}\par\hyperref[ans:c27q18]{解答见第~\pageref{ans:c27q18} 页。}
\end{example}
\begin{example}[由几何展开求有理函数的级数]\label{ex:c27q19}
把 $f(x)=1/(1+x)^2$ 展开成以0为中心的幂级数，指出收敛半径和端点行为。
\par\textit{改编训练。}\par\hyperref[ans:c27q19]{解答见第~\pageref{ans:c27q19} 页。}
\end{example}
\begin{example}[由几何展开求有理函数的级数]\label{ex:c27q20}
把 $f(x)=1/(2+x)^2$ 展开成以0为中心的幂级数，指出收敛半径和端点行为。
\par\textit{改编训练。}\par\hyperref[ans:c27q20]{解答见第~\pageref{ans:c27q20} 页。}
\end{example}
\begin{example}[由几何展开求有理函数的级数]\label{ex:c27q21}
把 $f(x)=1/(3+x)^2$ 展开成以0为中心的幂级数，指出收敛半径和端点行为。
\par\textit{改编训练。}\par\hyperref[ans:c27q21]{解答见第~\pageref{ans:c27q21} 页。}
\end{example}
\begin{example}[由几何展开求有理函数的级数]\label{ex:c27q22}
把 $f(x)=1/(4+x)^2$ 展开成以0为中心的幂级数，指出收敛半径和端点行为。
\par\textit{改编训练。}\par\hyperref[ans:c27q22]{解答见第~\pageref{ans:c27q22} 页。}
\end{example}
\begin{example}[由几何展开求有理函数的级数]\label{ex:c27q23}
把 $f(x)=1/(5+x)^2$ 展开成以0为中心的幂级数，指出收敛半径和端点行为。
\par\textit{改编训练。}\par\hyperref[ans:c27q23]{解答见第~\pageref{ans:c27q23} 页。}
\end{example}
\begin{example}[由几何展开求有理函数的级数]\label{ex:c27q24}
把 $f(x)=1/(6+x)^2$ 展开成以0为中心的幂级数，指出收敛半径和端点行为。
\par\textit{改编训练。}\par\hyperref[ans:c27q24]{解答见第~\pageref{ans:c27q24} 页。}
\end{example}
\subsection{解答}
\begin{solve}[习题~\ref{ex:c27q1}：平方系数的几何和]\phantomsection\label{ans:c27q1}
从 $\sum_{n=0}^\infty t^n=1/(1-t)$ 出发，先作用一次 $t\,d/dt$ 得 $\sum nt^n=t/(1-t)^2$，再作用一次，
\[
\sum_{n=1}^\infty n^2t^n=t\frac d{dt}\frac{t}{(1-t)^2}=\frac{t(1+t)}{(1-t)^3},\quad |t|<1.
\]
代入 $t=\frac{1}{2}$ 得总和 $6$。由于代入点严格位于收敛半径内部，两次逐项求导合法。幂次 $n^2$ 不会破坏几何衰减，因此该数值级数也绝对收敛。
\end{solve}

\begin{solve}[习题~\ref{ex:c27q2}：平方系数的几何和]\phantomsection\label{ans:c27q2}
从 $\sum_{n=0}^\infty t^n=1/(1-t)$ 出发，先作用一次 $t\,d/dt$ 得 $\sum nt^n=t/(1-t)^2$，再作用一次，
\[
\sum_{n=1}^\infty n^2t^n=t\frac d{dt}\frac{t}{(1-t)^2}=\frac{t(1+t)}{(1-t)^3},\quad |t|<1.
\]
代入 $t=\frac{1}{3}$ 得总和 $\frac{3}{2}$。由于代入点严格位于收敛半径内部，两次逐项求导合法。幂次 $n^2$ 不会破坏几何衰减，因此该数值级数也绝对收敛。
\end{solve}

\begin{solve}[习题~\ref{ex:c27q3}：平方系数的几何和]\phantomsection\label{ans:c27q3}
从 $\sum_{n=0}^\infty t^n=1/(1-t)$ 出发，先作用一次 $t\,d/dt$ 得 $\sum nt^n=t/(1-t)^2$，再作用一次，
\[
\sum_{n=1}^\infty n^2t^n=t\frac d{dt}\frac{t}{(1-t)^2}=\frac{t(1+t)}{(1-t)^3},\quad |t|<1.
\]
代入 $t=\frac{2}{3}$ 得总和 $30$。由于代入点严格位于收敛半径内部，两次逐项求导合法。幂次 $n^2$ 不会破坏几何衰减，因此该数值级数也绝对收敛。
\end{solve}

\begin{solve}[习题~\ref{ex:c27q4}：平方系数的几何和]\phantomsection\label{ans:c27q4}
从 $\sum_{n=0}^\infty t^n=1/(1-t)$ 出发，先作用一次 $t\,d/dt$ 得 $\sum nt^n=t/(1-t)^2$，再作用一次，
\[
\sum_{n=1}^\infty n^2t^n=t\frac d{dt}\frac{t}{(1-t)^2}=\frac{t(1+t)}{(1-t)^3},\quad |t|<1.
\]
代入 $t=\frac{1}{4}$ 得总和 $\frac{20}{27}$。由于代入点严格位于收敛半径内部，两次逐项求导合法。幂次 $n^2$ 不会破坏几何衰减，因此该数值级数也绝对收敛。
\end{solve}

\begin{solve}[习题~\ref{ex:c27q5}：平方系数的几何和]\phantomsection\label{ans:c27q5}
从 $\sum_{n=0}^\infty t^n=1/(1-t)$ 出发，先作用一次 $t\,d/dt$ 得 $\sum nt^n=t/(1-t)^2$，再作用一次，
\[
\sum_{n=1}^\infty n^2t^n=t\frac d{dt}\frac{t}{(1-t)^2}=\frac{t(1+t)}{(1-t)^3},\quad |t|<1.
\]
代入 $t=\frac{2}{5}$ 得总和 $\frac{70}{27}$。由于代入点严格位于收敛半径内部，两次逐项求导合法。幂次 $n^2$ 不会破坏几何衰减，因此该数值级数也绝对收敛。
\end{solve}

\begin{solve}[习题~\ref{ex:c27q6}：平方系数的几何和]\phantomsection\label{ans:c27q6}
从 $\sum_{n=0}^\infty t^n=1/(1-t)$ 出发，先作用一次 $t\,d/dt$ 得 $\sum nt^n=t/(1-t)^2$，再作用一次，
\[
\sum_{n=1}^\infty n^2t^n=t\frac d{dt}\frac{t}{(1-t)^2}=\frac{t(1+t)}{(1-t)^3},\quad |t|<1.
\]
代入 $t=\frac{3}{4}$ 得总和 $84$。由于代入点严格位于收敛半径内部，两次逐项求导合法。幂次 $n^2$ 不会破坏几何衰减，因此该数值级数也绝对收敛。
\end{solve}

\begin{solve}[习题~\ref{ex:c27q7}：立方系数的几何和]\phantomsection\label{ans:c27q7}
在 $|t|<1$ 内，对平方系数的和函数再作用 $t\,d/dt$，得到
\[
\sum_{n=1}^\infty n^3t^n=t\frac d{dt}\left[\frac{t(1+t)}{(1-t)^3}\right]=\frac{t(1+4t+t^2)}{(1-t)^4}.
\]
代入 $t=\frac{1}{2}$ 后，总和为 $26$。计算有理函数导数时可先把分母写成 $(1-t)^{-3}$，用乘积法则，再统一分母；这样比较容易避免商法则展开中的符号错误。
\end{solve}

\begin{solve}[习题~\ref{ex:c27q8}：立方系数的几何和]\phantomsection\label{ans:c27q8}
在 $|t|<1$ 内，对平方系数的和函数再作用 $t\,d/dt$，得到
\[
\sum_{n=1}^\infty n^3t^n=t\frac d{dt}\left[\frac{t(1+t)}{(1-t)^3}\right]=\frac{t(1+4t+t^2)}{(1-t)^4}.
\]
代入 $t=\frac{1}{3}$ 后，总和为 $\frac{33}{8}$。计算有理函数导数时可先把分母写成 $(1-t)^{-3}$，用乘积法则，再统一分母；这样比较容易避免商法则展开中的符号错误。
\end{solve}

\begin{solve}[习题~\ref{ex:c27q9}：立方系数的几何和]\phantomsection\label{ans:c27q9}
在 $|t|<1$ 内，对平方系数的和函数再作用 $t\,d/dt$，得到
\[
\sum_{n=1}^\infty n^3t^n=t\frac d{dt}\left[\frac{t(1+t)}{(1-t)^3}\right]=\frac{t(1+4t+t^2)}{(1-t)^4}.
\]
代入 $t=\frac{2}{3}$ 后，总和为 $222$。计算有理函数导数时可先把分母写成 $(1-t)^{-3}$，用乘积法则，再统一分母；这样比较容易避免商法则展开中的符号错误。
\end{solve}

\begin{solve}[习题~\ref{ex:c27q10}：立方系数的几何和]\phantomsection\label{ans:c27q10}
在 $|t|<1$ 内，对平方系数的和函数再作用 $t\,d/dt$，得到
\[
\sum_{n=1}^\infty n^3t^n=t\frac d{dt}\left[\frac{t(1+t)}{(1-t)^3}\right]=\frac{t(1+4t+t^2)}{(1-t)^4}.
\]
代入 $t=\frac{1}{4}$ 后，总和为 $\frac{44}{27}$。计算有理函数导数时可先把分母写成 $(1-t)^{-3}$，用乘积法则，再统一分母；这样比较容易避免商法则展开中的符号错误。
\end{solve}

\begin{solve}[习题~\ref{ex:c27q11}：立方系数的几何和]\phantomsection\label{ans:c27q11}
在 $|t|<1$ 内，对平方系数的和函数再作用 $t\,d/dt$，得到
\[
\sum_{n=1}^\infty n^3t^n=t\frac d{dt}\left[\frac{t(1+t)}{(1-t)^3}\right]=\frac{t(1+4t+t^2)}{(1-t)^4}.
\]
代入 $t=\frac{2}{5}$ 后，总和为 $\frac{230}{27}$。计算有理函数导数时可先把分母写成 $(1-t)^{-3}$，用乘积法则，再统一分母；这样比较容易避免商法则展开中的符号错误。
\end{solve}

\begin{solve}[习题~\ref{ex:c27q12}：立方系数的几何和]\phantomsection\label{ans:c27q12}
在 $|t|<1$ 内，对平方系数的和函数再作用 $t\,d/dt$，得到
\[
\sum_{n=1}^\infty n^3t^n=t\frac d{dt}\left[\frac{t(1+t)}{(1-t)^3}\right]=\frac{t(1+4t+t^2)}{(1-t)^4}.
\]
代入 $t=\frac{3}{4}$ 后，总和为 $876$。计算有理函数导数时可先把分母写成 $(1-t)^{-3}$，用乘积法则，再统一分母；这样比较容易避免商法则展开中的符号错误。
\end{solve}

\begin{solve}[习题~\ref{ex:c27q13}：两次积分产生的系数]\phantomsection\label{ans:c27q13}
令 $t=x/1$。先把级数乘以 $t^2$，其二阶导数为几何级数。因此 $H(t)=t^2S$ 满足 $H^{\prime\prime}=1/(1-t)$、$H(0)=H^{\prime}(0)=0$。两次积分给出
\[
H(t)=t+(1-t)\log(1-t).
\]
所以
\[
S=\frac{t+(1-t)\log(1-t)}{t^2}\quad(t\ne0),\qquad S(0)=\frac12.
\]
在 $t=0$，有理表达式表面为 $0/0$，但其分子二阶项为 $t^2/2$，所以与级数的常数项一致，是可去奇点。
\end{solve}

\begin{solve}[习题~\ref{ex:c27q14}：两次积分产生的系数]\phantomsection\label{ans:c27q14}
令 $t=x/2$。先把级数乘以 $t^2$，其二阶导数为几何级数。因此 $H(t)=t^2S$ 满足 $H^{\prime\prime}=1/(1-t)$、$H(0)=H^{\prime}(0)=0$。两次积分给出
\[
H(t)=t+(1-t)\log(1-t).
\]
所以
\[
S=\frac{t+(1-t)\log(1-t)}{t^2}\quad(t\ne0),\qquad S(0)=\frac12.
\]
在 $t=0$，有理表达式表面为 $0/0$，但其分子二阶项为 $t^2/2$，所以与级数的常数项一致，是可去奇点。
\end{solve}

\begin{solve}[习题~\ref{ex:c27q15}：两次积分产生的系数]\phantomsection\label{ans:c27q15}
令 $t=x/3$。先把级数乘以 $t^2$，其二阶导数为几何级数。因此 $H(t)=t^2S$ 满足 $H^{\prime\prime}=1/(1-t)$、$H(0)=H^{\prime}(0)=0$。两次积分给出
\[
H(t)=t+(1-t)\log(1-t).
\]
所以
\[
S=\frac{t+(1-t)\log(1-t)}{t^2}\quad(t\ne0),\qquad S(0)=\frac12.
\]
在 $t=0$，有理表达式表面为 $0/0$，但其分子二阶项为 $t^2/2$，所以与级数的常数项一致，是可去奇点。
\end{solve}

\begin{solve}[习题~\ref{ex:c27q16}：两次积分产生的系数]\phantomsection\label{ans:c27q16}
令 $t=x/4$。先把级数乘以 $t^2$，其二阶导数为几何级数。因此 $H(t)=t^2S$ 满足 $H^{\prime\prime}=1/(1-t)$、$H(0)=H^{\prime}(0)=0$。两次积分给出
\[
H(t)=t+(1-t)\log(1-t).
\]
所以
\[
S=\frac{t+(1-t)\log(1-t)}{t^2}\quad(t\ne0),\qquad S(0)=\frac12.
\]
在 $t=0$，有理表达式表面为 $0/0$，但其分子二阶项为 $t^2/2$，所以与级数的常数项一致，是可去奇点。
\end{solve}

\begin{solve}[习题~\ref{ex:c27q17}：两次积分产生的系数]\phantomsection\label{ans:c27q17}
令 $t=x/5$。先把级数乘以 $t^2$，其二阶导数为几何级数。因此 $H(t)=t^2S$ 满足 $H^{\prime\prime}=1/(1-t)$、$H(0)=H^{\prime}(0)=0$。两次积分给出
\[
H(t)=t+(1-t)\log(1-t).
\]
所以
\[
S=\frac{t+(1-t)\log(1-t)}{t^2}\quad(t\ne0),\qquad S(0)=\frac12.
\]
在 $t=0$，有理表达式表面为 $0/0$，但其分子二阶项为 $t^2/2$，所以与级数的常数项一致，是可去奇点。
\end{solve}

\begin{solve}[习题~\ref{ex:c27q18}：两次积分产生的系数]\phantomsection\label{ans:c27q18}
令 $t=x/6$。先把级数乘以 $t^2$，其二阶导数为几何级数。因此 $H(t)=t^2S$ 满足 $H^{\prime\prime}=1/(1-t)$、$H(0)=H^{\prime}(0)=0$。两次积分给出
\[
H(t)=t+(1-t)\log(1-t).
\]
所以
\[
S=\frac{t+(1-t)\log(1-t)}{t^2}\quad(t\ne0),\qquad S(0)=\frac12.
\]
在 $t=0$，有理表达式表面为 $0/0$，但其分子二阶项为 $t^2/2$，所以与级数的常数项一致，是可去奇点。
\end{solve}

\begin{solve}[习题~\ref{ex:c27q19}：由几何展开求有理函数的级数]\phantomsection\label{ans:c27q19}
先写 $1/(1+x)=1\sum_{n=0}^\infty(-x/1)^n$，在 $|x|<1$ 内逐项求导并改变符号，得到
\[
\frac1{(1+x)^2}=\sum_{n=0}^\infty\frac{(-1)^n(n+1)}{1^{n+2}}x^n.
\]
收敛半径为 $1$。在 $x=1$ 时通项为 $(-1)^n(n+1)/1$，在 $x=-1$ 时通项为 $(n+1)/1$，两者均不趋零。因此两端都不收敛，即使函数本身在正端点是有定义的，也不能据此推出此幂级数在那里收敛。
\end{solve}

\begin{solve}[习题~\ref{ex:c27q20}：由几何展开求有理函数的级数]\phantomsection\label{ans:c27q20}
先写 $1/(2+x)=\frac{1}{2}\sum_{n=0}^\infty(-x/2)^n$，在 $|x|<2$ 内逐项求导并改变符号，得到
\[
\frac1{(2+x)^2}=\sum_{n=0}^\infty\frac{(-1)^n(n+1)}{2^{n+2}}x^n.
\]
收敛半径为 $2$。在 $x=2$ 时通项为 $(-1)^n(n+1)/4$，在 $x=-2$ 时通项为 $(n+1)/4$，两者均不趋零。因此两端都不收敛，即使函数本身在正端点是有定义的，也不能据此推出此幂级数在那里收敛。
\end{solve}

\begin{solve}[习题~\ref{ex:c27q21}：由几何展开求有理函数的级数]\phantomsection\label{ans:c27q21}
先写 $1/(3+x)=\frac{1}{3}\sum_{n=0}^\infty(-x/3)^n$，在 $|x|<3$ 内逐项求导并改变符号，得到
\[
\frac1{(3+x)^2}=\sum_{n=0}^\infty\frac{(-1)^n(n+1)}{3^{n+2}}x^n.
\]
收敛半径为 $3$。在 $x=3$ 时通项为 $(-1)^n(n+1)/9$，在 $x=-3$ 时通项为 $(n+1)/9$，两者均不趋零。因此两端都不收敛，即使函数本身在正端点是有定义的，也不能据此推出此幂级数在那里收敛。
\end{solve}

\begin{solve}[习题~\ref{ex:c27q22}：由几何展开求有理函数的级数]\phantomsection\label{ans:c27q22}
先写 $1/(4+x)=\frac{1}{4}\sum_{n=0}^\infty(-x/4)^n$，在 $|x|<4$ 内逐项求导并改变符号，得到
\[
\frac1{(4+x)^2}=\sum_{n=0}^\infty\frac{(-1)^n(n+1)}{4^{n+2}}x^n.
\]
收敛半径为 $4$。在 $x=4$ 时通项为 $(-1)^n(n+1)/16$，在 $x=-4$ 时通项为 $(n+1)/16$，两者均不趋零。因此两端都不收敛，即使函数本身在正端点是有定义的，也不能据此推出此幂级数在那里收敛。
\end{solve}

\begin{solve}[习题~\ref{ex:c27q23}：由几何展开求有理函数的级数]\phantomsection\label{ans:c27q23}
先写 $1/(5+x)=\frac{1}{5}\sum_{n=0}^\infty(-x/5)^n$，在 $|x|<5$ 内逐项求导并改变符号，得到
\[
\frac1{(5+x)^2}=\sum_{n=0}^\infty\frac{(-1)^n(n+1)}{5^{n+2}}x^n.
\]
收敛半径为 $5$。在 $x=5$ 时通项为 $(-1)^n(n+1)/25$，在 $x=-5$ 时通项为 $(n+1)/25$，两者均不趋零。因此两端都不收敛，即使函数本身在正端点是有定义的，也不能据此推出此幂级数在那里收敛。
\end{solve}

\begin{solve}[习题~\ref{ex:c27q24}：由几何展开求有理函数的级数]\phantomsection\label{ans:c27q24}
先写 $1/(6+x)=\frac{1}{6}\sum_{n=0}^\infty(-x/6)^n$，在 $|x|<6$ 内逐项求导并改变符号，得到
\[
\frac1{(6+x)^2}=\sum_{n=0}^\infty\frac{(-1)^n(n+1)}{6^{n+2}}x^n.
\]
收敛半径为 $6$。在 $x=6$ 时通项为 $(-1)^n(n+1)/36$，在 $x=-6$ 时通项为 $(n+1)/36$，两者均不趋零。因此两端都不收敛，即使函数本身在正端点是有定义的，也不能据此推出此幂级数在那里收敛。
\end{solve}

\clearpage
\section{函数列、一致收敛与傅里叶入门}\label{sec:ch28}
逐点收敛只固定一个点后令下标趋于无穷，一致收敛（uniform convergence）要求误差对所有点同时变小。交换极限与积分、求导等运算需要适当条件。最后几题引入傅里叶级数（Fourier series）的正交系数。

\textbf{资源对应：}补充证明训练；CLP-2 \S~3.5、\S~3.6 提供幂级数背景。本章共 24 道编号习题。
\subsection{习题}
\begin{example}[边界层导致不一致收敛]\label{ex:c28q1}
讨论 $f_n(x)=x^{1n}$ 在 $[0,1]$ 上的逐点极限与一致收敛性，并讨论 $[0,r]$（$0<r<1$）的情形。
\par\textit{补充证明。}\par\hyperref[ans:c28q1]{解答见第~\pageref{ans:c28q1} 页。}
\end{example}
\begin{example}[边界层导致不一致收敛]\label{ex:c28q2}
讨论 $f_n(x)=x^{2n}$ 在 $[0,1]$ 上的逐点极限与一致收敛性，并讨论 $[0,r]$（$0<r<1$）的情形。
\par\textit{补充证明。}\par\hyperref[ans:c28q2]{解答见第~\pageref{ans:c28q2} 页。}
\end{example}
\begin{example}[边界层导致不一致收敛]\label{ex:c28q3}
讨论 $f_n(x)=x^{3n}$ 在 $[0,1]$ 上的逐点极限与一致收敛性，并讨论 $[0,r]$（$0<r<1$）的情形。
\par\textit{补充证明。}\par\hyperref[ans:c28q3]{解答见第~\pageref{ans:c28q3} 页。}
\end{example}
\begin{example}[边界层导致不一致收敛]\label{ex:c28q4}
讨论 $f_n(x)=x^{4n}$ 在 $[0,1]$ 上的逐点极限与一致收敛性，并讨论 $[0,r]$（$0<r<1$）的情形。
\par\textit{补充证明。}\par\hyperref[ans:c28q4]{解答见第~\pageref{ans:c28q4} 页。}
\end{example}
\begin{example}[边界层导致不一致收敛]\label{ex:c28q5}
讨论 $f_n(x)=x^{5n}$ 在 $[0,1]$ 上的逐点极限与一致收敛性，并讨论 $[0,r]$（$0<r<1$）的情形。
\par\textit{补充证明。}\par\hyperref[ans:c28q5]{解答见第~\pageref{ans:c28q5} 页。}
\end{example}
\begin{example}[边界层导致不一致收敛]\label{ex:c28q6}
讨论 $f_n(x)=x^{6n}$ 在 $[0,1]$ 上的逐点极限与一致收敛性，并讨论 $[0,r]$（$0<r<1$）的情形。
\par\textit{补充证明。}\par\hyperref[ans:c28q6]{解答见第~\pageref{ans:c28q6} 页。}
\end{example}
\begin{example}[移动峰值的函数列]\label{ex:c28q7}
在 $[0,\infty)$ 上设 $f_n(x)=n^{0}x e^{-n^{1}x}$。求逐点极限，判断是否一致收敛。
\par\textit{补充证明。}\par\hyperref[ans:c28q7]{解答见第~\pageref{ans:c28q7} 页。}
\end{example}
\begin{example}[移动峰值的函数列]\label{ex:c28q8}
在 $[0,\infty)$ 上设 $f_n(x)=n^{1}x e^{-n^{1}x}$。求逐点极限，判断是否一致收敛。
\par\textit{补充证明。}\par\hyperref[ans:c28q8]{解答见第~\pageref{ans:c28q8} 页。}
\end{example}
\begin{example}[移动峰值的函数列]\label{ex:c28q9}
在 $[0,\infty)$ 上设 $f_n(x)=n^{1}x e^{-n^{2}x}$。求逐点极限，判断是否一致收敛。
\par\textit{补充证明。}\par\hyperref[ans:c28q9]{解答见第~\pageref{ans:c28q9} 页。}
\end{example}
\begin{example}[移动峰值的函数列]\label{ex:c28q10}
在 $[0,\infty)$ 上设 $f_n(x)=n^{2}x e^{-n^{1}x}$。求逐点极限，判断是否一致收敛。
\par\textit{补充证明。}\par\hyperref[ans:c28q10]{解答见第~\pageref{ans:c28q10} 页。}
\end{example}
\begin{example}[移动峰值的函数列]\label{ex:c28q11}
在 $[0,\infty)$ 上设 $f_n(x)=n^{2}x e^{-n^{2}x}$。求逐点极限，判断是否一致收敛。
\par\textit{补充证明。}\par\hyperref[ans:c28q11]{解答见第~\pageref{ans:c28q11} 页。}
\end{example}
\begin{example}[移动峰值的函数列]\label{ex:c28q12}
在 $[0,\infty)$ 上设 $f_n(x)=n^{2}x e^{-n^{3}x}$。求逐点极限，判断是否一致收敛。
\par\textit{补充证明。}\par\hyperref[ans:c28q12]{解答见第~\pageref{ans:c28q12} 页。}
\end{example}
\begin{example}[锯齿函数的傅里叶系数]\label{ex:c28q13}
在 $(-\pi,\pi)$ 上令 $f(x)=1x$，并作 $2\pi$ 周期延拓。求傅里叶系数，说明级数在端点的和。
\par\textit{补充证明。}\par\hyperref[ans:c28q13]{解答见第~\pageref{ans:c28q13} 页。}
\end{example}
\begin{example}[锯齿函数的傅里叶系数]\label{ex:c28q14}
在 $(-\pi,\pi)$ 上令 $f(x)=2x$，并作 $2\pi$ 周期延拓。求傅里叶系数，说明级数在端点的和。
\par\textit{补充证明。}\par\hyperref[ans:c28q14]{解答见第~\pageref{ans:c28q14} 页。}
\end{example}
\begin{example}[锯齿函数的傅里叶系数]\label{ex:c28q15}
在 $(-\pi,\pi)$ 上令 $f(x)=3x$，并作 $2\pi$ 周期延拓。求傅里叶系数，说明级数在端点的和。
\par\textit{补充证明。}\par\hyperref[ans:c28q15]{解答见第~\pageref{ans:c28q15} 页。}
\end{example}
\begin{example}[锯齿函数的傅里叶系数]\label{ex:c28q16}
在 $(-\pi,\pi)$ 上令 $f(x)=4x$，并作 $2\pi$ 周期延拓。求傅里叶系数，说明级数在端点的和。
\par\textit{补充证明。}\par\hyperref[ans:c28q16]{解答见第~\pageref{ans:c28q16} 页。}
\end{example}
\begin{example}[锯齿函数的傅里叶系数]\label{ex:c28q17}
在 $(-\pi,\pi)$ 上令 $f(x)=5x$，并作 $2\pi$ 周期延拓。求傅里叶系数，说明级数在端点的和。
\par\textit{补充证明。}\par\hyperref[ans:c28q17]{解答见第~\pageref{ans:c28q17} 页。}
\end{example}
\begin{example}[锯齿函数的傅里叶系数]\label{ex:c28q18}
在 $(-\pi,\pi)$ 上令 $f(x)=6x$，并作 $2\pi$ 周期延拓。求傅里叶系数，说明级数在端点的和。
\par\textit{补充证明。}\par\hyperref[ans:c28q18]{解答见第~\pageref{ans:c28q18} 页。}
\end{example}
\begin{example}[平方函数的余弦展开]\label{ex:c28q19}
求 $f(x)=1x^2$ 在 $[-\pi,\pi]$ 上的傅里叶系数，并据此求 $\sum_{n=1}^\infty1/n^2$。
\par\textit{补充证明。}\par\hyperref[ans:c28q19]{解答见第~\pageref{ans:c28q19} 页。}
\end{example}
\begin{example}[平方函数的余弦展开]\label{ex:c28q20}
求 $f(x)=2x^2$ 在 $[-\pi,\pi]$ 上的傅里叶系数，并据此求 $\sum_{n=1}^\infty1/n^2$。
\par\textit{补充证明。}\par\hyperref[ans:c28q20]{解答见第~\pageref{ans:c28q20} 页。}
\end{example}
\begin{example}[平方函数的余弦展开]\label{ex:c28q21}
求 $f(x)=3x^2$ 在 $[-\pi,\pi]$ 上的傅里叶系数，并据此求 $\sum_{n=1}^\infty1/n^2$。
\par\textit{补充证明。}\par\hyperref[ans:c28q21]{解答见第~\pageref{ans:c28q21} 页。}
\end{example}
\begin{example}[平方函数的余弦展开]\label{ex:c28q22}
求 $f(x)=4x^2$ 在 $[-\pi,\pi]$ 上的傅里叶系数，并据此求 $\sum_{n=1}^\infty1/n^2$。
\par\textit{补充证明。}\par\hyperref[ans:c28q22]{解答见第~\pageref{ans:c28q22} 页。}
\end{example}
\begin{example}[平方函数的余弦展开]\label{ex:c28q23}
求 $f(x)=5x^2$ 在 $[-\pi,\pi]$ 上的傅里叶系数，并据此求 $\sum_{n=1}^\infty1/n^2$。
\par\textit{补充证明。}\par\hyperref[ans:c28q23]{解答见第~\pageref{ans:c28q23} 页。}
\end{example}
\begin{example}[平方函数的余弦展开]\label{ex:c28q24}
求 $f(x)=6x^2$ 在 $[-\pi,\pi]$ 上的傅里叶系数，并据此求 $\sum_{n=1}^\infty1/n^2$。
\par\textit{补充证明。}\par\hyperref[ans:c28q24]{解答见第~\pageref{ans:c28q24} 页。}
\end{example}
\subsection{解答}
\begin{solve}[习题~\ref{ex:c28q1}：边界层导致不一致收敛]\phantomsection\label{ans:c28q1}
对 $0\le x<1$，幂趋于零；在 $x=1$，函数值恒为1。因此极限函数在 $[0,1)$ 为0，在1处为1。它不连续，而连续函数列的一致极限必连续，所以原区间上不一致收敛。也可直接估计：
\[
\sup_{0\le x<1}|f_n(x)-0|=1.
\]
在任意固定的 $[0,r]$ 上，有 $\sup|f_n|=r^{1n}\to0$，因此一致收敛。去掉一个端点所得 $[0,1)$ 上仍不一致收敛，因为上确界依旧为1；必须把右端点固定在小于1的数处。
\end{solve}

\begin{solve}[习题~\ref{ex:c28q2}：边界层导致不一致收敛]\phantomsection\label{ans:c28q2}
对 $0\le x<1$，幂趋于零；在 $x=1$，函数值恒为1。因此极限函数在 $[0,1)$ 为0，在1处为1。它不连续，而连续函数列的一致极限必连续，所以原区间上不一致收敛。也可直接估计：
\[
\sup_{0\le x<1}|f_n(x)-0|=1.
\]
在任意固定的 $[0,r]$ 上，有 $\sup|f_n|=r^{2n}\to0$，因此一致收敛。去掉一个端点所得 $[0,1)$ 上仍不一致收敛，因为上确界依旧为1；必须把右端点固定在小于1的数处。
\end{solve}

\begin{solve}[习题~\ref{ex:c28q3}：边界层导致不一致收敛]\phantomsection\label{ans:c28q3}
对 $0\le x<1$，幂趋于零；在 $x=1$，函数值恒为1。因此极限函数在 $[0,1)$ 为0，在1处为1。它不连续，而连续函数列的一致极限必连续，所以原区间上不一致收敛。也可直接估计：
\[
\sup_{0\le x<1}|f_n(x)-0|=1.
\]
在任意固定的 $[0,r]$ 上，有 $\sup|f_n|=r^{3n}\to0$，因此一致收敛。去掉一个端点所得 $[0,1)$ 上仍不一致收敛，因为上确界依旧为1；必须把右端点固定在小于1的数处。
\end{solve}

\begin{solve}[习题~\ref{ex:c28q4}：边界层导致不一致收敛]\phantomsection\label{ans:c28q4}
对 $0\le x<1$，幂趋于零；在 $x=1$，函数值恒为1。因此极限函数在 $[0,1)$ 为0，在1处为1。它不连续，而连续函数列的一致极限必连续，所以原区间上不一致收敛。也可直接估计：
\[
\sup_{0\le x<1}|f_n(x)-0|=1.
\]
在任意固定的 $[0,r]$ 上，有 $\sup|f_n|=r^{4n}\to0$，因此一致收敛。去掉一个端点所得 $[0,1)$ 上仍不一致收敛，因为上确界依旧为1；必须把右端点固定在小于1的数处。
\end{solve}

\begin{solve}[习题~\ref{ex:c28q5}：边界层导致不一致收敛]\phantomsection\label{ans:c28q5}
对 $0\le x<1$，幂趋于零；在 $x=1$，函数值恒为1。因此极限函数在 $[0,1)$ 为0，在1处为1。它不连续，而连续函数列的一致极限必连续，所以原区间上不一致收敛。也可直接估计：
\[
\sup_{0\le x<1}|f_n(x)-0|=1.
\]
在任意固定的 $[0,r]$ 上，有 $\sup|f_n|=r^{5n}\to0$，因此一致收敛。去掉一个端点所得 $[0,1)$ 上仍不一致收敛，因为上确界依旧为1；必须把右端点固定在小于1的数处。
\end{solve}

\begin{solve}[习题~\ref{ex:c28q6}：边界层导致不一致收敛]\phantomsection\label{ans:c28q6}
对 $0\le x<1$，幂趋于零；在 $x=1$，函数值恒为1。因此极限函数在 $[0,1)$ 为0，在1处为1。它不连续，而连续函数列的一致极限必连续，所以原区间上不一致收敛。也可直接估计：
\[
\sup_{0\le x<1}|f_n(x)-0|=1.
\]
在任意固定的 $[0,r]$ 上，有 $\sup|f_n|=r^{6n}\to0$，因此一致收敛。去掉一个端点所得 $[0,1)$ 上仍不一致收敛，因为上确界依旧为1；必须把右端点固定在小于1的数处。
\end{solve}

\begin{solve}[习题~\ref{ex:c28q7}：移动峰值的函数列]\phantomsection\label{ans:c28q7}
当 $x=0$ 时函数值为0；当固定 $x>0$ 时，指数衰减压过任意固定次的 $n$ 幂，所以逐点极限为0。求关于 $x$ 的导数：
\[
f_n^{\prime}(x)=n^{0}e^{-n^{1}x}(1-n^{1}x).
\]
峰值在 $x=n^{-1}$ 处，因而
\[
\sup_{x\ge0}|f_n(x)|=e^{-1}n^{-1}.
\]
这个上确界趋于0，所以一致收敛。逐点极限看不到向原点移动的峰值，而上确界范数恰好能检测到这种现象。
\end{solve}

\begin{solve}[习题~\ref{ex:c28q8}：移动峰值的函数列]\phantomsection\label{ans:c28q8}
当 $x=0$ 时函数值为0；当固定 $x>0$ 时，指数衰减压过任意固定次的 $n$ 幂，所以逐点极限为0。求关于 $x$ 的导数：
\[
f_n^{\prime}(x)=n^{1}e^{-n^{1}x}(1-n^{1}x).
\]
峰值在 $x=n^{-1}$ 处，因而
\[
\sup_{x\ge0}|f_n(x)|=e^{-1}n^{0}.
\]
这个上确界不趋于0，所以不一致收敛。逐点极限看不到向原点移动的峰值，而上确界范数恰好能检测到这种现象。
\end{solve}

\begin{solve}[习题~\ref{ex:c28q9}：移动峰值的函数列]\phantomsection\label{ans:c28q9}
当 $x=0$ 时函数值为0；当固定 $x>0$ 时，指数衰减压过任意固定次的 $n$ 幂，所以逐点极限为0。求关于 $x$ 的导数：
\[
f_n^{\prime}(x)=n^{1}e^{-n^{2}x}(1-n^{2}x).
\]
峰值在 $x=n^{-2}$ 处，因而
\[
\sup_{x\ge0}|f_n(x)|=e^{-1}n^{-1}.
\]
这个上确界趋于0，所以一致收敛。逐点极限看不到向原点移动的峰值，而上确界范数恰好能检测到这种现象。
\end{solve}

\begin{solve}[习题~\ref{ex:c28q10}：移动峰值的函数列]\phantomsection\label{ans:c28q10}
当 $x=0$ 时函数值为0；当固定 $x>0$ 时，指数衰减压过任意固定次的 $n$ 幂，所以逐点极限为0。求关于 $x$ 的导数：
\[
f_n^{\prime}(x)=n^{2}e^{-n^{1}x}(1-n^{1}x).
\]
峰值在 $x=n^{-1}$ 处，因而
\[
\sup_{x\ge0}|f_n(x)|=e^{-1}n^{1}.
\]
这个上确界不趋于0，所以不一致收敛。逐点极限看不到向原点移动的峰值，而上确界范数恰好能检测到这种现象。
\end{solve}

\begin{solve}[习题~\ref{ex:c28q11}：移动峰值的函数列]\phantomsection\label{ans:c28q11}
当 $x=0$ 时函数值为0；当固定 $x>0$ 时，指数衰减压过任意固定次的 $n$ 幂，所以逐点极限为0。求关于 $x$ 的导数：
\[
f_n^{\prime}(x)=n^{2}e^{-n^{2}x}(1-n^{2}x).
\]
峰值在 $x=n^{-2}$ 处，因而
\[
\sup_{x\ge0}|f_n(x)|=e^{-1}n^{0}.
\]
这个上确界不趋于0，所以不一致收敛。逐点极限看不到向原点移动的峰值，而上确界范数恰好能检测到这种现象。
\end{solve}

\begin{solve}[习题~\ref{ex:c28q12}：移动峰值的函数列]\phantomsection\label{ans:c28q12}
当 $x=0$ 时函数值为0；当固定 $x>0$ 时，指数衰减压过任意固定次的 $n$ 幂，所以逐点极限为0。求关于 $x$ 的导数：
\[
f_n^{\prime}(x)=n^{2}e^{-n^{3}x}(1-n^{3}x).
\]
峰值在 $x=n^{-3}$ 处，因而
\[
\sup_{x\ge0}|f_n(x)|=e^{-1}n^{-1}.
\]
这个上确界趋于0，所以一致收敛。逐点极限看不到向原点移动的峰值，而上确界范数恰好能检测到这种现象。
\end{solve}

\begin{solve}[习题~\ref{ex:c28q13}：锯齿函数的傅里叶系数]\phantomsection\label{ans:c28q13}
采用 $f\sim a_0/2+\sum_{n\ge1}(a_n\cos nx+b_n\sin nx)$ 的约定。函数为奇函数，故 $a_0=a_n=0$。分部积分得
\[
b_n=\frac{2}\pi\int_0^\pi x\sin(nx)dx=\frac{2(-1)^{n+1}}n.
\]
因此在 $-\pi<x<\pi$ 上，
\[
1x=2\sum_{n=1}^\infty\frac{(-1)^{n+1}}n\sin(nx).
\]
周期延拓分段光滑，Dirichlet 收敛定理说明级数在每点收敛到左右极限的平均值。在 $x=\pm\pi$，两个单侧极限互为相反数，故级数和为0；不能把原区间端点的某个任意定义值直接当成级数和。
\end{solve}

\begin{solve}[习题~\ref{ex:c28q14}：锯齿函数的傅里叶系数]\phantomsection\label{ans:c28q14}
采用 $f\sim a_0/2+\sum_{n\ge1}(a_n\cos nx+b_n\sin nx)$ 的约定。函数为奇函数，故 $a_0=a_n=0$。分部积分得
\[
b_n=\frac{4}\pi\int_0^\pi x\sin(nx)dx=\frac{4(-1)^{n+1}}n.
\]
因此在 $-\pi<x<\pi$ 上，
\[
2x=4\sum_{n=1}^\infty\frac{(-1)^{n+1}}n\sin(nx).
\]
周期延拓分段光滑，Dirichlet 收敛定理说明级数在每点收敛到左右极限的平均值。在 $x=\pm\pi$，两个单侧极限互为相反数，故级数和为0；不能把原区间端点的某个任意定义值直接当成级数和。
\end{solve}

\begin{solve}[习题~\ref{ex:c28q15}：锯齿函数的傅里叶系数]\phantomsection\label{ans:c28q15}
采用 $f\sim a_0/2+\sum_{n\ge1}(a_n\cos nx+b_n\sin nx)$ 的约定。函数为奇函数，故 $a_0=a_n=0$。分部积分得
\[
b_n=\frac{6}\pi\int_0^\pi x\sin(nx)dx=\frac{6(-1)^{n+1}}n.
\]
因此在 $-\pi<x<\pi$ 上，
\[
3x=6\sum_{n=1}^\infty\frac{(-1)^{n+1}}n\sin(nx).
\]
周期延拓分段光滑，Dirichlet 收敛定理说明级数在每点收敛到左右极限的平均值。在 $x=\pm\pi$，两个单侧极限互为相反数，故级数和为0；不能把原区间端点的某个任意定义值直接当成级数和。
\end{solve}

\begin{solve}[习题~\ref{ex:c28q16}：锯齿函数的傅里叶系数]\phantomsection\label{ans:c28q16}
采用 $f\sim a_0/2+\sum_{n\ge1}(a_n\cos nx+b_n\sin nx)$ 的约定。函数为奇函数，故 $a_0=a_n=0$。分部积分得
\[
b_n=\frac{8}\pi\int_0^\pi x\sin(nx)dx=\frac{8(-1)^{n+1}}n.
\]
因此在 $-\pi<x<\pi$ 上，
\[
4x=8\sum_{n=1}^\infty\frac{(-1)^{n+1}}n\sin(nx).
\]
周期延拓分段光滑，Dirichlet 收敛定理说明级数在每点收敛到左右极限的平均值。在 $x=\pm\pi$，两个单侧极限互为相反数，故级数和为0；不能把原区间端点的某个任意定义值直接当成级数和。
\end{solve}

\begin{solve}[习题~\ref{ex:c28q17}：锯齿函数的傅里叶系数]\phantomsection\label{ans:c28q17}
采用 $f\sim a_0/2+\sum_{n\ge1}(a_n\cos nx+b_n\sin nx)$ 的约定。函数为奇函数，故 $a_0=a_n=0$。分部积分得
\[
b_n=\frac{10}\pi\int_0^\pi x\sin(nx)dx=\frac{10(-1)^{n+1}}n.
\]
因此在 $-\pi<x<\pi$ 上，
\[
5x=10\sum_{n=1}^\infty\frac{(-1)^{n+1}}n\sin(nx).
\]
周期延拓分段光滑，Dirichlet 收敛定理说明级数在每点收敛到左右极限的平均值。在 $x=\pm\pi$，两个单侧极限互为相反数，故级数和为0；不能把原区间端点的某个任意定义值直接当成级数和。
\end{solve}

\begin{solve}[习题~\ref{ex:c28q18}：锯齿函数的傅里叶系数]\phantomsection\label{ans:c28q18}
采用 $f\sim a_0/2+\sum_{n\ge1}(a_n\cos nx+b_n\sin nx)$ 的约定。函数为奇函数，故 $a_0=a_n=0$。分部积分得
\[
b_n=\frac{12}\pi\int_0^\pi x\sin(nx)dx=\frac{12(-1)^{n+1}}n.
\]
因此在 $-\pi<x<\pi$ 上，
\[
6x=12\sum_{n=1}^\infty\frac{(-1)^{n+1}}n\sin(nx).
\]
周期延拓分段光滑，Dirichlet 收敛定理说明级数在每点收敛到左右极限的平均值。在 $x=\pm\pi$，两个单侧极限互为相反数，故级数和为0；不能把原区间端点的某个任意定义值直接当成级数和。
\end{solve}

\begin{solve}[习题~\ref{ex:c28q19}：平方函数的余弦展开]\phantomsection\label{ans:c28q19}
函数为偶函数，所以所有正弦系数为0。常数项及余弦系数分别为
\[
\frac{a_0}2=\frac{1}\pi\int_0^\pi x^2dx=\frac{1\pi^2}3,\qquad a_n=\frac{2}\pi\int_0^\pi x^2\cos(nx)dx=\frac{4(-1)^n}{n^2}.
\]
后一式通过两次分部积分得到。周期延拓连续且分段光滑，傅里叶级数在端点也收敛到函数值。令 $x=\pi$，
\[
1\pi^2=\frac{1\pi^2}3+4\sum_{n=1}^\infty\frac1{n^2},\qquad\sum_{n=1}^\infty\frac1{n^2}=\frac{\pi^2}6.
\]
代入端点的合法性来自收敛定理；形式上写出系数后立即代值，仍缺少这一步依据。
\end{solve}

\begin{solve}[习题~\ref{ex:c28q20}：平方函数的余弦展开]\phantomsection\label{ans:c28q20}
函数为偶函数，所以所有正弦系数为0。常数项及余弦系数分别为
\[
\frac{a_0}2=\frac{2}\pi\int_0^\pi x^2dx=\frac{2\pi^2}3,\qquad a_n=\frac{4}\pi\int_0^\pi x^2\cos(nx)dx=\frac{8(-1)^n}{n^2}.
\]
后一式通过两次分部积分得到。周期延拓连续且分段光滑，傅里叶级数在端点也收敛到函数值。令 $x=\pi$，
\[
2\pi^2=\frac{2\pi^2}3+8\sum_{n=1}^\infty\frac1{n^2},\qquad\sum_{n=1}^\infty\frac1{n^2}=\frac{\pi^2}6.
\]
代入端点的合法性来自收敛定理；形式上写出系数后立即代值，仍缺少这一步依据。
\end{solve}

\begin{solve}[习题~\ref{ex:c28q21}：平方函数的余弦展开]\phantomsection\label{ans:c28q21}
函数为偶函数，所以所有正弦系数为0。常数项及余弦系数分别为
\[
\frac{a_0}2=\frac{3}\pi\int_0^\pi x^2dx=\frac{3\pi^2}3,\qquad a_n=\frac{6}\pi\int_0^\pi x^2\cos(nx)dx=\frac{12(-1)^n}{n^2}.
\]
后一式通过两次分部积分得到。周期延拓连续且分段光滑，傅里叶级数在端点也收敛到函数值。令 $x=\pi$，
\[
3\pi^2=\frac{3\pi^2}3+12\sum_{n=1}^\infty\frac1{n^2},\qquad\sum_{n=1}^\infty\frac1{n^2}=\frac{\pi^2}6.
\]
代入端点的合法性来自收敛定理；形式上写出系数后立即代值，仍缺少这一步依据。
\end{solve}

\begin{solve}[习题~\ref{ex:c28q22}：平方函数的余弦展开]\phantomsection\label{ans:c28q22}
函数为偶函数，所以所有正弦系数为0。常数项及余弦系数分别为
\[
\frac{a_0}2=\frac{4}\pi\int_0^\pi x^2dx=\frac{4\pi^2}3,\qquad a_n=\frac{8}\pi\int_0^\pi x^2\cos(nx)dx=\frac{16(-1)^n}{n^2}.
\]
后一式通过两次分部积分得到。周期延拓连续且分段光滑，傅里叶级数在端点也收敛到函数值。令 $x=\pi$，
\[
4\pi^2=\frac{4\pi^2}3+16\sum_{n=1}^\infty\frac1{n^2},\qquad\sum_{n=1}^\infty\frac1{n^2}=\frac{\pi^2}6.
\]
代入端点的合法性来自收敛定理；形式上写出系数后立即代值，仍缺少这一步依据。
\end{solve}

\begin{solve}[习题~\ref{ex:c28q23}：平方函数的余弦展开]\phantomsection\label{ans:c28q23}
函数为偶函数，所以所有正弦系数为0。常数项及余弦系数分别为
\[
\frac{a_0}2=\frac{5}\pi\int_0^\pi x^2dx=\frac{5\pi^2}3,\qquad a_n=\frac{10}\pi\int_0^\pi x^2\cos(nx)dx=\frac{20(-1)^n}{n^2}.
\]
后一式通过两次分部积分得到。周期延拓连续且分段光滑，傅里叶级数在端点也收敛到函数值。令 $x=\pi$，
\[
5\pi^2=\frac{5\pi^2}3+20\sum_{n=1}^\infty\frac1{n^2},\qquad\sum_{n=1}^\infty\frac1{n^2}=\frac{\pi^2}6.
\]
代入端点的合法性来自收敛定理；形式上写出系数后立即代值，仍缺少这一步依据。
\end{solve}

\begin{solve}[习题~\ref{ex:c28q24}：平方函数的余弦展开]\phantomsection\label{ans:c28q24}
函数为偶函数，所以所有正弦系数为0。常数项及余弦系数分别为
\[
\frac{a_0}2=\frac{6}\pi\int_0^\pi x^2dx=\frac{6\pi^2}3,\qquad a_n=\frac{12}\pi\int_0^\pi x^2\cos(nx)dx=\frac{24(-1)^n}{n^2}.
\]
后一式通过两次分部积分得到。周期延拓连续且分段光滑，傅里叶级数在端点也收敛到函数值。令 $x=\pi$，
\[
6\pi^2=\frac{6\pi^2}3+24\sum_{n=1}^\infty\frac1{n^2},\qquad\sum_{n=1}^\infty\frac1{n^2}=\frac{\pi^2}6.
\]
代入端点的合法性来自收敛定理；形式上写出系数后立即代值，仍缺少这一步依据。
\end{solve}

\clearpage
\section{级数、一致收敛与反例专题}\label{sec:ch29}
本章集中训练“必要条件”“充分条件”“反例”和“交换运算”的辨别。不能把对有限和成立的代数变形不加检查地推广到无穷级数。

\textbf{资源对应：}编者补充证明；与 CLP-2 \S~3.2--\S~3.6 衔接。本章共 20 道编号习题。
\subsection{习题}
\begin{example}[级数收敛的必要条件]\label{ex:c29q1}
证明 $\sum a_n$ 收敛必有 $a_n\to0$，并说明逆命题失败。
\par\textit{补充证明。来源：编者补充：本科数学分析标准命题与反例。}\par\hyperref[ans:c29q1]{解答见第~\pageref{ans:c29q1} 页。}
\end{example}
\begin{example}[级数的 Cauchy 判别]\label{ex:c29q2}
证明 $\sum a_n$ 收敛当且仅当对任意 $\varepsilon>0$，充分大的 $n$ 与任意 $p\ge1$ 都有 $|a_{n+1}+\cdots+a_{n+p}|<\varepsilon$。
\par\textit{补充证明。来源：编者补充：本科数学分析标准命题与反例。}\par\hyperref[ans:c29q2]{解答见第~\pageref{ans:c29q2} 页。}
\end{example}
\begin{example}[绝对收敛定理]\label{ex:c29q3}
证明 $\sum|a_n|$ 收敛蕴含 $\sum a_n$ 收敛。
\par\textit{补充证明。来源：编者补充：本科数学分析标准命题与反例。}\par\hyperref[ans:c29q3]{解答见第~\pageref{ans:c29q3} 页。}
\end{example}
\begin{example}[Leibniz 判别的证明]\label{ex:c29q4}
设 $b_n\ge0$、单调不增且趋于0。证明 $\sum(-1)^{n-1}b_n$ 收敛，并给出余项界。
\par\textit{补充证明。来源：编者补充：本科数学分析标准命题与反例。}\par\hyperref[ans:c29q4]{解答见第~\pageref{ans:c29q4} 页。}
\end{example}
\begin{example}[比值判别为何有效]\label{ex:c29q5}
若 $a_n>0$ 且 $a_{n+1}/a_n\to q<1$，证明 $\sum a_n$ 收敛。
\par\textit{补充证明。来源：编者补充：本科数学分析标准命题与反例。}\par\hyperref[ans:c29q5]{解答见第~\pageref{ans:c29q5} 页。}
\end{example}
\begin{example}[根值判别与统一指数界]\label{ex:c29q6}
若 $\limsup\sqrt[n]{|a_n|}<1$，证明级数绝对收敛。
\par\textit{补充证明。来源：编者补充：本科数学分析标准命题与反例。}\par\hyperref[ans:c29q6]{解答见第~\pageref{ans:c29q6} 页。}
\end{example}
\begin{example}[正项收敛与平方级数]\label{ex:c29q7}
若 $a_n\ge0$ 且 $\sum a_n$ 收敛，证明 $\sum a_n^2$ 收敛；逆命题是否成立？
\par\textit{补充证明。来源：编者补充：本科数学分析标准命题与反例。}\par\hyperref[ans:c29q7]{解答见第~\pageref{ans:c29q7} 页。}
\end{example}
\begin{example}[收敛级数乘有界因子]\label{ex:c29q8}
若 $\sum a_n$ 绝对收敛、$b_n$ 有界，证明 $\sum a_nb_n$ 绝对收敛。若只知条件收敛呢？
\par\textit{补充证明。来源：编者补充：本科数学分析标准命题与反例。}\par\hyperref[ans:c29q8]{解答见第~\pageref{ans:c29q8} 页。}
\end{example}
\begin{example}[绝对收敛级数的重排]\label{ex:c29q9}
给出绝对收敛级数任意重排后仍收敛到原和的证明。
\par\textit{补充证明。来源：编者补充：本科数学分析标准命题与反例。}\par\hyperref[ans:c29q9]{解答见第~\pageref{ans:c29q9} 页。}
\end{example}
\begin{example}[条件收敛级数的正负部分]\label{ex:c29q10}
若实级数 $\sum a_n$ 条件收敛，证明其正项部分和趋于正无穷、负项部分和趋于负无穷。
\par\textit{补充证明。来源：编者补充：本科数学分析标准命题与反例。}\par\hyperref[ans:c29q10]{解答见第~\pageref{ans:c29q10} 页。}
\end{example}
\begin{example}[Weierstrass 判别]\label{ex:c29q11}
若在集合 $E$ 上有 $|f_n(x)|\le M_n$，且 $\sum M_n$ 收敛，证明 $\sum f_n$ 一致且绝对收敛。
\par\textit{补充证明。来源：编者补充：本科数学分析标准命题与反例。}\par\hyperref[ans:c29q11]{解答见第~\pageref{ans:c29q11} 页。}
\end{example}
\begin{example}[幂级数在小闭区间上一致收敛]\label{ex:c29q12}
若幂级数 $\sum a_nx^n$ 的收敛半径为 $R>0$，证明它在每个 $[-r,r]$（$r<R$）上一致绝对收敛。
\par\textit{补充证明。来源：编者补充：本科数学分析标准命题与反例。}\par\hyperref[ans:c29q12]{解答见第~\pageref{ans:c29q12} 页。}
\end{example}
\begin{example}[一致极限保持连续性]\label{ex:c29q13}
若连续函数 $f_n$ 在 $E$ 上一致收敛到 $f$，证明 $f$ 在 $E$ 上连续。
\par\textit{补充证明。来源：编者补充：本科数学分析标准命题与反例。}\par\hyperref[ans:c29q13]{解答见第~\pageref{ans:c29q13} 页。}
\end{example}
\begin{example}[逐点收敛的极限可以不连续]\label{ex:c29q14}
用 $f_n(x)=x^n$ 说明连续函数列逐点收敛的极限可以不连续，并用上确界直接判定不一致收敛。
\par\textit{补充证明。来源：编者补充：本科数学分析标准命题与反例。}\par\hyperref[ans:c29q14]{解答见第~\pageref{ans:c29q14} 页。}
\end{example}
\begin{example}[幂级数半径不受有限项影响]\label{ex:c29q15}
证明增加、删除或改变幂级数的有限多个系数不改变其收敛半径。
\par\textit{补充证明。来源：编者补充：本科数学分析标准命题与反例。}\par\hyperref[ans:c29q15]{解答见第~\pageref{ans:c29q15} 页。}
\end{example}
\begin{example}[逐项求导后的半径]\label{ex:c29q16}
说明为何幂级数与其逐项导数具有相同收敛半径。
\par\textit{补充证明。来源：编者补充：本科数学分析标准命题与反例。}\par\hyperref[ans:c29q16]{解答见第~\pageref{ans:c29q16} 页。}
\end{example}
\begin{example}[正项级数与下标加权的通项]\label{ex:c29q17}
给出正项级数收敛，但 $na_n$ 不趋零的例子；说明加上单调性后为何不同。
\par\textit{补充证明。来源：编者补充：本科数学分析标准命题与反例。}\par\hyperref[ans:c29q17]{解答见第~\pageref{ans:c29q17} 页。}
\end{example}
\begin{example}[从幂级数唯一性比较系数]\label{ex:c29q18}
若两个幂级数在0的某邻域内表示同一函数，证明对应系数相等。
\par\textit{补充证明。来源：编者补充：本科数学分析标准命题与反例。}\par\hyperref[ans:c29q18]{解答见第~\pageref{ans:c29q18} 页。}
\end{example}
\begin{example}[一致收敛的线性组合]\label{ex:c29q19}
若 $f_n\to f$ 与 $g_n\to g$ 都一致收敛，证明任意常数 $a,b$ 下 $af_n+bg_n\to af+bg$ 一致收敛。
\par\textit{补充证明。来源：编者补充：本科数学分析标准命题与反例。}\par\hyperref[ans:c29q19]{解答见第~\pageref{ans:c29q19} 页。}
\end{example}
\begin{example}[一致收敛乘积所需的有界性]\label{ex:c29q20}
证明：若 $f_n,g_n$ 一致收敛到有界函数 $f,g$，则乘积也一致收敛。给出无界极限时的反例。
\par\textit{补充证明。来源：编者补充：本科数学分析标准命题与反例。}\par\hyperref[ans:c29q20]{解答见第~\pageref{ans:c29q20} 页。}
\end{example}
\subsection{解答}
\begin{solve}[习题~\ref{ex:c29q1}：级数收敛的必要条件]\phantomsection\label{ans:c29q1}
设部分和 $S_n\to S$，则 $a_n=S_n-S_{n-1}\to S-S=0$。反过来，$a_n=1/n\to0$，但调和级数发散：从 $n=2^{k-1}+1$ 到 $2^k$ 的一组共有 $2^{k-1}$ 项，每项不小于 $1/2^k$，故每组和至少为1/2。无限多个组使部分和无界。通项趋零是每项必须满足的局部要求，级数收敛还控制整个尾部的累积。
\end{solve}

\begin{solve}[习题~\ref{ex:c29q2}：级数的 Cauchy 判别]\phantomsection\label{ans:c29q2}
令 $S_n=\sum_{k=1}^na_k$，题中的尾和恰为 $S_{n+p}-S_n$。给定的统一条件就是部分和数列的 Cauchy 条件。实数或复数的完备性保证 Cauchy 数列与收敛数列等价，所以结论成立。“任意 $p$”不能省略：如果只固定一个 $p$，通项趋零就可使这有限个项的和趋零，却不能排除调和级数等发散情形。
\end{solve}

\begin{solve}[习题~\ref{ex:c29q3}：绝对收敛定理]\phantomsection\label{ans:c29q3}
绝对值级数的 Cauchy 条件说明，任意足够靠后的有限尾和 $\sum_{k=n+1}^{n+p}|a_k|$ 都能统一小于 $\varepsilon$。三角不等式给出 $|\sum_{k=n+1}^{n+p}a_k|\le\sum_{k=n+1}^{n+p}|a_k|$，故原级数也满足 Cauchy 条件，从而收敛。证明同时适用于复数项级数；它依赖范数三角不等式与标量域的完备性。
\end{solve}

\begin{solve}[习题~\ref{ex:c29q4}：Leibniz 判别的证明]\phantomsection\label{ans:c29q4}
偶数次部分和 $S_{2n}$ 每次增加 $b_{2n+1}-b_{2n+2}\ge0$，故单调不减；奇数次部分和 $S_{2n+1}$ 单调不增，且任意偶数次部分和不超过相应奇数次部分和。两列有界，所以各自收敛；它们之差 $b_{2n+1}\to0$，故极限相同，原部分和收敛。总和位于相邻两部分和之间，因此 $|S-S_n|\le b_{n+1}$，余项与第一舍去项同号。
\end{solve}

\begin{solve}[习题~\ref{ex:c29q5}：比值判别为何有效]\phantomsection\label{ans:c29q5}
选择 $r$ 满足 $q<r<1$。从某项 $N$ 起，相邻比值不超过 $r$，所以 $a_{N+k}\le a_Nr^k$。尾和被收敛几何级数 $a_N\sum r^k$ 控制，故收敛，前面有限项不影响结论。证明需要严格小于1来留出一个统一的几何比 $r$。当极限为1时，既有调和级数发散，又有平方倒数级数收敛，所以此判别不能下结论。
\end{solve}

\begin{solve}[习题~\ref{ex:c29q6}：根值判别与统一指数界]\phantomsection\label{ans:c29q6}
取一个严格介于该上极限与1之间的数 $r$。由上极限定义，从某项起所有 $n$ 都有 $\sqrt[n]{|a_n|}\le r$，于是 $|a_n|\le r^n$。几何级数比较给出绝对收敛。使用上极限而不是通常极限允许根值振荡；重要的是最终全部根值都被同一个小于1的常数控制，而非仅有某条子列满足该界。
\end{solve}

\begin{solve}[习题~\ref{ex:c29q7}：正项收敛与平方级数]\phantomsection\label{ans:c29q7}
原级数收敛使 $a_n\to0$，因此从某项起 $0\le a_n\le1$，于是 $a_n^2\le a_n$。比较判别加上有限项不影响收敛性，得到平方级数收敛。逆命题错误：取 $a_n=1/n$，平方级数收敛，而原级数发散。若只知道 $\sum a_n$ 条件收敛且允许变号，也不能直接使用同样比较；例如 $a_n=(-1)^n/\sqrt n$ 收敛，但平方级数是调和级数。
\end{solve}

\begin{solve}[习题~\ref{ex:c29q8}：收敛级数乘有界因子]\phantomsection\label{ans:c29q8}
取 $M$ 使 $|b_n|\le M$，则 $\sum|a_nb_n|\le M\sum|a_n|<\infty$，所以绝对收敛。若只知条件收敛，结论可能失败：取 $a_n=(-1)^{n-1}/n$、$b_n=(-1)^{n-1}$，前者级数收敛、后者有界，但乘积级数成为调和级数，发散。有界因子可能破坏原来依靠符号抵消获得的收敛性。
\end{solve}

\begin{solve}[习题~\ref{ex:c29q9}：绝对收敛级数的重排]\phantomsection\label{ans:c29q9}
给定 $\varepsilon>0$，取 $N$ 使原级数第 $N$ 项之后的绝对值总和小于 $\varepsilon$。任意重排的充分长部分和都会包含原来的前 $N$ 项，因为这只是有限个固定位置。此后，重排部分和与原总和的差只由原尾部中尚未纳入的项构成，其绝对值至多为原尾部绝对值和，小于 $\varepsilon$。这证明重排部分和趋于原和。绝对收敛提供了与项的顺序无关的尾部控制。
\end{solve}

\begin{solve}[习题~\ref{ex:c29q10}：条件收敛级数的正负部分]\phantomsection\label{ans:c29q10}
记 $a_n^+=\max(a_n,0)$、$a_n^-=-\min(a_n,0)$，则 $a_n=a_n^+-a_n^-$、$|a_n|=a_n^++a_n^-$。如果两个非负级数都收敛，绝对值级数就收敛，与条件收敛矛盾。如果恰有一个收敛而另一个发散，原级数部分和便趋于某个无穷方向，也与收敛矛盾。因此两个非负部分都发散到正无穷，原带符号的负项和则趋于负无穷。原级数的有限和来自两种无界累积的有序抵消。
\end{solve}

\begin{solve}[习题~\ref{ex:c29q11}：Weierstrass 判别]\phantomsection\label{ans:c29q11}
逐点绝对收敛直接来自比较判别。对任意有限尾和，$|\sum_{k=n+1}^{n+p}f_k(x)|\le\sum_{k=n+1}^{n+p}M_k$，右端与 $x$ 无关，且因数值级数收敛可以对所有 $p$ 同时任意小。这是函数级数的一致 Cauchy 条件，故一致收敛。这里控制序列必须形成收敛级数，只有每项分别有界不够。
\end{solve}

\begin{solve}[习题~\ref{ex:c29q12}：幂级数在小闭区间上一致收敛]\phantomsection\label{ans:c29q12}
选 $\rho$ 满足 $r<\rho<R$。由于级数在 $x=\rho$ 处收敛，其项 $a_n\rho^n$ 有界，设绝对值不超过 $M$。当 $|x|\le r$，有 $|a_nx^n|\le M(r/\rho)^n$，右侧为收敛几何级数。Weierstrass 判别给出一致绝对收敛。证明只处理严格位于半径内部的小闭区间，不能把 $r$ 自动取到 $R$。
\end{solve}

\begin{solve}[习题~\ref{ex:c29q13}：一致极限保持连续性]\phantomsection\label{ans:c29q13}
固定 $x_0\in E$ 和 $\varepsilon>0$。先选一个固定的 $N$，使所有 $x\in E$ 都有 $|f_N(x)-f(x)|<\varepsilon/3$。再用这一个连续函数 $f_N$ 在 $x_0$ 处的连续性，选 $\delta$ 使近点的 $|f_N(x)-f_N(x_0)|<\varepsilon/3$。三角不等式把 $|f(x)-f(x_0)|$ 分成两端逼近误差与中间连续性误差三部分，每部分都小于 $\varepsilon/3$。先固定 $N$ 再选 $\delta$ 的顺序是证明的关键。
\end{solve}

\begin{solve}[习题~\ref{ex:c29q14}：逐点收敛的极限可以不连续]\phantomsection\label{ans:c29q14}
在 $[0,1]$ 上，逐点极限为 $f(x)=0$（$x<1$）、$f(1)=1$，在1处不连续。对每个固定 $n$，虽然 $x^n<1$ 当 $x<1$，但当 $x\to1^-$ 时可任意接近1，所以 $\sup_{x\in[0,1]}|f_n(x)-f(x)|=1$。误差上确界不趋零，因此不一致收敛。上确界不一定由某一点取得，这不影响一致收敛判别。
\end{solve}

\begin{solve}[习题~\ref{ex:c29q15}：幂级数半径不受有限项影响]\phantomsection\label{ans:c29q15}
两个这样的幂级数之差是一个多项式，在每个有限 $x$ 处都取有限值。因此在任意给定 $x$ 处，一个级数收敛当且仅当另一个收敛，绝对收敛性也同样等价，因为只改变有限个项。于是它们的收敛域和半径相同。端点和函数值可以发生一个多项式大小的改变，但端点是否收敛不会改变。
\end{solve}

\begin{solve}[习题~\ref{ex:c29q16}：逐项求导后的半径]\phantomsection\label{ans:c29q16}
若原系数为 $a_n$，导数系数带来一个因子 $n$。根值判别中 $n^{1/n}\to1$，所以这个多项式增长因子不改变系数的指数增长率。更严谨地，若 $r<\rho<R$，以 $|a_n|\rho^n\le M$ 控制导数项，得到 $|na_nx^{n-1}|\le(M/\rho)n(r/\rho)^{n-1}$，其和收敛；反过来，对导数幂级数逐项积分恢复原级数加一个常数，同理得相反方向的半径不等式。端点敛散性则可能改变。
\end{solve}

\begin{solve}[习题~\ref{ex:c29q17}：正项级数与下标加权的通项]\phantomsection\label{ans:c29q17}
令 $a_n=1/n$ 当 $n=2^k$，其余位置令 $a_n=1/n^2$。它严格为正，级数被 $\sum_k2^{-k}+\sum_n1/n^2$ 控制而收敛，但在 $n=2^k$ 时 $na_n=1$。若另有 $a_n$ 单调不增，则对 $m=\lfloor n/2\rfloor$，有 $(n-m)a_n\le\sum_{k=m+1}^na_k$，右边是收敛级数的尾部，趋零，而 $n/(n-m)\le2$，因此 $na_n\to0$。单调性排除了稀疏尖峰。
\end{solve}

\begin{solve}[习题~\ref{ex:c29q18}：从幂级数唯一性比较系数]\phantomsection\label{ans:c29q18}
相减后得到 $\sum c_nx^n=0$。首先令 $x=0$，得 $c_0=0$。若已有 $c_0=\cdots=c_{m-1}=0$，对非零 $x$ 除以 $x^m$，得 $c_m+\sum_{n>m}c_nx^{n-m}=0$。余下幂级数在0附近连续且在0取0，令 $x\to0$，得到 $c_m=0$。由归纳全部系数为零。也可以利用半径内逐项求导，得到系数等于函数各阶导数除以阶乘。
\end{solve}

\begin{solve}[习题~\ref{ex:c29q19}：一致收敛的线性组合]\phantomsection\label{ans:c29q19}
对每个 $x$，误差不超过 $|a||f_n(x)-f(x)|+|b||g_n(x)-g(x)|$。取上确界得 $\|af_n+bg_n-(af+bg)\|_\infty\le|a|\|f_n-f\|_\infty+|b|\|g_n-g\|_\infty\to0$。两个数值误差都趋零，有限线性组合仍趋零。这一结论不要求函数本身有界，只要求给定的统一误差趋零。
\end{solve}

\begin{solve}[习题~\ref{ex:c29q20}：一致收敛乘积所需的有界性]\phantomsection\label{ans:c29q20}
一致收敛与 $f$ 有界保证充分大的 $n$ 有共同界 $\|f_n\|_\infty\le M$。于是 $\|f_ng_n-fg\|_\infty\le M\|g_n-g\|_\infty+\|g\|_\infty\|f_n-f\|_\infty\to0$。若删去有界性，取实轴上的 $f_n(x)=x$、$g_n(x)=1/n$，分别一致收敛到 $x$ 和0，但乘积 $x/n$ 不一致趋零，因为每个 $n$ 的误差在全实轴上无界。
\end{solve}

\clearpage
\section{黎曼积分与极限交换的证明专题}\label{sec:ch30}
积分的存在性、积分值的计算和极限交换是不同的问题。这里采用有界函数的 Riemann 积分，并明确区分正常积分、反常积分以及其他积分概念。

\textbf{资源对应：}编者补充证明；与 CLP-2 \S~1.1--\S~1.3、\S~1.12 衔接。本章共 20 道编号习题。
\subsection{习题}
\begin{example}[连续函数的黎曼可积性]\label{ex:c30q1}
证明闭区间上的连续函数 Riemann 可积。
\par\textit{补充证明。来源：编者补充：本科数学分析标准命题与反例。}\par\hyperref[ans:c30q1]{解答见第~\pageref{ans:c30q1} 页。}
\end{example}
\begin{example}[单调函数的可积性]\label{ex:c30q2}
证明闭区间上的单调函数 Riemann 可积，允许它有跳跃间断。
\par\textit{补充证明。来源：编者补充：本科数学分析标准命题与反例。}\par\hyperref[ans:c30q2]{解答见第~\pageref{ans:c30q2} 页。}
\end{example}
\begin{example}[有限个点的修改不改变积分]\label{ex:c30q3}
设有界可积函数 $f,g$ 在 $[a,b]$ 上只在有限个点不同。证明它们积分相等。
\par\textit{补充证明。来源：编者补充：本科数学分析标准命题与反例。}\par\hyperref[ans:c30q3]{解答见第~\pageref{ans:c30q3} 页。}
\end{example}
\begin{example}[Dirichlet 函数不可积]\label{ex:c30q4}
证明 $f(x)=1$（$x$ 为有理数）、$f(x)=0$（$x$ 为无理数）在 $[0,1]$ 上不 Riemann 可积。
\par\textit{补充证明。来源：编者补充：本科数学分析标准命题与反例。}\par\hyperref[ans:c30q4]{解答见第~\pageref{ans:c30q4} 页。}
\end{example}
\begin{example}[积分绝对值不等式]\label{ex:c30q5}
若 $f$ Riemann 可积，证明 $|f|$ 可积且 $|\int_a^bf|\le\int_a^b|f|$。
\par\textit{补充证明。来源：编者补充：本科数学分析标准命题与反例。}\par\hyperref[ans:c30q5]{解答见第~\pageref{ans:c30q5} 页。}
\end{example}
\begin{example}[非负连续函数积分为零]\label{ex:c30q6}
若 $f\ge0$ 且在 $[a,b]$ 连续，$\int_a^bf=0$，证明 $f$ 恒为零。
\par\textit{补充证明。来源：编者补充：本科数学分析标准命题与反例。}\par\hyperref[ans:c30q6]{解答见第~\pageref{ans:c30q6} 页。}
\end{example}
\begin{example}[积分平均值定理]\label{ex:c30q7}
设 $f$ 在 $[a,b]$ 连续、$a<b$。证明存在 $\xi\in[a,b]$ 使 $\int_a^bf=f(\xi)(b-a)$。
\par\textit{补充证明。来源：编者补充：本科数学分析标准命题与反例。}\par\hyperref[ans:c30q7]{解答见第~\pageref{ans:c30q7} 页。}
\end{example}
\begin{example}[基本定理的局部证明]\label{ex:c30q8}
若 $f$ 在 $[a,b]$ 连续，$F(x)=\int_a^xf(t)dt$，证明内点处 $F^{\prime}(x)=f(x)$。
\par\textit{补充证明。来源：编者补充：本科数学分析标准命题与反例。}\par\hyperref[ans:c30q8]{解答见第~\pageref{ans:c30q8} 页。}
\end{example}
\begin{example}[Newton--Leibniz 公式]\label{ex:c30q9}
若 $f$ 连续且 $G^{\prime}=f$，证明 $\int_a^bf=G(b)-G(a)$。
\par\textit{补充证明。来源：编者补充：本科数学分析标准命题与反例。}\par\hyperref[ans:c30q9]{解答见第~\pageref{ans:c30q9} 页。}
\end{example}
\begin{example}[一致收敛与积分交换]\label{ex:c30q10}
若 $f_n$ 在 $[a,b]$ 上连续，且一致收敛到 $f$，证明 $\int f_n\to\int f$。
\par\textit{补充证明。来源：编者补充：本科数学分析标准命题与反例。}\par\hyperref[ans:c30q10]{解答见第~\pageref{ans:c30q10} 页。}
\end{example}
\begin{example}[逐点收敛不足以交换积分]\label{ex:c30q11}
在 $[0,1]$ 上给出连续函数列逐点趋零，但积分不趋零的例子。
\par\textit{补充证明。来源：编者补充：本科数学分析标准命题与反例。}\par\hyperref[ans:c30q11]{解答见第~\pageref{ans:c30q11} 页。}
\end{example}
\begin{example}[一致收敛不自动保持导数]\label{ex:c30q12}
举出可微函数列一致收敛，但其导数列不在每点收敛的例子。
\par\textit{补充证明。来源：编者补充：本科数学分析标准命题与反例。}\par\hyperref[ans:c30q12]{解答见第~\pageref{ans:c30q12} 页。}
\end{example}
\begin{example}[逐项求导的一个充分条件]\label{ex:c30q13}
设 $f_n\in C^1[a,b]$，$f_n(x_0)$ 收敛，且 $f_n^{\prime}$ 一致收敛到连续函数 $g$。证明 $f_n$ 一致收敛到可微函数 $f$，且 $f^{\prime}=g$。
\par\textit{补充证明。来源：编者补充：本科数学分析标准命题与反例。}\par\hyperref[ans:c30q13]{解答见第~\pageref{ans:c30q13} 页。}
\end{example}
\begin{example}[反常积分的内部奇点]\label{ex:c30q14}
解释为什么 $\int_{-1}^1dx/x$ 不作为通常的反常积分收敛，而对称主值为0。
\par\textit{补充证明。来源：编者补充：本科数学分析标准命题与反例。}\par\hyperref[ans:c30q14]{解答见第~\pageref{ans:c30q14} 页。}
\end{example}
\begin{example}[绝对收敛蕴含积分收敛]\label{ex:c30q15}
设 $f$ 在每个 $[a,R]$ 上可积，且 $\int_a^\infty|f|$ 有限。证明 $\int_a^\infty f$ 收敛。
\par\textit{补充证明。来源：编者补充：本科数学分析标准命题与反例。}\par\hyperref[ans:c30q15]{解答见第~\pageref{ans:c30q15} 页。}
\end{example}
\begin{example}[正弦除以变量的条件收敛]\label{ex:c30q16}
证明 $\int_1^\infty\sin x\,dx/x$ 收敛，但不绝对收敛。
\par\textit{补充证明。来源：编者补充：本科数学分析标准命题与反例。}\par\hyperref[ans:c30q16]{解答见第~\pageref{ans:c30q16} 页。}
\end{example}
\begin{example}[积分比较的条件边界]\label{ex:c30q17}
若 $0\le f\le g$ 且 $\int_a^\infty g$ 收敛，证明 $\int_a^\infty f$ 收敛。为什么非负条件有作用？
\par\textit{补充证明。来源：编者补充：本科数学分析标准命题与反例。}\par\hyperref[ans:c30q17]{解答见第~\pageref{ans:c30q17} 页。}
\end{example}
\begin{example}[求和与积分的有限交换]\label{ex:c30q18}
设 $f_1,\ldots,f_m$ 都在 $[a,b]$ 可积，证明积分与有限求和可交换；说明无限和还缺少什么。
\par\textit{补充证明。来源：编者补充：本科数学分析标准命题与反例。}\par\hyperref[ans:c30q18]{解答见第~\pageref{ans:c30q18} 页。}
\end{example}
\begin{example}[积分定义的函数为 Lipschitz]\label{ex:c30q19}
若 $f$ 可积且 $|f|\le M$，证明 $F(x)=\int_a^xf$ 满足 Lipschitz 条件。是否能推出处处可导？
\par\textit{补充证明。来源：编者补充：本科数学分析标准命题与反例。}\par\hyperref[ans:c30q19]{解答见第~\pageref{ans:c30q19} 页。}
\end{example}
\begin{example}[非负积分与函数趋零]\label{ex:c30q20}
若连续函数 $f\ge0$ 且 $\int_0^\infty f<\infty$，是否必有 $f(x)\to0$？给出反例。
\par\textit{补充证明。来源：编者补充：本科数学分析标准命题与反例。}\par\hyperref[ans:c30q20]{解答见第~\pageref{ans:c30q20} 页。}
\end{example}
\subsection{解答}
\begin{solve}[习题~\ref{ex:c30q1}：连续函数的黎曼可积性]\phantomsection\label{ans:c30q1}
设区间长度为 $b-a>0$。连续函数在闭区间上一致连续。给定 $\varepsilon>0$，选 $\delta>0$，使距离小于 $\delta$ 的两点函数值相差小于 $\varepsilon/(b-a)$。取网格小于 $\delta$ 的有限分割，在每个小区间上振幅均不超过该界，因此上、下 Darboux 和之差不超过 $\sum\varepsilon\Delta x_i/(b-a)=\varepsilon$。由 Darboux 可积判别，函数可积。区间退化为一点时积分为0。证明真正使用的是统一的振幅控制，而非每点各自不同的连续性邻域。
\end{solve}

\begin{solve}[习题~\ref{ex:c30q2}：单调函数的可积性]\phantomsection\label{ans:c30q2}
以单调不减函数为例，它自动有界。将 $[a,b]$ 等分为 $n$ 份，每份长度 $h=(b-a)/n$。每个小区间的上、下确界分别为右、左端点函数值，因此上和减下和等于 $h\sum[f(x_i)-f(x_{i-1})]=h(f(b)-f(a))\to0$。由 Darboux 判别可积。单调不增情形对负函数应用即可。这个证明说明连续是可积的充分条件，单调函数即使不连续也仍可能可积。
\end{solve}

\begin{solve}[习题~\ref{ex:c30q3}：有限个点的修改不改变积分]\phantomsection\label{ans:c30q3}
令 $h=f-g$，它仅在有限点可能非零，并有统一上界 $|h|\le M$。给定 $\varepsilon>0$，把这些点包在有限个小区间中，使总长度小于 $\varepsilon/(M+1)$。补上区间端点形成分割；在其余小区间上 $h$ 恒为零，故 $|h|$ 的上和至多为 $M$ 乘这些小区间总长度，小于 $\varepsilon$。所以 $h$ 可积且积分为0，结论成立。有限点的取值影响函数本身，却不会积累成非零长度上的面积贡献。
\end{solve}

\begin{solve}[习题~\ref{ex:c30q4}：Dirichlet 函数不可积]\phantomsection\label{ans:c30q4}
每个非退化小区间都同时包含有理数与无理数，所以在任何分割的每个小区间上，上确界为1、下确界为0。因此任何上和都为1、任何下和都为0，其差恒为1，不能任意小。由 Darboux 判别不可积。函数虽然有界，但处处剧烈振荡，说明有界性本身不足以保证 Riemann 可积。这里的“不积分”特指 Riemann 积分，不能推广为任何积分理论中都不可积。
\end{solve}

\begin{solve}[习题~\ref{ex:c30q5}：积分绝对值不等式]\phantomsection\label{ans:c30q5}
绝对值函数满足 $||u|-|v||\le|u-v|$，所以每个小区间上 $|f|$ 的振幅不超过 $f$ 的振幅。由 $f$ 的 Darboux 可积判别，$|f|$ 同样可积。逐点有 $-|f|\le f\le|f|$，积分的保序性给出 $-\int|f|\le\int f\le\int|f|$，即所求绝对值不等式。反过来 $|f|$ 可积不一定推出 $f$ 可积，例如取有理点为1、无理点为负1，绝对值恒为1，但原函数不可积。
\end{solve}

\begin{solve}[习题~\ref{ex:c30q6}：非负连续函数积分为零]\phantomsection\label{ans:c30q6}
反设存在 $x_0$ 使 $f(x_0)>0$。连续性保证在 $x_0$ 的一个相对邻域内 $f(x)>f(x_0)/2$，该邻域在区间内有正长度，因而其积分为正。区间其他部分积分非负，所以全区间积分严格为正，矛盾。如果删去连续性，可取仅在一个点等于1、其余为0的函数，它可积且积分为零却不是处处为零。结论的“处处”依赖连续性。
\end{solve}

\begin{solve}[习题~\ref{ex:c30q7}：积分平均值定理]\phantomsection\label{ans:c30q7}
最值定理给出 $m\le f(x)\le M$。积分并除以正数 $b-a$ 得 $m\le A:=(b-a)^{-1}\int f\le M$。连续函数在最小值点与最大值点之间取遍中间值，所以存在 $\xi$ 使 $f(\xi)=A$。这一证明说明平均值是函数值域中的某个值。若函数不连续，平均值未必被取得，例如仅取0与1的阶跃函数可以具有1/2的平均值。
\end{solve}

\begin{solve}[习题~\ref{ex:c30q8}：基本定理的局部证明]\phantomsection\label{ans:c30q8}
对足够小的非零 $h$，差商为 $[F(x+h)-F(x)]/h=h^{-1}\int_x^{x+h}f(t)dt$。减去 $f(x)$ 后，其绝对值不超过连接 $x$ 与 $x+h$ 的小区间上 $|f(t)-f(x)|$ 的上确界；负 $h$ 的方向符号可通过绝对值与区间长度抵消。连续性使这个上确界趋零，所以差商趋于 $f(x)$。证明只用到 $f$ 在目标点附近的连续性。
\end{solve}

\begin{solve}[习题~\ref{ex:c30q9}：Newton--Leibniz 公式]\phantomsection\label{ans:c30q9}
令 $F(x)=\int_a^xf(t)dt$，由基本定理 $F^{\prime}=f$。于是 $G-F$ 的导数恒为零，在区间上为常数。代入 $x=a$ 得该常数为 $G(a)$，从而 $F(b)=G(b)-G(a)$，而 $F(b)$ 正是所求积分。这里先构造一个由积分定义的原函数，再利用原函数之差为常数，把它与给定原函数联系起来。
\end{solve}

\begin{solve}[习题~\ref{ex:c30q10}：一致收敛与积分交换]\phantomsection\label{ans:c30q10}
一致极限 $f$ 连续，因此所有积分都存在。记 $\|f_n-f\|_\infty=\sup_{[a,b]}|f_n-f|$，有 $|\int_a^bf_n-\int_a^bf|\le\int_a^b|f_n-f|\le(b-a)\|f_n-f\|_\infty\to0$。关键是误差的统一上界与有限区间长度。若区间无界，即使上确界误差趋零，这个估计也不能直接使用，因为区间长度不再有限。
\end{solve}

\begin{solve}[习题~\ref{ex:c30q11}：逐点收敛不足以交换积分]\phantomsection\label{ans:c30q11}
对 $n\ge2$，定义三角峰函数 $f_n(x)=n^2x$（$0\le x\le1/n$），$f_n(x)=2n-n^2x$（$1/n<x\le2/n$），其余为0。它在拼接点连续，且对任意固定 $x>0$，最终 $x>2/n$，所以 $f_n(x)=0$；在0处也恒为0。故逐点趋零。然而图像是底长 $2/n$、高 $n$ 的三角形，积分恒为1。峰越来越窄却越来越高，逐点观察会遗漏保留下来的总面积。
\end{solve}

\begin{solve}[习题~\ref{ex:c30q12}：一致收敛不自动保持导数]\phantomsection\label{ans:c30q12}
令 $f_n(x)=\sin(nx)/n$。在实轴上 $|f_n(x)|\le1/n$，所以一致收敛到0。但 $f_n^{\prime}(x)=\cos(nx)$，例如在 $x=\pi$ 处等于 $(-1)^n$，没有极限。即使函数本身的误差很小，其振荡频率也可能越来越高，因此导数误差不受函数值误差自动控制。要交换极限与求导，必须增加对导数列的统一收敛等条件。
\end{solve}

\begin{solve}[习题~\ref{ex:c30q13}：逐项求导的一个充分条件]\phantomsection\label{ans:c30q13}
由基本定理，$f_n(x)=f_n(x_0)+\int_{x_0}^xf_n^{\prime}(t)dt$。设点值极限为 $c$，定义 $f(x)=c+\int_{x_0}^xg(t)dt$。则 $\sup_x|f_n(x)-f(x)|\le|f_n(x_0)-c|+(b-a)\sup_t|f_n^{\prime}(t)-g(t)|\to0$，所以一致收敛。由于 $g$ 连续，基本定理给出 $f^{\prime}=g$。单点收敛控制积分常数，导数一致收敛控制区间内的变化，两者分工明确。
\end{solve}

\begin{solve}[习题~\ref{ex:c30q14}：反常积分的内部奇点]\phantomsection\label{ans:c30q14}
通常定义要求在内部奇点处分开，分别考察 $\int_{-1}^0dx/x$ 与 $\int_0^1dx/x$。前者趋于负无穷，后者趋于正无穷，两者都不收敛，所以整体不收敛。若强制以相同距离 $\varepsilon$ 截去两侧，则 $\int_{-1}^{-\varepsilon}dx/x+\int_\varepsilon^1dx/x=0$，其极限为0，称为 Cauchy 主值。这是一种额外规定的对称抵消，不等同于两个独立端点极限都存在。
\end{solve}

\begin{solve}[习题~\ref{ex:c30q15}：绝对收敛蕴含积分收敛]\phantomsection\label{ans:c30q15}
对 $S>R$，绝对值不等式给出 $|\int_R^Sf|\le\int_R^S|f|$。非负函数绝对积分收敛意味着右侧尾部积分可统一小于任意 $\varepsilon$，从而截断积分 $F(R)=\int_a^Rf$ 满足无穷远的 Cauchy 条件。实数完备性保证 $F(R)$ 有有限极限。这个论证不要求 $f$ 本身保持同号，只需要正负振荡的绝对总量有限。
\end{solve}

\begin{solve}[习题~\ref{ex:c30q16}：正弦除以变量的条件收敛]\phantomsection\label{ans:c30q16}
在 $[1,R]$ 分部积分：$\int_1^R\sin x\,dx/x=[-\cos x/x]_1^R-\int_1^R\cos x\,dx/x^2$。边界项有极限，后一积分由 $1/x^2$ 控制而绝对收敛，所以原积分收敛。对绝对值积分，在每个 $[k\pi+\pi/6,k\pi+5\pi/6]$ 上 $|\sin x|\ge1/2$，且 $x\le(k+1)\pi$，故该段贡献至少为 $1/[3(k+1)]$。这些下界形成发散的调和型级数，所以绝对积分发散。
\end{solve}

\begin{solve}[习题~\ref{ex:c30q17}：积分比较的条件边界]\phantomsection\label{ans:c30q17}
截断积分 $F(R)=\int_a^Rf$ 随 $R$ 单调不减，并满足 $0\le F(R)\le\int_a^\infty g<\infty$，所以有有限极限。非负性用于单调性，也使“大函数收敛控制小函数”有明确方向。若允许 $f$ 任意为负，仅有 $f\le g$ 无法控制负向大小，例如 $f=-1,g=0$，上面的大小关系成立，但 $f$ 的积分发散。对变号函数通常应比较绝对值。
\end{solve}

\begin{solve}[习题~\ref{ex:c30q18}：求和与积分的有限交换]\phantomsection\label{ans:c30q18}
有限和 $f=\sum_{j=1}^mc_jf_j$ 的任一黎曼和是各函数黎曼和的同一线性组合；随着网格趋零，每个黎曼和收敛，有限线性组合的极限等于各极限的组合，故 $\int f=\sum c_j\int f_j$。无限和则不能只重复这一步，因为需要把部分和极限移过积分号。一个充分条件是部分和一致收敛；其他积分理论中还有支配收敛等条件。缺少相应控制时，逐点收敛不够。
\end{solve}

\begin{solve}[习题~\ref{ex:c30q19}：积分定义的函数为 Lipschitz]\phantomsection\label{ans:c30q19}
对任意 $x,y$，$|F(x)-F(y)|=|\int_y^xf(t)dt|\le M|x-y|$，所以 $F$ 为 Lipschitz，特别是一致连续。但不一定处处可导。取在0两侧分别为负1与1的符号函数（0处任意），其从0到 $x$ 的积分为 $|x|$，在0不可导。若被积函数在某点连续，则基本定理保证积分函数在该点可导。
\end{solve}

\begin{solve}[习题~\ref{ex:c30q20}：非负积分与函数趋零]\phantomsection\label{ans:c30q20}
不一定。在每个正整数 $n$ 附近放置一个高度1、底边长度 $2^{-n}$ 的三角形尖峰，其余位置为0，各峰支撑区间互不相交。得到连续非负函数，积分为各三角形面积之和 $\sum2^{-n-1}<\infty$，但 $f(n)=1$，所以不趋零。积分控制总面积，不自动控制越来越窄的尖峰高度。若另加一致连续性，则可以证明趋零，因为固定高度的峰会由一致连续性强迫具有统一正宽度，从而破坏有限积分。
\end{solve}

\clearpage
\section{空间向量、平面与曲线}\label{sec:ch31}
空间微积分的计算建立在向量运算上。点积检测正交、叉积给出法向；参数曲线的速度、切向和曲率分别描述运动快慢、局部方向及弯曲程度。

\textbf{资源对应：}CLP-3 \S~1.1--\S~1.6；CLP-4 \S~1.1、\S~1.3。本章共 24 道编号习题。
\subsection{习题}
\begin{example}[点积、叉积与平行四边形]\label{ex:c31q1}
设 $u=(1,1,0)$、$v=(0,1,1)$。求点积、夹角余弦、叉积和以它们为邻边的平行四边形面积。
\par\textit{改编训练。}\par\hyperref[ans:c31q1]{解答见第~\pageref{ans:c31q1} 页。}
\end{example}
\begin{example}[点积、叉积与平行四边形]\label{ex:c31q2}
设 $u=(2,1,0)$、$v=(0,2,1)$。求点积、夹角余弦、叉积和以它们为邻边的平行四边形面积。
\par\textit{改编训练。}\par\hyperref[ans:c31q2]{解答见第~\pageref{ans:c31q2} 页。}
\end{example}
\begin{example}[点积、叉积与平行四边形]\label{ex:c31q3}
设 $u=(3,1,0)$、$v=(0,3,1)$。求点积、夹角余弦、叉积和以它们为邻边的平行四边形面积。
\par\textit{改编训练。}\par\hyperref[ans:c31q3]{解答见第~\pageref{ans:c31q3} 页。}
\end{example}
\begin{example}[点积、叉积与平行四边形]\label{ex:c31q4}
设 $u=(4,1,0)$、$v=(0,4,1)$。求点积、夹角余弦、叉积和以它们为邻边的平行四边形面积。
\par\textit{改编训练。}\par\hyperref[ans:c31q4]{解答见第~\pageref{ans:c31q4} 页。}
\end{example}
\begin{example}[点积、叉积与平行四边形]\label{ex:c31q5}
设 $u=(5,1,0)$、$v=(0,5,1)$。求点积、夹角余弦、叉积和以它们为邻边的平行四边形面积。
\par\textit{改编训练。}\par\hyperref[ans:c31q5]{解答见第~\pageref{ans:c31q5} 页。}
\end{example}
\begin{example}[点积、叉积与平行四边形]\label{ex:c31q6}
设 $u=(6,1,0)$、$v=(0,6,1)$。求点积、夹角余弦、叉积和以它们为邻边的平行四边形面积。
\par\textit{改编训练。}\par\hyperref[ans:c31q6]{解答见第~\pageref{ans:c31q6} 页。}
\end{example}
\begin{example}[平面方程与点面距离]\label{ex:c31q7}
求经过 $P=(1,0,2)$、法向量为 $n=(1,2,2)$ 的平面方程，并求原点到平面的距离和垂足。
\par\textit{改编训练。}\par\hyperref[ans:c31q7]{解答见第~\pageref{ans:c31q7} 页。}
\end{example}
\begin{example}[平面方程与点面距离]\label{ex:c31q8}
求经过 $P=(2,0,1)$、法向量为 $n=(1,2,2)$ 的平面方程，并求原点到平面的距离和垂足。
\par\textit{改编训练。}\par\hyperref[ans:c31q8]{解答见第~\pageref{ans:c31q8} 页。}
\end{example}
\begin{example}[平面方程与点面距离]\label{ex:c31q9}
求经过 $P=(2,0,3)$、法向量为 $n=(1,2,2)$ 的平面方程，并求原点到平面的距离和垂足。
\par\textit{改编训练。}\par\hyperref[ans:c31q9]{解答见第~\pageref{ans:c31q9} 页。}
\end{example}
\begin{example}[平面方程与点面距离]\label{ex:c31q10}
求经过 $P=(3,0,2)$、法向量为 $n=(1,2,2)$ 的平面方程，并求原点到平面的距离和垂足。
\par\textit{改编训练。}\par\hyperref[ans:c31q10]{解答见第~\pageref{ans:c31q10} 页。}
\end{example}
\begin{example}[平面方程与点面距离]\label{ex:c31q11}
求经过 $P=(4,0,1)$、法向量为 $n=(1,2,2)$ 的平面方程，并求原点到平面的距离和垂足。
\par\textit{改编训练。}\par\hyperref[ans:c31q11]{解答见第~\pageref{ans:c31q11} 页。}
\end{example}
\begin{example}[平面方程与点面距离]\label{ex:c31q12}
求经过 $P=(1,0,4)$、法向量为 $n=(1,2,2)$ 的平面方程，并求原点到平面的距离和垂足。
\par\textit{改编训练。}\par\hyperref[ans:c31q12]{解答见第~\pageref{ans:c31q12} 页。}
\end{example}
\begin{example}[圆柱螺线的弧长与曲率]\label{ex:c31q13}
设 $r(t)=(1\cos t,1\sin t,1t)$，$0\le t\le2\pi$。求弧长、单位切向量和曲率。
\par\textit{改编训练。}\par\hyperref[ans:c31q13]{解答见第~\pageref{ans:c31q13} 页。}
\end{example}
\begin{example}[圆柱螺线的弧长与曲率]\label{ex:c31q14}
设 $r(t)=(2\cos t,2\sin t,1t)$，$0\le t\le2\pi$。求弧长、单位切向量和曲率。
\par\textit{改编训练。}\par\hyperref[ans:c31q14]{解答见第~\pageref{ans:c31q14} 页。}
\end{example}
\begin{example}[圆柱螺线的弧长与曲率]\label{ex:c31q15}
设 $r(t)=(1\cos t,1\sin t,2t)$，$0\le t\le2\pi$。求弧长、单位切向量和曲率。
\par\textit{改编训练。}\par\hyperref[ans:c31q15]{解答见第~\pageref{ans:c31q15} 页。}
\end{example}
\begin{example}[圆柱螺线的弧长与曲率]\label{ex:c31q16}
设 $r(t)=(2\cos t,2\sin t,3t)$，$0\le t\le2\pi$。求弧长、单位切向量和曲率。
\par\textit{改编训练。}\par\hyperref[ans:c31q16]{解答见第~\pageref{ans:c31q16} 页。}
\end{example}
\begin{example}[圆柱螺线的弧长与曲率]\label{ex:c31q17}
设 $r(t)=(3\cos t,3\sin t,2t)$，$0\le t\le2\pi$。求弧长、单位切向量和曲率。
\par\textit{改编训练。}\par\hyperref[ans:c31q17]{解答见第~\pageref{ans:c31q17} 页。}
\end{example}
\begin{example}[圆柱螺线的弧长与曲率]\label{ex:c31q18}
设 $r(t)=(3\cos t,3\sin t,4t)$，$0\le t\le2\pi$。求弧长、单位切向量和曲率。
\par\textit{改编训练。}\par\hyperref[ans:c31q18]{解答见第~\pageref{ans:c31q18} 页。}
\end{example}
\begin{example}[圆的定向切向与曲率]\label{ex:c31q19}
对 $r(t)=(1\sin t,1\cos t)$，求 $t=0$ 处的单位切向量、主法向量、曲率和曲率半径。
\par\textit{改编训练。来源：CLP-4 \S~1.3，标签 prob\_s1.3:constantsa 的参数改编。}\par\hyperref[ans:c31q19]{解答见第~\pageref{ans:c31q19} 页。}
\end{example}
\begin{example}[圆的定向切向与曲率]\label{ex:c31q20}
对 $r(t)=(2\sin t,2\cos t)$，求 $t=0$ 处的单位切向量、主法向量、曲率和曲率半径。
\par\textit{改编训练。来源：CLP-4 \S~1.3，标签 prob\_s1.3:constantsa 的参数改编。}\par\hyperref[ans:c31q20]{解答见第~\pageref{ans:c31q20} 页。}
\end{example}
\begin{example}[圆的定向切向与曲率]\label{ex:c31q21}
对 $r(t)=(3\sin t,3\cos t)$，求 $t=0$ 处的单位切向量、主法向量、曲率和曲率半径。
\par\textit{改编训练。来源：CLP-4 \S~1.3，标签 prob\_s1.3:constantsa 的参数改编。}\par\hyperref[ans:c31q21]{解答见第~\pageref{ans:c31q21} 页。}
\end{example}
\begin{example}[圆的定向切向与曲率]\label{ex:c31q22}
对 $r(t)=(4\sin t,4\cos t)$，求 $t=0$ 处的单位切向量、主法向量、曲率和曲率半径。
\par\textit{改编训练。来源：CLP-4 \S~1.3，标签 prob\_s1.3:constantsa 的参数改编。}\par\hyperref[ans:c31q22]{解答见第~\pageref{ans:c31q22} 页。}
\end{example}
\begin{example}[圆的定向切向与曲率]\label{ex:c31q23}
对 $r(t)=(5\sin t,5\cos t)$，求 $t=0$ 处的单位切向量、主法向量、曲率和曲率半径。
\par\textit{改编训练。来源：CLP-4 \S~1.3，标签 prob\_s1.3:constantsa 的参数改编。}\par\hyperref[ans:c31q23]{解答见第~\pageref{ans:c31q23} 页。}
\end{example}
\begin{example}[圆的定向切向与曲率]\label{ex:c31q24}
对 $r(t)=(6\sin t,6\cos t)$，求 $t=0$ 处的单位切向量、主法向量、曲率和曲率半径。
\par\textit{改编训练。来源：CLP-4 \S~1.3，标签 prob\_s1.3:constantsa 的参数改编。}\par\hyperref[ans:c31q24]{解答见第~\pageref{ans:c31q24} 页。}
\end{example}
\subsection{解答}
\begin{solve}[习题~\ref{ex:c31q1}：点积、叉积与平行四边形]\phantomsection\label{ans:c31q1}
按坐标定义，$u\cdot v=1$，且两向量长度均为 $\sqrt{2}$，所以
\[
\cos\theta=\frac{1}{2}.
\]
叉积按行列式展开为
\[
u\times v=(1,-1,1),\qquad |u\times v|=\sqrt{3}.
\]
最后一个量就是所求面积。可以用恒等式 $|u\times v|^2=|u|^2|v|^2-(u\cdot v)^2$ 核对：两边都等于 $a^4+a^2+1$。
\end{solve}

\begin{solve}[习题~\ref{ex:c31q2}：点积、叉积与平行四边形]\phantomsection\label{ans:c31q2}
按坐标定义，$u\cdot v=2$，且两向量长度均为 $\sqrt{5}$，所以
\[
\cos\theta=\frac{2}{5}.
\]
叉积按行列式展开为
\[
u\times v=(1,-2,4),\qquad |u\times v|=\sqrt{21}.
\]
最后一个量就是所求面积。可以用恒等式 $|u\times v|^2=|u|^2|v|^2-(u\cdot v)^2$ 核对：两边都等于 $a^4+a^2+1$。
\end{solve}

\begin{solve}[习题~\ref{ex:c31q3}：点积、叉积与平行四边形]\phantomsection\label{ans:c31q3}
按坐标定义，$u\cdot v=3$，且两向量长度均为 $\sqrt{10}$，所以
\[
\cos\theta=\frac{3}{10}.
\]
叉积按行列式展开为
\[
u\times v=(1,-3,9),\qquad |u\times v|=\sqrt{91}.
\]
最后一个量就是所求面积。可以用恒等式 $|u\times v|^2=|u|^2|v|^2-(u\cdot v)^2$ 核对：两边都等于 $a^4+a^2+1$。
\end{solve}

\begin{solve}[习题~\ref{ex:c31q4}：点积、叉积与平行四边形]\phantomsection\label{ans:c31q4}
按坐标定义，$u\cdot v=4$，且两向量长度均为 $\sqrt{17}$，所以
\[
\cos\theta=\frac{4}{17}.
\]
叉积按行列式展开为
\[
u\times v=(1,-4,16),\qquad |u\times v|=\sqrt{273}.
\]
最后一个量就是所求面积。可以用恒等式 $|u\times v|^2=|u|^2|v|^2-(u\cdot v)^2$ 核对：两边都等于 $a^4+a^2+1$。
\end{solve}

\begin{solve}[习题~\ref{ex:c31q5}：点积、叉积与平行四边形]\phantomsection\label{ans:c31q5}
按坐标定义，$u\cdot v=5$，且两向量长度均为 $\sqrt{26}$，所以
\[
\cos\theta=\frac{5}{26}.
\]
叉积按行列式展开为
\[
u\times v=(1,-5,25),\qquad |u\times v|=\sqrt{651}.
\]
最后一个量就是所求面积。可以用恒等式 $|u\times v|^2=|u|^2|v|^2-(u\cdot v)^2$ 核对：两边都等于 $a^4+a^2+1$。
\end{solve}

\begin{solve}[习题~\ref{ex:c31q6}：点积、叉积与平行四边形]\phantomsection\label{ans:c31q6}
按坐标定义，$u\cdot v=6$，且两向量长度均为 $\sqrt{37}$，所以
\[
\cos\theta=\frac{6}{37}.
\]
叉积按行列式展开为
\[
u\times v=(1,-6,36),\qquad |u\times v|=\sqrt{1333}.
\]
最后一个量就是所求面积。可以用恒等式 $|u\times v|^2=|u|^2|v|^2-(u\cdot v)^2$ 核对：两边都等于 $a^4+a^2+1$。
\end{solve}

\begin{solve}[习题~\ref{ex:c31q7}：平面方程与点面距离]\phantomsection\label{ans:c31q7}
平面上的点 $X=(x,y,z)$ 满足 $n\cdot(X-P)=0$，于是平面为
\[
x+2y+2z=5.
\]
法向量长度为3，原点到平面的距离是 $\frac{5}{3}$。垂足在过原点且方向为 $n$ 的直线上，设为 $H=tn$，代入平面得 $9t=5$，所以
\[
H=\left(\frac{5}{9},\frac{10}{9},\frac{10}{9}\right).
\]
距离公式来自法向投影；用垂足的模长再次计算距离，得到同一结果。
\end{solve}

\begin{solve}[习题~\ref{ex:c31q8}：平面方程与点面距离]\phantomsection\label{ans:c31q8}
平面上的点 $X=(x,y,z)$ 满足 $n\cdot(X-P)=0$，于是平面为
\[
x+2y+2z=4.
\]
法向量长度为3，原点到平面的距离是 $\frac{4}{3}$。垂足在过原点且方向为 $n$ 的直线上，设为 $H=tn$，代入平面得 $9t=4$，所以
\[
H=\left(\frac{4}{9},\frac{8}{9},\frac{8}{9}\right).
\]
距离公式来自法向投影；用垂足的模长再次计算距离，得到同一结果。
\end{solve}

\begin{solve}[习题~\ref{ex:c31q9}：平面方程与点面距离]\phantomsection\label{ans:c31q9}
平面上的点 $X=(x,y,z)$ 满足 $n\cdot(X-P)=0$，于是平面为
\[
x+2y+2z=8.
\]
法向量长度为3，原点到平面的距离是 $\frac{8}{3}$。垂足在过原点且方向为 $n$ 的直线上，设为 $H=tn$，代入平面得 $9t=8$，所以
\[
H=\left(\frac{8}{9},\frac{16}{9},\frac{16}{9}\right).
\]
距离公式来自法向投影；用垂足的模长再次计算距离，得到同一结果。
\end{solve}

\begin{solve}[习题~\ref{ex:c31q10}：平面方程与点面距离]\phantomsection\label{ans:c31q10}
平面上的点 $X=(x,y,z)$ 满足 $n\cdot(X-P)=0$，于是平面为
\[
x+2y+2z=7.
\]
法向量长度为3，原点到平面的距离是 $\frac{7}{3}$。垂足在过原点且方向为 $n$ 的直线上，设为 $H=tn$，代入平面得 $9t=7$，所以
\[
H=\left(\frac{7}{9},\frac{14}{9},\frac{14}{9}\right).
\]
距离公式来自法向投影；用垂足的模长再次计算距离，得到同一结果。
\end{solve}

\begin{solve}[习题~\ref{ex:c31q11}：平面方程与点面距离]\phantomsection\label{ans:c31q11}
平面上的点 $X=(x,y,z)$ 满足 $n\cdot(X-P)=0$，于是平面为
\[
x+2y+2z=6.
\]
法向量长度为3，原点到平面的距离是 $2$。垂足在过原点且方向为 $n$ 的直线上，设为 $H=tn$，代入平面得 $9t=6$，所以
\[
H=\left(\frac{2}{3},\frac{4}{3},\frac{4}{3}\right).
\]
距离公式来自法向投影；用垂足的模长再次计算距离，得到同一结果。
\end{solve}

\begin{solve}[习题~\ref{ex:c31q12}：平面方程与点面距离]\phantomsection\label{ans:c31q12}
平面上的点 $X=(x,y,z)$ 满足 $n\cdot(X-P)=0$，于是平面为
\[
x+2y+2z=9.
\]
法向量长度为3，原点到平面的距离是 $3$。垂足在过原点且方向为 $n$ 的直线上，设为 $H=tn$，代入平面得 $9t=9$，所以
\[
H=\left(1,2,2\right).
\]
距离公式来自法向投影；用垂足的模长再次计算距离，得到同一结果。
\end{solve}

\begin{solve}[习题~\ref{ex:c31q13}：圆柱螺线的弧长与曲率]\phantomsection\label{ans:c31q13}
速度与加速度为 $r^{\prime}=(-1\sin t,1\cos t,1)$、$r^{\prime\prime}=(-1\cos t,-1\sin t,0)$，速度模恒为 $\sqrt{2}$。因此
\[
L=2\pi\sqrt{2},\qquad T=\frac{(-1\sin t,1\cos t,1)}{\sqrt{2}}.
\]
叉积模为 $1\sqrt{2}$，故
\[
\kappa=\frac{|r^{\prime}\times r^{\prime\prime}|}{|r^{\prime}|^3}=\frac{1}{2}.
\]
虽然半径固定，曲率仍受螺距参数影响；螺线向轴向拉长后，局部弯曲变小。速度模处处非零，曲率公式的正则性条件满足。
\end{solve}

\begin{solve}[习题~\ref{ex:c31q14}：圆柱螺线的弧长与曲率]\phantomsection\label{ans:c31q14}
速度与加速度为 $r^{\prime}=(-2\sin t,2\cos t,1)$、$r^{\prime\prime}=(-2\cos t,-2\sin t,0)$，速度模恒为 $\sqrt{5}$。因此
\[
L=2\pi\sqrt{5},\qquad T=\frac{(-2\sin t,2\cos t,1)}{\sqrt{5}}.
\]
叉积模为 $2\sqrt{5}$，故
\[
\kappa=\frac{|r^{\prime}\times r^{\prime\prime}|}{|r^{\prime}|^3}=\frac{2}{5}.
\]
虽然半径固定，曲率仍受螺距参数影响；螺线向轴向拉长后，局部弯曲变小。速度模处处非零，曲率公式的正则性条件满足。
\end{solve}

\begin{solve}[习题~\ref{ex:c31q15}：圆柱螺线的弧长与曲率]\phantomsection\label{ans:c31q15}
速度与加速度为 $r^{\prime}=(-1\sin t,1\cos t,2)$、$r^{\prime\prime}=(-1\cos t,-1\sin t,0)$，速度模恒为 $\sqrt{5}$。因此
\[
L=2\pi\sqrt{5},\qquad T=\frac{(-1\sin t,1\cos t,2)}{\sqrt{5}}.
\]
叉积模为 $1\sqrt{5}$，故
\[
\kappa=\frac{|r^{\prime}\times r^{\prime\prime}|}{|r^{\prime}|^3}=\frac{1}{5}.
\]
虽然半径固定，曲率仍受螺距参数影响；螺线向轴向拉长后，局部弯曲变小。速度模处处非零，曲率公式的正则性条件满足。
\end{solve}

\begin{solve}[习题~\ref{ex:c31q16}：圆柱螺线的弧长与曲率]\phantomsection\label{ans:c31q16}
速度与加速度为 $r^{\prime}=(-2\sin t,2\cos t,3)$、$r^{\prime\prime}=(-2\cos t,-2\sin t,0)$，速度模恒为 $\sqrt{13}$。因此
\[
L=2\pi\sqrt{13},\qquad T=\frac{(-2\sin t,2\cos t,3)}{\sqrt{13}}.
\]
叉积模为 $2\sqrt{13}$，故
\[
\kappa=\frac{|r^{\prime}\times r^{\prime\prime}|}{|r^{\prime}|^3}=\frac{2}{13}.
\]
虽然半径固定，曲率仍受螺距参数影响；螺线向轴向拉长后，局部弯曲变小。速度模处处非零，曲率公式的正则性条件满足。
\end{solve}

\begin{solve}[习题~\ref{ex:c31q17}：圆柱螺线的弧长与曲率]\phantomsection\label{ans:c31q17}
速度与加速度为 $r^{\prime}=(-3\sin t,3\cos t,2)$、$r^{\prime\prime}=(-3\cos t,-3\sin t,0)$，速度模恒为 $\sqrt{13}$。因此
\[
L=2\pi\sqrt{13},\qquad T=\frac{(-3\sin t,3\cos t,2)}{\sqrt{13}}.
\]
叉积模为 $3\sqrt{13}$，故
\[
\kappa=\frac{|r^{\prime}\times r^{\prime\prime}|}{|r^{\prime}|^3}=\frac{3}{13}.
\]
虽然半径固定，曲率仍受螺距参数影响；螺线向轴向拉长后，局部弯曲变小。速度模处处非零，曲率公式的正则性条件满足。
\end{solve}

\begin{solve}[习题~\ref{ex:c31q18}：圆柱螺线的弧长与曲率]\phantomsection\label{ans:c31q18}
速度与加速度为 $r^{\prime}=(-3\sin t,3\cos t,4)$、$r^{\prime\prime}=(-3\cos t,-3\sin t,0)$，速度模恒为 $\sqrt{25}$。因此
\[
L=2\pi\sqrt{25},\qquad T=\frac{(-3\sin t,3\cos t,4)}{\sqrt{25}}.
\]
叉积模为 $3\sqrt{25}$，故
\[
\kappa=\frac{|r^{\prime}\times r^{\prime\prime}|}{|r^{\prime}|^3}=\frac{3}{25}.
\]
虽然半径固定，曲率仍受螺距参数影响；螺线向轴向拉长后，局部弯曲变小。速度模处处非零，曲率公式的正则性条件满足。
\end{solve}

\begin{solve}[习题~\ref{ex:c31q19}：圆的定向切向与曲率]\phantomsection\label{ans:c31q19}
曲线是半径 $1$ 的圆。在 $t=0$ 时，位置为 $(0,1)$，速度为 $(1,0)$，故 $T(0)=(1,0)$。一般有 $T(t)=(\cos t,-\sin t)$，其导数在 $0$ 处为 $(0,-1)$，因此主法向量为 $N(0)=(0,-1)$，指向圆心。曲率为
\[
\kappa=\frac1{1},\qquad\rho=1.
\]
切向量由参数增加方向决定，不能仅凭图形任选方向；主法向量与加速度方向在匀速圆周运动中一致。
\end{solve}

\begin{solve}[习题~\ref{ex:c31q20}：圆的定向切向与曲率]\phantomsection\label{ans:c31q20}
曲线是半径 $2$ 的圆。在 $t=0$ 时，位置为 $(0,2)$，速度为 $(2,0)$，故 $T(0)=(1,0)$。一般有 $T(t)=(\cos t,-\sin t)$，其导数在 $0$ 处为 $(0,-1)$，因此主法向量为 $N(0)=(0,-1)$，指向圆心。曲率为
\[
\kappa=\frac1{2},\qquad\rho=2.
\]
切向量由参数增加方向决定，不能仅凭图形任选方向；主法向量与加速度方向在匀速圆周运动中一致。
\end{solve}

\begin{solve}[习题~\ref{ex:c31q21}：圆的定向切向与曲率]\phantomsection\label{ans:c31q21}
曲线是半径 $3$ 的圆。在 $t=0$ 时，位置为 $(0,3)$，速度为 $(3,0)$，故 $T(0)=(1,0)$。一般有 $T(t)=(\cos t,-\sin t)$，其导数在 $0$ 处为 $(0,-1)$，因此主法向量为 $N(0)=(0,-1)$，指向圆心。曲率为
\[
\kappa=\frac1{3},\qquad\rho=3.
\]
切向量由参数增加方向决定，不能仅凭图形任选方向；主法向量与加速度方向在匀速圆周运动中一致。
\end{solve}

\begin{solve}[习题~\ref{ex:c31q22}：圆的定向切向与曲率]\phantomsection\label{ans:c31q22}
曲线是半径 $4$ 的圆。在 $t=0$ 时，位置为 $(0,4)$，速度为 $(4,0)$，故 $T(0)=(1,0)$。一般有 $T(t)=(\cos t,-\sin t)$，其导数在 $0$ 处为 $(0,-1)$，因此主法向量为 $N(0)=(0,-1)$，指向圆心。曲率为
\[
\kappa=\frac1{4},\qquad\rho=4.
\]
切向量由参数增加方向决定，不能仅凭图形任选方向；主法向量与加速度方向在匀速圆周运动中一致。
\end{solve}

\begin{solve}[习题~\ref{ex:c31q23}：圆的定向切向与曲率]\phantomsection\label{ans:c31q23}
曲线是半径 $5$ 的圆。在 $t=0$ 时，位置为 $(0,5)$，速度为 $(5,0)$，故 $T(0)=(1,0)$。一般有 $T(t)=(\cos t,-\sin t)$，其导数在 $0$ 处为 $(0,-1)$，因此主法向量为 $N(0)=(0,-1)$，指向圆心。曲率为
\[
\kappa=\frac1{5},\qquad\rho=5.
\]
切向量由参数增加方向决定，不能仅凭图形任选方向；主法向量与加速度方向在匀速圆周运动中一致。
\end{solve}

\begin{solve}[习题~\ref{ex:c31q24}：圆的定向切向与曲率]\phantomsection\label{ans:c31q24}
曲线是半径 $6$ 的圆。在 $t=0$ 时，位置为 $(0,6)$，速度为 $(6,0)$，故 $T(0)=(1,0)$。一般有 $T(t)=(\cos t,-\sin t)$，其导数在 $0$ 处为 $(0,-1)$，因此主法向量为 $N(0)=(0,-1)$，指向圆心。曲率为
\[
\kappa=\frac1{6},\qquad\rho=6.
\]
切向量由参数增加方向决定，不能仅凭图形任选方向；主法向量与加速度方向在匀速圆周运动中一致。
\end{solve}

\clearpage
\section{多元极限、偏导数与可微性}\label{sec:ch32}
多元极限要求沿所有方式趋近时得到同一结果。两条路径得到不同极限足以否定极限，有限条路径相同却不足以证明存在。可微性要求一个统一的一阶线性近似，强于偏导数分别存在。

\textbf{资源对应：}CLP-3 \S~2.1--\S~2.3、\S~2.6。本章共 25 道编号习题。
\subsection{习题}
\begin{example}[沿坐标轴检验不存在]\label{ex:c32q1}
设 $f(x,y)=x^2/(x^2+y^2)$。求沿两条坐标轴趋近原点的极限，并判断二元极限是否存在。
\par\textit{选译。来源：CLP-3 \S~2.1，坐标轴路径题。}\par\hyperref[ans:c32q1]{解答见第~\pageref{ans:c32q1} 页。}
\end{example}
\begin{example}[极坐标估计证明极限]\label{ex:c32q2}
定义 $f(x,y)=\dfrac{1x^2y}{x^2+y^2}$（$(x,y)\ne(0,0)$）。求其在原点的极限，并给出适用所有路径的估计。
\par\textit{改编训练。}\par\hyperref[ans:c32q2]{解答见第~\pageref{ans:c32q2} 页。}
\end{example}
\begin{example}[极坐标估计证明极限]\label{ex:c32q3}
定义 $f(x,y)=\dfrac{2x^2y}{x^2+y^2}$（$(x,y)\ne(0,0)$）。求其在原点的极限，并给出适用所有路径的估计。
\par\textit{改编训练。}\par\hyperref[ans:c32q3]{解答见第~\pageref{ans:c32q3} 页。}
\end{example}
\begin{example}[极坐标估计证明极限]\label{ex:c32q4}
定义 $f(x,y)=\dfrac{3x^2y}{x^2+y^2}$（$(x,y)\ne(0,0)$）。求其在原点的极限，并给出适用所有路径的估计。
\par\textit{改编训练。}\par\hyperref[ans:c32q4]{解答见第~\pageref{ans:c32q4} 页。}
\end{example}
\begin{example}[极坐标估计证明极限]\label{ex:c32q5}
定义 $f(x,y)=\dfrac{4x^2y}{x^2+y^2}$（$(x,y)\ne(0,0)$）。求其在原点的极限，并给出适用所有路径的估计。
\par\textit{改编训练。}\par\hyperref[ans:c32q5]{解答见第~\pageref{ans:c32q5} 页。}
\end{example}
\begin{example}[极坐标估计证明极限]\label{ex:c32q6}
定义 $f(x,y)=\dfrac{5x^2y}{x^2+y^2}$（$(x,y)\ne(0,0)$）。求其在原点的极限，并给出适用所有路径的估计。
\par\textit{改编训练。}\par\hyperref[ans:c32q6]{解答见第~\pageref{ans:c32q6} 页。}
\end{example}
\begin{example}[极坐标估计证明极限]\label{ex:c32q7}
定义 $f(x,y)=\dfrac{6x^2y}{x^2+y^2}$（$(x,y)\ne(0,0)$）。求其在原点的极限，并给出适用所有路径的估计。
\par\textit{改编训练。}\par\hyperref[ans:c32q7]{解答见第~\pageref{ans:c32q7} 页。}
\end{example}
\begin{example}[直线路径一致仍可能无极限]\label{ex:c32q8}
判断 $f(x,y)=\dfrac{1x^2y}{x^4+y^2}$ 在原点的极限。比较沿直线与沿抛物线的情况。
\par\textit{改编训练。}\par\hyperref[ans:c32q8]{解答见第~\pageref{ans:c32q8} 页。}
\end{example}
\begin{example}[直线路径一致仍可能无极限]\label{ex:c32q9}
判断 $f(x,y)=\dfrac{2x^2y}{x^4+y^2}$ 在原点的极限。比较沿直线与沿抛物线的情况。
\par\textit{改编训练。}\par\hyperref[ans:c32q9]{解答见第~\pageref{ans:c32q9} 页。}
\end{example}
\begin{example}[直线路径一致仍可能无极限]\label{ex:c32q10}
判断 $f(x,y)=\dfrac{3x^2y}{x^4+y^2}$ 在原点的极限。比较沿直线与沿抛物线的情况。
\par\textit{改编训练。}\par\hyperref[ans:c32q10]{解答见第~\pageref{ans:c32q10} 页。}
\end{example}
\begin{example}[直线路径一致仍可能无极限]\label{ex:c32q11}
判断 $f(x,y)=\dfrac{4x^2y}{x^4+y^2}$ 在原点的极限。比较沿直线与沿抛物线的情况。
\par\textit{改编训练。}\par\hyperref[ans:c32q11]{解答见第~\pageref{ans:c32q11} 页。}
\end{example}
\begin{example}[直线路径一致仍可能无极限]\label{ex:c32q12}
判断 $f(x,y)=\dfrac{5x^2y}{x^4+y^2}$ 在原点的极限。比较沿直线与沿抛物线的情况。
\par\textit{改编训练。}\par\hyperref[ans:c32q12]{解答见第~\pageref{ans:c32q12} 页。}
\end{example}
\begin{example}[直线路径一致仍可能无极限]\label{ex:c32q13}
判断 $f(x,y)=\dfrac{6x^2y}{x^4+y^2}$ 在原点的极限。比较沿直线与沿抛物线的情况。
\par\textit{改编训练。}\par\hyperref[ans:c32q13]{解答见第~\pageref{ans:c32q13} 页。}
\end{example}
\begin{example}[二元多项式的偏导与全微分]\label{ex:c32q14}
设 $f(x,y)=1x^2y+2xy^2$。求一阶偏导、二阶偏导和在 $(1,1)$ 处的全微分。
\par\textit{改编训练。}\par\hyperref[ans:c32q14]{解答见第~\pageref{ans:c32q14} 页。}
\end{example}
\begin{example}[二元多项式的偏导与全微分]\label{ex:c32q15}
设 $f(x,y)=2x^2y+3xy^2$。求一阶偏导、二阶偏导和在 $(1,1)$ 处的全微分。
\par\textit{改编训练。}\par\hyperref[ans:c32q15]{解答见第~\pageref{ans:c32q15} 页。}
\end{example}
\begin{example}[二元多项式的偏导与全微分]\label{ex:c32q16}
设 $f(x,y)=3x^2y+1xy^2$。求一阶偏导、二阶偏导和在 $(1,1)$ 处的全微分。
\par\textit{改编训练。}\par\hyperref[ans:c32q16]{解答见第~\pageref{ans:c32q16} 页。}
\end{example}
\begin{example}[二元多项式的偏导与全微分]\label{ex:c32q17}
设 $f(x,y)=2x^2y+5xy^2$。求一阶偏导、二阶偏导和在 $(1,1)$ 处的全微分。
\par\textit{改编训练。}\par\hyperref[ans:c32q17]{解答见第~\pageref{ans:c32q17} 页。}
\end{example}
\begin{example}[二元多项式的偏导与全微分]\label{ex:c32q18}
设 $f(x,y)=4x^2y+3xy^2$。求一阶偏导、二阶偏导和在 $(1,1)$ 处的全微分。
\par\textit{改编训练。}\par\hyperref[ans:c32q18]{解答见第~\pageref{ans:c32q18} 页。}
\end{example}
\begin{example}[二元多项式的偏导与全微分]\label{ex:c32q19}
设 $f(x,y)=3x^2y+5xy^2$。求一阶偏导、二阶偏导和在 $(1,1)$ 处的全微分。
\par\textit{改编训练。}\par\hyperref[ans:c32q19]{解答见第~\pageref{ans:c32q19} 页。}
\end{example}
\begin{example}[偏导存在却不可微]\label{ex:c32q20}
设 $f(0,0)=0$，非原点处 $f(x,y)=1xy/\sqrt{x^2+y^2}$。求原点处的偏导数，并判断连续性与可微性。
\par\textit{改编训练。}\par\hyperref[ans:c32q20]{解答见第~\pageref{ans:c32q20} 页。}
\end{example}
\begin{example}[偏导存在却不可微]\label{ex:c32q21}
设 $f(0,0)=0$，非原点处 $f(x,y)=2xy/\sqrt{x^2+y^2}$。求原点处的偏导数，并判断连续性与可微性。
\par\textit{改编训练。}\par\hyperref[ans:c32q21]{解答见第~\pageref{ans:c32q21} 页。}
\end{example}
\begin{example}[偏导存在却不可微]\label{ex:c32q22}
设 $f(0,0)=0$，非原点处 $f(x,y)=3xy/\sqrt{x^2+y^2}$。求原点处的偏导数，并判断连续性与可微性。
\par\textit{改编训练。}\par\hyperref[ans:c32q22]{解答见第~\pageref{ans:c32q22} 页。}
\end{example}
\begin{example}[偏导存在却不可微]\label{ex:c32q23}
设 $f(0,0)=0$，非原点处 $f(x,y)=4xy/\sqrt{x^2+y^2}$。求原点处的偏导数，并判断连续性与可微性。
\par\textit{改编训练。}\par\hyperref[ans:c32q23]{解答见第~\pageref{ans:c32q23} 页。}
\end{example}
\begin{example}[偏导存在却不可微]\label{ex:c32q24}
设 $f(0,0)=0$，非原点处 $f(x,y)=5xy/\sqrt{x^2+y^2}$。求原点处的偏导数，并判断连续性与可微性。
\par\textit{改编训练。}\par\hyperref[ans:c32q24]{解答见第~\pageref{ans:c32q24} 页。}
\end{example}
\begin{example}[偏导存在却不可微]\label{ex:c32q25}
设 $f(0,0)=0$，非原点处 $f(x,y)=6xy/\sqrt{x^2+y^2}$。求原点处的偏导数，并判断连续性与可微性。
\par\textit{改编训练。}\par\hyperref[ans:c32q25]{解答见第~\pageref{ans:c32q25} 页。}
\end{example}
\subsection{解答}
\begin{solve}[习题~\ref{ex:c32q1}：沿坐标轴检验不存在]\phantomsection\label{ans:c32q1}
沿 $y=0$、$x\ne0$，函数恒等于1；沿 $x=0$、$y\ne0$，函数恒等于0。两条趋近路径的极限不同，所以二元极限不存在。改变原点处的函数值不能修复这一问题，因为二元极限只取决于去心邻域内的行为。
\end{solve}

\begin{solve}[习题~\ref{ex:c32q2}：极坐标估计证明极限]\phantomsection\label{ans:c32q2}
利用 $x^2/(x^2+y^2)\le1$，可得
\[
|f(x,y)|\le 1|y|\le 1\sqrt{x^2+y^2}.
\]
当 $(x,y)\to(0,0)$ 时，右侧趋零，故极限为0。给定 $\varepsilon>0$，取 $\delta=\varepsilon/1$；只要点到原点的距离小于 $\delta$，函数绝对值就小于 $\varepsilon$。这个估计对方向没有限制，因而构成二元极限存在的完整证明。仅验证直线 $y=kx$ 会留下非线性路径没有处理。
\end{solve}

\begin{solve}[习题~\ref{ex:c32q3}：极坐标估计证明极限]\phantomsection\label{ans:c32q3}
利用 $x^2/(x^2+y^2)\le1$，可得
\[
|f(x,y)|\le 2|y|\le 2\sqrt{x^2+y^2}.
\]
当 $(x,y)\to(0,0)$ 时，右侧趋零，故极限为0。给定 $\varepsilon>0$，取 $\delta=\varepsilon/2$；只要点到原点的距离小于 $\delta$，函数绝对值就小于 $\varepsilon$。这个估计对方向没有限制，因而构成二元极限存在的完整证明。仅验证直线 $y=kx$ 会留下非线性路径没有处理。
\end{solve}

\begin{solve}[习题~\ref{ex:c32q4}：极坐标估计证明极限]\phantomsection\label{ans:c32q4}
利用 $x^2/(x^2+y^2)\le1$，可得
\[
|f(x,y)|\le 3|y|\le 3\sqrt{x^2+y^2}.
\]
当 $(x,y)\to(0,0)$ 时，右侧趋零，故极限为0。给定 $\varepsilon>0$，取 $\delta=\varepsilon/3$；只要点到原点的距离小于 $\delta$，函数绝对值就小于 $\varepsilon$。这个估计对方向没有限制，因而构成二元极限存在的完整证明。仅验证直线 $y=kx$ 会留下非线性路径没有处理。
\end{solve}

\begin{solve}[习题~\ref{ex:c32q5}：极坐标估计证明极限]\phantomsection\label{ans:c32q5}
利用 $x^2/(x^2+y^2)\le1$，可得
\[
|f(x,y)|\le 4|y|\le 4\sqrt{x^2+y^2}.
\]
当 $(x,y)\to(0,0)$ 时，右侧趋零，故极限为0。给定 $\varepsilon>0$，取 $\delta=\varepsilon/4$；只要点到原点的距离小于 $\delta$，函数绝对值就小于 $\varepsilon$。这个估计对方向没有限制，因而构成二元极限存在的完整证明。仅验证直线 $y=kx$ 会留下非线性路径没有处理。
\end{solve}

\begin{solve}[习题~\ref{ex:c32q6}：极坐标估计证明极限]\phantomsection\label{ans:c32q6}
利用 $x^2/(x^2+y^2)\le1$，可得
\[
|f(x,y)|\le 5|y|\le 5\sqrt{x^2+y^2}.
\]
当 $(x,y)\to(0,0)$ 时，右侧趋零，故极限为0。给定 $\varepsilon>0$，取 $\delta=\varepsilon/5$；只要点到原点的距离小于 $\delta$，函数绝对值就小于 $\varepsilon$。这个估计对方向没有限制，因而构成二元极限存在的完整证明。仅验证直线 $y=kx$ 会留下非线性路径没有处理。
\end{solve}

\begin{solve}[习题~\ref{ex:c32q7}：极坐标估计证明极限]\phantomsection\label{ans:c32q7}
利用 $x^2/(x^2+y^2)\le1$，可得
\[
|f(x,y)|\le 6|y|\le 6\sqrt{x^2+y^2}.
\]
当 $(x,y)\to(0,0)$ 时，右侧趋零，故极限为0。给定 $\varepsilon>0$，取 $\delta=\varepsilon/6$；只要点到原点的距离小于 $\delta$，函数绝对值就小于 $\varepsilon$。这个估计对方向没有限制，因而构成二元极限存在的完整证明。仅验证直线 $y=kx$ 会留下非线性路径没有处理。
\end{solve}

\begin{solve}[习题~\ref{ex:c32q8}：直线路径一致仍可能无极限]\phantomsection\label{ans:c32q8}
沿 $y=kx$，当 $k\ne0$ 时，$f(x,kx)=1kx/(x^2+k^2)\to0$；当 $k=0$ 时函数恒为0，沿 $x=0$ 也为0。然而沿 $y=x^2$，
\[
f(x,x^2)=\frac{1x^4}{2x^4}=\frac{1}{2}\qquad(x\ne0).
\]
沿 $y=0$ 与 $y=x^2$ 的极限不同，所以二元极限不存在。这个反例表明，即便检查了所有过原点的直线，也不能代替对全部趋近方式的证明。分母中不同变量的次数往往提示应选择非线性路径。
\end{solve}

\begin{solve}[习题~\ref{ex:c32q9}：直线路径一致仍可能无极限]\phantomsection\label{ans:c32q9}
沿 $y=kx$，当 $k\ne0$ 时，$f(x,kx)=2kx/(x^2+k^2)\to0$；当 $k=0$ 时函数恒为0，沿 $x=0$ 也为0。然而沿 $y=x^2$，
\[
f(x,x^2)=\frac{2x^4}{2x^4}=1\qquad(x\ne0).
\]
沿 $y=0$ 与 $y=x^2$ 的极限不同，所以二元极限不存在。这个反例表明，即便检查了所有过原点的直线，也不能代替对全部趋近方式的证明。分母中不同变量的次数往往提示应选择非线性路径。
\end{solve}

\begin{solve}[习题~\ref{ex:c32q10}：直线路径一致仍可能无极限]\phantomsection\label{ans:c32q10}
沿 $y=kx$，当 $k\ne0$ 时，$f(x,kx)=3kx/(x^2+k^2)\to0$；当 $k=0$ 时函数恒为0，沿 $x=0$ 也为0。然而沿 $y=x^2$，
\[
f(x,x^2)=\frac{3x^4}{2x^4}=\frac{3}{2}\qquad(x\ne0).
\]
沿 $y=0$ 与 $y=x^2$ 的极限不同，所以二元极限不存在。这个反例表明，即便检查了所有过原点的直线，也不能代替对全部趋近方式的证明。分母中不同变量的次数往往提示应选择非线性路径。
\end{solve}

\begin{solve}[习题~\ref{ex:c32q11}：直线路径一致仍可能无极限]\phantomsection\label{ans:c32q11}
沿 $y=kx$，当 $k\ne0$ 时，$f(x,kx)=4kx/(x^2+k^2)\to0$；当 $k=0$ 时函数恒为0，沿 $x=0$ 也为0。然而沿 $y=x^2$，
\[
f(x,x^2)=\frac{4x^4}{2x^4}=2\qquad(x\ne0).
\]
沿 $y=0$ 与 $y=x^2$ 的极限不同，所以二元极限不存在。这个反例表明，即便检查了所有过原点的直线，也不能代替对全部趋近方式的证明。分母中不同变量的次数往往提示应选择非线性路径。
\end{solve}

\begin{solve}[习题~\ref{ex:c32q12}：直线路径一致仍可能无极限]\phantomsection\label{ans:c32q12}
沿 $y=kx$，当 $k\ne0$ 时，$f(x,kx)=5kx/(x^2+k^2)\to0$；当 $k=0$ 时函数恒为0，沿 $x=0$ 也为0。然而沿 $y=x^2$，
\[
f(x,x^2)=\frac{5x^4}{2x^4}=\frac{5}{2}\qquad(x\ne0).
\]
沿 $y=0$ 与 $y=x^2$ 的极限不同，所以二元极限不存在。这个反例表明，即便检查了所有过原点的直线，也不能代替对全部趋近方式的证明。分母中不同变量的次数往往提示应选择非线性路径。
\end{solve}

\begin{solve}[习题~\ref{ex:c32q13}：直线路径一致仍可能无极限]\phantomsection\label{ans:c32q13}
沿 $y=kx$，当 $k\ne0$ 时，$f(x,kx)=6kx/(x^2+k^2)\to0$；当 $k=0$ 时函数恒为0，沿 $x=0$ 也为0。然而沿 $y=x^2$，
\[
f(x,x^2)=\frac{6x^4}{2x^4}=3\qquad(x\ne0).
\]
沿 $y=0$ 与 $y=x^2$ 的极限不同，所以二元极限不存在。这个反例表明，即便检查了所有过原点的直线，也不能代替对全部趋近方式的证明。分母中不同变量的次数往往提示应选择非线性路径。
\end{solve}

\begin{solve}[习题~\ref{ex:c32q14}：二元多项式的偏导与全微分]\phantomsection\label{ans:c32q14}
对一个变量求偏导时，把另一个变量暂视为常数。因此
\[
f_x=2xy+2y^2,\qquad f_y=1x^2+4xy,
\]

\[
f_{xx}=2y,\quad f_{xy}=f_{yx}=2x+4y,\quad f_{yy}=4x.
\]
在 $(1,1)$ 处，
\[
df=(4)dx+(5)dy.
\]
多项式的一阶偏导连续，所以函数可微。混合偏导相等在这里可直接计算核对，也符合连续二阶偏导条件下的对称性定理。
\end{solve}

\begin{solve}[习题~\ref{ex:c32q15}：二元多项式的偏导与全微分]\phantomsection\label{ans:c32q15}
对一个变量求偏导时，把另一个变量暂视为常数。因此
\[
f_x=4xy+3y^2,\qquad f_y=2x^2+6xy,
\]

\[
f_{xx}=4y,\quad f_{xy}=f_{yx}=4x+6y,\quad f_{yy}=6x.
\]
在 $(1,1)$ 处，
\[
df=(7)dx+(8)dy.
\]
多项式的一阶偏导连续，所以函数可微。混合偏导相等在这里可直接计算核对，也符合连续二阶偏导条件下的对称性定理。
\end{solve}

\begin{solve}[习题~\ref{ex:c32q16}：二元多项式的偏导与全微分]\phantomsection\label{ans:c32q16}
对一个变量求偏导时，把另一个变量暂视为常数。因此
\[
f_x=6xy+1y^2,\qquad f_y=3x^2+2xy,
\]

\[
f_{xx}=6y,\quad f_{xy}=f_{yx}=6x+2y,\quad f_{yy}=2x.
\]
在 $(1,1)$ 处，
\[
df=(7)dx+(5)dy.
\]
多项式的一阶偏导连续，所以函数可微。混合偏导相等在这里可直接计算核对，也符合连续二阶偏导条件下的对称性定理。
\end{solve}

\begin{solve}[习题~\ref{ex:c32q17}：二元多项式的偏导与全微分]\phantomsection\label{ans:c32q17}
对一个变量求偏导时，把另一个变量暂视为常数。因此
\[
f_x=4xy+5y^2,\qquad f_y=2x^2+10xy,
\]

\[
f_{xx}=4y,\quad f_{xy}=f_{yx}=4x+10y,\quad f_{yy}=10x.
\]
在 $(1,1)$ 处，
\[
df=(9)dx+(12)dy.
\]
多项式的一阶偏导连续，所以函数可微。混合偏导相等在这里可直接计算核对，也符合连续二阶偏导条件下的对称性定理。
\end{solve}

\begin{solve}[习题~\ref{ex:c32q18}：二元多项式的偏导与全微分]\phantomsection\label{ans:c32q18}
对一个变量求偏导时，把另一个变量暂视为常数。因此
\[
f_x=8xy+3y^2,\qquad f_y=4x^2+6xy,
\]

\[
f_{xx}=8y,\quad f_{xy}=f_{yx}=8x+6y,\quad f_{yy}=6x.
\]
在 $(1,1)$ 处，
\[
df=(11)dx+(10)dy.
\]
多项式的一阶偏导连续，所以函数可微。混合偏导相等在这里可直接计算核对，也符合连续二阶偏导条件下的对称性定理。
\end{solve}

\begin{solve}[习题~\ref{ex:c32q19}：二元多项式的偏导与全微分]\phantomsection\label{ans:c32q19}
对一个变量求偏导时，把另一个变量暂视为常数。因此
\[
f_x=6xy+5y^2,\qquad f_y=3x^2+10xy,
\]

\[
f_{xx}=6y,\quad f_{xy}=f_{yx}=6x+10y,\quad f_{yy}=10x.
\]
在 $(1,1)$ 处，
\[
df=(11)dx+(13)dy.
\]
多项式的一阶偏导连续，所以函数可微。混合偏导相等在这里可直接计算核对，也符合连续二阶偏导条件下的对称性定理。
\end{solve}

\begin{solve}[习题~\ref{ex:c32q20}：偏导存在却不可微]\phantomsection\label{ans:c32q20}
在两条坐标轴上函数恒为0，因此 $f_x(0,0)=f_y(0,0)=0$。由 $2|xy|\le x^2+y^2$ 有
\[
|f(x,y)|\le \frac{1}{2}\sqrt{x^2+y^2}\to0,
\]
故函数连续。如果可微，其线性主部只能为0，因而应有 $f(x,y)/\sqrt{x^2+y^2}\to0$。但沿 $y=x\ne0$，
\[
\frac{f(x,x)}{\sqrt{2x^2}}=\frac{1}{2},
\]
所以不可微。这个例子同时满足连续与两偏导存在，仍不足以保证可微；余项必须比到原点的距离更高阶地趋零。
\end{solve}

\begin{solve}[习题~\ref{ex:c32q21}：偏导存在却不可微]\phantomsection\label{ans:c32q21}
在两条坐标轴上函数恒为0，因此 $f_x(0,0)=f_y(0,0)=0$。由 $2|xy|\le x^2+y^2$ 有
\[
|f(x,y)|\le 1\sqrt{x^2+y^2}\to0,
\]
故函数连续。如果可微，其线性主部只能为0，因而应有 $f(x,y)/\sqrt{x^2+y^2}\to0$。但沿 $y=x\ne0$，
\[
\frac{f(x,x)}{\sqrt{2x^2}}=1,
\]
所以不可微。这个例子同时满足连续与两偏导存在，仍不足以保证可微；余项必须比到原点的距离更高阶地趋零。
\end{solve}

\begin{solve}[习题~\ref{ex:c32q22}：偏导存在却不可微]\phantomsection\label{ans:c32q22}
在两条坐标轴上函数恒为0，因此 $f_x(0,0)=f_y(0,0)=0$。由 $2|xy|\le x^2+y^2$ 有
\[
|f(x,y)|\le \frac{3}{2}\sqrt{x^2+y^2}\to0,
\]
故函数连续。如果可微，其线性主部只能为0，因而应有 $f(x,y)/\sqrt{x^2+y^2}\to0$。但沿 $y=x\ne0$，
\[
\frac{f(x,x)}{\sqrt{2x^2}}=\frac{3}{2},
\]
所以不可微。这个例子同时满足连续与两偏导存在，仍不足以保证可微；余项必须比到原点的距离更高阶地趋零。
\end{solve}

\begin{solve}[习题~\ref{ex:c32q23}：偏导存在却不可微]\phantomsection\label{ans:c32q23}
在两条坐标轴上函数恒为0，因此 $f_x(0,0)=f_y(0,0)=0$。由 $2|xy|\le x^2+y^2$ 有
\[
|f(x,y)|\le 2\sqrt{x^2+y^2}\to0,
\]
故函数连续。如果可微，其线性主部只能为0，因而应有 $f(x,y)/\sqrt{x^2+y^2}\to0$。但沿 $y=x\ne0$，
\[
\frac{f(x,x)}{\sqrt{2x^2}}=2,
\]
所以不可微。这个例子同时满足连续与两偏导存在，仍不足以保证可微；余项必须比到原点的距离更高阶地趋零。
\end{solve}

\begin{solve}[习题~\ref{ex:c32q24}：偏导存在却不可微]\phantomsection\label{ans:c32q24}
在两条坐标轴上函数恒为0，因此 $f_x(0,0)=f_y(0,0)=0$。由 $2|xy|\le x^2+y^2$ 有
\[
|f(x,y)|\le \frac{5}{2}\sqrt{x^2+y^2}\to0,
\]
故函数连续。如果可微，其线性主部只能为0，因而应有 $f(x,y)/\sqrt{x^2+y^2}\to0$。但沿 $y=x\ne0$，
\[
\frac{f(x,x)}{\sqrt{2x^2}}=\frac{5}{2},
\]
所以不可微。这个例子同时满足连续与两偏导存在，仍不足以保证可微；余项必须比到原点的距离更高阶地趋零。
\end{solve}

\begin{solve}[习题~\ref{ex:c32q25}：偏导存在却不可微]\phantomsection\label{ans:c32q25}
在两条坐标轴上函数恒为0，因此 $f_x(0,0)=f_y(0,0)=0$。由 $2|xy|\le x^2+y^2$ 有
\[
|f(x,y)|\le 3\sqrt{x^2+y^2}\to0,
\]
故函数连续。如果可微，其线性主部只能为0，因而应有 $f(x,y)/\sqrt{x^2+y^2}\to0$。但沿 $y=x\ne0$，
\[
\frac{f(x,x)}{\sqrt{2x^2}}=3,
\]
所以不可微。这个例子同时满足连续与两偏导存在，仍不足以保证可微；余项必须比到原点的距离更高阶地趋零。
\end{solve}

\clearpage
\section{多元链式法则与 Jacobian}\label{sec:ch33}
Jacobian 矩阵（雅可比矩阵）把一阶偏导按输入、输出排列。复合映射的导数是矩阵乘积，行列式则描述局部面积或体积的缩放；使用逆映射时还需确认行列式非零。

\textbf{资源对应：}CLP-3 \S~2.4--\S~2.6；\S~3.2、\S~3.6 的换元背景。本章共 24 道编号习题。
\subsection{习题}
\begin{example}[线性换元的面积缩放]\label{ex:c33q1}
令 $x=1u+v$、$y=-u+2v$。求 Jacobian 行列式、逆变换，以及单位正方形在此变换下的面积。
\par\textit{改编训练。}\par\hyperref[ans:c33q1]{解答见第~\pageref{ans:c33q1} 页。}
\end{example}
\begin{example}[线性换元的面积缩放]\label{ex:c33q2}
令 $x=2u+v$、$y=-u+1v$。求 Jacobian 行列式、逆变换，以及单位正方形在此变换下的面积。
\par\textit{改编训练。}\par\hyperref[ans:c33q2]{解答见第~\pageref{ans:c33q2} 页。}
\end{example}
\begin{example}[线性换元的面积缩放]\label{ex:c33q3}
令 $x=2u+v$、$y=-u+3v$。求 Jacobian 行列式、逆变换，以及单位正方形在此变换下的面积。
\par\textit{改编训练。}\par\hyperref[ans:c33q3]{解答见第~\pageref{ans:c33q3} 页。}
\end{example}
\begin{example}[线性换元的面积缩放]\label{ex:c33q4}
令 $x=3u+v$、$y=-u+2v$。求 Jacobian 行列式、逆变换，以及单位正方形在此变换下的面积。
\par\textit{改编训练。}\par\hyperref[ans:c33q4]{解答见第~\pageref{ans:c33q4} 页。}
\end{example}
\begin{example}[线性换元的面积缩放]\label{ex:c33q5}
令 $x=4u+v$、$y=-u+1v$。求 Jacobian 行列式、逆变换，以及单位正方形在此变换下的面积。
\par\textit{改编训练。}\par\hyperref[ans:c33q5]{解答见第~\pageref{ans:c33q5} 页。}
\end{example}
\begin{example}[线性换元的面积缩放]\label{ex:c33q6}
令 $x=3u+v$、$y=-u+4v$。求 Jacobian 行列式、逆变换，以及单位正方形在此变换下的面积。
\par\textit{改编训练。}\par\hyperref[ans:c33q6]{解答见第~\pageref{ans:c33q6} 页。}
\end{example}
\begin{example}[极坐标下的链式法则]\label{ex:c33q7}
设 $z=x^2+1y^2$，$x=r\cos\theta$、$y=r\sin\theta$。求 $z_r$、$z_\theta$，并直接代入核对。
\par\textit{改编训练。}\par\hyperref[ans:c33q7]{解答见第~\pageref{ans:c33q7} 页。}
\end{example}
\begin{example}[极坐标下的链式法则]\label{ex:c33q8}
设 $z=x^2+2y^2$，$x=r\cos\theta$、$y=r\sin\theta$。求 $z_r$、$z_\theta$，并直接代入核对。
\par\textit{改编训练。}\par\hyperref[ans:c33q8]{解答见第~\pageref{ans:c33q8} 页。}
\end{example}
\begin{example}[极坐标下的链式法则]\label{ex:c33q9}
设 $z=x^2+3y^2$，$x=r\cos\theta$、$y=r\sin\theta$。求 $z_r$、$z_\theta$，并直接代入核对。
\par\textit{改编训练。}\par\hyperref[ans:c33q9]{解答见第~\pageref{ans:c33q9} 页。}
\end{example}
\begin{example}[极坐标下的链式法则]\label{ex:c33q10}
设 $z=x^2+4y^2$，$x=r\cos\theta$、$y=r\sin\theta$。求 $z_r$、$z_\theta$，并直接代入核对。
\par\textit{改编训练。}\par\hyperref[ans:c33q10]{解答见第~\pageref{ans:c33q10} 页。}
\end{example}
\begin{example}[极坐标下的链式法则]\label{ex:c33q11}
设 $z=x^2+5y^2$，$x=r\cos\theta$、$y=r\sin\theta$。求 $z_r$、$z_\theta$，并直接代入核对。
\par\textit{改编训练。}\par\hyperref[ans:c33q11]{解答见第~\pageref{ans:c33q11} 页。}
\end{example}
\begin{example}[极坐标下的链式法则]\label{ex:c33q12}
设 $z=x^2+6y^2$，$x=r\cos\theta$、$y=r\sin\theta$。求 $z_r$、$z_\theta$，并直接代入核对。
\par\textit{改编训练。}\par\hyperref[ans:c33q12]{解答见第~\pageref{ans:c33q12} 页。}
\end{example}
\begin{example}[隐函数的一阶偏导]\label{ex:c33q13}
由 $z+xz+yz=1x+2y$ 在原点附近定义 $z=z(x,y)$。求显式表达式以及 $z_x(0,0),z_y(0,0)$。
\par\textit{改编训练。}\par\hyperref[ans:c33q13]{解答见第~\pageref{ans:c33q13} 页。}
\end{example}
\begin{example}[隐函数的一阶偏导]\label{ex:c33q14}
由 $z+xz+yz=2x+1y$ 在原点附近定义 $z=z(x,y)$。求显式表达式以及 $z_x(0,0),z_y(0,0)$。
\par\textit{改编训练。}\par\hyperref[ans:c33q14]{解答见第~\pageref{ans:c33q14} 页。}
\end{example}
\begin{example}[隐函数的一阶偏导]\label{ex:c33q15}
由 $z+xz+yz=2x+3y$ 在原点附近定义 $z=z(x,y)$。求显式表达式以及 $z_x(0,0),z_y(0,0)$。
\par\textit{改编训练。}\par\hyperref[ans:c33q15]{解答见第~\pageref{ans:c33q15} 页。}
\end{example}
\begin{example}[隐函数的一阶偏导]\label{ex:c33q16}
由 $z+xz+yz=3x+2y$ 在原点附近定义 $z=z(x,y)$。求显式表达式以及 $z_x(0,0),z_y(0,0)$。
\par\textit{改编训练。}\par\hyperref[ans:c33q16]{解答见第~\pageref{ans:c33q16} 页。}
\end{example}
\begin{example}[隐函数的一阶偏导]\label{ex:c33q17}
由 $z+xz+yz=4x+1y$ 在原点附近定义 $z=z(x,y)$。求显式表达式以及 $z_x(0,0),z_y(0,0)$。
\par\textit{改编训练。}\par\hyperref[ans:c33q17]{解答见第~\pageref{ans:c33q17} 页。}
\end{example}
\begin{example}[隐函数的一阶偏导]\label{ex:c33q18}
由 $z+xz+yz=3x+4y$ 在原点附近定义 $z=z(x,y)$。求显式表达式以及 $z_x(0,0),z_y(0,0)$。
\par\textit{改编训练。}\par\hyperref[ans:c33q18]{解答见第~\pageref{ans:c33q18} 页。}
\end{example}
\begin{example}[二元函数的二阶泰勒式]\label{ex:c33q19}
求 $f(x,y)=e^{1x+2y}$ 在原点的二阶泰勒多项式，并给出余项的阶。
\par\textit{改编训练。}\par\hyperref[ans:c33q19]{解答见第~\pageref{ans:c33q19} 页。}
\end{example}
\begin{example}[二元函数的二阶泰勒式]\label{ex:c33q20}
求 $f(x,y)=e^{2x+1y}$ 在原点的二阶泰勒多项式，并给出余项的阶。
\par\textit{改编训练。}\par\hyperref[ans:c33q20]{解答见第~\pageref{ans:c33q20} 页。}
\end{example}
\begin{example}[二元函数的二阶泰勒式]\label{ex:c33q21}
求 $f(x,y)=e^{2x+3y}$ 在原点的二阶泰勒多项式，并给出余项的阶。
\par\textit{改编训练。}\par\hyperref[ans:c33q21]{解答见第~\pageref{ans:c33q21} 页。}
\end{example}
\begin{example}[二元函数的二阶泰勒式]\label{ex:c33q22}
求 $f(x,y)=e^{3x+2y}$ 在原点的二阶泰勒多项式，并给出余项的阶。
\par\textit{改编训练。}\par\hyperref[ans:c33q22]{解答见第~\pageref{ans:c33q22} 页。}
\end{example}
\begin{example}[二元函数的二阶泰勒式]\label{ex:c33q23}
求 $f(x,y)=e^{4x+1y}$ 在原点的二阶泰勒多项式，并给出余项的阶。
\par\textit{改编训练。}\par\hyperref[ans:c33q23]{解答见第~\pageref{ans:c33q23} 页。}
\end{example}
\begin{example}[二元函数的二阶泰勒式]\label{ex:c33q24}
求 $f(x,y)=e^{3x+4y}$ 在原点的二阶泰勒多项式，并给出余项的阶。
\par\textit{改编训练。}\par\hyperref[ans:c33q24]{解答见第~\pageref{ans:c33q24} 页。}
\end{example}
\subsection{解答}
\begin{solve}[习题~\ref{ex:c33q1}：线性换元的面积缩放]\phantomsection\label{ans:c33q1}
偏导矩阵为 $J=\begin{pmatrix}1&1\\-1&2\end{pmatrix}$，故 $\det J=3>0$。求逆矩阵得
\[
u=\frac{2x-y}{3},\qquad v=\frac{x+1y}{3}.
\]
面积元满足 $dxdy=3\,dudv$，因此单位正方形映成的平行四边形面积为 $3$。对边向量 $(1,-1)$ 与 $(1,2)$ 的二维叉积绝对值也给出同样面积。
\end{solve}

\begin{solve}[习题~\ref{ex:c33q2}：线性换元的面积缩放]\phantomsection\label{ans:c33q2}
偏导矩阵为 $J=\begin{pmatrix}2&1\\-1&1\end{pmatrix}$，故 $\det J=3>0$。求逆矩阵得
\[
u=\frac{1x-y}{3},\qquad v=\frac{x+2y}{3}.
\]
面积元满足 $dxdy=3\,dudv$，因此单位正方形映成的平行四边形面积为 $3$。对边向量 $(2,-1)$ 与 $(1,1)$ 的二维叉积绝对值也给出同样面积。
\end{solve}

\begin{solve}[习题~\ref{ex:c33q3}：线性换元的面积缩放]\phantomsection\label{ans:c33q3}
偏导矩阵为 $J=\begin{pmatrix}2&1\\-1&3\end{pmatrix}$，故 $\det J=7>0$。求逆矩阵得
\[
u=\frac{3x-y}{7},\qquad v=\frac{x+2y}{7}.
\]
面积元满足 $dxdy=7\,dudv$，因此单位正方形映成的平行四边形面积为 $7$。对边向量 $(2,-1)$ 与 $(1,3)$ 的二维叉积绝对值也给出同样面积。
\end{solve}

\begin{solve}[习题~\ref{ex:c33q4}：线性换元的面积缩放]\phantomsection\label{ans:c33q4}
偏导矩阵为 $J=\begin{pmatrix}3&1\\-1&2\end{pmatrix}$，故 $\det J=7>0$。求逆矩阵得
\[
u=\frac{2x-y}{7},\qquad v=\frac{x+3y}{7}.
\]
面积元满足 $dxdy=7\,dudv$，因此单位正方形映成的平行四边形面积为 $7$。对边向量 $(3,-1)$ 与 $(1,2)$ 的二维叉积绝对值也给出同样面积。
\end{solve}

\begin{solve}[习题~\ref{ex:c33q5}：线性换元的面积缩放]\phantomsection\label{ans:c33q5}
偏导矩阵为 $J=\begin{pmatrix}4&1\\-1&1\end{pmatrix}$，故 $\det J=5>0$。求逆矩阵得
\[
u=\frac{1x-y}{5},\qquad v=\frac{x+4y}{5}.
\]
面积元满足 $dxdy=5\,dudv$，因此单位正方形映成的平行四边形面积为 $5$。对边向量 $(4,-1)$ 与 $(1,1)$ 的二维叉积绝对值也给出同样面积。
\end{solve}

\begin{solve}[习题~\ref{ex:c33q6}：线性换元的面积缩放]\phantomsection\label{ans:c33q6}
偏导矩阵为 $J=\begin{pmatrix}3&1\\-1&4\end{pmatrix}$，故 $\det J=13>0$。求逆矩阵得
\[
u=\frac{4x-y}{13},\qquad v=\frac{x+3y}{13}.
\]
面积元满足 $dxdy=13\,dudv$，因此单位正方形映成的平行四边形面积为 $13$。对边向量 $(3,-1)$ 与 $(1,4)$ 的二维叉积绝对值也给出同样面积。
\end{solve}

\begin{solve}[习题~\ref{ex:c33q7}：极坐标下的链式法则]\phantomsection\label{ans:c33q7}
由链式法则，
\[
z_r=2x\cos\theta+2y\sin\theta=2r(\cos^2\theta+1\sin^2\theta),
\]

\[
z_\theta=2x(-r\sin\theta)+2y(r\cos\theta)=0r^2\sin\theta\cos\theta.
\]
直接代入得到 $z=r^2(\cos^2\theta+1\sin^2\theta)$，再求偏导与上述一致。本题系数为1，函数只依赖半径，所以角向导数恒为0。对 $r$ 求偏导与对 $\theta$ 求偏导代表不同的局部运动方式。
\end{solve}

\begin{solve}[习题~\ref{ex:c33q8}：极坐标下的链式法则]\phantomsection\label{ans:c33q8}
由链式法则，
\[
z_r=2x\cos\theta+4y\sin\theta=2r(\cos^2\theta+2\sin^2\theta),
\]

\[
z_\theta=2x(-r\sin\theta)+4y(r\cos\theta)=2r^2\sin\theta\cos\theta.
\]
直接代入得到 $z=r^2(\cos^2\theta+2\sin^2\theta)$，再求偏导与上述一致。两坐标方向系数不同，函数一般也随角度变化。对 $r$ 求偏导与对 $\theta$ 求偏导代表不同的局部运动方式。
\end{solve}

\begin{solve}[习题~\ref{ex:c33q9}：极坐标下的链式法则]\phantomsection\label{ans:c33q9}
由链式法则，
\[
z_r=2x\cos\theta+6y\sin\theta=2r(\cos^2\theta+3\sin^2\theta),
\]

\[
z_\theta=2x(-r\sin\theta)+6y(r\cos\theta)=4r^2\sin\theta\cos\theta.
\]
直接代入得到 $z=r^2(\cos^2\theta+3\sin^2\theta)$，再求偏导与上述一致。两坐标方向系数不同，函数一般也随角度变化。对 $r$ 求偏导与对 $\theta$ 求偏导代表不同的局部运动方式。
\end{solve}

\begin{solve}[习题~\ref{ex:c33q10}：极坐标下的链式法则]\phantomsection\label{ans:c33q10}
由链式法则，
\[
z_r=2x\cos\theta+8y\sin\theta=2r(\cos^2\theta+4\sin^2\theta),
\]

\[
z_\theta=2x(-r\sin\theta)+8y(r\cos\theta)=6r^2\sin\theta\cos\theta.
\]
直接代入得到 $z=r^2(\cos^2\theta+4\sin^2\theta)$，再求偏导与上述一致。两坐标方向系数不同，函数一般也随角度变化。对 $r$ 求偏导与对 $\theta$ 求偏导代表不同的局部运动方式。
\end{solve}

\begin{solve}[习题~\ref{ex:c33q11}：极坐标下的链式法则]\phantomsection\label{ans:c33q11}
由链式法则，
\[
z_r=2x\cos\theta+10y\sin\theta=2r(\cos^2\theta+5\sin^2\theta),
\]

\[
z_\theta=2x(-r\sin\theta)+10y(r\cos\theta)=8r^2\sin\theta\cos\theta.
\]
直接代入得到 $z=r^2(\cos^2\theta+5\sin^2\theta)$，再求偏导与上述一致。两坐标方向系数不同，函数一般也随角度变化。对 $r$ 求偏导与对 $\theta$ 求偏导代表不同的局部运动方式。
\end{solve}

\begin{solve}[习题~\ref{ex:c33q12}：极坐标下的链式法则]\phantomsection\label{ans:c33q12}
由链式法则，
\[
z_r=2x\cos\theta+12y\sin\theta=2r(\cos^2\theta+6\sin^2\theta),
\]

\[
z_\theta=2x(-r\sin\theta)+12y(r\cos\theta)=10r^2\sin\theta\cos\theta.
\]
直接代入得到 $z=r^2(\cos^2\theta+6\sin^2\theta)$，再求偏导与上述一致。两坐标方向系数不同，函数一般也随角度变化。对 $r$ 求偏导与对 $\theta$ 求偏导代表不同的局部运动方式。
\end{solve}

\begin{solve}[习题~\ref{ex:c33q13}：隐函数的一阶偏导]\phantomsection\label{ans:c33q13}
把方程整理为 $(1+x+y)z=1x+2y$。在原点附近分母非零，故
\[
z=\frac{1x+2y}{1+x+y}.
\]
商法则给出 $z_x=(1+(-1)y)/(1+x+y)^2$、$z_y=(2+(1)x)/(1+x+y)^2$，所以原点处偏导分别为 $1,2$。也可令 $F=(1+x+y)z-1x-2y$，使用 $z_x=-F_x/F_z$、$z_y=-F_y/F_z$；原点处 $F_z=1$，满足隐函数定理的非退化条件。
\end{solve}

\begin{solve}[习题~\ref{ex:c33q14}：隐函数的一阶偏导]\phantomsection\label{ans:c33q14}
把方程整理为 $(1+x+y)z=2x+1y$。在原点附近分母非零，故
\[
z=\frac{2x+1y}{1+x+y}.
\]
商法则给出 $z_x=(2+(1)y)/(1+x+y)^2$、$z_y=(1+(-1)x)/(1+x+y)^2$，所以原点处偏导分别为 $2,1$。也可令 $F=(1+x+y)z-2x-1y$，使用 $z_x=-F_x/F_z$、$z_y=-F_y/F_z$；原点处 $F_z=1$，满足隐函数定理的非退化条件。
\end{solve}

\begin{solve}[习题~\ref{ex:c33q15}：隐函数的一阶偏导]\phantomsection\label{ans:c33q15}
把方程整理为 $(1+x+y)z=2x+3y$。在原点附近分母非零，故
\[
z=\frac{2x+3y}{1+x+y}.
\]
商法则给出 $z_x=(2+(-1)y)/(1+x+y)^2$、$z_y=(3+(1)x)/(1+x+y)^2$，所以原点处偏导分别为 $2,3$。也可令 $F=(1+x+y)z-2x-3y$，使用 $z_x=-F_x/F_z$、$z_y=-F_y/F_z$；原点处 $F_z=1$，满足隐函数定理的非退化条件。
\end{solve}

\begin{solve}[习题~\ref{ex:c33q16}：隐函数的一阶偏导]\phantomsection\label{ans:c33q16}
把方程整理为 $(1+x+y)z=3x+2y$。在原点附近分母非零，故
\[
z=\frac{3x+2y}{1+x+y}.
\]
商法则给出 $z_x=(3+(1)y)/(1+x+y)^2$、$z_y=(2+(-1)x)/(1+x+y)^2$，所以原点处偏导分别为 $3,2$。也可令 $F=(1+x+y)z-3x-2y$，使用 $z_x=-F_x/F_z$、$z_y=-F_y/F_z$；原点处 $F_z=1$，满足隐函数定理的非退化条件。
\end{solve}

\begin{solve}[习题~\ref{ex:c33q17}：隐函数的一阶偏导]\phantomsection\label{ans:c33q17}
把方程整理为 $(1+x+y)z=4x+1y$。在原点附近分母非零，故
\[
z=\frac{4x+1y}{1+x+y}.
\]
商法则给出 $z_x=(4+(3)y)/(1+x+y)^2$、$z_y=(1+(-3)x)/(1+x+y)^2$，所以原点处偏导分别为 $4,1$。也可令 $F=(1+x+y)z-4x-1y$，使用 $z_x=-F_x/F_z$、$z_y=-F_y/F_z$；原点处 $F_z=1$，满足隐函数定理的非退化条件。
\end{solve}

\begin{solve}[习题~\ref{ex:c33q18}：隐函数的一阶偏导]\phantomsection\label{ans:c33q18}
把方程整理为 $(1+x+y)z=3x+4y$。在原点附近分母非零，故
\[
z=\frac{3x+4y}{1+x+y}.
\]
商法则给出 $z_x=(3+(-1)y)/(1+x+y)^2$、$z_y=(4+(1)x)/(1+x+y)^2$，所以原点处偏导分别为 $3,4$。也可令 $F=(1+x+y)z-3x-4y$，使用 $z_x=-F_x/F_z$、$z_y=-F_y/F_z$；原点处 $F_z=1$，满足隐函数定理的非退化条件。
\end{solve}

\begin{solve}[习题~\ref{ex:c33q19}：二元函数的二阶泰勒式]\phantomsection\label{ans:c33q19}
令 $u=1x+2y$，一元展开 $e^u=1+u+u^2/2+O(u^3)$ 给出
\[
T_2=1+1x+2y+\frac{1}{2}x^2+2xy+2y^2.
\]
因为 $|u|\le\sqrt{5}\sqrt{x^2+y^2}$，在原点附近余项为 $O((x^2+y^2)^{3/2})$。混合项系数为 $f_{xy}(0,0)$，而不是该数的一半：二次型 $\frac12h^THh$ 中有两个相等的交叉项，正好抵消外面的 $1/2$。
\end{solve}

\begin{solve}[习题~\ref{ex:c33q20}：二元函数的二阶泰勒式]\phantomsection\label{ans:c33q20}
令 $u=2x+1y$，一元展开 $e^u=1+u+u^2/2+O(u^3)$ 给出
\[
T_2=1+2x+1y+2x^2+2xy+\frac{1}{2}y^2.
\]
因为 $|u|\le\sqrt{5}\sqrt{x^2+y^2}$，在原点附近余项为 $O((x^2+y^2)^{3/2})$。混合项系数为 $f_{xy}(0,0)$，而不是该数的一半：二次型 $\frac12h^THh$ 中有两个相等的交叉项，正好抵消外面的 $1/2$。
\end{solve}

\begin{solve}[习题~\ref{ex:c33q21}：二元函数的二阶泰勒式]\phantomsection\label{ans:c33q21}
令 $u=2x+3y$，一元展开 $e^u=1+u+u^2/2+O(u^3)$ 给出
\[
T_2=1+2x+3y+2x^2+6xy+\frac{9}{2}y^2.
\]
因为 $|u|\le\sqrt{13}\sqrt{x^2+y^2}$，在原点附近余项为 $O((x^2+y^2)^{3/2})$。混合项系数为 $f_{xy}(0,0)$，而不是该数的一半：二次型 $\frac12h^THh$ 中有两个相等的交叉项，正好抵消外面的 $1/2$。
\end{solve}

\begin{solve}[习题~\ref{ex:c33q22}：二元函数的二阶泰勒式]\phantomsection\label{ans:c33q22}
令 $u=3x+2y$，一元展开 $e^u=1+u+u^2/2+O(u^3)$ 给出
\[
T_2=1+3x+2y+\frac{9}{2}x^2+6xy+2y^2.
\]
因为 $|u|\le\sqrt{13}\sqrt{x^2+y^2}$，在原点附近余项为 $O((x^2+y^2)^{3/2})$。混合项系数为 $f_{xy}(0,0)$，而不是该数的一半：二次型 $\frac12h^THh$ 中有两个相等的交叉项，正好抵消外面的 $1/2$。
\end{solve}

\begin{solve}[习题~\ref{ex:c33q23}：二元函数的二阶泰勒式]\phantomsection\label{ans:c33q23}
令 $u=4x+1y$，一元展开 $e^u=1+u+u^2/2+O(u^3)$ 给出
\[
T_2=1+4x+1y+8x^2+4xy+\frac{1}{2}y^2.
\]
因为 $|u|\le\sqrt{17}\sqrt{x^2+y^2}$，在原点附近余项为 $O((x^2+y^2)^{3/2})$。混合项系数为 $f_{xy}(0,0)$，而不是该数的一半：二次型 $\frac12h^THh$ 中有两个相等的交叉项，正好抵消外面的 $1/2$。
\end{solve}

\begin{solve}[习题~\ref{ex:c33q24}：二元函数的二阶泰勒式]\phantomsection\label{ans:c33q24}
令 $u=3x+4y$，一元展开 $e^u=1+u+u^2/2+O(u^3)$ 给出
\[
T_2=1+3x+4y+\frac{9}{2}x^2+12xy+8y^2.
\]
因为 $|u|\le\sqrt{25}\sqrt{x^2+y^2}$，在原点附近余项为 $O((x^2+y^2)^{3/2})$。混合项系数为 $f_{xy}(0,0)$，而不是该数的一半：二次型 $\frac12h^THh$ 中有两个相等的交叉项，正好抵消外面的 $1/2$。
\end{solve}

\clearpage
\section{链式法则、梯度与多元极值}\label{sec:ch34}
多元链式法则需要列出每一条依赖路径。方向导数以单位方向为标准；无约束极值看梯度与 Hessian，约束极值还须检查可行集与边界条件。

\textbf{资源对应：}CLP-3 \S~2.4、\S~2.5、\S~2.7、\S~2.9、\S~2.10。本章共 25 道编号习题。
\subsection{习题}
\begin{example}[方向向量必须单位化]\label{ex:c34q1}
求 $f(x,y,z)=e^{xyz}$ 在 $(0,1,1)$ 沿向量 $(0,1,1)$ 的方向导数。
\par\textit{选译并补充。来源：CLP-3 \S~2.7，M200 2008A。}\par\hyperref[ans:c34q1]{解答见第~\pageref{ans:c34q1} 页。}
\end{example}
\begin{example}[复合函数的两种求导路径]\label{ex:c34q2}
设 $z=x^2+1y^2$，$x=t$、$y=2t^2$。用链式法则与直接代入两种方法求 $dz/dt$。
\par\textit{改编训练。}\par\hyperref[ans:c34q2]{解答见第~\pageref{ans:c34q2} 页。}
\end{example}
\begin{example}[复合函数的两种求导路径]\label{ex:c34q3}
设 $z=x^2+2y^2$，$x=t$、$y=1t^2$。用链式法则与直接代入两种方法求 $dz/dt$。
\par\textit{改编训练。}\par\hyperref[ans:c34q3]{解答见第~\pageref{ans:c34q3} 页。}
\end{example}
\begin{example}[复合函数的两种求导路径]\label{ex:c34q4}
设 $z=x^2+3y^2$，$x=t$、$y=2t^2$。用链式法则与直接代入两种方法求 $dz/dt$。
\par\textit{改编训练。}\par\hyperref[ans:c34q4]{解答见第~\pageref{ans:c34q4} 页。}
\end{example}
\begin{example}[复合函数的两种求导路径]\label{ex:c34q5}
设 $z=x^2+2y^2$，$x=t$、$y=3t^2$。用链式法则与直接代入两种方法求 $dz/dt$。
\par\textit{改编训练。}\par\hyperref[ans:c34q5]{解答见第~\pageref{ans:c34q5} 页。}
\end{example}
\begin{example}[复合函数的两种求导路径]\label{ex:c34q6}
设 $z=x^2+4y^2$，$x=t$、$y=1t^2$。用链式法则与直接代入两种方法求 $dz/dt$。
\par\textit{改编训练。}\par\hyperref[ans:c34q6]{解答见第~\pageref{ans:c34q6} 页。}
\end{example}
\begin{example}[复合函数的两种求导路径]\label{ex:c34q7}
设 $z=x^2+3y^2$，$x=t$、$y=4t^2$。用链式法则与直接代入两种方法求 $dz/dt$。
\par\textit{改编训练。}\par\hyperref[ans:c34q7]{解答见第~\pageref{ans:c34q7} 页。}
\end{example}
\begin{example}[梯度与最陡方向]\label{ex:c34q8}
设 $f(x,y)=1x^2+2y^2$。求在 $(1,1)$ 沿向量 $(3,4)$ 的方向导数，并求最大方向导数及对应单位方向。
\par\textit{改编训练。}\par\hyperref[ans:c34q8]{解答见第~\pageref{ans:c34q8} 页。}
\end{example}
\begin{example}[梯度与最陡方向]\label{ex:c34q9}
设 $f(x,y)=2x^2+1y^2$。求在 $(1,1)$ 沿向量 $(3,4)$ 的方向导数，并求最大方向导数及对应单位方向。
\par\textit{改编训练。}\par\hyperref[ans:c34q9]{解答见第~\pageref{ans:c34q9} 页。}
\end{example}
\begin{example}[梯度与最陡方向]\label{ex:c34q10}
设 $f(x,y)=2x^2+3y^2$。求在 $(1,1)$ 沿向量 $(3,4)$ 的方向导数，并求最大方向导数及对应单位方向。
\par\textit{改编训练。}\par\hyperref[ans:c34q10]{解答见第~\pageref{ans:c34q10} 页。}
\end{example}
\begin{example}[梯度与最陡方向]\label{ex:c34q11}
设 $f(x,y)=3x^2+2y^2$。求在 $(1,1)$ 沿向量 $(3,4)$ 的方向导数，并求最大方向导数及对应单位方向。
\par\textit{改编训练。}\par\hyperref[ans:c34q11]{解答见第~\pageref{ans:c34q11} 页。}
\end{example}
\begin{example}[梯度与最陡方向]\label{ex:c34q12}
设 $f(x,y)=4x^2+1y^2$。求在 $(1,1)$ 沿向量 $(3,4)$ 的方向导数，并求最大方向导数及对应单位方向。
\par\textit{改编训练。}\par\hyperref[ans:c34q12]{解答见第~\pageref{ans:c34q12} 页。}
\end{example}
\begin{example}[梯度与最陡方向]\label{ex:c34q13}
设 $f(x,y)=1x^2+4y^2$。求在 $(1,1)$ 沿向量 $(3,4)$ 的方向导数，并求最大方向导数及对应单位方向。
\par\textit{改编训练。}\par\hyperref[ans:c34q13]{解答见第~\pageref{ans:c34q13} 页。}
\end{example}
\begin{example}[正定二次型的极小值]\label{ex:c34q14}
求 $f(x,y)=1(x-1)^2+1(y+2)^2+xy$ 的驻点，并判断是否为全局最小点。
\par\textit{改编训练。}\par\hyperref[ans:c34q14]{解答见第~\pageref{ans:c34q14} 页。}
\end{example}
\begin{example}[正定二次型的极小值]\label{ex:c34q15}
求 $f(x,y)=2(x-1)^2+1(y+2)^2+xy$ 的驻点，并判断是否为全局最小点。
\par\textit{改编训练。}\par\hyperref[ans:c34q15]{解答见第~\pageref{ans:c34q15} 页。}
\end{example}
\begin{example}[正定二次型的极小值]\label{ex:c34q16}
求 $f(x,y)=1(x-1)^2+3(y+2)^2+xy$ 的驻点，并判断是否为全局最小点。
\par\textit{改编训练。}\par\hyperref[ans:c34q16]{解答见第~\pageref{ans:c34q16} 页。}
\end{example}
\begin{example}[正定二次型的极小值]\label{ex:c34q17}
求 $f(x,y)=3(x-1)^2+2(y+2)^2+xy$ 的驻点，并判断是否为全局最小点。
\par\textit{改编训练。}\par\hyperref[ans:c34q17]{解答见第~\pageref{ans:c34q17} 页。}
\end{example}
\begin{example}[正定二次型的极小值]\label{ex:c34q18}
求 $f(x,y)=2(x-1)^2+4(y+2)^2+xy$ 的驻点，并判断是否为全局最小点。
\par\textit{改编训练。}\par\hyperref[ans:c34q18]{解答见第~\pageref{ans:c34q18} 页。}
\end{example}
\begin{example}[正定二次型的极小值]\label{ex:c34q19}
求 $f(x,y)=4(x-1)^2+3(y+2)^2+xy$ 的驻点，并判断是否为全局最小点。
\par\textit{改编训练。}\par\hyperref[ans:c34q19]{解答见第~\pageref{ans:c34q19} 页。}
\end{example}
\begin{example}[双曲线上的约束最值]\label{ex:c34q20}
在约束 $xy=1$ 下求 $f(x,y)=x^2+y^2$ 的最小值，并判断是否存在最大值。
\par\textit{改编训练。来源：CLP-3 \S~2.10，M200 2010A 的参数改编。}\par\hyperref[ans:c34q20]{解答见第~\pageref{ans:c34q20} 页。}
\end{example}
\begin{example}[双曲线上的约束最值]\label{ex:c34q21}
在约束 $xy=2$ 下求 $f(x,y)=x^2+y^2$ 的最小值，并判断是否存在最大值。
\par\textit{改编训练。来源：CLP-3 \S~2.10，M200 2010A 的参数改编。}\par\hyperref[ans:c34q21]{解答见第~\pageref{ans:c34q21} 页。}
\end{example}
\begin{example}[双曲线上的约束最值]\label{ex:c34q22}
在约束 $xy=3$ 下求 $f(x,y)=x^2+y^2$ 的最小值，并判断是否存在最大值。
\par\textit{改编训练。来源：CLP-3 \S~2.10，M200 2010A 的参数改编。}\par\hyperref[ans:c34q22]{解答见第~\pageref{ans:c34q22} 页。}
\end{example}
\begin{example}[双曲线上的约束最值]\label{ex:c34q23}
在约束 $xy=4$ 下求 $f(x,y)=x^2+y^2$ 的最小值，并判断是否存在最大值。
\par\textit{改编训练。来源：CLP-3 \S~2.10，M200 2010A 的参数改编。}\par\hyperref[ans:c34q23]{解答见第~\pageref{ans:c34q23} 页。}
\end{example}
\begin{example}[双曲线上的约束最值]\label{ex:c34q24}
在约束 $xy=5$ 下求 $f(x,y)=x^2+y^2$ 的最小值，并判断是否存在最大值。
\par\textit{改编训练。来源：CLP-3 \S~2.10，M200 2010A 的参数改编。}\par\hyperref[ans:c34q24]{解答见第~\pageref{ans:c34q24} 页。}
\end{example}
\begin{example}[双曲线上的约束最值]\label{ex:c34q25}
在约束 $xy=6$ 下求 $f(x,y)=x^2+y^2$ 的最小值，并判断是否存在最大值。
\par\textit{改编训练。来源：CLP-3 \S~2.10，M200 2010A 的参数改编。}\par\hyperref[ans:c34q25]{解答见第~\pageref{ans:c34q25} 页。}
\end{example}
\subsection{解答}
\begin{solve}[习题~\ref{ex:c34q1}：方向向量必须单位化]\phantomsection\label{ans:c34q1}
梯度为 $(yze^{xyz},xze^{xyz},xye^{xyz})$，在给定点等于 $(1,0,0)$。方向单位向量为 $(0,1,1)/\sqrt2$，所以方向导数为二者点积，等于0。本题即使漏除 $\sqrt2$ 也会偶然得到0，但这不说明单位化步骤可以普遍省略。
\end{solve}

\begin{solve}[习题~\ref{ex:c34q2}：复合函数的两种求导路径]\phantomsection\label{ans:c34q2}
链式法则给出
\[
\frac{dz}{dt}=2x\frac{dx}{dt}+2y\frac{dy}{dt}=2t+16t^3.
\]
直接代入得到 $z=t^2+4t^4$，再对 $t$ 求导也得到相同结果。在 $t=1$ 时变化率为 $18$。链式法则中两项对应通过 $x$ 与通过 $y$ 的两条变化途径，两项都应保留。直接代入是核验简单复合关系的有效方法。
\end{solve}

\begin{solve}[习题~\ref{ex:c34q3}：复合函数的两种求导路径]\phantomsection\label{ans:c34q3}
链式法则给出
\[
\frac{dz}{dt}=2x\frac{dx}{dt}+4y\frac{dy}{dt}=2t+8t^3.
\]
直接代入得到 $z=t^2+2t^4$，再对 $t$ 求导也得到相同结果。在 $t=1$ 时变化率为 $10$。链式法则中两项对应通过 $x$ 与通过 $y$ 的两条变化途径，两项都应保留。直接代入是核验简单复合关系的有效方法。
\end{solve}

\begin{solve}[习题~\ref{ex:c34q4}：复合函数的两种求导路径]\phantomsection\label{ans:c34q4}
链式法则给出
\[
\frac{dz}{dt}=2x\frac{dx}{dt}+6y\frac{dy}{dt}=2t+48t^3.
\]
直接代入得到 $z=t^2+12t^4$，再对 $t$ 求导也得到相同结果。在 $t=1$ 时变化率为 $50$。链式法则中两项对应通过 $x$ 与通过 $y$ 的两条变化途径，两项都应保留。直接代入是核验简单复合关系的有效方法。
\end{solve}

\begin{solve}[习题~\ref{ex:c34q5}：复合函数的两种求导路径]\phantomsection\label{ans:c34q5}
链式法则给出
\[
\frac{dz}{dt}=2x\frac{dx}{dt}+4y\frac{dy}{dt}=2t+72t^3.
\]
直接代入得到 $z=t^2+18t^4$，再对 $t$ 求导也得到相同结果。在 $t=1$ 时变化率为 $74$。链式法则中两项对应通过 $x$ 与通过 $y$ 的两条变化途径，两项都应保留。直接代入是核验简单复合关系的有效方法。
\end{solve}

\begin{solve}[习题~\ref{ex:c34q6}：复合函数的两种求导路径]\phantomsection\label{ans:c34q6}
链式法则给出
\[
\frac{dz}{dt}=2x\frac{dx}{dt}+8y\frac{dy}{dt}=2t+16t^3.
\]
直接代入得到 $z=t^2+4t^4$，再对 $t$ 求导也得到相同结果。在 $t=1$ 时变化率为 $18$。链式法则中两项对应通过 $x$ 与通过 $y$ 的两条变化途径，两项都应保留。直接代入是核验简单复合关系的有效方法。
\end{solve}

\begin{solve}[习题~\ref{ex:c34q7}：复合函数的两种求导路径]\phantomsection\label{ans:c34q7}
链式法则给出
\[
\frac{dz}{dt}=2x\frac{dx}{dt}+6y\frac{dy}{dt}=2t+192t^3.
\]
直接代入得到 $z=t^2+48t^4$，再对 $t$ 求导也得到相同结果。在 $t=1$ 时变化率为 $194$。链式法则中两项对应通过 $x$ 与通过 $y$ 的两条变化途径，两项都应保留。直接代入是核验简单复合关系的有效方法。
\end{solve}

\begin{solve}[习题~\ref{ex:c34q8}：梯度与最陡方向]\phantomsection\label{ans:c34q8}
在给定点梯度为 $g=(2,4)$，给定方向的单位向量为 $u=(3/5,4/5)$，所以
\[
D_uf(1,1)=g\cdot u=\frac{22}{5}.
\]
对任意单位向量 $v$，Cauchy--Schwarz 不等式给出 $g\cdot v\le|g|$，等号在 $v=g/|g|$ 时成立。故
\[
D_{\max}f=2\sqrt{5},\qquad v_{\max}=\frac{(1,2)}{\sqrt{5}}.
\]
梯度的模表示局部最大变化率，梯度的方向表示达到这个最大值的方向，二者是不同的信息。
\end{solve}

\begin{solve}[习题~\ref{ex:c34q9}：梯度与最陡方向]\phantomsection\label{ans:c34q9}
在给定点梯度为 $g=(4,2)$，给定方向的单位向量为 $u=(3/5,4/5)$，所以
\[
D_uf(1,1)=g\cdot u=4.
\]
对任意单位向量 $v$，Cauchy--Schwarz 不等式给出 $g\cdot v\le|g|$，等号在 $v=g/|g|$ 时成立。故
\[
D_{\max}f=2\sqrt{5},\qquad v_{\max}=\frac{(2,1)}{\sqrt{5}}.
\]
梯度的模表示局部最大变化率，梯度的方向表示达到这个最大值的方向，二者是不同的信息。
\end{solve}

\begin{solve}[习题~\ref{ex:c34q10}：梯度与最陡方向]\phantomsection\label{ans:c34q10}
在给定点梯度为 $g=(4,6)$，给定方向的单位向量为 $u=(3/5,4/5)$，所以
\[
D_uf(1,1)=g\cdot u=\frac{36}{5}.
\]
对任意单位向量 $v$，Cauchy--Schwarz 不等式给出 $g\cdot v\le|g|$，等号在 $v=g/|g|$ 时成立。故
\[
D_{\max}f=2\sqrt{13},\qquad v_{\max}=\frac{(2,3)}{\sqrt{13}}.
\]
梯度的模表示局部最大变化率，梯度的方向表示达到这个最大值的方向，二者是不同的信息。
\end{solve}

\begin{solve}[习题~\ref{ex:c34q11}：梯度与最陡方向]\phantomsection\label{ans:c34q11}
在给定点梯度为 $g=(6,4)$，给定方向的单位向量为 $u=(3/5,4/5)$，所以
\[
D_uf(1,1)=g\cdot u=\frac{34}{5}.
\]
对任意单位向量 $v$，Cauchy--Schwarz 不等式给出 $g\cdot v\le|g|$，等号在 $v=g/|g|$ 时成立。故
\[
D_{\max}f=2\sqrt{13},\qquad v_{\max}=\frac{(3,2)}{\sqrt{13}}.
\]
梯度的模表示局部最大变化率，梯度的方向表示达到这个最大值的方向，二者是不同的信息。
\end{solve}

\begin{solve}[习题~\ref{ex:c34q12}：梯度与最陡方向]\phantomsection\label{ans:c34q12}
在给定点梯度为 $g=(8,2)$，给定方向的单位向量为 $u=(3/5,4/5)$，所以
\[
D_uf(1,1)=g\cdot u=\frac{32}{5}.
\]
对任意单位向量 $v$，Cauchy--Schwarz 不等式给出 $g\cdot v\le|g|$，等号在 $v=g/|g|$ 时成立。故
\[
D_{\max}f=2\sqrt{17},\qquad v_{\max}=\frac{(4,1)}{\sqrt{17}}.
\]
梯度的模表示局部最大变化率，梯度的方向表示达到这个最大值的方向，二者是不同的信息。
\end{solve}

\begin{solve}[习题~\ref{ex:c34q13}：梯度与最陡方向]\phantomsection\label{ans:c34q13}
在给定点梯度为 $g=(2,8)$，给定方向的单位向量为 $u=(3/5,4/5)$，所以
\[
D_uf(1,1)=g\cdot u=\frac{38}{5}.
\]
对任意单位向量 $v$，Cauchy--Schwarz 不等式给出 $g\cdot v\le|g|$，等号在 $v=g/|g|$ 时成立。故
\[
D_{\max}f=2\sqrt{17},\qquad v_{\max}=\frac{(1,4)}{\sqrt{17}}.
\]
梯度的模表示局部最大变化率，梯度的方向表示达到这个最大值的方向，二者是不同的信息。
\end{solve}

\begin{solve}[习题~\ref{ex:c34q14}：正定二次型的极小值]\phantomsection\label{ans:c34q14}
驻点满足 $2x+y=2$、$x+2y=-4$。系数行列式为 $3>0$，因此解唯一：
\[
x_*=\frac{8}{3},\qquad y_*=-\frac{10}{3}.
\]
Hessian 为 $H=\begin{pmatrix}2&1\\1&2\end{pmatrix}$，其首个主子式 $2>0$、行列式 $3>0$，所以正定。对任意增量 $h$，二次函数恰有 $f(X_*+h)-f(X_*)=h^THh/2>0$（$h\ne0$），故该驻点是唯一全局最小点。正定性不只给出局部结论，还因为本题是整体二次函数而保证全局最优。
\end{solve}

\begin{solve}[习题~\ref{ex:c34q15}：正定二次型的极小值]\phantomsection\label{ans:c34q15}
驻点满足 $4x+y=4$、$x+2y=-4$。系数行列式为 $7>0$，因此解唯一：
\[
x_*=\frac{12}{7},\qquad y_*=-\frac{20}{7}.
\]
Hessian 为 $H=\begin{pmatrix}4&1\\1&2\end{pmatrix}$，其首个主子式 $4>0$、行列式 $7>0$，所以正定。对任意增量 $h$，二次函数恰有 $f(X_*+h)-f(X_*)=h^THh/2>0$（$h\ne0$），故该驻点是唯一全局最小点。正定性不只给出局部结论，还因为本题是整体二次函数而保证全局最优。
\end{solve}

\begin{solve}[习题~\ref{ex:c34q16}：正定二次型的极小值]\phantomsection\label{ans:c34q16}
驻点满足 $2x+y=2$、$x+6y=-12$。系数行列式为 $11>0$，因此解唯一：
\[
x_*=\frac{24}{11},\qquad y_*=-\frac{26}{11}.
\]
Hessian 为 $H=\begin{pmatrix}2&1\\1&6\end{pmatrix}$，其首个主子式 $2>0$、行列式 $11>0$，所以正定。对任意增量 $h$，二次函数恰有 $f(X_*+h)-f(X_*)=h^THh/2>0$（$h\ne0$），故该驻点是唯一全局最小点。正定性不只给出局部结论，还因为本题是整体二次函数而保证全局最优。
\end{solve}

\begin{solve}[习题~\ref{ex:c34q17}：正定二次型的极小值]\phantomsection\label{ans:c34q17}
驻点满足 $6x+y=6$、$x+4y=-8$。系数行列式为 $23>0$，因此解唯一：
\[
x_*=\frac{32}{23},\qquad y_*=-\frac{54}{23}.
\]
Hessian 为 $H=\begin{pmatrix}6&1\\1&4\end{pmatrix}$，其首个主子式 $6>0$、行列式 $23>0$，所以正定。对任意增量 $h$，二次函数恰有 $f(X_*+h)-f(X_*)=h^THh/2>0$（$h\ne0$），故该驻点是唯一全局最小点。正定性不只给出局部结论，还因为本题是整体二次函数而保证全局最优。
\end{solve}

\begin{solve}[习题~\ref{ex:c34q18}：正定二次型的极小值]\phantomsection\label{ans:c34q18}
驻点满足 $4x+y=4$、$x+8y=-16$。系数行列式为 $31>0$，因此解唯一：
\[
x_*=\frac{48}{31},\qquad y_*=-\frac{68}{31}.
\]
Hessian 为 $H=\begin{pmatrix}4&1\\1&8\end{pmatrix}$，其首个主子式 $4>0$、行列式 $31>0$，所以正定。对任意增量 $h$，二次函数恰有 $f(X_*+h)-f(X_*)=h^THh/2>0$（$h\ne0$），故该驻点是唯一全局最小点。正定性不只给出局部结论，还因为本题是整体二次函数而保证全局最优。
\end{solve}

\begin{solve}[习题~\ref{ex:c34q19}：正定二次型的极小值]\phantomsection\label{ans:c34q19}
驻点满足 $8x+y=8$、$x+6y=-12$。系数行列式为 $47>0$，因此解唯一：
\[
x_*=\frac{60}{47},\qquad y_*=-\frac{104}{47}.
\]
Hessian 为 $H=\begin{pmatrix}8&1\\1&6\end{pmatrix}$，其首个主子式 $8>0$、行列式 $47>0$，所以正定。对任意增量 $h$，二次函数恰有 $f(X_*+h)-f(X_*)=h^THh/2>0$（$h\ne0$），故该驻点是唯一全局最小点。正定性不只给出局部结论，还因为本题是整体二次函数而保证全局最优。
\end{solve}

\begin{solve}[习题~\ref{ex:c34q20}：双曲线上的约束最值]\phantomsection\label{ans:c34q20}
由平方非负有 $x^2+y^2\ge2|xy|=2$。等号要求 $|x|=|y|$；结合乘积为正，得到两个等号点 $x=y=\sqrt{1}$ 和 $x=y=-\sqrt{1}$。因此最小值为 $2$。令 $x\to\infty$、$y=1/x$，则
\[
f=x^2+\frac{1}{x^2}\to\infty,
\]
故不存在最大值。也可使用 Lagrange 乘子法找候选点，但找到驻点后仍须说明其全局性质和可行集的无界方向。
\end{solve}

\begin{solve}[习题~\ref{ex:c34q21}：双曲线上的约束最值]\phantomsection\label{ans:c34q21}
由平方非负有 $x^2+y^2\ge2|xy|=4$。等号要求 $|x|=|y|$；结合乘积为正，得到两个等号点 $x=y=\sqrt{2}$ 和 $x=y=-\sqrt{2}$。因此最小值为 $4$。令 $x\to\infty$、$y=2/x$，则
\[
f=x^2+\frac{4}{x^2}\to\infty,
\]
故不存在最大值。也可使用 Lagrange 乘子法找候选点，但找到驻点后仍须说明其全局性质和可行集的无界方向。
\end{solve}

\begin{solve}[习题~\ref{ex:c34q22}：双曲线上的约束最值]\phantomsection\label{ans:c34q22}
由平方非负有 $x^2+y^2\ge2|xy|=6$。等号要求 $|x|=|y|$；结合乘积为正，得到两个等号点 $x=y=\sqrt{3}$ 和 $x=y=-\sqrt{3}$。因此最小值为 $6$。令 $x\to\infty$、$y=3/x$，则
\[
f=x^2+\frac{9}{x^2}\to\infty,
\]
故不存在最大值。也可使用 Lagrange 乘子法找候选点，但找到驻点后仍须说明其全局性质和可行集的无界方向。
\end{solve}

\begin{solve}[习题~\ref{ex:c34q23}：双曲线上的约束最值]\phantomsection\label{ans:c34q23}
由平方非负有 $x^2+y^2\ge2|xy|=8$。等号要求 $|x|=|y|$；结合乘积为正，得到两个等号点 $x=y=\sqrt{4}$ 和 $x=y=-\sqrt{4}$。因此最小值为 $8$。令 $x\to\infty$、$y=4/x$，则
\[
f=x^2+\frac{16}{x^2}\to\infty,
\]
故不存在最大值。也可使用 Lagrange 乘子法找候选点，但找到驻点后仍须说明其全局性质和可行集的无界方向。
\end{solve}

\begin{solve}[习题~\ref{ex:c34q24}：双曲线上的约束最值]\phantomsection\label{ans:c34q24}
由平方非负有 $x^2+y^2\ge2|xy|=10$。等号要求 $|x|=|y|$；结合乘积为正，得到两个等号点 $x=y=\sqrt{5}$ 和 $x=y=-\sqrt{5}$。因此最小值为 $10$。令 $x\to\infty$、$y=5/x$，则
\[
f=x^2+\frac{25}{x^2}\to\infty,
\]
故不存在最大值。也可使用 Lagrange 乘子法找候选点，但找到驻点后仍须说明其全局性质和可行集的无界方向。
\end{solve}

\begin{solve}[习题~\ref{ex:c34q25}：双曲线上的约束最值]\phantomsection\label{ans:c34q25}
由平方非负有 $x^2+y^2\ge2|xy|=12$。等号要求 $|x|=|y|$；结合乘积为正，得到两个等号点 $x=y=\sqrt{6}$ 和 $x=y=-\sqrt{6}$。因此最小值为 $12$。令 $x\to\infty$、$y=6/x$，则
\[
f=x^2+\frac{36}{x^2}\to\infty,
\]
故不存在最大值。也可使用 Lagrange 乘子法找候选点，但找到驻点后仍须说明其全局性质和可行集的无界方向。
\end{solve}

\clearpage
\section{约束最值与边界检查}\label{sec:ch35}
紧集上的连续函数一定取得最值，但找出最值还需要检查内部、边界和角点。Lagrange 乘子法是寻找正规约束下候选点的方法，不能代替存在性与全局比较。

\textbf{资源对应：}CLP-3 \S~2.9、\S~2.10。本章共 24 道编号习题。
\subsection{习题}
\begin{example}[椭圆上的线性函数最值]\label{ex:c35q1}
在 $x^2/1+y^2/4=1$ 上求 $x+y$ 的最大、最小值及取值点。
\par\textit{改编训练。}\par\hyperref[ans:c35q1]{解答见第~\pageref{ans:c35q1} 页。}
\end{example}
\begin{example}[椭圆上的线性函数最值]\label{ex:c35q2}
在 $x^2/4+y^2/1=1$ 上求 $x+y$ 的最大、最小值及取值点。
\par\textit{改编训练。}\par\hyperref[ans:c35q2]{解答见第~\pageref{ans:c35q2} 页。}
\end{example}
\begin{example}[椭圆上的线性函数最值]\label{ex:c35q3}
在 $x^2/4+y^2/9=1$ 上求 $x+y$ 的最大、最小值及取值点。
\par\textit{改编训练。}\par\hyperref[ans:c35q3]{解答见第~\pageref{ans:c35q3} 页。}
\end{example}
\begin{example}[椭圆上的线性函数最值]\label{ex:c35q4}
在 $x^2/9+y^2/4=1$ 上求 $x+y$ 的最大、最小值及取值点。
\par\textit{改编训练。}\par\hyperref[ans:c35q4]{解答见第~\pageref{ans:c35q4} 页。}
\end{example}
\begin{example}[椭圆上的线性函数最值]\label{ex:c35q5}
在 $x^2/16+y^2/1=1$ 上求 $x+y$ 的最大、最小值及取值点。
\par\textit{改编训练。}\par\hyperref[ans:c35q5]{解答见第~\pageref{ans:c35q5} 页。}
\end{example}
\begin{example}[椭圆上的线性函数最值]\label{ex:c35q6}
在 $x^2/9+y^2/16=1$ 上求 $x+y$ 的最大、最小值及取值点。
\par\textit{改编训练。}\par\hyperref[ans:c35q6]{解答见第~\pageref{ans:c35q6} 页。}
\end{example}
\begin{example}[正数定和时的三元乘积]\label{ex:c35q7}
若 $x,y,z\ge0$ 且 $x+y+z=3$，求 $xyz$ 的最大、最小值。
\par\textit{改编训练。}\par\hyperref[ans:c35q7]{解答见第~\pageref{ans:c35q7} 页。}
\end{example}
\begin{example}[正数定和时的三元乘积]\label{ex:c35q8}
若 $x,y,z\ge0$ 且 $x+y+z=6$，求 $xyz$ 的最大、最小值。
\par\textit{改编训练。}\par\hyperref[ans:c35q8]{解答见第~\pageref{ans:c35q8} 页。}
\end{example}
\begin{example}[正数定和时的三元乘积]\label{ex:c35q9}
若 $x,y,z\ge0$ 且 $x+y+z=9$，求 $xyz$ 的最大、最小值。
\par\textit{改编训练。}\par\hyperref[ans:c35q9]{解答见第~\pageref{ans:c35q9} 页。}
\end{example}
\begin{example}[正数定和时的三元乘积]\label{ex:c35q10}
若 $x,y,z\ge0$ 且 $x+y+z=12$，求 $xyz$ 的最大、最小值。
\par\textit{改编训练。}\par\hyperref[ans:c35q10]{解答见第~\pageref{ans:c35q10} 页。}
\end{example}
\begin{example}[正数定和时的三元乘积]\label{ex:c35q11}
若 $x,y,z\ge0$ 且 $x+y+z=15$，求 $xyz$ 的最大、最小值。
\par\textit{改编训练。}\par\hyperref[ans:c35q11]{解答见第~\pageref{ans:c35q11} 页。}
\end{example}
\begin{example}[正数定和时的三元乘积]\label{ex:c35q12}
若 $x,y,z\ge0$ 且 $x+y+z=18$，求 $xyz$ 的最大、最小值。
\par\textit{改编训练。}\par\hyperref[ans:c35q12]{解答见第~\pageref{ans:c35q12} 页。}
\end{example}
\begin{example}[圆盘上的二次型]\label{ex:c35q13}
在闭圆盘 $x^2+y^2\le1$ 上求 $f(x,y)=x^2+2y^2$ 的最大、最小值及取值位置。
\par\textit{改编训练。}\par\hyperref[ans:c35q13]{解答见第~\pageref{ans:c35q13} 页。}
\end{example}
\begin{example}[圆盘上的二次型]\label{ex:c35q14}
在闭圆盘 $x^2+y^2\le1$ 上求 $f(x,y)=x^2+3y^2$ 的最大、最小值及取值位置。
\par\textit{改编训练。}\par\hyperref[ans:c35q14]{解答见第~\pageref{ans:c35q14} 页。}
\end{example}
\begin{example}[圆盘上的二次型]\label{ex:c35q15}
在闭圆盘 $x^2+y^2\le1$ 上求 $f(x,y)=x^2+4y^2$ 的最大、最小值及取值位置。
\par\textit{改编训练。}\par\hyperref[ans:c35q15]{解答见第~\pageref{ans:c35q15} 页。}
\end{example}
\begin{example}[圆盘上的二次型]\label{ex:c35q16}
在闭圆盘 $x^2+y^2\le1$ 上求 $f(x,y)=x^2+5y^2$ 的最大、最小值及取值位置。
\par\textit{改编训练。}\par\hyperref[ans:c35q16]{解答见第~\pageref{ans:c35q16} 页。}
\end{example}
\begin{example}[圆盘上的二次型]\label{ex:c35q17}
在闭圆盘 $x^2+y^2\le1$ 上求 $f(x,y)=x^2+6y^2$ 的最大、最小值及取值位置。
\par\textit{改编训练。}\par\hyperref[ans:c35q17]{解答见第~\pageref{ans:c35q17} 页。}
\end{example}
\begin{example}[圆盘上的二次型]\label{ex:c35q18}
在闭圆盘 $x^2+y^2\le1$ 上求 $f(x,y)=x^2+7y^2$ 的最大、最小值及取值位置。
\par\textit{改编训练。}\par\hyperref[ans:c35q18]{解答见第~\pageref{ans:c35q18} 页。}
\end{example}
\begin{example}[Hessian 退化时继续判别]\label{ex:c35q19}
判断 $f(x,y)=x^4+1y^4$ 与 $g(x,y)=x^4-1y^4$ 在原点的极值性质。
\par\textit{改编训练。}\par\hyperref[ans:c35q19]{解答见第~\pageref{ans:c35q19} 页。}
\end{example}
\begin{example}[Hessian 退化时继续判别]\label{ex:c35q20}
判断 $f(x,y)=x^4+2y^4$ 与 $g(x,y)=x^4-2y^4$ 在原点的极值性质。
\par\textit{改编训练。}\par\hyperref[ans:c35q20]{解答见第~\pageref{ans:c35q20} 页。}
\end{example}
\begin{example}[Hessian 退化时继续判别]\label{ex:c35q21}
判断 $f(x,y)=x^4+3y^4$ 与 $g(x,y)=x^4-3y^4$ 在原点的极值性质。
\par\textit{改编训练。}\par\hyperref[ans:c35q21]{解答见第~\pageref{ans:c35q21} 页。}
\end{example}
\begin{example}[Hessian 退化时继续判别]\label{ex:c35q22}
判断 $f(x,y)=x^4+4y^4$ 与 $g(x,y)=x^4-4y^4$ 在原点的极值性质。
\par\textit{改编训练。}\par\hyperref[ans:c35q22]{解答见第~\pageref{ans:c35q22} 页。}
\end{example}
\begin{example}[Hessian 退化时继续判别]\label{ex:c35q23}
判断 $f(x,y)=x^4+5y^4$ 与 $g(x,y)=x^4-5y^4$ 在原点的极值性质。
\par\textit{改编训练。}\par\hyperref[ans:c35q23]{解答见第~\pageref{ans:c35q23} 页。}
\end{example}
\begin{example}[Hessian 退化时继续判别]\label{ex:c35q24}
判断 $f(x,y)=x^4+6y^4$ 与 $g(x,y)=x^4-6y^4$ 在原点的极值性质。
\par\textit{改编训练。}\par\hyperref[ans:c35q24]{解答见第~\pageref{ans:c35q24} 页。}
\end{example}
\subsection{解答}
\begin{solve}[习题~\ref{ex:c35q1}：椭圆上的线性函数最值]\phantomsection\label{ans:c35q1}
令 $u=x/1$、$v=y/2$，则 $u^2+v^2=1$。由 Cauchy--Schwarz 不等式，$|1u+2v|\le\sqrt{5}$。因此
\[
M=\sqrt{5},\qquad m=-\sqrt{5}.
\]
最大值在 $(x,y)=(1,4)/\sqrt{5}$ 处取得，最小值在其相反点取得。椭圆紧且函数连续，存在性有保证；等号条件明确给出了达到界的点。
\end{solve}

\begin{solve}[习题~\ref{ex:c35q2}：椭圆上的线性函数最值]\phantomsection\label{ans:c35q2}
令 $u=x/2$、$v=y/1$，则 $u^2+v^2=1$。由 Cauchy--Schwarz 不等式，$|2u+1v|\le\sqrt{5}$。因此
\[
M=\sqrt{5},\qquad m=-\sqrt{5}.
\]
最大值在 $(x,y)=(4,1)/\sqrt{5}$ 处取得，最小值在其相反点取得。椭圆紧且函数连续，存在性有保证；等号条件明确给出了达到界的点。
\end{solve}

\begin{solve}[习题~\ref{ex:c35q3}：椭圆上的线性函数最值]\phantomsection\label{ans:c35q3}
令 $u=x/2$、$v=y/3$，则 $u^2+v^2=1$。由 Cauchy--Schwarz 不等式，$|2u+3v|\le\sqrt{13}$。因此
\[
M=\sqrt{13},\qquad m=-\sqrt{13}.
\]
最大值在 $(x,y)=(4,9)/\sqrt{13}$ 处取得，最小值在其相反点取得。椭圆紧且函数连续，存在性有保证；等号条件明确给出了达到界的点。
\end{solve}

\begin{solve}[习题~\ref{ex:c35q4}：椭圆上的线性函数最值]\phantomsection\label{ans:c35q4}
令 $u=x/3$、$v=y/2$，则 $u^2+v^2=1$。由 Cauchy--Schwarz 不等式，$|3u+2v|\le\sqrt{13}$。因此
\[
M=\sqrt{13},\qquad m=-\sqrt{13}.
\]
最大值在 $(x,y)=(9,4)/\sqrt{13}$ 处取得，最小值在其相反点取得。椭圆紧且函数连续，存在性有保证；等号条件明确给出了达到界的点。
\end{solve}

\begin{solve}[习题~\ref{ex:c35q5}：椭圆上的线性函数最值]\phantomsection\label{ans:c35q5}
令 $u=x/4$、$v=y/1$，则 $u^2+v^2=1$。由 Cauchy--Schwarz 不等式，$|4u+1v|\le\sqrt{17}$。因此
\[
M=\sqrt{17},\qquad m=-\sqrt{17}.
\]
最大值在 $(x,y)=(16,1)/\sqrt{17}$ 处取得，最小值在其相反点取得。椭圆紧且函数连续，存在性有保证；等号条件明确给出了达到界的点。
\end{solve}

\begin{solve}[习题~\ref{ex:c35q6}：椭圆上的线性函数最值]\phantomsection\label{ans:c35q6}
令 $u=x/3$、$v=y/4$，则 $u^2+v^2=1$。由 Cauchy--Schwarz 不等式，$|3u+4v|\le\sqrt{25}$。因此
\[
M=\sqrt{25},\qquad m=-\sqrt{25}.
\]
最大值在 $(x,y)=(9,16)/\sqrt{25}$ 处取得，最小值在其相反点取得。椭圆紧且函数连续，存在性有保证；等号条件明确给出了达到界的点。
\end{solve}

\begin{solve}[习题~\ref{ex:c35q7}：正数定和时的三元乘积]\phantomsection\label{ans:c35q7}
可行集为紧三角形，乘积连续。边界上至少一坐标为零，所以最小值为0。内部使用算术—几何平均不等式得到
\[
xyz\le\left(\frac{x+y+z}3\right)^3=1,
\]
等号当且仅当 $x=y=z=1$。因此最大值为 $1$。也可用乘子法得到内部三个坐标相等，但必须补查边界；仅给出内部候选点无法回答最小值。
\end{solve}

\begin{solve}[习题~\ref{ex:c35q8}：正数定和时的三元乘积]\phantomsection\label{ans:c35q8}
可行集为紧三角形，乘积连续。边界上至少一坐标为零，所以最小值为0。内部使用算术—几何平均不等式得到
\[
xyz\le\left(\frac{x+y+z}3\right)^3=8,
\]
等号当且仅当 $x=y=z=2$。因此最大值为 $8$。也可用乘子法得到内部三个坐标相等，但必须补查边界；仅给出内部候选点无法回答最小值。
\end{solve}

\begin{solve}[习题~\ref{ex:c35q9}：正数定和时的三元乘积]\phantomsection\label{ans:c35q9}
可行集为紧三角形，乘积连续。边界上至少一坐标为零，所以最小值为0。内部使用算术—几何平均不等式得到
\[
xyz\le\left(\frac{x+y+z}3\right)^3=27,
\]
等号当且仅当 $x=y=z=3$。因此最大值为 $27$。也可用乘子法得到内部三个坐标相等，但必须补查边界；仅给出内部候选点无法回答最小值。
\end{solve}

\begin{solve}[习题~\ref{ex:c35q10}：正数定和时的三元乘积]\phantomsection\label{ans:c35q10}
可行集为紧三角形，乘积连续。边界上至少一坐标为零，所以最小值为0。内部使用算术—几何平均不等式得到
\[
xyz\le\left(\frac{x+y+z}3\right)^3=64,
\]
等号当且仅当 $x=y=z=4$。因此最大值为 $64$。也可用乘子法得到内部三个坐标相等，但必须补查边界；仅给出内部候选点无法回答最小值。
\end{solve}

\begin{solve}[习题~\ref{ex:c35q11}：正数定和时的三元乘积]\phantomsection\label{ans:c35q11}
可行集为紧三角形，乘积连续。边界上至少一坐标为零，所以最小值为0。内部使用算术—几何平均不等式得到
\[
xyz\le\left(\frac{x+y+z}3\right)^3=125,
\]
等号当且仅当 $x=y=z=5$。因此最大值为 $125$。也可用乘子法得到内部三个坐标相等，但必须补查边界；仅给出内部候选点无法回答最小值。
\end{solve}

\begin{solve}[习题~\ref{ex:c35q12}：正数定和时的三元乘积]\phantomsection\label{ans:c35q12}
可行集为紧三角形，乘积连续。边界上至少一坐标为零，所以最小值为0。内部使用算术—几何平均不等式得到
\[
xyz\le\left(\frac{x+y+z}3\right)^3=216,
\]
等号当且仅当 $x=y=z=6$。因此最大值为 $216$。也可用乘子法得到内部三个坐标相等，但必须补查边界；仅给出内部候选点无法回答最小值。
\end{solve}

\begin{solve}[习题~\ref{ex:c35q13}：圆盘上的二次型]\phantomsection\label{ans:c35q13}
函数非负，且在原点取0，所以最小值为0。另一方面，
\[
f(x,y)\le 2(x^2+y^2)\le 2.
\]
要同时取等号，需要 $x=0$ 且 $x^2+y^2=1$，故最大值为 $2$，在 $(0,1)$、$(0,-1)$ 两点取得。内部只有原点一个驻点，它仅提供最小值；最大值来自边界。这是“只解梯度等于零”遗漏全局最大值的简单例子。
\end{solve}

\begin{solve}[习题~\ref{ex:c35q14}：圆盘上的二次型]\phantomsection\label{ans:c35q14}
函数非负，且在原点取0，所以最小值为0。另一方面，
\[
f(x,y)\le 3(x^2+y^2)\le 3.
\]
要同时取等号，需要 $x=0$ 且 $x^2+y^2=1$，故最大值为 $3$，在 $(0,1)$、$(0,-1)$ 两点取得。内部只有原点一个驻点，它仅提供最小值；最大值来自边界。这是“只解梯度等于零”遗漏全局最大值的简单例子。
\end{solve}

\begin{solve}[习题~\ref{ex:c35q15}：圆盘上的二次型]\phantomsection\label{ans:c35q15}
函数非负，且在原点取0，所以最小值为0。另一方面，
\[
f(x,y)\le 4(x^2+y^2)\le 4.
\]
要同时取等号，需要 $x=0$ 且 $x^2+y^2=1$，故最大值为 $4$，在 $(0,1)$、$(0,-1)$ 两点取得。内部只有原点一个驻点，它仅提供最小值；最大值来自边界。这是“只解梯度等于零”遗漏全局最大值的简单例子。
\end{solve}

\begin{solve}[习题~\ref{ex:c35q16}：圆盘上的二次型]\phantomsection\label{ans:c35q16}
函数非负，且在原点取0，所以最小值为0。另一方面，
\[
f(x,y)\le 5(x^2+y^2)\le 5.
\]
要同时取等号，需要 $x=0$ 且 $x^2+y^2=1$，故最大值为 $5$，在 $(0,1)$、$(0,-1)$ 两点取得。内部只有原点一个驻点，它仅提供最小值；最大值来自边界。这是“只解梯度等于零”遗漏全局最大值的简单例子。
\end{solve}

\begin{solve}[习题~\ref{ex:c35q17}：圆盘上的二次型]\phantomsection\label{ans:c35q17}
函数非负，且在原点取0，所以最小值为0。另一方面，
\[
f(x,y)\le 6(x^2+y^2)\le 6.
\]
要同时取等号，需要 $x=0$ 且 $x^2+y^2=1$，故最大值为 $6$，在 $(0,1)$、$(0,-1)$ 两点取得。内部只有原点一个驻点，它仅提供最小值；最大值来自边界。这是“只解梯度等于零”遗漏全局最大值的简单例子。
\end{solve}

\begin{solve}[习题~\ref{ex:c35q18}：圆盘上的二次型]\phantomsection\label{ans:c35q18}
函数非负，且在原点取0，所以最小值为0。另一方面，
\[
f(x,y)\le 7(x^2+y^2)\le 7.
\]
要同时取等号，需要 $x=0$ 且 $x^2+y^2=1$，故最大值为 $7$，在 $(0,1)$、$(0,-1)$ 两点取得。内部只有原点一个驻点，它仅提供最小值；最大值来自边界。这是“只解梯度等于零”遗漏全局最大值的简单例子。
\end{solve}

\begin{solve}[习题~\ref{ex:c35q19}：Hessian 退化时继续判别]\phantomsection\label{ans:c35q19}
两函数在原点梯度均为零，Hessian 也都为零矩阵，所以二阶充分判别法无法给出结论。对第一个函数，
\[
f(x,y)=x^4+1y^4>0\quad((x,y)\ne(0,0)),
\]
故原点是严格全局最小点。第二个函数沿 $y=0$ 为正，沿 $x=0$ 为负，所以原点不是极值点，而是鞍点。相同的退化 Hessian 可以对应不同性质，说明“二阶判别失效”不等于“没有极值”。
\end{solve}

\begin{solve}[习题~\ref{ex:c35q20}：Hessian 退化时继续判别]\phantomsection\label{ans:c35q20}
两函数在原点梯度均为零，Hessian 也都为零矩阵，所以二阶充分判别法无法给出结论。对第一个函数，
\[
f(x,y)=x^4+2y^4>0\quad((x,y)\ne(0,0)),
\]
故原点是严格全局最小点。第二个函数沿 $y=0$ 为正，沿 $x=0$ 为负，所以原点不是极值点，而是鞍点。相同的退化 Hessian 可以对应不同性质，说明“二阶判别失效”不等于“没有极值”。
\end{solve}

\begin{solve}[习题~\ref{ex:c35q21}：Hessian 退化时继续判别]\phantomsection\label{ans:c35q21}
两函数在原点梯度均为零，Hessian 也都为零矩阵，所以二阶充分判别法无法给出结论。对第一个函数，
\[
f(x,y)=x^4+3y^4>0\quad((x,y)\ne(0,0)),
\]
故原点是严格全局最小点。第二个函数沿 $y=0$ 为正，沿 $x=0$ 为负，所以原点不是极值点，而是鞍点。相同的退化 Hessian 可以对应不同性质，说明“二阶判别失效”不等于“没有极值”。
\end{solve}

\begin{solve}[习题~\ref{ex:c35q22}：Hessian 退化时继续判别]\phantomsection\label{ans:c35q22}
两函数在原点梯度均为零，Hessian 也都为零矩阵，所以二阶充分判别法无法给出结论。对第一个函数，
\[
f(x,y)=x^4+4y^4>0\quad((x,y)\ne(0,0)),
\]
故原点是严格全局最小点。第二个函数沿 $y=0$ 为正，沿 $x=0$ 为负，所以原点不是极值点，而是鞍点。相同的退化 Hessian 可以对应不同性质，说明“二阶判别失效”不等于“没有极值”。
\end{solve}

\begin{solve}[习题~\ref{ex:c35q23}：Hessian 退化时继续判别]\phantomsection\label{ans:c35q23}
两函数在原点梯度均为零，Hessian 也都为零矩阵，所以二阶充分判别法无法给出结论。对第一个函数，
\[
f(x,y)=x^4+5y^4>0\quad((x,y)\ne(0,0)),
\]
故原点是严格全局最小点。第二个函数沿 $y=0$ 为正，沿 $x=0$ 为负，所以原点不是极值点，而是鞍点。相同的退化 Hessian 可以对应不同性质，说明“二阶判别失效”不等于“没有极值”。
\end{solve}

\begin{solve}[习题~\ref{ex:c35q24}：Hessian 退化时继续判别]\phantomsection\label{ans:c35q24}
两函数在原点梯度均为零，Hessian 也都为零矩阵，所以二阶充分判别法无法给出结论。对第一个函数，
\[
f(x,y)=x^4+6y^4>0\quad((x,y)\ne(0,0)),
\]
故原点是严格全局最小点。第二个函数沿 $y=0$ 为正，沿 $x=0$ 为负，所以原点不是极值点，而是鞍点。相同的退化 Hessian 可以对应不同性质，说明“二阶判别失效”不等于“没有极值”。
\end{solve}

\clearpage
\section{二重积分与极坐标}\label{sec:ch36}
二重积分（double integral）的积分域与被积函数同样重要。交换积分次序时先描述同一个区域；极坐标换元的面积元为 $r\,dr\,d\theta$，其中的因子 $r$ 不可遗漏。

\textbf{资源对应：}CLP-3 \S~3.1--\S~3.3。本章共 25 道编号习题。
\subsection{习题}
\begin{example}[不作累次积分的几何计算]\label{ex:c36q1}
用几何意义计算 $\int_{-1}^{3}\int_{-4}^{1}dy\,dx$ 和 $\int_0^2\int_0^{\sqrt{4-y^2}}dx\,dy$。
\par\textit{选译。来源：CLP-3 \S~3.1，首个几何解释题的(a)(b)。}\par\hyperref[ans:c36q1]{解答见第~\pageref{ans:c36q1} 页。}
\end{example}
\begin{example}[矩形区域上的乘积分离]\label{ex:c36q2}
计算 $\displaystyle\iint_Dx^{0}y^{1}dx\,dy$，其中 $D=[0,2]\times[0,3]$；同时求区域内该函数的平均值。
\par\textit{改编训练。}\par\hyperref[ans:c36q2]{解答见第~\pageref{ans:c36q2} 页。}
\end{example}
\begin{example}[矩形区域上的乘积分离]\label{ex:c36q3}
计算 $\displaystyle\iint_Dx^{1}y^{1}dx\,dy$，其中 $D=[0,2]\times[0,3]$；同时求区域内该函数的平均值。
\par\textit{改编训练。}\par\hyperref[ans:c36q3]{解答见第~\pageref{ans:c36q3} 页。}
\end{example}
\begin{example}[矩形区域上的乘积分离]\label{ex:c36q4}
计算 $\displaystyle\iint_Dx^{1}y^{2}dx\,dy$，其中 $D=[0,2]\times[0,3]$；同时求区域内该函数的平均值。
\par\textit{改编训练。}\par\hyperref[ans:c36q4]{解答见第~\pageref{ans:c36q4} 页。}
\end{example}
\begin{example}[矩形区域上的乘积分离]\label{ex:c36q5}
计算 $\displaystyle\iint_Dx^{2}y^{1}dx\,dy$，其中 $D=[0,2]\times[0,3]$；同时求区域内该函数的平均值。
\par\textit{改编训练。}\par\hyperref[ans:c36q5]{解答见第~\pageref{ans:c36q5} 页。}
\end{example}
\begin{example}[矩形区域上的乘积分离]\label{ex:c36q6}
计算 $\displaystyle\iint_Dx^{2}y^{3}dx\,dy$，其中 $D=[0,2]\times[0,3]$；同时求区域内该函数的平均值。
\par\textit{改编训练。}\par\hyperref[ans:c36q6]{解答见第~\pageref{ans:c36q6} 页。}
\end{example}
\begin{example}[矩形区域上的乘积分离]\label{ex:c36q7}
计算 $\displaystyle\iint_Dx^{3}y^{2}dx\,dy$，其中 $D=[0,2]\times[0,3]$；同时求区域内该函数的平均值。
\par\textit{改编训练。}\par\hyperref[ans:c36q7]{解答见第~\pageref{ans:c36q7} 页。}
\end{example}
\begin{example}[三角形上的累次积分]\label{ex:c36q8}
设 $D=\{(x,y):x\ge0,y\ge0,x+y\le1\}$，计算 $\displaystyle\iint_Dx^{1}y\,dx\,dy$。
\par\textit{改编训练。}\par\hyperref[ans:c36q8]{解答见第~\pageref{ans:c36q8} 页。}
\end{example}
\begin{example}[三角形上的累次积分]\label{ex:c36q9}
设 $D=\{(x,y):x\ge0,y\ge0,x+y\le1\}$，计算 $\displaystyle\iint_Dx^{2}y\,dx\,dy$。
\par\textit{改编训练。}\par\hyperref[ans:c36q9]{解答见第~\pageref{ans:c36q9} 页。}
\end{example}
\begin{example}[三角形上的累次积分]\label{ex:c36q10}
设 $D=\{(x,y):x\ge0,y\ge0,x+y\le1\}$，计算 $\displaystyle\iint_Dx^{3}y\,dx\,dy$。
\par\textit{改编训练。}\par\hyperref[ans:c36q10]{解答见第~\pageref{ans:c36q10} 页。}
\end{example}
\begin{example}[三角形上的累次积分]\label{ex:c36q11}
设 $D=\{(x,y):x\ge0,y\ge0,x+y\le1\}$，计算 $\displaystyle\iint_Dx^{4}y\,dx\,dy$。
\par\textit{改编训练。}\par\hyperref[ans:c36q11]{解答见第~\pageref{ans:c36q11} 页。}
\end{example}
\begin{example}[三角形上的累次积分]\label{ex:c36q12}
设 $D=\{(x,y):x\ge0,y\ge0,x+y\le1\}$，计算 $\displaystyle\iint_Dx^{5}y\,dx\,dy$。
\par\textit{改编训练。}\par\hyperref[ans:c36q12]{解答见第~\pageref{ans:c36q12} 页。}
\end{example}
\begin{example}[三角形上的累次积分]\label{ex:c36q13}
设 $D=\{(x,y):x\ge0,y\ge0,x+y\le1\}$，计算 $\displaystyle\iint_Dx^{6}y\,dx\,dy$。
\par\textit{改编训练。}\par\hyperref[ans:c36q13]{解答见第~\pageref{ans:c36q13} 页。}
\end{example}
\begin{example}[圆盘上的径向函数]\label{ex:c36q14}
在圆盘 $D:x^2+y^2\le 1$ 上计算 $\displaystyle\iint_D(x^2+y^2)^{0}dx\,dy$。
\par\textit{改编训练。}\par\hyperref[ans:c36q14]{解答见第~\pageref{ans:c36q14} 页。}
\end{example}
\begin{example}[圆盘上的径向函数]\label{ex:c36q15}
在圆盘 $D:x^2+y^2\le 4$ 上计算 $\displaystyle\iint_D(x^2+y^2)^{1}dx\,dy$。
\par\textit{改编训练。}\par\hyperref[ans:c36q15]{解答见第~\pageref{ans:c36q15} 页。}
\end{example}
\begin{example}[圆盘上的径向函数]\label{ex:c36q16}
在圆盘 $D:x^2+y^2\le 1$ 上计算 $\displaystyle\iint_D(x^2+y^2)^{2}dx\,dy$。
\par\textit{改编训练。}\par\hyperref[ans:c36q16]{解答见第~\pageref{ans:c36q16} 页。}
\end{example}
\begin{example}[圆盘上的径向函数]\label{ex:c36q17}
在圆盘 $D:x^2+y^2\le 9$ 上计算 $\displaystyle\iint_D(x^2+y^2)^{1}dx\,dy$。
\par\textit{改编训练。}\par\hyperref[ans:c36q17]{解答见第~\pageref{ans:c36q17} 页。}
\end{example}
\begin{example}[圆盘上的径向函数]\label{ex:c36q18}
在圆盘 $D:x^2+y^2\le 4$ 上计算 $\displaystyle\iint_D(x^2+y^2)^{2}dx\,dy$。
\par\textit{改编训练。}\par\hyperref[ans:c36q18]{解答见第~\pageref{ans:c36q18} 页。}
\end{example}
\begin{example}[圆盘上的径向函数]\label{ex:c36q19}
在圆盘 $D:x^2+y^2\le 1$ 上计算 $\displaystyle\iint_D(x^2+y^2)^{3}dx\,dy$。
\par\textit{改编训练。}\par\hyperref[ans:c36q19]{解答见第~\pageref{ans:c36q19} 页。}
\end{example}
\begin{example}[椭圆换元与面积矩]\label{ex:c36q20}
设 $D:x^2/1+y^2/4\le1$。求面积和 $\displaystyle\iint_D(x^2/1+y^2/4)dx\,dy$。
\par\textit{改编训练。}\par\hyperref[ans:c36q20]{解答见第~\pageref{ans:c36q20} 页。}
\end{example}
\begin{example}[椭圆换元与面积矩]\label{ex:c36q21}
设 $D:x^2/4+y^2/9\le1$。求面积和 $\displaystyle\iint_D(x^2/4+y^2/9)dx\,dy$。
\par\textit{改编训练。}\par\hyperref[ans:c36q21]{解答见第~\pageref{ans:c36q21} 页。}
\end{example}
\begin{example}[椭圆换元与面积矩]\label{ex:c36q22}
设 $D:x^2/9+y^2/1\le1$。求面积和 $\displaystyle\iint_D(x^2/9+y^2/1)dx\,dy$。
\par\textit{改编训练。}\par\hyperref[ans:c36q22]{解答见第~\pageref{ans:c36q22} 页。}
\end{example}
\begin{example}[椭圆换元与面积矩]\label{ex:c36q23}
设 $D:x^2/4+y^2/25\le1$。求面积和 $\displaystyle\iint_D(x^2/4+y^2/25)dx\,dy$。
\par\textit{改编训练。}\par\hyperref[ans:c36q23]{解答见第~\pageref{ans:c36q23} 页。}
\end{example}
\begin{example}[椭圆换元与面积矩]\label{ex:c36q24}
设 $D:x^2/16+y^2/9\le1$。求面积和 $\displaystyle\iint_D(x^2/16+y^2/9)dx\,dy$。
\par\textit{改编训练。}\par\hyperref[ans:c36q24]{解答见第~\pageref{ans:c36q24} 页。}
\end{example}
\begin{example}[椭圆换元与面积矩]\label{ex:c36q25}
设 $D:x^2/9+y^2/25\le1$。求面积和 $\displaystyle\iint_D(x^2/9+y^2/25)dx\,dy$。
\par\textit{改编训练。}\par\hyperref[ans:c36q25]{解答见第~\pageref{ans:c36q25} 页。}
\end{example}
\subsection{解答}
\begin{solve}[习题~\ref{ex:c36q1}：不作累次积分的几何计算]\phantomsection\label{ans:c36q1}
第一式的被积函数为1，积分域为边长4与5的矩形，所以值为20。第二式的区域为第一象限内半径2的四分之一圆盘，面积为 $\pi\cdot2^2/4=\pi$。几何解释仍然依赖准确识别积分限，不能把第二式误看成半圆。
\end{solve}

\begin{solve}[习题~\ref{ex:c36q2}：矩形区域上的乘积分离]\phantomsection\label{ans:c36q2}
被积函数连续且区域为直积，Fubini 定理允许把二重积分分解为两个一维积分的乘积：
\[
I=\left(\int_0^2x^{0}dx\right)\left(\int_0^3y^{1}dy\right)=2\cdot\frac{9}{2}=9.
\]
区域面积为6，故平均值为 $\frac{3}{2}$。这种分离依赖矩形区域和乘积形式；若积分上限依赖另一个变量，就不能直接把积分写成两个独立积分相乘。
\end{solve}

\begin{solve}[习题~\ref{ex:c36q3}：矩形区域上的乘积分离]\phantomsection\label{ans:c36q3}
被积函数连续且区域为直积，Fubini 定理允许把二重积分分解为两个一维积分的乘积：
\[
I=\left(\int_0^2x^{1}dx\right)\left(\int_0^3y^{1}dy\right)=2\cdot\frac{9}{2}=9.
\]
区域面积为6，故平均值为 $\frac{3}{2}$。这种分离依赖矩形区域和乘积形式；若积分上限依赖另一个变量，就不能直接把积分写成两个独立积分相乘。
\end{solve}

\begin{solve}[习题~\ref{ex:c36q4}：矩形区域上的乘积分离]\phantomsection\label{ans:c36q4}
被积函数连续且区域为直积，Fubini 定理允许把二重积分分解为两个一维积分的乘积：
\[
I=\left(\int_0^2x^{1}dx\right)\left(\int_0^3y^{2}dy\right)=2\cdot9=18.
\]
区域面积为6，故平均值为 $3$。这种分离依赖矩形区域和乘积形式；若积分上限依赖另一个变量，就不能直接把积分写成两个独立积分相乘。
\end{solve}

\begin{solve}[习题~\ref{ex:c36q5}：矩形区域上的乘积分离]\phantomsection\label{ans:c36q5}
被积函数连续且区域为直积，Fubini 定理允许把二重积分分解为两个一维积分的乘积：
\[
I=\left(\int_0^2x^{2}dx\right)\left(\int_0^3y^{1}dy\right)=\frac{8}{3}\cdot\frac{9}{2}=12.
\]
区域面积为6，故平均值为 $2$。这种分离依赖矩形区域和乘积形式；若积分上限依赖另一个变量，就不能直接把积分写成两个独立积分相乘。
\end{solve}

\begin{solve}[习题~\ref{ex:c36q6}：矩形区域上的乘积分离]\phantomsection\label{ans:c36q6}
被积函数连续且区域为直积，Fubini 定理允许把二重积分分解为两个一维积分的乘积：
\[
I=\left(\int_0^2x^{2}dx\right)\left(\int_0^3y^{3}dy\right)=\frac{8}{3}\cdot\frac{81}{4}=54.
\]
区域面积为6，故平均值为 $9$。这种分离依赖矩形区域和乘积形式；若积分上限依赖另一个变量，就不能直接把积分写成两个独立积分相乘。
\end{solve}

\begin{solve}[习题~\ref{ex:c36q7}：矩形区域上的乘积分离]\phantomsection\label{ans:c36q7}
被积函数连续且区域为直积，Fubini 定理允许把二重积分分解为两个一维积分的乘积：
\[
I=\left(\int_0^2x^{3}dx\right)\left(\int_0^3y^{2}dy\right)=4\cdot9=36.
\]
区域面积为6，故平均值为 $6$。这种分离依赖矩形区域和乘积形式；若积分上限依赖另一个变量，就不能直接把积分写成两个独立积分相乘。
\end{solve}

\begin{solve}[习题~\ref{ex:c36q8}：三角形上的累次积分]\phantomsection\label{ans:c36q8}
先固定 $x$，则 $0\le y\le1-x$、$0\le x\le1$，所以
\[
I=\int_0^1x^{1}\left[\frac{y^2}2\right]_0^{1-x}dx=\frac12\int_0^1x^{1}(1-x)^2dx.
\]
展开并逐项积分，
\[
I=\frac12\left(\frac1{2}-\frac2{3}+\frac1{4}\right)=\frac{1}{24}.
\]
这个结果也可先对 $x$ 积分得到。积分限描述的是同一个三角形，交换次序后外层变量的范围与内层上限都要一起改变。
\end{solve}

\begin{solve}[习题~\ref{ex:c36q9}：三角形上的累次积分]\phantomsection\label{ans:c36q9}
先固定 $x$，则 $0\le y\le1-x$、$0\le x\le1$，所以
\[
I=\int_0^1x^{2}\left[\frac{y^2}2\right]_0^{1-x}dx=\frac12\int_0^1x^{2}(1-x)^2dx.
\]
展开并逐项积分，
\[
I=\frac12\left(\frac1{3}-\frac2{4}+\frac1{5}\right)=\frac{1}{60}.
\]
这个结果也可先对 $x$ 积分得到。积分限描述的是同一个三角形，交换次序后外层变量的范围与内层上限都要一起改变。
\end{solve}

\begin{solve}[习题~\ref{ex:c36q10}：三角形上的累次积分]\phantomsection\label{ans:c36q10}
先固定 $x$，则 $0\le y\le1-x$、$0\le x\le1$，所以
\[
I=\int_0^1x^{3}\left[\frac{y^2}2\right]_0^{1-x}dx=\frac12\int_0^1x^{3}(1-x)^2dx.
\]
展开并逐项积分，
\[
I=\frac12\left(\frac1{4}-\frac2{5}+\frac1{6}\right)=\frac{1}{120}.
\]
这个结果也可先对 $x$ 积分得到。积分限描述的是同一个三角形，交换次序后外层变量的范围与内层上限都要一起改变。
\end{solve}

\begin{solve}[习题~\ref{ex:c36q11}：三角形上的累次积分]\phantomsection\label{ans:c36q11}
先固定 $x$，则 $0\le y\le1-x$、$0\le x\le1$，所以
\[
I=\int_0^1x^{4}\left[\frac{y^2}2\right]_0^{1-x}dx=\frac12\int_0^1x^{4}(1-x)^2dx.
\]
展开并逐项积分，
\[
I=\frac12\left(\frac1{5}-\frac2{6}+\frac1{7}\right)=\frac{1}{210}.
\]
这个结果也可先对 $x$ 积分得到。积分限描述的是同一个三角形，交换次序后外层变量的范围与内层上限都要一起改变。
\end{solve}

\begin{solve}[习题~\ref{ex:c36q12}：三角形上的累次积分]\phantomsection\label{ans:c36q12}
先固定 $x$，则 $0\le y\le1-x$、$0\le x\le1$，所以
\[
I=\int_0^1x^{5}\left[\frac{y^2}2\right]_0^{1-x}dx=\frac12\int_0^1x^{5}(1-x)^2dx.
\]
展开并逐项积分，
\[
I=\frac12\left(\frac1{6}-\frac2{7}+\frac1{8}\right)=\frac{1}{336}.
\]
这个结果也可先对 $x$ 积分得到。积分限描述的是同一个三角形，交换次序后外层变量的范围与内层上限都要一起改变。
\end{solve}

\begin{solve}[习题~\ref{ex:c36q13}：三角形上的累次积分]\phantomsection\label{ans:c36q13}
先固定 $x$，则 $0\le y\le1-x$、$0\le x\le1$，所以
\[
I=\int_0^1x^{6}\left[\frac{y^2}2\right]_0^{1-x}dx=\frac12\int_0^1x^{6}(1-x)^2dx.
\]
展开并逐项积分，
\[
I=\frac12\left(\frac1{7}-\frac2{8}+\frac1{9}\right)=\frac{1}{504}.
\]
这个结果也可先对 $x$ 积分得到。积分限描述的是同一个三角形，交换次序后外层变量的范围与内层上限都要一起改变。
\end{solve}

\begin{solve}[习题~\ref{ex:c36q14}：圆盘上的径向函数]\phantomsection\label{ans:c36q14}
取 $x=r\cos\theta$、$y=r\sin\theta$，范围为 $0\le r\le 1$、$0\le\theta\le2\pi$，Jacobian 绝对值为 $r$。因此
\[
I=\int_0^{2\pi}d\theta\int_0^{1}r^{1}dr=1\pi.
\]
在圆心处坐标变换的 Jacobian 为零，但这一单点对积分没有影响；换元在其余区域分片应用即可。结果的长度次数为 $2m+2$，与被积函数次数加上面积元的次数一致，可作为尺度检验。
\end{solve}

\begin{solve}[习题~\ref{ex:c36q15}：圆盘上的径向函数]\phantomsection\label{ans:c36q15}
取 $x=r\cos\theta$、$y=r\sin\theta$，范围为 $0\le r\le 2$、$0\le\theta\le2\pi$，Jacobian 绝对值为 $r$。因此
\[
I=\int_0^{2\pi}d\theta\int_0^{2}r^{3}dr=8\pi.
\]
在圆心处坐标变换的 Jacobian 为零，但这一单点对积分没有影响；换元在其余区域分片应用即可。结果的长度次数为 $2m+2$，与被积函数次数加上面积元的次数一致，可作为尺度检验。
\end{solve}

\begin{solve}[习题~\ref{ex:c36q16}：圆盘上的径向函数]\phantomsection\label{ans:c36q16}
取 $x=r\cos\theta$、$y=r\sin\theta$，范围为 $0\le r\le 1$、$0\le\theta\le2\pi$，Jacobian 绝对值为 $r$。因此
\[
I=\int_0^{2\pi}d\theta\int_0^{1}r^{5}dr=\frac{1}{3}\pi.
\]
在圆心处坐标变换的 Jacobian 为零，但这一单点对积分没有影响；换元在其余区域分片应用即可。结果的长度次数为 $2m+2$，与被积函数次数加上面积元的次数一致，可作为尺度检验。
\end{solve}

\begin{solve}[习题~\ref{ex:c36q17}：圆盘上的径向函数]\phantomsection\label{ans:c36q17}
取 $x=r\cos\theta$、$y=r\sin\theta$，范围为 $0\le r\le 3$、$0\le\theta\le2\pi$，Jacobian 绝对值为 $r$。因此
\[
I=\int_0^{2\pi}d\theta\int_0^{3}r^{3}dr=\frac{81}{2}\pi.
\]
在圆心处坐标变换的 Jacobian 为零，但这一单点对积分没有影响；换元在其余区域分片应用即可。结果的长度次数为 $2m+2$，与被积函数次数加上面积元的次数一致，可作为尺度检验。
\end{solve}

\begin{solve}[习题~\ref{ex:c36q18}：圆盘上的径向函数]\phantomsection\label{ans:c36q18}
取 $x=r\cos\theta$、$y=r\sin\theta$，范围为 $0\le r\le 2$、$0\le\theta\le2\pi$，Jacobian 绝对值为 $r$。因此
\[
I=\int_0^{2\pi}d\theta\int_0^{2}r^{5}dr=\frac{64}{3}\pi.
\]
在圆心处坐标变换的 Jacobian 为零，但这一单点对积分没有影响；换元在其余区域分片应用即可。结果的长度次数为 $2m+2$，与被积函数次数加上面积元的次数一致，可作为尺度检验。
\end{solve}

\begin{solve}[习题~\ref{ex:c36q19}：圆盘上的径向函数]\phantomsection\label{ans:c36q19}
取 $x=r\cos\theta$、$y=r\sin\theta$，范围为 $0\le r\le 1$、$0\le\theta\le2\pi$，Jacobian 绝对值为 $r$。因此
\[
I=\int_0^{2\pi}d\theta\int_0^{1}r^{7}dr=\frac{1}{4}\pi.
\]
在圆心处坐标变换的 Jacobian 为零，但这一单点对积分没有影响；换元在其余区域分片应用即可。结果的长度次数为 $2m+2$，与被积函数次数加上面积元的次数一致，可作为尺度检验。
\end{solve}

\begin{solve}[习题~\ref{ex:c36q20}：椭圆换元与面积矩]\phantomsection\label{ans:c36q20}
令 $x=1u$、$y=2v$，椭圆变成单位圆盘，Jacobian 为 $2$。于是面积为 $2\pi$。对于第二式，再在 $uv$ 平面用极坐标，得到
\[
I=2\int_0^{2\pi}\int_0^1r^2r\,dr\,d\theta=1\pi.
\]
积分平均值等于 $1/2$，与两条半轴长度无关，因为被积函数使用了同样的归一化坐标。连续进行两次换元时，两次 Jacobian 因子必须都保留。
\end{solve}

\begin{solve}[习题~\ref{ex:c36q21}：椭圆换元与面积矩]\phantomsection\label{ans:c36q21}
令 $x=2u$、$y=3v$，椭圆变成单位圆盘，Jacobian 为 $6$。于是面积为 $6\pi$。对于第二式，再在 $uv$ 平面用极坐标，得到
\[
I=6\int_0^{2\pi}\int_0^1r^2r\,dr\,d\theta=3\pi.
\]
积分平均值等于 $1/2$，与两条半轴长度无关，因为被积函数使用了同样的归一化坐标。连续进行两次换元时，两次 Jacobian 因子必须都保留。
\end{solve}

\begin{solve}[习题~\ref{ex:c36q22}：椭圆换元与面积矩]\phantomsection\label{ans:c36q22}
令 $x=3u$、$y=1v$，椭圆变成单位圆盘，Jacobian 为 $3$。于是面积为 $3\pi$。对于第二式，再在 $uv$ 平面用极坐标，得到
\[
I=3\int_0^{2\pi}\int_0^1r^2r\,dr\,d\theta=\frac{3}{2}\pi.
\]
积分平均值等于 $1/2$，与两条半轴长度无关，因为被积函数使用了同样的归一化坐标。连续进行两次换元时，两次 Jacobian 因子必须都保留。
\end{solve}

\begin{solve}[习题~\ref{ex:c36q23}：椭圆换元与面积矩]\phantomsection\label{ans:c36q23}
令 $x=2u$、$y=5v$，椭圆变成单位圆盘，Jacobian 为 $10$。于是面积为 $10\pi$。对于第二式，再在 $uv$ 平面用极坐标，得到
\[
I=10\int_0^{2\pi}\int_0^1r^2r\,dr\,d\theta=5\pi.
\]
积分平均值等于 $1/2$，与两条半轴长度无关，因为被积函数使用了同样的归一化坐标。连续进行两次换元时，两次 Jacobian 因子必须都保留。
\end{solve}

\begin{solve}[习题~\ref{ex:c36q24}：椭圆换元与面积矩]\phantomsection\label{ans:c36q24}
令 $x=4u$、$y=3v$，椭圆变成单位圆盘，Jacobian 为 $12$。于是面积为 $12\pi$。对于第二式，再在 $uv$ 平面用极坐标，得到
\[
I=12\int_0^{2\pi}\int_0^1r^2r\,dr\,d\theta=6\pi.
\]
积分平均值等于 $1/2$，与两条半轴长度无关，因为被积函数使用了同样的归一化坐标。连续进行两次换元时，两次 Jacobian 因子必须都保留。
\end{solve}

\begin{solve}[习题~\ref{ex:c36q25}：椭圆换元与面积矩]\phantomsection\label{ans:c36q25}
令 $x=3u$、$y=5v$，椭圆变成单位圆盘，Jacobian 为 $15$。于是面积为 $15\pi$。对于第二式，再在 $uv$ 平面用极坐标，得到
\[
I=15\int_0^{2\pi}\int_0^1r^2r\,dr\,d\theta=\frac{15}{2}\pi.
\]
积分平均值等于 $1/2$，与两条半轴长度无关，因为被积函数使用了同样的归一化坐标。连续进行两次换元时，两次 Jacobian 因子必须都保留。
\end{solve}

\clearpage
\section{交换积分次序与二重积分换元}\label{sec:ch37}
交换积分次序首先是一道区域描述题。直接积分出现非初等原函数时，改变积分顺序可能使问题简化；换元除了替换被积函数，还要替换区域和面积元。

\textbf{资源对应：}CLP-3 \S~3.1、\S~3.2 的综合训练。本章共 24 道编号习题。
\subsection{习题}
\begin{example}[非初等内积分的换序]\label{ex:c37q1}
计算 $I=\displaystyle\int_0^{1}\int_x^{1}e^{y^2}dy\,dx$。
\par\textit{改编训练。}\par\hyperref[ans:c37q1]{解答见第~\pageref{ans:c37q1} 页。}
\end{example}
\begin{example}[非初等内积分的换序]\label{ex:c37q2}
计算 $I=\displaystyle\int_0^{2}\int_x^{2}e^{y^2}dy\,dx$。
\par\textit{改编训练。}\par\hyperref[ans:c37q2]{解答见第~\pageref{ans:c37q2} 页。}
\end{example}
\begin{example}[非初等内积分的换序]\label{ex:c37q3}
计算 $I=\displaystyle\int_0^{3}\int_x^{3}e^{y^2}dy\,dx$。
\par\textit{改编训练。}\par\hyperref[ans:c37q3]{解答见第~\pageref{ans:c37q3} 页。}
\end{example}
\begin{example}[非初等内积分的换序]\label{ex:c37q4}
计算 $I=\displaystyle\int_0^{4}\int_x^{4}e^{y^2}dy\,dx$。
\par\textit{改编训练。}\par\hyperref[ans:c37q4]{解答见第~\pageref{ans:c37q4} 页。}
\end{example}
\begin{example}[非初等内积分的换序]\label{ex:c37q5}
计算 $I=\displaystyle\int_0^{5}\int_x^{5}e^{y^2}dy\,dx$。
\par\textit{改编训练。}\par\hyperref[ans:c37q5]{解答见第~\pageref{ans:c37q5} 页。}
\end{example}
\begin{example}[非初等内积分的换序]\label{ex:c37q6}
计算 $I=\displaystyle\int_0^{6}\int_x^{6}e^{y^2}dy\,dx$。
\par\textit{改编训练。}\par\hyperref[ans:c37q6]{解答见第~\pageref{ans:c37q6} 页。}
\end{example}
\begin{example}[两条曲线之间的换序]\label{ex:c37q7}
对区域 $D:0\le x\le1, x^{2}\le y\le x$，写出两种积分次序，并计算 $\iint_Dy\,dxdy$。
\par\textit{改编训练。}\par\hyperref[ans:c37q7]{解答见第~\pageref{ans:c37q7} 页。}
\end{example}
\begin{example}[两条曲线之间的换序]\label{ex:c37q8}
对区域 $D:0\le x\le1, x^{3}\le y\le x$，写出两种积分次序，并计算 $\iint_Dy\,dxdy$。
\par\textit{改编训练。}\par\hyperref[ans:c37q8]{解答见第~\pageref{ans:c37q8} 页。}
\end{example}
\begin{example}[两条曲线之间的换序]\label{ex:c37q9}
对区域 $D:0\le x\le1, x^{4}\le y\le x$，写出两种积分次序，并计算 $\iint_Dy\,dxdy$。
\par\textit{改编训练。}\par\hyperref[ans:c37q9]{解答见第~\pageref{ans:c37q9} 页。}
\end{example}
\begin{example}[两条曲线之间的换序]\label{ex:c37q10}
对区域 $D:0\le x\le1, x^{5}\le y\le x$，写出两种积分次序，并计算 $\iint_Dy\,dxdy$。
\par\textit{改编训练。}\par\hyperref[ans:c37q10]{解答见第~\pageref{ans:c37q10} 页。}
\end{example}
\begin{example}[两条曲线之间的换序]\label{ex:c37q11}
对区域 $D:0\le x\le1, x^{6}\le y\le x$，写出两种积分次序，并计算 $\iint_Dy\,dxdy$。
\par\textit{改编训练。}\par\hyperref[ans:c37q11]{解答见第~\pageref{ans:c37q11} 页。}
\end{example}
\begin{example}[两条曲线之间的换序]\label{ex:c37q12}
对区域 $D:0\le x\le1, x^{7}\le y\le x$，写出两种积分次序，并计算 $\iint_Dy\,dxdy$。
\par\textit{改编训练。}\par\hyperref[ans:c37q12]{解答见第~\pageref{ans:c37q12} 页。}
\end{example}
\begin{example}[圆环上的径向负幂]\label{ex:c37q13}
在圆环 $D:1\le x^2+y^2\le4$ 上计算 $\iint_D(x^2+y^2)^{-0}dxdy$。
\par\textit{改编训练。}\par\hyperref[ans:c37q13]{解答见第~\pageref{ans:c37q13} 页。}
\end{example}
\begin{example}[圆环上的径向负幂]\label{ex:c37q14}
在圆环 $D:1\le x^2+y^2\le4$ 上计算 $\iint_D(x^2+y^2)^{-\frac{1}{2}}dxdy$。
\par\textit{改编训练。}\par\hyperref[ans:c37q14]{解答见第~\pageref{ans:c37q14} 页。}
\end{example}
\begin{example}[圆环上的径向负幂]\label{ex:c37q15}
在圆环 $D:1\le x^2+y^2\le4$ 上计算 $\iint_D(x^2+y^2)^{-1}dxdy$。
\par\textit{改编训练。}\par\hyperref[ans:c37q15]{解答见第~\pageref{ans:c37q15} 页。}
\end{example}
\begin{example}[圆环上的径向负幂]\label{ex:c37q16}
在圆环 $D:1\le x^2+y^2\le4$ 上计算 $\iint_D(x^2+y^2)^{-\frac{3}{2}}dxdy$。
\par\textit{改编训练。}\par\hyperref[ans:c37q16]{解答见第~\pageref{ans:c37q16} 页。}
\end{example}
\begin{example}[圆环上的径向负幂]\label{ex:c37q17}
在圆环 $D:1\le x^2+y^2\le4$ 上计算 $\iint_D(x^2+y^2)^{-2}dxdy$。
\par\textit{改编训练。}\par\hyperref[ans:c37q17]{解答见第~\pageref{ans:c37q17} 页。}
\end{example}
\begin{example}[圆环上的径向负幂]\label{ex:c37q18}
在圆环 $D:1\le x^2+y^2\le4$ 上计算 $\iint_D(x^2+y^2)^{-\frac{5}{2}}dxdy$。
\par\textit{改编训练。}\par\hyperref[ans:c37q18]{解答见第~\pageref{ans:c37q18} 页。}
\end{example}
\begin{example}[沿斜边界的线性换元]\label{ex:c37q19}
设 $D$ 由 $0\le x+y\le 1$、$0\le x-y\le 2$ 给出。求 $\displaystyle\iint_D(x+y)^2dxdy$。
\par\textit{改编训练。}\par\hyperref[ans:c37q19]{解答见第~\pageref{ans:c37q19} 页。}
\end{example}
\begin{example}[沿斜边界的线性换元]\label{ex:c37q20}
设 $D$ 由 $0\le x+y\le 2$、$0\le x-y\le 1$ 给出。求 $\displaystyle\iint_D(x+y)^2dxdy$。
\par\textit{改编训练。}\par\hyperref[ans:c37q20]{解答见第~\pageref{ans:c37q20} 页。}
\end{example}
\begin{example}[沿斜边界的线性换元]\label{ex:c37q21}
设 $D$ 由 $0\le x+y\le 2$、$0\le x-y\le 3$ 给出。求 $\displaystyle\iint_D(x+y)^2dxdy$。
\par\textit{改编训练。}\par\hyperref[ans:c37q21]{解答见第~\pageref{ans:c37q21} 页。}
\end{example}
\begin{example}[沿斜边界的线性换元]\label{ex:c37q22}
设 $D$ 由 $0\le x+y\le 3$、$0\le x-y\le 2$ 给出。求 $\displaystyle\iint_D(x+y)^2dxdy$。
\par\textit{改编训练。}\par\hyperref[ans:c37q22]{解答见第~\pageref{ans:c37q22} 页。}
\end{example}
\begin{example}[沿斜边界的线性换元]\label{ex:c37q23}
设 $D$ 由 $0\le x+y\le 4$、$0\le x-y\le 1$ 给出。求 $\displaystyle\iint_D(x+y)^2dxdy$。
\par\textit{改编训练。}\par\hyperref[ans:c37q23]{解答见第~\pageref{ans:c37q23} 页。}
\end{example}
\begin{example}[沿斜边界的线性换元]\label{ex:c37q24}
设 $D$ 由 $0\le x+y\le 3$、$0\le x-y\le 4$ 给出。求 $\displaystyle\iint_D(x+y)^2dxdy$。
\par\textit{改编训练。}\par\hyperref[ans:c37q24]{解答见第~\pageref{ans:c37q24} 页。}
\end{example}
\subsection{解答}
\begin{solve}[习题~\ref{ex:c37q1}：非初等内积分的换序]\phantomsection\label{ans:c37q1}
积分域是 $0\le x\le y\le 1$。固定 $y$ 后，$x$ 从0到 $y$，于是
\[
I=\int_0^{1}\int_0^y e^{y^2}dx\,dy=\int_0^{1}y e^{y^2}dy=\frac{e^{1}-1}2.
\]
连续函数在有界闭三角形上可积，因此换序合法。原来的内积分没有初等表达式，并不妨碍二重积分有简洁结果；新次序产生的因子 $y$ 恰好匹配指数内层的导数。
\end{solve}

\begin{solve}[习题~\ref{ex:c37q2}：非初等内积分的换序]\phantomsection\label{ans:c37q2}
积分域是 $0\le x\le y\le 2$。固定 $y$ 后，$x$ 从0到 $y$，于是
\[
I=\int_0^{2}\int_0^y e^{y^2}dx\,dy=\int_0^{2}y e^{y^2}dy=\frac{e^{4}-1}2.
\]
连续函数在有界闭三角形上可积，因此换序合法。原来的内积分没有初等表达式，并不妨碍二重积分有简洁结果；新次序产生的因子 $y$ 恰好匹配指数内层的导数。
\end{solve}

\begin{solve}[习题~\ref{ex:c37q3}：非初等内积分的换序]\phantomsection\label{ans:c37q3}
积分域是 $0\le x\le y\le 3$。固定 $y$ 后，$x$ 从0到 $y$，于是
\[
I=\int_0^{3}\int_0^y e^{y^2}dx\,dy=\int_0^{3}y e^{y^2}dy=\frac{e^{9}-1}2.
\]
连续函数在有界闭三角形上可积，因此换序合法。原来的内积分没有初等表达式，并不妨碍二重积分有简洁结果；新次序产生的因子 $y$ 恰好匹配指数内层的导数。
\end{solve}

\begin{solve}[习题~\ref{ex:c37q4}：非初等内积分的换序]\phantomsection\label{ans:c37q4}
积分域是 $0\le x\le y\le 4$。固定 $y$ 后，$x$ 从0到 $y$，于是
\[
I=\int_0^{4}\int_0^y e^{y^2}dx\,dy=\int_0^{4}y e^{y^2}dy=\frac{e^{16}-1}2.
\]
连续函数在有界闭三角形上可积，因此换序合法。原来的内积分没有初等表达式，并不妨碍二重积分有简洁结果；新次序产生的因子 $y$ 恰好匹配指数内层的导数。
\end{solve}

\begin{solve}[习题~\ref{ex:c37q5}：非初等内积分的换序]\phantomsection\label{ans:c37q5}
积分域是 $0\le x\le y\le 5$。固定 $y$ 后，$x$ 从0到 $y$，于是
\[
I=\int_0^{5}\int_0^y e^{y^2}dx\,dy=\int_0^{5}y e^{y^2}dy=\frac{e^{25}-1}2.
\]
连续函数在有界闭三角形上可积，因此换序合法。原来的内积分没有初等表达式，并不妨碍二重积分有简洁结果；新次序产生的因子 $y$ 恰好匹配指数内层的导数。
\end{solve}

\begin{solve}[习题~\ref{ex:c37q6}：非初等内积分的换序]\phantomsection\label{ans:c37q6}
积分域是 $0\le x\le y\le 6$。固定 $y$ 后，$x$ 从0到 $y$，于是
\[
I=\int_0^{6}\int_0^y e^{y^2}dx\,dy=\int_0^{6}y e^{y^2}dy=\frac{e^{36}-1}2.
\]
连续函数在有界闭三角形上可积，因此换序合法。原来的内积分没有初等表达式，并不妨碍二重积分有简洁结果；新次序产生的因子 $y$ 恰好匹配指数内层的导数。
\end{solve}

\begin{solve}[习题~\ref{ex:c37q7}：两条曲线之间的换序]\phantomsection\label{ans:c37q7}
从 $x^{2}\le y\le x$ 得到 $y\le x\le y^{1/2}$，而 $0\le y\le1$。因此
\[
I=\int_0^1\int_{x^{2}}^x y\,dy\,dx=\int_0^1\int_y^{y^{1/2}}y\,dx\,dy.
\]
按第一种次序，$I=\frac12\int_0^1(x^2-x^{4})dx$，所以
\[
I=\frac12\left(\frac13-\frac1{5}\right)=\frac{1}{15}.
\]
反解边界时，指数的倒数只对非负区间使用，避免了多值根的问题。
\end{solve}

\begin{solve}[习题~\ref{ex:c37q8}：两条曲线之间的换序]\phantomsection\label{ans:c37q8}
从 $x^{3}\le y\le x$ 得到 $y\le x\le y^{1/3}$，而 $0\le y\le1$。因此
\[
I=\int_0^1\int_{x^{3}}^x y\,dy\,dx=\int_0^1\int_y^{y^{1/3}}y\,dx\,dy.
\]
按第一种次序，$I=\frac12\int_0^1(x^2-x^{6})dx$，所以
\[
I=\frac12\left(\frac13-\frac1{7}\right)=\frac{2}{21}.
\]
反解边界时，指数的倒数只对非负区间使用，避免了多值根的问题。
\end{solve}

\begin{solve}[习题~\ref{ex:c37q9}：两条曲线之间的换序]\phantomsection\label{ans:c37q9}
从 $x^{4}\le y\le x$ 得到 $y\le x\le y^{1/4}$，而 $0\le y\le1$。因此
\[
I=\int_0^1\int_{x^{4}}^x y\,dy\,dx=\int_0^1\int_y^{y^{1/4}}y\,dx\,dy.
\]
按第一种次序，$I=\frac12\int_0^1(x^2-x^{8})dx$，所以
\[
I=\frac12\left(\frac13-\frac1{9}\right)=\frac{1}{9}.
\]
反解边界时，指数的倒数只对非负区间使用，避免了多值根的问题。
\end{solve}

\begin{solve}[习题~\ref{ex:c37q10}：两条曲线之间的换序]\phantomsection\label{ans:c37q10}
从 $x^{5}\le y\le x$ 得到 $y\le x\le y^{1/5}$，而 $0\le y\le1$。因此
\[
I=\int_0^1\int_{x^{5}}^x y\,dy\,dx=\int_0^1\int_y^{y^{1/5}}y\,dx\,dy.
\]
按第一种次序，$I=\frac12\int_0^1(x^2-x^{10})dx$，所以
\[
I=\frac12\left(\frac13-\frac1{11}\right)=\frac{4}{33}.
\]
反解边界时，指数的倒数只对非负区间使用，避免了多值根的问题。
\end{solve}

\begin{solve}[习题~\ref{ex:c37q11}：两条曲线之间的换序]\phantomsection\label{ans:c37q11}
从 $x^{6}\le y\le x$ 得到 $y\le x\le y^{1/6}$，而 $0\le y\le1$。因此
\[
I=\int_0^1\int_{x^{6}}^x y\,dy\,dx=\int_0^1\int_y^{y^{1/6}}y\,dx\,dy.
\]
按第一种次序，$I=\frac12\int_0^1(x^2-x^{12})dx$，所以
\[
I=\frac12\left(\frac13-\frac1{13}\right)=\frac{5}{39}.
\]
反解边界时，指数的倒数只对非负区间使用，避免了多值根的问题。
\end{solve}

\begin{solve}[习题~\ref{ex:c37q12}：两条曲线之间的换序]\phantomsection\label{ans:c37q12}
从 $x^{7}\le y\le x$ 得到 $y\le x\le y^{1/7}$，而 $0\le y\le1$。因此
\[
I=\int_0^1\int_{x^{7}}^x y\,dy\,dx=\int_0^1\int_y^{y^{1/7}}y\,dx\,dy.
\]
按第一种次序，$I=\frac12\int_0^1(x^2-x^{14})dx$，所以
\[
I=\frac12\left(\frac13-\frac1{15}\right)=\frac{2}{15}.
\]
反解边界时，指数的倒数只对非负区间使用，避免了多值根的问题。
\end{solve}

\begin{solve}[习题~\ref{ex:c37q13}：圆环上的径向负幂]\phantomsection\label{ans:c37q13}
极坐标中 $1\le r\le2$，角度覆盖一周，所以
\[
I=2\pi\int_1^2r^{1}dr.
\]

\[
I=\frac{2\pi}{2}(2^{2}-1)=3\pi.
\]
圆环与原点保持正距离，故即使被积函数在原点奇异，本题也只是普通的有界区域积分。判断奇点是否造成反常积分，必须结合实际积分域。
\end{solve}

\begin{solve}[习题~\ref{ex:c37q14}：圆环上的径向负幂]\phantomsection\label{ans:c37q14}
极坐标中 $1\le r\le2$，角度覆盖一周，所以
\[
I=2\pi\int_1^2r^{0}dr.
\]

\[
I=\frac{2\pi}{1}(2^{1}-1)=2\pi.
\]
圆环与原点保持正距离，故即使被积函数在原点奇异，本题也只是普通的有界区域积分。判断奇点是否造成反常积分，必须结合实际积分域。
\end{solve}

\begin{solve}[习题~\ref{ex:c37q15}：圆环上的径向负幂]\phantomsection\label{ans:c37q15}
极坐标中 $1\le r\le2$，角度覆盖一周，所以
\[
I=2\pi\int_1^2r^{-1}dr.
\]

\[
I=2\pi\log2.
\]
圆环与原点保持正距离，故即使被积函数在原点奇异，本题也只是普通的有界区域积分。判断奇点是否造成反常积分，必须结合实际积分域。
\end{solve}

\begin{solve}[习题~\ref{ex:c37q16}：圆环上的径向负幂]\phantomsection\label{ans:c37q16}
极坐标中 $1\le r\le2$，角度覆盖一周，所以
\[
I=2\pi\int_1^2r^{-2}dr.
\]

\[
I=\frac{2\pi}{-1}(2^{-1}-1)=1\pi.
\]
圆环与原点保持正距离，故即使被积函数在原点奇异，本题也只是普通的有界区域积分。判断奇点是否造成反常积分，必须结合实际积分域。
\end{solve}

\begin{solve}[习题~\ref{ex:c37q17}：圆环上的径向负幂]\phantomsection\label{ans:c37q17}
极坐标中 $1\le r\le2$，角度覆盖一周，所以
\[
I=2\pi\int_1^2r^{-3}dr.
\]

\[
I=\frac{2\pi}{-2}(2^{-2}-1)=\frac{3}{4}\pi.
\]
圆环与原点保持正距离，故即使被积函数在原点奇异，本题也只是普通的有界区域积分。判断奇点是否造成反常积分，必须结合实际积分域。
\end{solve}

\begin{solve}[习题~\ref{ex:c37q18}：圆环上的径向负幂]\phantomsection\label{ans:c37q18}
极坐标中 $1\le r\le2$，角度覆盖一周，所以
\[
I=2\pi\int_1^2r^{-4}dr.
\]

\[
I=\frac{2\pi}{-3}(2^{-3}-1)=\frac{7}{12}\pi.
\]
圆环与原点保持正距离，故即使被积函数在原点奇异，本题也只是普通的有界区域积分。判断奇点是否造成反常积分，必须结合实际积分域。
\end{solve}

\begin{solve}[习题~\ref{ex:c37q19}：沿斜边界的线性换元]\phantomsection\label{ans:c37q19}
取 $u=x+y$、$v=x-y$，则 $x=(u+v)/2$、$y=(u-v)/2$，Jacobian 行列式为 $-1/2$，面积元取其绝对值。因此
\[
I=\frac12\int_0^{1}\int_0^{2}u^2dv\,du=\frac{1}{3}.
\]
负行列式表示坐标变换反转了定向，但通常的面积积分使用绝对值，不会因此变为负数。新变量同时把两组斜直线边界变成坐标平行的矩形。
\end{solve}

\begin{solve}[习题~\ref{ex:c37q20}：沿斜边界的线性换元]\phantomsection\label{ans:c37q20}
取 $u=x+y$、$v=x-y$，则 $x=(u+v)/2$、$y=(u-v)/2$，Jacobian 行列式为 $-1/2$，面积元取其绝对值。因此
\[
I=\frac12\int_0^{2}\int_0^{1}u^2dv\,du=\frac{4}{3}.
\]
负行列式表示坐标变换反转了定向，但通常的面积积分使用绝对值，不会因此变为负数。新变量同时把两组斜直线边界变成坐标平行的矩形。
\end{solve}

\begin{solve}[习题~\ref{ex:c37q21}：沿斜边界的线性换元]\phantomsection\label{ans:c37q21}
取 $u=x+y$、$v=x-y$，则 $x=(u+v)/2$、$y=(u-v)/2$，Jacobian 行列式为 $-1/2$，面积元取其绝对值。因此
\[
I=\frac12\int_0^{2}\int_0^{3}u^2dv\,du=4.
\]
负行列式表示坐标变换反转了定向，但通常的面积积分使用绝对值，不会因此变为负数。新变量同时把两组斜直线边界变成坐标平行的矩形。
\end{solve}

\begin{solve}[习题~\ref{ex:c37q22}：沿斜边界的线性换元]\phantomsection\label{ans:c37q22}
取 $u=x+y$、$v=x-y$，则 $x=(u+v)/2$、$y=(u-v)/2$，Jacobian 行列式为 $-1/2$，面积元取其绝对值。因此
\[
I=\frac12\int_0^{3}\int_0^{2}u^2dv\,du=9.
\]
负行列式表示坐标变换反转了定向，但通常的面积积分使用绝对值，不会因此变为负数。新变量同时把两组斜直线边界变成坐标平行的矩形。
\end{solve}

\begin{solve}[习题~\ref{ex:c37q23}：沿斜边界的线性换元]\phantomsection\label{ans:c37q23}
取 $u=x+y$、$v=x-y$，则 $x=(u+v)/2$、$y=(u-v)/2$，Jacobian 行列式为 $-1/2$，面积元取其绝对值。因此
\[
I=\frac12\int_0^{4}\int_0^{1}u^2dv\,du=\frac{32}{3}.
\]
负行列式表示坐标变换反转了定向，但通常的面积积分使用绝对值，不会因此变为负数。新变量同时把两组斜直线边界变成坐标平行的矩形。
\end{solve}

\begin{solve}[习题~\ref{ex:c37q24}：沿斜边界的线性换元]\phantomsection\label{ans:c37q24}
取 $u=x+y$、$v=x-y$，则 $x=(u+v)/2$、$y=(u-v)/2$，Jacobian 行列式为 $-1/2$，面积元取其绝对值。因此
\[
I=\frac12\int_0^{3}\int_0^{4}u^2dv\,du=18.
\]
负行列式表示坐标变换反转了定向，但通常的面积积分使用绝对值，不会因此变为负数。新变量同时把两组斜直线边界变成坐标平行的矩形。
\end{solve}

\clearpage
\section{三重积分与坐标变换}\label{sec:ch38}
三重积分（triple integral）可表示体积、质量和空间矩。柱坐标的体积元为 $r\,dr\,d\theta\,dz$；球坐标采用 $\rho$ 为径向距离、$\varphi$ 为与正 $z$ 轴夹角，此时体积元为 $\rho^2\sin\varphi\,d\rho\,d\varphi\,d\theta$。

\textbf{资源对应：}CLP-3 \S~3.5--\S~3.7。本章共 24 道编号习题。
\subsection{习题}
\begin{example}[长方体上的体积矩]\label{ex:c38q1}
设 $V=[0,1]\times[0,2]\times[0,3]$。求体积和 $\displaystyle\iiint_Vxyz\,dV$。
\par\textit{改编训练。}\par\hyperref[ans:c38q1]{解答见第~\pageref{ans:c38q1} 页。}
\end{example}
\begin{example}[长方体上的体积矩]\label{ex:c38q2}
设 $V=[0,2]\times[0,2]\times[0,1]$。求体积和 $\displaystyle\iiint_Vxyz\,dV$。
\par\textit{改编训练。}\par\hyperref[ans:c38q2]{解答见第~\pageref{ans:c38q2} 页。}
\end{example}
\begin{example}[长方体上的体积矩]\label{ex:c38q3}
设 $V=[0,2]\times[0,3]\times[0,4]$。求体积和 $\displaystyle\iiint_Vxyz\,dV$。
\par\textit{改编训练。}\par\hyperref[ans:c38q3]{解答见第~\pageref{ans:c38q3} 页。}
\end{example}
\begin{example}[长方体上的体积矩]\label{ex:c38q4}
设 $V=[0,3]\times[0,1]\times[0,2]$。求体积和 $\displaystyle\iiint_Vxyz\,dV$。
\par\textit{改编训练。}\par\hyperref[ans:c38q4]{解答见第~\pageref{ans:c38q4} 页。}
\end{example}
\begin{example}[长方体上的体积矩]\label{ex:c38q5}
设 $V=[0,4]\times[0,2]\times[0,3]$。求体积和 $\displaystyle\iiint_Vxyz\,dV$。
\par\textit{改编训练。}\par\hyperref[ans:c38q5]{解答见第~\pageref{ans:c38q5} 页。}
\end{example}
\begin{example}[长方体上的体积矩]\label{ex:c38q6}
设 $V=[0,3]\times[0,4]\times[0,2]$。求体积和 $\displaystyle\iiint_Vxyz\,dV$。
\par\textit{改编训练。}\par\hyperref[ans:c38q6]{解答见第~\pageref{ans:c38q6} 页。}
\end{example}
\begin{example}[四面体体积与形心]\label{ex:c38q7}
求 $V=\{x,y,z\ge0:x+y+z\le 1\}$ 的体积，以及均匀实体的形心。
\par\textit{改编训练。}\par\hyperref[ans:c38q7]{解答见第~\pageref{ans:c38q7} 页。}
\end{example}
\begin{example}[四面体体积与形心]\label{ex:c38q8}
求 $V=\{x,y,z\ge0:x+y+z\le 2\}$ 的体积，以及均匀实体的形心。
\par\textit{改编训练。}\par\hyperref[ans:c38q8]{解答见第~\pageref{ans:c38q8} 页。}
\end{example}
\begin{example}[四面体体积与形心]\label{ex:c38q9}
求 $V=\{x,y,z\ge0:x+y+z\le 3\}$ 的体积，以及均匀实体的形心。
\par\textit{改编训练。}\par\hyperref[ans:c38q9]{解答见第~\pageref{ans:c38q9} 页。}
\end{example}
\begin{example}[四面体体积与形心]\label{ex:c38q10}
求 $V=\{x,y,z\ge0:x+y+z\le 4\}$ 的体积，以及均匀实体的形心。
\par\textit{改编训练。}\par\hyperref[ans:c38q10]{解答见第~\pageref{ans:c38q10} 页。}
\end{example}
\begin{example}[四面体体积与形心]\label{ex:c38q11}
求 $V=\{x,y,z\ge0:x+y+z\le 5\}$ 的体积，以及均匀实体的形心。
\par\textit{改编训练。}\par\hyperref[ans:c38q11]{解答见第~\pageref{ans:c38q11} 页。}
\end{example}
\begin{example}[四面体体积与形心]\label{ex:c38q12}
求 $V=\{x,y,z\ge0:x+y+z\le 6\}$ 的体积，以及均匀实体的形心。
\par\textit{改编训练。}\par\hyperref[ans:c38q12]{解答见第~\pageref{ans:c38q12} 页。}
\end{example}
\begin{example}[球内的径向积分]\label{ex:c38q13}
设 $B:x^2+y^2+z^2\le 1$。计算 $\displaystyle\iiint_B(x^2+y^2+z^2)^{0}dV$。
\par\textit{改编训练。}\par\hyperref[ans:c38q13]{解答见第~\pageref{ans:c38q13} 页。}
\end{example}
\begin{example}[球内的径向积分]\label{ex:c38q14}
设 $B:x^2+y^2+z^2\le 4$。计算 $\displaystyle\iiint_B(x^2+y^2+z^2)^{0}dV$。
\par\textit{改编训练。}\par\hyperref[ans:c38q14]{解答见第~\pageref{ans:c38q14} 页。}
\end{example}
\begin{example}[球内的径向积分]\label{ex:c38q15}
设 $B:x^2+y^2+z^2\le 1$。计算 $\displaystyle\iiint_B(x^2+y^2+z^2)^{1}dV$。
\par\textit{改编训练。}\par\hyperref[ans:c38q15]{解答见第~\pageref{ans:c38q15} 页。}
\end{example}
\begin{example}[球内的径向积分]\label{ex:c38q16}
设 $B:x^2+y^2+z^2\le 4$。计算 $\displaystyle\iiint_B(x^2+y^2+z^2)^{1}dV$。
\par\textit{改编训练。}\par\hyperref[ans:c38q16]{解答见第~\pageref{ans:c38q16} 页。}
\end{example}
\begin{example}[球内的径向积分]\label{ex:c38q17}
设 $B:x^2+y^2+z^2\le 9$。计算 $\displaystyle\iiint_B(x^2+y^2+z^2)^{1}dV$。
\par\textit{改编训练。}\par\hyperref[ans:c38q17]{解答见第~\pageref{ans:c38q17} 页。}
\end{example}
\begin{example}[球内的径向积分]\label{ex:c38q18}
设 $B:x^2+y^2+z^2\le 1$。计算 $\displaystyle\iiint_B(x^2+y^2+z^2)^{2}dV$。
\par\textit{改编训练。}\par\hyperref[ans:c38q18]{解答见第~\pageref{ans:c38q18} 页。}
\end{example}
\begin{example}[圆柱体内的径向质量]\label{ex:c38q19}
设 $V:x^2+y^2\le 1$、$0\le z\le 2$，密度 $\delta=x^2+y^2$。求总质量和形心的 $z$ 坐标。
\par\textit{改编训练。}\par\hyperref[ans:c38q19]{解答见第~\pageref{ans:c38q19} 页。}
\end{example}
\begin{example}[圆柱体内的径向质量]\label{ex:c38q20}
设 $V:x^2+y^2\le 4$、$0\le z\le 3$，密度 $\delta=x^2+y^2$。求总质量和形心的 $z$ 坐标。
\par\textit{改编训练。}\par\hyperref[ans:c38q20]{解答见第~\pageref{ans:c38q20} 页。}
\end{example}
\begin{example}[圆柱体内的径向质量]\label{ex:c38q21}
设 $V:x^2+y^2\le 9$、$0\le z\le 1$，密度 $\delta=x^2+y^2$。求总质量和形心的 $z$ 坐标。
\par\textit{改编训练。}\par\hyperref[ans:c38q21]{解答见第~\pageref{ans:c38q21} 页。}
\end{example}
\begin{example}[圆柱体内的径向质量]\label{ex:c38q22}
设 $V:x^2+y^2\le 4$、$0\le z\le 5$，密度 $\delta=x^2+y^2$。求总质量和形心的 $z$ 坐标。
\par\textit{改编训练。}\par\hyperref[ans:c38q22]{解答见第~\pageref{ans:c38q22} 页。}
\end{example}
\begin{example}[圆柱体内的径向质量]\label{ex:c38q23}
设 $V:x^2+y^2\le 16$、$0\le z\le 3$，密度 $\delta=x^2+y^2$。求总质量和形心的 $z$ 坐标。
\par\textit{改编训练。}\par\hyperref[ans:c38q23]{解答见第~\pageref{ans:c38q23} 页。}
\end{example}
\begin{example}[圆柱体内的径向质量]\label{ex:c38q24}
设 $V:x^2+y^2\le 9$、$0\le z\le 5$，密度 $\delta=x^2+y^2$。求总质量和形心的 $z$ 坐标。
\par\textit{改编训练。}\par\hyperref[ans:c38q24]{解答见第~\pageref{ans:c38q24} 页。}
\end{example}
\subsection{解答}
\begin{solve}[习题~\ref{ex:c38q1}：长方体上的体积矩]\phantomsection\label{ans:c38q1}
长方体体积为 $6$。被积函数连续，且区域是三个区间的直积，可以分离积分：
\[
\iiint_Vxyz\,dV=\left(\int_0^{1}x\,dx\right)\left(\int_0^{2}y\,dy\right)\left(\int_0^{3}z\,dz\right)=\frac{9}{2}.
\]
如果把 $xyz$ 看成空间中的非负密度，则后一数值是总质量；若把它看成某个物理量，则体积平均值为 $\frac{3}{4}$。同一个积分的解释依赖被积函数的单位与含义，不能只凭形式决定。
\end{solve}

\begin{solve}[习题~\ref{ex:c38q2}：长方体上的体积矩]\phantomsection\label{ans:c38q2}
长方体体积为 $4$。被积函数连续，且区域是三个区间的直积，可以分离积分：
\[
\iiint_Vxyz\,dV=\left(\int_0^{2}x\,dx\right)\left(\int_0^{2}y\,dy\right)\left(\int_0^{1}z\,dz\right)=2.
\]
如果把 $xyz$ 看成空间中的非负密度，则后一数值是总质量；若把它看成某个物理量，则体积平均值为 $\frac{1}{2}$。同一个积分的解释依赖被积函数的单位与含义，不能只凭形式决定。
\end{solve}

\begin{solve}[习题~\ref{ex:c38q3}：长方体上的体积矩]\phantomsection\label{ans:c38q3}
长方体体积为 $24$。被积函数连续，且区域是三个区间的直积，可以分离积分：
\[
\iiint_Vxyz\,dV=\left(\int_0^{2}x\,dx\right)\left(\int_0^{3}y\,dy\right)\left(\int_0^{4}z\,dz\right)=72.
\]
如果把 $xyz$ 看成空间中的非负密度，则后一数值是总质量；若把它看成某个物理量，则体积平均值为 $3$。同一个积分的解释依赖被积函数的单位与含义，不能只凭形式决定。
\end{solve}

\begin{solve}[习题~\ref{ex:c38q4}：长方体上的体积矩]\phantomsection\label{ans:c38q4}
长方体体积为 $6$。被积函数连续，且区域是三个区间的直积，可以分离积分：
\[
\iiint_Vxyz\,dV=\left(\int_0^{3}x\,dx\right)\left(\int_0^{1}y\,dy\right)\left(\int_0^{2}z\,dz\right)=\frac{9}{2}.
\]
如果把 $xyz$ 看成空间中的非负密度，则后一数值是总质量；若把它看成某个物理量，则体积平均值为 $\frac{3}{4}$。同一个积分的解释依赖被积函数的单位与含义，不能只凭形式决定。
\end{solve}

\begin{solve}[习题~\ref{ex:c38q5}：长方体上的体积矩]\phantomsection\label{ans:c38q5}
长方体体积为 $24$。被积函数连续，且区域是三个区间的直积，可以分离积分：
\[
\iiint_Vxyz\,dV=\left(\int_0^{4}x\,dx\right)\left(\int_0^{2}y\,dy\right)\left(\int_0^{3}z\,dz\right)=72.
\]
如果把 $xyz$ 看成空间中的非负密度，则后一数值是总质量；若把它看成某个物理量，则体积平均值为 $3$。同一个积分的解释依赖被积函数的单位与含义，不能只凭形式决定。
\end{solve}

\begin{solve}[习题~\ref{ex:c38q6}：长方体上的体积矩]\phantomsection\label{ans:c38q6}
长方体体积为 $24$。被积函数连续，且区域是三个区间的直积，可以分离积分：
\[
\iiint_Vxyz\,dV=\left(\int_0^{3}x\,dx\right)\left(\int_0^{4}y\,dy\right)\left(\int_0^{2}z\,dz\right)=72.
\]
如果把 $xyz$ 看成空间中的非负密度，则后一数值是总质量；若把它看成某个物理量，则体积平均值为 $3$。同一个积分的解释依赖被积函数的单位与含义，不能只凭形式决定。
\end{solve}

\begin{solve}[习题~\ref{ex:c38q7}：四面体体积与形心]\phantomsection\label{ans:c38q7}
先按 $z,y,x$ 的顺序积分：
\[
|V|=\int_0^{1}\int_0^{1-x}\int_0^{1-x-y}dz\,dy\,dx=\frac{1}{6}.
\]
横坐标的一阶矩为
\[
\iiint_Vx\,dV=\frac12\int_0^{1}x(1-x)^2dx=\frac{1}{24}.
\]
除以体积得 $\bar x=\frac{1}{4}$。区域对三个坐标置换对称，所以 $\bar y=\bar z=\frac{1}{4}$。形心不是三个截距点的坐标平均，而是整个均匀实体的一阶矩除以体积。
\end{solve}

\begin{solve}[习题~\ref{ex:c38q8}：四面体体积与形心]\phantomsection\label{ans:c38q8}
先按 $z,y,x$ 的顺序积分：
\[
|V|=\int_0^{2}\int_0^{2-x}\int_0^{2-x-y}dz\,dy\,dx=\frac{4}{3}.
\]
横坐标的一阶矩为
\[
\iiint_Vx\,dV=\frac12\int_0^{2}x(2-x)^2dx=\frac{2}{3}.
\]
除以体积得 $\bar x=\frac{1}{2}$。区域对三个坐标置换对称，所以 $\bar y=\bar z=\frac{1}{2}$。形心不是三个截距点的坐标平均，而是整个均匀实体的一阶矩除以体积。
\end{solve}

\begin{solve}[习题~\ref{ex:c38q9}：四面体体积与形心]\phantomsection\label{ans:c38q9}
先按 $z,y,x$ 的顺序积分：
\[
|V|=\int_0^{3}\int_0^{3-x}\int_0^{3-x-y}dz\,dy\,dx=\frac{9}{2}.
\]
横坐标的一阶矩为
\[
\iiint_Vx\,dV=\frac12\int_0^{3}x(3-x)^2dx=\frac{27}{8}.
\]
除以体积得 $\bar x=\frac{3}{4}$。区域对三个坐标置换对称，所以 $\bar y=\bar z=\frac{3}{4}$。形心不是三个截距点的坐标平均，而是整个均匀实体的一阶矩除以体积。
\end{solve}

\begin{solve}[习题~\ref{ex:c38q10}：四面体体积与形心]\phantomsection\label{ans:c38q10}
先按 $z,y,x$ 的顺序积分：
\[
|V|=\int_0^{4}\int_0^{4-x}\int_0^{4-x-y}dz\,dy\,dx=\frac{32}{3}.
\]
横坐标的一阶矩为
\[
\iiint_Vx\,dV=\frac12\int_0^{4}x(4-x)^2dx=\frac{32}{3}.
\]
除以体积得 $\bar x=1$。区域对三个坐标置换对称，所以 $\bar y=\bar z=1$。形心不是三个截距点的坐标平均，而是整个均匀实体的一阶矩除以体积。
\end{solve}

\begin{solve}[习题~\ref{ex:c38q11}：四面体体积与形心]\phantomsection\label{ans:c38q11}
先按 $z,y,x$ 的顺序积分：
\[
|V|=\int_0^{5}\int_0^{5-x}\int_0^{5-x-y}dz\,dy\,dx=\frac{125}{6}.
\]
横坐标的一阶矩为
\[
\iiint_Vx\,dV=\frac12\int_0^{5}x(5-x)^2dx=\frac{625}{24}.
\]
除以体积得 $\bar x=\frac{5}{4}$。区域对三个坐标置换对称，所以 $\bar y=\bar z=\frac{5}{4}$。形心不是三个截距点的坐标平均，而是整个均匀实体的一阶矩除以体积。
\end{solve}

\begin{solve}[习题~\ref{ex:c38q12}：四面体体积与形心]\phantomsection\label{ans:c38q12}
先按 $z,y,x$ 的顺序积分：
\[
|V|=\int_0^{6}\int_0^{6-x}\int_0^{6-x-y}dz\,dy\,dx=36.
\]
横坐标的一阶矩为
\[
\iiint_Vx\,dV=\frac12\int_0^{6}x(6-x)^2dx=54.
\]
除以体积得 $\bar x=\frac{3}{2}$。区域对三个坐标置换对称，所以 $\bar y=\bar z=\frac{3}{2}$。形心不是三个截距点的坐标平均，而是整个均匀实体的一阶矩除以体积。
\end{solve}

\begin{solve}[习题~\ref{ex:c38q13}：球内的径向积分]\phantomsection\label{ans:c38q13}
球坐标下被积函数为 $\rho^{0}$，且 $0\le\rho\le 1$、$0\le\varphi\le\pi$、$0\le\theta\le2\pi$。因此
\[
I=\int_0^{2\pi}d\theta\int_0^\pi\sin\varphi\,d\varphi\int_0^{1}\rho^{2}d\rho=\frac{4}{3}\pi.
\]
当幂指数为0时，这一公式退化为球体积；一般情形可以看成径向密度产生的总质量。$\rho$ 的积分幂次需要加2，正是球坐标体积元带来的贡献。
\end{solve}

\begin{solve}[习题~\ref{ex:c38q14}：球内的径向积分]\phantomsection\label{ans:c38q14}
球坐标下被积函数为 $\rho^{0}$，且 $0\le\rho\le 2$、$0\le\varphi\le\pi$、$0\le\theta\le2\pi$。因此
\[
I=\int_0^{2\pi}d\theta\int_0^\pi\sin\varphi\,d\varphi\int_0^{2}\rho^{2}d\rho=\frac{32}{3}\pi.
\]
当幂指数为0时，这一公式退化为球体积；一般情形可以看成径向密度产生的总质量。$\rho$ 的积分幂次需要加2，正是球坐标体积元带来的贡献。
\end{solve}

\begin{solve}[习题~\ref{ex:c38q15}：球内的径向积分]\phantomsection\label{ans:c38q15}
球坐标下被积函数为 $\rho^{2}$，且 $0\le\rho\le 1$、$0\le\varphi\le\pi$、$0\le\theta\le2\pi$。因此
\[
I=\int_0^{2\pi}d\theta\int_0^\pi\sin\varphi\,d\varphi\int_0^{1}\rho^{4}d\rho=\frac{4}{5}\pi.
\]
当幂指数为0时，这一公式退化为球体积；一般情形可以看成径向密度产生的总质量。$\rho$ 的积分幂次需要加2，正是球坐标体积元带来的贡献。
\end{solve}

\begin{solve}[习题~\ref{ex:c38q16}：球内的径向积分]\phantomsection\label{ans:c38q16}
球坐标下被积函数为 $\rho^{2}$，且 $0\le\rho\le 2$、$0\le\varphi\le\pi$、$0\le\theta\le2\pi$。因此
\[
I=\int_0^{2\pi}d\theta\int_0^\pi\sin\varphi\,d\varphi\int_0^{2}\rho^{4}d\rho=\frac{128}{5}\pi.
\]
当幂指数为0时，这一公式退化为球体积；一般情形可以看成径向密度产生的总质量。$\rho$ 的积分幂次需要加2，正是球坐标体积元带来的贡献。
\end{solve}

\begin{solve}[习题~\ref{ex:c38q17}：球内的径向积分]\phantomsection\label{ans:c38q17}
球坐标下被积函数为 $\rho^{2}$，且 $0\le\rho\le 3$、$0\le\varphi\le\pi$、$0\le\theta\le2\pi$。因此
\[
I=\int_0^{2\pi}d\theta\int_0^\pi\sin\varphi\,d\varphi\int_0^{3}\rho^{4}d\rho=\frac{972}{5}\pi.
\]
当幂指数为0时，这一公式退化为球体积；一般情形可以看成径向密度产生的总质量。$\rho$ 的积分幂次需要加2，正是球坐标体积元带来的贡献。
\end{solve}

\begin{solve}[习题~\ref{ex:c38q18}：球内的径向积分]\phantomsection\label{ans:c38q18}
球坐标下被积函数为 $\rho^{4}$，且 $0\le\rho\le 1$、$0\le\varphi\le\pi$、$0\le\theta\le2\pi$。因此
\[
I=\int_0^{2\pi}d\theta\int_0^\pi\sin\varphi\,d\varphi\int_0^{1}\rho^{6}d\rho=\frac{4}{7}\pi.
\]
当幂指数为0时，这一公式退化为球体积；一般情形可以看成径向密度产生的总质量。$\rho$ 的积分幂次需要加2，正是球坐标体积元带来的贡献。
\end{solve}

\begin{solve}[习题~\ref{ex:c38q19}：圆柱体内的径向质量]\phantomsection\label{ans:c38q19}
柱坐标下密度为 $r^2$，体积元为 $r\,dr\,d\theta\,dz$，所以
\[
M=\int_0^{2}dz\int_0^{2\pi}d\theta\int_0^{1}r^3dr=1\pi.
\]
关于 $xy$ 平面的一阶矩为
\[
M_{xy}=\int_0^{2}z\,dz\int_0^{2\pi}d\theta\int_0^{1}r^3dr=1\pi.
\]
因此 $\bar z=M_{xy}/M=1$。由于径向对称，$\bar x=\bar y=0$。密度并非常数，但与 $z$ 无关，所以沿高度的质量分布仍均匀。
\end{solve}

\begin{solve}[习题~\ref{ex:c38q20}：圆柱体内的径向质量]\phantomsection\label{ans:c38q20}
柱坐标下密度为 $r^2$，体积元为 $r\,dr\,d\theta\,dz$，所以
\[
M=\int_0^{3}dz\int_0^{2\pi}d\theta\int_0^{2}r^3dr=24\pi.
\]
关于 $xy$ 平面的一阶矩为
\[
M_{xy}=\int_0^{3}z\,dz\int_0^{2\pi}d\theta\int_0^{2}r^3dr=36\pi.
\]
因此 $\bar z=M_{xy}/M=\frac{3}{2}$。由于径向对称，$\bar x=\bar y=0$。密度并非常数，但与 $z$ 无关，所以沿高度的质量分布仍均匀。
\end{solve}

\begin{solve}[习题~\ref{ex:c38q21}：圆柱体内的径向质量]\phantomsection\label{ans:c38q21}
柱坐标下密度为 $r^2$，体积元为 $r\,dr\,d\theta\,dz$，所以
\[
M=\int_0^{1}dz\int_0^{2\pi}d\theta\int_0^{3}r^3dr=\frac{81}{2}\pi.
\]
关于 $xy$ 平面的一阶矩为
\[
M_{xy}=\int_0^{1}z\,dz\int_0^{2\pi}d\theta\int_0^{3}r^3dr=\frac{81}{4}\pi.
\]
因此 $\bar z=M_{xy}/M=\frac{1}{2}$。由于径向对称，$\bar x=\bar y=0$。密度并非常数，但与 $z$ 无关，所以沿高度的质量分布仍均匀。
\end{solve}

\begin{solve}[习题~\ref{ex:c38q22}：圆柱体内的径向质量]\phantomsection\label{ans:c38q22}
柱坐标下密度为 $r^2$，体积元为 $r\,dr\,d\theta\,dz$，所以
\[
M=\int_0^{5}dz\int_0^{2\pi}d\theta\int_0^{2}r^3dr=40\pi.
\]
关于 $xy$ 平面的一阶矩为
\[
M_{xy}=\int_0^{5}z\,dz\int_0^{2\pi}d\theta\int_0^{2}r^3dr=100\pi.
\]
因此 $\bar z=M_{xy}/M=\frac{5}{2}$。由于径向对称，$\bar x=\bar y=0$。密度并非常数，但与 $z$ 无关，所以沿高度的质量分布仍均匀。
\end{solve}

\begin{solve}[习题~\ref{ex:c38q23}：圆柱体内的径向质量]\phantomsection\label{ans:c38q23}
柱坐标下密度为 $r^2$，体积元为 $r\,dr\,d\theta\,dz$，所以
\[
M=\int_0^{3}dz\int_0^{2\pi}d\theta\int_0^{4}r^3dr=384\pi.
\]
关于 $xy$ 平面的一阶矩为
\[
M_{xy}=\int_0^{3}z\,dz\int_0^{2\pi}d\theta\int_0^{4}r^3dr=576\pi.
\]
因此 $\bar z=M_{xy}/M=\frac{3}{2}$。由于径向对称，$\bar x=\bar y=0$。密度并非常数，但与 $z$ 无关，所以沿高度的质量分布仍均匀。
\end{solve}

\begin{solve}[习题~\ref{ex:c38q24}：圆柱体内的径向质量]\phantomsection\label{ans:c38q24}
柱坐标下密度为 $r^2$，体积元为 $r\,dr\,d\theta\,dz$，所以
\[
M=\int_0^{5}dz\int_0^{2\pi}d\theta\int_0^{3}r^3dr=\frac{405}{2}\pi.
\]
关于 $xy$ 平面的一阶矩为
\[
M_{xy}=\int_0^{5}z\,dz\int_0^{2\pi}d\theta\int_0^{3}r^3dr=\frac{2025}{4}\pi.
\]
因此 $\bar z=M_{xy}/M=\frac{5}{2}$。由于径向对称，$\bar x=\bar y=0$。密度并非常数，但与 $z$ 无关，所以沿高度的质量分布仍均匀。
\end{solve}

\clearpage
\section{空间积分的几何综合}\label{sec:ch39}
选取坐标应当适应区域的几何结构：圆锥适合柱坐标，球与椭球适合球坐标及线性缩放，单纯形适合按平面边界逐层积分。对称性可以减少计算，但需要同时检查密度是否对称。

\textbf{资源对应：}CLP-3 \S~3.5--\S~3.7 的扩展训练。本章共 24 道编号习题。
\subsection{习题}
\begin{example}[圆锥的分层积分]\label{ex:c39q1}
设 $V:0\le z\le 1, x^2+y^2\le z^2$。计算 $\iiint_Vz^{0}dV$。
\par\textit{改编训练。}\par\hyperref[ans:c39q1]{解答见第~\pageref{ans:c39q1} 页。}
\end{example}
\begin{example}[圆锥的分层积分]\label{ex:c39q2}
设 $V:0\le z\le 2, x^2+y^2\le z^2$。计算 $\iiint_Vz^{1}dV$。
\par\textit{改编训练。}\par\hyperref[ans:c39q2]{解答见第~\pageref{ans:c39q2} 页。}
\end{example}
\begin{example}[圆锥的分层积分]\label{ex:c39q3}
设 $V:0\le z\le 3, x^2+y^2\le z^2$。计算 $\iiint_Vz^{0}dV$。
\par\textit{改编训练。}\par\hyperref[ans:c39q3]{解答见第~\pageref{ans:c39q3} 页。}
\end{example}
\begin{example}[圆锥的分层积分]\label{ex:c39q4}
设 $V:0\le z\le 2, x^2+y^2\le z^2$。计算 $\iiint_Vz^{2}dV$。
\par\textit{改编训练。}\par\hyperref[ans:c39q4]{解答见第~\pageref{ans:c39q4} 页。}
\end{example}
\begin{example}[圆锥的分层积分]\label{ex:c39q5}
设 $V:0\le z\le 3, x^2+y^2\le z^2$。计算 $\iiint_Vz^{1}dV$。
\par\textit{改编训练。}\par\hyperref[ans:c39q5]{解答见第~\pageref{ans:c39q5} 页。}
\end{example}
\begin{example}[圆锥的分层积分]\label{ex:c39q6}
设 $V:0\le z\le 1, x^2+y^2\le z^2$。计算 $\iiint_Vz^{3}dV$。
\par\textit{改编训练。}\par\hyperref[ans:c39q6]{解答见第~\pageref{ans:c39q6} 页。}
\end{example}
\begin{example}[球内坐标平方的积分]\label{ex:c39q7}
设 $B:x^2+y^2+z^2\le 1$。求 $\iiint_Bz^2dV$，进而求均匀球体绕 $z$ 轴的转动惯量。
\par\textit{改编训练。}\par\hyperref[ans:c39q7]{解答见第~\pageref{ans:c39q7} 页。}
\end{example}
\begin{example}[球内坐标平方的积分]\label{ex:c39q8}
设 $B:x^2+y^2+z^2\le 4$。求 $\iiint_Bz^2dV$，进而求均匀球体绕 $z$ 轴的转动惯量。
\par\textit{改编训练。}\par\hyperref[ans:c39q8]{解答见第~\pageref{ans:c39q8} 页。}
\end{example}
\begin{example}[球内坐标平方的积分]\label{ex:c39q9}
设 $B:x^2+y^2+z^2\le 9$。求 $\iiint_Bz^2dV$，进而求均匀球体绕 $z$ 轴的转动惯量。
\par\textit{改编训练。}\par\hyperref[ans:c39q9]{解答见第~\pageref{ans:c39q9} 页。}
\end{example}
\begin{example}[球内坐标平方的积分]\label{ex:c39q10}
设 $B:x^2+y^2+z^2\le 16$。求 $\iiint_Bz^2dV$，进而求均匀球体绕 $z$ 轴的转动惯量。
\par\textit{改编训练。}\par\hyperref[ans:c39q10]{解答见第~\pageref{ans:c39q10} 页。}
\end{example}
\begin{example}[球内坐标平方的积分]\label{ex:c39q11}
设 $B:x^2+y^2+z^2\le 25$。求 $\iiint_Bz^2dV$，进而求均匀球体绕 $z$ 轴的转动惯量。
\par\textit{改编训练。}\par\hyperref[ans:c39q11]{解答见第~\pageref{ans:c39q11} 页。}
\end{example}
\begin{example}[球内坐标平方的积分]\label{ex:c39q12}
设 $B:x^2+y^2+z^2\le 36$。求 $\iiint_Bz^2dV$，进而求均匀球体绕 $z$ 轴的转动惯量。
\par\textit{改编训练。}\par\hyperref[ans:c39q12]{解答见第~\pageref{ans:c39q12} 页。}
\end{example}
\begin{example}[椭球体积与二次矩]\label{ex:c39q13}
设 $E:x^2/1+y^2/4+z^2/9\le1$。求体积和 $\iiint_Ex^2dV$。
\par\textit{改编训练。}\par\hyperref[ans:c39q13]{解答见第~\pageref{ans:c39q13} 页。}
\end{example}
\begin{example}[椭球体积与二次矩]\label{ex:c39q14}
设 $E:x^2/4+y^2/4+z^2/1\le1$。求体积和 $\iiint_Ex^2dV$。
\par\textit{改编训练。}\par\hyperref[ans:c39q14]{解答见第~\pageref{ans:c39q14} 页。}
\end{example}
\begin{example}[椭球体积与二次矩]\label{ex:c39q15}
设 $E:x^2/4+y^2/9+z^2/16\le1$。求体积和 $\iiint_Ex^2dV$。
\par\textit{改编训练。}\par\hyperref[ans:c39q15]{解答见第~\pageref{ans:c39q15} 页。}
\end{example}
\begin{example}[椭球体积与二次矩]\label{ex:c39q16}
设 $E:x^2/9+y^2/1+z^2/4\le1$。求体积和 $\iiint_Ex^2dV$。
\par\textit{改编训练。}\par\hyperref[ans:c39q16]{解答见第~\pageref{ans:c39q16} 页。}
\end{example}
\begin{example}[椭球体积与二次矩]\label{ex:c39q17}
设 $E:x^2/16+y^2/4+z^2/9\le1$。求体积和 $\iiint_Ex^2dV$。
\par\textit{改编训练。}\par\hyperref[ans:c39q17]{解答见第~\pageref{ans:c39q17} 页。}
\end{example}
\begin{example}[椭球体积与二次矩]\label{ex:c39q18}
设 $E:x^2/9+y^2/16+z^2/4\le1$。求体积和 $\iiint_Ex^2dV$。
\par\textit{改编训练。}\par\hyperref[ans:c39q18]{解答见第~\pageref{ans:c39q18} 页。}
\end{example}
\begin{example}[单纯形上的乘积密度]\label{ex:c39q19}
在 $V:x,y,z\ge0, x+y+z\le 1$ 上，密度为 $xyz$。求总质量。
\par\textit{改编训练。}\par\hyperref[ans:c39q19]{解答见第~\pageref{ans:c39q19} 页。}
\end{example}
\begin{example}[单纯形上的乘积密度]\label{ex:c39q20}
在 $V:x,y,z\ge0, x+y+z\le 2$ 上，密度为 $xyz$。求总质量。
\par\textit{改编训练。}\par\hyperref[ans:c39q20]{解答见第~\pageref{ans:c39q20} 页。}
\end{example}
\begin{example}[单纯形上的乘积密度]\label{ex:c39q21}
在 $V:x,y,z\ge0, x+y+z\le 3$ 上，密度为 $xyz$。求总质量。
\par\textit{改编训练。}\par\hyperref[ans:c39q21]{解答见第~\pageref{ans:c39q21} 页。}
\end{example}
\begin{example}[单纯形上的乘积密度]\label{ex:c39q22}
在 $V:x,y,z\ge0, x+y+z\le 4$ 上，密度为 $xyz$。求总质量。
\par\textit{改编训练。}\par\hyperref[ans:c39q22]{解答见第~\pageref{ans:c39q22} 页。}
\end{example}
\begin{example}[单纯形上的乘积密度]\label{ex:c39q23}
在 $V:x,y,z\ge0, x+y+z\le 5$ 上，密度为 $xyz$。求总质量。
\par\textit{改编训练。}\par\hyperref[ans:c39q23]{解答见第~\pageref{ans:c39q23} 页。}
\end{example}
\begin{example}[单纯形上的乘积密度]\label{ex:c39q24}
在 $V:x,y,z\ge0, x+y+z\le 6$ 上，密度为 $xyz$。求总质量。
\par\textit{改编训练。}\par\hyperref[ans:c39q24]{解答见第~\pageref{ans:c39q24} 页。}
\end{example}
\subsection{解答}
\begin{solve}[习题~\ref{ex:c39q1}：圆锥的分层积分]\phantomsection\label{ans:c39q1}
在固定高度 $z$ 处，横截面是半径 $z$ 的圆盘，面积为 $\pi z^2$。所以
\[
I=\pi\int_0^{1}z^{2}dz=\frac{1}{3}\pi.
\]
柱坐标也给出 $I=\int_0^{1}\int_0^{2\pi}\int_0^zr z^{0}drd\theta dz$，与分层解释一致。积分指数由密度的次数和横截面积的二次增长共同决定，不能把圆锥当成截面面积固定的圆柱。
\end{solve}

\begin{solve}[习题~\ref{ex:c39q2}：圆锥的分层积分]\phantomsection\label{ans:c39q2}
在固定高度 $z$ 处，横截面是半径 $z$ 的圆盘，面积为 $\pi z^2$。所以
\[
I=\pi\int_0^{2}z^{3}dz=4\pi.
\]
柱坐标也给出 $I=\int_0^{2}\int_0^{2\pi}\int_0^zr z^{1}drd\theta dz$，与分层解释一致。积分指数由密度的次数和横截面积的二次增长共同决定，不能把圆锥当成截面面积固定的圆柱。
\end{solve}

\begin{solve}[习题~\ref{ex:c39q3}：圆锥的分层积分]\phantomsection\label{ans:c39q3}
在固定高度 $z$ 处，横截面是半径 $z$ 的圆盘，面积为 $\pi z^2$。所以
\[
I=\pi\int_0^{3}z^{2}dz=9\pi.
\]
柱坐标也给出 $I=\int_0^{3}\int_0^{2\pi}\int_0^zr z^{0}drd\theta dz$，与分层解释一致。积分指数由密度的次数和横截面积的二次增长共同决定，不能把圆锥当成截面面积固定的圆柱。
\end{solve}

\begin{solve}[习题~\ref{ex:c39q4}：圆锥的分层积分]\phantomsection\label{ans:c39q4}
在固定高度 $z$ 处，横截面是半径 $z$ 的圆盘，面积为 $\pi z^2$。所以
\[
I=\pi\int_0^{2}z^{4}dz=\frac{32}{5}\pi.
\]
柱坐标也给出 $I=\int_0^{2}\int_0^{2\pi}\int_0^zr z^{2}drd\theta dz$，与分层解释一致。积分指数由密度的次数和横截面积的二次增长共同决定，不能把圆锥当成截面面积固定的圆柱。
\end{solve}

\begin{solve}[习题~\ref{ex:c39q5}：圆锥的分层积分]\phantomsection\label{ans:c39q5}
在固定高度 $z$ 处，横截面是半径 $z$ 的圆盘，面积为 $\pi z^2$。所以
\[
I=\pi\int_0^{3}z^{3}dz=\frac{81}{4}\pi.
\]
柱坐标也给出 $I=\int_0^{3}\int_0^{2\pi}\int_0^zr z^{1}drd\theta dz$，与分层解释一致。积分指数由密度的次数和横截面积的二次增长共同决定，不能把圆锥当成截面面积固定的圆柱。
\end{solve}

\begin{solve}[习题~\ref{ex:c39q6}：圆锥的分层积分]\phantomsection\label{ans:c39q6}
在固定高度 $z$ 处，横截面是半径 $z$ 的圆盘，面积为 $\pi z^2$。所以
\[
I=\pi\int_0^{1}z^{5}dz=\frac{1}{6}\pi.
\]
柱坐标也给出 $I=\int_0^{1}\int_0^{2\pi}\int_0^zr z^{3}drd\theta dz$，与分层解释一致。积分指数由密度的次数和横截面积的二次增长共同决定，不能把圆锥当成截面面积固定的圆柱。
\end{solve}

\begin{solve}[习题~\ref{ex:c39q7}：球内坐标平方的积分]\phantomsection\label{ans:c39q7}
由对称性，三个坐标平方的体积积分相等。它们的和是球半径变量的平方，因此
\[
3\iiint_Bz^2dV=4\pi\int_0^{1}\rho^4d\rho=\frac{4}{5}\pi.
\]
所以 $\iiint_Bz^2dV=\frac{4}{15}\pi$。密度为1时，绕 $z$ 轴的转动惯量为
\[
I_z=\iiint_B(x^2+y^2)dV=\frac{8}{15}\pi.
\]
球质量为 $\frac{4}{3}\pi$，因此也可写成 $I_z=\frac25M1$，符合惯量为质量乘长度平方的量纲。
\end{solve}

\begin{solve}[习题~\ref{ex:c39q8}：球内坐标平方的积分]\phantomsection\label{ans:c39q8}
由对称性，三个坐标平方的体积积分相等。它们的和是球半径变量的平方，因此
\[
3\iiint_Bz^2dV=4\pi\int_0^{2}\rho^4d\rho=\frac{128}{5}\pi.
\]
所以 $\iiint_Bz^2dV=\frac{128}{15}\pi$。密度为1时，绕 $z$ 轴的转动惯量为
\[
I_z=\iiint_B(x^2+y^2)dV=\frac{256}{15}\pi.
\]
球质量为 $\frac{32}{3}\pi$，因此也可写成 $I_z=\frac25M4$，符合惯量为质量乘长度平方的量纲。
\end{solve}

\begin{solve}[习题~\ref{ex:c39q9}：球内坐标平方的积分]\phantomsection\label{ans:c39q9}
由对称性，三个坐标平方的体积积分相等。它们的和是球半径变量的平方，因此
\[
3\iiint_Bz^2dV=4\pi\int_0^{3}\rho^4d\rho=\frac{972}{5}\pi.
\]
所以 $\iiint_Bz^2dV=\frac{324}{5}\pi$。密度为1时，绕 $z$ 轴的转动惯量为
\[
I_z=\iiint_B(x^2+y^2)dV=\frac{648}{5}\pi.
\]
球质量为 $36\pi$，因此也可写成 $I_z=\frac25M9$，符合惯量为质量乘长度平方的量纲。
\end{solve}

\begin{solve}[习题~\ref{ex:c39q10}：球内坐标平方的积分]\phantomsection\label{ans:c39q10}
由对称性，三个坐标平方的体积积分相等。它们的和是球半径变量的平方，因此
\[
3\iiint_Bz^2dV=4\pi\int_0^{4}\rho^4d\rho=\frac{4096}{5}\pi.
\]
所以 $\iiint_Bz^2dV=\frac{4096}{15}\pi$。密度为1时，绕 $z$ 轴的转动惯量为
\[
I_z=\iiint_B(x^2+y^2)dV=\frac{8192}{15}\pi.
\]
球质量为 $\frac{256}{3}\pi$，因此也可写成 $I_z=\frac25M16$，符合惯量为质量乘长度平方的量纲。
\end{solve}

\begin{solve}[习题~\ref{ex:c39q11}：球内坐标平方的积分]\phantomsection\label{ans:c39q11}
由对称性，三个坐标平方的体积积分相等。它们的和是球半径变量的平方，因此
\[
3\iiint_Bz^2dV=4\pi\int_0^{5}\rho^4d\rho=2500\pi.
\]
所以 $\iiint_Bz^2dV=\frac{2500}{3}\pi$。密度为1时，绕 $z$ 轴的转动惯量为
\[
I_z=\iiint_B(x^2+y^2)dV=\frac{5000}{3}\pi.
\]
球质量为 $\frac{500}{3}\pi$，因此也可写成 $I_z=\frac25M25$，符合惯量为质量乘长度平方的量纲。
\end{solve}

\begin{solve}[习题~\ref{ex:c39q12}：球内坐标平方的积分]\phantomsection\label{ans:c39q12}
由对称性，三个坐标平方的体积积分相等。它们的和是球半径变量的平方，因此
\[
3\iiint_Bz^2dV=4\pi\int_0^{6}\rho^4d\rho=\frac{31104}{5}\pi.
\]
所以 $\iiint_Bz^2dV=\frac{10368}{5}\pi$。密度为1时，绕 $z$ 轴的转动惯量为
\[
I_z=\iiint_B(x^2+y^2)dV=\frac{20736}{5}\pi.
\]
球质量为 $288\pi$，因此也可写成 $I_z=\frac25M36$，符合惯量为质量乘长度平方的量纲。
\end{solve}

\begin{solve}[习题~\ref{ex:c39q13}：椭球体积与二次矩]\phantomsection\label{ans:c39q13}
作线性缩放 $x=1u$、$y=2v$、$z=3w$，Jacobian 为 $6$，新区域是单位球。故
\[
|E|=8\pi.
\]
对于二次矩，被积函数还带来因子 $1$，于是
\[
\iiint_Ex^2dV=6\iiint_{u^2+v^2+w^2\le1}u^2dudvdw=\frac{8}{5}\pi.
\]
换元的体积缩放与被积函数的坐标缩放是两件事，二者都要计入。
\end{solve}

\begin{solve}[习题~\ref{ex:c39q14}：椭球体积与二次矩]\phantomsection\label{ans:c39q14}
作线性缩放 $x=2u$、$y=2v$、$z=1w$，Jacobian 为 $4$，新区域是单位球。故
\[
|E|=\frac{16}{3}\pi.
\]
对于二次矩，被积函数还带来因子 $4$，于是
\[
\iiint_Ex^2dV=16\iiint_{u^2+v^2+w^2\le1}u^2dudvdw=\frac{64}{15}\pi.
\]
换元的体积缩放与被积函数的坐标缩放是两件事，二者都要计入。
\end{solve}

\begin{solve}[习题~\ref{ex:c39q15}：椭球体积与二次矩]\phantomsection\label{ans:c39q15}
作线性缩放 $x=2u$、$y=3v$、$z=4w$，Jacobian 为 $24$，新区域是单位球。故
\[
|E|=32\pi.
\]
对于二次矩，被积函数还带来因子 $4$，于是
\[
\iiint_Ex^2dV=96\iiint_{u^2+v^2+w^2\le1}u^2dudvdw=\frac{128}{5}\pi.
\]
换元的体积缩放与被积函数的坐标缩放是两件事，二者都要计入。
\end{solve}

\begin{solve}[习题~\ref{ex:c39q16}：椭球体积与二次矩]\phantomsection\label{ans:c39q16}
作线性缩放 $x=3u$、$y=1v$、$z=2w$，Jacobian 为 $6$，新区域是单位球。故
\[
|E|=8\pi.
\]
对于二次矩，被积函数还带来因子 $9$，于是
\[
\iiint_Ex^2dV=54\iiint_{u^2+v^2+w^2\le1}u^2dudvdw=\frac{72}{5}\pi.
\]
换元的体积缩放与被积函数的坐标缩放是两件事，二者都要计入。
\end{solve}

\begin{solve}[习题~\ref{ex:c39q17}：椭球体积与二次矩]\phantomsection\label{ans:c39q17}
作线性缩放 $x=4u$、$y=2v$、$z=3w$，Jacobian 为 $24$，新区域是单位球。故
\[
|E|=32\pi.
\]
对于二次矩，被积函数还带来因子 $16$，于是
\[
\iiint_Ex^2dV=384\iiint_{u^2+v^2+w^2\le1}u^2dudvdw=\frac{512}{5}\pi.
\]
换元的体积缩放与被积函数的坐标缩放是两件事，二者都要计入。
\end{solve}

\begin{solve}[习题~\ref{ex:c39q18}：椭球体积与二次矩]\phantomsection\label{ans:c39q18}
作线性缩放 $x=3u$、$y=4v$、$z=2w$，Jacobian 为 $24$，新区域是单位球。故
\[
|E|=32\pi.
\]
对于二次矩，被积函数还带来因子 $9$，于是
\[
\iiint_Ex^2dV=216\iiint_{u^2+v^2+w^2\le1}u^2dudvdw=\frac{288}{5}\pi.
\]
换元的体积缩放与被积函数的坐标缩放是两件事，二者都要计入。
\end{solve}

\begin{solve}[习题~\ref{ex:c39q19}：单纯形上的乘积密度]\phantomsection\label{ans:c39q19}
先对 $z$ 积分，再对 $y$ 积分：
\[
M=\frac12\int_0^{1}\int_0^{1-x}xy(1-x-y)^2dy\,dx.
\]
对内层作 $y=(a-x)t$ 的缩放，其中 $\int_0^1t(1-t)^2dt=1/12$，所以
\[
M=\frac1{24}\int_0^{1}x(1-x)^4dx=\frac{1}{720}.
\]
最后一步可以展开四次幂或再次使用有限区间 Beta 型积分。密度在区域边界上可能为零，但在内部为正，总质量为正且随尺度的六次方变化。
\end{solve}

\begin{solve}[习题~\ref{ex:c39q20}：单纯形上的乘积密度]\phantomsection\label{ans:c39q20}
先对 $z$ 积分，再对 $y$ 积分：
\[
M=\frac12\int_0^{2}\int_0^{2-x}xy(2-x-y)^2dy\,dx.
\]
对内层作 $y=(a-x)t$ 的缩放，其中 $\int_0^1t(1-t)^2dt=1/12$，所以
\[
M=\frac1{24}\int_0^{2}x(2-x)^4dx=\frac{4}{45}.
\]
最后一步可以展开四次幂或再次使用有限区间 Beta 型积分。密度在区域边界上可能为零，但在内部为正，总质量为正且随尺度的六次方变化。
\end{solve}

\begin{solve}[习题~\ref{ex:c39q21}：单纯形上的乘积密度]\phantomsection\label{ans:c39q21}
先对 $z$ 积分，再对 $y$ 积分：
\[
M=\frac12\int_0^{3}\int_0^{3-x}xy(3-x-y)^2dy\,dx.
\]
对内层作 $y=(a-x)t$ 的缩放，其中 $\int_0^1t(1-t)^2dt=1/12$，所以
\[
M=\frac1{24}\int_0^{3}x(3-x)^4dx=\frac{81}{80}.
\]
最后一步可以展开四次幂或再次使用有限区间 Beta 型积分。密度在区域边界上可能为零，但在内部为正，总质量为正且随尺度的六次方变化。
\end{solve}

\begin{solve}[习题~\ref{ex:c39q22}：单纯形上的乘积密度]\phantomsection\label{ans:c39q22}
先对 $z$ 积分，再对 $y$ 积分：
\[
M=\frac12\int_0^{4}\int_0^{4-x}xy(4-x-y)^2dy\,dx.
\]
对内层作 $y=(a-x)t$ 的缩放，其中 $\int_0^1t(1-t)^2dt=1/12$，所以
\[
M=\frac1{24}\int_0^{4}x(4-x)^4dx=\frac{256}{45}.
\]
最后一步可以展开四次幂或再次使用有限区间 Beta 型积分。密度在区域边界上可能为零，但在内部为正，总质量为正且随尺度的六次方变化。
\end{solve}

\begin{solve}[习题~\ref{ex:c39q23}：单纯形上的乘积密度]\phantomsection\label{ans:c39q23}
先对 $z$ 积分，再对 $y$ 积分：
\[
M=\frac12\int_0^{5}\int_0^{5-x}xy(5-x-y)^2dy\,dx.
\]
对内层作 $y=(a-x)t$ 的缩放，其中 $\int_0^1t(1-t)^2dt=1/12$，所以
\[
M=\frac1{24}\int_0^{5}x(5-x)^4dx=\frac{3125}{144}.
\]
最后一步可以展开四次幂或再次使用有限区间 Beta 型积分。密度在区域边界上可能为零，但在内部为正，总质量为正且随尺度的六次方变化。
\end{solve}

\begin{solve}[习题~\ref{ex:c39q24}：单纯形上的乘积密度]\phantomsection\label{ans:c39q24}
先对 $z$ 积分，再对 $y$ 积分：
\[
M=\frac12\int_0^{6}\int_0^{6-x}xy(6-x-y)^2dy\,dx.
\]
对内层作 $y=(a-x)t$ 的缩放，其中 $\int_0^1t(1-t)^2dt=1/12$，所以
\[
M=\frac1{24}\int_0^{6}x(6-x)^4dx=\frac{324}{5}.
\]
最后一步可以展开四次幂或再次使用有限区间 Beta 型积分。密度在区域边界上可能为零，但在内部为正，总质量为正且随尺度的六次方变化。
\end{solve}

\clearpage
\section{多重积分的参数与反常情形}\label{sec:ch40}
无界区域或奇异被积函数需要说明截断方式及收敛性。对非负被积函数可以利用单调极限与非负积分换序；若采用 Riemann 反常积分，应把区域截断后再计算极限。

\textbf{资源对应：}CLP-3 \S~3.1、\S~3.5、\S~3.7 的扩展训练。本章共 24 道编号习题。
\subsection{习题}
\begin{example}[第一象限的指数积分]\label{ex:c40q1}
计算 $\displaystyle\int_0^\infty\int_0^\infty e^{-1x-2y}dy\,dx$，说明收敛的依据。
\par\textit{改编训练。}\par\hyperref[ans:c40q1]{解答见第~\pageref{ans:c40q1} 页。}
\end{example}
\begin{example}[第一象限的指数积分]\label{ex:c40q2}
计算 $\displaystyle\int_0^\infty\int_0^\infty e^{-2x-1y}dy\,dx$，说明收敛的依据。
\par\textit{改编训练。}\par\hyperref[ans:c40q2]{解答见第~\pageref{ans:c40q2} 页。}
\end{example}
\begin{example}[第一象限的指数积分]\label{ex:c40q3}
计算 $\displaystyle\int_0^\infty\int_0^\infty e^{-2x-3y}dy\,dx$，说明收敛的依据。
\par\textit{改编训练。}\par\hyperref[ans:c40q3]{解答见第~\pageref{ans:c40q3} 页。}
\end{example}
\begin{example}[第一象限的指数积分]\label{ex:c40q4}
计算 $\displaystyle\int_0^\infty\int_0^\infty e^{-3x-2y}dy\,dx$，说明收敛的依据。
\par\textit{改编训练。}\par\hyperref[ans:c40q4]{解答见第~\pageref{ans:c40q4} 页。}
\end{example}
\begin{example}[第一象限的指数积分]\label{ex:c40q5}
计算 $\displaystyle\int_0^\infty\int_0^\infty e^{-4x-1y}dy\,dx$，说明收敛的依据。
\par\textit{改编训练。}\par\hyperref[ans:c40q5]{解答见第~\pageref{ans:c40q5} 页。}
\end{example}
\begin{example}[第一象限的指数积分]\label{ex:c40q6}
计算 $\displaystyle\int_0^\infty\int_0^\infty e^{-3x-4y}dy\,dx$，说明收敛的依据。
\par\textit{改编训练。}\par\hyperref[ans:c40q6]{解答见第~\pageref{ans:c40q6} 页。}
\end{example}
\begin{example}[三角形中的非均匀密度]\label{ex:c40q7}
薄片区域为 $x,y\ge0$、$x+y\le 1$，密度为 $x+y$。求质量和形心。
\par\textit{改编训练。}\par\hyperref[ans:c40q7]{解答见第~\pageref{ans:c40q7} 页。}
\end{example}
\begin{example}[三角形中的非均匀密度]\label{ex:c40q8}
薄片区域为 $x,y\ge0$、$x+y\le 2$，密度为 $x+y$。求质量和形心。
\par\textit{改编训练。}\par\hyperref[ans:c40q8]{解答见第~\pageref{ans:c40q8} 页。}
\end{example}
\begin{example}[三角形中的非均匀密度]\label{ex:c40q9}
薄片区域为 $x,y\ge0$、$x+y\le 3$，密度为 $x+y$。求质量和形心。
\par\textit{改编训练。}\par\hyperref[ans:c40q9]{解答见第~\pageref{ans:c40q9} 页。}
\end{example}
\begin{example}[三角形中的非均匀密度]\label{ex:c40q10}
薄片区域为 $x,y\ge0$、$x+y\le 4$，密度为 $x+y$。求质量和形心。
\par\textit{改编训练。}\par\hyperref[ans:c40q10]{解答见第~\pageref{ans:c40q10} 页。}
\end{example}
\begin{example}[三角形中的非均匀密度]\label{ex:c40q11}
薄片区域为 $x,y\ge0$、$x+y\le 5$，密度为 $x+y$。求质量和形心。
\par\textit{改编训练。}\par\hyperref[ans:c40q11]{解答见第~\pageref{ans:c40q11} 页。}
\end{example}
\begin{example}[三角形中的非均匀密度]\label{ex:c40q12}
薄片区域为 $x,y\ge0$、$x+y\le 6$，密度为 $x+y$。求质量和形心。
\par\textit{改编训练。}\par\hyperref[ans:c40q12]{解答见第~\pageref{ans:c40q12} 页。}
\end{example}
\begin{example}[球心奇点的维数阈值]\label{ex:c40q13}
判断单位球内 $\displaystyle\iiint_B(x^2+y^2+z^2)^{-0}dV$ 的收敛性；收敛时求值。
\par\textit{改编训练。}\par\hyperref[ans:c40q13]{解答见第~\pageref{ans:c40q13} 页。}
\end{example}
\begin{example}[球心奇点的维数阈值]\label{ex:c40q14}
判断单位球内 $\displaystyle\iiint_B(x^2+y^2+z^2)^{-\frac{1}{2}}dV$ 的收敛性；收敛时求值。
\par\textit{改编训练。}\par\hyperref[ans:c40q14]{解答见第~\pageref{ans:c40q14} 页。}
\end{example}
\begin{example}[球心奇点的维数阈值]\label{ex:c40q15}
判断单位球内 $\displaystyle\iiint_B(x^2+y^2+z^2)^{-1}dV$ 的收敛性；收敛时求值。
\par\textit{改编训练。}\par\hyperref[ans:c40q15]{解答见第~\pageref{ans:c40q15} 页。}
\end{example}
\begin{example}[球心奇点的维数阈值]\label{ex:c40q16}
判断单位球内 $\displaystyle\iiint_B(x^2+y^2+z^2)^{-\frac{3}{2}}dV$ 的收敛性；收敛时求值。
\par\textit{改编训练。}\par\hyperref[ans:c40q16]{解答见第~\pageref{ans:c40q16} 页。}
\end{example}
\begin{example}[球心奇点的维数阈值]\label{ex:c40q17}
判断单位球内 $\displaystyle\iiint_B(x^2+y^2+z^2)^{-2}dV$ 的收敛性；收敛时求值。
\par\textit{改编训练。}\par\hyperref[ans:c40q17]{解答见第~\pageref{ans:c40q17} 页。}
\end{example}
\begin{example}[球心奇点的维数阈值]\label{ex:c40q18}
判断单位球内 $\displaystyle\iiint_B(x^2+y^2+z^2)^{-\frac{5}{2}}dV$ 的收敛性；收敛时求值。
\par\textit{改编训练。}\par\hyperref[ans:c40q18]{解答见第~\pageref{ans:c40q18} 页。}
\end{example}
\begin{example}[三角形上的指数密度]\label{ex:c40q19}
计算 $\displaystyle\iint_D e^{x+y}dxdy$，其中 $D:x,y\ge0,x+y\le 1$。
\par\textit{改编训练。}\par\hyperref[ans:c40q19]{解答见第~\pageref{ans:c40q19} 页。}
\end{example}
\begin{example}[三角形上的指数密度]\label{ex:c40q20}
计算 $\displaystyle\iint_D e^{x+y}dxdy$，其中 $D:x,y\ge0,x+y\le 2$。
\par\textit{改编训练。}\par\hyperref[ans:c40q20]{解答见第~\pageref{ans:c40q20} 页。}
\end{example}
\begin{example}[三角形上的指数密度]\label{ex:c40q21}
计算 $\displaystyle\iint_D e^{x+y}dxdy$，其中 $D:x,y\ge0,x+y\le 3$。
\par\textit{改编训练。}\par\hyperref[ans:c40q21]{解答见第~\pageref{ans:c40q21} 页。}
\end{example}
\begin{example}[三角形上的指数密度]\label{ex:c40q22}
计算 $\displaystyle\iint_D e^{x+y}dxdy$，其中 $D:x,y\ge0,x+y\le 4$。
\par\textit{改编训练。}\par\hyperref[ans:c40q22]{解答见第~\pageref{ans:c40q22} 页。}
\end{example}
\begin{example}[三角形上的指数密度]\label{ex:c40q23}
计算 $\displaystyle\iint_D e^{x+y}dxdy$，其中 $D:x,y\ge0,x+y\le 5$。
\par\textit{改编训练。}\par\hyperref[ans:c40q23]{解答见第~\pageref{ans:c40q23} 页。}
\end{example}
\begin{example}[三角形上的指数密度]\label{ex:c40q24}
计算 $\displaystyle\iint_D e^{x+y}dxdy$，其中 $D:x,y\ge0,x+y\le 6$。
\par\textit{改编训练。}\par\hyperref[ans:c40q24]{解答见第~\pageref{ans:c40q24} 页。}
\end{example}
\subsection{解答}
\begin{solve}[习题~\ref{ex:c40q1}：第一象限的指数积分]\phantomsection\label{ans:c40q1}
先在矩形 $[0,R]\times[0,S]$ 上积分，连续被积函数允许分离，得
\[
I_{R,S}=\frac{1-e^{-1R}}{1}\frac{1-e^{-2S}}{2}.
\]
两端分别趋于无穷时，该值单调增加且趋于 $\frac{1}{2}$，故反常积分收敛。非负性保证扩大截断区域时不会发生正负抵消；所得有限界也能控制其他位于第一象限内的截断方式。
\end{solve}

\begin{solve}[习题~\ref{ex:c40q2}：第一象限的指数积分]\phantomsection\label{ans:c40q2}
先在矩形 $[0,R]\times[0,S]$ 上积分，连续被积函数允许分离，得
\[
I_{R,S}=\frac{1-e^{-2R}}{2}\frac{1-e^{-1S}}{1}.
\]
两端分别趋于无穷时，该值单调增加且趋于 $\frac{1}{2}$，故反常积分收敛。非负性保证扩大截断区域时不会发生正负抵消；所得有限界也能控制其他位于第一象限内的截断方式。
\end{solve}

\begin{solve}[习题~\ref{ex:c40q3}：第一象限的指数积分]\phantomsection\label{ans:c40q3}
先在矩形 $[0,R]\times[0,S]$ 上积分，连续被积函数允许分离，得
\[
I_{R,S}=\frac{1-e^{-2R}}{2}\frac{1-e^{-3S}}{3}.
\]
两端分别趋于无穷时，该值单调增加且趋于 $\frac{1}{6}$，故反常积分收敛。非负性保证扩大截断区域时不会发生正负抵消；所得有限界也能控制其他位于第一象限内的截断方式。
\end{solve}

\begin{solve}[习题~\ref{ex:c40q4}：第一象限的指数积分]\phantomsection\label{ans:c40q4}
先在矩形 $[0,R]\times[0,S]$ 上积分，连续被积函数允许分离，得
\[
I_{R,S}=\frac{1-e^{-3R}}{3}\frac{1-e^{-2S}}{2}.
\]
两端分别趋于无穷时，该值单调增加且趋于 $\frac{1}{6}$，故反常积分收敛。非负性保证扩大截断区域时不会发生正负抵消；所得有限界也能控制其他位于第一象限内的截断方式。
\end{solve}

\begin{solve}[习题~\ref{ex:c40q5}：第一象限的指数积分]\phantomsection\label{ans:c40q5}
先在矩形 $[0,R]\times[0,S]$ 上积分，连续被积函数允许分离，得
\[
I_{R,S}=\frac{1-e^{-4R}}{4}\frac{1-e^{-1S}}{1}.
\]
两端分别趋于无穷时，该值单调增加且趋于 $\frac{1}{4}$，故反常积分收敛。非负性保证扩大截断区域时不会发生正负抵消；所得有限界也能控制其他位于第一象限内的截断方式。
\end{solve}

\begin{solve}[习题~\ref{ex:c40q6}：第一象限的指数积分]\phantomsection\label{ans:c40q6}
先在矩形 $[0,R]\times[0,S]$ 上积分，连续被积函数允许分离，得
\[
I_{R,S}=\frac{1-e^{-3R}}{3}\frac{1-e^{-4S}}{4}.
\]
两端分别趋于无穷时，该值单调增加且趋于 $\frac{1}{12}$，故反常积分收敛。非负性保证扩大截断区域时不会发生正负抵消；所得有限界也能控制其他位于第一象限内的截断方式。
\end{solve}

\begin{solve}[习题~\ref{ex:c40q7}：三角形中的非均匀密度]\phantomsection\label{ans:c40q7}
令 $s=x+y$、$u=x$，则 $0\le s\le 1$、$0\le u\le s$，面积 Jacobian 为1。质量为
\[
M=\int_0^{1}\int_0^s s\,du\,ds=\frac{1}{3}.
\]
横坐标的一阶矩为
\[
M_y=\int_0^{1}\int_0^s us\,du\,ds=\frac{1}{8}.
\]
所以 $\bar x=\frac{3}{8}$，由对称性 $\bar y=\frac{3}{8}$。密度随到斜边方向的坐标和增大，形心比均匀三角形的 $(a/3,a/3)$ 更靠近斜边。
\end{solve}

\begin{solve}[习题~\ref{ex:c40q8}：三角形中的非均匀密度]\phantomsection\label{ans:c40q8}
令 $s=x+y$、$u=x$，则 $0\le s\le 2$、$0\le u\le s$，面积 Jacobian 为1。质量为
\[
M=\int_0^{2}\int_0^s s\,du\,ds=\frac{8}{3}.
\]
横坐标的一阶矩为
\[
M_y=\int_0^{2}\int_0^s us\,du\,ds=2.
\]
所以 $\bar x=\frac{3}{4}$，由对称性 $\bar y=\frac{3}{4}$。密度随到斜边方向的坐标和增大，形心比均匀三角形的 $(a/3,a/3)$ 更靠近斜边。
\end{solve}

\begin{solve}[习题~\ref{ex:c40q9}：三角形中的非均匀密度]\phantomsection\label{ans:c40q9}
令 $s=x+y$、$u=x$，则 $0\le s\le 3$、$0\le u\le s$，面积 Jacobian 为1。质量为
\[
M=\int_0^{3}\int_0^s s\,du\,ds=9.
\]
横坐标的一阶矩为
\[
M_y=\int_0^{3}\int_0^s us\,du\,ds=\frac{81}{8}.
\]
所以 $\bar x=\frac{9}{8}$，由对称性 $\bar y=\frac{9}{8}$。密度随到斜边方向的坐标和增大，形心比均匀三角形的 $(a/3,a/3)$ 更靠近斜边。
\end{solve}

\begin{solve}[习题~\ref{ex:c40q10}：三角形中的非均匀密度]\phantomsection\label{ans:c40q10}
令 $s=x+y$、$u=x$，则 $0\le s\le 4$、$0\le u\le s$，面积 Jacobian 为1。质量为
\[
M=\int_0^{4}\int_0^s s\,du\,ds=\frac{64}{3}.
\]
横坐标的一阶矩为
\[
M_y=\int_0^{4}\int_0^s us\,du\,ds=32.
\]
所以 $\bar x=\frac{3}{2}$，由对称性 $\bar y=\frac{3}{2}$。密度随到斜边方向的坐标和增大，形心比均匀三角形的 $(a/3,a/3)$ 更靠近斜边。
\end{solve}

\begin{solve}[习题~\ref{ex:c40q11}：三角形中的非均匀密度]\phantomsection\label{ans:c40q11}
令 $s=x+y$、$u=x$，则 $0\le s\le 5$、$0\le u\le s$，面积 Jacobian 为1。质量为
\[
M=\int_0^{5}\int_0^s s\,du\,ds=\frac{125}{3}.
\]
横坐标的一阶矩为
\[
M_y=\int_0^{5}\int_0^s us\,du\,ds=\frac{625}{8}.
\]
所以 $\bar x=\frac{15}{8}$，由对称性 $\bar y=\frac{15}{8}$。密度随到斜边方向的坐标和增大，形心比均匀三角形的 $(a/3,a/3)$ 更靠近斜边。
\end{solve}

\begin{solve}[习题~\ref{ex:c40q12}：三角形中的非均匀密度]\phantomsection\label{ans:c40q12}
令 $s=x+y$、$u=x$，则 $0\le s\le 6$、$0\le u\le s$，面积 Jacobian 为1。质量为
\[
M=\int_0^{6}\int_0^s s\,du\,ds=72.
\]
横坐标的一阶矩为
\[
M_y=\int_0^{6}\int_0^s us\,du\,ds=162.
\]
所以 $\bar x=\frac{9}{4}$，由对称性 $\bar y=\frac{9}{4}$。密度随到斜边方向的坐标和增大，形心比均匀三角形的 $(a/3,a/3)$ 更靠近斜边。
\end{solve}

\begin{solve}[习题~\ref{ex:c40q13}：球心奇点的维数阈值]\phantomsection\label{ans:c40q13}
在球心去掉半径 $\varepsilon$ 的小球，用球坐标得到
\[
I_\varepsilon=4\pi\int_\varepsilon^1\rho^{2}d\rho.
\]
指数 $2>-1$，因此积分收敛，值为 $\frac{4}{3}\pi$。三维体积元自带 $\rho^2$ 因子，所以奇异强度的阈值是 $p<3$；不能直接套用一维的 $p<1$。
\end{solve}

\begin{solve}[习题~\ref{ex:c40q14}：球心奇点的维数阈值]\phantomsection\label{ans:c40q14}
在球心去掉半径 $\varepsilon$ 的小球，用球坐标得到
\[
I_\varepsilon=4\pi\int_\varepsilon^1\rho^{1}d\rho.
\]
指数 $1>-1$，因此积分收敛，值为 $2\pi$。三维体积元自带 $\rho^2$ 因子，所以奇异强度的阈值是 $p<3$；不能直接套用一维的 $p<1$。
\end{solve}

\begin{solve}[习题~\ref{ex:c40q15}：球心奇点的维数阈值]\phantomsection\label{ans:c40q15}
在球心去掉半径 $\varepsilon$ 的小球，用球坐标得到
\[
I_\varepsilon=4\pi\int_\varepsilon^1\rho^{0}d\rho.
\]
指数 $0>-1$，因此积分收敛，值为 $4\pi$。三维体积元自带 $\rho^2$ 因子，所以奇异强度的阈值是 $p<3$；不能直接套用一维的 $p<1$。
\end{solve}

\begin{solve}[习题~\ref{ex:c40q16}：球心奇点的维数阈值]\phantomsection\label{ans:c40q16}
在球心去掉半径 $\varepsilon$ 的小球，用球坐标得到
\[
I_\varepsilon=4\pi\int_\varepsilon^1\rho^{-1}d\rho.
\]
此时 $I_\varepsilon=-4\pi\log\varepsilon\to\infty$，故发散。三维体积元自带 $\rho^2$ 因子，所以奇异强度的阈值是 $p<3$；不能直接套用一维的 $p<1$。
\end{solve}

\begin{solve}[习题~\ref{ex:c40q17}：球心奇点的维数阈值]\phantomsection\label{ans:c40q17}
在球心去掉半径 $\varepsilon$ 的小球，用球坐标得到
\[
I_\varepsilon=4\pi\int_\varepsilon^1\rho^{-2}d\rho.
\]
幂指数 $-2<-1$，积分在0附近发散到正无穷。三维体积元自带 $\rho^2$ 因子，所以奇异强度的阈值是 $p<3$；不能直接套用一维的 $p<1$。
\end{solve}

\begin{solve}[习题~\ref{ex:c40q18}：球心奇点的维数阈值]\phantomsection\label{ans:c40q18}
在球心去掉半径 $\varepsilon$ 的小球，用球坐标得到
\[
I_\varepsilon=4\pi\int_\varepsilon^1\rho^{-3}d\rho.
\]
幂指数 $-3<-1$，积分在0附近发散到正无穷。三维体积元自带 $\rho^2$ 因子，所以奇异强度的阈值是 $p<3$；不能直接套用一维的 $p<1$。
\end{solve}

\begin{solve}[习题~\ref{ex:c40q19}：三角形上的指数密度]\phantomsection\label{ans:c40q19}
取 $u=x+y$、$v=x$，区域变为 $0\le u\le 1,0\le v\le u$，Jacobian 绝对值为1。于是
\[
I=\int_0^{1}\int_0^ue^u\,dv\,du=\int_0^{1}ue^u\,du=(0)e^{1}+1.
\]
也可先按直角坐标求内积分，得到同样结果。新的变量把被积函数中反复出现的组合 $x+y$ 作为独立变量，让积分结构与区域边界同时简化。
\end{solve}

\begin{solve}[习题~\ref{ex:c40q20}：三角形上的指数密度]\phantomsection\label{ans:c40q20}
取 $u=x+y$、$v=x$，区域变为 $0\le u\le 2,0\le v\le u$，Jacobian 绝对值为1。于是
\[
I=\int_0^{2}\int_0^ue^u\,dv\,du=\int_0^{2}ue^u\,du=(1)e^{2}+1.
\]
也可先按直角坐标求内积分，得到同样结果。新的变量把被积函数中反复出现的组合 $x+y$ 作为独立变量，让积分结构与区域边界同时简化。
\end{solve}

\begin{solve}[习题~\ref{ex:c40q21}：三角形上的指数密度]\phantomsection\label{ans:c40q21}
取 $u=x+y$、$v=x$，区域变为 $0\le u\le 3,0\le v\le u$，Jacobian 绝对值为1。于是
\[
I=\int_0^{3}\int_0^ue^u\,dv\,du=\int_0^{3}ue^u\,du=(2)e^{3}+1.
\]
也可先按直角坐标求内积分，得到同样结果。新的变量把被积函数中反复出现的组合 $x+y$ 作为独立变量，让积分结构与区域边界同时简化。
\end{solve}

\begin{solve}[习题~\ref{ex:c40q22}：三角形上的指数密度]\phantomsection\label{ans:c40q22}
取 $u=x+y$、$v=x$，区域变为 $0\le u\le 4,0\le v\le u$，Jacobian 绝对值为1。于是
\[
I=\int_0^{4}\int_0^ue^u\,dv\,du=\int_0^{4}ue^u\,du=(3)e^{4}+1.
\]
也可先按直角坐标求内积分，得到同样结果。新的变量把被积函数中反复出现的组合 $x+y$ 作为独立变量，让积分结构与区域边界同时简化。
\end{solve}

\begin{solve}[习题~\ref{ex:c40q23}：三角形上的指数密度]\phantomsection\label{ans:c40q23}
取 $u=x+y$、$v=x$，区域变为 $0\le u\le 5,0\le v\le u$，Jacobian 绝对值为1。于是
\[
I=\int_0^{5}\int_0^ue^u\,dv\,du=\int_0^{5}ue^u\,du=(4)e^{5}+1.
\]
也可先按直角坐标求内积分，得到同样结果。新的变量把被积函数中反复出现的组合 $x+y$ 作为独立变量，让积分结构与区域边界同时简化。
\end{solve}

\begin{solve}[习题~\ref{ex:c40q24}：三角形上的指数密度]\phantomsection\label{ans:c40q24}
取 $u=x+y$、$v=x$，区域变为 $0\le u\le 6,0\le v\le u$，Jacobian 绝对值为1。于是
\[
I=\int_0^{6}\int_0^ue^u\,dv\,du=\int_0^{6}ue^u\,du=(5)e^{6}+1.
\]
也可先按直角坐标求内积分，得到同样结果。新的变量把被积函数中反复出现的组合 $x+y$ 作为独立变量，让积分结构与区域边界同时简化。
\end{solve}

\clearpage
\section{曲线积分与保守场}\label{sec:ch41}
第一类曲线积分对弧长积分，不随曲线反向而改变；第二类曲线积分对有向位移积分，反向时改变符号。若向量场有势函数，路径积分可化为势函数的端点差。

\textbf{资源对应：}CLP-4 \S~1.6、\S~2.3、\S~2.4、\S~4.3。本章共 25 道编号习题。
\subsection{习题}
\begin{example}[正方形边界的直接积分]\label{ex:c41q1}
沿正方形 $(0,0),(1,0),(1,1),(0,1)$ 的逆时针边界，计算 $\oint_Cx^2y^2dx+x^3y\,dy$。
\par\textit{选译。来源：CLP-4 \S~2.4，首道正方形边界题。}\par\hyperref[ans:c41q1]{解答见第~\pageref{ans:c41q1} 页。}
\end{example}
\begin{example}[圆上的第一类曲线积分]\label{ex:c41q2}
设 $C:x^2+y^2=1$。求 $\displaystyle\int_Cx^2ds$，并说明曲线反向时结果如何变化。
\par\textit{改编训练。}\par\hyperref[ans:c41q2]{解答见第~\pageref{ans:c41q2} 页。}
\end{example}
\begin{example}[圆上的第一类曲线积分]\label{ex:c41q3}
设 $C:x^2+y^2=4$。求 $\displaystyle\int_Cx^2ds$，并说明曲线反向时结果如何变化。
\par\textit{改编训练。}\par\hyperref[ans:c41q3]{解答见第~\pageref{ans:c41q3} 页。}
\end{example}
\begin{example}[圆上的第一类曲线积分]\label{ex:c41q4}
设 $C:x^2+y^2=9$。求 $\displaystyle\int_Cx^2ds$，并说明曲线反向时结果如何变化。
\par\textit{改编训练。}\par\hyperref[ans:c41q4]{解答见第~\pageref{ans:c41q4} 页。}
\end{example}
\begin{example}[圆上的第一类曲线积分]\label{ex:c41q5}
设 $C:x^2+y^2=16$。求 $\displaystyle\int_Cx^2ds$，并说明曲线反向时结果如何变化。
\par\textit{改编训练。}\par\hyperref[ans:c41q5]{解答见第~\pageref{ans:c41q5} 页。}
\end{example}
\begin{example}[圆上的第一类曲线积分]\label{ex:c41q6}
设 $C:x^2+y^2=25$。求 $\displaystyle\int_Cx^2ds$，并说明曲线反向时结果如何变化。
\par\textit{改编训练。}\par\hyperref[ans:c41q6]{解答见第~\pageref{ans:c41q6} 页。}
\end{example}
\begin{example}[圆上的第一类曲线积分]\label{ex:c41q7}
设 $C:x^2+y^2=36$。求 $\displaystyle\int_Cx^2ds$，并说明曲线反向时结果如何变化。
\par\textit{改编训练。}\par\hyperref[ans:c41q7]{解答见第~\pageref{ans:c41q7} 页。}
\end{example}
\begin{example}[椭圆上的环流]\label{ex:c41q8}
沿椭圆 $x=1\cos t$、$y=2\sin t$（$0\le t\le2\pi$）计算 $\oint_C(-y\,dx+x\,dy)$。
\par\textit{改编训练。}\par\hyperref[ans:c41q8]{解答见第~\pageref{ans:c41q8} 页。}
\end{example}
\begin{example}[椭圆上的环流]\label{ex:c41q9}
沿椭圆 $x=2\cos t$、$y=1\sin t$（$0\le t\le2\pi$）计算 $\oint_C(-y\,dx+x\,dy)$。
\par\textit{改编训练。}\par\hyperref[ans:c41q9]{解答见第~\pageref{ans:c41q9} 页。}
\end{example}
\begin{example}[椭圆上的环流]\label{ex:c41q10}
沿椭圆 $x=2\cos t$、$y=3\sin t$（$0\le t\le2\pi$）计算 $\oint_C(-y\,dx+x\,dy)$。
\par\textit{改编训练。}\par\hyperref[ans:c41q10]{解答见第~\pageref{ans:c41q10} 页。}
\end{example}
\begin{example}[椭圆上的环流]\label{ex:c41q11}
沿椭圆 $x=3\cos t$、$y=2\sin t$（$0\le t\le2\pi$）计算 $\oint_C(-y\,dx+x\,dy)$。
\par\textit{改编训练。}\par\hyperref[ans:c41q11]{解答见第~\pageref{ans:c41q11} 页。}
\end{example}
\begin{example}[椭圆上的环流]\label{ex:c41q12}
沿椭圆 $x=4\cos t$、$y=1\sin t$（$0\le t\le2\pi$）计算 $\oint_C(-y\,dx+x\,dy)$。
\par\textit{改编训练。}\par\hyperref[ans:c41q12]{解答见第~\pageref{ans:c41q12} 页。}
\end{example}
\begin{example}[椭圆上的环流]\label{ex:c41q13}
沿椭圆 $x=3\cos t$、$y=4\sin t$（$0\le t\le2\pi$）计算 $\oint_C(-y\,dx+x\,dy)$。
\par\textit{改编训练。}\par\hyperref[ans:c41q13]{解答见第~\pageref{ans:c41q13} 页。}
\end{example}
\begin{example}[势函数与路径无关]\label{ex:c41q14}
设 $F=(2xy+1,x^2+2)$，计算从 $(0,0)$ 到 $(1,2)$ 的曲线积分，并构造势函数。
\par\textit{改编训练。}\par\hyperref[ans:c41q14]{解答见第~\pageref{ans:c41q14} 页。}
\end{example}
\begin{example}[势函数与路径无关]\label{ex:c41q15}
设 $F=(2xy+2,x^2+1)$，计算从 $(0,0)$ 到 $(1,2)$ 的曲线积分，并构造势函数。
\par\textit{改编训练。}\par\hyperref[ans:c41q15]{解答见第~\pageref{ans:c41q15} 页。}
\end{example}
\begin{example}[势函数与路径无关]\label{ex:c41q16}
设 $F=(2xy+2,x^2+3)$，计算从 $(0,0)$ 到 $(1,2)$ 的曲线积分，并构造势函数。
\par\textit{改编训练。}\par\hyperref[ans:c41q16]{解答见第~\pageref{ans:c41q16} 页。}
\end{example}
\begin{example}[势函数与路径无关]\label{ex:c41q17}
设 $F=(2xy+3,x^2+2)$，计算从 $(0,0)$ 到 $(1,2)$ 的曲线积分，并构造势函数。
\par\textit{改编训练。}\par\hyperref[ans:c41q17]{解答见第~\pageref{ans:c41q17} 页。}
\end{example}
\begin{example}[势函数与路径无关]\label{ex:c41q18}
设 $F=(2xy+4,x^2+1)$，计算从 $(0,0)$ 到 $(1,2)$ 的曲线积分，并构造势函数。
\par\textit{改编训练。}\par\hyperref[ans:c41q18]{解答见第~\pageref{ans:c41q18} 页。}
\end{example}
\begin{example}[势函数与路径无关]\label{ex:c41q19}
设 $F=(2xy+3,x^2+4)$，计算从 $(0,0)$ 到 $(1,2)$ 的曲线积分，并构造势函数。
\par\textit{改编训练。}\par\hyperref[ans:c41q19]{解答见第~\pageref{ans:c41q19} 页。}
\end{example}
\begin{example}[有孔区域与无旋场]\label{ex:c41q20}
在去掉原点的平面上设 $F=1(-y,x)/(x^2+y^2)$。证明其二维旋度为零，但沿单位圆逆时针积分不为零。
\par\textit{改编训练。}\par\hyperref[ans:c41q20]{解答见第~\pageref{ans:c41q20} 页。}
\end{example}
\begin{example}[有孔区域与无旋场]\label{ex:c41q21}
在去掉原点的平面上设 $F=2(-y,x)/(x^2+y^2)$。证明其二维旋度为零，但沿单位圆逆时针积分不为零。
\par\textit{改编训练。}\par\hyperref[ans:c41q21]{解答见第~\pageref{ans:c41q21} 页。}
\end{example}
\begin{example}[有孔区域与无旋场]\label{ex:c41q22}
在去掉原点的平面上设 $F=3(-y,x)/(x^2+y^2)$。证明其二维旋度为零，但沿单位圆逆时针积分不为零。
\par\textit{改编训练。}\par\hyperref[ans:c41q22]{解答见第~\pageref{ans:c41q22} 页。}
\end{example}
\begin{example}[有孔区域与无旋场]\label{ex:c41q23}
在去掉原点的平面上设 $F=4(-y,x)/(x^2+y^2)$。证明其二维旋度为零，但沿单位圆逆时针积分不为零。
\par\textit{改编训练。}\par\hyperref[ans:c41q23]{解答见第~\pageref{ans:c41q23} 页。}
\end{example}
\begin{example}[有孔区域与无旋场]\label{ex:c41q24}
在去掉原点的平面上设 $F=5(-y,x)/(x^2+y^2)$。证明其二维旋度为零，但沿单位圆逆时针积分不为零。
\par\textit{改编训练。}\par\hyperref[ans:c41q24]{解答见第~\pageref{ans:c41q24} 页。}
\end{example}
\begin{example}[有孔区域与无旋场]\label{ex:c41q25}
在去掉原点的平面上设 $F=6(-y,x)/(x^2+y^2)$。证明其二维旋度为零，但沿单位圆逆时针积分不为零。
\par\textit{改编训练。}\par\hyperref[ans:c41q25]{解答见第~\pageref{ans:c41q25} 页。}
\end{example}
\subsection{解答}
\begin{solve}[习题~\ref{ex:c41q1}：正方形边界的直接积分]\phantomsection\label{ans:c41q1}
依次计算下、右、上、左四条边。下边 $y=0$，贡献0；右边 $x=1$、$y:0\to1$，贡献 $\int_0^1y\,dy=1/2$；上边 $y=1$、$x:1\to0$，贡献 $\int_1^0x^2dx=-1/3$；左边 $x=0$，贡献0。总和为 $1/6$。Green 公式也给出区域积分 $\iint(3x^2y-2x^2y)dxdy=1/6$。
\end{solve}

\begin{solve}[习题~\ref{ex:c41q2}：圆上的第一类曲线积分]\phantomsection\label{ans:c41q2}
参数化为 $x=1\cos t$、$y=1\sin t$，$0\le t\le2\pi$。速度模为 $1$，所以
\[
\int_Cx^2ds=1\int_0^{2\pi}\cos^2t\,dt=1\pi.
\]
第一类积分中的 $ds$ 是非负弧长元，反向参数化不改变结果。还可利用圆的对称性得 $\int_Cx^2ds=\int_Cy^2ds$，两者之和等于常数 $a^2$ 乘周长，验证相同结果。
\end{solve}

\begin{solve}[习题~\ref{ex:c41q3}：圆上的第一类曲线积分]\phantomsection\label{ans:c41q3}
参数化为 $x=2\cos t$、$y=2\sin t$，$0\le t\le2\pi$。速度模为 $2$，所以
\[
\int_Cx^2ds=8\int_0^{2\pi}\cos^2t\,dt=8\pi.
\]
第一类积分中的 $ds$ 是非负弧长元，反向参数化不改变结果。还可利用圆的对称性得 $\int_Cx^2ds=\int_Cy^2ds$，两者之和等于常数 $a^2$ 乘周长，验证相同结果。
\end{solve}

\begin{solve}[习题~\ref{ex:c41q4}：圆上的第一类曲线积分]\phantomsection\label{ans:c41q4}
参数化为 $x=3\cos t$、$y=3\sin t$，$0\le t\le2\pi$。速度模为 $3$，所以
\[
\int_Cx^2ds=27\int_0^{2\pi}\cos^2t\,dt=27\pi.
\]
第一类积分中的 $ds$ 是非负弧长元，反向参数化不改变结果。还可利用圆的对称性得 $\int_Cx^2ds=\int_Cy^2ds$，两者之和等于常数 $a^2$ 乘周长，验证相同结果。
\end{solve}

\begin{solve}[习题~\ref{ex:c41q5}：圆上的第一类曲线积分]\phantomsection\label{ans:c41q5}
参数化为 $x=4\cos t$、$y=4\sin t$，$0\le t\le2\pi$。速度模为 $4$，所以
\[
\int_Cx^2ds=64\int_0^{2\pi}\cos^2t\,dt=64\pi.
\]
第一类积分中的 $ds$ 是非负弧长元，反向参数化不改变结果。还可利用圆的对称性得 $\int_Cx^2ds=\int_Cy^2ds$，两者之和等于常数 $a^2$ 乘周长，验证相同结果。
\end{solve}

\begin{solve}[习题~\ref{ex:c41q6}：圆上的第一类曲线积分]\phantomsection\label{ans:c41q6}
参数化为 $x=5\cos t$、$y=5\sin t$，$0\le t\le2\pi$。速度模为 $5$，所以
\[
\int_Cx^2ds=125\int_0^{2\pi}\cos^2t\,dt=125\pi.
\]
第一类积分中的 $ds$ 是非负弧长元，反向参数化不改变结果。还可利用圆的对称性得 $\int_Cx^2ds=\int_Cy^2ds$，两者之和等于常数 $a^2$ 乘周长，验证相同结果。
\end{solve}

\begin{solve}[习题~\ref{ex:c41q7}：圆上的第一类曲线积分]\phantomsection\label{ans:c41q7}
参数化为 $x=6\cos t$、$y=6\sin t$，$0\le t\le2\pi$。速度模为 $6$，所以
\[
\int_Cx^2ds=216\int_0^{2\pi}\cos^2t\,dt=216\pi.
\]
第一类积分中的 $ds$ 是非负弧长元，反向参数化不改变结果。还可利用圆的对称性得 $\int_Cx^2ds=\int_Cy^2ds$，两者之和等于常数 $a^2$ 乘周长，验证相同结果。
\end{solve}

\begin{solve}[习题~\ref{ex:c41q8}：椭圆上的环流]\phantomsection\label{ans:c41q8}
该参数方向为逆时针，$dx=-1\sin t\,dt$、$dy=2\cos t\,dt$，所以
\[
-y\,dx+x\,dy=2(\sin^2t+\cos^2t)dt=2\,dt.
\]
积分为 $4\pi$。用 Green 公式验证：$P=-y,Q=x$，故 $Q_x-P_y=2$，环流等于椭圆面积 $2\pi$ 的两倍。若改为顺时针方向，答案变为负数。
\end{solve}

\begin{solve}[习题~\ref{ex:c41q9}：椭圆上的环流]\phantomsection\label{ans:c41q9}
该参数方向为逆时针，$dx=-2\sin t\,dt$、$dy=1\cos t\,dt$，所以
\[
-y\,dx+x\,dy=2(\sin^2t+\cos^2t)dt=2\,dt.
\]
积分为 $4\pi$。用 Green 公式验证：$P=-y,Q=x$，故 $Q_x-P_y=2$，环流等于椭圆面积 $2\pi$ 的两倍。若改为顺时针方向，答案变为负数。
\end{solve}

\begin{solve}[习题~\ref{ex:c41q10}：椭圆上的环流]\phantomsection\label{ans:c41q10}
该参数方向为逆时针，$dx=-2\sin t\,dt$、$dy=3\cos t\,dt$，所以
\[
-y\,dx+x\,dy=6(\sin^2t+\cos^2t)dt=6\,dt.
\]
积分为 $12\pi$。用 Green 公式验证：$P=-y,Q=x$，故 $Q_x-P_y=2$，环流等于椭圆面积 $6\pi$ 的两倍。若改为顺时针方向，答案变为负数。
\end{solve}

\begin{solve}[习题~\ref{ex:c41q11}：椭圆上的环流]\phantomsection\label{ans:c41q11}
该参数方向为逆时针，$dx=-3\sin t\,dt$、$dy=2\cos t\,dt$，所以
\[
-y\,dx+x\,dy=6(\sin^2t+\cos^2t)dt=6\,dt.
\]
积分为 $12\pi$。用 Green 公式验证：$P=-y,Q=x$，故 $Q_x-P_y=2$，环流等于椭圆面积 $6\pi$ 的两倍。若改为顺时针方向，答案变为负数。
\end{solve}

\begin{solve}[习题~\ref{ex:c41q12}：椭圆上的环流]\phantomsection\label{ans:c41q12}
该参数方向为逆时针，$dx=-4\sin t\,dt$、$dy=1\cos t\,dt$，所以
\[
-y\,dx+x\,dy=4(\sin^2t+\cos^2t)dt=4\,dt.
\]
积分为 $8\pi$。用 Green 公式验证：$P=-y,Q=x$，故 $Q_x-P_y=2$，环流等于椭圆面积 $4\pi$ 的两倍。若改为顺时针方向，答案变为负数。
\end{solve}

\begin{solve}[习题~\ref{ex:c41q13}：椭圆上的环流]\phantomsection\label{ans:c41q13}
该参数方向为逆时针，$dx=-3\sin t\,dt$、$dy=4\cos t\,dt$，所以
\[
-y\,dx+x\,dy=12(\sin^2t+\cos^2t)dt=12\,dt.
\]
积分为 $24\pi$。用 Green 公式验证：$P=-y,Q=x$，故 $Q_x-P_y=2$，环流等于椭圆面积 $12\pi$ 的两倍。若改为顺时针方向，答案变为负数。
\end{solve}

\begin{solve}[习题~\ref{ex:c41q14}：势函数与路径无关]\phantomsection\label{ans:c41q14}
先对第一分量关于 $x$ 积分，势函数候选为 $\varphi=x^2y+ax+h(y)$。令其对 $y$ 的导数等于第二分量，得到 $h^{\prime}(y)=b$。因此
\[
\varphi(x,y)=x^2y+1x+2y+C.
\]
该势函数在全平面上有定义，故沿任意分段光滑路径，
\[
\int_C F\cdot dr=\varphi(1,2)-\varphi(0,0)=7.
\]
这里实际构造了势函数，所以路径无关已经得到证明；只检查混合偏导相等时，还须考虑定义域的拓扑条件。
\end{solve}

\begin{solve}[习题~\ref{ex:c41q15}：势函数与路径无关]\phantomsection\label{ans:c41q15}
先对第一分量关于 $x$ 积分，势函数候选为 $\varphi=x^2y+ax+h(y)$。令其对 $y$ 的导数等于第二分量，得到 $h^{\prime}(y)=b$。因此
\[
\varphi(x,y)=x^2y+2x+1y+C.
\]
该势函数在全平面上有定义，故沿任意分段光滑路径，
\[
\int_C F\cdot dr=\varphi(1,2)-\varphi(0,0)=6.
\]
这里实际构造了势函数，所以路径无关已经得到证明；只检查混合偏导相等时，还须考虑定义域的拓扑条件。
\end{solve}

\begin{solve}[习题~\ref{ex:c41q16}：势函数与路径无关]\phantomsection\label{ans:c41q16}
先对第一分量关于 $x$ 积分，势函数候选为 $\varphi=x^2y+ax+h(y)$。令其对 $y$ 的导数等于第二分量，得到 $h^{\prime}(y)=b$。因此
\[
\varphi(x,y)=x^2y+2x+3y+C.
\]
该势函数在全平面上有定义，故沿任意分段光滑路径，
\[
\int_C F\cdot dr=\varphi(1,2)-\varphi(0,0)=10.
\]
这里实际构造了势函数，所以路径无关已经得到证明；只检查混合偏导相等时，还须考虑定义域的拓扑条件。
\end{solve}

\begin{solve}[习题~\ref{ex:c41q17}：势函数与路径无关]\phantomsection\label{ans:c41q17}
先对第一分量关于 $x$ 积分，势函数候选为 $\varphi=x^2y+ax+h(y)$。令其对 $y$ 的导数等于第二分量，得到 $h^{\prime}(y)=b$。因此
\[
\varphi(x,y)=x^2y+3x+2y+C.
\]
该势函数在全平面上有定义，故沿任意分段光滑路径，
\[
\int_C F\cdot dr=\varphi(1,2)-\varphi(0,0)=9.
\]
这里实际构造了势函数，所以路径无关已经得到证明；只检查混合偏导相等时，还须考虑定义域的拓扑条件。
\end{solve}

\begin{solve}[习题~\ref{ex:c41q18}：势函数与路径无关]\phantomsection\label{ans:c41q18}
先对第一分量关于 $x$ 积分，势函数候选为 $\varphi=x^2y+ax+h(y)$。令其对 $y$ 的导数等于第二分量，得到 $h^{\prime}(y)=b$。因此
\[
\varphi(x,y)=x^2y+4x+1y+C.
\]
该势函数在全平面上有定义，故沿任意分段光滑路径，
\[
\int_C F\cdot dr=\varphi(1,2)-\varphi(0,0)=8.
\]
这里实际构造了势函数，所以路径无关已经得到证明；只检查混合偏导相等时，还须考虑定义域的拓扑条件。
\end{solve}

\begin{solve}[习题~\ref{ex:c41q19}：势函数与路径无关]\phantomsection\label{ans:c41q19}
先对第一分量关于 $x$ 积分，势函数候选为 $\varphi=x^2y+ax+h(y)$。令其对 $y$ 的导数等于第二分量，得到 $h^{\prime}(y)=b$。因此
\[
\varphi(x,y)=x^2y+3x+4y+C.
\]
该势函数在全平面上有定义，故沿任意分段光滑路径，
\[
\int_C F\cdot dr=\varphi(1,2)-\varphi(0,0)=13.
\]
这里实际构造了势函数，所以路径无关已经得到证明；只检查混合偏导相等时，还须考虑定义域的拓扑条件。
\end{solve}

\begin{solve}[习题~\ref{ex:c41q20}：有孔区域与无旋场]\phantomsection\label{ans:c41q20}
写 $P=-1y/(x^2+y^2)$、$Q=1x/(x^2+y^2)$。直接求偏导可得 $Q_x=P_y=1(y^2-x^2)/(x^2+y^2)^2$，所以 $Q_x-P_y=0$。沿单位圆 $r(t)=(\cos t,\sin t)$，有
\[
F(r(t))\cdot r^{\prime}(t)=1,\qquad\oint_C F\cdot dr=2\pi\ne0.
\]
因此该场在整个去心平面上没有单值势函数。若套用 Green 公式到整个单位圆盘，便会遗漏向量场在原点无定义、偏导不连续这一不满足假设的事实。
\end{solve}

\begin{solve}[习题~\ref{ex:c41q21}：有孔区域与无旋场]\phantomsection\label{ans:c41q21}
写 $P=-2y/(x^2+y^2)$、$Q=2x/(x^2+y^2)$。直接求偏导可得 $Q_x=P_y=2(y^2-x^2)/(x^2+y^2)^2$，所以 $Q_x-P_y=0$。沿单位圆 $r(t)=(\cos t,\sin t)$，有
\[
F(r(t))\cdot r^{\prime}(t)=2,\qquad\oint_C F\cdot dr=4\pi\ne0.
\]
因此该场在整个去心平面上没有单值势函数。若套用 Green 公式到整个单位圆盘，便会遗漏向量场在原点无定义、偏导不连续这一不满足假设的事实。
\end{solve}

\begin{solve}[习题~\ref{ex:c41q22}：有孔区域与无旋场]\phantomsection\label{ans:c41q22}
写 $P=-3y/(x^2+y^2)$、$Q=3x/(x^2+y^2)$。直接求偏导可得 $Q_x=P_y=3(y^2-x^2)/(x^2+y^2)^2$，所以 $Q_x-P_y=0$。沿单位圆 $r(t)=(\cos t,\sin t)$，有
\[
F(r(t))\cdot r^{\prime}(t)=3,\qquad\oint_C F\cdot dr=6\pi\ne0.
\]
因此该场在整个去心平面上没有单值势函数。若套用 Green 公式到整个单位圆盘，便会遗漏向量场在原点无定义、偏导不连续这一不满足假设的事实。
\end{solve}

\begin{solve}[习题~\ref{ex:c41q23}：有孔区域与无旋场]\phantomsection\label{ans:c41q23}
写 $P=-4y/(x^2+y^2)$、$Q=4x/(x^2+y^2)$。直接求偏导可得 $Q_x=P_y=4(y^2-x^2)/(x^2+y^2)^2$，所以 $Q_x-P_y=0$。沿单位圆 $r(t)=(\cos t,\sin t)$，有
\[
F(r(t))\cdot r^{\prime}(t)=4,\qquad\oint_C F\cdot dr=8\pi\ne0.
\]
因此该场在整个去心平面上没有单值势函数。若套用 Green 公式到整个单位圆盘，便会遗漏向量场在原点无定义、偏导不连续这一不满足假设的事实。
\end{solve}

\begin{solve}[习题~\ref{ex:c41q24}：有孔区域与无旋场]\phantomsection\label{ans:c41q24}
写 $P=-5y/(x^2+y^2)$、$Q=5x/(x^2+y^2)$。直接求偏导可得 $Q_x=P_y=5(y^2-x^2)/(x^2+y^2)^2$，所以 $Q_x-P_y=0$。沿单位圆 $r(t)=(\cos t,\sin t)$，有
\[
F(r(t))\cdot r^{\prime}(t)=5,\qquad\oint_C F\cdot dr=10\pi\ne0.
\]
因此该场在整个去心平面上没有单值势函数。若套用 Green 公式到整个单位圆盘，便会遗漏向量场在原点无定义、偏导不连续这一不满足假设的事实。
\end{solve}

\begin{solve}[习题~\ref{ex:c41q25}：有孔区域与无旋场]\phantomsection\label{ans:c41q25}
写 $P=-6y/(x^2+y^2)$、$Q=6x/(x^2+y^2)$。直接求偏导可得 $Q_x=P_y=6(y^2-x^2)/(x^2+y^2)^2$，所以 $Q_x-P_y=0$。沿单位圆 $r(t)=(\cos t,\sin t)$，有
\[
F(r(t))\cdot r^{\prime}(t)=6,\qquad\oint_C F\cdot dr=12\pi\ne0.
\]
因此该场在整个去心平面上没有单值势函数。若套用 Green 公式到整个单位圆盘，便会遗漏向量场在原点无定义、偏导不连续这一不满足假设的事实。
\end{solve}

\clearpage
\section{曲线与曲面积分的计算综合}\label{sec:ch42}
参数化需要同时给出参数范围和方向。第一类积分用速度模或面积因子，第二类积分用有向速度或有向面积向量；把两类微元混淆会引起系统性的符号和系数错误。

\textbf{资源对应：}CLP-4 \S~1.6、\S~2.4、\S~3.3。本章共 24 道编号习题。
\subsection{习题}
\begin{example}[抛物线弧长]\label{ex:c42q1}
求曲线 $y=\frac{1}{2}x^2$ 从 $x=0$ 到 $x=1$ 的弧长。
\par\textit{改编训练。}\par\hyperref[ans:c42q1]{解答见第~\pageref{ans:c42q1} 页。}
\end{example}
\begin{example}[抛物线弧长]\label{ex:c42q2}
求曲线 $y=1x^2$ 从 $x=0$ 到 $x=1$ 的弧长。
\par\textit{改编训练。}\par\hyperref[ans:c42q2]{解答见第~\pageref{ans:c42q2} 页。}
\end{example}
\begin{example}[抛物线弧长]\label{ex:c42q3}
求曲线 $y=\frac{3}{2}x^2$ 从 $x=0$ 到 $x=1$ 的弧长。
\par\textit{改编训练。}\par\hyperref[ans:c42q3]{解答见第~\pageref{ans:c42q3} 页。}
\end{example}
\begin{example}[抛物线弧长]\label{ex:c42q4}
求曲线 $y=2x^2$ 从 $x=0$ 到 $x=1$ 的弧长。
\par\textit{改编训练。}\par\hyperref[ans:c42q4]{解答见第~\pageref{ans:c42q4} 页。}
\end{example}
\begin{example}[抛物线弧长]\label{ex:c42q5}
求曲线 $y=\frac{5}{2}x^2$ 从 $x=0$ 到 $x=1$ 的弧长。
\par\textit{改编训练。}\par\hyperref[ans:c42q5]{解答见第~\pageref{ans:c42q5} 页。}
\end{example}
\begin{example}[抛物线弧长]\label{ex:c42q6}
求曲线 $y=3x^2$ 从 $x=0$ 到 $x=1$ 的弧长。
\par\textit{改编训练。}\par\hyperref[ans:c42q6]{解答见第~\pageref{ans:c42q6} 页。}
\end{example}
\begin{example}[平面片的向上通量]\label{ex:c42q7}
设 $S:z=1x+2y$，$0\le x,y\le1$，取向上法向。求 $F=(0,0,z)$ 的通量。
\par\textit{改编训练。}\par\hyperref[ans:c42q7]{解答见第~\pageref{ans:c42q7} 页。}
\end{example}
\begin{example}[平面片的向上通量]\label{ex:c42q8}
设 $S:z=2x+1y$，$0\le x,y\le1$，取向上法向。求 $F=(0,0,z)$ 的通量。
\par\textit{改编训练。}\par\hyperref[ans:c42q8]{解答见第~\pageref{ans:c42q8} 页。}
\end{example}
\begin{example}[平面片的向上通量]\label{ex:c42q9}
设 $S:z=2x+3y$，$0\le x,y\le1$，取向上法向。求 $F=(0,0,z)$ 的通量。
\par\textit{改编训练。}\par\hyperref[ans:c42q9]{解答见第~\pageref{ans:c42q9} 页。}
\end{example}
\begin{example}[平面片的向上通量]\label{ex:c42q10}
设 $S:z=3x+2y$，$0\le x,y\le1$，取向上法向。求 $F=(0,0,z)$ 的通量。
\par\textit{改编训练。}\par\hyperref[ans:c42q10]{解答见第~\pageref{ans:c42q10} 页。}
\end{example}
\begin{example}[平面片的向上通量]\label{ex:c42q11}
设 $S:z=4x+1y$，$0\le x,y\le1$，取向上法向。求 $F=(0,0,z)$ 的通量。
\par\textit{改编训练。}\par\hyperref[ans:c42q11]{解答见第~\pageref{ans:c42q11} 页。}
\end{example}
\begin{example}[平面片的向上通量]\label{ex:c42q12}
设 $S:z=3x+4y$，$0\le x,y\le1$，取向上法向。求 $F=(0,0,z)$ 的通量。
\par\textit{改编训练。}\par\hyperref[ans:c42q12]{解答见第~\pageref{ans:c42q12} 页。}
\end{example}
\begin{example}[上半球上的常向量通量]\label{ex:c42q13}
设 $S$ 为球面 $x^2+y^2+z^2=1$ 的上半部，取外法向。求 $F=(0,0,1)$ 的通量。
\par\textit{改编训练。}\par\hyperref[ans:c42q13]{解答见第~\pageref{ans:c42q13} 页。}
\end{example}
\begin{example}[上半球上的常向量通量]\label{ex:c42q14}
设 $S$ 为球面 $x^2+y^2+z^2=4$ 的上半部，取外法向。求 $F=(0,0,1)$ 的通量。
\par\textit{改编训练。}\par\hyperref[ans:c42q14]{解答见第~\pageref{ans:c42q14} 页。}
\end{example}
\begin{example}[上半球上的常向量通量]\label{ex:c42q15}
设 $S$ 为球面 $x^2+y^2+z^2=9$ 的上半部，取外法向。求 $F=(0,0,1)$ 的通量。
\par\textit{改编训练。}\par\hyperref[ans:c42q15]{解答见第~\pageref{ans:c42q15} 页。}
\end{example}
\begin{example}[上半球上的常向量通量]\label{ex:c42q16}
设 $S$ 为球面 $x^2+y^2+z^2=16$ 的上半部，取外法向。求 $F=(0,0,1)$ 的通量。
\par\textit{改编训练。}\par\hyperref[ans:c42q16]{解答见第~\pageref{ans:c42q16} 页。}
\end{example}
\begin{example}[上半球上的常向量通量]\label{ex:c42q17}
设 $S$ 为球面 $x^2+y^2+z^2=25$ 的上半部，取外法向。求 $F=(0,0,1)$ 的通量。
\par\textit{改编训练。}\par\hyperref[ans:c42q17]{解答见第~\pageref{ans:c42q17} 页。}
\end{example}
\begin{example}[上半球上的常向量通量]\label{ex:c42q18}
设 $S$ 为球面 $x^2+y^2+z^2=36$ 的上半部，取外法向。求 $F=(0,0,1)$ 的通量。
\par\textit{改编训练。}\par\hyperref[ans:c42q18]{解答见第~\pageref{ans:c42q18} 页。}
\end{example}
\begin{example}[两条路径给出不同功]\label{ex:c42q19}
设 $F=(y,x^2)$。从 $(0,0)$ 到 $(1,1)$，分别沿抛物线 $y=1x^2$ 和折线 $(0,0)\to(1,0)\to(1,1)$ 计算功。
\par\textit{改编训练。}\par\hyperref[ans:c42q19]{解答见第~\pageref{ans:c42q19} 页。}
\end{example}
\begin{example}[两条路径给出不同功]\label{ex:c42q20}
设 $F=(y,x^2)$。从 $(0,0)$ 到 $(1,2)$，分别沿抛物线 $y=2x^2$ 和折线 $(0,0)\to(1,0)\to(1,2)$ 计算功。
\par\textit{改编训练。}\par\hyperref[ans:c42q20]{解答见第~\pageref{ans:c42q20} 页。}
\end{example}
\begin{example}[两条路径给出不同功]\label{ex:c42q21}
设 $F=(y,x^2)$。从 $(0,0)$ 到 $(1,3)$，分别沿抛物线 $y=3x^2$ 和折线 $(0,0)\to(1,0)\to(1,3)$ 计算功。
\par\textit{改编训练。}\par\hyperref[ans:c42q21]{解答见第~\pageref{ans:c42q21} 页。}
\end{example}
\begin{example}[两条路径给出不同功]\label{ex:c42q22}
设 $F=(y,x^2)$。从 $(0,0)$ 到 $(1,4)$，分别沿抛物线 $y=4x^2$ 和折线 $(0,0)\to(1,0)\to(1,4)$ 计算功。
\par\textit{改编训练。}\par\hyperref[ans:c42q22]{解答见第~\pageref{ans:c42q22} 页。}
\end{example}
\begin{example}[两条路径给出不同功]\label{ex:c42q23}
设 $F=(y,x^2)$。从 $(0,0)$ 到 $(1,5)$，分别沿抛物线 $y=5x^2$ 和折线 $(0,0)\to(1,0)\to(1,5)$ 计算功。
\par\textit{改编训练。}\par\hyperref[ans:c42q23]{解答见第~\pageref{ans:c42q23} 页。}
\end{example}
\begin{example}[两条路径给出不同功]\label{ex:c42q24}
设 $F=(y,x^2)$。从 $(0,0)$ 到 $(1,6)$，分别沿抛物线 $y=6x^2$ 和折线 $(0,0)\to(1,0)\to(1,6)$ 计算功。
\par\textit{改编训练。}\par\hyperref[ans:c42q24]{解答见第~\pageref{ans:c42q24} 页。}
\end{example}
\subsection{解答}
\begin{solve}[习题~\ref{ex:c42q1}：抛物线弧长]\phantomsection\label{ans:c42q1}
图形曲线的弧长为 $L=\int_0^1\sqrt{1+(y^{\prime})^2}dx=\int_0^1\sqrt{1+1x^2}dx$。令 $u=1x$，使用原函数
\[
\int\sqrt{1+u^2}\,du=\frac12\left[u\sqrt{1+u^2}+\log(u+\sqrt{1+u^2})\right]+C,
\]
得到
\[
L=\frac{\sqrt{2}}2+\frac{1}{2}\log(1+\sqrt{2}).
\]
弧长至少等于两端点直线距离，这可以作为数值核验；所用对数真数在积分范围内为正。
\end{solve}

\begin{solve}[习题~\ref{ex:c42q2}：抛物线弧长]\phantomsection\label{ans:c42q2}
图形曲线的弧长为 $L=\int_0^1\sqrt{1+(y^{\prime})^2}dx=\int_0^1\sqrt{1+4x^2}dx$。令 $u=2x$，使用原函数
\[
\int\sqrt{1+u^2}\,du=\frac12\left[u\sqrt{1+u^2}+\log(u+\sqrt{1+u^2})\right]+C,
\]
得到
\[
L=\frac{\sqrt{5}}2+\frac{1}{4}\log(2+\sqrt{5}).
\]
弧长至少等于两端点直线距离，这可以作为数值核验；所用对数真数在积分范围内为正。
\end{solve}

\begin{solve}[习题~\ref{ex:c42q3}：抛物线弧长]\phantomsection\label{ans:c42q3}
图形曲线的弧长为 $L=\int_0^1\sqrt{1+(y^{\prime})^2}dx=\int_0^1\sqrt{1+9x^2}dx$。令 $u=3x$，使用原函数
\[
\int\sqrt{1+u^2}\,du=\frac12\left[u\sqrt{1+u^2}+\log(u+\sqrt{1+u^2})\right]+C,
\]
得到
\[
L=\frac{\sqrt{10}}2+\frac{1}{6}\log(3+\sqrt{10}).
\]
弧长至少等于两端点直线距离，这可以作为数值核验；所用对数真数在积分范围内为正。
\end{solve}

\begin{solve}[习题~\ref{ex:c42q4}：抛物线弧长]\phantomsection\label{ans:c42q4}
图形曲线的弧长为 $L=\int_0^1\sqrt{1+(y^{\prime})^2}dx=\int_0^1\sqrt{1+16x^2}dx$。令 $u=4x$，使用原函数
\[
\int\sqrt{1+u^2}\,du=\frac12\left[u\sqrt{1+u^2}+\log(u+\sqrt{1+u^2})\right]+C,
\]
得到
\[
L=\frac{\sqrt{17}}2+\frac{1}{8}\log(4+\sqrt{17}).
\]
弧长至少等于两端点直线距离，这可以作为数值核验；所用对数真数在积分范围内为正。
\end{solve}

\begin{solve}[习题~\ref{ex:c42q5}：抛物线弧长]\phantomsection\label{ans:c42q5}
图形曲线的弧长为 $L=\int_0^1\sqrt{1+(y^{\prime})^2}dx=\int_0^1\sqrt{1+25x^2}dx$。令 $u=5x$，使用原函数
\[
\int\sqrt{1+u^2}\,du=\frac12\left[u\sqrt{1+u^2}+\log(u+\sqrt{1+u^2})\right]+C,
\]
得到
\[
L=\frac{\sqrt{26}}2+\frac{1}{10}\log(5+\sqrt{26}).
\]
弧长至少等于两端点直线距离，这可以作为数值核验；所用对数真数在积分范围内为正。
\end{solve}

\begin{solve}[习题~\ref{ex:c42q6}：抛物线弧长]\phantomsection\label{ans:c42q6}
图形曲线的弧长为 $L=\int_0^1\sqrt{1+(y^{\prime})^2}dx=\int_0^1\sqrt{1+36x^2}dx$。令 $u=6x$，使用原函数
\[
\int\sqrt{1+u^2}\,du=\frac12\left[u\sqrt{1+u^2}+\log(u+\sqrt{1+u^2})\right]+C,
\]
得到
\[
L=\frac{\sqrt{37}}2+\frac{1}{12}\log(6+\sqrt{37}).
\]
弧长至少等于两端点直线距离，这可以作为数值核验；所用对数真数在积分范围内为正。
\end{solve}

\begin{solve}[习题~\ref{ex:c42q7}：平面片的向上通量]\phantomsection\label{ans:c42q7}
对于图形曲面，向上有向面积元为 $(-f_x,-f_y,1)dxdy=(-1,-2,1)dxdy$。代入向量场后点积为 $z=1x+2y$，于是
\[
\iint_SF\cdot n\,dS=\int_0^1\int_0^1(1x+2y)dxdy=\frac{3}{2}.
\]
这里已经把单位法向与面积因子相乘，不能再额外乘一次 $\sqrt{1+f_x^2+f_y^2}$。若使用单位法向，则必须另外配套使用实际面积元，二者最终相同。
\end{solve}

\begin{solve}[习题~\ref{ex:c42q8}：平面片的向上通量]\phantomsection\label{ans:c42q8}
对于图形曲面，向上有向面积元为 $(-f_x,-f_y,1)dxdy=(-2,-1,1)dxdy$。代入向量场后点积为 $z=2x+1y$，于是
\[
\iint_SF\cdot n\,dS=\int_0^1\int_0^1(2x+1y)dxdy=\frac{3}{2}.
\]
这里已经把单位法向与面积因子相乘，不能再额外乘一次 $\sqrt{1+f_x^2+f_y^2}$。若使用单位法向，则必须另外配套使用实际面积元，二者最终相同。
\end{solve}

\begin{solve}[习题~\ref{ex:c42q9}：平面片的向上通量]\phantomsection\label{ans:c42q9}
对于图形曲面，向上有向面积元为 $(-f_x,-f_y,1)dxdy=(-2,-3,1)dxdy$。代入向量场后点积为 $z=2x+3y$，于是
\[
\iint_SF\cdot n\,dS=\int_0^1\int_0^1(2x+3y)dxdy=\frac{5}{2}.
\]
这里已经把单位法向与面积因子相乘，不能再额外乘一次 $\sqrt{1+f_x^2+f_y^2}$。若使用单位法向，则必须另外配套使用实际面积元，二者最终相同。
\end{solve}

\begin{solve}[习题~\ref{ex:c42q10}：平面片的向上通量]\phantomsection\label{ans:c42q10}
对于图形曲面，向上有向面积元为 $(-f_x,-f_y,1)dxdy=(-3,-2,1)dxdy$。代入向量场后点积为 $z=3x+2y$，于是
\[
\iint_SF\cdot n\,dS=\int_0^1\int_0^1(3x+2y)dxdy=\frac{5}{2}.
\]
这里已经把单位法向与面积因子相乘，不能再额外乘一次 $\sqrt{1+f_x^2+f_y^2}$。若使用单位法向，则必须另外配套使用实际面积元，二者最终相同。
\end{solve}

\begin{solve}[习题~\ref{ex:c42q11}：平面片的向上通量]\phantomsection\label{ans:c42q11}
对于图形曲面，向上有向面积元为 $(-f_x,-f_y,1)dxdy=(-4,-1,1)dxdy$。代入向量场后点积为 $z=4x+1y$，于是
\[
\iint_SF\cdot n\,dS=\int_0^1\int_0^1(4x+1y)dxdy=\frac{5}{2}.
\]
这里已经把单位法向与面积因子相乘，不能再额外乘一次 $\sqrt{1+f_x^2+f_y^2}$。若使用单位法向，则必须另外配套使用实际面积元，二者最终相同。
\end{solve}

\begin{solve}[习题~\ref{ex:c42q12}：平面片的向上通量]\phantomsection\label{ans:c42q12}
对于图形曲面，向上有向面积元为 $(-f_x,-f_y,1)dxdy=(-3,-4,1)dxdy$。代入向量场后点积为 $z=3x+4y$，于是
\[
\iint_SF\cdot n\,dS=\int_0^1\int_0^1(3x+4y)dxdy=\frac{7}{2}.
\]
这里已经把单位法向与面积因子相乘，不能再额外乘一次 $\sqrt{1+f_x^2+f_y^2}$。若使用单位法向，则必须另外配套使用实际面积元，二者最终相同。
\end{solve}

\begin{solve}[习题~\ref{ex:c42q13}：上半球上的常向量通量]\phantomsection\label{ans:c42q13}
将上半球补上 $z=0$ 的底圆盘形成闭合边界。常向量场散度为0，所以两部分通量之和为0。底圆盘的外法向朝下，与 $F$ 点积为负1，其通量为 $-1\pi$。因此上半球通量为
\[
\iint_SF\cdot n\,dS=1\pi.
\]
它等于上半球在水平面上的投影面积，而不是半球本身的面积。法向倾斜带来的余弦权重正是投影面积关系的来源。
\end{solve}

\begin{solve}[习题~\ref{ex:c42q14}：上半球上的常向量通量]\phantomsection\label{ans:c42q14}
将上半球补上 $z=0$ 的底圆盘形成闭合边界。常向量场散度为0，所以两部分通量之和为0。底圆盘的外法向朝下，与 $F$ 点积为负1，其通量为 $-4\pi$。因此上半球通量为
\[
\iint_SF\cdot n\,dS=4\pi.
\]
它等于上半球在水平面上的投影面积，而不是半球本身的面积。法向倾斜带来的余弦权重正是投影面积关系的来源。
\end{solve}

\begin{solve}[习题~\ref{ex:c42q15}：上半球上的常向量通量]\phantomsection\label{ans:c42q15}
将上半球补上 $z=0$ 的底圆盘形成闭合边界。常向量场散度为0，所以两部分通量之和为0。底圆盘的外法向朝下，与 $F$ 点积为负1，其通量为 $-9\pi$。因此上半球通量为
\[
\iint_SF\cdot n\,dS=9\pi.
\]
它等于上半球在水平面上的投影面积，而不是半球本身的面积。法向倾斜带来的余弦权重正是投影面积关系的来源。
\end{solve}

\begin{solve}[习题~\ref{ex:c42q16}：上半球上的常向量通量]\phantomsection\label{ans:c42q16}
将上半球补上 $z=0$ 的底圆盘形成闭合边界。常向量场散度为0，所以两部分通量之和为0。底圆盘的外法向朝下，与 $F$ 点积为负1，其通量为 $-16\pi$。因此上半球通量为
\[
\iint_SF\cdot n\,dS=16\pi.
\]
它等于上半球在水平面上的投影面积，而不是半球本身的面积。法向倾斜带来的余弦权重正是投影面积关系的来源。
\end{solve}

\begin{solve}[习题~\ref{ex:c42q17}：上半球上的常向量通量]\phantomsection\label{ans:c42q17}
将上半球补上 $z=0$ 的底圆盘形成闭合边界。常向量场散度为0，所以两部分通量之和为0。底圆盘的外法向朝下，与 $F$ 点积为负1，其通量为 $-25\pi$。因此上半球通量为
\[
\iint_SF\cdot n\,dS=25\pi.
\]
它等于上半球在水平面上的投影面积，而不是半球本身的面积。法向倾斜带来的余弦权重正是投影面积关系的来源。
\end{solve}

\begin{solve}[习题~\ref{ex:c42q18}：上半球上的常向量通量]\phantomsection\label{ans:c42q18}
将上半球补上 $z=0$ 的底圆盘形成闭合边界。常向量场散度为0，所以两部分通量之和为0。底圆盘的外法向朝下，与 $F$ 点积为负1，其通量为 $-36\pi$。因此上半球通量为
\[
\iint_SF\cdot n\,dS=36\pi.
\]
它等于上半球在水平面上的投影面积，而不是半球本身的面积。法向倾斜带来的余弦权重正是投影面积关系的来源。
\end{solve}

\begin{solve}[习题~\ref{ex:c42q19}：两条路径给出不同功]\phantomsection\label{ans:c42q19}
抛物线上取 $x=t,y=1t^2$，$0\le t\le1$，则
\[
W_1=\int_0^1(1t^2+2t^3)dt=\frac{5}{6}.
\]
折线第一段 $y=0,dy=0$，贡献0；第二段 $x=1,dx=0$，贡献 $\int_0^{1}dy=1$，所以 $W_2=1$。两值不同，说明该场不保守。其 $Q_x-P_y=2x-1$ 不恒为0，也与这一结论一致。
\end{solve}

\begin{solve}[习题~\ref{ex:c42q20}：两条路径给出不同功]\phantomsection\label{ans:c42q20}
抛物线上取 $x=t,y=2t^2$，$0\le t\le1$，则
\[
W_1=\int_0^1(2t^2+4t^3)dt=\frac{5}{3}.
\]
折线第一段 $y=0,dy=0$，贡献0；第二段 $x=1,dx=0$，贡献 $\int_0^{2}dy=2$，所以 $W_2=2$。两值不同，说明该场不保守。其 $Q_x-P_y=2x-1$ 不恒为0，也与这一结论一致。
\end{solve}

\begin{solve}[习题~\ref{ex:c42q21}：两条路径给出不同功]\phantomsection\label{ans:c42q21}
抛物线上取 $x=t,y=3t^2$，$0\le t\le1$，则
\[
W_1=\int_0^1(3t^2+6t^3)dt=\frac{5}{2}.
\]
折线第一段 $y=0,dy=0$，贡献0；第二段 $x=1,dx=0$，贡献 $\int_0^{3}dy=3$，所以 $W_2=3$。两值不同，说明该场不保守。其 $Q_x-P_y=2x-1$ 不恒为0，也与这一结论一致。
\end{solve}

\begin{solve}[习题~\ref{ex:c42q22}：两条路径给出不同功]\phantomsection\label{ans:c42q22}
抛物线上取 $x=t,y=4t^2$，$0\le t\le1$，则
\[
W_1=\int_0^1(4t^2+8t^3)dt=\frac{10}{3}.
\]
折线第一段 $y=0,dy=0$，贡献0；第二段 $x=1,dx=0$，贡献 $\int_0^{4}dy=4$，所以 $W_2=4$。两值不同，说明该场不保守。其 $Q_x-P_y=2x-1$ 不恒为0，也与这一结论一致。
\end{solve}

\begin{solve}[习题~\ref{ex:c42q23}：两条路径给出不同功]\phantomsection\label{ans:c42q23}
抛物线上取 $x=t,y=5t^2$，$0\le t\le1$，则
\[
W_1=\int_0^1(5t^2+10t^3)dt=\frac{25}{6}.
\]
折线第一段 $y=0,dy=0$，贡献0；第二段 $x=1,dx=0$，贡献 $\int_0^{5}dy=5$，所以 $W_2=5$。两值不同，说明该场不保守。其 $Q_x-P_y=2x-1$ 不恒为0，也与这一结论一致。
\end{solve}

\begin{solve}[习题~\ref{ex:c42q24}：两条路径给出不同功]\phantomsection\label{ans:c42q24}
抛物线上取 $x=t,y=6t^2$，$0\le t\le1$，则
\[
W_1=\int_0^1(6t^2+12t^3)dt=5.
\]
折线第一段 $y=0,dy=0$，贡献0；第二段 $x=1,dx=0$，贡献 $\int_0^{6}dy=6$，所以 $W_2=6$。两值不同，说明该场不保守。其 $Q_x-P_y=2x-1$ 不恒为0，也与这一结论一致。
\end{solve}

\clearpage
\section{曲面积分与积分定理}\label{sec:ch43}
第一类曲面积分对面积元积分，第二类曲面积分计算通量并依赖法向。Gauss 散度定理适用于闭合曲面外法向；Stokes 定理连接有向曲面上的旋度通量与其边界的环流。

\textbf{资源对应：}CLP-4 \S~3.3、\S~4.1--\S~4.4。本章共 24 道编号习题。
\subsection{习题}
\begin{example}[图形曲面的面积]\label{ex:c43q1}
求平面片 $z=1x+2y$ 在 $0\le x\le1$、$0\le y\le2$ 上方的面积。
\par\textit{改编训练。}\par\hyperref[ans:c43q1]{解答见第~\pageref{ans:c43q1} 页。}
\end{example}
\begin{example}[图形曲面的面积]\label{ex:c43q2}
求平面片 $z=2x+1y$ 在 $0\le x\le1$、$0\le y\le2$ 上方的面积。
\par\textit{改编训练。}\par\hyperref[ans:c43q2]{解答见第~\pageref{ans:c43q2} 页。}
\end{example}
\begin{example}[图形曲面的面积]\label{ex:c43q3}
求平面片 $z=2x+3y$ 在 $0\le x\le1$、$0\le y\le2$ 上方的面积。
\par\textit{改编训练。}\par\hyperref[ans:c43q3]{解答见第~\pageref{ans:c43q3} 页。}
\end{example}
\begin{example}[图形曲面的面积]\label{ex:c43q4}
求平面片 $z=3x+2y$ 在 $0\le x\le1$、$0\le y\le2$ 上方的面积。
\par\textit{改编训练。}\par\hyperref[ans:c43q4]{解答见第~\pageref{ans:c43q4} 页。}
\end{example}
\begin{example}[图形曲面的面积]\label{ex:c43q5}
求平面片 $z=4x+1y$ 在 $0\le x\le1$、$0\le y\le2$ 上方的面积。
\par\textit{改编训练。}\par\hyperref[ans:c43q5]{解答见第~\pageref{ans:c43q5} 页。}
\end{example}
\begin{example}[图形曲面的面积]\label{ex:c43q6}
求平面片 $z=3x+4y$ 在 $0\le x\le1$、$0\le y\le2$ 上方的面积。
\par\textit{改编训练。}\par\hyperref[ans:c43q6]{解答见第~\pageref{ans:c43q6} 页。}
\end{example}
\begin{example}[球面的二次函数积分]\label{ex:c43q7}
设 $S:x^2+y^2+z^2=1$。求 $\displaystyle\iint_Sz^2dS$。
\par\textit{改编训练。}\par\hyperref[ans:c43q7]{解答见第~\pageref{ans:c43q7} 页。}
\end{example}
\begin{example}[球面的二次函数积分]\label{ex:c43q8}
设 $S:x^2+y^2+z^2=4$。求 $\displaystyle\iint_Sz^2dS$。
\par\textit{改编训练。}\par\hyperref[ans:c43q8]{解答见第~\pageref{ans:c43q8} 页。}
\end{example}
\begin{example}[球面的二次函数积分]\label{ex:c43q9}
设 $S:x^2+y^2+z^2=9$。求 $\displaystyle\iint_Sz^2dS$。
\par\textit{改编训练。}\par\hyperref[ans:c43q9]{解答见第~\pageref{ans:c43q9} 页。}
\end{example}
\begin{example}[球面的二次函数积分]\label{ex:c43q10}
设 $S:x^2+y^2+z^2=16$。求 $\displaystyle\iint_Sz^2dS$。
\par\textit{改编训练。}\par\hyperref[ans:c43q10]{解答见第~\pageref{ans:c43q10} 页。}
\end{example}
\begin{example}[球面的二次函数积分]\label{ex:c43q11}
设 $S:x^2+y^2+z^2=25$。求 $\displaystyle\iint_Sz^2dS$。
\par\textit{改编训练。}\par\hyperref[ans:c43q11]{解答见第~\pageref{ans:c43q11} 页。}
\end{example}
\begin{example}[球面的二次函数积分]\label{ex:c43q12}
设 $S:x^2+y^2+z^2=36$。求 $\displaystyle\iint_Sz^2dS$。
\par\textit{改编训练。}\par\hyperref[ans:c43q12]{解答见第~\pageref{ans:c43q12} 页。}
\end{example}
\begin{example}[散度定理计算球面通量]\label{ex:c43q13}
设 $F=(x,y,z)$，求穿过球面 $x^2+y^2+z^2=1$ 的向外通量。
\par\textit{改编训练。}\par\hyperref[ans:c43q13]{解答见第~\pageref{ans:c43q13} 页。}
\end{example}
\begin{example}[散度定理计算球面通量]\label{ex:c43q14}
设 $F=(x,y,z)$，求穿过球面 $x^2+y^2+z^2=4$ 的向外通量。
\par\textit{改编训练。}\par\hyperref[ans:c43q14]{解答见第~\pageref{ans:c43q14} 页。}
\end{example}
\begin{example}[散度定理计算球面通量]\label{ex:c43q15}
设 $F=(x,y,z)$，求穿过球面 $x^2+y^2+z^2=9$ 的向外通量。
\par\textit{改编训练。}\par\hyperref[ans:c43q15]{解答见第~\pageref{ans:c43q15} 页。}
\end{example}
\begin{example}[散度定理计算球面通量]\label{ex:c43q16}
设 $F=(x,y,z)$，求穿过球面 $x^2+y^2+z^2=16$ 的向外通量。
\par\textit{改编训练。}\par\hyperref[ans:c43q16]{解答见第~\pageref{ans:c43q16} 页。}
\end{example}
\begin{example}[散度定理计算球面通量]\label{ex:c43q17}
设 $F=(x,y,z)$，求穿过球面 $x^2+y^2+z^2=25$ 的向外通量。
\par\textit{改编训练。}\par\hyperref[ans:c43q17]{解答见第~\pageref{ans:c43q17} 页。}
\end{example}
\begin{example}[散度定理计算球面通量]\label{ex:c43q18}
设 $F=(x,y,z)$，求穿过球面 $x^2+y^2+z^2=36$ 的向外通量。
\par\textit{改编训练。}\par\hyperref[ans:c43q18]{解答见第~\pageref{ans:c43q18} 页。}
\end{example}
\begin{example}[Stokes 定理与平面圆盘]\label{ex:c43q19}
设 $F=(-1y,1x,2z)$。沿圆 $x^2+y^2=1,z=0$ 从正 $z$ 轴看逆时针计算环流，并用 Stokes 定理验证。
\par\textit{改编训练。}\par\hyperref[ans:c43q19]{解答见第~\pageref{ans:c43q19} 页。}
\end{example}
\begin{example}[Stokes 定理与平面圆盘]\label{ex:c43q20}
设 $F=(-2y,2x,1z)$。沿圆 $x^2+y^2=1,z=0$ 从正 $z$ 轴看逆时针计算环流，并用 Stokes 定理验证。
\par\textit{改编训练。}\par\hyperref[ans:c43q20]{解答见第~\pageref{ans:c43q20} 页。}
\end{example}
\begin{example}[Stokes 定理与平面圆盘]\label{ex:c43q21}
设 $F=(-2y,2x,3z)$。沿圆 $x^2+y^2=1,z=0$ 从正 $z$ 轴看逆时针计算环流，并用 Stokes 定理验证。
\par\textit{改编训练。}\par\hyperref[ans:c43q21]{解答见第~\pageref{ans:c43q21} 页。}
\end{example}
\begin{example}[Stokes 定理与平面圆盘]\label{ex:c43q22}
设 $F=(-3y,3x,2z)$。沿圆 $x^2+y^2=1,z=0$ 从正 $z$ 轴看逆时针计算环流，并用 Stokes 定理验证。
\par\textit{改编训练。}\par\hyperref[ans:c43q22]{解答见第~\pageref{ans:c43q22} 页。}
\end{example}
\begin{example}[Stokes 定理与平面圆盘]\label{ex:c43q23}
设 $F=(-4y,4x,1z)$。沿圆 $x^2+y^2=1,z=0$ 从正 $z$ 轴看逆时针计算环流，并用 Stokes 定理验证。
\par\textit{改编训练。}\par\hyperref[ans:c43q23]{解答见第~\pageref{ans:c43q23} 页。}
\end{example}
\begin{example}[Stokes 定理与平面圆盘]\label{ex:c43q24}
设 $F=(-3y,3x,4z)$。沿圆 $x^2+y^2=1,z=0$ 从正 $z$ 轴看逆时针计算环流，并用 Stokes 定理验证。
\par\textit{改编训练。}\par\hyperref[ans:c43q24]{解答见第~\pageref{ans:c43q24} 页。}
\end{example}
\subsection{解答}
\begin{solve}[习题~\ref{ex:c43q1}：图形曲面的面积]\phantomsection\label{ans:c43q1}
将曲面写成 $z=f(x,y)$，有 $f_x=1$、$f_y=2$。因此面积元为
\[
dS=\sqrt{1+f_x^2+f_y^2}dxdy=\sqrt{6}dxdy.
\]
投影区域面积为2，故曲面面积为 $2\sqrt{6}$。也可用参数向量 $r_x=(1,0,1)$、$r_y=(0,1,2)$ 的叉积模得到同样的面积因子。实际曲面比平面投影大，这与因子至少为1相符。
\end{solve}

\begin{solve}[习题~\ref{ex:c43q2}：图形曲面的面积]\phantomsection\label{ans:c43q2}
将曲面写成 $z=f(x,y)$，有 $f_x=2$、$f_y=1$。因此面积元为
\[
dS=\sqrt{1+f_x^2+f_y^2}dxdy=\sqrt{6}dxdy.
\]
投影区域面积为2，故曲面面积为 $2\sqrt{6}$。也可用参数向量 $r_x=(1,0,2)$、$r_y=(0,1,1)$ 的叉积模得到同样的面积因子。实际曲面比平面投影大，这与因子至少为1相符。
\end{solve}

\begin{solve}[习题~\ref{ex:c43q3}：图形曲面的面积]\phantomsection\label{ans:c43q3}
将曲面写成 $z=f(x,y)$，有 $f_x=2$、$f_y=3$。因此面积元为
\[
dS=\sqrt{1+f_x^2+f_y^2}dxdy=\sqrt{14}dxdy.
\]
投影区域面积为2，故曲面面积为 $2\sqrt{14}$。也可用参数向量 $r_x=(1,0,2)$、$r_y=(0,1,3)$ 的叉积模得到同样的面积因子。实际曲面比平面投影大，这与因子至少为1相符。
\end{solve}

\begin{solve}[习题~\ref{ex:c43q4}：图形曲面的面积]\phantomsection\label{ans:c43q4}
将曲面写成 $z=f(x,y)$，有 $f_x=3$、$f_y=2$。因此面积元为
\[
dS=\sqrt{1+f_x^2+f_y^2}dxdy=\sqrt{14}dxdy.
\]
投影区域面积为2，故曲面面积为 $2\sqrt{14}$。也可用参数向量 $r_x=(1,0,3)$、$r_y=(0,1,2)$ 的叉积模得到同样的面积因子。实际曲面比平面投影大，这与因子至少为1相符。
\end{solve}

\begin{solve}[习题~\ref{ex:c43q5}：图形曲面的面积]\phantomsection\label{ans:c43q5}
将曲面写成 $z=f(x,y)$，有 $f_x=4$、$f_y=1$。因此面积元为
\[
dS=\sqrt{1+f_x^2+f_y^2}dxdy=\sqrt{18}dxdy.
\]
投影区域面积为2，故曲面面积为 $2\sqrt{18}$。也可用参数向量 $r_x=(1,0,4)$、$r_y=(0,1,1)$ 的叉积模得到同样的面积因子。实际曲面比平面投影大，这与因子至少为1相符。
\end{solve}

\begin{solve}[习题~\ref{ex:c43q6}：图形曲面的面积]\phantomsection\label{ans:c43q6}
将曲面写成 $z=f(x,y)$，有 $f_x=3$、$f_y=4$。因此面积元为
\[
dS=\sqrt{1+f_x^2+f_y^2}dxdy=\sqrt{26}dxdy.
\]
投影区域面积为2，故曲面面积为 $2\sqrt{26}$。也可用参数向量 $r_x=(1,0,3)$、$r_y=(0,1,4)$ 的叉积模得到同样的面积因子。实际曲面比平面投影大，这与因子至少为1相符。
\end{solve}

\begin{solve}[习题~\ref{ex:c43q7}：球面的二次函数积分]\phantomsection\label{ans:c43q7}
由球面的旋转对称性，三个坐标平方的面积积分相等；而它们的和在球面上是常数。因此
\[
3\iint_Sz^2dS=\iint_S(x^2+y^2+z^2)dS=1\cdot4\pi1.
\]
所以积分为 $\frac{4}{3}\pi$。若用球面参数化，也可写为
\[
1\int_0^{2\pi}d\theta\int_0^\pi\cos^2\varphi\sin\varphi\,d\varphi=\frac{4}{3}\pi.
\]
第一类曲面积分与法向选择无关。对称法成立的前提是整个积分域具有对应的对称性，不能随意用于只有部分球面的区域。
\end{solve}

\begin{solve}[习题~\ref{ex:c43q8}：球面的二次函数积分]\phantomsection\label{ans:c43q8}
由球面的旋转对称性，三个坐标平方的面积积分相等；而它们的和在球面上是常数。因此
\[
3\iint_Sz^2dS=\iint_S(x^2+y^2+z^2)dS=4\cdot4\pi4.
\]
所以积分为 $\frac{64}{3}\pi$。若用球面参数化，也可写为
\[
16\int_0^{2\pi}d\theta\int_0^\pi\cos^2\varphi\sin\varphi\,d\varphi=\frac{64}{3}\pi.
\]
第一类曲面积分与法向选择无关。对称法成立的前提是整个积分域具有对应的对称性，不能随意用于只有部分球面的区域。
\end{solve}

\begin{solve}[习题~\ref{ex:c43q9}：球面的二次函数积分]\phantomsection\label{ans:c43q9}
由球面的旋转对称性，三个坐标平方的面积积分相等；而它们的和在球面上是常数。因此
\[
3\iint_Sz^2dS=\iint_S(x^2+y^2+z^2)dS=9\cdot4\pi9.
\]
所以积分为 $108\pi$。若用球面参数化，也可写为
\[
81\int_0^{2\pi}d\theta\int_0^\pi\cos^2\varphi\sin\varphi\,d\varphi=108\pi.
\]
第一类曲面积分与法向选择无关。对称法成立的前提是整个积分域具有对应的对称性，不能随意用于只有部分球面的区域。
\end{solve}

\begin{solve}[习题~\ref{ex:c43q10}：球面的二次函数积分]\phantomsection\label{ans:c43q10}
由球面的旋转对称性，三个坐标平方的面积积分相等；而它们的和在球面上是常数。因此
\[
3\iint_Sz^2dS=\iint_S(x^2+y^2+z^2)dS=16\cdot4\pi16.
\]
所以积分为 $\frac{1024}{3}\pi$。若用球面参数化，也可写为
\[
256\int_0^{2\pi}d\theta\int_0^\pi\cos^2\varphi\sin\varphi\,d\varphi=\frac{1024}{3}\pi.
\]
第一类曲面积分与法向选择无关。对称法成立的前提是整个积分域具有对应的对称性，不能随意用于只有部分球面的区域。
\end{solve}

\begin{solve}[习题~\ref{ex:c43q11}：球面的二次函数积分]\phantomsection\label{ans:c43q11}
由球面的旋转对称性，三个坐标平方的面积积分相等；而它们的和在球面上是常数。因此
\[
3\iint_Sz^2dS=\iint_S(x^2+y^2+z^2)dS=25\cdot4\pi25.
\]
所以积分为 $\frac{2500}{3}\pi$。若用球面参数化，也可写为
\[
625\int_0^{2\pi}d\theta\int_0^\pi\cos^2\varphi\sin\varphi\,d\varphi=\frac{2500}{3}\pi.
\]
第一类曲面积分与法向选择无关。对称法成立的前提是整个积分域具有对应的对称性，不能随意用于只有部分球面的区域。
\end{solve}

\begin{solve}[习题~\ref{ex:c43q12}：球面的二次函数积分]\phantomsection\label{ans:c43q12}
由球面的旋转对称性，三个坐标平方的面积积分相等；而它们的和在球面上是常数。因此
\[
3\iint_Sz^2dS=\iint_S(x^2+y^2+z^2)dS=36\cdot4\pi36.
\]
所以积分为 $1728\pi$。若用球面参数化，也可写为
\[
1296\int_0^{2\pi}d\theta\int_0^\pi\cos^2\varphi\sin\varphi\,d\varphi=1728\pi.
\]
第一类曲面积分与法向选择无关。对称法成立的前提是整个积分域具有对应的对称性，不能随意用于只有部分球面的区域。
\end{solve}

\begin{solve}[习题~\ref{ex:c43q13}：散度定理计算球面通量]\phantomsection\label{ans:c43q13}
向量场在闭球及其邻域内有连续一阶偏导，球面光滑闭合，因此可使用 Gauss 散度定理。散度为3，所以
\[
\iint_SF\cdot n\,dS=\iiint_B3\,dV=3\cdot\frac{4\pi1}3=4\pi.
\]
直接验证时，外单位法向为 $(x,y,z)/1$，故球面上 $F\cdot n=1$，乘球面积 $4\pi$ 也得到相同结果。若选择内法向，通量符号相反；闭合曲面方向不是可以省略的附注。
\end{solve}

\begin{solve}[习题~\ref{ex:c43q14}：散度定理计算球面通量]\phantomsection\label{ans:c43q14}
向量场在闭球及其邻域内有连续一阶偏导，球面光滑闭合，因此可使用 Gauss 散度定理。散度为3，所以
\[
\iint_SF\cdot n\,dS=\iiint_B3\,dV=3\cdot\frac{4\pi8}3=32\pi.
\]
直接验证时，外单位法向为 $(x,y,z)/2$，故球面上 $F\cdot n=2$，乘球面积 $16\pi$ 也得到相同结果。若选择内法向，通量符号相反；闭合曲面方向不是可以省略的附注。
\end{solve}

\begin{solve}[习题~\ref{ex:c43q15}：散度定理计算球面通量]\phantomsection\label{ans:c43q15}
向量场在闭球及其邻域内有连续一阶偏导，球面光滑闭合，因此可使用 Gauss 散度定理。散度为3，所以
\[
\iint_SF\cdot n\,dS=\iiint_B3\,dV=3\cdot\frac{4\pi27}3=108\pi.
\]
直接验证时，外单位法向为 $(x,y,z)/3$，故球面上 $F\cdot n=3$，乘球面积 $36\pi$ 也得到相同结果。若选择内法向，通量符号相反；闭合曲面方向不是可以省略的附注。
\end{solve}

\begin{solve}[习题~\ref{ex:c43q16}：散度定理计算球面通量]\phantomsection\label{ans:c43q16}
向量场在闭球及其邻域内有连续一阶偏导，球面光滑闭合，因此可使用 Gauss 散度定理。散度为3，所以
\[
\iint_SF\cdot n\,dS=\iiint_B3\,dV=3\cdot\frac{4\pi64}3=256\pi.
\]
直接验证时，外单位法向为 $(x,y,z)/4$，故球面上 $F\cdot n=4$，乘球面积 $64\pi$ 也得到相同结果。若选择内法向，通量符号相反；闭合曲面方向不是可以省略的附注。
\end{solve}

\begin{solve}[习题~\ref{ex:c43q17}：散度定理计算球面通量]\phantomsection\label{ans:c43q17}
向量场在闭球及其邻域内有连续一阶偏导，球面光滑闭合，因此可使用 Gauss 散度定理。散度为3，所以
\[
\iint_SF\cdot n\,dS=\iiint_B3\,dV=3\cdot\frac{4\pi125}3=500\pi.
\]
直接验证时，外单位法向为 $(x,y,z)/5$，故球面上 $F\cdot n=5$，乘球面积 $100\pi$ 也得到相同结果。若选择内法向，通量符号相反；闭合曲面方向不是可以省略的附注。
\end{solve}

\begin{solve}[习题~\ref{ex:c43q18}：散度定理计算球面通量]\phantomsection\label{ans:c43q18}
向量场在闭球及其邻域内有连续一阶偏导，球面光滑闭合，因此可使用 Gauss 散度定理。散度为3，所以
\[
\iint_SF\cdot n\,dS=\iiint_B3\,dV=3\cdot\frac{4\pi216}3=864\pi.
\]
直接验证时，外单位法向为 $(x,y,z)/6$，故球面上 $F\cdot n=6$，乘球面积 $144\pi$ 也得到相同结果。若选择内法向，通量符号相反；闭合曲面方向不是可以省略的附注。
\end{solve}

\begin{solve}[习题~\ref{ex:c43q19}：Stokes 定理与平面圆盘]\phantomsection\label{ans:c43q19}
取 $r(t)=(\cos t,\sin t,0)$，则 $F(r(t))\cdot r^{\prime}(t)=1$，直接积分得 $2\pi$。另一方面，
\[
\nabla\times F=(0,0,2).
\]
选择平面单位圆盘，向上法向与边界的逆时针方向满足右手规则，故
\[
\oint_C F\cdot dr=\iint_D(\nabla\times F)\cdot(0,0,1)dS=2\pi.
\]
第三分量 $2z$ 没有影响该结果，因为它在边界上为0，且不贡献本题的旋度。
\end{solve}

\begin{solve}[习题~\ref{ex:c43q20}：Stokes 定理与平面圆盘]\phantomsection\label{ans:c43q20}
取 $r(t)=(\cos t,\sin t,0)$，则 $F(r(t))\cdot r^{\prime}(t)=2$，直接积分得 $4\pi$。另一方面，
\[
\nabla\times F=(0,0,4).
\]
选择平面单位圆盘，向上法向与边界的逆时针方向满足右手规则，故
\[
\oint_C F\cdot dr=\iint_D(\nabla\times F)\cdot(0,0,1)dS=4\pi.
\]
第三分量 $1z$ 没有影响该结果，因为它在边界上为0，且不贡献本题的旋度。
\end{solve}

\begin{solve}[习题~\ref{ex:c43q21}：Stokes 定理与平面圆盘]\phantomsection\label{ans:c43q21}
取 $r(t)=(\cos t,\sin t,0)$，则 $F(r(t))\cdot r^{\prime}(t)=2$，直接积分得 $4\pi$。另一方面，
\[
\nabla\times F=(0,0,4).
\]
选择平面单位圆盘，向上法向与边界的逆时针方向满足右手规则，故
\[
\oint_C F\cdot dr=\iint_D(\nabla\times F)\cdot(0,0,1)dS=4\pi.
\]
第三分量 $3z$ 没有影响该结果，因为它在边界上为0，且不贡献本题的旋度。
\end{solve}

\begin{solve}[习题~\ref{ex:c43q22}：Stokes 定理与平面圆盘]\phantomsection\label{ans:c43q22}
取 $r(t)=(\cos t,\sin t,0)$，则 $F(r(t))\cdot r^{\prime}(t)=3$，直接积分得 $6\pi$。另一方面，
\[
\nabla\times F=(0,0,6).
\]
选择平面单位圆盘，向上法向与边界的逆时针方向满足右手规则，故
\[
\oint_C F\cdot dr=\iint_D(\nabla\times F)\cdot(0,0,1)dS=6\pi.
\]
第三分量 $2z$ 没有影响该结果，因为它在边界上为0，且不贡献本题的旋度。
\end{solve}

\begin{solve}[习题~\ref{ex:c43q23}：Stokes 定理与平面圆盘]\phantomsection\label{ans:c43q23}
取 $r(t)=(\cos t,\sin t,0)$，则 $F(r(t))\cdot r^{\prime}(t)=4$，直接积分得 $8\pi$。另一方面，
\[
\nabla\times F=(0,0,8).
\]
选择平面单位圆盘，向上法向与边界的逆时针方向满足右手规则，故
\[
\oint_C F\cdot dr=\iint_D(\nabla\times F)\cdot(0,0,1)dS=8\pi.
\]
第三分量 $1z$ 没有影响该结果，因为它在边界上为0，且不贡献本题的旋度。
\end{solve}

\begin{solve}[习题~\ref{ex:c43q24}：Stokes 定理与平面圆盘]\phantomsection\label{ans:c43q24}
取 $r(t)=(\cos t,\sin t,0)$，则 $F(r(t))\cdot r^{\prime}(t)=3$，直接积分得 $6\pi$。另一方面，
\[
\nabla\times F=(0,0,6).
\]
选择平面单位圆盘，向上法向与边界的逆时针方向满足右手规则，故
\[
\oint_C F\cdot dr=\iint_D(\nabla\times F)\cdot(0,0,1)dS=6\pi.
\]
第三分量 $4z$ 没有影响该结果，因为它在边界上为0，且不贡献本题的旋度。
\end{solve}

\clearpage
\section{散度、旋度与通量综合}\label{sec:ch44}
散度（divergence）是向量场局部源汇强度，旋度（curl）描述局部旋转。势函数、散度定理和 Stokes 定理各有不同假设；计算前先决定所求的是线积分、面积通量还是体积分。

\textbf{资源对应：}CLP-4 \S~2.3、\S~4.1--\S~4.4。本章共 24 道编号习题。
\subsection{习题}
\begin{example}[立方体上的多项式通量]\label{ex:c44q1}
设 $F=(1x^2,2y^2,3z^2)$，求其散度、旋度和单位立方体边界的外向通量。
\par\textit{改编训练。}\par\hyperref[ans:c44q1]{解答见第~\pageref{ans:c44q1} 页。}
\end{example}
\begin{example}[立方体上的多项式通量]\label{ex:c44q2}
设 $F=(2x^2,1y^2,2z^2)$，求其散度、旋度和单位立方体边界的外向通量。
\par\textit{改编训练。}\par\hyperref[ans:c44q2]{解答见第~\pageref{ans:c44q2} 页。}
\end{example}
\begin{example}[立方体上的多项式通量]\label{ex:c44q3}
设 $F=(3x^2,2y^2,1z^2)$，求其散度、旋度和单位立方体边界的外向通量。
\par\textit{改编训练。}\par\hyperref[ans:c44q3]{解答见第~\pageref{ans:c44q3} 页。}
\end{example}
\begin{example}[立方体上的多项式通量]\label{ex:c44q4}
设 $F=(2x^2,3y^2,4z^2)$，求其散度、旋度和单位立方体边界的外向通量。
\par\textit{改编训练。}\par\hyperref[ans:c44q4]{解答见第~\pageref{ans:c44q4} 页。}
\end{example}
\begin{example}[立方体上的多项式通量]\label{ex:c44q5}
设 $F=(4x^2,2y^2,3z^2)$，求其散度、旋度和单位立方体边界的外向通量。
\par\textit{改编训练。}\par\hyperref[ans:c44q5]{解答见第~\pageref{ans:c44q5} 页。}
\end{example}
\begin{example}[立方体上的多项式通量]\label{ex:c44q6}
设 $F=(3x^2,4y^2,2z^2)$，求其散度、旋度和单位立方体边界的外向通量。
\par\textit{改编训练。}\par\hyperref[ans:c44q6]{解答见第~\pageref{ans:c44q6} 页。}
\end{example}
\begin{example}[一个非零旋度的场]\label{ex:c44q7}
设 $F=(-y^{1},x^{1},0)$。求散度与旋度，并求单位正方形 $[0,1]^2$ 边界逆时针环流。
\par\textit{改编训练。}\par\hyperref[ans:c44q7]{解答见第~\pageref{ans:c44q7} 页。}
\end{example}
\begin{example}[一个非零旋度的场]\label{ex:c44q8}
设 $F=(-y^{2},x^{2},0)$。求散度与旋度，并求单位正方形 $[0,1]^2$ 边界逆时针环流。
\par\textit{改编训练。}\par\hyperref[ans:c44q8]{解答见第~\pageref{ans:c44q8} 页。}
\end{example}
\begin{example}[一个非零旋度的场]\label{ex:c44q9}
设 $F=(-y^{3},x^{3},0)$。求散度与旋度，并求单位正方形 $[0,1]^2$ 边界逆时针环流。
\par\textit{改编训练。}\par\hyperref[ans:c44q9]{解答见第~\pageref{ans:c44q9} 页。}
\end{example}
\begin{example}[一个非零旋度的场]\label{ex:c44q10}
设 $F=(-y^{4},x^{4},0)$。求散度与旋度，并求单位正方形 $[0,1]^2$ 边界逆时针环流。
\par\textit{改编训练。}\par\hyperref[ans:c44q10]{解答见第~\pageref{ans:c44q10} 页。}
\end{example}
\begin{example}[一个非零旋度的场]\label{ex:c44q11}
设 $F=(-y^{5},x^{5},0)$。求散度与旋度，并求单位正方形 $[0,1]^2$ 边界逆时针环流。
\par\textit{改编训练。}\par\hyperref[ans:c44q11]{解答见第~\pageref{ans:c44q11} 页。}
\end{example}
\begin{example}[一个非零旋度的场]\label{ex:c44q12}
设 $F=(-y^{6},x^{6},0)$。求散度与旋度，并求单位正方形 $[0,1]^2$ 边界逆时针环流。
\par\textit{改编训练。}\par\hyperref[ans:c44q12]{解答见第~\pageref{ans:c44q12} 页。}
\end{example}
\begin{example}[势函数给出的空间路径积分]\label{ex:c44q13}
设 $\varphi(x,y,z)=x^2y+1yz^2+2x$，$F=\nabla\varphi$。写出 $F$，计算从 $(0,0,0)$ 到 $(1,2,1)$ 的路径积分。
\par\textit{改编训练。}\par\hyperref[ans:c44q13]{解答见第~\pageref{ans:c44q13} 页。}
\end{example}
\begin{example}[势函数给出的空间路径积分]\label{ex:c44q14}
设 $\varphi(x,y,z)=x^2y+2yz^2+1x$，$F=\nabla\varphi$。写出 $F$，计算从 $(0,0,0)$ 到 $(1,2,1)$ 的路径积分。
\par\textit{改编训练。}\par\hyperref[ans:c44q14]{解答见第~\pageref{ans:c44q14} 页。}
\end{example}
\begin{example}[势函数给出的空间路径积分]\label{ex:c44q15}
设 $\varphi(x,y,z)=x^2y+2yz^2+3x$，$F=\nabla\varphi$。写出 $F$，计算从 $(0,0,0)$ 到 $(1,2,1)$ 的路径积分。
\par\textit{改编训练。}\par\hyperref[ans:c44q15]{解答见第~\pageref{ans:c44q15} 页。}
\end{example}
\begin{example}[势函数给出的空间路径积分]\label{ex:c44q16}
设 $\varphi(x,y,z)=x^2y+3yz^2+2x$，$F=\nabla\varphi$。写出 $F$，计算从 $(0,0,0)$ 到 $(1,2,1)$ 的路径积分。
\par\textit{改编训练。}\par\hyperref[ans:c44q16]{解答见第~\pageref{ans:c44q16} 页。}
\end{example}
\begin{example}[势函数给出的空间路径积分]\label{ex:c44q17}
设 $\varphi(x,y,z)=x^2y+4yz^2+1x$，$F=\nabla\varphi$。写出 $F$，计算从 $(0,0,0)$ 到 $(1,2,1)$ 的路径积分。
\par\textit{改编训练。}\par\hyperref[ans:c44q17]{解答见第~\pageref{ans:c44q17} 页。}
\end{example}
\begin{example}[势函数给出的空间路径积分]\label{ex:c44q18}
设 $\varphi(x,y,z)=x^2y+3yz^2+4x$，$F=\nabla\varphi$。写出 $F$，计算从 $(0,0,0)$ 到 $(1,2,1)$ 的路径积分。
\par\textit{改编训练。}\par\hyperref[ans:c44q18]{解答见第~\pageref{ans:c44q18} 页。}
\end{example}
\begin{example}[封闭曲面与缺失底面]\label{ex:c44q19}
抛物面 $S:z=1-x^2-y^2$ 位于 $z\ge0$，取向上法向。求 $F=(0,0,z)$ 穿过 $S$ 的通量。
\par\textit{改编训练。}\par\hyperref[ans:c44q19]{解答见第~\pageref{ans:c44q19} 页。}
\end{example}
\begin{example}[封闭曲面与缺失底面]\label{ex:c44q20}
抛物面 $S:z=2-x^2-y^2$ 位于 $z\ge0$，取向上法向。求 $F=(0,0,z)$ 穿过 $S$ 的通量。
\par\textit{改编训练。}\par\hyperref[ans:c44q20]{解答见第~\pageref{ans:c44q20} 页。}
\end{example}
\begin{example}[封闭曲面与缺失底面]\label{ex:c44q21}
抛物面 $S:z=3-x^2-y^2$ 位于 $z\ge0$，取向上法向。求 $F=(0,0,z)$ 穿过 $S$ 的通量。
\par\textit{改编训练。}\par\hyperref[ans:c44q21]{解答见第~\pageref{ans:c44q21} 页。}
\end{example}
\begin{example}[封闭曲面与缺失底面]\label{ex:c44q22}
抛物面 $S:z=4-x^2-y^2$ 位于 $z\ge0$，取向上法向。求 $F=(0,0,z)$ 穿过 $S$ 的通量。
\par\textit{改编训练。}\par\hyperref[ans:c44q22]{解答见第~\pageref{ans:c44q22} 页。}
\end{example}
\begin{example}[封闭曲面与缺失底面]\label{ex:c44q23}
抛物面 $S:z=5-x^2-y^2$ 位于 $z\ge0$，取向上法向。求 $F=(0,0,z)$ 穿过 $S$ 的通量。
\par\textit{改编训练。}\par\hyperref[ans:c44q23]{解答见第~\pageref{ans:c44q23} 页。}
\end{example}
\begin{example}[封闭曲面与缺失底面]\label{ex:c44q24}
抛物面 $S:z=6-x^2-y^2$ 位于 $z\ge0$，取向上法向。求 $F=(0,0,z)$ 穿过 $S$ 的通量。
\par\textit{改编训练。}\par\hyperref[ans:c44q24]{解答见第~\pageref{ans:c44q24} 页。}
\end{example}
\subsection{解答}
\begin{solve}[习题~\ref{ex:c44q1}：立方体上的多项式通量]\phantomsection\label{ans:c44q1}
各分量只依赖对应坐标，故旋度为零，散度为
\[
\nabla\cdot F=2x+4y+6z.
\]
在单位立方体内积分得到通量 $6$。直接按六个面计算时，$x=0,y=0,z=0$ 三面通量为0，另外三面分别贡献 $1,2,3$，也得到相同结果。旋度为零并不意味着通量为零；它控制的是局部环流结构，与散度描述的量不同。
\end{solve}

\begin{solve}[习题~\ref{ex:c44q2}：立方体上的多项式通量]\phantomsection\label{ans:c44q2}
各分量只依赖对应坐标，故旋度为零，散度为
\[
\nabla\cdot F=4x+2y+4z.
\]
在单位立方体内积分得到通量 $5$。直接按六个面计算时，$x=0,y=0,z=0$ 三面通量为0，另外三面分别贡献 $2,1,2$，也得到相同结果。旋度为零并不意味着通量为零；它控制的是局部环流结构，与散度描述的量不同。
\end{solve}

\begin{solve}[习题~\ref{ex:c44q3}：立方体上的多项式通量]\phantomsection\label{ans:c44q3}
各分量只依赖对应坐标，故旋度为零，散度为
\[
\nabla\cdot F=6x+4y+2z.
\]
在单位立方体内积分得到通量 $6$。直接按六个面计算时，$x=0,y=0,z=0$ 三面通量为0，另外三面分别贡献 $3,2,1$，也得到相同结果。旋度为零并不意味着通量为零；它控制的是局部环流结构，与散度描述的量不同。
\end{solve}

\begin{solve}[习题~\ref{ex:c44q4}：立方体上的多项式通量]\phantomsection\label{ans:c44q4}
各分量只依赖对应坐标，故旋度为零，散度为
\[
\nabla\cdot F=4x+6y+8z.
\]
在单位立方体内积分得到通量 $9$。直接按六个面计算时，$x=0,y=0,z=0$ 三面通量为0，另外三面分别贡献 $2,3,4$，也得到相同结果。旋度为零并不意味着通量为零；它控制的是局部环流结构，与散度描述的量不同。
\end{solve}

\begin{solve}[习题~\ref{ex:c44q5}：立方体上的多项式通量]\phantomsection\label{ans:c44q5}
各分量只依赖对应坐标，故旋度为零，散度为
\[
\nabla\cdot F=8x+4y+6z.
\]
在单位立方体内积分得到通量 $9$。直接按六个面计算时，$x=0,y=0,z=0$ 三面通量为0，另外三面分别贡献 $4,2,3$，也得到相同结果。旋度为零并不意味着通量为零；它控制的是局部环流结构，与散度描述的量不同。
\end{solve}

\begin{solve}[习题~\ref{ex:c44q6}：立方体上的多项式通量]\phantomsection\label{ans:c44q6}
各分量只依赖对应坐标，故旋度为零，散度为
\[
\nabla\cdot F=6x+8y+4z.
\]
在单位立方体内积分得到通量 $9$。直接按六个面计算时，$x=0,y=0,z=0$ 三面通量为0，另外三面分别贡献 $3,4,2$，也得到相同结果。旋度为零并不意味着通量为零；它控制的是局部环流结构，与散度描述的量不同。
\end{solve}

\begin{solve}[习题~\ref{ex:c44q7}：一个非零旋度的场]\phantomsection\label{ans:c44q7}
因第一分量与 $x$ 无关、第二分量与 $y$ 无关，散度为0。旋度为
\[
\nabla\times F=(0,0,1x^{0}+1y^{0}).
\]
由 Green 公式，环流为
\[
\int_0^1\int_0^1(1x^{0}+1y^{0})dxdy=2.
\]
所有偏导在平面上连续，正方形边界分段光滑，定理条件满足。结果与幂指数无关，是因为两个幂导数在 $[0,1]$ 上的积分都等于1。
\end{solve}

\begin{solve}[习题~\ref{ex:c44q8}：一个非零旋度的场]\phantomsection\label{ans:c44q8}
因第一分量与 $x$ 无关、第二分量与 $y$ 无关，散度为0。旋度为
\[
\nabla\times F=(0,0,2x^{1}+2y^{1}).
\]
由 Green 公式，环流为
\[
\int_0^1\int_0^1(2x^{1}+2y^{1})dxdy=2.
\]
所有偏导在平面上连续，正方形边界分段光滑，定理条件满足。结果与幂指数无关，是因为两个幂导数在 $[0,1]$ 上的积分都等于1。
\end{solve}

\begin{solve}[习题~\ref{ex:c44q9}：一个非零旋度的场]\phantomsection\label{ans:c44q9}
因第一分量与 $x$ 无关、第二分量与 $y$ 无关，散度为0。旋度为
\[
\nabla\times F=(0,0,3x^{2}+3y^{2}).
\]
由 Green 公式，环流为
\[
\int_0^1\int_0^1(3x^{2}+3y^{2})dxdy=2.
\]
所有偏导在平面上连续，正方形边界分段光滑，定理条件满足。结果与幂指数无关，是因为两个幂导数在 $[0,1]$ 上的积分都等于1。
\end{solve}

\begin{solve}[习题~\ref{ex:c44q10}：一个非零旋度的场]\phantomsection\label{ans:c44q10}
因第一分量与 $x$ 无关、第二分量与 $y$ 无关，散度为0。旋度为
\[
\nabla\times F=(0,0,4x^{3}+4y^{3}).
\]
由 Green 公式，环流为
\[
\int_0^1\int_0^1(4x^{3}+4y^{3})dxdy=2.
\]
所有偏导在平面上连续，正方形边界分段光滑，定理条件满足。结果与幂指数无关，是因为两个幂导数在 $[0,1]$ 上的积分都等于1。
\end{solve}

\begin{solve}[习题~\ref{ex:c44q11}：一个非零旋度的场]\phantomsection\label{ans:c44q11}
因第一分量与 $x$ 无关、第二分量与 $y$ 无关，散度为0。旋度为
\[
\nabla\times F=(0,0,5x^{4}+5y^{4}).
\]
由 Green 公式，环流为
\[
\int_0^1\int_0^1(5x^{4}+5y^{4})dxdy=2.
\]
所有偏导在平面上连续，正方形边界分段光滑，定理条件满足。结果与幂指数无关，是因为两个幂导数在 $[0,1]$ 上的积分都等于1。
\end{solve}

\begin{solve}[习题~\ref{ex:c44q12}：一个非零旋度的场]\phantomsection\label{ans:c44q12}
因第一分量与 $x$ 无关、第二分量与 $y$ 无关，散度为0。旋度为
\[
\nabla\times F=(0,0,6x^{5}+6y^{5}).
\]
由 Green 公式，环流为
\[
\int_0^1\int_0^1(6x^{5}+6y^{5})dxdy=2.
\]
所有偏导在平面上连续，正方形边界分段光滑，定理条件满足。结果与幂指数无关，是因为两个幂导数在 $[0,1]$ 上的积分都等于1。
\end{solve}

\begin{solve}[习题~\ref{ex:c44q13}：势函数给出的空间路径积分]\phantomsection\label{ans:c44q13}
按坐标求偏导，
\[
F=(2xy+2,x^2+1z^2,2yz).
\]
沿任意连接两点的分段光滑路径，链式法则与一元基本定理给出
\[
\int_CF\cdot dr=\varphi(1,2,1)-\varphi(0,0,0)=6.
\]
这一证明不依赖路径的具体形状，只要求势函数在路径附近可微。场的旋度为零可以通过混合偏导相等验证；但这里路径无关直接来自已给出的全局势函数。
\end{solve}

\begin{solve}[习题~\ref{ex:c44q14}：势函数给出的空间路径积分]\phantomsection\label{ans:c44q14}
按坐标求偏导，
\[
F=(2xy+1,x^2+2z^2,4yz).
\]
沿任意连接两点的分段光滑路径，链式法则与一元基本定理给出
\[
\int_CF\cdot dr=\varphi(1,2,1)-\varphi(0,0,0)=7.
\]
这一证明不依赖路径的具体形状，只要求势函数在路径附近可微。场的旋度为零可以通过混合偏导相等验证；但这里路径无关直接来自已给出的全局势函数。
\end{solve}

\begin{solve}[习题~\ref{ex:c44q15}：势函数给出的空间路径积分]\phantomsection\label{ans:c44q15}
按坐标求偏导，
\[
F=(2xy+3,x^2+2z^2,4yz).
\]
沿任意连接两点的分段光滑路径，链式法则与一元基本定理给出
\[
\int_CF\cdot dr=\varphi(1,2,1)-\varphi(0,0,0)=9.
\]
这一证明不依赖路径的具体形状，只要求势函数在路径附近可微。场的旋度为零可以通过混合偏导相等验证；但这里路径无关直接来自已给出的全局势函数。
\end{solve}

\begin{solve}[习题~\ref{ex:c44q16}：势函数给出的空间路径积分]\phantomsection\label{ans:c44q16}
按坐标求偏导，
\[
F=(2xy+2,x^2+3z^2,6yz).
\]
沿任意连接两点的分段光滑路径，链式法则与一元基本定理给出
\[
\int_CF\cdot dr=\varphi(1,2,1)-\varphi(0,0,0)=10.
\]
这一证明不依赖路径的具体形状，只要求势函数在路径附近可微。场的旋度为零可以通过混合偏导相等验证；但这里路径无关直接来自已给出的全局势函数。
\end{solve}

\begin{solve}[习题~\ref{ex:c44q17}：势函数给出的空间路径积分]\phantomsection\label{ans:c44q17}
按坐标求偏导，
\[
F=(2xy+1,x^2+4z^2,8yz).
\]
沿任意连接两点的分段光滑路径，链式法则与一元基本定理给出
\[
\int_CF\cdot dr=\varphi(1,2,1)-\varphi(0,0,0)=11.
\]
这一证明不依赖路径的具体形状，只要求势函数在路径附近可微。场的旋度为零可以通过混合偏导相等验证；但这里路径无关直接来自已给出的全局势函数。
\end{solve}

\begin{solve}[习题~\ref{ex:c44q18}：势函数给出的空间路径积分]\phantomsection\label{ans:c44q18}
按坐标求偏导，
\[
F=(2xy+4,x^2+3z^2,6yz).
\]
沿任意连接两点的分段光滑路径，链式法则与一元基本定理给出
\[
\int_CF\cdot dr=\varphi(1,2,1)-\varphi(0,0,0)=12.
\]
这一证明不依赖路径的具体形状，只要求势函数在路径附近可微。场的旋度为零可以通过混合偏导相等验证；但这里路径无关直接来自已给出的全局势函数。
\end{solve}

\begin{solve}[习题~\ref{ex:c44q19}：封闭曲面与缺失底面]\phantomsection\label{ans:c44q19}
给曲面补上 $z=0$ 的底圆盘，围成区域 $V:0\le z\le 1-r^2$、$0\le r\le\sqrt{1}$。闭合边界的外法向在抛物面上朝上，在底面上朝下。底面处 $F=0$，所以底面通量为0。由 $\nabla\cdot F=1$，
\[
\iint_SF\cdot n\,dS=|V|=2\pi\int_0^{\sqrt{1}}(1-r^2)rdr=\frac{1}{2}\pi.
\]
散度定理最初只能直接作用于闭合曲面；补面后必须计算并扣除新增面的通量。本题新增面恰好为零，并非所有题都会如此。
\end{solve}

\begin{solve}[习题~\ref{ex:c44q20}：封闭曲面与缺失底面]\phantomsection\label{ans:c44q20}
给曲面补上 $z=0$ 的底圆盘，围成区域 $V:0\le z\le 2-r^2$、$0\le r\le\sqrt{2}$。闭合边界的外法向在抛物面上朝上，在底面上朝下。底面处 $F=0$，所以底面通量为0。由 $\nabla\cdot F=1$，
\[
\iint_SF\cdot n\,dS=|V|=2\pi\int_0^{\sqrt{2}}(2-r^2)rdr=2\pi.
\]
散度定理最初只能直接作用于闭合曲面；补面后必须计算并扣除新增面的通量。本题新增面恰好为零，并非所有题都会如此。
\end{solve}

\begin{solve}[习题~\ref{ex:c44q21}：封闭曲面与缺失底面]\phantomsection\label{ans:c44q21}
给曲面补上 $z=0$ 的底圆盘，围成区域 $V:0\le z\le 3-r^2$、$0\le r\le\sqrt{3}$。闭合边界的外法向在抛物面上朝上，在底面上朝下。底面处 $F=0$，所以底面通量为0。由 $\nabla\cdot F=1$，
\[
\iint_SF\cdot n\,dS=|V|=2\pi\int_0^{\sqrt{3}}(3-r^2)rdr=\frac{9}{2}\pi.
\]
散度定理最初只能直接作用于闭合曲面；补面后必须计算并扣除新增面的通量。本题新增面恰好为零，并非所有题都会如此。
\end{solve}

\begin{solve}[习题~\ref{ex:c44q22}：封闭曲面与缺失底面]\phantomsection\label{ans:c44q22}
给曲面补上 $z=0$ 的底圆盘，围成区域 $V:0\le z\le 4-r^2$、$0\le r\le\sqrt{4}$。闭合边界的外法向在抛物面上朝上，在底面上朝下。底面处 $F=0$，所以底面通量为0。由 $\nabla\cdot F=1$，
\[
\iint_SF\cdot n\,dS=|V|=2\pi\int_0^{\sqrt{4}}(4-r^2)rdr=8\pi.
\]
散度定理最初只能直接作用于闭合曲面；补面后必须计算并扣除新增面的通量。本题新增面恰好为零，并非所有题都会如此。
\end{solve}

\begin{solve}[习题~\ref{ex:c44q23}：封闭曲面与缺失底面]\phantomsection\label{ans:c44q23}
给曲面补上 $z=0$ 的底圆盘，围成区域 $V:0\le z\le 5-r^2$、$0\le r\le\sqrt{5}$。闭合边界的外法向在抛物面上朝上，在底面上朝下。底面处 $F=0$，所以底面通量为0。由 $\nabla\cdot F=1$，
\[
\iint_SF\cdot n\,dS=|V|=2\pi\int_0^{\sqrt{5}}(5-r^2)rdr=\frac{25}{2}\pi.
\]
散度定理最初只能直接作用于闭合曲面；补面后必须计算并扣除新增面的通量。本题新增面恰好为零，并非所有题都会如此。
\end{solve}

\begin{solve}[习题~\ref{ex:c44q24}：封闭曲面与缺失底面]\phantomsection\label{ans:c44q24}
给曲面补上 $z=0$ 的底圆盘，围成区域 $V:0\le z\le 6-r^2$、$0\le r\le\sqrt{6}$。闭合边界的外法向在抛物面上朝上，在底面上朝下。底面处 $F=0$，所以底面通量为0。由 $\nabla\cdot F=1$，
\[
\iint_SF\cdot n\,dS=|V|=2\pi\int_0^{\sqrt{6}}(6-r^2)rdr=18\pi.
\]
散度定理最初只能直接作用于闭合曲面；补面后必须计算并扣除新增面的通量。本题新增面恰好为零，并非所有题都会如此。
\end{solve}

\clearpage
\section{开放题库原题选译与综合自测}\label{sec:ch45}
本章集中列出可按原章节追溯的 CLP 题目。中文题面保留其核心数学内容，答案以中文重述；少数题补全建模或光滑性条件。建议先独立完成，再用答案检查对概念和定理边界的理解。

\textbf{资源对应：}CLP-1 至 CLP-4，逐题列出源文件中的顺序号。本章共 25 道编号习题。
\subsection{习题}
\begin{example}[求导的线性性]\label{ex:c45q1}
若 $f,g$ 可导，判断 $(f+g)^{\prime}=f^{\prime}+g^{\prime}$ 是否成立。
\par\textit{选译并补充说明。来源：CLP-1 \S~2.4，题块顺序 1（question 与 Mquestion 合并计数）。}\par\hyperref[ans:c45q1]{解答见第~\pageref{ans:c45q1} 页。}
\end{example}
\begin{example}[乘积求导的错误公式]\label{ex:c45q2}
若 $f,g$ 可导，判断 $(fg)^{\prime}=f^{\prime}g^{\prime}$ 是否成立，并纠正。
\par\textit{选译并补充说明。来源：CLP-1 \S~2.4，题块顺序 2（question 与 Mquestion 合并计数）。}\par\hyperref[ans:c45q2]{解答见第~\pageref{ans:c45q2} 页。}
\end{example}
\begin{example}[最高速度给出时间下界]\label{ex:c45q3}
某动物迁移路程为5000千米，运动速度始终不超过每小时70千米。能断言其累计运动时间至少多少？
\par\textit{选译并补充说明。来源：CLP-1 \S~2.13，题块顺序 1（question 与 Mquestion 合并计数）。}\par\hyperref[ans:c45q3]{解答见第~\pageref{ans:c45q3} 页。}
\end{example}
\begin{example}[平均变化率对应某一时刻]\label{ex:c45q4}
设累计飞行路程 $s(t)$ 在 $[0,24]$ 连续、在 $(0,24)$ 可导，且 $s(24)-s(0)=240$ 千米。中值定理保证什么？
\par\textit{选译并补充说明。来源：CLP-1 \S~2.13，题块顺序 2（question 与 Mquestion 合并计数）。}\par\hyperref[ans:c45q4]{解答见第~\pageref{ans:c45q4} 页。}
\end{example}
\begin{example}[误差上界不必最紧]\label{ex:c45q5}
真实值 $f(10)=-3$，近似值 $F(10)=5$。判断绝对误差不超过7、8、9、100这四句话。
\par\textit{选译并补充说明。来源：CLP-1 \S~3.4.8，题块顺序 1（question 与 Mquestion 合并计数）。}\par\hyperref[ans:c45q5]{解答见第~\pageref{ans:c45q5} 页。}
\end{example}
\begin{example}[多项式的高阶 Taylor 近似]\label{ex:c45q6}
设 $f(x)=5x^3-24x^2+ex-\pi^4$，$T_5$ 为它在1处的五次 Taylor 多项式。求 $|f(37)-T_5(37)|$。
\par\textit{选译并补充说明。来源：CLP-1 \S~3.4.8，题块顺序 4（question 与 Mquestion 合并计数）。}\par\hyperref[ans:c45q6]{解答见第~\pageref{ans:c45q6} 页。}
\end{example}
\begin{example}[定积分区间的可加性]\label{ex:c45q7}
已知 $\int_0^b\cos x\,dx=\sin b$，求 $\int_a^b\cos x\,dx$。
\par\textit{选译并补充说明。来源：CLP-2 \S~1.2，题块顺序 2（question 与 Mquestion 合并计数）。}\par\hyperref[ans:c45q7]{解答见第~\pageref{ans:c45q7} 页。}
\end{example}
\begin{example}[反向积分中的符号]\label{ex:c45q8}
正函数 $f$ 的积分 $\int_5^0f(x)dx$ 采用100个等步长的有向黎曼和。步长、函数值和最终和的符号分别是什么？
\par\textit{选译并补充说明。来源：CLP-2 \S~1.2，题块顺序 4（question 与 Mquestion 合并计数）。}\par\hyperref[ans:c45q8]{解答见第~\pageref{ans:c45q8} 页。}
\end{example}
\begin{example}[连续被积函数何时产生反常积分]\label{ex:c45q9}
对 $\int_0^b dx/(1+x^2)$，哪些 $b$ 会使积分成为反常积分？允许 $b=\pm\infty$。
\par\textit{选译并补充说明。来源：CLP-2 \S~1.12，题块顺序 2（question 与 Mquestion 合并计数）。}\par\hyperref[ans:c45q9]{解答见第~\pageref{ans:c45q9} 页。}
\end{example}
\begin{example}[不能直接跨过对数奇点]\label{ex:c45q10}
计算 $\int_0^1x^4/(x^5-1)dx$，或者说明它发散。
\par\textit{选译并补充说明。来源：CLP-2 \S~1.12，题块顺序 6（question 与 Mquestion 合并计数）。}\par\hyperref[ans:c45q10]{解答见第~\pageref{ans:c45q10} 页。}
\end{example}
\begin{example}[非零通项极限]\label{ex:c45q11}
已知 $a_n\to1/2$，判断 $\sum_{n=7}^\infty a_n$ 是否收敛。
\par\textit{选译并补充说明。来源：CLP-2 \S~3.3，题块顺序 6（question 与 Mquestion 合并计数）。}\par\hyperref[ans:c45q11]{解答见第~\pageref{ans:c45q11} 页。}
\end{example}
\begin{example}[比值判别无结论的收敛例子]\label{ex:c45q12}
给出一个收敛级数，其相邻项比值极限为1。
\par\textit{选译并补充说明。来源：CLP-2 \S~3.3，题块顺序 11（question 与 Mquestion 合并计数）。}\par\hyperref[ans:c45q12]{解答见第~\pageref{ans:c45q12} 页。}
\end{example}
\begin{example}[幂级数一定收敛的点]\label{ex:c45q13}
设 $m\ge1$ 为整数，系数 $A_n$ 任意。对 $\sum_{n=m}^\infty A_n(x-c)^n$，仅凭这些信息能保证在哪些点收敛？
\par\textit{选译并补充说明。来源：CLP-2 \S~3.5，题块顺序 3（question 与 Mquestion 合并计数）。}\par\hyperref[ans:c45q13]{解答见第~\pageref{ans:c45q13} 页。}
\end{example}
\begin{example}[几何级数乘幂]\label{ex:c45q14}
给出 $x^3/(1-x)$ 的幂级数表达式及其收敛区间。
\par\textit{选译并补充说明。来源：CLP-2 \S~3.5，题块顺序 16（question 与 Mquestion 合并计数）。}\par\hyperref[ans:c45q14]{解答见第~\pageref{ans:c45q14} 页。}
\end{example}
\begin{example}[非零偏导与山顶]\label{ex:c45q15}
可微地形高度为 $z=f(x,y)$。若 $f_y(0,0)<0$，当前位置能否是局部最高点？
\par\textit{选译并补充说明。来源：CLP-3 \S~2.2，题块顺序 2（question 与 Mquestion 合并计数）。}\par\hyperref[ans:c45q15]{解答见第~\pageref{ans:c45q15} 页。}
\end{example}
\begin{example}[偏导数不能冒充全方向链式法则]\label{ex:c45q16}
设 $f(x,y)=x^2y/(x^2+y^2)$（非原点），$f(0,0)=0$。按定义求 $f_x(0,0)$、$f_y(0,0)$ 和 $d f(t,t)/dt$ 在0处的值。
\par\textit{选译并补充说明。来源：CLP-3 \S~2.2，题块顺序 3（question 与 Mquestion 合并计数）。}\par\hyperref[ans:c45q16]{解答见第~\pageref{ans:c45q16} 页。}
\end{example}
\begin{example}[圆锥面在正方形上的最值]\label{ex:c45q17}
求 $z=\sqrt{x^2+y^2}$ 在 $|x|,|y|\le1$ 上的最高和最低点，并解释偏导是否必须为零。
\par\textit{选译并补充说明。来源：CLP-3 \S~2.9，题块顺序 2（question 与 Mquestion 合并计数）。}\par\hyperref[ans:c45q17]{解答见第~\pageref{ans:c45q17} 页。}
\end{example}
\begin{example}[直线约束下的极值必要条件]\label{ex:c45q18}
设可微函数 $g(X)$ 限制在直线 $X=A+td$（$d\ne0$）上。若 $X_0$ 是约束局部极值点，给出梯度必须满足的条件。
\par\textit{选译并补充说明。来源：CLP-3 \S~2.9，题块顺序 3（question 与 Mquestion 合并计数）。}\par\hyperref[ans:c45q18]{解答见第~\pageref{ans:c45q18} 页。}
\end{example}
\begin{example}[球面上的线性函数]\label{ex:c45q19}
求 $f(x,y,z)=x+y-z$ 在单位球面上的最大值和最小值。
\par\textit{选译并补充说明。来源：CLP-3 \S~2.10，题块顺序 3（question 与 Mquestion 合并计数）。}\par\hyperref[ans:c45q19]{解答见第~\pageref{ans:c45q19} 页。}
\end{example}
\begin{example}[保守场加常向量]\label{ex:c45q20}
若 $F$ 为保守场，常向量为 $(a,b,c)$，判断 $F+(a,b,c)$ 是否仍保守并给出势函数。
\par\textit{选译并补充说明。来源：CLP-4 \S~2.3，题块顺序 3（question 与 Mquestion 合并计数）。}\par\hyperref[ans:c45q20]{解答见第~\pageref{ans:c45q20} 页。}
\end{example}
\begin{example}[限制定义域后恢复势函数]\label{ex:c45q21}
在 $D=\{(x,y):x>1\}$ 上，判断 $F=(-y,x)/(x^2+y^2)$ 是否保守。
\par\textit{选译并补充说明。来源：CLP-4 \S~2.3，题块顺序 5（question 与 Mquestion 合并计数）。}\par\hyperref[ans:c45q21]{解答见第~\pageref{ans:c45q21} 页。}
\end{example}
\begin{example}[旋度未必垂直原场]\label{ex:c45q22}
判断 $\nabla\times F$ 是否一定与 $F$ 垂直。
\par\textit{选译并补充说明。来源：CLP-4 \S~4.1，题块顺序 2（question 与 Mquestion 合并计数）。}\par\hyperref[ans:c45q22]{解答见第~\pageref{ans:c45q22} 页。}
\end{example}
\begin{example}[由对称性得到零通量]\label{ex:c45q23}
设 $F=(z+\sin y,zy,\sin x\cos y)$。求其散度和通过半径3、中心在原点的球面的向外通量。
\par\textit{选译并补充说明。来源：CLP-4 \S~4.2，题块顺序 7（question 与 Mquestion 合并计数）。}\par\hyperref[ans:c45q23]{解答见第~\pageref{ans:c45q23} 页。}
\end{example}
\begin{example}[平移球的线性散度积分]\label{ex:c45q24}
球体积为 $V$、球心为 $(x_0,y_0,z_0)$。求 $F=(x^2,xy,3z-yz)$ 通过边界球面的向外通量。
\par\textit{选译并补充说明。来源：CLP-4 \S~4.2，题块顺序 9（question 与 Mquestion 合并计数）。}\par\hyperref[ans:c45q24]{解答见第~\pageref{ans:c45q24} 页。}
\end{example}
\begin{example}[选择使环流最大的曲线]\label{ex:c45q25}
在所有光滑简单闭曲线的逆时针定向中，找出使 $\oint_C(y^3-y)dx-2x^3dy$ 最大的曲线，并求最大值。
\par\textit{选译并补充说明。来源：CLP-4 \S~4.3，题块顺序 25（question 与 Mquestion 合并计数）。}\par\hyperref[ans:c45q25]{解答见第~\pageref{ans:c45q25} 页。}
\end{example}
\subsection{解答}
\begin{solve}[习题~\ref{ex:c45q1}：求导的线性性]\phantomsection\label{ans:c45q1}
成立。由差商的和等于差商之和，再对两个存在的有限极限使用加法法则，就得到和的求导公式。可导条件确保两个导数都存在；若只有函数值存在，不能直接使用这一公式。
\end{solve}

\begin{solve}[习题~\ref{ex:c45q2}：乘积求导的错误公式]\phantomsection\label{ans:c45q2}
一般不成立。正确公式为 $(fg)^{\prime}=f^{\prime}g+fg^{\prime}$。取 $f=g=x$，错误公式给出1，但实际导数为 $2x$。乘积的变化包含两个因子分别变化的贡献，不是两个导数直接相乘。
\end{solve}

\begin{solve}[习题~\ref{ex:c45q3}：最高速度给出时间下界]\phantomsection\label{ans:c45q3}
路程不超过最高速度乘累计运动时间，所以 $T\ge5000/70=500/7$ 小时，约71.43小时。若动物大部分时间低于最高速度，实际时间会更长；这些信息不能确定精确时间。这里把速度上界作为题设模型，不视为对现实动物的已核实生物学数据。
\end{solve}

\begin{solve}[习题~\ref{ex:c45q4}：平均变化率对应某一时刻]\phantomsection\label{ans:c45q4}
存在 $c\in(0,24)$ 使 $s^{\prime}(c)=240/24=10$ 千米每小时。这里使用累计路程而非未经说明的空间位置坐标，且明确假设可导；在这一模型下导数才表示速率。定理只保证至少一个这样的时刻，不确定其具体位置，也不说明全程匀速。
\end{solve}

\begin{solve}[习题~\ref{ex:c45q5}：误差上界不必最紧]\phantomsection\label{ans:c45q5}
绝对误差为 $|-3-5|=8$，所以不超过7为假，其余三句为真。8是这组给定界中最紧的界，9和100虽然正确，但提供的信息较弱。误差上界只需是一个有效的保证，不必恰好等于实际误差。
\end{solve}

\begin{solve}[习题~\ref{ex:c45q6}：多项式的高阶 Taylor 近似]\phantomsection\label{ans:c45q6}
原函数次数只有3，四阶及更高阶导数恒为0。因此五次 Taylor 多项式与原函数完全相同，所谓四次、五次系数都为0。误差精确等于0，而不是仅仅“很小”。距离展开中心较远在本题中不造成误差，因为展开已经是一个恒等式。
\end{solve}

\begin{solve}[习题~\ref{ex:c45q7}：定积分区间的可加性]\phantomsection\label{ans:c45q7}
由区间可加性与反向积分定义，\[\int_a^b\cos x\,dx=\int_0^b\cos x\,dx-\int_0^a\cos x\,dx=\sin b-\sin a.\]该公式不要求 $a<b$；反向区间通过符号约定自动包含。
\end{solve}

\begin{solve}[习题~\ref{ex:c45q8}：反向积分中的符号]\phantomsection\label{ans:c45q8}
有向步长为 $(0-5)/100=-1/20$，每个函数值为正，因此每项的乘积为负，和也为负。普通几何小区间长度仍非负；这里负号表达从5到0的积分方向。相反，$\int_0^5f$ 为正。
\end{solve}

\begin{solve}[习题~\ref{ex:c45q9}：连续被积函数何时产生反常积分]\phantomsection\label{ans:c45q9}
被积函数在整个实轴连续，因此每个有限 $b$ 都给出正常定积分。只有 $b=+\infty$ 或 $b=-\infty$ 时，因为区间无界而成为反常积分。成为反常积分并不意味着发散；本题两个无穷端点的积分都收敛。
\end{solve}

\begin{solve}[习题~\ref{ex:c45q10}：不能直接跨过对数奇点]\phantomsection\label{ans:c45q10}
在上端点分母为0，必须截断到 $t<1$。令 $u=x^5$，有 $\int_0^t x^4/(x^5-1)dx=\frac15\log(1-t^5)$。当 $t\to1^-$ 时趋于负无穷，所以反常积分发散。原函数形式存在并不保证端点极限有限。
\end{solve}

\begin{solve}[习题~\ref{ex:c45q11}：非零通项极限]\phantomsection\label{ans:c45q11}
发散，因为级数收敛的必要条件是通项趋于0。把起点从1改为7只删除有限项，不影响尾部通项极限，也不影响敛散性。无需知道每个 $a_n$ 的具体表达式。
\end{solve}

\begin{solve}[习题~\ref{ex:c45q12}：比值判别无结论的收敛例子]\phantomsection\label{ans:c45q12}
取 $\sum_{n=1}^\infty1/n^2$，它为收敛的 $p$ 级数，但 $a_{n+1}/a_n=n^2/(n+1)^2\to1$。因此“比值极限等于1”不能推出发散，也不能推出收敛；需要另用适合的判别。
\end{solve}

\begin{solve}[习题~\ref{ex:c45q13}：幂级数一定收敛的点]\phantomsection\label{ans:c45q13}
可以保证在 $x=c$ 收敛，因为每一项都为0。除此以外不能保证：取 $A_n=n!$，在任何 $x\ne c$ 处相邻绝对值比为 $(n+1)|x-c|\to\infty$，所以发散。若起点为0，中心点的值就是常数项 $A_0$；不需要把通常的数值表达式 $0^0$ 作为推导依据。
\end{solve}

\begin{solve}[习题~\ref{ex:c45q14}：几何级数乘幂]\phantomsection\label{ans:c45q14}
在 $|x|<1$ 上，$1/(1-x)=\sum_{n=0}^\infty x^n$，乘以 $x^3$ 得 $x^3/(1-x)=\sum_{n=3}^\infty x^n$。在 $x=1$ 与 $x=-1$ 处通项均不趋零，所以收敛区间为 $(-1,1)$。起点从0变为3反映的是乘以 $x^3$ 后幂次整体平移。
\end{solve}

\begin{solve}[习题~\ref{ex:c45q15}：非零偏导与山顶]\phantomsection\label{ans:c45q15}
不能。沿负 $y$ 方向取足够小的负增量 $h$，由偏导定义 $[f(0,h)-f(0,0)]/h<0$，乘以负数 $h$ 后得到高度差为正，所以附近存在更高点。正 $y$ 方向则局部下降。论证只需目标点偏导为负，不必额外假设偏导在邻域内处处为负。
\end{solve}

\begin{solve}[习题~\ref{ex:c45q16}：偏导数不能冒充全方向链式法则]\phantomsection\label{ans:c45q16}
两条坐标轴上的函数值恒为0，所以两偏导都为0。沿对角线，非零 $t$ 时 $f(t,t)=t/2$，在0也等于0，故复合函数导数为1/2。若误把两偏导代入可微函数的链式法则会得到0，矛盾说明函数在原点不可微。偏导存在不够保证链式法则所需的全可微条件。
\end{solve}

\begin{solve}[习题~\ref{ex:c45q17}：圆锥面在正方形上的最值]\phantomsection\label{ans:c45q17}
高度等于平面点到原点的距离，所以最小值0在原点取得，最大值 $\sqrt2$ 在四个角点取得。原点处沿坐标轴的截面为绝对值函数，两偏导不存在。角点处函数可微但梯度不为零，因为它们是约束区域的边界点。可导内点极值的必要条件不能套到这两类例外上。
\end{solve}

\begin{solve}[习题~\ref{ex:c45q18}：直线约束下的极值必要条件]\phantomsection\label{ans:c45q18}
令 $h(t)=g(A+td)$，选择 $t_0$ 使 $X_0=A+t_0d$。链式法则给出 $h^{\prime}(t_0)=\nabla g(X_0)\cdot d$。一元 Fermat 条件要求它为0，所以梯度与直线方向正交，同时 $X_0$ 必须确实在直线上。这是必要条件，不能单凭它断言极值存在。
\end{solve}

\begin{solve}[习题~\ref{ex:c45q19}：球面上的线性函数]\phantomsection\label{ans:c45q19}
把函数写成 $(1,1,-1)\cdot(x,y,z)$。Cauchy--Schwarz 不等式给出绝对值不超过 $\sqrt3$。最大值 $\sqrt3$ 在 $(1,1,-1)/\sqrt3$ 处取得，最小值负 $\sqrt3$ 在相反点取得。球面紧，线性函数连续，所以最值存在；等号条件给出全部取值点。
\end{solve}

\begin{solve}[习题~\ref{ex:c45q20}：保守场加常向量]\phantomsection\label{ans:c45q20}
设 $F=\nabla\varphi$，令 $\psi=\varphi+ax+by+cz$，则 $\nabla\psi=F+(a,b,c)$。所以仍为保守场。给出的势函数在原势函数的同一域内成立，不需要重新证明定义域的单连通性，因为势函数已经实际构造出来。
\end{solve}

\begin{solve}[习题~\ref{ex:c45q21}：限制定义域后恢复势函数]\phantomsection\label{ans:c45q21}
取 $\varphi(x,y)=\arctan(y/x)$。因为 $x>1$，该函数单值且光滑，求偏导得 $\varphi_x=-y/(x^2+y^2)$、$\varphi_y=x/(x^2+y^2)$，所以 $F=\nabla\varphi$。同样的向量场在整个去心平面上没有单值势函数，但限制到这里后障碍消失；定义域是问题的实质部分。
\end{solve}

\begin{solve}[习题~\ref{ex:c45q22}：旋度未必垂直原场]\phantomsection\label{ans:c45q22}
不一定。取 $F=(1,0,y)$，则 $\nabla\times F=(1,0,0)$，两者点积恒为1，并不垂直。叉积符号中的 $\nabla$ 是微分算子，不能把它当成一个普通固定向量，机械套用两个普通向量叉积垂直每个因子的几何性质。
\end{solve}

\begin{solve}[习题~\ref{ex:c45q23}：由对称性得到零通量]\phantomsection\label{ans:c45q23}
散度为 $0+z+0=z$。向量场在整个球内光滑，由散度定理通量等于球内的 $\iiint z\,dV$。球关于 $z=0$ 对称，而 $z$ 关于反射变号，故积分为0。通量为零并不要求每个点的法向分量为零，而可以来自不同位置的抵消。
\end{solve}

\begin{solve}[习题~\ref{ex:c45q24}：平移球的线性散度积分]\phantomsection\label{ans:c45q24}
散度为 $2x+x+3-y=3+3x-y$。球体的形心为球心，因此 $\iiint_Bx\,dV=x_0V$、$\iiint_By\,dV=y_0V$。散度定理给出通量 $(3+3x_0-y_0)V$。不需要知道球半径，因为题目已经给出体积，且散度只是坐标的一次函数。
\end{solve}

\begin{solve}[习题~\ref{ex:c45q25}：选择使环流最大的曲线]\phantomsection\label{ans:c45q25}
Green 公式把环流化为内部区域上 $\iint(1-6x^2-3y^2)dxdy$。为使积分最大，应包含被积函数为正的椭圆内部，并排除为负的外部，所以最优边界为 $6x^2+3y^2=1$。令 $u=\sqrt6x,v=\sqrt3y$，Jacobian 为 $1/\sqrt{18}$，所以最大值为\[\frac1{\sqrt{18}}\int_0^{2\pi}\int_0^1(1-r^2)r\,dr\,d\theta=\frac{\pi}{2\sqrt{18}}.\]与任何其他区域比较，其遗漏的正部分和加入的负部分都只会减小积分。
\end{solve}

\clearpage
\section*{来源、许可与改编索引}
\addcontentsline{toc}{section}{来源、许可与改编索引}
\label{sec:sources}
CLP 系列原作者为 Joel Feldman、Andrew Rechnitzer、Elyse Yeager。原书把题目、提示、简答和详解分开编排，适用于本科微积分课程。这里的中文稿改变了选题顺序和解答组织，并加入参数练习与证明补充。章节中的“资源对应”表示题型与知识主题的对应，不表示该章每一道参数变式都原封不动出现于原书。具体选译题在题后给出章节、原标签或原题描述，便于追溯。
\begin{description}
\item[CLP-1：一元微分] \href{https://personal.math.ubc.ca/~CLP/CLP1/}{Differential Calculus 官方主页}；\href{https://personal.math.ubc.ca/~CLP/CLP1/clp_1_dc_problems.pdf}{题目与解答 PDF}；\href{https://github.com/arechnitzer/CLP1}{LaTeX 源码}。本书的函数极限、连续、求导、中值定理、泰勒近似与导数应用参照对应章节。
\item[CLP-2：积分与级数] \href{https://personal.math.ubc.ca/~CLP/CLP2/}{Integral Calculus 官方主页}；\href{https://github.com/arechnitzer/CLP2}{LaTeX 源码}。积分方法、反常积分、积分应用、可分离方程、数列和级数参照对应章节。
\item[CLP-3：多元微积分] \href{https://personal.math.ubc.ca/~CLP/CLP3/}{Multivariable Calculus 官方主页}；\href{https://personal.math.ubc.ca/~CLP/CLP3/clp_3_mc_problems.pdf}{题目与解答 PDF}；\href{https://github.com/arechnitzer/CLP3}{LaTeX 源码}。空间几何、多元极限、偏导数、极值、二重与三重积分参照对应章节。
\item[CLP-4：向量微积分] \href{https://personal.math.ubc.ca/~CLP/CLP4/}{Vector Calculus 官方主页}；\href{https://personal.math.ubc.ca/~CLP/CLP4/clp_4_vc_problems.pdf}{题目与解答 PDF}；\href{https://github.com/arechnitzer/CLP4}{LaTeX 源码}。曲率、曲线积分、保守场、曲面积分及积分定理参照对应章节。
\item[许可核对] \href{https://personal.math.ubc.ca/~CLP/src/}{CLP 官方源码与许可汇总页}；\href{https://creativecommons.org/licenses/by-nc-sa/4.0/}{CC BY-NC-SA 4.0 许可}。原资料版权归原作者所有。本书保留署名，明确标记改编，并按同一许可提供。
\end{description}
\subsection*{检索与选择记录}
检索同时发现 OpenStax Calculus 的开放教材、MIT OpenCourseWare 的 Analysis I 作业及解答页面，以及若干中文教材习题解答仓库。考虑到中文仓库的完整程度、上游题目权利与解答可核查程度差异，本书未直接拼接这些中文商业教材的习题解答全文，主要选取 CLP 作为可追溯的开放资源基础。进一步学习可访问\href{https://openstax.org/details/books/calculus-volume-1}{OpenStax Calculus Volume 1}和\href{https://ocw.mit.edu/courses/18-100b-analysis-i-fall-2010/download/}{MIT 18.100B Analysis I 课程资源页}；这些链接是延伸阅读，不表示本书题目都来自这些资源。
\end{document}
